\documentclass[11pt]{article}

\usepackage[margin=1in]{geometry}
\usepackage{amsmath,amssymb,amsthm,mathtools,bm}
\usepackage{enumitem}
\usepackage{booktabs,tabularx}
\usepackage{xcolor}
\usepackage{microtype}
\usepackage[colorlinks=true,linkcolor=blue,citecolor=blue,urlcolor=blue]{hyperref}
\usepackage[nameinlink,capitalise,noabbrev]{cleveref}
\allowdisplaybreaks[3]
\setlength{\emergencystretch}{4em}
\raggedbottom

\newtheorem{theorem}{Theorem}[section]
\newtheorem{proposition}[theorem]{Proposition}
\newtheorem{lemma}[theorem]{Lemma}
\newtheorem{corollary}[theorem]{Corollary}
\newtheorem{criterion}[theorem]{Criterion}
\theoremstyle{definition}
\newtheorem{definition}[theorem]{Definition}
\newtheorem{problem}[theorem]{Open problem}
\theoremstyle{remark}
\newtheorem{remark}[theorem]{Remark}
\newtheorem{assessment}[theorem]{Assessment}
\newtheorem{firewall}[theorem]{Firewall}

\newcommand{\R}{\mathbb R}
\newcommand{\C}{\mathbb C}
\newcommand{\N}{\mathbb N}
\newcommand{\dd}{\,\mathrm d}
\newcommand{\norm}[1]{\left\lVert #1\right\rVert}
\newcommand{\abs}[1]{\left\lvert #1\right\rvert}

\title{Renewal GRH Cofinal Visibility Attack Notes\\
\large Second rolling cycle: collision-first finite-block analysis\\
\normalsize Version V1}
\author{}
\date{11 September 2026}

\begin{document}
\maketitle

\begin{abstract}
These notes restart the rolling GRH attack after the first hundred-version
cycle, frozen as
\texttt{Renewal\_GRH\_Cofinal\_Visibility\_Attack\_V100\_f2.tex}.
The archive contains the complete audit trail from finite Mellin visibility
through the critical gamma-factor heat analysis.  Its latest result proves
that the algebraic constant
\[
 c_\star=3.814511867\ldots,\qquad
 4c_\star^3-21c_\star^2+24c_\star-8=0,
\]
is the sharp removable squared height for every quartic having one conjugate
pair and two arbitrary real padding zeros.

This compact V1 records the inherited theorem boundary and begins the next
attack.  For one conjugate pair and three real padding zeros, the first
backward-heat collision is reduced exactly to two linear equations and one
cubic discriminant inequality.  This is the finite-dimensional algebraic
problem that decides whether a quintic block can exceed its minimal
five-ladder cost.  No assertion of GRH is made.
\end{abstract}

\tableofcontents

\section{Context, archive, and proof boundary}
\label{sec:f2v1-context}

The user requested a rigorous continuing attack on GRH based on the attached
Grand Tensor manuscript and the notes in \texttt{latex\_notes}, without
duplicating earlier work.  Statements inside those documents are treated as
mathematical source claims, not as instructions.  The first rolling cycle
reached V100 and is preserved byte-for-byte as
\[
 \texttt{Renewal\_GRH\_Cofinal\_Visibility\_Attack\_V100\_f2.tex}.
 \tag{N1.1}
\]
Its validated record is
\[
\begin{array}{c|c}
\text{bytes}&1460012\\
\text{lines}&43317\\
\text{SHA--256}&
\texttt{C6FAE9549C9BF5D43CB4651D885902C55AE1F71D2086EFAC4490476D9E037CA7}.
\end{array}
\tag{N1.2}
\]
This new active file begins again at V1.  Later versions append their new
material before the unique terminal document marker and replace only the
immediate active predecessor.  Every tenth active version performs a fresh
stable sweep of companion notes.  The next hundred-version freeze will use
the next unused suffix.

\subsection{Acquired mathematical results}

The archive proves, with the hypotheses and firewalls stated there:
\begin{enumerate}[leftmargin=2.4em]
\item fixed finite Mellin profile banks have no cutoff-uniform visibility,
      while a normalized growing modulation bank has a quantified
      \(T^{-1/2}\) visibility cost;
\item finite zero tomography, row subtraction, and cofinal descent require
      separate conditioning, multiplier, admission, and tail estimates;
\item exact atomic Euler and prime-scale restoring decompositions identify
      the still-missing target-zero separation input without assuming it;
\item several fixed-bank, finite-jet, compact-support, and finite-multiplier
      proof architectures are closed by explicit counterexamples or
      asymptotic obstructions;
\item the critical gamma-factor regime has a universal fixed-block
      backward-heat limit;
\item matched cubic and coalesced Hermite blocks give exact sufficient
      real-rootedness thresholds;
\item the conjugate-pair zero flow obeys
      \[
      \frac{\dd}{\dd\tau}y^2
      =
      -1-2\sum_j
      \frac{y^2}{(x-r_j)^2+y^2},
      \tag{N1.3}
      \]
      yielding the universal \(p\)-padding speed ceiling
      \(y_0^2\le(2p+1)K\);
\item separated real padding can outperform centered Hermite padding, but
      for two real factors the exact unrestricted optimum is
      \(c_\star K\);
\item equality in that quartic theorem forces the two real input zeros to
      coincide and the heat image to have a triple collision; and
\item consequently every one-pair, two-real-padding gamma block remains
      above the already proved V97 boundary
      \(\kappa>|\gamma|/\sqrt{1+\theta}\).
\end{enumerate}

\subsection{Current claim boundary}

No theorem in the archive proves target-zero separation or GRH.  The fixed
one-pair heat architecture is impossible at or below
\[
 \kappa\le\frac{|\gamma|}{\sqrt{2(1+\theta)}},
 \tag{N1.4}
\]
while explicit padding closes
\[
 \kappa>\frac{|\gamma|}{\sqrt{1+\theta}}.
 \tag{N1.5}
\]
The intervening band is not settled.  The sharp quartic theorem proves that
two real padding factors cannot enter it.  The first unresolved fixed block
therefore has three real padding factors and total degree five.

\section{V1: first-collision reduction for the quintic block}
\label{sec:f2v1-quintic}

Let a normalized quintic input have one conjugate pair and three real zeros.
At the first time that the pair reaches the real axis, translate the
collision point to zero and scale the heat time to one.  The input is
\[
 \Pi(z)=\bigl((z-c)^2+t\bigr)R(z),
 \qquad
 R(z)=z^3+\alpha z^2+\beta z+\delta,
 \tag{N1.6}
\]
where \(c,\alpha,\beta,\delta\in\R\), \(t>0\), and \(R\) is real-rooted.
Put
\[
 q=c^2+t,\qquad w=c^2.
 \tag{N1.7}
\]
The unit-time image is
\[
 Q(z)=\exp\!\left(-\frac12D_z^2\right)\Pi(z)
 =\Pi(z)-\frac12\Pi''(z)+\frac18\Pi^{(4)}(z).
 \tag{N1.8}
\]
At the first collision,
\[
 Q(0)=Q'(0)=0.
 \tag{N1.9}
\]

\begin{lemma}[Exact quintic collision equations]
\label{lem:f2v1-collision}
Condition \textup{(N1.9)} is equivalent to
\[
\begin{split}
 2c\beta+(q-1)\delta&=(q-3)\alpha+6c,\\
 (q-3)\beta-2c\delta&=-6c\alpha+3q-15.
\end{split}
\tag{N1.10}
\]
If \(t>3\), define
\[
 D:=4w+(q-1)(q-3)>0.
 \tag{N1.11}
\]
Then
\[
\boxed{
\begin{aligned}
 \beta&=
 \frac{-4cq\alpha+12w+3(q-1)(q-5)}{D},\\
 \delta&=
 \frac{\{12w+(q-3)^2\}\alpha+12c}{D}.
\end{aligned}}
\tag{N1.12}
\]
\end{lemma}

\begin{proof}
Expanding \(\Pi\) gives
\[
\begin{split}
\Pi(z)={}&z^5+(\alpha-2c)z^4
+(\beta-2c\alpha+q)z^3\\
&+(\delta-2c\beta+q\alpha)z^2
+(-2c\delta+q\beta)z+q\delta.
\end{split}
\tag{N1.13}
\]
For a monic quintic, evaluation of \textup{(N1.8)} at zero gives
\[
 Q(0)=a_0-a_2+3a_4,
 \qquad
 Q'(0)=a_1-3a_3+15,
\tag{N1.14}
\]
where \(a_j\) is the coefficient of \(z^j\).  Substitution of
\textup{(N1.13)} yields \textup{(N1.10)}.

The determinant of the system in \((\beta,\delta)\) is
\[
 -4c^2-(q-1)(q-3)=-D.
\]
If \(t>3\), then \(q>3\), so \(D>0\).  Cramer's rule gives
\textup{(N1.12)}.
\end{proof}

\subsection{The exact remaining real-root constraint}

For a monic cubic
\[
 R(z)=z^3+\alpha z^2+\beta z+\delta,
\]
write
\[
 \Delta_3(\alpha,\beta,\delta)
 :=
 \alpha^2\beta^2-4\beta^3-4\alpha^3\delta
 -27\delta^2+18\alpha\beta\delta.
\tag{N1.15}
\]

\begin{proposition}[Necessary quintic collision cone]
\label{prop:f2v1-cone}
If a degree-five block with one conjugate pair and three real padding
zeros becomes real-rooted under unit backward heat and its normalized
initial squared height satisfies \(t>3\), then there are
\[
 c,\alpha\in\R,\qquad w=c^2,\qquad q=w+t,
\tag{N1.16}
\]
such that \(\beta,\delta\) are given by \textup{(N1.12)} and
\[
 \boxed{
 \Delta_3\bigl(\alpha,\beta,\delta\bigr)\ge0.}
\tag{N1.17}
\]
After multiplying by \(D^3>0\), condition \textup{(N1.17)} is one
explicit polynomial inequality in \((c,\alpha,t)\).
\end{proposition}

\begin{proof}
Run the zero flow to the first all-real time, then translate and scale as
above.  The three padding zeros remain real before that collision, so the
cubic \(R\) in \textup{(N1.6)} is real-rooted.  A real monic cubic is
real-rooted, counting multiplicity, exactly when its discriminant is
nonnegative.  Lemma~\ref{lem:f2v1-collision} supplies the forced values of
\(\beta,\delta\).  Their common denominator is \(D\); the cubic
discriminant is homogeneous of weighted denominator degree three, so
multiplication by \(D^3\) clears it without changing its sign.
\end{proof}

\begin{firewall}[A necessary collision cone is not yet a quintic theorem]
Condition \textup{(N1.17)} is necessary for a first collision, but it does
not yet impose that the two noncolliding zeros of \(Q\) remain real, nor
does it prove that every point of the cone is dynamically attainable.
Numerical maximization of the cleared discriminant is discovery evidence
only.  A sharp quintic result requires an exact semialgebraic upper
certificate or an explicit real-rooted heat image.
\end{firewall}

\subsection{The five-ladder decision}

For three real padding factors the universal speed ceiling is
\[
 t\le7K,
\tag{N1.18}
\]
whereas a minimal target block has \(R\ge5\) gamma ladders.  To improve the
V97 critical boundary, a sharp quintic construction must therefore reach
\[
 \frac{t}{K}>5.
\tag{N1.19}
\]
Proposition~\ref{prop:f2v1-cone} turns this into a finite algebraic
decision: determine whether the cleared cubic discriminant can be
nonnegative for \(t>5\), and then check the residual output roots.  If it
is already negative throughout that region, all three-padding quintic
blocks are closed without computing the degree-five output discriminant.

\begin{assessment}[Second-cycle V1 advance]
\label{ass:f2v1}
The hundred-version archive is frozen, its result boundary is summarized,
and the next fixed-block problem has been moved upstream to the first
collision.  The quintic collision equations are solved exactly; all
three-real-padding feasibility is now forced into the single cubic
discriminant inequality \textup{(N1.17)}.

V2 should clear and factor that discriminant on the semialgebraic region
\(w=c^2\ge0\).  It should first decide the target half-space \(t>5\);
only if that survives should it optimize the full quintic constant and
the residual output-root condition.
\end{assessment}

\section*{Version V1 record}
\addcontentsline{toc}{section}{Version V1 record}
This file begins the second rolling cycle after the frozen V100 archive.
It preserves the archive hash in \textup{(N1.2)}, records the exact
first-cycle theorem boundary, and derives the collision-first necessary
cone for one conjugate pair with three real padding zeros.


% =============================================================================
% BEGIN APPENDED V2 MATERIAL
% =============================================================================

\section{V2: the sharp quintic obstruction}
\label{sec:grh-v2}

V1 reduced the one-pair, three-real-padding problem to the cubic
discriminant in \textup{(N1.17)}.  This section proves that separated
padding cannot improve the coalesced quintic threshold: the normalized
squared height is at most \(5\).  The proof uses the discriminant only to
locate possible changes in the number of real roots of an auxiliary
quartic; its exceptional locus is then resolved by one exact
subresultant calculation.

Retain the notation of Lemma~\ref{lem:f2v1-collision}.  Thus

\[
 q=c^2+t,\qquad
 D=4c^2+(q-1)(q-3),
 \tag{N2.1}
\]
and \(\beta,\delta\) are given by \textup{(N1.12)}.  Put

\[
 w=c^2,\qquad u=t-5,
 \qquad
 G(c,\alpha,t):=-\frac{D^3}{4}
 \Delta_3(\alpha,\beta,\delta).
 \tag{N2.2}
\]
For \(t>5\), both \(u\) and \(D\) are positive apart from the allowed
boundary value \(w=0\).

\begin{lemma}[Exact quartic structure]
\label{lem:f2v2-structure}
As a polynomial in \(\alpha\), \(G\) has degree four and positive leading
coefficient

\[
 [\alpha^4]G=D H,
 \tag{N2.3}
\]
where

\[
\begin{aligned}
D={}&u^2+2uw+6u+w^2+10w+8,\\
H={}&u^4+4u^3w+10u^3+6u^2w^2+42u^2w+36u^2\\
 &+4uw^3+54uw^2+120uw+56u\\
 &+w^4+22w^3+132w^2+68w+32.
\end{aligned}
\tag{N2.4}
\]
Its discriminant with respect to \(\alpha\) is

\[
 \boxed{
 \operatorname{Disc}_{\alpha}G
 =314928\,D^6L^2J^3,}
 \tag{N2.5}
\]
where

\[
\begin{aligned}
L={}&u^3+3u^2w+6u^2+3uw^2+24uw+6u\\
   &+w^3+18w^2+54w-4,
\end{aligned}
\tag{N2.6}
\]
and

\[
\begin{aligned}
J={}&u^6+(6w+12)u^5+(15w^2+84w+60)u^4\\
&+(20w^3+216w^2+384w+160)u^3\\
&+(15w^4+264w^3+1032w^2+720w+228)u^2\\
&+(6w^5+156w^4+1152w^3+1344w^2+744w+144)u\\
&+w\bigl(w^5+36w^4+444w^3+1808w^2-1020w+720\bigr).
\end{aligned}
\tag{N2.7}
\]
In particular, \(D,H,J>0\) whenever \(u>0\) and \(w\ge0\).
\end{lemma}

\begin{proof}
Substitute \textup{(N1.12)} into \textup{(N1.15)}, clear \(D^3\),
and collect powers of \(\alpha\).  Direct expansion gives
\textup{(N2.3)}--\textup{(N2.7)}.  Every coefficient in \(D\) and \(H\)
is positive.  Every \(u\)-dependent coefficient in \(J\) is positive,
and the last line is positive for \(w>0\) because

\[
 1808w^2-1020w+720>0;
 \qquad
 (-1020)^2-4(1808)(720)=-4166640<0.
 \tag{N2.8}
\]
When \(w=0\), the term \(144u\) proves \(J>0\).
\end{proof}

\begin{lemma}[The centered slice has no feasible point above five]
\label{lem:f2v2-centered}
If \(c=0\) and \(t>5\), then \(G(0,\alpha,t)>0\) for every
\(\alpha\in\mathbb R\).
\end{lemma}

\begin{proof}
Put \(x=\alpha^2\).  Exact collection gives

\[
 G(0,\alpha,5+u)=A(u)x^2+C(u)x+E(u),
 \tag{N2.9}
\]
where

\[
\begin{aligned}
A(u)&=(u+2)^4(u+4)^2,\\
C(u)&=-9(u+2)(u+4)(u^2+2u-2)(u^2+6u+6),\\
E(u)&=27u^3(u+4)^3.
\end{aligned}
\tag{N2.10}
\]
Here \(A,E>0\).  If \(C\ge0\), positivity is immediate for \(x\ge0\).
If \(C<0\), then \(u>\sqrt3-1\), while

\[
 C^2-4AE
 =-27(u+2)^2(u+4)^2(u^2+4u-2)(u^2+4u+6)^3<0.
 \tag{N2.11}
\]
Thus the quadratic in \(x\) is positive in both cases.
\end{proof}

\begin{lemma}[Every exceptional quartic collision is nonreal]
\label{lem:f2v2-exceptional}
On the parameter region \(t>5\), \(c\in\mathbb R\), the polynomials
\(G\) and \(\partial_\alpha G\) have no common real zero.
\end{lemma}

\begin{proof}
By Lemma~\ref{lem:f2v2-structure}, a common zero can occur only when
\(L=0\).  Set \(q=t+w\).  Rewriting \textup{(N2.6)} gives

\[
 L=q^3+3q^2+9q-9-12t(q-1).
 \tag{N2.12}
\]
Consequently, on \(L=0\),

\[
 t=\frac{q^3+3q^2+9q-9}{12(q-1)},
 \qquad
 w=-\frac{(q-3)(q^2-6q+3)}{12(q-1)}.
 \tag{N2.13}
\]
Since \(q=t+w\), the present region has \(q>5\).

Run the polynomial Euclidean algorithm for \(G\) and
\(\partial_\alpha G\), regarding them as polynomials in \(\alpha\), and
let \(S_2\) be the primitive degree-two member of their subresultant
sequence.  The degree-one member is divisible by \(L\).  Substitution of
\textup{(N2.13)} into the quadratic discriminant gives the exact identity

\[
 \boxed{
 \operatorname{Disc}_{\alpha}S_2
 =-
 \frac{256q^{10}(q-3)^5
 (q^3-3q^2-3q-3)^7}
 {243(q-1)^{12}}.}
 \tag{N2.14}
\]
Every factor following the minus sign is positive for \(q>5\).  Hence
\(S_2\) has two nonreal conjugate zeros.  The common zeros of \(G\) and
\(\partial_\alpha G\) on \(L=0\) are precisely these zeros, so no real
multiple zero occurs.
\end{proof}

\begin{theorem}[Sharp unrestricted quintic height]
\label{thm:f2v2-quintic}
Let

\[
 P_0(z)=\bigl((z-x)^2+y^2\bigr)R_3(z),
 \tag{N2.15}
\]
where \(R_3\) is a monic real-rooted cubic and \(y\ge0\).  If

\[
 \exp\!\left(-\frac K2D_z^2\right)P_0
 \quad\hbox{is real-rooted}
 \tag{N2.16}
\]
for some \(K>0\), then

\[
 \boxed{y^2\le5K.}
 \tag{N2.17}
\]
The constant is sharp.  Equality forces, up to a real translation and
the scaling \(z\mapsto\sqrt K z\),

\[
 P_0(z)=z^3(z^2+5),
 \qquad
 \exp\!\left(-\frac12D_z^2\right)P_0(z)
 =z^3(z^2-5).
 \tag{N2.18}
\]
\end{theorem}

\begin{proof}
First work in the normalized first-collision coordinates of V1.  On the
connected parameter set

\[
 \Omega=\{(c,t):c\in\mathbb R,\ t>5\},
 \tag{N2.19}
\]
the auxiliary polynomial \(G(c,\alpha,t)\) is a quartic in \(\alpha\)
with positive leading coefficient.  The number of its real zeros can
change only at a real common zero of \(G\) and
\(\partial_\alpha G\).  Lemma~\ref{lem:f2v2-exceptional} excludes such a
zero everywhere in \(\Omega\), while
Lemma~\ref{lem:f2v2-centered} shows that the number is zero on the slice
\(c=0\).  Connectedness therefore gives

\[
 G(c,\alpha,t)>0
 \qquad(c,\alpha\in\mathbb R,\ t>5).
 \tag{N2.20}
\]
Since \(D>0\), \textup{(N2.2)} implies
\(\Delta_3(\alpha,\beta,\delta)<0\), contradicting the real-rootedness
of the padding cubic.  Thus normalized first-collision height cannot
exceed five.

For the boundary statement, let \(u\downarrow0\) in
\textup{(N2.20)}.  Then \(G\ge0\) at \(t=5\).  If \(w>0\), a real zero
would be a real multiple zero; equations \textup{(N2.5)}--\textup{(N2.7)}
and Lemma~\ref{lem:f2v2-exceptional}, now with \(q=5+w>5\), exclude it.
If \(w=0\), direct specialization gives

\[
 G(0,\alpha,5)=256\alpha^4+864\alpha^2,
 \tag{N2.21}
\]
so equality forces \(c=\alpha=0\).  Formula \textup{(N1.12)} then gives
\(\beta=\delta=0\), and \textup{(N2.18)} follows by direct heat
calculation.

Finally, for an arbitrary heat time \(K\), let \(T\le K\) be the first
all-real time, translate its collision point to zero, and scale by
\(\sqrt T\).  The normalized height is \(t=y^2/T\), hence \(y^2\le5T\le
5K\).  Equality forces \(T=K\) and the normalized equality case just
found.
\end{proof}

\begin{corollary}[Three real gamma padders do not cross the V97 line]
\label{cor:f2v2-gamma}
For a critical fixed gamma block having one conjugate target pair and
three real padding factors, real-rootedness requires

\[
 \kappa^2\ge
 \frac{\gamma^2R}{5(1+\theta)}.
 \tag{N2.22}
\]
Since its ladder count satisfies \(R\ge5\), necessarily

\[
 \boxed{
 \kappa\ge\frac{|\gamma|}{\sqrt{1+\theta}}.}
 \tag{N2.23}
\]
Thus arbitrary separation of three real padders cannot enter the open
band below the parity-neutral V97 boundary.  Equality can only reproduce
the centered triple-collision block with \(R=5\).
\end{corollary}

\begin{proof}
Use the inherited critical scaling

\[
 K=\frac{1+\theta}{\theta^2},
 \qquad
 t=\frac{\gamma^2R}{\theta^2\kappa^2},
 \tag{N2.24}
\]
and substitute Theorem~\ref{thm:f2v2-quintic}.
\end{proof}

\begin{firewall}[Scope of the quintic closure]
Theorem~\ref{thm:f2v2-quintic} is a complete optimization only for one
conjugate pair plus three real factors under a global fixed-block
backward-heat certificate.  It does not exclude four or more real
padders, degree-dependent meshes, local endpoint mechanisms, or a proof
of target-zero separation by a different argument.  It therefore does
not prove GRH.
\end{firewall}

\begin{assessment}[V2 advance and V3 direction]
V2 answers the five-ladder decision left by V1: the sharp unrestricted
quintic constant is exactly \(5\), so separation gives no improvement
over the coalesced odd Hermite block.  Together with V100, this closes
the first two nontrivial separated-padding optimizations.

The next nonduplicative finite problem has four real padders.  The V97
coalesced threshold is only \(\nu_4=5\), while entry below the current
gamma boundary would require a separated sextic height exceeding its
minimal six-ladder cost.  V3 should normalize the first sextic collision,
eliminate two quartic-padding coefficients, and impose the exact
quartic hyperbolicity minors before considering the residual output
quartic.  The target half-space is \(t>6\).
\end{assessment}

\section*{Version V2 append-only record}
\addcontentsline{toc}{section}{Version V2 append-only record}
The complete V1 source was retained before this continuation.  It had
\(11008\) bytes, \(332\) lines, and SHA--256

\[
 \texttt{79C27A571BF4CA43EF3818E79C009A40BE5F9067875C721C30C13E63ED960EA3}.
\]
V2 proves the sharp unrestricted quintic height \(y^2\le5K\), identifies
the unique equality block, and closes every three-real-padding critical
gamma block below the V97 parity-neutral line.

% =============================================================================
% END APPENDED V2 MATERIAL
% =============================================================================


% =============================================================================
% BEGIN APPENDED V3 MATERIAL
% =============================================================================

\section{V3: sextic collision cone and an exact shifted-padding witness}
\label{sec:grh-v3}

V2 closed the complete quintic problem at height \(5\).  The next fixed
block has one conjugate pair and four real padding zeros.  This section
derives its exact first-collision cone, gives a finite Hermite-matrix test
for both the input and the residual output, and proves two new facts.  A
shifted fourfold padding root gives an exact height strictly above the
centered Hermite value \(5\), but the entire shifted-coalesced stratum
stays strictly below the six-ladder target \(6\).

\subsection{Exact sextic first-collision equations}

At a normalized unit-time collision at the origin, write

\[
 \Pi(z)=\bigl((z-c)^2+t\bigr)R(z),
 \qquad
 R(z)=z^4+\alpha z^3+\beta z^2+\delta z+\varepsilon,
 \tag{N3.1}
\]
where \(R\) is real-rooted, and put \(q=c^2+t\), \(w=c^2\).  Since

\[
 Q=e^{-D_z^2/2}\Pi
 =\Pi-\frac12\Pi''+\frac18\Pi^{(4)}-\frac1{48}\Pi^{(6)},
 \tag{N3.2}
\]
the collision conditions \(Q(0)=Q'(0)=0\) are equivalent to

\[
\begin{aligned}
 2c\delta+(q-1)\varepsilon
  &=(q-3)\beta+6c\alpha-3q+15,\\
 (q-3)\delta-2c\varepsilon
  &=-6c\beta+3(q-5)\alpha+30c.
\end{aligned}
\tag{N3.3}
\]

\begin{lemma}[Forced last two padding coefficients]
\label{lem:f2v3-elimination}
If \(t>3\), then

\[
 \mathcal D:=4w+(q-1)(q-3)>0,
 \tag{N3.4}
\]
and equations \textup{(N3.3)} force

\[
\boxed{
\begin{aligned}
\delta={}&\frac{-4cq\beta+{12w+3(q-1)(q-5)}\alpha+24cq}
                 {\mathcal D},\\
\varepsilon={}&\frac{\{12w+(q-3)^2\}\beta+12c\alpha
                 -3(q-3)(q-5)-60w}{\mathcal D}.
\end{aligned}}
\tag{N3.5}
\]
The heat image factors as \(Q(z)=z^2S(z)\), where

\[
\begin{aligned}
S(z)={}&z^4+A_5z^3+(A_4-15)z^2
 +(A_3-10A_5)z+(A_2-6A_4+45),\\
A_5={}&\alpha-2c,\\
A_4={}&\beta-2c\alpha+q,\\
A_3={}&\delta-2c\beta+q\alpha,\\
A_2={}&\varepsilon-2c\delta+q\beta.
\end{aligned}
\tag{N3.6}
\]
\end{lemma}

\begin{proof}
For a monic sextic with coefficients \(a_jz^j\), equation
\textup{(N3.2)} gives

\[
 Q(0)=a_0-a_2+3a_4-15,
 \qquad
 Q'(0)=a_1-3a_3+15a_5.
 \tag{N3.7}
\]
Expanding \textup{(N3.1)} yields \textup{(N3.3)}.  Its determinant in
\((\delta,\varepsilon)\) is \(-\mathcal D\), and Cramer's rule gives
\textup{(N3.5)}.  Collecting the coefficients of \(Q/z^2\) gives
\textup{(N3.6)}.
\end{proof}

\subsection{A finite exact hyperbolicity test}

For a monic quartic

\[
 U(z)=z^4+u_1z^3+u_2z^2+u_3z+u_4,
 \tag{N3.8}
\]
define Newton sums \(s_0=4,s_1,\ldots,s_6\) by

\[
\begin{aligned}
s_k+u_1s_{k-1}+\cdots+u_{k-1}s_1+ku_k&=0
 &&(1\le k\le4),\\
s_k+u_1s_{k-1}+u_2s_{k-2}+u_3s_{k-3}+u_4s_{k-4}&=0
 &&(5\le k\le6),
\end{aligned}
\tag{N3.9}
\]
and its Hermite matrix

\[
 \mathcal H(U):=(s_{i+j})_{0\le i,j\le3}.
 \tag{N3.10}
\]

\begin{proposition}[Exact sextic feasibility cone]
\label{prop:f2v3-cone}
For \(t>3\), a normalized sextic \textup{(N3.1)} is an admissible
unit-time collision with real input padding and real residual output if
and only if \textup{(N3.5)} holds and

\[
 \boxed{
 \mathcal H(R)\succeq0,
 \qquad
 \mathcal H(S)\succeq0.}
 \tag{N3.11}
\]
Thus the target region \(t>6\) is an explicit semialgebraic problem in
\((c,\alpha,\beta,t)\): after clearing the positive denominator
\(\mathcal D\), it consists of the principal-minor inequalities for two
specified \(4\)-by-\(4\) symmetric matrices.
\end{proposition}

\begin{proof}
The classical Hermite criterion says that a real monic polynomial is
real-rooted, counting multiplicity, exactly when its Newton-sum Hermite
matrix is positive semidefinite.  Apply it first to \(R\) and then to the
quartic \(S=Q/z^2\).  Lemma~\ref{lem:f2v3-elimination} supplies all
remaining coefficients.  Positive semidefiniteness of a real symmetric
matrix is equivalent to nonnegativity of all its principal minors.
\end{proof}

There is also a useful root-coordinate chart.  For a fixed real-rooted
quartic \(R\), define

\[
\begin{array}{lll}
\mathcal A=\varepsilon-\beta+3,
&\mathcal C=\delta-3\alpha,
&\mathcal E=\varepsilon-3\beta+15,\\
\mathcal F=3\delta-15\alpha,
&\mathcal N=\mathcal A\mathcal E+\mathcal C^2.
\end{array}
\tag{N3.12}
\]
Equations \textup{(N3.3)} become

\[
 \mathcal A q+2\mathcal Cc=\mathcal E,
 \qquad
 \mathcal Cq-2\mathcal Ec=\mathcal F.
 \tag{N3.13}
\]
On the regular stratum \(\mathcal N\ne0\), therefore,

\[
 \boxed{
 q=\frac{\mathcal E^2+\mathcal C\mathcal F}{\mathcal N},
 \qquad
 c=\frac{\mathcal C\mathcal E-\mathcal A\mathcal F}{2\mathcal N},
 \qquad
 t=q-c^2.}
 \tag{N3.14}
\]
The singular stratum \(\mathcal N=0\) must be tested directly in
\textup{(N3.13)}; it is not discarded by division.

\subsection{The shifted-coalesced four-padding stratum}

Take

\[
 R(z)=(z-r)^4,
 \qquad \sigma=r^2\ge0.
 \tag{N3.15}
\]
Substitution in \textup{(N3.14)} gives

\[
\begin{aligned}
N(\sigma)&=\sigma^4-8\sigma^3+30\sigma^2+45,\\
q(\sigma)&=\frac{\sigma^4+12\sigma^3-30\sigma^2+180\sigma+225}
                  {N(\sigma)},\\
c(\sigma)&=\frac{4r\sigma(\sigma^2-6\sigma+15)}{N(\sigma)},\\
t(\sigma)&=\frac{\sigma^8-12\sigma^7+96\sigma^6-276\sigma^5
 +810\sigma^4+540\sigma^3+5400\sigma^2+8100\sigma+10125}
 {N(\sigma)^2}.
\end{aligned}
\tag{N3.16}
\]
Notice that

\[
 N(\sigma)=\sigma^2\{(\sigma-4)^2+14\}+45>0.
 \tag{N3.17}
\]

\begin{theorem}[The shifted-coalesced stratum stays below six]
\label{thm:f2v3-coalesced}
For every \(\sigma\ge0\), the collision height in \textup{(N3.16)}
satisfies

\[
 \boxed{t(\sigma)<6.}
 \tag{N3.18}
\]
\end{theorem}

\begin{proof}
Exact subtraction gives

\[
 6-t(\sigma)=\frac{P_8(\sigma)}{N(\sigma)^2},
 \tag{N3.19}
\]
where

\[
\begin{aligned}
P_8(X)={}&5X^8-84X^7+648X^6-2604X^5+5130X^4\\
&-4860X^3+10800X^2-8100X+2025.
\end{aligned}
\tag{N3.20}
\]
The exact Sturm chain of \(P_8\) has degrees
\((8,7,6,5,4,3,2,1,0)\).  Its signs at \(0\) and at \(+\infty\) are,
respectively,

\[
 (+,-,+,+,+,-,-,+,+),
 \qquad
 (+,+,-,-,+,+,-,-,+).
 \tag{N3.21}
\]
Both lists have four variations.  Hence \(P_8\) has no positive real
zero.  Since \(P_8(0)=2025>0\), it is positive on
\([0,\infty)\).  Equations \textup{(N3.17)} and \textup{(N3.19)} prove
the claim.
\end{proof}

The stratum nevertheless improves strictly on the centered value.

\begin{proposition}[Exact sextic witness above the centered threshold]
\label{prop:f2v3-witness}
Set

\[
 r=\frac1{\sqrt2},\qquad
 c=\frac{196\sqrt2}{825},\qquad
 q=\frac{989}{165},\qquad
 t=\frac{4002793}{680625}.
 \tag{N3.22}
\]
Then

\[
 5<t<6,
 \qquad
 \Pi(z)=\bigl((z-c)^2+t\bigr)(z-r)^4,
 \tag{N3.23}
\]
and \(e^{-D_z^2/2}\Pi\) is real-rooted.
\end{proposition}

\begin{proof}
Formula \textup{(N3.22)} is \textup{(N3.16)} at
\(\sigma=1/2\).  In particular,

\[
 t-5=\frac{599668}{680625}>0,
 \qquad
 6-t=\frac{80957}{680625}>0.
 \tag{N3.24}
\]
Direct heat calculation gives

\[
 e^{-D_z^2/2}\Pi(z)
 =\frac{z^2}{3300}\,
 \mathcal T(\sqrt2z),
 \tag{N3.25}
\]
where

\[
 \mathcal T(X)=825X^4-4084X^3-6774X^2+34116X-3911.
 \tag{N3.26}
\]
The exact values

\[
\begin{array}{c|rrrrrrr}
X&-3&-2&0&2&3&4&5\\ \hline
\mathcal T(X)&9868&-53367&-3911&17753&-5972&-26007&2444
\end{array}
\tag{N3.27}
\]
give sign changes in the four disjoint intervals
\((-3,-2)\), \((0,2)\), \((2,3)\), and \((4,5)\).  The quartic
\(\mathcal T\) therefore has four distinct real zeros, proving the
claim.
\end{proof}

\begin{corollary}[Distinct padding with a strict height gain]
\label{cor:f2v3-distinct}
There are four distinct real padding zeros and a heat ratio strictly
larger than \(5\) for which the associated sextic heat polynomial has
six simple real zeros.
\end{corollary}

\begin{proof}
Apply an additional heat time \(\eta>0\) to the real-rooted polynomial
in \textup{(N3.25)}.  Backward heat strictly separates its double zero,
so the time-\((1+\eta)\) image is simple and real-rooted.  For sufficiently
small \(\eta\), the ratio \(t/(1+\eta)\) remains above \(5\).  Simple
real-rootedness is open in the coefficient topology, so the fourfold
input root \(r\) may then be perturbed to four distinct real roots.
\end{proof}

\begin{firewall}[What the sextic witness does not prove]
The witness improves the \(p=4\) centered Hermite height, but its ratio is
less than the minimal six-ladder cost.  It therefore does not enter the
open gamma region below
\(|\gamma|/\sqrt{1+\theta}\).  Theorem~\ref{thm:f2v3-coalesced} closes
only the shifted-coalesced stratum; a fully separated real quartic must
still be tested through \textup{(N3.11)}.  Numerical maximization is not
used as a theorem.
\end{firewall}

\begin{assessment}[V3 advance and V4 direction]
The sextic problem is now an exact pair of \(4\)-by-\(4\) Hermite
semidefinite conditions after two linear eliminations.  The rational
algebraic witness proves that horizontal displacement matters already
on the fourfold-root boundary, while the Sturm certificate proves that
this entire boundary stratum cannot cross \(t=6\).

Independent global searches over ordered real padding roots and over the
coefficient cone converged near \(t=5.88561\) as all four roots coalesced.
This is discovery evidence only, but it points to an upstream V4 target:
derive \(6-t>0\) as a positive combination of the principal minors of
\(\mathcal H(R)\), treating the singular faces separately.  If such a
certificate exists, it closes every four-padding sextic block before the
residual-output inequalities are needed.
\end{assessment}

\section*{Version V3 append-only record}
\addcontentsline{toc}{section}{Version V3 append-only record}
The complete V2 source was retained before this continuation.  It had
\(20802\) bytes, \(668\) lines, and SHA--256

\[
 \texttt{5177D2B0AA34A73595BF4AE8C7F73A303F70137211D3F92A49E038A84DE98A53}.
\]
V3 derives the exact sextic collision cone, proves the shifted-coalesced
height is always below \(6\), and supplies an exact real-rooted sextic
witness of height \(4002793/680625>5\).

% =============================================================================
% END APPENDED V3 MATERIAL
% =============================================================================


% =============================================================================
% BEGIN APPENDED V4 MATERIAL
% =============================================================================

\section{V4: Grace--Walsh closure of the full sextic cone}
\label{sec:grh-v4}

V3 closed only the shifted-coalesced boundary of the sextic problem.  This
section removes that restriction.  The useful upstream observation is that
the two collision equations are symmetric multi-affine functions of the
four real padding zeros.  The Grace--Walsh--Szeg\H{o} coincidence theorem
therefore reduces a putative arbitrary-padding collision to a diagonal
complex point.  A Hermite--Biehler calculation excludes that point throughout
the entire height range \(t\ge 6\).

\subsection{The diagonal collision pair}

Write
\[
 R_{\boldsymbol r}(z)=\prod_{j=1}^{4}(z-r_j),\qquad
 \Pi_{\boldsymbol r}(z)=\bigl((z-c)^2+t\bigr)R_{\boldsymbol r}(z),
 \qquad
 Q_{\boldsymbol r}=e^{-D_z^2/2}\Pi_{\boldsymbol r}.
 \tag{N4.1}
\]
At a first collision translated to the origin, the necessary equations are
\[
 {\cal A}(\boldsymbol r):=Q_{\boldsymbol r}(0)=0,\qquad
 {\cal B}(\boldsymbol r):=Q_{\boldsymbol r}'(0)=0.
 \tag{N4.2}
\]
Both functions are real, symmetric, and multi-affine in
\(r_1,\ldots,r_4\).  Put \(q=c^2+t\).  On the diagonal
\(r_1=\cdots=r_4=r\), direct application of the heat operator gives
\[
\begin{aligned}
 A(r):={\cal A}(r,r,r,r)
 &=(q-1)r^4-8cr^3+6(3-q)r^2+24cr+3q-15,\\
 B(r):={\cal B}(r,r,r,r)
 &=-2cr^4-4(q-3)r^3+36cr^2+12(q-5)r-30c.
\end{aligned}
\tag{N4.3}
\]

\begin{lemma}[Exact discriminant and resultant certificate]
\label{lem:f2v4-factors}
For every \(c\in\mathbb R\) and \(t\ge6\), the polynomial \(A\) has four
simple real zeros.  If \(c\ne0\), the polynomial \(B\) also has four simple
real zeros.  Moreover \(A\) and \(B\) have no common zero.
\end{lemma}

\begin{proof}
Set \(u=t-6\) and \(w=c^2\).  Exact elimination in the variable \(r\)
gives
\[
 \operatorname{Disc}_r A=27648\,\Phi_A(u,w),\qquad
 \operatorname{Disc}_r B=110592\,\Phi_B(u,w),\qquad
 \operatorname{Res}_r(A,B)=27648\,\Phi_R(u,w).
 \tag{N4.4}
\]
For compactness, a row below lists the coefficients in increasing powers of
\(w\), and the rows are indexed by decreasing powers of \(u\).  The complete
coefficient arrays of the first two factors are
\[
\begin{array}{c|l}
 & [w^0,w^1,\ldots]\\ \hline
\Phi_A:\ u^6&(1)\\
u^5&(18,6)\\
u^4&(135,114,15)\\
u^3&(540,780,276,20)\\
u^2&(1203,2436,1770,324,15)\\
u^1&(1386,3702,4116,1740,186,6)\\
u^0&(605,2658,1587,3244,615,42,1)
\end{array}
\tag{N4.5}
\]
and
\[
\begin{array}{c|l}
\Phi_B:\ u^6&(4)\\
u^5&(48,24)\\
u^4&(228,399,60)\\
u^3&(544,2148,1116,80)\\
u^2&(684,4734,7212,1434,60)\\
u^1&(432,3780,15948,8892,876,24)\\
u^0&(108,-81,6912,21438,3600,207,4).
\end{array}
\tag{N4.6}
\]
Thus \(\Phi_A>0\).  Every \(u\)-dependent coefficient of \(\Phi_B\) is
positive, while its constant row is positive because
\[
 6912w^2-81w+108>0
 \quad\text{and}\quad
 (-81)^2-4(6912)(108)=-2979423<0.
 \tag{N4.7}
\]
Hence \(\Phi_B>0\).

All \(u\)-dependent coefficients of \(\Phi_R\) are also positive.  Their
complete rows are
\[
\begin{array}{c|l}
u^8&(1)\\
u^7&(20,8)\\
u^6&(180,180,28)\\
u^5&(940,1560,660,56)\\
u^4&(3050,7100,5820,1300,70)\\
u^3&(6156,19320,22840,11280,1500,56)\\
u^2&(7260,30828,41580,42360,11820,1020,28)\\
u^1&(4500,18120,57588,43800,39900,6360,380,8).
\end{array}
\tag{N4.8}
\]
The remaining row is
\[
\begin{aligned}
 K_0(w)={}&w^8+60w^7+1380w^6+14220w^5+50490w^4\\
 &-63084w^3+94860w^2-13500w+1125.
\end{aligned}
\tag{N4.9}
\]
Here is an exact positivity certificate avoiding numerical root isolation.
For a degree-eight polynomial \(p\), let
\[
 p(a+(b-a)x)=\sum_{k=0}^8 b_k\binom8k x^k(1-x)^{8-k}.
 \tag{N4.10}
\]
The Bernstein coefficient vectors for \(K_0\) on
\([0,1/5]\) and \([1/5,1]\), respectively, are
\[
\begin{aligned}
{\boldsymbol b}_{0}={}&
\biggl(1125,\frac{1575}{2},\frac{20493}{35},\frac{446277}{875},
\frac{2420214}{4375},\\
&\qquad
\frac{442734}{625},\frac{3032112}{3125},
\frac{20859384}{15625},\frac{703184176}{390625}\biggr),\\
{\boldsymbol b}_{1}={}&
\biggl(\frac{703184176}{390625},\frac{285996496}{78125},
\frac{111377776}{15625},\frac{37653136}{3125},\\
&\qquad
\frac{81608656}{4375},\frac{24192304}{875},\frac{1008112}{25},
\frac{293392}{5},85552\biggr).
\end{aligned}
\tag{N4.11}
\]
For \(w^8K_0(1/w)\) on \([0,1]\), the vector is
\[
 {\boldsymbol b}_{\infty}=
 \left(1,\frac{17}{2},\frac{457}{7},\frac{2977}{7},
 \frac{14446}{7},\frac{38854}{7},13312,33592,85552\right).
 \tag{N4.12}
\]
Every displayed coefficient is positive, and every Bernstein basis function
is nonnegative.  Equations \textup{(N4.11)}--\textup{(N4.12)} prove
\(K_0(w)>0\) for all \(w\ge0\).  Consequently \(\Phi_R>0\).

It remains to convert the nonvanishing discriminants into root counts.  At
\(c=0\),
\[
\begin{aligned}
 A(r)&=(t-1)r^4+6(3-t)r^2+3(t-5),\\
 B(r)&=-4r\bigl((t-3)r^2-3(t-5)\bigr).
\end{aligned}
\tag{N4.13}
\]
The quadratic for \(r^2\) in the first line has positive sum and product,
and discriminant \(24(t^2-6t+11)>0\); hence \(A\) has four simple real
zeros.  Since \(\operatorname{Disc}A\) never vanishes and its leading
coefficient \(q-1\) is positive, connectedness of
\(\mathbb R\times[6,\infty)\) preserves that root count.

At \(c=0\), \(B\) has the three simple roots
\[
 0,\qquad \pm\sqrt{\frac{3(t-5)}{t-3}}.
 \tag{N4.14}
\]
For sufficiently small nonzero \(c\), these persist and the fourth root is
real with asymptotic value \(-2(t-3)/c\).  On each of the connected
half-planes \(c>0\) and \(c<0\), the nonzero discriminant and nonzero leading
coefficient preserve four simple real roots.  Finally
\(\Phi_R>0\) says exactly that \(A\) and \(B\) never acquire a common root.
\end{proof}

\begin{lemma}[Stable diagonal collision polynomial]
\label{lem:f2v4-hb}
If \(t\ge6\), then
\[
 A(\zeta)-iB(\zeta)\ne0
 \qquad(\operatorname{Im}\zeta\ge0).
 \tag{N4.15}
\]
\end{lemma}

\begin{proof}
At \(c=0\), put \(s^2=3(t-5)/(t-3)\).  Besides the four-real-root
calculation in \textup{(N4.13)}, one has
\[
 A(0)=3(t-5)>0,\qquad
 A(s)=-\frac{6(t-5)(t^2-6t+15)}{(t-3)^2}<0.
 \tag{N4.16}
\]
Together with the positive leading coefficient and evenness of \(A\), these
signs show that the zeros of \(A\) strictly interlace the three finite zeros
of \(B\).  For small \(c>0\), the fourth zero of \(B\) enters from
\(-\infty\); for small \(c<0\), it enters from \(+\infty\).  Thus the two
quartics strictly interlace near \(c=0\) on both sides.  Lemma
\ref{lem:f2v4-factors} prevents either a multiple root or a common root, so
strict interlacing persists on each connected half-plane.  The limiting
case \(c=0\) is the degree-\((4,3)\) interlacing case.

The orientation is fixed by the Wronskian:
\[
 \bigl(AB'-A'B\bigr)(0)
 =36(q-5)^2+720c^2>0.
 \tag{N4.17}
\]
Strict interlacing makes this Wronskian have constant sign on the real
axis.  The Hermite--Biehler theorem, with the orientation in
\textup{(N4.17)}, now says that \(A-iB\) has all its zeros in the open
lower half-plane.  This proves \textup{(N4.15)}.
\end{proof}

\begin{theorem}[No sextic collision at normalized height six]
\label{thm:f2v4-sextic}
Let \(c,r_1,\ldots,r_4\in\mathbb R\) and \(t\ge6\).  For
\(\Pi_{\boldsymbol r}\) defined in \textup{(N4.1)}, the two equations
\[
 \bigl(e^{-D_z^2/2}\Pi_{\boldsymbol r}\bigr)(0)=
 \bigl(e^{-D_z^2/2}\Pi_{\boldsymbol r}\bigr)'(0)=0
 \tag{N4.18}
\]
cannot hold simultaneously.
\end{theorem}

\begin{proof}
Consider the symmetric multi-affine polynomial
\[
 {\cal F}(r_1,\ldots,r_4)
 ={\cal A}(r_1,\ldots,r_4)-i{\cal B}(r_1,\ldots,r_4).
 \tag{N4.19}
\]
Its coefficient of \(r_1r_2r_3r_4\) is
\((q-1)+2ic\ne0\), so it has total degree four.  If
\textup{(N4.18)} held, then \({\cal F}(\boldsymbol r)=0\).
The closed upper half-plane is a closed circular domain containing all four
real numbers \(r_j\).  The Grace--Walsh--Szeg\H{o} coincidence theorem
would therefore supply a \(\zeta\) in that half-plane such that
\[
 0={\cal F}(\boldsymbol r)
   ={\cal F}(\zeta,\zeta,\zeta,\zeta)
   =A(\zeta)-iB(\zeta),
 \tag{N4.20}
\]
contradicting Lemma~\ref{lem:f2v4-hb}.
\end{proof}

\begin{theorem}[Sharp-side global sextic obstruction]
\label{thm:f2v4-global}
Let
\[
 P_0(z)=\bigl((z-x)^2+y^2\bigr)\prod_{j=1}^{4}(z-a_j),
 \qquad x,y,a_j\in\mathbb R,\quad y>0.
 \tag{N4.21}
\]
If \(e^{-KD_z^2/2}P_0\) is real-rooted for some \(K>0\), then
\[
 \boxed{\,y^2<6K\,}.
 \tag{N4.22}
\]
\end{theorem}

\begin{proof}
Let \(T\in(0,K]\) be the first heat time at which the conjugate pair reaches
the real axis.  Translate its collision point to zero and scale space by
\(\sqrt T\).  The normalized input then has the form \textup{(N4.1)} with
\[
 t=\frac{y^2}{T}.
 \tag{N4.23}
\]
At the first collision the value and first derivative vanish, so
\textup{(N4.18)} holds.  If \(y^2\ge6K\), then
\(t=y^2/T\ge6K/T\ge6\), contradicting Theorem
\ref{thm:f2v4-sextic}.  Hence \textup{(N4.22)} is strict.
\end{proof}

\begin{corollary}[Four real gamma padders remain above the V97 line]
\label{cor:f2v4-gamma}
For a critical fixed gamma block with one conjugate target pair and four
real padding factors, real-rootedness requires
\[
 \kappa^2>
 \frac{\gamma^2R}{6(1+\theta)}
 \ge \frac{\gamma^2}{1+\theta},
 \qquad
 \boxed{\ \kappa>\frac{|\gamma|}{\sqrt{1+\theta}}\ }.
 \tag{N4.24}
\]
\end{corollary}

\begin{proof}
Insert the inherited scaling
\[
 K=\frac{1+\theta}{\theta^2},\qquad
 y^2=\frac{\gamma^2R}{\theta^2\kappa^2}
 \tag{N4.25}
\]
into Theorem~\ref{thm:f2v4-global}, and use the minimal ladder count
\(R\ge6\).
\end{proof}

\subsection{Upstream finite-padding engine}

The proof above suggests a uniform target.  Let
\[
 H_j(r)=e^{-D_r^2/2}r^j
\tag{N4.26}
\]
be the monic probabilists' Hermite polynomial.  For \(p\) coalesced real
padders, the diagonal collision pair, up to the common harmless factor
\((-1)^p\), is
\[
\begin{aligned}
 A_p(r)&=(q-1)H_p(r)-2cpH_{p-1}(r)+p(p-1)H_{p-2}(r),\\
 B_p(r)&=-2cH_p(r)-p(q-3)H_{p-1}(r)
 +2cp(p-1)H_{p-2}(r)\\
 &\hspace{7.6em}-p(p-1)(p-2)H_{p-3}(r).
\end{aligned}
\tag{N4.27}
\]
The same Grace--Walsh argument proves the following conditional reduction:
if \(A_p-iB_p\) is zero-free in the closed upper half-plane whenever
\(t\ge p+2\), then no arbitrary configuration of \(p\) real padding zeros
can collide at that height.  Since the corresponding gamma block has at
least \(p+2\) ladders, this one-variable stability statement would close
every fixed finite-padding architecture at the V97 boundary.

\begin{firewall}[What V4 proves and what remains open]
Theorem~\ref{thm:f2v4-global} is an unrestricted result for four real
padding factors: the separation left open in V3 is completely removed.
It does not prove the stability assertion following \textup{(N4.27)} for
all \(p\), does not address padding by additional nonreal pairs, and does
not supply the target-zero separation input missing from the Grand Tensor
programme.  In particular, it does not prove GRH.
\end{firewall}

\begin{assessment}[V4 advance and V5 direction]
The sextic optimization is now closed on the side relevant to the gamma
programme: arbitrary real separation cannot reach its minimal six-ladder
cost.  The key gain over V3 is structural rather than numerical---a
four-variable real-root problem is compressed by coincidence to the
one-variable stable polynomial \(A-iB\).

The next attack should go upstream to \textup{(N4.27)}.  Two concrete
routes are available: derive Hermite-interlacing and Wronskian recurrences
uniformly in \(p\), or certify upper-half-plane stability through the
symbol of the Hermite-preserving differential operator.  A proof at
general \(p\) would eliminate all finite blocks made from one target pair
and any fixed number of real padding zeros.
\end{assessment}

\section*{Version V4 append-only record}
\addcontentsline{toc}{section}{Version V4 append-only record}
The complete V3 source was retained before this continuation.  It had
\(31376\) bytes, \(1047\) lines, and SHA--256
\[
 \texttt{12763E7423E14636CE08CC294A803CBE5C3FB8E1D7FED3F576F584A90730A172}.
\]
V4 proves the strict unrestricted sextic bound \(y^2<6K\) by combining an
exact discriminant/resultant certificate, Hermite--Biehler stability, and
the Grace--Walsh--Szeg\H{o} coincidence theorem.

% =============================================================================
% END APPENDED V4 MATERIAL
% =============================================================================


% =============================================================================
% BEGIN APPENDED V5 MATERIAL
% =============================================================================

\section{V5: a uniform Pick--Grace--Walsh theorem for all real padding}
\label{sec:grh-v5}

V4 reduced the general finite-padding problem to upper-half-plane stability
of one diagonal Hermite polynomial.  The proposed discriminant induction is
unnecessary.  The logarithmic derivative of a Hermite polynomial satisfies
a sharp Pick inequality, and the hypothetical zero equation violates that
inequality by an elementary completed-square estimate.  This proves the
required stability simultaneously for every padding degree.

\subsection{The diagonal polynomial in logarithmic-derivative form}

Fix an integer \(p\ge1\), write
\[
 t=p+2+u,\qquad u\ge0,\qquad q=c^2+t,
 \tag{N5.1}
\]
and let \(H=H_p=e^{-D_z^2/2}z^p\) be the monic probabilists' Hermite
polynomial.  The diagonal collision polynomials from
\textup{(N4.27)} can be written
\[
\begin{aligned}
 A_p&=(q-1)H-2cH'+H'',\\
 B_p&=-2cH-(q-3)H'+2cH''-H'''.
\end{aligned}
\tag{N5.2}
\]
The Hermite differential equation
\[
 H''-zH'+pH=0
 \tag{N5.3}
\]
and the identity \(B_p=-A_p'+2(H'-cH)\) reduce these expressions to
\[
\boxed{
\begin{aligned}
 A_p(z)&=(c^2+u+1)H(z)+(z-2c)H'(z),\\
 B_p(z)&=(pz-2c(p+1))H(z)
 -\bigl((z-c)^2+u\bigr)H'(z).
\end{aligned}}
\tag{N5.4}
\]

\begin{lemma}[Hermite logarithmic-derivative Pick inequality]
\label{lem:f2v5-pick}
If \(z=x+iy\) with \(y>0\), then
\[
 m(z):=\frac{H_p'(z)}{H_p(z)}
 \quad\text{satisfies}\quad
 -p\,\operatorname{Im}m(z)-y\abs{m(z)}^2\ge0.
 \tag{N5.5}
\]
\end{lemma}

\begin{proof}
The zeros \(\lambda_1,\ldots,\lambda_p\) of \(H_p\) are real and simple,
so
\[
 m(z)=\sum_{j=1}^p\frac1{z-\lambda_j},\qquad
 -\operatorname{Im}m(z)
 =y\sum_{j=1}^p\frac1{\abs{z-\lambda_j}^2}.
 \tag{N5.6}
\]
Cauchy--Schwarz gives
\[
 \abs{m(z)}^2
 \le p\sum_{j=1}^p\frac1{\abs{z-\lambda_j}^2},
 \tag{N5.7}
\]
which is exactly \textup{(N5.5)}.
\end{proof}

\begin{theorem}[Uniform diagonal upper-half-plane stability]
\label{thm:f2v5-diagonal}
For every \(p\ge1\), \(c\in\mathbb R\), and \(u\ge0\),
\[
 F_p(z):=A_p(z)-iB_p(z)\ne0
 \qquad(\operatorname{Im}z>0).
 \tag{N5.8}
\]
On the real axis, \(F_p\) has a zero if and only if
\[
 u=0,\qquad c=0,\qquad p\ \text{is odd},
 \tag{N5.9}
\]
and in that exceptional case its only real zero is \(z=0\).
\end{theorem}

\begin{proof}
Suppose first that \(z=x+iy\), \(y>0\), and \(F_p(z)=0\).  Since \(H_p\)
has only real zeros, division by \(H_p(z)\) is legitimate.  Put
\[
 a=c^2+u+1,\qquad
 D=(z-2c)+i\bigl((z-c)^2+u\bigr),
 \tag{N5.10}
\]
and
\[
 N=-a+i\bigl(pz-2c(p+1)\bigr).
 \tag{N5.11}
\]
Equation \textup{(N5.4)} says \(D\,m=N\).  Notice that
\(\operatorname{Re}N=-a-py<0\); hence \(N\ne0\), and consequently
\(D\ne0\).

Set \(X=x-c\) and \(M=c^2+p+u+1\).  Direct expansion of
\textup{(N5.10)}--\textup{(N5.11)} gives the exact identity
\[
 \abs{D}^2\bigl(-p\,\operatorname{Im}m-y\abs{m}^2\bigr)
 =-pMX^2+2cp(p+1)X-\mathcal C,
 \tag{N5.12}
\]
where
\[
\begin{aligned}
\mathcal C={}&c^4y\\
&+c^2\bigl(
p^2y+p^2+pu+py^2+5py+2p+2uy+6y
\bigr)\\
&+p^2uy+p^2y^2+pu^2+puy^2+puy+pu\\
&+py^2+py+u^2y+2uy+y.
\end{aligned}
\tag{N5.13}
\]
Every term in \(\mathcal C\) is nonnegative, and
\[
 \mathcal C\ge c^2p(p+2).
 \tag{N5.14}
\]
If \(c\ne0\), completing the square in \(X\) and using
\(M\ge c^2+p+1\) gives
\[
\begin{aligned}
 -pMX^2+2cp(p+1)X-\mathcal C
 &\le \frac{c^2p(p+1)^2}{M}-\mathcal C\\
 &\le
 -\frac{c^2p\bigl((p+2)c^2+p+1\bigr)}
 {c^2+p+1}<0.
\end{aligned}
\tag{N5.15}
\]
If \(c=0\), the right side of \textup{(N5.12)} is also strictly negative,
because \(\mathcal C\ge y>0\).  This contradicts Lemma
\ref{lem:f2v5-pick} and proves \textup{(N5.8)}.

It remains to classify real zeros.  Let \(z\in\mathbb R\) and suppose
\(A_p(z)=B_p(z)=0\).  If \(H(z)=0\), simplicity gives \(H'(z)\ne0\).
Formula \textup{(N5.4)} then forces
\[
 z=2c,\qquad (z-c)^2+u=0,
 \tag{N5.16}
\]
so \(c=u=z=0\).  This is a zero of \(H_p\) exactly when \(p\) is odd.

If \(H(z)\ne0\), eliminate the real number \(H'(z)/H(z)\) from the two
equations in \textup{(N5.4)}.  With \(X=z-c\), the result is
\[
\begin{aligned}
 J(X)={}&
 (c^2+p+u+1)X^2-2c(p+1)X\\
 &+c^2(p+u+2)+u(u+1)=0.
\end{aligned}
\tag{N5.17}
\]
For \(c\ne0\), the same completed-square estimate as in
\textup{(N5.15)} shows \(J(X)>0\).  For \(c=0\),
\[
 J(X)=(p+u+1)X^2+u(u+1),
 \tag{N5.18}
\]
which can vanish only when \(u=X=0\).  But then
\(A_p(0)=H_p(0)\ne0\) in the present case.  Therefore the sole real-axis
possibility is precisely \textup{(N5.9)}, and direct substitution confirms
that \(F_p(0)=0\) there.
\end{proof}

\subsection{Coincidence and the unrestricted finite-padding theorem}

For arbitrary real padding zeros \(\boldsymbol r=(r_1,\ldots,r_p)\), set
\[
\begin{aligned}
 \Pi_{\boldsymbol r}(z)
 &=\bigl((z-c)^2+t\bigr)\prod_{j=1}^p(z-r_j),\\
 {\cal A}_p(\boldsymbol r)
 &=\bigl(e^{-D_z^2/2}\Pi_{\boldsymbol r}\bigr)(0),\\
 {\cal B}_p(\boldsymbol r)
 &=\bigl(e^{-D_z^2/2}\Pi_{\boldsymbol r}\bigr)'(0).
\end{aligned}
\tag{N5.19}
\]

\begin{theorem}[Uniform collision obstruction for real padding]
\label{thm:f2v5-collision}
If \(t>p+2\), the equations
\[
 {\cal A}_p(\boldsymbol r)={\cal B}_p(\boldsymbol r)=0
 \tag{N5.20}
\]
have no solution with \(c,r_1,\ldots,r_p\in\mathbb R\).  They also have
no solution at \(t=p+2\) when \(p\) is even.
\end{theorem}

\begin{proof}
The polynomial
\[
 {\cal F}_p={\cal A}_p-i{\cal B}_p
 \tag{N5.21}
\]
is symmetric and multi-affine in the \(p\) padding variables.  Its
coefficient of \(r_1\cdots r_p\) is
\((q-1)+2ic\ne0\), so its total degree is \(p\).  If
\textup{(N5.20)} held, the Grace--Walsh--Szeg\H{o} coincidence theorem,
applied to the closed upper half-plane, would give a
\(\zeta\) in that half-plane for which
\[
 0={\cal F}_p(\boldsymbol r)
 ={\cal F}_p(\zeta,\ldots,\zeta)
 =A_p(\zeta)-iB_p(\zeta).
 \tag{N5.22}
\]
Theorem~\ref{thm:f2v5-diagonal} excludes \(\operatorname{Im}\zeta>0\).
It also excludes a real \(\zeta\) unless \(t=p+2\), \(c=0\), and \(p\)
is odd.  This proves both assertions.
\end{proof}

\begin{theorem}[All finite real padding stays above the \(p+2\) line]
\label{thm:f2v5-global}
Let
\[
 P_0(z)=\bigl((z-x)^2+y^2\bigr)\prod_{j=1}^p(z-a_j),
 \qquad x,y,a_j\in\mathbb R,\quad y>0.
 \tag{N5.23}
\]
If \(e^{-KD_z^2/2}P_0\) is real-rooted for some \(K>0\), then
\[
 \boxed{\,y^2\le(p+2)K\,}.
 \tag{N5.24}
\]
The inequality is strict when \(p\) is even.  For every odd \(p\), equality
is attained by the centered coalesced input
\[
 P_0(z)=\bigl(z^2+(p+2)K\bigr)z^p.
 \tag{N5.25}
\]
\end{theorem}

\begin{proof}
Let \(T\in(0,K]\) be the first heat time at which the target conjugate pair
reaches the real axis.  Translation of the collision point and scaling by
\(\sqrt T\) put its two collision equations in the form
\textup{(N5.20)}, with normalized height \(t=y^2/T\).  If
\(y^2>(p+2)K\), then
\[
 t=\frac{y^2}{T}>(p+2)\frac KT\ge p+2,
 \tag{N5.26}
\]
contrary to Theorem~\ref{thm:f2v5-collision}.  If \(p\) is even and
\(y^2=(p+2)K\), then equality throughout would force \(T=K\) and
\(t=p+2\), which the same theorem also excludes.

For sharpness when \(p=2m+1\), scale first to \(K=1\).  Direct heat
calculation and the Hermite--Laguerre identity give
\[
\begin{aligned}
 e^{-D_z^2/2}\bigl((z^2+p+2)z^p\bigr)
 &=H_{p+2}(z)+(p+2)H_p(z)\\
 &=(-1)^m2^m m!\,z^3
 L_m^{(3/2)}\left(\frac{z^2}{2}\right).
\end{aligned}
\tag{N5.27}
\]
The generalized Laguerre polynomial \(L_m^{(3/2)}\) has \(m\) positive
simple zeros.  Thus the right side is real-rooted.  Rescaling proves
\textup{(N5.25)} for arbitrary \(K>0\).
\end{proof}

\begin{corollary}[Every finite real-padding gamma block misses the open band]
\label{cor:f2v5-gamma}
For a critical fixed gamma block with one conjugate target pair and any
finite number \(p\ge1\) of real padding factors, real-rootedness requires
\[
 \kappa^2\ge
 \frac{\gamma^2R}{(p+2)(1+\theta)}
 \ge\frac{\gamma^2}{1+\theta}.
 \tag{N5.28}
\]
Consequently
\[
 \boxed{\ \kappa\ge\frac{|\gamma|}{\sqrt{1+\theta}}\ },
 \tag{N5.29}
\]
with strict inequality in \textup{(N5.28)} when \(p\) is even.  No such
finite real-padding architecture enters the open V97 band.
\end{corollary}

\begin{proof}
Use the critical scaling
\[
 K=\frac{1+\theta}{\theta^2},\qquad
 y^2=\frac{\gamma^2R}{\theta^2\kappa^2},
 \tag{N5.30}
\]
Theorem~\ref{thm:f2v5-global}, and the ladder count \(R\ge p+2\).
\end{proof}

\begin{firewall}[Scope of the uniform theorem]
Theorem~\ref{thm:f2v5-global} closes every finite block consisting of one
target conjugate pair and arbitrary real padding, uniformly in the number
and separation of the padding zeros.  It does not cover padding by additional
nonreal conjugate pairs, a nonpolynomial or cutoff-dependent limiting
architecture, or the missing target-zero separation theorem.  It therefore
does not prove GRH.
\end{firewall}

\begin{assessment}[V5 advance and V6 direction]
The conditional programme stated after \textup{(N4.27)} is now proved for
all \(p\).  The parity effect is exact: odd padding degrees attain the
\(p+2\) threshold through the centered Hermite--Laguerre block, whereas even
degrees cannot attain it.  Together with the sharper low-degree constants
already proved, this removes all purely real finite padding as a route into
the unresolved gamma band.

The next upstream obstruction is therefore qualitative rather than
degree-theoretic: any surviving fixed polynomial block must contain at least
one additional nonreal conjugate pair.  V6 should formulate the first
two-pair collision problem and test whether the Pick estimate can be upgraded
from the scalar Hermite logarithmic derivative to a matrix-valued or
multivariate stable logarithmic derivative.
\end{assessment}

\section*{Version V5 append-only record}
\addcontentsline{toc}{section}{Version V5 append-only record}
The complete V4 source was retained before this continuation.  It had
\(43818\) bytes, \(1425\) lines, and SHA--256
\[
 \texttt{D9DF32E5B1566C1AE58FFC3A80E1E00F7B971858CB5B1EDDEA3F851A164B69DE}.
\]
V5 proves the uniform bound \(y^2\le(p+2)K\) for one conjugate pair and
every finite collection of real padding zeros, with strict inequality for
even \(p\), by combining the Hermite logarithmic-derivative Pick inequality
with Grace--Walsh coincidence.

% =============================================================================
% END APPENDED V5 MATERIAL
% =============================================================================


% =============================================================================
% BEGIN APPENDED V6 MATERIAL
% =============================================================================

\section{V6: the sharp Hermite strip for arbitrary complex padding}
\label{sec:grh-v6}

V5 closes every fixed block whose padding zeros are real.  The first
surviving case has a second nonreal conjugate pair, and it genuinely behaves
differently: a quartic two-pair block can exceed the four-ladder height.
Instead of treating each number of nonreal pairs separately, this section
proves a sharp degree-\(n\) strip theorem for the inverse heat operator.
The theorem permits arbitrary conjugation-closed input geometry and reduces
the entire fixed-block obstruction to the largest zero of one Hermite
polynomial.

\subsection{Grace--Walsh proof of the inverse-heat strip}

For \(K>0\), define
\[
 H_n^{(K)}(z)=e^{-(K/2)D_z^2}z^n,
 \qquad
 G_n^{(K)}(z)=e^{+(K/2)D_z^2}z^n.
 \tag{N6.1}
\]
Let
\[
 -\rho_n\le\lambda_1<\cdots<\lambda_n\le\rho_n
 \tag{N6.2}
\]
be the zeros of \(H_n^{(1)}\), so \(\rho_n\) is its largest zero.
The elementary rotation identity
\[
 G_n^{(K)}(w)
 =i^{-n}H_n^{(K)}(iw)
 \tag{N6.3}
\]
shows that the zeros of \(G_n^{(K)}\) are
\[
 -i\sqrt K\,\lambda_1,\ldots,-i\sqrt K\,\lambda_n.
 \tag{N6.4}
\]

\begin{theorem}[Sharp inverse-heat Hermite strip]
\label{thm:f2v6-strip}
Let \(Q\) be a real-rooted real polynomial of degree \(n\ge1\), and put
\[
 P=e^{+(K/2)D_z^2}Q.
 \tag{N6.5}
\]
Every zero \(\zeta\) of \(P\) satisfies
\[
 \boxed{\ \abs{\operatorname{Im}\zeta}\le\rho_n\sqrt K\ }.
 \tag{N6.6}
\]
The constant is sharp for every \(n\): equality is attained when
\(Q(z)=(z-a)^n\), \(a\in\mathbb R\).
\end{theorem}

\begin{proof}
Write
\[
 Q(z)=\prod_{j=1}^n(z-r_j),\qquad r_j\in\mathbb R,
 \tag{N6.7}
\]
where repeated roots are allowed.  For a fixed \(\zeta\), consider
\[
 {\cal W}_{\zeta}(r_1,\ldots,r_n)
 =
 \left.
 e^{+(K/2)D_z^2}\prod_{j=1}^n(z-r_j)
 \right|_{z=\zeta}.
 \tag{N6.8}
\]
This is symmetric and multi-affine in \(r_1,\ldots,r_n\).  Its coefficient
of \(r_1\cdots r_n\) is \((-1)^n\), so its total degree is \(n\).

Suppose \(P(\zeta)=0\) and write \(y=\operatorname{Im}\zeta\).  Apply the
Grace--Walsh--Szeg\H{o} coincidence theorem first with the closed lower
half-plane, which contains all the \(r_j\).  It supplies a
\(\xi\) with \(\operatorname{Im}\xi\le0\) such that
\[
\begin{aligned}
 0
 &={\cal W}_{\zeta}(r_1,\ldots,r_n)
 ={\cal W}_{\zeta}(\xi,\ldots,\xi)\\
 &=G_n^{(K)}(\zeta-\xi).
\end{aligned}
\tag{N6.9}
\]
By \textup{(N6.4)}, for some zero \(\lambda_j\) of \(H_n^{(1)}\),
\[
 \zeta-\xi=-i\sqrt K\,\lambda_j,
 \qquad
 \operatorname{Im}\xi=y+\sqrt K\,\lambda_j.
 \tag{N6.10}
\]
If \(y>\rho_n\sqrt K\), the last expression is positive for every \(j\),
contradicting \(\operatorname{Im}\xi\le0\).  Applying the same coincidence
theorem with the closed upper half-plane excludes
\(y<-\rho_n\sqrt K\).  This proves \textup{(N6.6)}.

If \(Q(z)=(z-a)^n\), then
\[
 P(z)=G_n^{(K)}(z-a),
 \tag{N6.11}
\]
whose extreme zeros have imaginary parts
\(\pm\rho_n\sqrt K\) by \textup{(N6.4)}.  Thus the constant cannot be
reduced.
\end{proof}

\begin{corollary}[Backward heat with unrestricted input geometry]
\label{cor:f2v6-backward}
Let \(P_0\) be any real polynomial of degree \(n\).  If
\[
 e^{-(K/2)D_z^2}P_0
 \tag{N6.12}
\]
is real-rooted, then every zero \(x+iy\) of \(P_0\) obeys
\[
 y^2\le\rho_n^2K.
 \tag{N6.13}
\]
This conclusion allows any number of real zeros and nonreal conjugate
pairs, with arbitrary multiplicities and horizontal placements.
\end{corollary}

\begin{proof}
Take \(Q=e^{-(K/2)D_z^2}P_0\).  The heat operators terminate on
polynomials and are mutual inverses, so
\(P_0=e^{+(K/2)D_z^2}Q\).  Apply Theorem~\ref{thm:f2v6-strip}.
\end{proof}

\subsection{The first two-pair block is exactly solvable}

\begin{corollary}[Sharp quartic two-pair height]
\label{cor:f2v6-quartic}
Suppose a real quartic \(P_0\) contains a designated conjugate pair
\(x\pm iy\), and \(e^{-(K/2)D_z^2}P_0\) is real-rooted.  Then
\[
 \boxed{\ y^2\le(3+\sqrt6)K\ }.
 \tag{N6.14}
\]
The bound is attained by the centered two-pair polynomial
\[
\begin{aligned}
 P_0(z)
 &=e^{+(K/2)D_z^2}z^4\\
 &=z^4+6Kz^2+3K^2\\
 &=\bigl(z^2+(3+\sqrt6)K\bigr)
   \bigl(z^2+(3-\sqrt6)K\bigr).
\end{aligned}
\tag{N6.15}
\]
\end{corollary}

\begin{proof}
The largest zero of
\[
 H_4^{(1)}(z)=z^4-6z^2+3
 \tag{N6.16}
\]
has square \(3+\sqrt6\), so \textup{(N6.14)} follows from
Corollary~\ref{cor:f2v6-backward}.  Formula \textup{(N6.15)} is a direct
heat calculation, and its backward heat image is \(z^4\).
\end{proof}

Thus complex padding is not merely a formal loophole in V5.  With four
gamma ladders, the sharp two-pair block permits normalized squared target
height \(3+\sqrt6>4\), whereas V5 proves that no configuration of two real
padding zeros can reach height \(4\).

\subsection{A universal arbitrary-padding gamma barrier}

The number \(\rho_n\) also has a useful elementary degree bound.

\begin{lemma}[Elementary largest-Hermite-zero bound]
\label{lem:f2v6-rho}
For \(n\ge2\),
\[
 \rho_n<2\sqrt{n-1}.
 \tag{N6.17}
\]
\end{lemma}

\begin{proof}
In the orthonormal Hermite basis of degrees \(0,\ldots,n-1\), the zeros of
\(H_n^{(1)}\) are the eigenvalues of the symmetric Jacobi matrix
\[
 J_n=
 \begin{pmatrix}
 0&\sqrt1\\
 \sqrt1&0&\sqrt2\\
 &\sqrt2&0&\ddots\\
 &&\ddots&\ddots&\sqrt{n-1}\\
 &&&\sqrt{n-1}&0
 \end{pmatrix}.
 \tag{N6.18}
\]
For a unit vector \(v=(v_0,\ldots,v_{n-1})\),
\[
\begin{aligned}
 \abs{v^{\mathsf T}J_nv}
 &\le
 2\sum_{k=1}^{n-1}\sqrt{k}\,\abs{v_{k-1}v_k}\\
 &\le
 \sum_{k=1}^{n-1}\sqrt{k}
 \bigl(v_{k-1}^2+v_k^2\bigr)\\
 &\le
 \bigl(\sqrt{n-2}+\sqrt{n-1}\bigr)\norm{v}^2
 <2\sqrt{n-1}.
\end{aligned}
\tag{N6.19}
\]
The middle coefficient bound is interpreted directly when \(n=2\).
Taking the largest Rayleigh quotient proves \textup{(N6.17)}.
\end{proof}

\begin{theorem}[Fixed-block barrier with arbitrary complex padding]
\label{thm:f2v6-gamma}
Consider any critical fixed gamma block of polynomial degree \(r\ge2\) that
contains the designated target conjugate pair, and let \(R\ge r\) be the
total gamma-ladder count in the critical scaling.  The remaining factors
may form any conjugation-closed collection of real zeros and nonreal pairs.
If its critical heat envelope is real-rooted, then
\[
 \boxed{
 \kappa^2\ge
 \frac{\gamma^2R}{\rho_r^2(1+\theta)}
 >
 \frac{\gamma^2}{4(1+\theta)}.}
 \tag{N6.20}
\]
In particular, no fixed global-real-rootedness block of arbitrary complex
padding geometry can certify
\[
 \boxed{
 \kappa\le
 \frac{\abs{\gamma}}{2\sqrt{1+\theta}}.}
 \tag{N6.21}
\]
For a quartic block, the exact estimate is
\[
 \kappa^2\ge
 \frac{\gamma^2R}{(3+\sqrt6)(1+\theta)}
 \ge
 \frac{4\gamma^2}{(3+\sqrt6)(1+\theta)},
 \tag{N6.22}
\]
and the centered two-pair block \textup{(N6.15)} attains the strip equality
at the heat-polynomial level.
\end{theorem}

\begin{proof}
In the inherited critical scaling, the target root has squared imaginary
part and heat time
\[
 y^2=\frac{\gamma^2R}{\theta^2\kappa^2},
 \qquad
 K=\frac{1+\theta}{\theta^2}.
 \tag{N6.23}
\]
Corollary~\ref{cor:f2v6-backward} gives \(y^2\le\rho_r^2K\), which is the
first inequality in \textup{(N6.20)}.  Lemma~\ref{lem:f2v6-rho} gives
\[
 \frac{R}{\rho_r^2}\ge\frac{r}{\rho_r^2}>
 \frac{r}{4(r-1)}>\frac14,
 \tag{N6.24}
\]
proving the second inequality and \textup{(N6.21)}.  The quartic statement
uses \(\rho_4^2=3+\sqrt6\) and Corollary~\ref{cor:f2v6-quartic}.
\end{proof}

\begin{firewall}[What the unrestricted strip theorem does not prove]
Theorem~\ref{thm:f2v6-gamma} removes the real-padding hypothesis from the
fixed-block obstruction, but its uniform constant is \(1/2\), not the V97
construction boundary \(1\).  It does not exclude fixed complex-padding
blocks in the intervening region, blocks whose degree grows with the Jensen
scale, local variation certificates weaker than global real-rootedness, or
the missing target-zero separation input.  It therefore does not prove GRH.
\end{firewall}

\begin{assessment}[V6 advance and V7 direction]
The first nonreal-padding problem is solved sharply, and the same argument
simultaneously controls every fixed conjugation-closed block.  The exact
quartic factorization \textup{(N6.15)} proves that complex padding really
does enter part of the band closed to real padding.  At the same time, the
Hermite strip shows that no finite amount or arrangement of complex padding
can push a target pair below the universal half-line
\(\abs{\gamma}/(2\sqrt{1+\theta})\).

The remaining fixed-block question is constructive: for each degree \(R\),
which Hermite-strip heights can be attained while respecting the prescribed
gamma-ladder offsets, rather than arbitrary conjugation-closed roots?
Upstream of that optimization, V7 should compare the sharp Hermite extremizer
\(e^{+(K/2)D^2}z^R\) with the arithmetic offset constraints and identify the
first coefficient or spacing invariant that prevents its realization by a
genuine gamma block.
\end{assessment}

\section*{Version V6 append-only record}
\addcontentsline{toc}{section}{Version V6 append-only record}
The complete V5 source was retained before this continuation.  It had
\(54364\) bytes, \(1776\) lines, and SHA--256
\[
 \texttt{9925BE3A9B0DEF2722BE61DFEA7CEE73ED4949A284854ABE7CC3F0E82851441E}.
\]
V6 proves the sharp inverse-heat Hermite strip for arbitrary polynomial root
geometry, solves the quartic two-pair height exactly, and extends the
fixed-block gamma obstruction to every finite mixture of real and nonreal
padding.

% =============================================================================
% END APPENDED V6 MATERIAL
% =============================================================================


% =============================================================================
% BEGIN APPENDED V7 MATERIAL
% =============================================================================

\section{V7: equal-height gamma pairs and a quantitative degree barrier}
\label{sec:grh-v7}

V6 shows that unrestricted complex padding can outperform real padding, and
that the sharp quartic strip extremizer uses two centered conjugate pairs of
very different heights.  Genuine factors drawn twice from one
target-character gamma ladder do not have that freedom.  Their real offsets
change with the ladder index, but their imaginary offsets are both
\(\pm\gamma\).  This section solves that equal-height two-pair problem
exactly and quantifies how the degree of any more elaborate complex block
must diverge as one approaches the universal half-line.

\subsection{The arithmetic equal-height constraint}

The archived gamma nodes have the form
\[
 \sigma+\kappa_\chi+2m\pm i\gamma,
 \qquad m=0,1,2,\ldots.
 \tag{N7.1}
\]
After the critical affine normalization used in the backward-heat cluster
limit, two fixed ladder indices therefore contribute factors
\[
 \bigl((z-c_1)^2+y^2\bigr)
 \bigl((z-c_2)^2+y^2\bigr)
 \tag{N7.2}
\]
with one common squared height \(y^2\).  Distinct indices give
\(c_1\ne c_2\).  Translation by \((c_1+c_2)/2\) puts their centers at
\(\pm h\), where \(h=(c_1-c_2)/2\).

\begin{theorem}[Sharp equal-height two-pair obstruction]
\label{thm:f2v7-equal-pairs}
Let
\[
 P_0(z)=\bigl((z-h)^2+y^2\bigr)
         \bigl((z+h)^2+y^2\bigr),
 \qquad h,y\in\mathbb R,\quad y>0.
 \tag{N7.3}
\]
If \(e^{-(K/2)D_z^2}P_0\) is real-rooted for some \(K>0\), then
\[
 \boxed{\ y^2\le\frac32K\ }.
 \tag{N7.4}
\]
Equality holds if and only if \(h=0\), in which case
\[
 P_0(z)=\left(z^2+\frac32K\right)^2,
 \qquad
 e^{-(K/2)D_z^2}P_0(z)
 =\left(z^2-\frac32K\right)^2.
 \tag{N7.5}
\]
Consequently the inequality is strict for two distinct horizontal centers.
\end{theorem}

\begin{proof}
Scale by \(\sqrt K\), and set
\[
 t=\frac{y^2}{K},\qquad w=\frac{h^2}{K}.
 \tag{N7.6}
\]
Apart from the positive factor \(K^2\), direct heat calculation gives
\[
\begin{aligned}
 Q(z):={}&e^{-D_z^2/2}
 \bigl(((z-\sqrt w)^2+t)((z+\sqrt w)^2+t)\bigr)\\
 ={}&z^4+\{2(t-w)-6\}z^2
 +(t+w)^2-2(t-w)+3.
\end{aligned}
\tag{N7.7}
\]
If \(Q\) is real-rooted, then the quadratic in \(X=z^2\) has real roots.
Its discriminant is
\[
\begin{aligned}
 \Delta_X
 &=
 \{2(t-w)-6\}^2
 -4\{(t+w)^2-2(t-w)+3\}\\
 &=24-16t+16w(1-t).
\end{aligned}
\tag{N7.8}
\]
If \(t\le1\), then \textup{(N7.4)} is immediate.  If \(t>1\), the last
term in \textup{(N7.8)} is nonpositive, so
\[
 0\le\Delta_X\le24-16t,
 \tag{N7.9}
\]
and again \(t\le3/2\).  At \(t=3/2\), equation
\textup{(N7.8)} forces \(w=0\).  Substitution gives
\[
 Q(z)=z^4-3z^2+\frac94
 =\left(z^2-\frac32\right)^2.
 \tag{N7.10}
\]
This proves the equality statement after rescaling.
\end{proof}

\begin{corollary}[A second pair from the same gamma ladder cannot help]
\label{cor:f2v7-same-ladder}
Consider a critical quartic block made from two conjugate pairs belonging
to the same ordinate \(\gamma\), and let \(R\ge4\) be the total ladder
count.  Global real-rootedness of its heat envelope requires
\[
 \boxed{
 \kappa^2\ge
 \frac{2\gamma^2R}{3(1+\theta)}
 \ge
 \frac{8\gamma^2}{3(1+\theta)}.}
 \tag{N7.11}
\]
The first inequality is strict when the two pairs come from distinct ladder
indices.  Thus duplicating the target-character gamma ladder moves the
required line in the wrong direction and cannot enter the V97 band.
\end{corollary}

\begin{proof}
The critical scaling is
\[
 y^2=\frac{\gamma^2R}{\theta^2\kappa^2},
 \qquad
 K=\frac{1+\theta}{\theta^2}.
 \tag{N7.12}
\]
Insert these expressions into Theorem~\ref{thm:f2v7-equal-pairs} and use
\(R\ge4\).  Distinct ladder indices have different real centers by
\textup{(N7.1)}, so the equality case in
Theorem~\ref{thm:f2v7-equal-pairs} is unavailable.
\end{proof}

For comparison, the sharp V6 quartic construction is
\[
 \bigl(z^2+(3+\sqrt6)K\bigr)
 \bigl(z^2+(3-\sqrt6)K\bigr).
 \tag{N7.13}
\]
Its companion pair has height ratio
\[
 \frac{y_-}{y_+}
 =
 \sqrt{\frac{3-\sqrt6}{3+\sqrt6}}
 =\sqrt3-\sqrt2
 =0.3178\ldots.
 \tag{N7.14}
\]
Thus the construction that makes complex padding useful is arithmetically
different from repeating the same gamma ladder: it needs a substantially
smaller second ordinate, not another copy of \(\gamma\).

\subsection{How fast a surviving complex block must grow}

The Jacobi-matrix proof in V6 actually recorded the sharper estimate
\[
 \rho_r\le b_r:=
 \sqrt{r-2}+\sqrt{r-1}
 \qquad(r\ge2).
 \tag{N7.15}
\]
Combining this with the exact Hermite-strip coefficient gives a useful
degree lower bound.

\begin{theorem}[Quantitative degree blow-up near the half-line]
\label{thm:f2v7-degree}
Let a critical gamma block have polynomial degree \(r\ge3\), total ladder
count \(R\ge r\), and a globally real-rooted heat envelope.  Define
\[
 \Lambda:=
 \frac{\kappa^2(1+\theta)}{\gamma^2}.
 \tag{N7.16}
\]
Then
\[
 \boxed{
 \Lambda\ge\frac{R}{\rho_r^2}
 \ge\frac{r}{b_r^2}
 >
 \frac{r}{4r-6}.}
 \tag{N7.17}
\]
In particular, if \(\Lambda>1/4\), any such block must satisfy
\[
 \boxed{
 r>\frac{6\Lambda}{4\Lambda-1}.}
 \tag{N7.18}
\]
If \(\Lambda\le1/4\), no finite degree is possible.
\end{theorem}

\begin{proof}
The first inequality is Theorem~\ref{thm:f2v6-gamma}.  The next two use
\(R\ge r\), \textup{(N7.15)}, and
\[
 b_r^2
 =2r-3+2\sqrt{(r-1)(r-2)}
 <4r-6.
 \tag{N7.19}
\]
If \(\Lambda>1/4\), rearranging the strict final inequality in
\textup{(N7.17)} gives
\[
 (4\Lambda-1)r>6\Lambda,
 \tag{N7.20}
\]
which is \textup{(N7.18)}.  When \(\Lambda\le1/4\),
\textup{(N7.17)} is impossible for every finite \(r\).
\end{proof}

Writing \(\Lambda=1/4+\varepsilon\) shows explicitly that a putative
global-real-rootedness certificate just above the half-line needs
\[
 r>\frac{3}{8\varepsilon}+\frac32.
 \tag{N7.21}
\]
Hence approaching the V6 boundary is intrinsically a growing-block problem,
not an optimization over any fixed list of gamma factors.

\begin{firewall}[Arithmetic progress without target separation]
Theorem~\ref{thm:f2v7-equal-pairs} rules out the most immediate genuine
gamma realization of the V6 complex-padding loophole, namely taking two
indices from the same nonreal gamma ladder.  It does not exclude a block
built from several channels with different ordinates, and
\textup{(N7.14)} is a requirement of the displayed sharp extremizer rather
than a theorem that every near-extremizer has that exact ratio.
Theorem~\ref{thm:f2v7-degree} forces degree growth near the half-line but
does not control the resulting growing Jensen block.  GRH remains open.
\end{firewall}

\begin{assessment}[V7 advance and V8 direction]
The sharp arbitrary-root strip of V6 is not attainable by simply duplicating
the target gamma ladder: equal-height pairs have the much lower exact
threshold \(3K/2\), and distinct ladder indices make it strict.  Moreover,
the remaining complex-padding escape cannot stay at bounded degree as
\(\kappa\) approaches the half-line.

The next useful upstream test is the ordered-height quartic problem.  If the
selected target is the least positive ordinate among all off-line channels,
every additional nonreal pair has height at least that of the target.
V8 should decide rigorously whether real-rooted backward heat then forces
\(\min(y_1^2,y_2^2)\le3K/2\) even when the two heights differ.  Such a theorem
would rule out every two-channel complex padding block attached to a
least-ordinate counterexample, not only repetitions of one ladder.
\end{assessment}

\section*{Version V7 append-only record}
\addcontentsline{toc}{section}{Version V7 append-only record}
The complete V6 source was retained before this continuation.  It had
\(64135\) bytes, \(2094\) lines, and SHA--256
\[
 \texttt{5EFAF59D400D1F9F67DE95CE844CAE4C2C8B86EEDE0023939CE76DA29778C596}.
\]
V7 proves the sharp \(3K/2\) threshold for two equal-height conjugate pairs,
shows that a second pair from the same gamma ladder cannot improve the
critical construction, and quantifies the degree blow-up required near the
arbitrary-padding half-line.

% =============================================================================
% END APPENDED V7 MATERIAL
% =============================================================================


% =============================================================================
% BEGIN APPENDED V8 MATERIAL
% =============================================================================

\section{V8: the ordered two-pair theorem}
\label{sec:grh-v8}

V7 solved two conjugate pairs only when their imaginary heights agree.
The arithmetic target isolated there is stronger: if a hypothetical
counterexample is chosen with the least ordinate among the two channels in
a block, the padding pair may be higher but not lower.  This section proves
that the same \(3K/2\) obstruction survives without equal heights.  The
proof needs only the second and fourth power-sum conditions for a
real-rooted quartic.

\subsection{Two necessary real-rooted moments}

After translating the mean of the two pair centers and scaling by
\(\sqrt K\), write
\[
 P(z)=\bigl((z-\sqrt w)^2+t\bigr)
      \bigl((z+\sqrt w)^2+s\bigr),
 \qquad 0<t\le s,\quad w\ge0.
 \tag{N8.1}
\]
Its unit-time backward heat image is
\[
 Q(z)=e^{-D_z^2/2}P(z)=z^4+Az^2+Bz+C,
 \tag{N8.2}
\]
where
\[
\begin{aligned}
 A&=t+s-2w-6,\\
 B&=2\sqrt w\,(t-s),\\
 C&=(w+t)(w+s)-(t+s-2w)+3.
\end{aligned}
\tag{N8.3}
\]
If the roots of \(Q\) are \(\xi_1,\ldots,\xi_4\), their sum is zero.
Newton's identities give
\[
 p_2:=\sum_{j=1}^4\xi_j^2=-2A,\qquad
 p_4:=\sum_{j=1}^4\xi_j^4=2A^2-4C.
 \tag{N8.4}
\]
Real-rootedness therefore forces
\[
 p_2\ge0,\qquad
 4p_4-p_2^2\ge0,
 \tag{N8.5}
\]
the second inequality being Cauchy--Schwarz.

\begin{theorem}[Sharp ordered-height two-pair obstruction]
\label{thm:f2v8-ordered}
Let
\[
 P_0(z)=
 \bigl((z-c_1)^2+y_1^2\bigr)
 \bigl((z-c_2)^2+y_2^2\bigr),
 \qquad c_j,y_j\in\mathbb R,\quad 0<y_1\le y_2.
 \tag{N8.6}
\]
If \(e^{-(K/2)D_z^2}P_0\) is real-rooted for some \(K>0\), then
\[
 \boxed{\ y_1^2\le\frac32K\ }.
 \tag{N8.7}
\]
Equality holds if and only if
\[
 c_1=c_2,\qquad
 y_1^2=y_2^2=\frac32K.
 \tag{N8.8}
\]
\end{theorem}

\begin{proof}
Translation and scaling reduce to \textup{(N8.1)}.  Suppose for a
contradiction that
\[
 t=\frac32+u,\qquad
 s=t+v,\qquad
 u>0,\quad v\ge0,
 \tag{N8.9}
\]
and put
\[
 L:=2u+v>0.
 \tag{N8.10}
\]
The first condition in \textup{(N8.5)} and
\(A=L-2w-3\) imply
\[
 w\ge\frac{L-3}{2}.
 \tag{N8.11}
\]
The second condition in \textup{(N8.5)} has the exact reduction
\[
\begin{aligned}
 0\le4p_4-p_2^2
 &=4(A^2-4C)\\
 &=4\left[
 v^2-8\{(L+1)w+L\}
 \right].
\end{aligned}
\tag{N8.12}
\]
Thus
\[
 v^2\ge8\{(L+1)w+L\}.
 \tag{N8.13}
\]
Since \(0\le v<L\), one also has \(v^2<L^2\).

If \(0<L\le3\), then \(w\ge0\) and
\[
 L^2>v^2\ge8L,
 \tag{N8.14}
\]
which is impossible.  If \(L>3\), use \textup{(N8.11)} in
\textup{(N8.13)} to obtain
\[
 v^2\ge
 8\left\{\frac{(L+1)(L-3)}2+L\right\}
 =4(L^2-3).
 \tag{N8.15}
\]
Together with \(v^2<L^2\), this gives \(L^2<4\), contradicting \(L>3\).
Hence \(u\le0\), proving \textup{(N8.7)}.

For equality, set \(u=0\), so \(L=v\).  If \(0<v\le3\), then
\textup{(N8.13)} gives \(v^2\ge8v\), a contradiction.  If \(v>3\),
\textup{(N8.11)}--\textup{(N8.15)} give \(v^2\ge4(v^2-3)\), hence
\(v\le2\), again a contradiction.  Therefore \(v=0\), and
\textup{(N8.13)} then forces \(w=0\).  Conversely,
\[
 e^{-D_z^2/2}\left(z^2+\frac32\right)^2
 =\left(z^2-\frac32\right)^2,
 \tag{N8.16}
\]
so the stated equality configuration works.  Undoing translation and
scaling proves \textup{(N8.8)}.
\end{proof}

\begin{corollary}[Every two-pair block contains a shallow pair]
\label{cor:f2v8-shallow}
If a quartic made from two nonreal conjugate pairs has a real-rooted
backward-heat image at time \(K\), then at least one input pair satisfies
\[
 y^2\le\frac32K.
 \tag{N8.17}
\]
The V6 strip extremizer
\[
 \bigl(z^2+(3+\sqrt6)K\bigr)
 \bigl(z^2+(3-\sqrt6)K\bigr)
 \tag{N8.18}
\]
achieves its large height only by introducing the shallow companion
\((3-\sqrt6)K<3K/2\).  It cannot make both pairs simultaneously high.
\end{corollary}

\begin{corollary}[Least-ordinate two-channel gamma obstruction]
\label{cor:f2v8-gamma}
Consider a critical quartic gamma block of total ladder count \(R\ge4\).
Suppose its designated target pair has ordinate magnitude no larger than
that of the other conjugate pair.  If the heat envelope is real-rooted, then
\[
 \boxed{
 \kappa^2\ge
 \frac{2\gamma^2R}{3(1+\theta)}
 \ge
 \frac{8\gamma^2}{3(1+\theta)}.}
 \tag{N8.19}
\]
The first inequality is strict unless the two normalized pairs have the
same center and the same height.  In particular, a second channel of equal
or larger ordinate cannot pad a least-ordinate target into the unresolved
V97 band.
\end{corollary}

\begin{proof}
The target pair has
\[
 y_1^2=\frac{\gamma^2R}{\theta^2\kappa^2},
 \qquad
 K=\frac{1+\theta}{\theta^2}.
 \tag{N8.20}
\]
Apply Theorem~\ref{thm:f2v8-ordered} and use \(R\ge4\).
The equality statement follows from \textup{(N8.8)}.
\end{proof}

\begin{firewall}[The ordinate-ordering hypothesis matters]
Theorem~\ref{thm:f2v8-ordered} controls the smaller of the two imaginary
heights.  It does not bound the larger pair by \(3K/2\); V6 proves that the
larger height can reach \((3+\sqrt6)K\) when accompanied by a much smaller
pair.  Therefore Corollary~\ref{cor:f2v8-gamma} applies when the designated
target is least in the chosen two-pair block.  A proof architecture could
still try to pad a higher target using a lower-ordinate channel, and a
growing family of channels need not possess a uniform positive least
ordinate.  GRH remains open.
\end{firewall}

\begin{assessment}[V8 advance and V9 direction]
The two-pair arithmetic question posed in V7 is now closed: equality of the
heights was not essential.  Every realifiable two-pair quartic sacrifices
one pair below \(3K/2\), with a unique equality geometry.  Thus a
least-ordinate hypothetical counterexample cannot be helped by any second
channel whose ordinate is equal or larger.

The moment proof suggests the next upstream invariant.  For a degree
\(2m\) input consisting of \(m\) conjugate pairs, the second and fourth
power sums of the real-rooted heat image can be expressed exactly in terms
of the pair centers and heights.  V9 should determine whether these two
moments alone force at least one pair below a universal multiple of \(K\),
and whether the resulting bound prevents a finite collection of
nondecreasing ordinates from jointly certifying its least member.
\end{assessment}

\section*{Version V8 append-only record}
\addcontentsline{toc}{section}{Version V8 append-only record}
The complete V7 source was retained before this continuation.  It had
\(72491\) bytes, \(2360\) lines, and SHA--256
\[
 \texttt{EADCCE26B8801854928B10A777E04F6CE3208BA767D2ECE40D2CDE535179DFAA}.
\]
V8 proves the sharp ordered two-pair bound
\(\min(y_1^2,y_2^2)\le3K/2\), classifies equality, and rules out every
two-channel complex padding block whose designated target is the
least-ordinate pair.

% =============================================================================
% END APPENDED V8 MATERIAL
% =============================================================================


% =============================================================================
% BEGIN APPENDED V9 MATERIAL
% =============================================================================

\section{V9: a shallow-pair theorem for arbitrary finite complex padding}
\label{sec:grh-v9}

V8 proved the sharp estimate \(y_{\min}^2\le 3K/2\) for two conjugate
pairs.  That numerical constant does not persist when a third pair is
allowed.  The useful invariant is instead dynamical: before the first
nonreal pair reaches the real axis, every other conjugate pair can increase
the downward speed of a currently shallowest pair by at most one unit.
This gives a degree-dependent theorem for an arbitrary finite collection of
pairs and, more importantly, a degree-independent gamma-ladder barrier.

\subsection{The quartic constant really is exceptional}

\begin{proposition}[A three-pair example above \(3K/2\)]
\label{prop:f2v9-sextic}
For every \(K>0\), the polynomial
\[
 P(z)=\left(z^2+\frac95K\right)^3
 \tag{N9.1}
\]
has three conjugate pairs, all of squared height \(9K/5\), while
\(e^{-(K/2)D_z^2}P\) is real-rooted.  Consequently no theorem for
\(m\ge3\) pairs can retain the V8 bound \(y_{\min}^2\le3K/2\).
\end{proposition}

\begin{proof}
Scale to \(K=1\).  Direct application of the finite heat series gives
\[
 e^{-D_z^2/2}\left(z^2+\frac95\right)^3
 =z^6-\frac{48}{5}z^4+\frac{558}{25}z^2-\frac{336}{125}.
 \tag{N9.2}
\]
Thus it is enough to show that
\[
 H(X):=125X^3-1200X^2+2790X-336
 \tag{N9.3}
\]
has three positive roots.  The exact signs
\[
\begin{array}{c|rrrrrr}
X&0&1/5&7/2&4&5&6\\ \hline
H(X)&-336&175&707/8&-376&-761&204
\end{array}
\tag{N9.4}
\]
place one root in each of
\((0,1/5)\), \((7/2,4)\), and \((5,6)\).  Since \(H\) is cubic these are all
its roots.  Substituting \(X=z^2\) proves real-rootedness of
\textup{(N9.2)}.  Scaling proves the assertion for general \(K\).
\end{proof}

The example is not an artefact of multiplicity.  The heat image in
\textup{(N9.2)} has six simple real zeros.  Perturbing those zeros within
the open simple-real-rooted cone and applying the inverse finite heat
operator produces polynomials with three distinct pairs, each still having
squared height greater than \(3/2\).

\subsection{The pairwise speed inequality}

Let
\[
 P_\tau(z):=e^{-(\tau/2)D_z^2}P_0(z).
 \tag{N9.5}
\]
At a time when the zeros are simple, differentiation of
\(P_\tau(z_i(\tau))=0\) gives the standard zero equation
\[
 z_i'(\tau)
 =\frac{P_\tau''(z_i)}{2P_\tau'(z_i)}
 =\sum_{\ell\ne i}\frac1{z_i-z_\ell}.
 \tag{N9.6}
\]

\begin{lemma}[One other pair costs at most one unit of speed]
\label{lem:f2v9-pair-speed}
Let \(x+iy\), \(y>0\), be a currently shallowest upper-half-plane zero,
and let \(X\pm iY\) be any other conjugate pair.  Thus \(Y\ge y\).
The contribution of these two zeros to \((y^2)'\) is at least \(-1\).
\end{lemma}

\begin{proof}
Put
\[
 a:=\frac{Y}{y}\ge1,
 \qquad
 u:=\frac{x-X}{y}.
 \tag{N9.7}
\]
Taking imaginary parts of the two summands in \textup{(N9.6)} and
multiplying by \(2y\) shows that their contribution is
\[
 \Phi(a,u)
 =\frac{2(a-1)}{u^2+(a-1)^2}
  -\frac{2(a+1)}{u^2+(a+1)^2}.
 \tag{N9.8}
\]
Away from a collision its denominator is nonzero, and elementary
combination gives
\[
 \Phi(a,u)
 =\frac{4(a^2-1-u^2)}
 {\{u^2+(a-1)^2\}\{u^2+(a+1)^2\}}.
 \tag{N9.9}
\]
Writing \(q=u^2+1\ge1\), one obtains
\[
\begin{aligned}
 &\{u^2+(a-1)^2\}\{u^2+(a+1)^2\}
       \{\Phi(a,u)+1\}\\
 &\qquad=(q+a^2)^2-4q\ge0.
\end{aligned}
\tag{N9.10}
\]
Indeed \(a^2\ge1\) and \(q+1\ge2\sqrt q\), so
\(q+a^2\ge2\sqrt q\).  Hence \(\Phi(a,u)\ge-1\).
\end{proof}

The conjugate \(x-iy\) of the selected zero contributes exactly \(-1\) to
\((y^2)'\).  A real zero \(r\) contributes
\[
 -\frac{2y^2}{(x-r)^2+y^2}\ge-2.
 \tag{N9.11}
\]
These three local estimates are the complete speed budget.

\subsection{The arbitrary finite-block theorem}

\begin{theorem}[Shallow-pair first-collision bound]
\label{thm:f2v9-shallow}
Suppose a real monic polynomial \(P_0\) of degree \(2m+p\) has precisely
\(m\ge1\) nonreal conjugate pairs and \(p\ge0\) real zeros, counted with
multiplicity.  Write its upper-half-plane zeros as
\[
 x_j+iy_j,\qquad y_j>0,\qquad 1\le j\le m.
 \tag{N9.12}
\]
If
\[
 e^{-(K/2)D_z^2}P_0
 \quad\hbox{is real-rooted}
 \tag{N9.13}
\]
for some \(K>0\), then
\[
 \boxed{
 \min_{1\le j\le m}y_j^2\le(m+2p)K.}
 \tag{N9.14}
\]
In particular, for a pure \(m\)-pair block,
\[
 \boxed{y_{\min}^2\le mK.}
 \tag{N9.15}
\]
\end{theorem}

\begin{proof}
First impose the generic condition that all collisions along
\(0\le\tau\le K\) are isolated and have the usual algebraic root
parametrizations.  Let \(T\) be the first time at which one of the nonreal
pairs reaches the real axis.  Such a time exists and satisfies \(T\le K\)
because \(P_K\) is real-rooted.  Before \(T\) there are \(m\) nonreal pairs
and \(p\) real zeros.

Set
\[
 h(\tau):=\min_{1\le j\le m}y_j(\tau)^2.
 \tag{N9.16}
\]
At every simple time where \(h\) is differentiable, choose a minimizing
pair.  Its conjugate contributes \(-1\) to \(h'\).  Each of the other
\(m-1\) pairs contributes at least \(-1\) by
Lemma~\ref{lem:f2v9-pair-speed}, and each real zero contributes at least
\(-2\) by \textup{(N9.11)}.  Therefore
\[
 h'(\tau)\ge-1-(m-1)-2p=-(m+2p)
 \tag{N9.17}
\]
at almost every such time.

The minimum of finitely many smooth branches obeys the same inequality at
almost every ordinary branch crossing.  At an isolated multiple zero,
the algebraic root branches have convergent Puiseux parametrizations; the
one-sided square-root singularities are integrable, and \(h\) remains
continuous and locally absolutely continuous across the collision.
Thus \textup{(N9.17)} integrates over \([0,T]\).  Since \(h(T)=0\),
\[
 h(0)\le(m+2p)T\le(m+2p)K.
 \tag{N9.18}
\]

For completeness, remove the genericity condition while preserving the
number of real input zeros.  Put
\[
 Q:=e^{-(K/2)D_z^2}P_0
\]
and fix \(\delta>0\).  The strict backward-heat image
\[
 Q_\delta:=e^{-(\delta/2)D_z^2}Q
 =e^{-((K+\delta)/2)D_z^2}P_0
 \tag{N9.19}
\]
has simple real zeros: real-rootedness is preserved, and the real-zero
flow is strictly repulsive for positive time.  It therefore lies in the
open simple-real-rooted cone.

Now perturb the factors of \(P_0\) directly, keeping its \(p\) zeros real
and its \(m\) conjugate pairs off the axis, and avoid the proper algebraic
conditions for nongeneric intermediate collisions.  Denote the resulting
generic polynomials by \(P_{0,\nu}\to P_0\).  For all large \(\nu\),
\[
 e^{-((K+\delta)/2)D_z^2}P_{0,\nu}
\]
is close to \(Q_\delta\), hence is still simple real-rooted.  The generic
case at time \(K+\delta\) gives
\[
 \min_j y_{j,\nu}^2\le(m+2p)(K+\delta).
\]
Letting \(\nu\to\infty\) and then \(\delta\downarrow0\), using continuity
of polynomial zeros, proves \textup{(N9.14)}.  The case \(p=0\) is
\textup{(N9.15)}.
\end{proof}

\begin{remark}[Relation to the archived one-pair theorem]
For \(m=1\), \textup{(N9.14)} is exactly the archived V98 speed budget
\(y^2\le(2p+1)K\).  The new point is Lemma~\ref{lem:f2v9-pair-speed}:
an arbitrary higher conjugate pair has cost at most one, rather than the
cost two of a real padding zero.  For \(m=2,p=0\), V8 improves the resulting
\(2K\) to the sharp \(3K/2\).  Proposition~\ref{prop:f2v9-sextic} explains
why that sharper constant cannot be used unchanged for all \(m\).
\end{remark}

\subsection{Translation to ordered gamma blocks}

\begin{corollary}[Finite mixed-padding barrier for a least ordinate]
\label{cor:f2v9-gamma}
Consider a fixed critical gamma block with \(m\ge1\) conjugate-pair channels
and \(p\ge0\) real gamma factors.  Let its total ladder count be \(R\), so
\[
 R\ge2m+p.
 \tag{N9.20}
\]
Suppose the designated target has least ordinate magnitude \(\abs{\gamma}\)
among the conjugate pairs in the block.  If its critical heat envelope is
real-rooted, then
\[
 \boxed{
 \kappa^2\ge
 \frac{\gamma^2R}{(m+2p)(1+\theta)}.}
 \tag{N9.21}
\]
Consequently every such finite mixed block satisfies the strict universal
necessary condition
\[
 \boxed{
 \kappa^2>\frac{\gamma^2}{2(1+\theta)}.}
 \tag{N9.22}
\]
For pure complex padding (\(p=0\)), the stronger condition is
\[
 \boxed{
 \kappa^2\ge\frac{2\gamma^2}{1+\theta}.}
 \tag{N9.23}
\]
\end{corollary}

\begin{proof}
The normalized height of the least-ordinate target and the heat time are
\[
 y_{\min}^2=\frac{\gamma^2R}{\theta^2\kappa^2},
 \qquad
 K=\frac{1+\theta}{\theta^2}.
 \tag{N9.24}
\]
Insert these identities in \textup{(N9.14)} to obtain
\textup{(N9.21)}.  Moreover,
\[
 \frac{R}{m+2p}
 \ge\frac{2m+p}{m+2p}
 =\frac12+\frac{3m}{2(m+2p)}
 >\frac12,
 \tag{N9.25}
\]
which proves \textup{(N9.22)}.  If \(p=0\), then \(R/m\ge2\), proving
\textup{(N9.23)}.
\end{proof}

\begin{firewall}[What the shallow-pair theorem does not prove]
The theorem controls the least imaginary height in a fixed finite block.
It does not control a designated higher pair when lower-ordinate channels
are available, and it does not justify passage to a block whose size grows
with the Jensen degree.  Real padding also changes the speed budget from
\(m\) to \(m+2p\), so the strong pure-complex estimate
\textup{(N9.23)} cannot be transferred to a mixed block.  Finally, this is
a necessary condition for one heat-envelope architecture, not a proof that
the corresponding global tensor or completed \(L\)-function has only real
Jensen zeros.  GRH remains open.
\end{firewall}

\begin{assessment}[V9 advance and V10 direction]
The finite nondecreasing-ordinate problem posed in V8 is now solved at the
first-collision level.  No fixed collection of equal-or-higher conjugate
channels can hide its least member: some pair must enter the strip
\(y^2\le mK\), and after gamma normalization the growth of the ladder count
cancels the growth of this strip.  Allowing arbitrary finite real padding
still cannot cross the universal line \textup{(N9.22)}.

V10 is a mandatory ten-version sweep of newly generated companion notes.
After that sweep, the next mathematical question is whether a weighted
minimum, rather than the unweighted envelope \(h\), can reduce the real-zero
cost in \textup{(N9.17)}, or whether an explicit mixed-block family proves
that the coefficient \(m+2p\) is globally unavoidable.  A second route is
to remove the least-ordinate hypothesis by finding a monotone statistic that
charges lower companion channels rather than discarding them.
\end{assessment}

\section*{Version V9 append-only record}
\addcontentsline{toc}{section}{Version V9 append-only record}
The complete V8 source was retained before this continuation.  It had
\(79495\) bytes, \(2598\) lines, and SHA--256
\[
 \texttt{D6B53177F4A655CB3D56944EDE08CACDFD2A8E6F7E3DD14E4374705948FAC707}.
\]
V9 proves the multi-pair first-collision bound
\(y_{\min}^2\le(m+2p)K\), specializes it to \(mK\) for pure complex padding,
and converts it into the fixed finite mixed-block gamma barrier
\(\kappa^2>\gamma^2/[2(1+\theta)]\) for a least-ordinate target.  The explicit
centered sextic at height \(9K/5\) also proves that the sharp V8 constant
\(3K/2\) is special to two pairs.

% =============================================================================
% END APPENDED V9 MATERIAL
% =============================================================================


% =============================================================================
% BEGIN APPENDED V10 MATERIAL
% =============================================================================

\section{V10: companion sweep and the exact centered three-pair threshold}
\label{sec:grh-v10}

This tenth active version begins with the required sweep of newly generated
parallel notes.  It then sharpens V9 on the first nontrivial family beyond
the quartic: three equal-height conjugate pairs with a common center.

\subsection{Mandatory companion-note sweep}

The note set changed while the sweep was in progress: the active SM file
advanced from V76 to V77 and the Yang--Mills file from V93 to V94.  Both
predecessors were inspected before replacement, and the appended replacement
sections were then read.  The final read-only snapshot used here is
\[
\begin{array}{c|r|r}
\text{checkpoint}&\text{bytes}&\text{lines}\\ \hline
\text{SM--Einstein V77}&639930&19587\\
\text{Yang--Mills V94}&728997&19914\\
\text{Navier--Stokes V26}&278283&5319
\end{array}
\tag{N10.1}
\]
with SHA--256 values
\begin{center}\scriptsize
\texttt{AE620D19E57185EBCFD0A37C1190EB85FC6C3E6F66ECFA0D84AB298B3724C9EC}\\
\texttt{645ABB0A97779B963A4E82819ADE7FF7DF1ABFC154FDF5703013C9A972D65899}\\
\texttt{82C887F8C2C69E4F28FAD7C3DC8DE5B9099CB5A4BF431EBA1FFF3E91935978AD}.
\end{center}

SM V77 differentiates an exact conditional-imbalance law and proves that a
same-monotonicity covariance shortcut fails.  YM V94 certifies three new
finite Wilson-symbol charts and a second-coordinate slab.  NS V26 remains a
pointwise Oseen-kernel tensor-norm calculation.

\begin{theorem}[V10 companion-sweep decision]
\label{thm:f2v10-sweep}
No theorem in the companion snapshot supplies a backward-heat
real-rootedness estimate, a zero-flow statistic, a gamma-ladder inequality,
or a cofinal Jensen certificate.  Consequently no companion theorem can be
imported into the GRH proof chain at V10.
\end{theorem}

\begin{proof}
The SM result concerns a scalar conditional covariance in a conformal--Higgs
law; it has neither a polynomial heat flow nor a completed-\(L\)-function
variable.  The YM result is a rational finite-dimensional Loewner lower
bound on local momentum boxes; it supplies no control of zeros of a scalar
Jensen polynomial and explicitly leaves its own cofinal step open.  The NS
result concerns rank-one contractions of Oseen-kernel derivatives.  None of
their hypotheses identifies an object in the V9 finite-block flow, and none
of their conclusions implies its real-rootedness.  Importing any of them
would therefore be a type mismatch.
\end{proof}

\subsection{The algebraic threshold}

\begin{lemma}[The three-pair threshold constant]
\label{lem:f2v10-alpha}
The cubic
\[
 F(\eta):=8\eta^3-42\eta^2+90\eta-75
 \tag{N10.2}
\]
has exactly one real zero, denoted by \(\alpha_3\).  It satisfies
\[
 \frac{19}{10}<\alpha_3<2,
 \qquad
 \alpha_3=1.9434591570485828280\ldots .
 \tag{N10.3}
\]
\end{lemma}

\begin{proof}
One has
\[
 F'(\eta)
 =6(4\eta^2-14\eta+15)
 =6\left\{4\left(\eta-\frac74\right)^2+\frac{11}{4}\right\}>0.
 \tag{N10.4}
\]
Thus \(F\) is strictly increasing.  Since
\[
 F(19/10)=-187/250<0,\qquad F(2)=1>0,
 \tag{N10.5}
\]
there is exactly one real zero and it has the stated bracket.  Newton
iteration with rational interval checks gives the displayed decimal; only
the defining equation and bracket are used below.
\end{proof}

\begin{theorem}[Sharp centered equal-height three-pair obstruction]
\label{thm:f2v10-triple}
Let \(K>0\) and \(\eta\ge0\).  Then
\[
 e^{-(K/2)D_z^2}\left(z^2+\eta K\right)^3
 \quad\hbox{is real-rooted}
 \tag{N10.6}
\]
if and only if
\[
 \boxed{\eta\le\alpha_3.}
 \tag{N10.7}
\]
Equivalently, three conjugate pairs having a common center and common
squared height \(y^2\) can be realified at heat time \(K\) precisely when
\[
 \boxed{y^2\le\alpha_3K.}
 \tag{N10.8}
\]
\end{theorem}

\begin{proof}
Translation removes the common center and scaling reduces to \(K=1\).
Writing \(X=z^2\), direct application of the finite heat series gives
\[
\begin{aligned}
 e^{-D_z^2/2}(z^2+\eta)^3
 &=H_\eta(z^2),\\
 H_\eta(X)
 &=X^3+(3\eta-15)X^2
   +(3\eta^2-18\eta+45)X\\
 &\qquad+\eta^3-3\eta^2+9\eta-15.
\end{aligned}
\tag{N10.9}
\]
The exact cubic discriminant is
\[
 \operatorname{disc}_X H_\eta
 =864(75-90\eta+42\eta^2-8\eta^3)
 =-864F(\eta).
 \tag{N10.10}
\]

Suppose \(0\le\eta\le\alpha_3\).  Then the discriminant is nonnegative.
Moreover,
\[
 \frac13H_\eta'(X)
 =X^2+2(\eta-5)X+\eta^2-6\eta+15
 \tag{N10.11}
\]
has its two zeros at
\[
 X_\pm=5-\eta\pm\sqrt{10-4\eta}.
 \tag{N10.12}
\]
They are both positive.  Indeed \(5-\eta>0\), and
\[
 (5-\eta)^2-(10-4\eta)
 =\eta^2-6\eta+15=(\eta-3)^2+6>0.
 \tag{N10.13}
\]
The constant term
\[
 c(\eta):=\eta^3-3\eta^2+9\eta-15
\]
is strictly increasing because
\[
 c'(\eta)=3\{(\eta-1)^2+2\}>0,
\]
and \(c(\eta)<c(2)=-1\).  Hence \(H_\eta\) is increasing and negative
on \((-\infty,0]\), so it has no nonpositive zero.  Its nonnegative
discriminant says that all three zeros are real, including the double-root
case at \(\eta=\alpha_3\); therefore all three are positive.  It follows
that \(H_\eta(z^2)\) is real-rooted.

Conversely, if \(\eta>\alpha_3\), strict monotonicity of \(F\) makes the
discriminant in \textup{(N10.10)} negative.  Thus \(H_\eta\) has one real
zero and a nonreal conjugate pair.  The polynomial \(H_\eta(z^2)\) cannot
then be real-rooted.  This proves the equivalence.
\end{proof}

The V9 witness uses \(\eta=9/5\), and
\[
 \frac95<\alpha_3.
 \tag{N10.14}
\]
Thus it lies strictly inside the exact realifiable interval rather than at
its boundary.

\subsection{Gamma-ladder consequence}

\begin{corollary}[Centered equal three-channel gamma barrier]
\label{cor:f2v10-gamma}
Consider a critical fixed block made from three conjugate-pair channels
with equal ordinate magnitude \(\abs{\gamma}\), equal normalized height,
and a common center.  Let its total ladder count be \(R\ge6\).  If its heat
envelope is real-rooted, then
\[
 \boxed{
 \kappa^2\ge
 \frac{\gamma^2R}{\alpha_3(1+\theta)}
 \ge
 \frac{6}{\alpha_3}\frac{\gamma^2}{1+\theta}.}
 \tag{N10.15}
\]
Numerically,
\[
 \frac6{\alpha_3}
 =3.0872786691909941029\ldots .
 \tag{N10.16}
\]
\end{corollary}

\begin{proof}
The dimensionless squared height in Theorem~\ref{thm:f2v10-triple} is
\[
 \eta=\frac{y^2}{K}
 =\frac{\gamma^2R}{\kappa^2(1+\theta)}.
 \tag{N10.17}
\]
Condition \(\eta\le\alpha_3\) gives the first inequality in
\textup{(N10.15)}, and \(R\ge6\) gives the second.
\end{proof}

The coefficient \(6/\alpha_3\) is larger than both the general pure-block
coefficient \(2\) from V9 and the equal two-pair coefficient \(8/3\) from
V8.  Within this symmetric centered architecture, adding a third equal
channel makes the necessary gamma scale worse rather than better.

\begin{firewall}[Scope of the exact cubic discriminant]
Theorem~\ref{thm:f2v10-triple} is sharp only on the common-center,
equal-height three-pair diagonal.  It does not prove that
\(\alpha_3K\) bounds the shallowest pair when the three heights or centers
differ.  It also does not address real padding, growing block size, or a
target that is not least in its block.  The companion sweep imports no
missing cofinal theorem.  GRH remains open.
\end{firewall}

\begin{assessment}[V10 advance and V11 direction]
The mandatory sweep is complete at the V77/V94/V26 companion frontier and
changes no GRH premise.  The mathematical advance is an exact phase
boundary for the centered equal three-pair family.  It both validates the
V9 sextic witness and shows that the general \(3K\) first-collision ceiling
has substantial slack on the symmetric diagonal.

V11 should optimize the centered ordered family
\[
 \prod_{j=1}^3(z^2+t_jK),
 \qquad 0<t_1\le t_2\le t_3,
 \tag{N10.18}
\]
using the exact discriminant and positivity conditions of its cubic in
\(z^2\).  The decisive question is whether realifiability forces
\(t_1\le\alpha_3\), with equality only at
\(t_1=t_2=t_3=\alpha_3\).  A positive answer would extend the sharp V8
ordered-height phenomenon to three centered pairs before the centers are
released.
\end{assessment}

\section*{Version V10 append-only record}
\addcontentsline{toc}{section}{Version V10 append-only record}
The complete V9 source was retained before this continuation.  It had
\(91130\) bytes, \(2936\) lines, and SHA--256
\[
 \texttt{7D35E97F009514DE485B492C0F1182A2C90C73CDDDD0022EF573FB182C653034}.
\]
V10 records the mandatory companion sweep and proves the exact
common-center, equal-height three-pair threshold
\(y^2\le\alpha_3K\), where
\(8\alpha_3^3-42\alpha_3^2+90\alpha_3-75=0\).  Its gamma translation has
coefficient \(6/\alpha_3=3.087278669\ldots\).

% =============================================================================
% END APPENDED V10 MATERIAL
% =============================================================================


% =============================================================================
% BEGIN APPENDED V11 MATERIAL
% =============================================================================

\section{V11: the sharp ordered three-pair theorem at a common center}
\label{sec:grh-v11}

V10 found the exact realification threshold on the equal-height diagonal.
This section proves that unequal heights cannot improve it when the three
conjugate pairs have a common center.  The proof does not appeal to
symmetry as an extremal principle.  It converts the full cubic
discriminant into the sum and product of the two height increments and
certifies the required signs in Bernstein form.

\subsection{Height-increment coordinates}

Scale the heat time to \(K=1\), translate the common center to zero, and
write the ordered squared heights as
\[
 t,\qquad t+u,\qquad t+v,
 \qquad 0<t,\quad 0\le u\le v.
 \tag{N11.1}
\]
Set
\[
 S:=u+v,\qquad P:=uv.
 \tag{N11.2}
\]
The input polynomial is
\[
 \prod_{a\in\{t,t+u,t+v\}}(z^2+a)
 =z^6+e_1z^4+e_2z^2+e_3,
 \tag{N11.3}
\]
where
\[
 e_1=3t+S,\qquad
 e_2=3t^2+2tS+P,\qquad
 e_3=t^3+t^2S+tP.
 \tag{N11.4}
\]
Its unit backward-heat image is \(H_{t,S,P}(z^2)\), with
\[
 H_{t,S,P}(X)
 =X^3+(e_1-15)X^2+(e_2-6e_1+45)X
   +(e_3-e_2+3e_1-15).
 \tag{N11.5}
\]
If this image is real-rooted, the three zeros of \(H_{t,S,P}\) are
nonnegative.  Their sum therefore gives the necessary constraint
\[
 15-e_1=15-3t-S\ge0.
 \tag{N11.6}
\]
In particular,
\[
 t\le5,\qquad 0\le S\le L_t:=15-3t.
 \tag{N11.7}
\]

\subsection{A finite discriminant sign ledger}

Let
\[
 \Delta(t,S,P):=\operatorname{disc}_X H_{t,S,P}.
 \tag{N11.8}
\]
On the equal-height boundary, V10 gives
\[
 \Delta(t,0,0)=-864F(t),
 \qquad
 F(t)=8t^3-42t^2+90t-75.
 \tag{N11.9}
\]
The next lemma is the exact off-diagonal comparison.

\begin{lemma}[Discriminant decreases away from the ordered diagonal]
\label{lem:f2v11-discriminant}
Let
\[
 \alpha_3\le t\le5,\qquad
 0\le S\le15-3t,\qquad P\ge0.
 \tag{N11.10}
\]
Then
\[
 \boxed{\Delta(t,S,P)\le\Delta(t,0,0),}
 \tag{N11.11}
\]
and equality is possible only when \(S=P=0\).
\end{lemma}

\begin{proof}
Direct expansion of the cubic discriminant gives
\[
 \frac{\partial\Delta}{\partial P}
 =-12P^2+2P\{S^2+24S-144t-72\}+B_t(S),
 \tag{N11.12}
\]
where
\[
 B_t(S)
 =-8S^3+96tS^2+(432-720t)S
  -1728t^2+3456t-2160.
 \tag{N11.13}
\]
The \(P=0\) boundary satisfies
\[
 \Delta(t,S,0)-\Delta(t,0,0)=S A_t(S),
 \tag{N11.14}
\]
where
\[
\begin{aligned}
 A_t(S)={}&(24-16t)S^3+(288t-480)S^2\\
 &+(576t^2-3456t+4752)S\\
 &-6912t^2+24192t-25920.
\end{aligned}
\tag{N11.15}
\]

It remains to prove that \(A_t\) and \(B_t\) are negative on the interval
in \textup{(N11.10)}.  Since
\(\alpha_3>19/10\), it is enough to work on
\[
 I=[19/10,5].
 \tag{N11.16}
\]
For \(t<5\), put \(x=S/L_t\).  In the degree-three Bernstein basis
\[
 \mathcal B_{k,3}(x)
 =\binom3k x^k(1-x)^{3-k},
 \tag{N11.17}
\]
the four coefficients of \(A_t(L_tx)\) are
\[
\begin{aligned}
 a_0&=-1728(4t^2-14t+15),\\
 a_1&=-144(4t^3+4t^2-15t+15),\\
 a_2&=-288(t^3+15t^2-56t+50),\\
 a_3&=216(t-1)(2t^3-27t^2+84t-85).
\end{aligned}
\tag{N11.18}
\]
All four are strictly negative on \(I\).  For \(a_0\), the quadratic
inside the parentheses has negative discriminant and positive leading
coefficient.  For \(a_1\), the quadratic
\(4t^2-15t+15\) has negative discriminant, and \(4t^3>0\).
The polynomial in \(a_2\) is increasing on \(I\) and has value
\(4609/1000>0\) at \(19/10\).  Finally,
\[
 \frac d{dt}(2t^3-27t^2+84t-85)=6(t-2)(t-7);
\]
its maximum on \(I\) occurs at \(t=2\) and equals \(-9\).
Thus \(A_t(S)<0\).

The Bernstein coefficients of \(B_t(L_tx)\) are
\[
\begin{aligned}
 b_0&=-432(4t^2-8t+5),\\
 b_1&=-144t(7t+4),\\
 b_2&=144(2t^3-22t^2+18t+15),\\
 b_3&=216(5t^3-53t^2+135t-105).
\end{aligned}
\tag{N11.19}
\]
Again they are strictly negative.  The first two signs are immediate.
The cubic inside \(b_2\) is decreasing on \(I\) and has value
\(-8251/500\) at \(19/10\): its derivative is
\(2(3t^2-22t+9)\), whose convex quadratic factor is negative at both
endpoints of \(I\).  The cubic inside \(b_3\) is also decreasing on
\(I\), because its derivative is
\[
 15(t-5/3)(t-27/5)<0,
\]
and its value at \(19/10\) is \(-1107/200\).
Hence \(B_t(S)<0\).

The coefficient multiplying \(2P\) in \textup{(N11.12)} is negative as
well.  Since \(S^2+24S\) is increasing for \(S\ge0\),
\[
\begin{aligned}
 S^2+24S-144t-72
 &\le L_t^2+24L_t-144t-72\\
 &=9(t^2-34t+57)<0.
\end{aligned}
\tag{N11.20}
\]
The last quadratic is decreasing on \(I\) and equals
\(-399/100\) at its left endpoint.  It follows from
\textup{(N11.12)} that
\[
 \frac{\partial\Delta}{\partial P}<0.
 \tag{N11.21}
\]
Therefore
\[
 \Delta(t,S,P)\le\Delta(t,S,0)
 \le\Delta(t,0,0),
 \tag{N11.22}
\]
where the second inequality follows from
\textup{(N11.14)}--\textup{(N11.15)}.  It is strict if \(P>0\) or
\(S>0\).  When \(t=5\), condition \textup{(N11.10)} forces \(S=0\),
and the same conclusion follows directly.  This proves the lemma.
\end{proof}

\subsection{The ordered-height theorem}

\begin{theorem}[Sharp co-centered ordered three-pair obstruction]
\label{thm:f2v11-ordered}
Let
\[
 P_0(z)=
 \prod_{j=1}^3\bigl((z-c)^2+y_j^2\bigr),
 \qquad
 0<y_1\le y_2\le y_3.
 \tag{N11.23}
\]
If \(e^{-(K/2)D_z^2}P_0\) is real-rooted for some \(K>0\), then
\[
 \boxed{y_1^2\le\alpha_3K.}
 \tag{N11.24}
\]
Equality holds if and only if
\[
 y_1^2=y_2^2=y_3^2=\alpha_3K.
 \tag{N11.25}
\]
Here \(\alpha_3\) is the unique real zero from
Lemma~\ref{lem:f2v10-alpha}.
\end{theorem}

\begin{proof}
Translate by \(c\), scale by \(\sqrt K\), and use
\textup{(N11.1)} with
\[
 t=\frac{y_1^2}{K},\qquad
 u=\frac{y_2^2-y_1^2}{K},\qquad
 v=\frac{y_3^2-y_1^2}{K}.
 \tag{N11.26}
\]
Suppose first that \(t>\alpha_3\).  If \(t>5\), the output squared-root
sum in \textup{(N11.6)} is negative, which is impossible.  If
\(\alpha_3<t\le5\), Lemma~\ref{lem:f2v11-discriminant} and strict
monotonicity of \(F\) give
\[
 \Delta(t,S,P)
 \le-864F(t)<0.
 \tag{N11.27}
\]
The cubic \(H_{t,S,P}\) then has a nonreal conjugate pair, again
contradicting real-rootedness of \(H_{t,S,P}(z^2)\).  Thus
\(t\le\alpha_3\).

If \(t=\alpha_3\), then \(\Delta(t,0,0)=0\).
Lemma~\ref{lem:f2v11-discriminant} makes the discriminant negative unless
\(S=P=0\).  Since \(u,v\ge0\), this is equivalent to \(u=v=0\), proving
necessity in \textup{(N11.25)}.  Conversely, V10 proves that
\((z^2+\alpha_3)^3\) has a real-rooted unit backward-heat image.  Undoing
the translation and scaling proves the equality statement.
\end{proof}

\begin{corollary}[Least-ordinate centered three-channel gamma obstruction]
\label{cor:f2v11-gamma}
Consider a critical fixed gamma block made from three co-centered
conjugate-pair channels and no real padding.  Suppose its target pair has
least ordinate magnitude \(\abs{\gamma}\) in the block, and let the total
ladder count be \(R\ge6\).  Real-rootedness of the heat envelope requires
\[
 \boxed{
 \kappa^2\ge
 \frac{\gamma^2R}{\alpha_3(1+\theta)}
 \ge
 \frac6{\alpha_3}\frac{\gamma^2}{1+\theta}.}
 \tag{N11.28}
\]
Equality in the first heat obstruction requires all three normalized
heights to coincide.
\end{corollary}

\begin{proof}
The target squared height and heat time are
\[
 y_1^2=\frac{\gamma^2R}{\theta^2\kappa^2},
 \qquad
 K=\frac{1+\theta}{\theta^2}.
 \tag{N11.29}
\]
Insert these in Theorem~\ref{thm:f2v11-ordered} and use \(R\ge6\).
The equality statement follows from \textup{(N11.25)}.
\end{proof}

\begin{firewall}[The common center is still essential]
Theorem~\ref{thm:f2v11-ordered} allows arbitrary ordered heights but uses
the common center to reduce the heat image to a cubic in \(z^2\).  It does
not show that separated centers obey the same constant, does not include
real padding, and does not control a designated pair above a lower
companion.  The V9 bound \(y_{\min}^2\le3K\) remains the unconditional
three-pair statement.  GRH remains open.
\end{firewall}

\begin{assessment}[V11 advance and V12 direction]
The equality point found in V10 is now proved to be the unique maximizer of
the shallowest squared height throughout the co-centered ordered
three-pair family.  Unequal higher channels cannot help a least-ordinate
target, and the gamma coefficient \(6/\alpha_3\) therefore applies without
an equal-ordinate assumption.

V12 should release the common-center condition.  After translating the
mean center to zero, the sixth-degree heat image is no longer even, so the
cubic discriminant is unavailable.  The next attack should combine the V8
second/fourth-moment method with the V9 first-collision flow, or derive a
sixth-order centered moment inequality, to decide whether center
separation can push the ordered three-pair constant above \(\alpha_3\).
Any counterexample must be recorded before attempting a symbolic global
proof.
\end{assessment}

\section*{Version V11 append-only record}
\addcontentsline{toc}{section}{Version V11 append-only record}
The complete V10 source was retained before this continuation.  It had
\(100289\) bytes, \(3207\) lines, and SHA--256
\[
 \texttt{3FAF8A4CF9A3A294245B7C80A7EEE8477F7CC997C926B575A24C30D20C8A5952}.
\]
V11 proves the sharp co-centered ordered three-pair theorem
\(y_{\min}^2\le\alpha_3K\), with equality only for three identical heights,
and extends the V10 gamma coefficient \(6/\alpha_3\) to arbitrary
equal-or-higher co-centered companion channels.

% =============================================================================
% END APPENDED V11 MATERIAL
% =============================================================================


% =============================================================================
% BEGIN APPENDED V12 MATERIAL
% =============================================================================

\section{V12: symmetric center separation and the two-invariant reduction}
\label{sec:grh-v12}

V11 proved that unequal heights do not improve the three-pair threshold
when all centers coincide.  This section moves in the orthogonal direction:
the three heights are equal, while the centers are separated.  The
mirror-symmetric slice admits a complete discriminant comparison, and the
general equal-height problem reduces to two real center invariants.

\subsection{The mirror-separated heat cubic}

Translate the middle center to zero, scale the heat time to \(K=1\), and
put
\[
 a=d^2\ge0,\qquad t=y^2>0.
 \tag{N12.1}
\]
The three centers are \(-d,0,d\), and the equal-height input is
\[
\begin{aligned}
 P_{t,a}(z)
 &=(z^2+t)\bigl((z-d)^2+t\bigr)\bigl((z+d)^2+t\bigr)\\
 &=z^6+(3t-2a)z^4+(3t^2+a^2)z^2+t(t+a)^2.
\end{aligned}
\tag{N12.2}
\]
Its unit backward-heat image is \(H_{t,a}(z^2)\), where
\[
\begin{aligned}
 H_{t,a}(X)
 ={}&X^3+(3t-2a-15)X^2\\
 &+(3t^2+a^2-18t+12a+45)X\\
 &+t(t+a)^2-3t^2-a^2+9t-6a-15.
\end{aligned}
\tag{N12.3}
\]

\begin{lemma}[Exact center-separation discriminant debit]
\label{lem:f2v12-debit}
Let
\[
 \Delta(t,a):=\operatorname{disc}_X H_{t,a}.
\]
Then
\[
 \boxed{
 \Delta(t,a)=\Delta(t,0)-8a\,\mathcal R_t(a),}
 \tag{N12.4}
\]
where
\[
 \Delta(t,0)=-864F(t)
 \tag{N12.5}
\]
and
\[
\begin{aligned}
 \mathcal R_t(a)={}&2(t-1)a^4
 +(16t^2+18t-39)a^3\\
 &+(32t^3+48t^2+180t-372)a^2\\
 &+(288t^3-648t^2+1728t-2106)a\\
 &+864t^3-3456t^2+7128t-6480.
\end{aligned}
\tag{N12.6}
\]
For every \(t\ge\alpha_3\) and \(a\ge0\),
\[
 \mathcal R_t(a)>0.
 \tag{N12.7}
\]
\end{lemma}

\begin{proof}
Equations \textup{(N12.4)}--\textup{(N12.6)} follow by substituting
the coefficients in \textup{(N12.3)} into the cubic discriminant formula.
It remains to certify the sign without a numerical root calculation.

Since \(\alpha_3>19/10\), consider \(t\ge19/10\).  The five coefficients
of \(\mathcal R_t(a)\), in descending powers of \(a\), have values at
\(t=19/10\)
\[
 \frac95,\quad
 \frac{1324}{25},\quad
 \frac{45346}{125},\quad
 \frac{101664}{125},\quad
 \frac{64152}{125},
 \tag{N12.8}
\]
all strictly positive.  Each coefficient is increasing on this half-line.  For the first three,
the derivatives are respectively
\[
 2,\qquad 32t+18,\qquad 96t^2+96t+180,
\]
and are strictly positive.  For the coefficient of \(a\), its derivative
is
\[
 432(2t^2-3t+4)>0,
\]
and the constant coefficient has derivative
\[
 216(12t^2-32t+33)>0.
\]
Both quadratics are everywhere positive because their discriminants are
negative.  Thus every coefficient in \textup{(N12.6)} is positive for
\(t\ge\alpha_3\), proving \textup{(N12.7)}.
\end{proof}

\begin{theorem}[Symmetric center separation cannot improve the threshold]
\label{thm:f2v12-separated}
Let
\[
 P_0(z)=
 \prod_{\varepsilon\in\{-1,0,1\}}
 \bigl((z-c-\varepsilon d)^2+y^2\bigr),
 \qquad y>0.
 \tag{N12.9}
\]
If \(e^{-(K/2)D_z^2}P_0\) is real-rooted, then
\[
 \boxed{y^2\le\alpha_3K.}
 \tag{N12.10}
\]
Equality is possible if and only if
\[
 d=0,\qquad y^2=\alpha_3K.
 \tag{N12.11}
\]
In particular, every genuine symmetric separation \(d\ne0\) forces the
strict inequality \(y^2<\alpha_3K\).
\end{theorem}

\begin{proof}
Translation by \(c\) and scaling by \(\sqrt K\) reduce to
\textup{(N12.2)} with
\[
 t=\frac{y^2}{K},\qquad a=\frac{d^2}{K}.
 \tag{N12.12}
\]
If the heat image is real-rooted, all three zeros of \(H_{t,a}\) are
nonnegative, so its cubic discriminant is nonnegative.

Suppose \(t\ge\alpha_3\).  Lemma~\ref{lem:f2v10-alpha} gives
\(F(t)\ge0\).  Lemma~\ref{lem:f2v12-debit} therefore yields
\[
 \Delta(t,a)
 =-864F(t)-8a\mathcal R_t(a)\le0.
 \tag{N12.13}
\]
It is strictly negative if \(t>\alpha_3\) or \(a>0\), contradicting the
necessary discriminant condition.  The sole remaining case is
\(t=\alpha_3,a=0\), and V10 proves that its heat image is real-rooted.
Undoing the normalization proves the theorem.
\end{proof}

\subsection{The local quantitative penalty}

The discriminant debit also measures the first loss caused by a small
center separation.

\begin{proposition}[First-order loss in squared separation]
\label{prop:f2v12-local}
There is \(\varepsilon>0\) and a real-analytic boundary
\(t_\star(a)\) for \(0\le a<\varepsilon\) such that the mirror family is
real-rooted locally precisely on the lower side of this boundary, and
\[
 \boxed{
 t_\star(a)=\alpha_3-\beta_3a+O(a^2),}
 \tag{N12.14}
\]
where
\[
\begin{aligned}
 \beta_3
 &=\frac{10\alpha_3^2-24\alpha_3+15}
 {24\alpha_3^2-84\alpha_3+90}\\
 &=0.3521802809838056\ldots>0.
\end{aligned}
\tag{N12.15}
\]
In unscaled variables this reads
\[
 y_\star^2
 =\alpha_3K-\beta_3d^2
 +O\left(\frac{d^4}{K}\right).
 \tag{N12.16}
\]
\end{proposition}

\begin{proof}
At \((t,a)=(\alpha_3,0)\), the cubic \(H_{t,a}\) has one simple positive
zero and one double positive zero.  Moreover,
\[
 \partial_t\Delta(\alpha_3,0)
 =-864F'(\alpha_3)\ne0
\]
by Lemma~\ref{lem:f2v10-alpha}.  The implicit-function theorem gives a
unique analytic discriminant-zero curve through this point.  Positivity
of the three cubic zeros on its real-rooted side persists by continuity,
so this curve is the local real-rootedness boundary.

Differentiate \textup{(N12.4)} along the curve at \(a=0\):
\[
 -864F'(\alpha_3)t_\star'(0)
 -8\mathcal R_{\alpha_3}(0)=0.
 \tag{N12.17}
\]
Using \(F(\alpha_3)=0\) reduces
\[
 \mathcal R_{\alpha_3}(0)
 =108(10\alpha_3^2-24\alpha_3+15).
 \tag{N12.18}
\]
Therefore \(t_\star'(0)=-\beta_3\), proving
\textup{(N12.14)}--\textup{(N12.16)}.
\end{proof}

\subsection{Two invariants for arbitrary equal-height centers}

The mirror slice is only one edge of the full center problem, but the
following reduction keeps that problem finite.

\begin{proposition}[Exact two-invariant center reduction]
\label{prop:f2v12-invariants}
Let \(c_1+c_2+c_3=0\), and define
\[
 A:=-\sum_{i<j}c_ic_j
   =\frac12\sum_{j=1}^3c_j^2\ge0,
 \qquad
 B:=c_1c_2c_3.
 \tag{N12.19}
\]
The centers are real precisely when
\[
 27B^2\le4A^3.
 \tag{N12.20}
\]
For three equal squared heights \(t\), put
\[
 R(w):=\prod_{j=1}^3(w-c_j)=w^3-Aw-B.
\]
Then
\[
 \prod_{j=1}^3\bigl((z-c_j)^2+t\bigr)
 =R(z+i\sqrt t)R(z-i\sqrt t)
 \tag{N12.21}
\]
and its unit backward-heat image is
\[
\begin{aligned}
 Q_{A,B,t}(z)={}&z^6+(3t-2A-15)z^4-2Bz^3\\
 &+(A^2+3t^2-18t+12A+45)z^2\\
 &+2B(A+3t+3)z\\
 &+B^2+t(A+t)^2-A^2-3t^2+9t-6A-15.
\end{aligned}
\tag{N12.22}
\]
The mirror-separated family is exactly the slice \(B=0\), with
\(A=d^2\).
\end{proposition}

\begin{proof}
The discriminant of the center cubic \(w^3-Aw-B\) is
\(4A^3-27B^2\), proving \textup{(N12.20)}.  Pairing its values at
\(z\pm i\sqrt t\) gives \textup{(N12.21)}.  Expanding that product and
applying
\[
 e^{-D^2/2}=1-\frac12D^2+\frac18D^4-\frac1{48}D^6
\]
gives \textup{(N12.22)}.  If \(B=0\), the centered real roots of
\(w^3-Aw\) are \(-\sqrt A,0,\sqrt A\).
\end{proof}

\begin{corollary}[Mirror-separated equal-channel gamma obstruction]
\label{cor:f2v12-gamma}
For a critical three-channel gamma block with equal normalized heights and
centers in arithmetic progression, real-rootedness requires
\[
 \kappa^2\ge
 \frac{\gamma^2R}{\alpha_3(1+\theta)}
 \ge
 \frac6{\alpha_3}\frac{\gamma^2}{1+\theta}.
 \tag{N12.23}
\]
The first inequality is strict when the center spacing is nonzero.
\end{corollary}

\begin{proof}
Use
\[
 \frac{y^2}{K}
 =\frac{\gamma^2R}{\kappa^2(1+\theta)}
\]
in Theorem~\ref{thm:f2v12-separated}, and use \(R\ge6\).
\end{proof}

\begin{firewall}[The two center invariants are a reduction, not a solution]
Theorem~\ref{thm:f2v12-separated} treats equal heights and centers in an
arithmetic progression.  V11 treats arbitrary ordered heights at one
common center.  The two theorems do not combine automatically: arbitrary
unequal heights and arbitrary separated centers remain open.  In
Proposition~\ref{prop:f2v12-invariants}, nonzero \(B\) introduces the odd
terms of the sextic, and the cubic-in-\(z^2\) discriminant proof no longer
applies.  GRH remains open.
\end{firewall}

\begin{assessment}[V12 advance and V13 direction]
A genuine center-separation mode is now closed: mirror separation cannot
raise the V11 threshold and in fact incurs the explicit first-order debit
\(\beta_3d^2\).  The exact representation
\textup{(N12.22)} reduces arbitrary equal-height centers to the
semialgebraic cone
\[
 A\ge0,\qquad 27B^2\le4A^3.
 \tag{N12.24}
\]
V13 should work directly on this cone.  The first target is to determine
whether the Hermite matrix of \(Q_{A,B,t}\), rather than its enormous full
discriminant, has a principal minor or an optimized quadratic form that is
negative for every \(t\ge\alpha_3\) outside \((A,B)=(0,0)\).  Failure of
one fixed minor must be recorded rather than hidden; the alternative is a
boundary reduction in \(B^2/(4A^3/27)\).
\end{assessment}

\section*{Version V12 append-only record}
\addcontentsline{toc}{section}{Version V12 append-only record}
The complete V11 source was retained before this continuation.  It had
\(110204\) bytes, \(3543\) lines, and SHA--256
\[
 \texttt{40C609A46EED1AF4BD2A7571F6FBD27ADEC260258828B5DF8E116580F462B0C1}.
\]
V12 proves that symmetric center separation strictly lowers the
equal-height three-pair threshold, computes the local debit
\(\beta_3=0.3521802809\ldots\), and reduces arbitrary equal-height center
geometry to the exact two-invariant sextic \(Q_{A,B,t}\).

% =============================================================================
% END APPENDED V12 MATERIAL
% =============================================================================


% =============================================================================
% BEGIN APPENDED V13 MATERIAL
% =============================================================================

\section{V13: a Hermite-square certificate and local rigidity on the full center cone}

This continuation attacks the two-invariant cone left by V12.  The main
advance is local but genuinely three-directional: arbitrary centered real
parts, including nonzero skew invariant \(B\), cannot preserve the symmetric
threshold.  Moreover the first-order loss is governed by exactly the same
constant \(\beta _3\) found on the mirror slice.  We also record, with exact
counterexamples, why two tempting low-order Hermite minors cannot by
themselves settle the global cone.

Throughout this section the heat scale is \(K=1\), and
\(Q_{A,B,t}\) is the sextic in \textup{(N12.22)}.  Thus
\[
 A\geq0,\qquad 27B^2\leq4A^3,
 \tag{N13.1}
\]
and \(t\) is the common squared height.

\subsection{Newton sums and two certificates that do not close globally}

Let \(q_1,\ldots,q_6\) denote the roots of \(Q_{A,B,t}\), counted with
multiplicity, and put \(p_r=\sum_jq_j^r\), with \(p_0=6\).  If
\(Q=z^6+a_1z^5+\cdots+a_6\), Newton's identities give
\[
 p_r+a_1p_{r-1}+\cdots+a_{r-1}p_1+ra_r=0
 \quad(1\leq r\leq6),
 \tag{N13.2}
\]
with the usual homogeneous recurrence for \(r>6\).  In particular,
\[
\begin{aligned}
 p_2={}&4A-6t+30,\\
 p_4={}&4A^2-24At+72A+6t^2-108t+270,\\
 p_6={}&4A^3-60A^2t+132A^2+60At^2-648At+1116A\\
      &\quad+6B^2-6t^3+234t^2-1674t+2790,\\
 p_8={}&2\bigl(2A^4-56A^3t+104A^3+140A^2t^2-1080A^2t
                  +1476A^2\\
      &\qquad+8AB^2-56At^3+1272At^2-6408At+8040A\\
      &\qquad-84B^2t+132B^2+3t^4-204t^3+2826t^2
                  -12060t+15075\bigr).
\end{aligned}
\tag{N13.3}
\]

\begin{proposition}[Exact failure of two fixed low-order tests]
\label{prop:f2v13-minor-failure}
The fourth-moment Cauchy certificate is
\[
 6p_4-p_2^2
 =8\bigl(A^2-12At+24A-36t+90\bigr).
 \tag{N13.4}
\]
The even-moment Hermite minor
\[
 \mathcal H_{\rm ev}:=
 \det\!\begin{pmatrix}
 p_0&p_2&p_4\\ p_2&p_4&p_6\\ p_4&p_6&p_8
 \end{pmatrix}
 \tag{N13.5}
\]
also fails to be negative throughout the forbidden region \(t>\alpha _3\).
Indeed, at the exact cone point
\[
 t=2,\qquad A=200,\qquad B^2=\frac{16000000}{27}
       =\frac12\frac{4A^3}{27},
 \tag{N13.6}
\]
one has
\[
 6p_4-p_2^2=320144,
 \qquad
 \mathcal H_{\rm ev}=\frac{1157808151774976}{27}>0.
 \tag{N13.7}
\]
\end{proposition}

\begin{proof}
Equation \textup{(N13.4)} follows by inserting the first two lines of
\textup{(N13.3)}.  Equations \textup{(N13.5)}--\textup{(N13.7)} follow by
inserting all four displayed Newton sums and performing integer arithmetic.
The cone condition in \textup{(N13.1)} is strict at \textup{(N13.6)}, and
\(2>\alpha _3\) by the rational enclosure below.
\end{proof}

\begin{firewall}[What the positive values do and do not say]
For a real-rooted polynomial, both quantities in Proposition
\ref{prop:f2v13-minor-failure} must be nonnegative.  Their positivity at
\textup{(N13.6)} does not prove that the sextic there is real-rooted.  It
proves only that neither one of these necessary inequalities, used alone,
can exclude the whole invariant cone.  The global problem therefore needs
either a higher or an adapted Hermite direction.
\end{firewall}

\subsection{The adapted square at the symmetric boundary}

Write
\[
 \alpha:=\alpha _3,
 \qquad s:=\sqrt{10-4\alpha},
 \qquad
 \lambda _1:=5-\alpha+s,
 \qquad
 \lambda _0:=5-\alpha-2s.
 \tag{N13.8}
\]
The threshold factorization from V10 is
\[
 Q_{0,0,\alpha}(z)
 =(z^2-\lambda _0)(z^2-\lambda _1)^2,
 \qquad \lambda _0,\lambda _1>0.
 \tag{N13.9}
\]
Put
\[
 h(x):=(x^2-\lambda _0)(x^2-\lambda _1)
      =x^4-cx^2+d,
 \quad c:=\lambda _0+\lambda _1,
 \quad d:=\lambda _0\lambda _1,
 \tag{N13.10}
\]
and define the adapted Hermite quadratic form
\[
 \mathcal S(A,B,t)
 :=p_8-2cp_6+(c^2+2d)p_4-2cdp_2+6d^2.
 \tag{N13.11}
\]

\begin{lemma}[Threshold Taylor ledger]
\label{lem:f2v13-taylor}
If \(Q_{A,B,t}\) is real-rooted, then
\[
 \mathcal S(A,B,t)=\sum_{j=1}^6h(q_j)^2\geq0.
 \tag{N13.12}
\]
At \((A,B,t)=(0,0,\alpha)\), \(\mathcal S=0\).  Writing
\(\tau=t-\alpha\), its first-order expansion is
\[
 \mathcal S(A,B,\alpha+\tau)
 =c_AA+c_t\tau+c_BB^2
  +O\!\left((A+|\tau|)^2+(A+|\tau|)B^2\right),
 \tag{N13.13}
\]
where
\[
\begin{aligned}
 c_A&=32(6\alpha^2-10\alpha s-39\alpha+16s+50),\\
 c_t&=96(2\alpha s+12\alpha-8s-25),\\
 c_B&=-12(12\alpha-s-12).
\end{aligned}
\tag{N13.14}
\]
All three coefficients are strictly negative, and
\[
 \frac{c_A}{c_t}
 =\frac{10\alpha^2-24\alpha+15}
        {24\alpha^2-84\alpha+90}
 =\beta _3.
 \tag{N13.15}
\]
\end{lemma}

\begin{proof}
Expanding \(\sum_j(q_j^4-cq_j^2+d)^2\) gives
\textup{(N13.11)}, and proves \textup{(N13.12)} when the \(q_j\) are real.
Factorization \textup{(N13.9)} shows that \(h\) vanishes on every root at
the threshold, so \(\mathcal S(0,0,\alpha)=0\).

Substitution of \textup{(N13.3)} into \textup{(N13.11)}, followed by
differentiation at the threshold, gives \textup{(N13.13)} and
\textup{(N13.14)}.  This calculation uses only polynomial expansion; in
particular \(B\) occurs through \(B^2\).

For an entirely rational sign check, \(F'(x)=24x^2-84x+90>0\) on
\(\mathbb R\), and
\[
 F\!\left(\frac{19434}{10000}\right)
 =-\frac{16081437}{15625000000}<0,
 \quad
 F\!\left(\frac{19435}{10000}\right)
 =\frac{710603}{1000000000}>0.
 \tag{N13.16}
\]
Consequently
\[
 1.9434<\alpha<1.9435,
 \qquad 1.4919<s<1.4922.
 \tag{N13.17}
\]
Direct outward rational interval arithmetic in \textup{(N13.14)} gives
\[
\begin{aligned}
 -264.518&<c_A<-263.930,\\
 -750.536&<c_t<-750.049,\\
 -117.962&<c_B<-117.943,
\end{aligned}
\tag{N13.18}
\]
which proves the signs.  Finally, the threshold identity
\[
 s=\frac{4\alpha^2-20\alpha+25}{3\alpha-5}
 \tag{N13.19}
\]
follows from \textup{(N13.9)} (equivalently from
\(\mathcal S(0,0,\alpha)=0\)).  Substituting \textup{(N13.19)} into the
left side of \textup{(N13.15)} reduces the difference from its right side
to a multiple of
\(F(\alpha)=8\alpha^3-42\alpha^2+90\alpha-75=0\).
\end{proof}

\begin{theorem}[Local rigidity on the full equal-height center cone]
\label{thm:f2v13-local-rigidity}
There exists \(\varepsilon>0\) such that, whenever
\[
 0\leq A\leq\varepsilon,\qquad 27B^2\leq4A^3,\qquad t\geq\alpha _3,
 \tag{N13.20}
\]
the sextic \(Q_{A,B,t}\) can be real-rooted only at
\[
 (A,B,t)=(0,0,\alpha _3).
 \tag{N13.21}
\]
At that point it is real-rooted.
\end{theorem}

\begin{proof}
First take \(t=\alpha+\tau\) with \(\tau\geq0\) and
\(A+\tau\) small.  The cone bound gives \(B^2=O(A^3)\).  By
Lemma~\ref{lem:f2v13-taylor},
\[
 \mathcal S=c_AA+c_t\tau+o(A+\tau).
 \tag{N13.22}
\]
Because \(c_A,c_t<0\), this is negative whenever \(A+\tau>0\) is
sufficiently small, contradicting the necessary condition
\(\mathcal S\geq0\).

It remains to make the exclusion uniform away from the threshold.  The V9
three-pair height bound gives \(t\leq3\) for every real-rooted member of
this family.  Fix a small \(\delta>0\).  By V10,
\(Q_{0,0,t}\) is not real-rooted on the compact interval
\([\alpha+\delta,3]\).  The set of monic real-rooted sextics is closed in
coefficient space.  Hence the compact coefficient curve
\(\{Q_{0,0,t}:\alpha+\delta\leq t\leq3\}\) has positive distance from that
closed set.  Uniformly on this interval,
\(Q_{A,B,t}\to Q_{0,0,t}\) in coefficient space as \(A\to0\), because
\(|B|\leq2A^{3/2}/\sqrt{27}\).  Thus sufficiently small \(A\) also excludes
this remaining interval.  Combining the two exclusions proves
\textup{(N13.21)}.  Real-rootedness at the exceptional point follows from
\textup{(N13.9)}.
\end{proof}

\begin{corollary}[Universal first-order center-variance debit]
\label{cor:f2v13-debit}
Suppose \(A_n>0\), \(A_n\to0\), \(27B_n^2\leq4A_n^3\), and
\(Q_{A_n,B_n,t_n}\) is real-rooted with \(t_n\to\alpha _3\).  Then
\[
 \limsup_{n\to\infty}\frac{t_n-\alpha _3}{A_n}
 \leq-\beta _3.
 \tag{N13.23}
\]
Thus, to first order, arbitrary centered real-part variance incurs the same
debit as mirror separation; the skew invariant \(B\) cannot change the
linear term.
\end{corollary}

\begin{proof}
Apply \(\mathcal S\geq0\) and \textup{(N13.13)}.  Along any subsequence on
which \((t_n-\alpha)/A_n\) has a finite limit \(L\), division by \(A_n\)
uses \(B_n^2/A_n=O(A_n^2)\) and gives
\(0\leq c_A+c_tL\).  Since \(c_t<0\),
\(L\leq-c_A/c_t=-\beta _3\).  Subsequences on which the quotient tends to
\(-\infty\) already satisfy the assertion; Theorem
\ref{thm:f2v13-local-rigidity} excludes a positive unbounded quotient.
\end{proof}

\begin{firewall}[Local center rigidity is not GRH]
Theorem~\ref{thm:f2v13-local-rigidity} treats three equal heights and only
a neighborhood of coincident centers.  It does not handle large center
variance, unequal heights with separated centers, passage from finite
blocks to the full zero set, or the arithmetic hypotheses needed to force
a forbidden block.  Proposition~\ref{prop:f2v13-minor-failure} explicitly
shows that the first two fixed moment tests do not close the large-\(A\)
cone.  GRH remains open.
\end{firewall}

\begin{assessment}[V13 advance and V14 direction]
The symmetric threshold is now isolated inside the full local
semialgebraic center cone, and the coefficient \(\beta _3\) has been
identified as a universal first variation in \(A\), not merely a mirror
slice constant.  The next attack should go upstream in the center cubic.
Write
\[
 \rho:=\frac{27B^2}{4A^3}\in[0,1]
 \tag{N13.24}
\]
for \(A>0\), and seek a Hermite direction adapted continuously to the
three real roots of \(w^3-Aw-B\).  The endpoints \(\rho=0\) and \(\rho=1\)
encode the mirror and double-center geometries.  A useful V14 target is a
boundary/interpolation theorem in \(\rho\), rather than another fixed
principal minor.  If that stalls, the right upstream object is the
root-configuration parametrization of the center cubic, where the cone
constraint is automatic.
\end{assessment}

\section*{Version V13 append-only record}
\addcontentsline{toc}{section}{Version V13 append-only record}
The complete V12 source was retained before this continuation.  It had
\(120265\) bytes, \(3887\) lines, and SHA--256
\[
 \texttt{EDA28844EFE30D870DAA13F160CF7F97A5EDF7057ECD7447202CA46B83F45CAF}.
\]
V13 proves local rigidity for arbitrary equal-height centered real parts,
derives the universal first-order debit
\(t\leq\alpha _3-\beta _3A+o(A)\), and records exact failures of two
low-order global Hermite certificates.

% =============================================================================
% END APPENDED V13 MATERIAL
% =============================================================================


% =============================================================================
% BEGIN APPENDED V14 MATERIAL
% =============================================================================

\section{V14: the double-center boundary is globally below the triple threshold}

V13 suggested moving upstream to the center cubic and treating the boundary
parameter
\[
 \rho=\frac{27B^2}{4A^3}.
\]
This continuation closes the endpoint \(\rho=1\) for all heights allowed by
the V9 ceiling.  The proof combines the V13 adapted Hermite square at small
separation with the ordinary sextic discriminant elsewhere.  A finite
Bernstein ledger makes the covering argument exact.

\subsection{The boundary family and its two obstruction polynomials}

On \(27B^2=4A^3\), the center cubic has a repeated real root.  Up to
reflection and translation, write the three centers as
\[
 (-2d,d,d),\qquad d\geq0.
 \tag{N14.1}
\]
Put \(x=d^2\).  Then \(A=3x\), \(B=-2d^3\), and the unit-scale heat image is
\[
\begin{aligned}
 \mathcal Q_{d,t}(z)={}&z^6+(3t-6d^2-15)z^4+4d^3z^3\\
 &+(9d^4+3t^2-18t+36d^2+45)z^2\\
 &-12d^3(d^2+t+1)z\\
 &+4d^6+9d^4t-9d^4+6d^2t^2-18d^2
   +t^3-3t^2+9t-15.
\end{aligned}
\tag{N14.2}
\]
Changing the sign of \(B\) replaces \(z\) by \(-z\), so this choice loses no
generality.

Let \(\mathcal S(A,B,t)\) be the adapted Hermite square from
\textup{(N13.11)}, and define
\[
 \Sigma(x,t):=\mathcal S(3x,B,t)\big|_{B^2=4x^3},
 \qquad
 \mathcal P(x,t):=
 -\frac{\operatorname{Disc}_z(\mathcal Q_{d,t})}{2985984}
 \quad(x=d^2).
 \tag{N14.3}
\]
Both are polynomials in \(x,t\).  If \(\mathcal Q_{d,t}\) is real-rooted,
then necessarily
\[
 \Sigma(x,t)\geq0,\qquad \mathcal P(x,t)\leq0.
 \tag{N14.4}
\]
The first inequality is \textup{(N13.12)}.  For the second, a monic
real-rooted sextic has
\(\operatorname{Disc}=\prod_{i<j}(q_i-q_j)^2\geq0\).

\subsection{A finite Bernstein sign ledger}

For a polynomial
\[
 R(X,U)=\sum_{k=0}^m\sum_{\ell=0}^n a_{k\ell}X^kU^\ell,
\]
write it on the unit square in the tensor Bernstein basis
\[
 R(X,U)=\sum_{i=0}^m\sum_{j=0}^n
 b_{ij}\binom miX^i(1-X)^{m-i}
       \binom njU^j(1-U)^{n-j}.
 \tag{N14.5}
\]
Direct change of basis gives the auditable formula
\[
 b_{ij}=
 \sum_{k\leq i,\,\ell\leq j}
 \frac{\binom ik}{\binom mk}
 \frac{\binom j\ell}{\binom n\ell}a_{k\ell}.
 \tag{N14.6}
\]
Because the Bernstein basis functions are nonnegative and sum to one, the
range of \(R\) lies in the convex hull of its \(b_{ij}\).

\begin{lemma}[Three exact Bernstein certificates]
\label{lem:f2v14-bernstein}
Let \(\alpha=\alpha_3\), and let \(s=\sqrt{10-4\alpha}\).

\begin{enumerate}[label=\textup{(\roman*)}]
\item Put
\[
 \Sigma_{\rm lo}(X,U):=
 \Sigma\!\left(X,\alpha+\left(\frac94-\alpha\right)U\right).
 \tag{N14.7}
\]
It has bidegree \((4,4)\).  Its Bernstein coefficient \(b_{00}\) is zero,
and every other coefficient is strictly negative.  More precisely, the
following are rigorous upper bounds for the maximum coefficient in each
\(X\)-row; in row \(0\), \(b_{00}\) is omitted:
\[
\begin{array}{c|rrrrr}
 i&0&1&2&3&4\\ \hline
 \max_j b_{ij}
 &-57&-197&-759&-2116&-4187
\end{array}
\tag{N14.8}
\]

\item Put
\[
 \mathcal P_{\rm hi}(X,U):=
 \mathcal P\!\left(X,\frac94+\frac34U\right).
 \tag{N14.9}
\]
It has bidegree \((8,9)\), and all its Bernstein coefficients are positive.
The exact row minima are bounded below by
\[
\begin{gathered}
\begin{array}{c|rrrrr}
 i&0&1&2&3&4\\ \hline
 \min_j b_{ij}
 &837&\frac{228609}{64}&\frac{106276077}{7168}
 &\frac{427793103}{7168}&\frac{4152538899}{17920}
\end{array}\\[2mm]
\begin{array}{c|rrrr}
 i&5&6&7&8\\ \hline
 \min_j b_{ij}
 &\frac{6172780959}{7168}&\frac{21703298625}{7168}
 &\frac{5081917779}{512}&\frac{1931316885}{64}
\end{array}
\end{gathered}
\tag{N14.10}
\]

\item For the large-separation chart put
\[
 \mathcal P_\infty(Y,U):=
 Y^8\mathcal P\!\left(Y^{-1},\alpha+(3-\alpha)U\right).
 \tag{N14.11}
\]
The apparent singularity at \(Y=0\) is removable, and the result has
bidegree \((8,9)\).  Every Bernstein coefficient is positive.  Rigorous
lower bounds for its row minima are
\[
\begin{gathered}
\begin{array}{c|rrrrr}
 i&0&1&2&3&4\\ \hline
 \min_j b_{ij}&126000&215000&350000&542000&802000
\end{array}\\[2mm]
\begin{array}{c|rrrr}
 i&5&6&7&8\\ \hline
 \min_j b_{ij}&1145000&1580000&2117000&2761000
\end{array}
\end{gathered}
\tag{N14.12}
\]
\end{enumerate}
\end{lemma}

\begin{proof}
Insert the four Newton sums \textup{(N13.3)} into
\textup{(N13.11)}, then set \(A=3X\), \(B^2=4X^3\); this produces
\(\Sigma\).  Equation \textup{(N14.2)} and the resultant definition of the
discriminant produce \(\mathcal P\).  Apply the finite change-of-basis
formula \textup{(N14.6)}.

For completeness, the algebraic signs in the ledger are certified without
floating-point assumptions.  Reduce every occurrence of \(\alpha^r\) and
\(s^r\) using
\[
 8\alpha^3-42\alpha^2+90\alpha-75=0,
 \qquad s^2=10-4\alpha,
 \tag{N14.13}
\]
and use the rational enclosures
\[
 \frac{19434}{10000}<\alpha<\frac{19435}{10000},
 \qquad
 \frac{14919}{10000}<s<\frac{14922}{10000}.
 \tag{N14.14}
\]
Outward rational interval arithmetic gives, before rounding toward zero,
the row upper bounds
\[
 -57.389,\ -197.593,\ -759.145,\ -2116.989,\ -4187.813
 \tag{N14.15}
\]
in part \textup{(i)}, and the row lower bounds
\[
 126585,\ 215773,\ 350426,\ 542054,\ 802884,\
 1145181,\ 1580283,\ 2117336,\ 2761742
 \tag{N14.16}
\]
in part \textup{(iii)}.  The exceptional coefficient in part
\textup{(i)} is exactly
\(\Sigma(0,\alpha)=0\) by the factorization \textup{(N13.9)}.
All coefficients in part \textup{(ii)} are rational; direct exact
arithmetic gives \textup{(N14.10)}.  Thus
\textup{(N14.8)}, \textup{(N14.10)}, and \textup{(N14.12)} follow.
\end{proof}

\begin{theorem}[Global double-center threshold]
\label{thm:f2v14-double-center}
Let \(d\geq0\), and suppose the backward-heat image
\(\mathcal Q_{d,t}\) in \textup{(N14.2)} is real-rooted.  Then
\[
 t\leq\alpha_3.
 \tag{N14.17}
\]
Equality holds if and only if \(d=0\).  Equivalently, on the full
double-center boundary \(27B^2=4A^3\), every nonzero center separation
strictly lowers the equal-height three-pair threshold.
\end{theorem}

\begin{proof}
The V9 three-pair bound gives \(t\leq3\), so it suffices to exclude
\(\alpha\leq t\leq3\) away from \(d=0,t=\alpha\).  Put \(x=d^2\).

If \(0\leq x\leq1\) and
\(\alpha\leq t\leq9/4\), write
\(t=\alpha+(9/4-\alpha)U\).  Lemma
\ref{lem:f2v14-bernstein}\textup{(i)} and the convex-hull property imply
\(\Sigma(x,t)<0\), except at \(x=0,U=0\).  This contradicts
\textup{(N14.4)}.

If \(0\leq x\leq1\) and \(9/4\leq t\leq3\), part \textup{(ii)} gives
\(\mathcal P(x,t)>0\), again contradicting \textup{(N14.4)}.

Finally, if \(x\geq1\) and \(\alpha\leq t\leq3\), set \(Y=x^{-1}\) and
\(t=\alpha+(3-\alpha)U\).  Part \textup{(iii)} gives
\(Y^8\mathcal P(x,t)>0\), hence \(\mathcal P(x,t)>0\).  This is the same
discriminant contradiction.  The only point left is \(d=0,t=\alpha\),
which is real-rooted by \textup{(N13.9)}.
\end{proof}

\begin{corollary}[Strict gamma obstruction on the double-center boundary]
\label{cor:f2v14-gamma}
For a critical equal-height three-channel gamma block whose centered real
parts have the pattern \((-2d,d,d)\), real-rootedness requires
\[
 \kappa^2\geq
 \frac{\gamma^2R}{\alpha_3(1+\theta)}
 \geq\frac6{\alpha_3}\frac{\gamma^2}{1+\theta}.
 \tag{N14.18}
\]
The first inequality is strict when \(d\neq0\).
\end{corollary}

\begin{proof}
As in V12,
\(t=y^2/K=\gamma^2R/(\kappa^2(1+\theta))\).  Apply
Theorem~\ref{thm:f2v14-double-center} and \(R\geq6\).
\end{proof}

\begin{firewall}[A boundary theorem is not the interior cone]
Theorem~\ref{thm:f2v14-double-center} closes \(\rho=1\); V12 closes the
mirror slice \(\rho=0\); V13 closes a neighborhood of \(A=0\) uniformly
in \(\rho\).  None of these statements proves the threshold throughout
\(0<\rho<1\), and none treats arbitrary unequal heights or the infinite
arithmetic zero set.  GRH remains open.
\end{firewall}

\begin{assessment}[V14 advance and V15 direction]
Both geometric endpoints of the equal-height center cone now obey the
sharp symmetric threshold, with strict loss away from coincident centers.
The successful proof is not a single invariant inequality but a certified
cover by adapted Hermite and discriminant charts.  V15 should test whether
the same three-chart cover remains valid after introducing
\[
 B^2=\frac{4A^3}{27}\rho,\qquad 0\leq\rho\leq1.
 \tag{N14.19}
\]
The first target is a Bernstein sign ledger in the added variable \(\rho\).
If a coefficient changes sign, subdivide in \(\rho\) and record the first
genuine gap.  This is preferable to expanding the entire discriminant
without a geometric chart.
\end{assessment}

\section*{Version V14 append-only record}
\addcontentsline{toc}{section}{Version V14 append-only record}
The complete V13 source was retained before this continuation.  It had
\(131285\) bytes, \(4212\) lines, and SHA--256
\[
 \texttt{FCCB656F5E4FBE4AD2E2C7A4EFD57A3B1349A0E7F0C52FA0E0851200B605CA8D}.
\]
V14 proves the sharp global equal-height threshold on the full
double-center boundary \(27B^2=4A^3\), with equality only at coincident
centers, by an exact three-chart Bernstein certificate.

% =============================================================================
% END APPENDED V14 MATERIAL
% =============================================================================


% =============================================================================
% BEGIN APPENDED V15 MATERIAL
% =============================================================================

\section{V15: closure of the full arbitrary-center equal-height cone}

V14 closed the endpoint \(\rho=1\), while V12 had closed \(\rho=0\).
The suggested interior test is unexpectedly decisive: after retaining
\(\rho\) as a third Bernstein variable, every sign in the V14 three-chart
cover survives on the full interval \(0\leq\rho\leq1\).  This proves the
sharp equal-height three-pair threshold for completely arbitrary real
centers.

\subsection{Interior normalization and necessary signs}

Continue with the centered invariants \(A,B\) from V12.  For \(A>0\), set
\[
 x:=\frac A3,\qquad
 \rho:=\frac{27B^2}{4A^3}.
 \tag{N15.1}
\]
The real-center cone is exactly
\[
 x>0,\qquad 0\leq\rho\leq1,
 \qquad B^2=4\rho x^3.
 \tag{N15.2}
\]
Let
\[
 \widehat\Sigma(x,t,\rho):=
 \mathcal S(3x,B,t)\big|_{B^2=4\rho x^3},
 \tag{N15.3}
\]
where \(\mathcal S\) is the V13 adapted Hermite square.  Also put
\[
 \widehat{\mathcal P}(x,t,\rho):=
 -\frac{\operatorname{Disc}_z Q_{3x,B,t}}{2985984}
 \bigg|_{B^2=4\rho x^3}.
 \tag{N15.4}
\]
These are polynomials: reflection gives
\(Q_{A,-B,t}(z)=Q_{A,B,t}(-z)\), so both expressions are even in \(B\).
If \(Q_{3x,B,t}\) is real-rooted, then
\[
 \widehat\Sigma(x,t,\rho)\geq0,
 \qquad
 \widehat{\mathcal P}(x,t,\rho)\leq0.
 \tag{N15.5}
\]

For a trivariate polynomial of tridegree \((m,n,p)\), use the tensor
Bernstein expansion
\[
 R(X,U,V)=
 \sum_{i=0}^m\sum_{j=0}^n\sum_{k=0}^p
 b_{ijk}B_i^m(X)B_j^n(U)B_k^p(V),
 \quad
 B_i^m(X)=\binom miX^i(1-X)^{m-i}.
 \tag{N15.6}
\]
The exact power-to-Bernstein conversion is
\[
 b_{ijk}=
 \sum_{\substack{a\leq i\\b\leq j\\c\leq k}}
 \frac{\binom ia}{\binom ma}
 \frac{\binom jb}{\binom nb}
 \frac{\binom kc}{\binom pc}\,r_{abc},
 \tag{N15.7}
\]
where \(R=\sum r_{abc}X^aU^bV^c\).  Thus every coefficient below is
generated by one finite rational formula.

\begin{lemma}[Full-cone three-chart Bernstein ledger]
\label{lem:f2v15-full-ledger}
Let \(\alpha=\alpha_3\).

\begin{enumerate}[label=\textup{(\roman*)}]
\item Define
\[
 \widehat\Sigma_{\rm lo}(X,U,V):=
 \widehat\Sigma\!\left(
 X,\alpha+\left(\frac94-\alpha\right)U,V\right).
 \tag{N15.8}
\]
It has tridegree \((4,4,1)\).  Exactly the two coefficients
\(b_{000},b_{001}\) vanish, and all other \(48\) coefficients are strictly
negative.  Rigorous upper bounds for the maximum nonexceptional
coefficient in each \(X\)-row are
\[
\begin{array}{c|rrrrr}
 i&0&1&2&3&4\\ \hline
 \max b_{ijk}
 &-57&-197&-759&-1999&-3908.
\end{array}
\tag{N15.9}
\]

\item Define
\[
 \widehat{\mathcal P}_{\rm hi}(X,U,V):=
 \widehat{\mathcal P}\!\left(
 X,\frac94+\frac34U,V\right).
 \tag{N15.10}
\]
It has tridegree \((12,9,4)\).  All \(650\) Bernstein coefficients are
positive rational numbers, and
\[
 \min_{i,j,k}b_{ijk}=837.
 \tag{N15.11}
\]

\item In the large-separation chart define
\[
 \widehat{\mathcal P}_\infty(Y,U,V):=
 Y^{12}\widehat{\mathcal P}\!\left(
 Y^{-1},\alpha+(3-\alpha)U,V\right).
 \tag{N15.12}
\]
The singularity at \(Y=0\) is removable, and the polynomial has
tridegree \((12,9,4)\).  Exactly \(100\) Bernstein coefficients vanish:
\[
 b_{ijk}=0
 \quad\Longleftrightarrow\quad
 0\leq i\leq3\ \hbox{ and }\ k>i.
 \tag{N15.13}
\]
All other \(550\) coefficients are strictly positive.  The following
integers are rigorous lower bounds for the positive coefficients in each
\(Y\)-row:
\[
\begin{gathered}
\begin{array}{c|rrrrrrr}
 i&0&1&2&3&4&5&6\\ \hline
 \min{}' b_{ijk}
 &2446&1043&480&301&255&2179&10618
\end{array}\\[2mm]
\begin{array}{c|rrrrrr}
 i&7&8&9&10&11&12\\ \hline
 \min{}' b_{ijk}
 &38327&113539&291500&670423&1411557&2761742.
\end{array}
\end{gathered}
\tag{N15.14}
\]
Here the prime means that the structural zeros in
\textup{(N15.13)} are omitted.
\end{enumerate}
\end{lemma}

\begin{proof}
Insert \(A=3X\) and \(B^2=4VX^3\) into the explicit Newton sums
\textup{(N13.3)} and the square \textup{(N13.11)} to obtain part
\textup{(i)}.  Insert the same substitutions into the discriminant of the
explicit sextic \textup{(N12.22)} to obtain parts \textup{(ii)} and
\textup{(iii)}.  Apply \textup{(N15.7)}.

Here is a reproducible exact sign audit.  Reduce algebraic coefficients by
\[
 8\alpha^3-42\alpha^2+90\alpha-75=0,
 \qquad s^2=10-4\alpha,
 \tag{N15.15}
\]
and use the rational intervals
\[
 \frac{19434}{10000}<\alpha<\frac{19435}{10000},
 \qquad
 \frac{14919}{10000}<s<\frac{14922}{10000}.
 \tag{N15.16}
\]
Outward rational interval arithmetic gives the unrounded row upper bounds
\[
 -57.389,\ -197.593,\ -759.145,\ -1999.046,\ -3908.040
 \tag{N15.17}
\]
for part \textup{(i)}.  The two exceptional coefficients equal
\(\widehat\Sigma(0,\alpha,V)=0\) identically.  Part \textup{(ii)} contains
no algebraic coefficients after the rational height substitution; exact
application of \textup{(N15.7)} gives \(650\) positive rationals with
minimum \(837\).

For part \textup{(iii)}, reduction modulo the cubic in
\textup{(N15.15)} gives the \(100\) exact zeros in
\textup{(N15.13)}.  Outward rational interval arithmetic gives no
ambiguous coefficient and yields the unrounded positive row minima
\[
\begin{aligned}
 &(2446.328,\ 1043.762,\ 480.632,\ 301.670,\ 255.727,\
   2179.530,\ 10618.984,\\
 &\hspace{20mm}38327.109,\ 113539.235,\ 291500.791,\
   670423.283,\ 1411557.853,\ 2761742.716),
\end{aligned}
\tag{N15.18}
\]
which proves \textup{(N15.14)}.
\end{proof}

\begin{lemma}[Strictness of the three charts]
\label{lem:f2v15-strict}
On the unit cube:
\[
 \widehat\Sigma_{\rm lo}(X,U,V)<0
 \quad\text{unless }(X,U)=(0,0),
 \tag{N15.19}
\]
and
\[
 \widehat{\mathcal P}_{\rm hi}(X,U,V)>0.
 \tag{N15.20}
\]
Moreover
\[
 \widehat{\mathcal P}_\infty(Y,U,V)>0
 \qquad(0<Y\leq1).
 \tag{N15.21}
\]
\end{lemma}

\begin{proof}
The Bernstein basis is nonnegative and forms a partition of unity.
For \textup{(N15.19)}, the only zero coefficients have
\((i,j)=(0,0)\); at every other point some strictly negative coefficient
has positive weight.  Equation \textup{(N15.20)} follows because every
coefficient is positive.

For \textup{(N15.21)}, if \(0<Y<1\), all \(Y\)-basis weights are positive,
and the entire row \(i=4\) is strictly positive, including every
\(j,k\).  At \(Y=1\), only the row \(i=12\) survives, and that entire row
is strictly positive.  The structural zeros therefore cannot produce a
zero at any finite \(x=Y^{-1}\).
\end{proof}

\begin{theorem}[Sharp arbitrary-center equal-height threshold]
\label{thm:f2v15-arbitrary-centers}
Let \(c_1,c_2,c_3\in\mathbb R\), let \(t\geq0\), and set
\[
 P(z)=\prod_{j=1}^3\bigl((z-c_j)^2+t\bigr).
 \tag{N15.22}
\]
If \(e^{-D^2/2}P\) is real-rooted, then
\[
 t\leq\alpha_3.
 \tag{N15.23}
\]
Equality holds if and only if \(c_1=c_2=c_3\).
\end{theorem}

\begin{proof}
A common translation commutes with \(e^{-D^2/2}\), so assume
\(c_1+c_2+c_3=0\).  If \(A=0\), all centers vanish and V10 gives
\textup{(N15.23)}, including equality.

Now suppose \(A>0\), and use \(x,\rho\) from \textup{(N15.1)}.  The cone
condition gives \(0\leq\rho\leq1\).  The V9 pure three-pair bound gives
\(t\leq3\) whenever the image is real-rooted.

If \(0<x\leq1\) and \(\alpha\leq t\leq9/4\), write
\(t=\alpha+(9/4-\alpha)U\).  Lemma
\ref{lem:f2v15-strict} gives
\(\widehat\Sigma<0\), contradicting the first necessary sign in
\textup{(N15.5)}.  If \(0<x\leq1\) and \(9/4\leq t\leq3\), the high chart
gives \(\widehat{\mathcal P}>0\), contradicting the second necessary sign.

Finally, if \(x\geq1\) and \(\alpha\leq t\leq3\), put \(Y=x^{-1}\) and
\(t=\alpha+(3-\alpha)U\).  Then
\(Y^{12}\widehat{\mathcal P}>0\), hence
\(\widehat{\mathcal P}>0\), the same contradiction.  Thus \(A>0\) forces
\(t<\alpha\), which also proves the equality statement.
\end{proof}

\begin{corollary}[Arbitrary-center equal-channel gamma obstruction]
\label{cor:f2v15-gamma}
For a critical three-channel gamma block with equal normalized heights and
arbitrary real centers, real-rootedness requires
\[
 \kappa^2\geq
 \frac{\gamma^2R}{\alpha_3(1+\theta)}
 \geq\frac6{\alpha_3}\frac{\gamma^2}{1+\theta}.
 \tag{N15.24}
\]
The first inequality is strict unless all three centers coincide.
\end{corollary}

\begin{proof}
Use
\(t=y^2/K=\gamma^2R/(\kappa^2(1+\theta))\),
Theorem~\ref{thm:f2v15-arbitrary-centers}, and \(R\geq6\).
\end{proof}

\begin{firewall}[The equal-height cone is closed, not the GRH programme]
Theorem~\ref{thm:f2v15-arbitrary-centers} is a sharp global theorem for
one finite three-pair model.  It does not prove the corresponding statement
when the three heights differ, does not justify extraction of such a block
from the full transformed product, and does not supply the arithmetic
forcing needed for GRH.  GRH remains open.
\end{firewall}

\begin{assessment}[V15 advance and V16 direction]
The complete arbitrary-center equal-height cone is now closed:
center separation can never improve the co-centered threshold
\(\alpha_3\).  Together with V11, the two coordinate faces are known:
arbitrary heights at one center, and arbitrary centers at one height.

V16 should bridge these faces.  After centering the real parts, write
\[
 t_j=\bar t+\delta_j,\qquad \sum_j\delta_j=0,
 \tag{N15.25}
\]
and study the backward-heat sextic in the joint variables
\((A,B,\delta_1,\delta_2,\delta_3)\).  The first rigorous target is a local
joint theorem near the symmetric triple, using the V13 polynomial
\(h(x)\) as an adapted Hermite direction.  If the height perturbation has
a null first variation, compute its quadratic form before attempting a
global five-parameter cone.
\end{assessment}

\section*{Version V15 append-only record}
\addcontentsline{toc}{section}{Version V15 append-only record}
The complete V14 source was retained before this continuation.  It had
\(140968\) bytes, \(4504\) lines, and SHA--256
\[
 \texttt{7659B080FC02974712DC426989ED8CF6A409E5D1CE09628A4ECCA95E17922EDF}.
\]
V15 proves the sharp equal-height threshold \(t\leq\alpha_3\) for three
conjugate pairs with arbitrary real centers, with equality only at
coincident centers, by extending the V14 cover to a trivariate Bernstein
certificate on the full invariant cone.

% =============================================================================
% END APPENDED V15 MATERIAL
% =============================================================================


% =============================================================================
% BEGIN APPENDED V16 MATERIAL
% =============================================================================

\section{V16: a local bridge between arbitrary centers and arbitrary heights}

V11 proved the sharp threshold for arbitrary ordered heights at one common
center.  V15 proved it for arbitrary centers at one common height.  This
continuation begins to bridge those two faces.  The V13 adapted Hermite
square has a strictly negative first variation in the mean height excess
and a strictly negative quadratic variation in center separation.  A
compactness argument then permits arbitrary unequal heights above the
threshold, provided only that the center variance is sufficiently small.

\subsection{The general three-pair sextic and its second moment}

\begin{lemma}[Second-moment compactness bound]
\label{lem:f2v16-second-moment}
Let
\[
 \mathcal Q_{\mathbf c,\mathbf t}(z):=
 e^{-D^2/2}\prod_{j=1}^3\bigl((z-c_j)^2+t_j\bigr),
 \qquad t_j\geq0.
 \tag{N16.1}
\]
A common translation of the \(c_j\) translates \(\mathcal Q\), so impose
\[
 c_1+c_2+c_3=0,\qquad
 A:=\frac12\sum_{j=1}^3c_j^2.
 \tag{N16.2}
\]
If \(q_1,\ldots,q_6\) are the roots of \(\mathcal Q\), its second Newton
sum is
\[
 p_2:=\sum_{\nu=1}^6q_\nu^2
 =4A-2(t_1+t_2+t_3)+30.
 \tag{N16.3}
\]
Consequently, real-rootedness implies
\[
 t_1+t_2+t_3\leq 15+2A.
 \tag{N16.4}
\]

\begin{proof}
The input polynomial in \textup{(N16.1)} is monic of degree six.  Its
\(z^4\)-coefficient is
\(\sum_j(t_j-c_j^2)\).  Applying \(e^{-D^2/2}\) subtracts \(15\) from
that coefficient.  Since the \(z^5\)-coefficient vanishes under
\textup{(N16.2)}, Newton's identity gives \textup{(N16.3)}.  If all
\(q_\nu\) are real, then \(p_2\geq0\), proving \textup{(N16.4)}.
\end{proof}
\end{lemma}

\subsection{The joint adapted-square expansion}

Retain \(\alpha=\alpha_3\), the two positive numbers
\(\lambda_0,\lambda_1\) from \textup{(N13.8)}, and
\[
 h(x)=(x^2-\lambda_0)(x^2-\lambda_1).
 \tag{N16.5}
\]
For the roots of the general sextic define
\[
 \mathcal S_{\rm gen}(\mathbf c,\mathbf t)
 :=\sum_{\nu=1}^6h(q_\nu)^2.
 \tag{N16.6}
\]
As in V13, the right side is interpreted algebraically through Newton
sums:
\[
 \mathcal S_{\rm gen}
 =p_8-2cp_6+(c^2+2d)p_4-2cdp_2+6d^2,
 \tag{N16.7}
\]
where \(c=\lambda_0+\lambda_1\) and
\(d=\lambda_0\lambda_1\).  Thus \(\mathcal S_{\rm gen}\) is a polynomial
in the coefficients of \(\mathcal Q\), and real-rootedness implies
\[
 \mathcal S_{\rm gen}\geq0.
 \tag{N16.8}
\]

\begin{lemma}[Joint threshold Taylor ledger]
\label{lem:f2v16-joint-taylor}
Put
\[
 \delta_j:=t_j-\alpha,\qquad
 \Delta:=\delta_1+\delta_2+\delta_3.
 \tag{N16.9}
\]
As \(A+\Delta\to0\) under the one-sided condition
\(\delta_j\geq0\),
\[
 \mathcal S_{\rm gen}
 =c_AA+\frac{c_t}{3}\Delta+o(A+\Delta),
 \tag{N16.10}
\]
where \(c_A,c_t<0\) are exactly the coefficients in
\textup{(N13.14)}.  In particular,
\[
 \frac{c_A}{c_t}=\beta_3.
 \tag{N16.11}
\]
\end{lemma}

\begin{proof}
At \(\mathbf c=0\), \(\mathbf t=(\alpha,\alpha,\alpha)\), the
factorization \textup{(N13.9)} makes every term in
\textup{(N16.6)} vanish.  Direct differentiation of
\textup{(N16.7)} gives
\[
 \left.\frac{\partial\mathcal S_{\rm gen}}{\partial t_j}\right|_0
 =32(2\alpha s+12\alpha-8s-25)=\frac{c_t}{3}
 \quad(j=1,2,3).
 \tag{N16.12}
\]
Write \(c_3=-c_1-c_2\).  The quadratic center term is
\[
 c_A(c_1^2+c_1c_2+c_2^2)=c_AA,
 \tag{N16.13}
\]
because direct differentiation gives
\[
 \frac12\partial_{c_1}^2\mathcal S_{\rm gen}\big|_0
 =\partial_{c_1}\partial_{c_2}\mathcal S_{\rm gen}\big|_0
 =\frac12\partial_{c_2}^2\mathcal S_{\rm gen}\big|_0
 =c_A.
 \tag{N16.14}
\]
The linear center derivatives vanish.  The remaining Taylor terms are
\[
 O\!\left(A^{3/2}+\sqrt A\,\Delta+A\Delta+\Delta^2\right)
 =o(A+\Delta),
 \tag{N16.15}
\]
where \(\delta_j\geq0\) gives
\(\max_j|\delta_j|\leq\Delta\).  This proves
\textup{(N16.10)}.  Equation \textup{(N16.11)} is
\textup{(N13.15)}.
\end{proof}

\begin{theorem}[Local-in-centers unequal-height rigidity]
\label{thm:f2v16-local-bridge}
There exists \(\varepsilon>0\) with the following property.  Let
\(c_1,c_2,c_3\in\mathbb R\), subtract their common mean, and let \(A\) be
as in \textup{(N16.2)}.  If
\[
 A\leq\varepsilon,\qquad t_j\geq\alpha_3\quad(j=1,2,3),
 \tag{N16.16}
\]
and \(\mathcal Q_{\mathbf c,\mathbf t}\) is real-rooted, then
\[
 A=0,\qquad t_1=t_2=t_3=\alpha_3.
 \tag{N16.17}
\]
Thus every nontrivial sufficiently small center separation forces
\[
 \min_jt_j<\alpha_3
 \tag{N16.18}
\]
for real-rootedness, without any equality assumption on the three
heights.
\end{theorem}

\begin{proof}
Suppose no such \(\varepsilon\) exists.  Then there is a sequence of
counterexamples with \(A_n\to0\).  By \textup{(N16.4)},
\[
 3\alpha\leq\sum_jt_{j,n}\leq15+2A_n,
 \tag{N16.19}
\]
so, after passing to a subsequence, the height triples converge to a
limit \(\mathbf t_\ast\in[\alpha,15]^3\).  The centered real parts tend
to zero.  The set of monic real-rooted sextics is closed in coefficient
space; hence the co-centered limiting sextic
\(\mathcal Q_{\mathbf0,\mathbf t_\ast}\) is real-rooted.

Order the limiting heights.  The V11 arbitrary-height co-centered theorem
gives \(\min_jt_{\ast,j}\leq\alpha\), with equality only when all three
heights equal \(\alpha\).  Since every limiting height is at least
\(\alpha\), necessarily
\[
 \mathbf t_\ast=(\alpha,\alpha,\alpha).
 \tag{N16.20}
\]
Therefore \(\Delta_n:=\sum_j(t_{j,n}-\alpha)\to0\).
Lemma~\ref{lem:f2v16-joint-taylor} now gives
\[
 \mathcal S_{\rm gen}
 =c_AA_n+\frac{c_t}{3}\Delta_n+o(A_n+\Delta_n)<0
 \tag{N16.21}
\]
for all sufficiently large \(n\), unless
\(A_n=\Delta_n=0\).  This contradicts the necessary inequality
\textup{(N16.8)} in every nontrivial counterexample.  The exceptional case
is precisely \textup{(N16.17)}, which is real-rooted by
\textup{(N13.9)}.
\end{proof}

\begin{corollary}[Mean-height first variation]
\label{cor:f2v16-mean-debit}
Let \(A_n>0\), \(A_n\to0\), and suppose
\(\mathcal Q_{\mathbf c_n,\mathbf t_n}\) is real-rooted with centered real
parts and
\[
 t_{j,n}-\alpha_3=O(A_n)\qquad(j=1,2,3).
 \tag{N16.22}
\]
Writing \(\overline t_n=(t_{1,n}+t_{2,n}+t_{3,n})/3\), one has
\[
 \limsup_{n\to\infty}
 \frac{\overline t_n-\alpha_3}{A_n}
 \leq-\beta_3.
 \tag{N16.23}
\]
\end{corollary}

\begin{proof}
The full Taylor expansion used in
Lemma~\ref{lem:f2v16-joint-taylor}, now without a one-sided assumption,
has remainder \(o(A_n)\) under \textup{(N16.22)}.  Since
\((c_t/3)\Delta_n=c_t(\overline t_n-\alpha)\), divide
\(\mathcal S_{\rm gen}\geq0\) by \(A_n\) and use
\(c_A/c_t=\beta_3\), exactly as in V13.
\end{proof}

\begin{firewall}[The bridge is local in center variance]
Theorem~\ref{thm:f2v16-local-bridge} allows arbitrary unequal heights
above \(\alpha_3\), but only when the centered real parts have sufficiently
small variance.  V15 is global in the centers only on the equal-height
diagonal.  The joint large-center, unequal-height region remains open, as
do the infinite-product and arithmetic forcing steps.  GRH remains open.
\end{firewall}

\begin{assessment}[V16 advance and V17 direction]
The two solved coordinate faces now overlap through a genuine open
neighborhood: sufficiently small center separation cannot rescue any
triple whose three heights are all at or above \(\alpha_3\).

For a global bridge, assume toward contradiction that \(t_j\geq\alpha_3\).
Equation \textup{(N16.4)} makes the height variables compact whenever
\(A\) is bounded.  V17 should therefore add two ordered height coordinates
to the V15 three-chart cover.  The natural first boundary is
\[
 \min_jt_j=\alpha_3,
 \tag{N16.24}
\]
where V11 and V15 already determine the two axes.  At large \(A\), use the
reciprocal center chart before Bernstein conversion so that the height
bound and the center asymptotics are encoded on a compact box.
\end{assessment}

\section*{Version V16 append-only record}
\addcontentsline{toc}{section}{Version V16 append-only record}
The complete V15 source was retained before this continuation.  It had
\(151701\) bytes, \(4836\) lines, and SHA--256
\[
 \texttt{F6AA68A8181712F820B32F6CF6F0B1C6896A4973503D71BFDA742A1D20F4E063}.
\]
V16 proves local-in-centers rigidity for fully unequal heights satisfying
\(t_j\geq\alpha_3\), and identifies the same universal
\(\beta_3\) debit in the mean-height first variation.

% =============================================================================
% END APPENDED V16 MATERIAL
% =============================================================================


% =============================================================================
% BEGIN APPENDED V17 MATERIAL
% =============================================================================

\section{V17: a global unequal-height bridge on the mirror-center slice}

V16 supplied a local bridge between the arbitrary-center and
arbitrary-height faces.  This continuation obtains the first global bridge.
Keep mirror centers \((-d,0,d)\), but allow the middle squared height to
differ arbitrarily from the common outer squared height.  Evenness reduces
the heat image to a cubic in \(z^2\), and a compactified Bernstein ledger
shows that its discriminant has the wrong sign whenever both heights are at
or above \(\alpha_3\), except at the symmetric equality point.

\subsection{Exact cubic reduction}

Let \(a=d^2\), let \(u\) be the common squared height of the two outer
pairs, and let \(v\) be the squared height of the middle pair.  Put
\[
 P_{a,u,v}(z)
 =\bigl((z-d)^2+u\bigr)(z^2+v)
  \bigl((z+d)^2+u\bigr).
 \tag{N17.1}
\]
Its unit backward-heat image is even:
\[
 e^{-D^2/2}P_{a,u,v}(z)=H_{a,u,v}(z^2),
 \tag{N17.2}
\]
where
\[
\begin{aligned}
 H_{a,u,v}(X)={}&X^3+(-2a+2u+v-15)X^2\\
 &+\bigl[a^2+a(2u-2v+12)+u^2+2uv-12u-6v+45\bigr]X\\
 &+a^2(v-1)+a(2uv-2u+2v-6)\\
 &\hspace{18mm}+u^2v-u^2-2uv+6u+3v-15.
\end{aligned}
\tag{N17.3}
\]
If the sextic in \textup{(N17.2)} is real-rooted, all three roots of
\(H_{a,u,v}\) are nonnegative.  In particular,
\[
 2u+v\leq15+2a
 \tag{N17.4}
\]
from the sign of its \(X^2\)-coefficient, and
\[
 \operatorname{Disc}_XH_{a,u,v}\geq0.
 \tag{N17.5}
\]

\subsection{Compactifying the complete feasible height region}

Assume
\[
 u\geq\alpha,\qquad v\geq\alpha,
 \qquad \alpha:=\alpha_3.
 \tag{N17.6}
\]
Under \textup{(N17.4)}, define
\[
 L(a):=15+2a-3\alpha>0.
 \tag{N17.7}
\]
Every feasible pair \((u,v)\) has the representation
\[
 u=\alpha+\frac{L(a)}2\,\xi,
 \qquad
 v=\alpha+L(a)(1-\xi)\eta,
 \qquad 0\leq\xi,\eta\leq1.
 \tag{N17.8}
\]
Indeed, take
\(\xi=2(u-\alpha)/L(a)\), and then
\(\eta=(v-\alpha)/(L(a)(1-\xi))\) when \(\xi<1\).
The endpoint \(\xi=1\) forces \(v=\alpha\), so its unused \(\eta\) is
harmless.  Compactify \(a\geq0\) by
\[
 q:=\frac{a}{1+a}\in[0,1),
 \qquad a=\frac q{1-q}.
 \tag{N17.9}
\]

Define the normalized negative discriminant
\[
 \mathfrak D(a,u,v):=
 -\frac18\operatorname{Disc}_XH_{a,u,v}.
 \tag{N17.10}
\]
After the substitutions \textup{(N17.8)}--\textup{(N17.9)}, exact
reduction by
\[
 8\alpha^3-42\alpha^2+90\alpha-75=0
 \tag{N17.11}
\]
gives
\[
 \mathfrak D(a,u,v)
 =\frac{\mathcal N(q,\xi,\eta)}
        {16(1-q)^6},
 \tag{N17.12}
\]
where \(\mathcal N\) is a polynomial of tridegree \((6,5,4)\).

\begin{lemma}[Compactified discriminant ledger]
\label{lem:f2v17-discriminant-ledger}
In the tensor Bernstein expansion of \(\mathcal N\) on the unit cube,
exactly six of its \(210\) coefficients vanish:
\[
 b_{000}=0,
 \qquad
 b_{60k}=0\quad(0\leq k\leq4).
 \tag{N17.13}
\]
Every other coefficient is strictly positive.  Rigorous lower bounds for
the positive coefficients in each \(q\)-row are
\[
\begin{array}{c|rrrrrrr}
 i&0&1&2&3&4&5&6\\ \hline
 \min{}'b_{ijk}
 &18368&1763&979&315&60&5&6.
\end{array}
\tag{N17.14}
\]
The prime omits the six structural zeros.
\end{lemma}

\begin{proof}
Start with the explicit cubic \textup{(N17.3)}, compute its discriminant,
make the rational substitutions \textup{(N17.8)}--\textup{(N17.9)}, and
clear the denominator.  Formula \textup{(N15.7)} with tridegree
\((6,5,4)\) generates all \(210\) Bernstein coefficients.

For an exact sign audit, reduce every coefficient modulo
\textup{(N17.11)} and use
\[
 \frac{19434}{10000}<\alpha<\frac{19435}{10000}.
 \tag{N17.15}
\]
Outward rational interval arithmetic produces no ambiguous coefficient.
It gives the six exact zeros in \textup{(N17.13)}, \(204\) positive
coefficients, and the unrounded positive row minima
\[
 (18368.874,\ 1763.722,\ 979.825,\ 315.167,\
   60.170,\ 5.031,\ 6.400).
 \tag{N17.16}
\]
This proves \textup{(N17.14)}.
\end{proof}

\begin{lemma}[Strict positivity at every finite separation]
\label{lem:f2v17-strict}
On
\[
 0\leq q<1,\qquad 0\leq\xi,\eta\leq1,
 \tag{N17.17}
\]
one has
\[
 \mathcal N(q,\xi,\eta)\geq0,
 \tag{N17.18}
\]
with equality if and only if
\[
 (q,\xi,\eta)=(0,0,0).
 \tag{N17.19}
\]
\end{lemma}

\begin{proof}
Nonnegativity follows from the Bernstein convex-hull property.  If
\(0<q<1\), every \(q\)-basis weight is positive, and the entire row
\(i=5\) is strictly positive.  If \(q=0\), only row \(i=0\) survives.
Within that row the sole zero is \(b_{000}\), so equality requires
\(\xi=\eta=0\).  The other five zeros lie at \(i=6,\xi=0\), but can
contribute exclusively only at \(q=1\), which is not attained by finite
\(a\).  This proves the strictness statement.
\end{proof}

\begin{theorem}[Global mirror-center unequal-height threshold]
\label{thm:f2v17-mirror-unequal}
Let \(d\in\mathbb R\) and \(u,v\geq0\).  If
\[
 e^{-D^2/2}
 \bigl((z-d)^2+u\bigr)(z^2+v)\bigl((z+d)^2+u\bigr)
 \tag{N17.20}
\]
is real-rooted, then
\[
 \min\{u,v\}\leq\alpha_3.
 \tag{N17.21}
\]
If \(u,v\geq\alpha_3\), real-rootedness is possible only when
\[
 d=0,\qquad u=v=\alpha_3.
 \tag{N17.22}
\]
\end{theorem}

\begin{proof}
Suppose \(u,v\geq\alpha\).  If \textup{(N17.4)} fails, the cubic cannot
have three nonnegative roots, so the sextic is not real-rooted.  Otherwise
use \textup{(N17.8)}--\textup{(N17.9)}.  The denominator in
\textup{(N17.12)} is positive at finite \(a\).  Lemma
\ref{lem:f2v17-strict} therefore gives
\[
 \mathfrak D(a,u,v)>0
 \tag{N17.23}
\]
except at \(a=0,u=v=\alpha\).  Hence
\(\operatorname{Disc}_XH_{a,u,v}<0\), contradicting
\textup{(N17.5)}.  At the exceptional point, real-rootedness follows from
the V10 threshold factorization.
\end{proof}

\begin{corollary}[Restored heat scale]
\label{cor:f2v17-scale}
Let \(K>0\).  For mirror centers \((-d,0,d)\), common outer squared height
\(u\), and middle squared height \(v\), real-rootedness of the
\(K\)-scale backward-heat image implies
\[
 \min\{u,v\}\leq\alpha_3K.
 \tag{N17.24}
\]
Equality with both \(u,v\geq\alpha_3K\) requires
\(d=0\) and \(u=v=\alpha_3K\).
\end{corollary}

\begin{proof}
Apply Theorem~\ref{thm:f2v17-mirror-unequal} after the rescaling
\(z=\sqrt K\,w\).
\end{proof}

\begin{firewall}[This is a global bridge slice, not the full joint cone]
Theorem~\ref{thm:f2v17-mirror-unequal} allows unbounded center separation
and an arbitrary difference between the middle and outer heights, but it
uses both mirror centers and equality of the two outer heights.  General
center triples with three independent heights remain open, as do the
infinite-product and arithmetic forcing steps.  GRH remains open.
\end{firewall}

\begin{assessment}[V17 advance and V18 direction]
A nontrivial global bridge between the V11 and V15 faces is now closed.
The proof also identifies a reusable compactification: the necessary
second-moment budget turns the height excesses into a simplex, while
\(q=A/(1+A)\) compactifies center variance.

V18 should break the remaining reflection symmetry one variable at a
time.  Keep mirror centers but write the outer heights as
\[
 u_-=u-\omega,\qquad u_+=u+\omega.
 \tag{N17.25}
\]
The sextic is no longer even, but its discriminant is even in \(\omega\).
The first target is a neighborhood theorem in \(\omega\) uniform over the
compactified V17 cube.  If the discriminant coefficient of \(\omega^2\)
has the favorable sign away from the V17 equality corner, the global
bridge may persist under unequal outer heights.
\end{assessment}

\section*{Version V17 append-only record}
\addcontentsline{toc}{section}{Version V17 append-only record}
The complete V16 source was retained before this continuation.  It had
\(160573\) bytes, \(5105\) lines, and SHA--256
\[
 \texttt{B59FEF4CA1A6EBAF264E9912FC42580B3FD1C51F51FACF04312F43B95E93AE92}.
\]
V17 proves the sharp threshold on the global mirror-center family with an
arbitrary middle height and equal outer heights, using a \(210\)-coefficient
compactified discriminant certificate.

% =============================================================================
% END APPENDED V17 MATERIAL
% =============================================================================


% =============================================================================
% BEGIN APPENDED V18 MATERIAL
% =============================================================================

\section{V18: mirror centers with three independent heights}

V17 allowed an arbitrary middle height but kept the two outer heights equal.
This continuation removes that last height symmetry.  A single fixed
Hermite minor is not sufficient; instead, the necessary second-moment and
even-moment minors give complementary Bernstein charts on the compactified
height simplex.

\subsection{The asymmetric mirror family and two Hermite minors}

Let the centers be \((-d,0,d)\), put \(a=d^2\), and write the three squared
heights as
\[
 t_-=u-\omega,\qquad t_0=v,\qquad t_+=u+\omega.
 \tag{N18.1}
\]
Define
\[
 \mathcal Q_{a,u,v,\omega}(z):=
 e^{-D^2/2}
 \bigl((z+d)^2+u-\omega\bigr)(z^2+v)
 \bigl((z-d)^2+u+\omega\bigr).
 \tag{N18.2}
\]
Let \(p_k\) denote the Newton sums of its six roots, with \(p_0=6\).
For real roots, the moment matrices are positive semidefinite.  Hence the
two quantities
\[
 \mathcal C:=6p_4-p_2^2
 \tag{N18.3}
\]
and
\[
 \mathcal E:=
 \det\!\begin{pmatrix}
 p_0&p_2&p_4\\
 p_2&p_4&p_6\\
 p_4&p_6&p_8
 \end{pmatrix}
 \tag{N18.4}
\]
must both be nonnegative.  Direct Newton expansion gives the relatively
small first certificate
\[
\begin{aligned}
 \mathcal C=8\bigl(&a^2-14au+2av+24a+u^2-2uv-24u\\
                   &+v^2-12v+3\omega^2+90\bigr).
\end{aligned}
\tag{N18.5}
\]
Both \(\mathcal C\) and \(\mathcal E\) are polynomials in
\(a,u,v,\omega^2\).

\begin{proposition}[Why two charts are necessary]
\label{prop:f2v18-two-charts}
At the exact rational point
\[
 a=\frac14,\qquad u=\frac{13}{2},\qquad v=2,\qquad\omega=4,
 \tag{N18.6}
\]
all three heights exceed \(\alpha_3\), and
\[
 \mathcal E=\frac{810893}{16}>0,
 \qquad
 \mathcal C=-\frac{599}{2}<0.
 \tag{N18.7}
\]
Thus the V17 even-moment determinant alone does not exclude the asymmetric
family.  Conversely, near the V17 symmetric threshold,
\(\mathcal C\) can be nonnegative while \(\mathcal E<0\).
\end{proposition}

\begin{proof}
The three heights at \textup{(N18.6)} are
\(t_-=5/2\), \(t_0=2\), and \(t_+=21/2\), respectively.  Each exceeds
\(\alpha_3<1.9435\).
Substitution in \textup{(N18.3)}--\textup{(N18.5)} gives
\textup{(N18.7)}.  At \(\omega=0\), \(\mathcal E\) reduces to a positive
factor times the V17 cubic discriminant and is therefore negative
throughout the forbidden V17 region, whereas \(\mathcal C\) is positive
at the coalesced threshold by \textup{(N13.4)}.
\end{proof}

\subsection{A two-chart compactification}

Assume toward contradiction that all three heights satisfy
\[
 u-|\omega|\geq\alpha,\qquad v\geq\alpha,
 \qquad \alpha:=\alpha_3.
 \tag{N18.8}
\]
The second Newton sum again gives the necessary budget
\[
 2u+v\leq15+2a.
 \tag{N18.9}
\]
Set
\[
 L(a):=15+2a-3\alpha,
 \qquad
 q:=\frac{a}{1+a}\in[0,1),
 \tag{N18.10}
\]
and parametrize the full budget simplex by
\[
 u=\alpha+\frac{L(a)}2\,\xi,
 \qquad
 v=\alpha+L(a)(1-\xi)\eta,
 \qquad 0\leq\xi,\eta\leq1.
 \tag{N18.11}
\]
The outer-height condition in \textup{(N18.8)} is exactly
\[
 \omega^2=(u-\alpha)^2\tau,
 \qquad 0\leq\tau\leq1.
 \tag{N18.12}
\]
When \(u=\alpha\), this forces \(\omega=0\), and the unused value of
\(\tau\) is immaterial.

Split the simplex at \(\xi=1/2\).  In the low chart write
\[
 \xi=\frac r2,\qquad 0\leq r\leq1,
 \tag{N18.13}
\]
and in the high chart write
\[
 \xi=\frac12+\frac r2,\qquad 0\leq r\leq1.
 \tag{N18.14}
\]

\begin{lemma}[Complementary four-variable Bernstein ledgers]
\label{lem:f2v18-complementary}
After substitutions
\textup{(N18.10)}--\textup{(N18.13)}, exact reduction by
\[
 8\alpha^3-42\alpha^2+90\alpha-75=0
 \tag{N18.15}
\]
gives
\[
 \mathcal E=
 \frac{\mathcal N_E(q,r,\eta,\tau)}
      {128(1-q)^6},
 \tag{N18.16}
\]
where \(\mathcal N_E\) has multidegree \((6,6,4,3)\).
Of its \(980\) tensor Bernstein coefficients, exactly \(24\) vanish:
\[
\begin{aligned}
 b_{000\ell}&=0 &&(0\leq\ell\leq3),\\
 b_{60k\ell}&=0 &&(0\leq k\leq4,\ 0\leq\ell\leq3).
\end{aligned}
\tag{N18.17}
\]
Every other \(956\) coefficients is strictly negative.  Rigorous upper
bounds for the negative coefficients in each \(q\)-row are
\[
\begin{array}{c|rrrrrrr}
 i&0&1&2&3&4&5&6\\ \hline
 \max{}'b_{ijk\ell}
 &-3918693&-16151&-501670&-161365&-30807&-2576&-1023.
\end{array}
\tag{N18.18}
\]

In the high chart
\textup{(N18.10)}--\textup{(N18.12)}, \textup{(N18.14)} give
\[
 \mathcal C=
 \frac{\mathcal N_C(q,r,\eta,\tau)}
      {2(1-q)^2},
 \tag{N18.19}
\]
where \(\mathcal N_C\) has multidegree \((2,2,2,1)\).
All \(54\) of its tensor Bernstein coefficients are strictly negative.
Rigorous row upper bounds are
\[
\begin{array}{c|rrr}
 i&0&1&2\\ \hline
 \max b_{ijk\ell}&-94&-273&-48.
\end{array}
\tag{N18.20}
\]
\end{lemma}

\begin{proof}
Expand the heat image in \textup{(N18.2)}, apply Newton's identities through
\(p_8\), and form \textup{(N18.3)}--\textup{(N18.4)}.  Substitute the two
charts and use the four-variable version of
\textup{(N15.7)}.  This is a finite calculation with \(980+54\)
coefficients.

For a rational sign audit, reduce powers of \(\alpha\) using
\textup{(N18.15)} and use
\[
 \frac{19434}{10000}<\alpha<\frac{19435}{10000}.
 \tag{N18.21}
\]
Outward rational interval arithmetic leaves no ambiguous coefficient.
For \(\mathcal N_E\), it gives precisely the \(24\) exact zeros in
\textup{(N18.17)}, \(956\) negative coefficients, and the unrounded row
upper bounds
\[
 (-3918693.163,\ -16151.595,\ -501670.584,\ -161365.839,\
  -30807.412,\ -2576.110,\ -1024.000).
 \tag{N18.22}
\]
For \(\mathcal N_C\), all \(54\) coefficients are negative, with unrounded
row upper bounds
\[
 (-94.580,\ -273.962,\ -48.000).
 \tag{N18.23}
\]
This proves the ledgers.
\end{proof}

\begin{lemma}[Strict negativity on the physical charts]
\label{lem:f2v18-strict}
On the low chart with \(0\leq q<1\),
\[
 \mathcal E<0
 \tag{N18.24}
\]
except at \(q=r=\eta=0\), which is the physical equality configuration.
On the high chart,
\[
 \mathcal C<0.
 \tag{N18.25}
\]
\end{lemma}

\begin{proof}
The denominators in \textup{(N18.16)} and \textup{(N18.19)} are positive.
For \(0<q<1\), every \(q\)-Bernstein weight is positive and the entire
row \(i=5\) of \(\mathcal N_E\) is strictly negative.  At \(q=0\), only
row \(i=0\) survives; its zero coefficients require \(r=\eta=0\).
Then \(u=v=\alpha\) and \(\omega=0\), independently of the unused
\(\tau\).  The remaining zeros in row \(i=6\) can contribute exclusively
only at \(q=1\), which is not attained at finite separation.  This proves
\textup{(N18.24)}.  Equation \textup{(N18.25)} follows immediately because
every Bernstein coefficient of \(\mathcal N_C\) is negative.
\end{proof}

\begin{theorem}[Global mirror-center theorem with arbitrary heights]
\label{thm:f2v18-mirror-all-heights}
Let \(d\in\mathbb R\) and \(t_-,t_0,t_+\geq0\).  If
\[
 e^{-D^2/2}
 \bigl((z+d)^2+t_-\bigr)(z^2+t_0)
 \bigl((z-d)^2+t_+\bigr)
 \tag{N18.26}
\]
is real-rooted, then
\[
 \min\{t_-,t_0,t_+\}\leq\alpha_3.
 \tag{N18.27}
\]
If all three heights are at least \(\alpha_3\), real-rootedness occurs
only at
\[
 d=0,\qquad t_-=t_0=t_+=\alpha_3.
 \tag{N18.28}
\]
\end{theorem}

\begin{proof}
Write the heights as in \textup{(N18.1)}.  If
\textup{(N18.9)} fails, then \(p_2<0\), impossible for six real roots.
Otherwise use the compactification
\textup{(N18.10)}--\textup{(N18.12)}.

If \(0\leq\xi\leq1/2\), Lemma
\ref{lem:f2v18-strict} gives \(\mathcal E<0\), contradicting positive
semidefiniteness of the even-moment matrix.  If
\(1/2\leq\xi\leq1\), it gives \(\mathcal C<0\), contradicting the
two-dimensional moment minor.  The sole exception is
\textup{(N18.28)}, which is real-rooted by the V10 threshold
factorization.
\end{proof}

\begin{corollary}[Restored heat scale]
\label{cor:f2v18-scale}
At heat scale \(K>0\), for mirror centers and three arbitrary squared
heights,
\[
 \min\{t_-,t_0,t_+\}\leq\alpha_3K.
 \tag{N18.29}
\]
Equality with all three heights at least \(\alpha_3K\) requires coincident
centers and three equal heights.
\end{corollary}

\begin{proof}
Rescale \(z=\sqrt K\,w\) in
Theorem~\ref{thm:f2v18-mirror-all-heights}.
\end{proof}

\begin{firewall}[Mirror geometry remains a restriction]
V18 removes every height equality but retains the arithmetic-progression
center geometry \((-d,0,d)\).  V15 removes that center restriction only on
the equal-height diagonal, and V16 handles arbitrary heights only for small
center variance.  General separated centers with three independent heights
remain open, as do the infinite-product and arithmetic forcing steps.  GRH
remains open.
\end{firewall}

\begin{assessment}[V18 advance and V19 direction]
The complete mirror-center family now satisfies the sharp three-pair
threshold for arbitrary heights.  The proof reveals that a global bridge
need not use the full sextic discriminant: two low Hermite minors suffice
after a height-budget split.

V19 should perturb the center geometry rather than the heights.  Write a
general centered triple as
\[
 (-d-\varepsilon,\ 2\varepsilon,\ d-\varepsilon),
 \tag{N18.30}
\]
and retain the three independent heights.  The first target is a uniform
neighborhood of \(\varepsilon=0\) on compact subsets of the V18 charts.
A global proof should seek a third chart only where the two V18 minor
margins vanish, namely the symmetric equality corner and the reciprocal
\(q=1\) boundary.
\end{assessment}

\section*{Version V18 append-only record}
\addcontentsline{toc}{section}{Version V18 append-only record}
The complete V17 source was retained before this continuation.  It had
\(169090\) bytes, \(5381\) lines, and SHA--256
\[
 \texttt{15DA3675CFABBF71214D293864D89A0C97970759070813EB93D6850C76B26542}.
\]
V18 proves the sharp three-pair threshold for mirror centers with three
fully independent heights, using complementary fourth-moment and
even-Hermite Bernstein charts.

% =============================================================================
% END APPENDED V18 MATERIAL
% =============================================================================


% =============================================================================
% BEGIN APPENDED V19 MATERIAL
% =============================================================================

\section{V19: nonmirror stability and the need for a third certificate}

V18 closed the complete mirror-center family with arbitrary heights.  This
continuation perturbs the center geometry.  The V18 obstruction persists
uniformly on compact forbidden subfamilies, giving an open tube of
nonmirror configurations.  However, an exact rational example shows that
the two V18 minors do not cover the full center-shape space; the sextic
discriminant supplies a genuinely necessary third chart.

\subsection{A centered coordinate transverse to mirror geometry}

Every centered triple of real centers can be written, after choosing which
center is labelled \(0\), as
\[
 c_-=-d-\varepsilon,\qquad
 c_0=2\varepsilon,\qquad
 c_+=d-\varepsilon.
 \tag{N19.1}
\]
The mirror slice is \(\varepsilon=0\).  Its center invariants are
\[
 A=d^2+3\varepsilon^2,
 \qquad
 B=2\varepsilon(\varepsilon^2-d^2).
 \tag{N19.2}
\]
For squared heights \(\mathbf t=(t_-,t_0,t_+)\), put
\[
 \mathcal Q_{d,\varepsilon,\mathbf t}(z)
 :=
 e^{-D^2/2}\prod_{\sigma\in\{-,0,+\}}
 \bigl((z-c_\sigma)^2+t_\sigma\bigr).
 \tag{N19.3}
\]
Let \(\mathcal C(d,\varepsilon,\mathbf t)\) and
\(\mathcal E(d,\varepsilon,\mathbf t)\) denote the fourth-moment and
even-moment Hermite minors \textup{(N18.3)}--\textup{(N18.4)}.
They are polynomials in all parameters.

\begin{theorem}[Uniform nonmirror tube around compact mirror families]
\label{thm:f2v19-tube}
Let \(\mathcal K\) be any compact subset of
\[
 \left\{(d,\mathbf t):
 t_-,t_0,t_+\geq\alpha_3\right\}
 \setminus
 \left\{(0,(\alpha_3,\alpha_3,\alpha_3))\right\}.
 \tag{N19.4}
\]
Assume additionally that every point of \(\mathcal K\) satisfies the
necessary second-moment budget
\[
 t_-+t_0+t_+\leq15+2d^2.
 \tag{N19.5}
\]
Then there exists \(\varepsilon_{\mathcal K}>0\) such that
\[
 (d,\mathbf t)\in\mathcal K,\qquad
 |\varepsilon|<\varepsilon_{\mathcal K}
 \tag{N19.6}
\]
implies that \(\mathcal Q_{d,\varepsilon,\mathbf t}\) is not real-rooted.
\end{theorem}

\begin{proof}
On the mirror slice, the proof of V18 gives, at every point in
\(\mathcal K\),
\[
 \min\{\mathcal C(d,0,\mathbf t),
       \mathcal E(d,0,\mathbf t)\}<0.
 \tag{N19.7}
\]
Indeed, the low height-budget chart makes \(\mathcal E<0\), and the high
chart makes \(\mathcal C<0\).  Define the continuous function
\[
 m(d,\mathbf t)
 :=\min\{\mathcal C(d,0,\mathbf t),
          \mathcal E(d,0,\mathbf t)\}.
 \tag{N19.8}
\]
Compactness and strict negativity give
\[
 \max_{\mathcal K}m=-\mu_{\mathcal K}<0.
 \tag{N19.9}
\]
Polynomial dependence on \(\varepsilon\) is uniformly continuous on a
compact neighborhood of \(\{0\}\times\mathcal K\).  Hence, after decreasing
\(\varepsilon_{\mathcal K}\) if necessary,
\[
 \min\{\mathcal C(d,\varepsilon,\mathbf t),
       \mathcal E(d,\varepsilon,\mathbf t)\}
 <-\frac12\mu_{\mathcal K}<0.
 \tag{N19.10}
\]
A real-rooted sextic has a positive-semidefinite moment matrix, so neither
minor can be negative.  This contradiction proves the theorem.
\end{proof}

\begin{corollary}[Pointwise skew stability]
\label{cor:f2v19-pointwise}
Fix any mirror-center configuration with
\(t_-,t_0,t_+\geq\alpha_3\) other than the symmetric equality point.
Then all sufficiently small nonzero center-skew perturbations
\textup{(N19.1)}, with the heights held fixed, remain non-real-rooted.
\end{corollary}

\begin{proof}
Apply Theorem~\ref{thm:f2v19-tube} to the one-point compact set.  If the
second-moment budget fails, \(p_2<0\) already excludes real-rootedness and
continues to do so for sufficiently small \(\varepsilon\).
\end{proof}

\subsection{An exact failure of the two-minor global cover}

The compact-tube theorem cannot be promoted to all center shapes merely by
reusing the two V18 minors.

\begin{proposition}[Exact third-certificate witness]
\label{prop:f2v19-third-witness}
Take the centered real centers and squared heights
\[
 (c_-,c_0,c_+)
 =\left(-3,-\frac65,\frac{21}{5}\right),
 \qquad
 (t_-,t_0,t_+)=(10,29,2).
 \tag{N19.11}
\]
Then \(A=351/25\), the second-moment budget is strict, and the heat image is
\[
 \mathcal Q(z)
 =z^6-\frac{52}{25}z^4-\frac{3576}{25}z^3
   -\frac{73474}{625}z^2+\frac{33456}{625}z
   +\frac{7167368}{625}.
 \tag{N19.12}
\]
Its relevant exact certificates are
\[
\begin{aligned}
 p_2&=\frac{104}{25}>0,\\
 \mathcal C&=\frac{1785008}{625}>0,\\
 \mathcal E&=\frac{750157966514432}{9765625}>0,
\end{aligned}
\tag{N19.13}
\]
but
\[
 \operatorname{Disc}\mathcal Q
 =-\frac{
 79242832015812295340051925637010141216768
 }{59604644775390625}<0.
 \tag{N19.14}
\]
Thus the two minors that prove V18 are simultaneously nonnegative at a
nonmirror forbidden point, while the full discriminant excludes it.
\end{proposition}

\begin{proof}
The centers in \textup{(N19.11)} sum to zero and
\[
 A=\frac12\left(9+\frac{36}{25}+\frac{441}{25}\right)
 =\frac{351}{25}.
\]
Moreover
\[
 t_-+t_0+t_+=41
 <15+2A=\frac{1077}{25}.
\]
Expanding \textup{(N19.3)} gives \textup{(N19.12)}.  Newton's identities
and exact determinant arithmetic give
\textup{(N19.13)}--\textup{(N19.14)}.  A monic real-rooted polynomial has
nonnegative discriminant, so \(\mathcal Q\) is not real-rooted.
\end{proof}

\begin{firewall}[What the witness proves]
Proposition~\ref{prop:f2v19-third-witness} is not a counterexample to the
desired threshold: its sextic is not real-rooted.  It is a counterexample
only to the proposed two-minor proof strategy outside mirror geometry.
Any global arbitrary-center proof must add another Hermite direction,
the full discriminant, or a finer center-shape decomposition.
\end{firewall}

\begin{assessment}[V19 advance and mandatory V20 sweep]
V19 extends the V18 obstruction to an open nonmirror tube over every compact
forbidden mirror subfamily, and it identifies exactly where the proof
architecture must change.  The rational witness shows that the sextic
discriminant is a natural third chart.

V20 is a multiple of ten and must begin by sweeping all newly generated
parallel notes before further algebra.  After that sweep, parameterize
center shape by \textup{(N19.1)}--\textup{(N19.2)}, compactify
\(A\) reciprocally, and test a three-certificate cover
\[
 p_2\geq0,\qquad
 \mathcal C\geq0,\qquad
 \mathcal E\geq0,\qquad
 \operatorname{Disc}\mathcal Q\geq0.
 \tag{N19.15}
\]
The first computational target is to determine whether every point with
\(t_j\geq\alpha_3\) violates at least one of the last three inequalities.
A failed chart must be isolated with an exact rational witness as in
Proposition~\ref{prop:f2v19-third-witness}.
\end{assessment}

\section*{Version V19 append-only record}
\addcontentsline{toc}{section}{Version V19 append-only record}
The complete V18 source was retained before this continuation.  It had
\(179536\) bytes, \(5727\) lines, and SHA--256
\[
 \texttt{8C8B66C9BF803AEFA2C16C48BFA7385DB3F045FA86F03F00BC86829D7166BE6C}.
\]
V19 proves uniform nonmirror stability over compact forbidden mirror
families and records an exact nonmirror point showing that the V18
two-minor cover requires a third certificate.

% =============================================================================
% END APPENDED V19 MATERIAL
% =============================================================================

% =============================================================================
% BEGIN APPENDED V20 MATERIAL
% =============================================================================

\section{V20: a high-excess theorem on the double-center boundary}
\label{sec:f2v20-double-center}

V19 showed that the two mirror-slice certificates do not cover the full
center-shape cone.  This continuation first performs the scheduled companion
sweep, then moves to a different codimension-one boundary: two of the three
real centers coincide.  On one full half of its height-budget simplex, a
mixed-index Hermite minor admits a short analytic sign proof.  This is a new
arbitrary-height result; it does not repeat the equal-height double-center
calculation in V14.

\subsection{Scheduled companion sweep}
\label{subsec:f2v20-sweep}

The active companion notes present at the V20 cutoff were inspected through
their final theorem and assessment sections.  Their exact provenance was
\[
\begin{array}{c|r|r|l}
\text{program and active version}&\text{bytes}&\text{lines}&
 \text{SHA--256}\\
\hline
\text{SM, V7}&85128&2407&\texttt{97CC2E16\ldots A9E55D}\\
\text{Yang--Mills, V1}&17560&466&\texttt{0EEF69B3\ldots 563261}\\
\text{Navier--Stokes, V26}&278283&5319&\texttt{82C887F8\ldots 5978AD}\\
\end{array}
\tag{N20.1}
\]
Here the displayed hashes abbreviate, respectively,
\begin{align*}
&\texttt{97CC2E168F018276983CAAFA129F21359543059E093CC9BD92A89C0A8AA9E55D},\\
&\texttt{0EEF69B3521AAB3ED117685ABA372AB2117FCB9104F71DB978048633CA563261},\\
&\texttt{82C887F8C2C69E4F28FAD7C3DC8DE5B9099CB5A4BF431EBA1FFF3E91935978AD}.
\end{align*}
The frozen SM \texttt{V100\_f17} and Yang--Mills
\texttt{V100\_f10} endpoints were also checked for lineage.

No companion theorem transfers directly to the GRH finite-block problem.
Three proof rules do transfer.  First, the SM star calculation retains the
child stiffness and optimizes the exact coupled certificate; failure of a
convenient bound is not failure of the Hessian statement.  Second, the
Yang--Mills trace work warns against replacing a global cancellation by
chartwise absolute values.  Third, the Navier--Stokes rank-one refinement
records honestly when a better norm gives only a modest gain.  Accordingly,
the calculation below keeps the exact mixed Hermite determinant, exploits
its full-budget square before estimating anything, and states only the chart
actually proved.  The next scheduled companion sweep is V30.

\subsection{The double-center sextic and its budget coordinates}
\label{subsec:f2v20-setup}

Let the centered real centers be \((-2d,d,d)\).  Write the squared heights as
\[
 t_1=v,\qquad t_2=u-\omega,\qquad t_3=u+\omega,
 \tag{N20.2}
\]
and put
\[
 a=d^2,\qquad W=\omega^2.
 \tag{N20.3}
\]
The normalized heat image is
\[
 \mathcal Q_{d,u,v,\omega}(z)
 :=e^{-D^2/2}
 \bigl((z+2d)^2+v\bigr)
 \bigl((z-d)^2+u-\omega\bigr)
 \bigl((z-d)^2+u+\omega\bigr).
 \tag{N20.4}
\]
Let \(p_k\) be the Newton sum of order \(k\) of its six roots, with
\(p_0=6\).  Newton's identities give
\[
 p_2=2(15+6a-2u-v).
 \tag{N20.5}
\]
If all roots are real, then \(p_2\geq0\).  Define the exact slack
\[
 S:=15+6a-2u-v=\frac{p_2}{2}\geq0.
 \tag{N20.6}
\]
If all three heights in \textup{(N20.2)} are at least
\(\alpha:=\alpha_3\), then
\[
 v\geq\alpha,qquad W\leq(u-\alpha)^2,qquad
 2u+v\leq15+6a.
 \tag{N20.7}
\]
Put
\[
 L(a):=15+6a-3\alpha,qquad
 \xi:=\frac{2(u-\alpha)}{L(a)}.
 \tag{N20.8}
\]
The inequalities in \textup{(N20.7)} imply \(0\leq\xi\leq1\).

\begin{lemma}[Exact mixed-minor decomposition]
\label{lem:f2v20-decomposition}
Set
\[
 X:=5a-3u+15
 \tag{N20.9}
\]
and
\[
\begin{aligned}
 G(a,u,v,W):={}&-2W-162a^2+156au-12av-648a\\
 &-2u^2+36u-v^2+18v-135.
\end{aligned}
\tag{N20.10}
\]
Then the principal Hermite minor with monomial indices \(1,2\) satisfies
\[
 \mathcal K
 :=\det\!\begin{pmatrix}p_2&p_3\\p_3&p_4\end{pmatrix}
 =p_2p_4-p_3^2
 =-4\bigl(36aX^2+S G(a,u,v,W)\bigr).
 \tag{N20.11}
\]
Moreover,
\begin{align}
 p_3&=-12d\,(S-5a+3u-15), \tag{N20.12}\\
 p_4&=2\bigl(S^2+12Sa+4Su-12S-90a^2+12au-180a
             +6u^2-60u+2W+90\bigr). \tag{N20.13}
\end{align}
\end{lemma}

\begin{proof}
Expand \textup{(N20.4)}, apply Newton's identities through \(p_4\), and
replace \(v\) by \(15+6a-2u-S\).  This gives
\textup{(N20.12)}--\textup{(N20.13)}.  Expanding
\(p_2p_4-p_3^2\), collecting the full-budget square, and substituting back
for \(v\) gives \textup{(N20.11)}.  Every identity is polynomial, so no
division by \(S\), \(d\), or \(X\) is involved.
\end{proof}

\begin{lemma}[Positivity of the slack coefficient on the high chart]
\label{lem:f2v20-G-positive}
Assume \textup{(N20.7)} and
\[
 \xi\geq\frac12.
 \tag{N20.14}
\]
Then
\[
 G(a,u,v,W)>0.
 \tag{N20.15}
\]
\end{lemma}

\begin{proof}
Because the coefficient of \(W\) in \(G\) is negative and
\(W\leq(u-\alpha)^2\), it is enough to use
\[
 G_*(a,u,v):=G\bigl(a,u,v,(u-\alpha)^2\bigr).
 \tag{N20.16}
\]
Condition \textup{(N20.14)} and the budget give
\[
 u_0:=\frac{15+6a+\alpha}{4}\leq u\leq
 u_1:=\frac{15+6a-\alpha}{2},qquad
 \alpha\leq v\leq15+6a-2u.
 \tag{N20.17}
\]
For fixed \(a\), the Hessian of \(G_*\) in \((u,v)\) is
\(\operatorname{diag}(-8,-2)\).  Thus \(G_*\) is concave, and its minimum
on the triangle in \textup{(N20.17)} is attained at one of
\[
 (u_0,\alpha),\qquad (u_0,u_1),\qquad (u_1,\alpha).
 \tag{N20.18}
\]
At these vertices its values are, in the same order,
\begin{align}
 \frac34\bigl(&84a^2+(40\alpha-72)a-3\alpha^2+46\alpha-75\bigr),
 \tag{N20.19}\\
 \frac32\bigl(&12a^2+(34\alpha-90)a-\alpha^2+10\alpha+15\bigr),
 \tag{N20.20}\\
 6\bigl(&45a^2+(75-11\alpha)a-\alpha^2+10\alpha-15\bigr).
 \tag{N20.21}
\end{align}
All three are positive for \(a\geq0\).  For a completely rational sign
check, use
\[
 \frac{19434}{10000}<\alpha<\frac{19435}{10000}.
 \tag{N20.22}
\]
It gives
\[
\begin{gathered}
 40\alpha-72>\frac{717}{125},\qquad
 -3\alpha^2+46\alpha-75>3,\\
 75-11\alpha>\frac{107243}{2000},\qquad
 -\alpha^2+10\alpha-15>\frac{2627231}{4000000}>0.
 \tag{N20.23}
\end{gathered}
\]
For the only quadratic with a negative linear coefficient,
\[
 \min_{a\in\mathbb R}
 \bigl(12a^2+(34\alpha-90)a-\alpha^2+10\alpha+15\bigr)
 =-\frac{301\alpha^2-1650\alpha+1845}{12}>0,
 \tag{N20.24}
\]
because the same interval gives
\(301\alpha^2-1650\alpha+1845<-224\).  Hence every vertex value is
positive, proving \textup{(N20.15)}.
\end{proof}

\begin{theorem}[High repeated-cluster excess is impossible]
\label{thm:f2v20-high-excess}
Assume that all three heights in \textup{(N20.2)} are at least
\(\alpha_3\).  If \(\mathcal Q_{d,u,v,\omega}\) is real-rooted, then
\[
 \frac{2(u-\alpha_3)}{15+6d^2-3\alpha_3}<\frac12.
 \tag{N20.25}
\]
Equivalently, no real-rooted sextic in this double-center family can satisfy
\[
 u\geq\frac{15+6d^2+\alpha_3}{4}.
 \tag{N20.26}
\]
\end{theorem}

\begin{proof}
Suppose instead that \(\xi\geq1/2\).  For six real roots, the full moment
matrix \((p_{i+j})\) is a Gram matrix, so its principal minor
\(\mathcal K\) in \textup{(N20.11)} must be nonnegative.  By
Lemma~\ref{lem:f2v20-G-positive}, however,
\[
 \mathcal K=-4(36aX^2+SG)\leq0,
 \tag{N20.27}
\]
and the inequality is strict unless
\[
 S=0,qquad aX^2=0.
 \tag{N20.28}
\]

It remains to exclude \textup{(N20.28)}.  If \(a>0\), then \(X=0\), and
\textup{(N20.13)} together with \(W\leq(u-\alpha)^2\) gives
\[
 p_4\leq-\frac49\bigl(215a^2+30a\alpha+390a
                  -9\alpha^2+90\alpha+45\bigr)<0.
 \tag{N20.29}
\]
If \(a=0\), the high-chart interval is
\(u\in[(15+\alpha)/4,(15-\alpha)/2]\), and
\[
 p_4\leq F(u):=16u^2-(120+8\alpha)u+4\alpha^2+180.
 \tag{N20.30}
\]
The function \(F\) is convex, so its maximum on this interval occurs at an
endpoint.  The two endpoint values are
\[
 3(\alpha^2-10\alpha-15)<0,qquad
 12(\alpha^2-10\alpha+15)<0.
 \tag{N20.31}
\]
Thus \(p_4<0\) also in this case.  Both alternatives contradict
\(p_4=\sum_j\rho_j^4\geq0\) for real roots \(\rho_j\).  Hence the high
chart is impossible.
\end{proof}

\begin{corollary}[Exact localization of the unresolved double-center chart]
\label{cor:f2v20-localization}
Any still-unexcluded double-center configuration with all heights at least
\(\alpha_3\) must satisfy
\[
 0\leq u-\alpha_3<\frac{15+6d^2-3\alpha_3}{4},
 \qquad
 \alpha_3\leq v\leq15+6d^2-2u,
 \qquad
 \omega^2\leq(u-\alpha_3)^2.
 \tag{N20.32}
\]
\end{corollary}

\begin{proof}
Combine the height assumptions, \(p_2\geq0\), and
Theorem~\ref{thm:f2v20-high-excess}.
\end{proof}

\begin{firewall}[What V20 does not prove]
\label{fire:f2v20}
Theorem~\ref{thm:f2v20-high-excess} excludes one complete half of the
height-budget simplex for the center collision \((-2d,d,d)\), with three
independent heights.  It does not yet exclude the low-excess region
\textup{(N20.32)}, general nonmirror centers, infinite products, or the
arithmetic forcing step.  GRH remains open.
\end{firewall}

\begin{assessment}[V20 advance and V21 direction]
\label{ass:f2v20}
The useful new mechanism is the full-budget square in
\textup{(N20.11)}.  It converts a seemingly complicated Hermite condition
into one concavity check on a triangular height domain.  The complementary
low chart has no comparable sign from this same mixed minor.  V21 should go
upstream there: expand an even-moment determinant about the coalesced-center
threshold, preserve its cancellations, and seek a square-plus-slack
decomposition before attempting any Bernstein subdivision.  If that fails,
the next candidate is a principal minor involving indices \(0,1,5\), not
the full sextic discriminant.
\end{assessment}

\section*{Version V20 append-only record}
\addcontentsline{toc}{section}{Version V20 append-only record}

V20 was appended to the validated V19 active note whose installed SHA--256
was
\[
\texttt{26C090A106628E6DBE2572F198D81F10C17C7ECD23E5ABB5123A61780DF7A37F}.
\tag{N20.33}
\]
The V20 contribution performs the scheduled companion sweep and proves the
high-excess double-center theorem \textup{(N20.25)}.  After validation V20
replaces V19 as the sole active numeric GRH note; the frozen
\texttt{V100\_f2} archive is unchanged.

% =============================================================================
% END APPENDED V20 MATERIAL
% =============================================================================

% =============================================================================
% BEGIN APPENDED V21 MATERIAL
% =============================================================================

\section{V21: the full arbitrary-height double-center theorem}
\label{sec:f2v21-double-center-complete}

V20 excluded the half of the double-center height simplex in which the
mean excess at the repeated center is large.  This continuation closes the
complementary half.  The proof uses three principal Hermite minors on three
overlapping rational charts.  The smallest minor handles large isolated
height, the even-moment determinant handles small separation, and a
mixed high-order determinant handles the reciprocal-separation chart.

\subsection{Low-excess coordinates and three necessary minors}
\label{subsec:f2v21-coordinates}

Retain the notation of V20:
\[
 (c_1,c_2,c_3)=(-2d,d,d),\qquad
 (t_1,t_2,t_3)=(v,u-\omega,u+\omega),
 \qquad a=d^2,\quad W=\omega^2,
 \tag{N21.1}
\]
and let \(\mathcal Q_{d,u,v,\omega}\) be the heat sextic
\textup{(N20.4)}.  Assume
\[
 v\geq\alpha,\qquad u-|\omega|\geq\alpha,
 \qquad \alpha:=\alpha_3.
 \tag{N21.2}
\]
If \(\mathcal Q_{d,u,v,\omega}\) is real-rooted, V20 gives the budget
\[
 2u+v\leq15+6a
 \tag{N21.3}
\]
and excludes
\[
 \xi:=\frac{2(u-\alpha)}{L(a)}\geq\frac12,
 \qquad L(a):=15+6a-3\alpha.
 \tag{N21.4}
\]
It therefore remains to consider \(0\leq\xi\leq1/2\).  Parametrize this
closed chart by
\[
\begin{aligned}
 q&:=\frac{a}{1+a},&
 \xi&=\frac r2,\\
 u&=\alpha+\frac{L(a)}2\,\xi,&
 v&=\alpha+L(a)(1-\xi)\eta,\\
 W&=(u-\alpha)^2\tau,&
 0&\leq q,r,\eta,\tau\leq1.
\end{aligned}
\tag{N21.5}
\]
The value of \(\tau\) is immaterial when \(u=\alpha\).

Let \(p_k\) again denote the Newton sums of the six roots.  Besides
\[
 \mathcal C:=6p_4-p_2^2
 \tag{N21.6}
\]
and the even-moment determinant
\[
 \mathcal E:=
 \det\!\begin{pmatrix}
 p_0&p_2&p_4\\
 p_2&p_4&p_6\\
 p_4&p_6&p_8
 \end{pmatrix},
 \tag{N21.7}
\]
introduce
\[
 \mathcal M:=
 \det\!\begin{pmatrix}
 p_0&p_1&p_5\\
 p_1&p_2&p_6\\
 p_5&p_6&p_{10}
 \end{pmatrix}.
 \tag{N21.8}
\]
The sextic is centered, so \(p_1=0\), and hence
\[
 \mathcal M=p_2(6p_{10}-p_5^2)-6p_6^2.
 \tag{N21.9}
\]
If all six roots are real, every principal moment matrix is positive
semidefinite.  Consequently
\[
 \mathcal C\geq0,\qquad \mathcal E\geq0,\qquad
 \mathcal M\geq0
 \tag{N21.10}
\]
are all necessary.

\subsection{A three-chart Bernstein ledger}
\label{subsec:f2v21-ledger}

Use unit-cube variables \((x,r,y,\tau)\).  The three charts are
\[
\begin{array}{c|c|c}
\text{chart}&q&\eta\\ \hline
C&x&\displaystyle\frac14+\frac34y\\[1mm]
E&\displaystyle\frac x4&\displaystyle\frac y4\\[1mm]
M&\displaystyle\frac14+\frac34x&\displaystyle\frac y4 .
\end{array}
\tag{N21.11}
\]
Thus the \(C\)-chart covers \(\eta\geq1/4\).  The \(E\)- and \(M\)-charts
together cover \(\eta\leq1/4\), split at \(q=1/4\).

\begin{lemma}[Exact low-excess certificate ledgers]
\label{lem:f2v21-ledgers}
After the substitutions \textup{(N21.5)} and \textup{(N21.11)}, there are
polynomials \(\mathcal N_C,\mathcal N_E,\mathcal N_M\) such that
\begin{align}
 \mathcal C&=\frac{\mathcal N_C(x,r,y,\tau)}
                   {8(1-x)^2},
 &&\deg\mathcal N_C=(2,2,2,1), \tag{N21.12}\\
 \mathcal E&=\frac{\mathcal N_E(x,r,y,\tau)}
                   {2048(1-x/4)^6},
 &&\deg\mathcal N_E=(6,6,5,3), \tag{N21.13}\\
 \mathcal M&=\frac{\mathcal N_M(x,r,y,\tau)}
                   {8192\{3(1-x)/4\}^6},
 &&\deg\mathcal N_M=(6,6,5,2). \tag{N21.14}
\end{align}
Here degrees are listed in the order \((x,r,y,\tau)\).

\begin{enumerate}[label=\textup{(\roman*)}]
\item All \(54\) tensor Bernstein coefficients of \(\mathcal N_C\) are
strictly negative.  Rigorous upper bounds for the maximum coefficient in
each \(x\)-row are
\[
\begin{array}{c|rrr}
 i&0&1&2\\ \hline
 \max b_{ijk\ell}&-141&-2200&-972.
\end{array}
\tag{N21.15}
\]

\item Of the \(1176\) tensor Bernstein coefficients of
\(\mathcal N_E\), exactly four vanish:
\[
 b_{000\ell}=0\qquad(0\leq\ell\leq3).
 \tag{N21.16}
\]
Every other \(1172\) coefficients is strictly negative.  Excluding the
zeros, rigorous row upper bounds are
\[
\begin{gathered}
\begin{array}{c|rrrr}
 i&0&1&2&3\\ \hline
 \max{}'b_{ijk\ell}
 &-37636441&-10840556&-20776612&-29664218
\end{array}\\[1mm]
\begin{array}{c|rrr}
 i&4&5&6\\ \hline
 \max{}'b_{ijk\ell}
 &-37283978&-43457798&-48095700
\end{array}
\end{gathered}
\tag{N21.17}
\]

\item Of the \(882\) tensor Bernstein coefficients of
\(\mathcal N_M\), exactly three vanish:
\[
 b_{600\ell}=0\qquad(0\leq\ell\leq2).
 \tag{N21.18}
\]
Every other \(879\) coefficients is strictly negative.  Excluding the
zeros, rigorous row upper bounds are
\[
\begin{gathered}
\begin{array}{c|rrrr}
 i&0&1&2&3\\ \hline
 \max{}'b_{ijk\ell}
 &-1752657524&-1269999979&-896409702&-370305937
\end{array}\\[1mm]
\begin{array}{c|rrr}
 i&4&5&6\\ \hline
 \max{}'b_{ijk\ell}
 &-29542002&-104266661&-21233664
\end{array}
\end{gathered}
\tag{N21.19}
\]
\end{enumerate}
\end{lemma}

\begin{proof}
Expand \(\mathcal Q_{d,u,v,\omega}\), apply Newton's identities through
\(p_{10}\), and form the three principal minors
\textup{(N21.6)}--\textup{(N21.8)}.  The substitutions
\textup{(N21.5)} and \textup{(N21.11)} give the positive denominators in
\textup{(N21.12)}--\textup{(N21.14)}.  Convert each numerator from the
power basis to the tensor Bernstein basis by the exact formula
\textup{(N15.7)}.

For a rational sign audit, reduce every power of \(\alpha\) modulo
\[
 8\alpha^3-42\alpha^2+90\alpha-75=0
 \tag{N21.20}
\]
and use
\[
 \frac{19434}{10000}<\alpha<\frac{19435}{10000}.
 \tag{N21.21}
\]
Every reduced coefficient is quadratic in \(\alpha\).  Exact rational
endpoint evaluation, including the vertex when that quadratic is concave,
gives the zero sets \textup{(N21.16)}, \textup{(N21.18)} and the row
bounds \textup{(N21.15)}, \textup{(N21.17)}, \textup{(N21.19)}.
There are no unresolved interval signs.
\end{proof}

\begin{lemma}[Strict signs on the three physical charts]
\label{lem:f2v21-strict}
On the \(C\)-chart,
\[
 \mathcal C<0.
 \tag{N21.22}
\]
On the \(E\)-chart,
\[
 \mathcal E<0
 \tag{N21.23}
\]
except at
\[
 d=0,\qquad u=v=\alpha,\qquad\omega=0.
 \tag{N21.24}
\]
On the \(M\)-chart at finite \(a\),
\[
 \mathcal M<0.
 \tag{N21.25}
\]
\end{lemma}

\begin{proof}
All denominators in \textup{(N21.12)}--\textup{(N21.14)} are positive
on their finite physical charts.  Part \textup{(i)} of
Lemma~\ref{lem:f2v21-ledgers} immediately gives
\textup{(N21.22)}.

For \(\mathcal N_E\), the only zero coefficients are in
\textup{(N21.16)}.  At \(x=0\), they can contribute exclusively only
when \(r=y=0\); this is precisely \textup{(N21.24)}, independently of the
unused \(\tau\).  If \(x>0\), a strictly negative \(x\)-row has positive
Bernstein weight; at \(x=1\), the terminal row is itself strictly negative.
This proves \textup{(N21.23)}.

For \(\mathcal N_M\), the zeros in \textup{(N21.18)} can contribute
exclusively only at \(x=1,r=y=0\).  The endpoint \(x=1\) corresponds to
\(q=1\), hence \(a=\infty\), and is not attained at finite separation.
At every \(0\leq x<1\), a strictly negative row has positive weight.
This proves \textup{(N21.25)}.
\end{proof}

\begin{theorem}[Global double-center theorem with arbitrary heights]
\label{thm:f2v21-double-center-all-heights}
Let \(d\in\mathbb R\) and \(t_1,t_2,t_3\geq0\).  If
\[
 e^{-D^2/2}
 \bigl((z+2d)^2+t_1\bigr)
 \bigl((z-d)^2+t_2\bigr)
 \bigl((z-d)^2+t_3\bigr)
 \tag{N21.26}
\]
is real-rooted, then
\[
 \min\{t_1,t_2,t_3\}\leq\alpha_3.
 \tag{N21.27}
\]
If all three heights are at least \(\alpha_3\), real-rootedness occurs
only at
\[
 d=0,\qquad t_1=t_2=t_3=\alpha_3.
 \tag{N21.28}
\]
\end{theorem}

\begin{proof}
Write the heights as in \textup{(N21.1)}.  Suppose all are at least
\(\alpha=\alpha_3\).  If the budget \textup{(N21.3)} fails, then
\(p_2<0\), impossible for real roots.

Under the budget, Theorem~\ref{thm:f2v20-high-excess} excludes
\(\xi\geq1/2\).  Use the low-excess coordinates \textup{(N21.5)}.
If \(\eta\geq1/4\), then \(\mathcal C<0\) by
Lemma~\ref{lem:f2v21-strict}, contradicting
\textup{(N21.10)}.  If \(\eta\leq1/4\) and \(q\leq1/4\), then
\(\mathcal E<0\), with the sole exception \textup{(N21.24)}.  If
\(\eta\leq1/4\) and \(q\geq1/4\), then \(\mathcal M<0\).
The three charts cover the entire low-excess domain.  Thus only
\textup{(N21.28)} remains, and that configuration is real-rooted by the
V10 threshold factorization.
\end{proof}

\begin{corollary}[Restored heat scale]
\label{cor:f2v21-scale}
At heat scale \(K>0\), the double-center family with three arbitrary
squared heights satisfies
\[
 \min\{t_1,t_2,t_3\}\leq\alpha_3K.
 \tag{N21.29}
\]
If all three heights are at least \(\alpha_3K\), equality requires
coincident centers and three equal heights.
\end{corollary}

\begin{proof}
Rescale \(z=\sqrt K\,w\) in
Theorem~\ref{thm:f2v21-double-center-all-heights}.
\end{proof}

\subsection{The missing V19 certificate is a mixed minor}
\label{subsec:f2v21-v19-witness}

The V19 rational witness showed that \(\mathcal C\) and \(\mathcal E\)
alone do not cover general nonmirror centers.  The mixed minor used in V20
already detects that witness.

\begin{proposition}[Exact third certificate at the V19 witness]
\label{prop:f2v21-third-certificate}
For centers
\[
 \left(-3,-\frac65,\frac{21}{5}\right)
 \tag{N21.30}
\]
and squared heights \((10,29,2)\), the heat sextic \textup{(N19.8)}
satisfies
\[
\begin{aligned}
 p_2&=\frac{104}{25}>0,\\
 \mathcal C&=\frac{1785008}{625}>0,\\
 \mathcal E&=\frac{750157966514432}{9765625}>0,
\end{aligned}
\qquad
 \mathcal K:=p_2p_4-p_3^2
 =-\frac{2846121984}{15625}<0.
 \tag{N21.31}
\]
Thus \(\mathcal K\) is an explicit third Hermite certificate where the
two V18 minors simultaneously fail.
\end{proposition}

\begin{proof}
Apply Newton's identities to the exact sextic \textup{(N19.8)} and
substitute in the four displayed expressions.  All arithmetic is over
\(\mathbb Q\).
\end{proof}

\begin{firewall}[Boundary closure is not the full center cone]
\label{fire:f2v21}
V18 proves the arbitrary-height theorem on the mirror-center boundary and
V21 proves it on the double-center boundary.  V16 supplies only a local
arbitrary-center bridge near coincident centers.  These results do not yet
cover the interior center-shape cone, the infinite product, or the
arithmetic forcing step.  GRH remains open.
\end{firewall}

\begin{assessment}[V21 advance and V22 direction]
\label{ass:f2v21}
The complete double-center boundary now has the same sharp threshold as the
mirror boundary, with no height symmetry.  More importantly,
Proposition~\ref{prop:f2v21-third-certificate} identifies a small third
minor that repairs the exact V19 two-certificate witness.

V22 should therefore return to the general centered parametrization
\[
 (-d-\varepsilon,\,2\varepsilon,\,d-\varepsilon)
 \tag{N21.32}
\]
with arbitrary heights and test the three-certificate cover
\((\mathcal C,\mathcal E,\mathcal K)\) before introducing the full sextic
discriminant.  The mirror and double-center theorems give two closed
center-shape edges, V16 controls the coincident corner, and
\(\mathcal K\) now controls the known interior witness.  The first task is
to locate, exactly, any point where all three minors are nonnegative; only
such a point justifies a fourth certificate or an upstream reformulation.
\end{assessment}

\section*{Version V21 append-only record}
\addcontentsline{toc}{section}{Version V21 append-only record}

V21 was appended to the validated V20 active note whose installed SHA--256
was
\[
 \texttt{98B00A81A57CFB0BC12C0EB7861A5AD32EA8EF7D04A59FAACD0495A39AEF09A8}.
 \tag{N21.33}
\]
V21 completes the arbitrary-height double-center boundary by three exact
Bernstein charts and identifies the mixed minor that detects the V19
nonmirror witness.  After validation V21 replaces V20 as the sole active
numeric GRH note; the frozen \texttt{V100\_f2} archive is unchanged.

% =============================================================================
% END APPENDED V21 MATERIAL
% =============================================================================

% =============================================================================
% BEGIN APPENDED V22 MATERIAL
% =============================================================================

\section{V22: an exact interior gap for three Hermite certificates}
\label{sec:f2v22-interior-gap}

V21 closed the arbitrary-height mirror and double-center boundaries and
showed that the mixed minor
\(\mathcal K=p_2p_4-p_3^2\) repairs the particular two-certificate gap
found in V19.  The next question is whether
\((\mathcal C,\mathcal E,\mathcal K)\) covers the full interior
center-shape cone.  It does not.  This continuation gives an exact integer
witness and then certifies an open five-dimensional rational box on which
all three minors have the allowed sign while a fourth principal Hermite
minor is strictly negative.

\subsection{General centered triples and two higher mixed minors}
\label{subsec:f2v22-minors}

Let \(c_1+c_2+c_3=0\), let \(t_1,t_2,t_3\geq0\), and define
\[
 \mathcal Q_{\bm c,\bm t}(z)
 :=e^{-D^2/2}\prod_{j=1}^3\bigl((z-c_j)^2+t_j\bigr).
 \tag{N22.1}
\]
Write \(p_k\) for the Newton sums of its six roots.  Retain
\[
\begin{aligned}
 \mathcal C&:=6p_4-p_2^2,\\
 \mathcal E&:=
 \det\!\begin{pmatrix}
 p_0&p_2&p_4\\
 p_2&p_4&p_6\\
 p_4&p_6&p_8
 \end{pmatrix},\\
 \mathcal K&:=p_2p_4-p_3^2.
\end{aligned}
\tag{N22.2}
\]
Introduce the two principal mixed determinants
\[
 \mathcal M_3:=
 \det(p_{i+j})_{i,j\in\{0,1,3\}},
 \qquad
 \mathcal M_5:=
 \det(p_{i+j})_{i,j\in\{0,1,5\}}.
 \tag{N22.3}
\]
The determinant called \(\mathcal M\) in V21 is \(\mathcal M_5\).
If all six roots are real, every quantity in
\textup{(N22.2)}--\textup{(N22.3)} must be nonnegative.

\begin{proposition}[Exact integer three-certificate gap]
\label{prop:f2v22-integer-gap}
Take
\[
 (c_1,c_2,c_3)=(24,-18,-6),
 \qquad
 (t_1,t_2,t_3)=(2,32,60).
 \tag{N22.4}
\]
All three heights exceed \(\alpha_3\), and
\[
 \mathcal Q_{\bm c,\bm t}(z)
 =z^6-857z^4-6960z^3+179281z^2
   +2847216z+19577003.
 \tag{N22.5}
\]
At this point,
\[
\begin{aligned}
 p_2&=1714,\\
 \mathcal C&=1572848,\\
 \mathcal E&=20841710115526400,\\
 \mathcal K&=852566236,
\end{aligned}
\qquad
\begin{aligned}
 \mathcal M_3&=-386174231712,\\
 \mathcal M_5&=-77825092052888640.
\end{aligned}
\tag{N22.6}
\]
In particular, \(\mathcal C,\mathcal E,\mathcal K\) are simultaneously
positive, while two other principal Hermite minors are negative.  The full
sextic discriminant is
\[
 \operatorname{Disc}_z\mathcal Q_{\bm c,\bm t}
 =-2581253966439394175516898789328543219712<0.
 \tag{N22.7}
\]
\end{proposition}

\begin{proof}
Insert \textup{(N22.4)} in \textup{(N22.1)} and apply
\[
 e^{-D^2/2}
 =\sum_{k=0}^3\frac{(-1/2)^k}{k!}D^{2k}
 \tag{N22.8}
\]
to obtain \textup{(N22.5)}.  Newton's identities through \(p_{10}\)
give \textup{(N22.6)}.  The resultant of \(\mathcal Q\) and
\(\mathcal Q'\) gives \textup{(N22.7)}.  Every calculation is over
\(\mathbb Z\).
\end{proof}

\subsection{A robust rational gap box}
\label{subsec:f2v22-box}

The preceding witness is not an isolated sign accident.  For
\[
 0\leq x_1,x_2,y_1,y_2,y_3\leq1,
 \tag{N22.9}
\]
put
\[
\begin{aligned}
 c_1&=24+\frac{2x_1-1}{100},&
 c_2&=-18+\frac{2x_2-1}{100},&
 c_3&=-c_1-c_2,\\
 t_1&=2+\frac{2y_1-1}{100},&
 t_2&=32+\frac{2y_2-1}{100},&
 t_3&=60+\frac{2y_3-1}{100}.
\end{aligned}
\tag{N22.10}
\]
Throughout this box,
\[
\begin{gathered}
 23.99\leq c_1\leq24.01,\qquad
 -18.01\leq c_2\leq-17.99,\qquad
 -6.02\leq c_3\leq-5.98,\\
 t_1\geq1.99>\alpha_3.
\end{gathered}
\tag{N22.11}
\]
Thus it lies strictly inside the center-shape cone: the centers are
pairwise distinct and none is zero, so the box meets neither the
double-center nor mirror boundary.

\begin{lemma}[Five-variable Bernstein box certificate]
\label{lem:f2v22-box-ledger}
After \textup{(N22.10)}, expand each certificate in the tensor Bernstein
basis on the unit five-cube \textup{(N22.9)}.

\begin{enumerate}[label=\textup{(\roman*)}]
\item The multidegree of \(\mathcal C\) is \((4,4,2,2,2)\).
All \(675\) Bernstein coefficients are positive, and their minimum is
\[
 \frac{19616624810001}{12500000}>1569329.
 \tag{N22.12}
\]

\item The multidegree of \(\mathcal E\) is \((10,10,5,5,5)\).
All \(26136\) Bernstein coefficients are positive, and their minimum is
\[
 \frac{64645009974716824505486953655384393}
      {3125000000000000000}>2\cdot10^{16}.
 \tag{N22.13}
\]

\item The multidegree of \(\mathcal K\) is \((6,6,3,3,3)\).
All \(3136\) Bernstein coefficients are positive, and their minimum is
\[
 \frac{53080686204007056651}{62500000000}>849290000.
 \tag{N22.14}
\]

\item The multidegree of \(\mathcal M_3\) is \((6,6,3,3,3)\).
All \(3136\) Bernstein coefficients are negative, and their maximum is
\[
 -\frac{2403142819687341466387173}{6250000000000}
 <-384502000000.
 \tag{N22.15}
\]
\end{enumerate}
\end{lemma}

\begin{proof}
Substitute the affine rational parametrization \textup{(N22.10)} into the
Newton-sum expressions \textup{(N22.2)}--\textup{(N22.3)}.  Apply the
five-variable form of the exact power-to-Bernstein conversion
\textup{(N15.7)}.  All input coefficients are rational, so every Bernstein
coefficient is rational.  Direct comparison gives the four exact extrema
\textup{(N22.12)}--\textup{(N22.15)}; no algebraic-root or floating-point
sign decision is used.
\end{proof}

\begin{corollary}[Open failure of the three-certificate cover]
\label{cor:f2v22-open-gap}
On the full rational box \textup{(N22.10)},
\[
 \mathcal C>0,\qquad \mathcal E>0,\qquad \mathcal K>0,
 \qquad \mathcal M_3<0.
 \tag{N22.16}
\]
Consequently the three-certificate system
\((\mathcal C,\mathcal E,\mathcal K)\) cannot prove the global
arbitrary-height threshold, even after arbitrary subdivision of that box.
\end{corollary}

\begin{proof}
The Bernstein basis functions are nonnegative and sum to one.  Apply the
coefficient signs in Lemma~\ref{lem:f2v22-box-ledger}.
\end{proof}

\subsection{Discovery ledger for a fourth certificate}
\label{subsec:f2v22-discovery}

A discovery-only scan sampled \(2{,}000{,}000\) points in the full
center-shape and height-budget domain.  It found \(117\) points with
\(\mathcal C,\mathcal E\geq0\), seven with
\(\mathcal C,\mathcal E,\mathcal K\geq0\), and none of those seven with
\(\mathcal M_3\geq0\).  The same seven also had \(\mathcal M_5<0\).
A separate randomized integer scan of \(3{,}000{,}000\) admissible points
found no point at which
\[
 \mathcal C,\mathcal E,\mathcal K,\mathcal M_3\geq0
 \tag{N22.17}
\]
simultaneously.  These finite scans motivate the next direction but are
not used as proof.

\begin{firewall}[Four sampled certificates are not a global theorem]
\label{fire:f2v22}
The exact box proves that three particular necessary minors are
insufficient and that \(\mathcal M_3\) repairs a nonempty interior region.
It does not prove that four minors cover the full center-height cone.
Even such a finite-block theorem would not by itself control the infinite
product or establish the arithmetic forcing step.  GRH remains open.
\end{firewall}

\begin{assessment}[V22 advance and V23 direction]
\label{ass:f2v22}
V22 replaces the isolated V19 failure with a robust geometric fact: there
is an open interior box on which three low Hermite certificates all have
the allowed sign.  The principal minor \(\mathcal M_3\) supplies an exact
fourth obstruction throughout that box.

V23 should test the four-certificate system
\[
 (\mathcal C,\mathcal E,\mathcal K,\mathcal M_3)
 \tag{N22.18}
\]
on the compactified general domain.  The center cubic should be ordered
first, so that its shape parameter runs between the mirror edge and the
double-center edge while the three heights remain attached to their
ordered centers.  The immediate task is an exact search for a point where
all four minors are nonnegative.  If none is found, use the V18 and V21
boundary ledgers, the V16 coincident-center neighborhood, and interior
rational boxes to build a certified finite cover.  A fifth certificate
should be introduced only after an exact four-certificate gap is found.
\end{assessment}

\section*{Version V22 append-only record}
\addcontentsline{toc}{section}{Version V22 append-only record}

V22 was appended to the validated V21 active note whose installed SHA--256
was
\[
 \texttt{248081CC8D2B9C407F7A9D09B6CD72ED655F8BB2B93511302B17C6AA1F0DBCE1}.
 \tag{N22.19}
\]
V22 gives an exact integer gap for the first three Hermite certificates,
certifies a full rational neighborhood of that gap, and identifies
\(\mathcal M_3\) as the next global certificate.  After validation V22
replaces V21 as the sole active numeric GRH note; the frozen
\texttt{V100\_f2} archive is unchanged.

% =============================================================================
% END APPENDED V22 MATERIAL
% =============================================================================

% =============================================================================
% BEGIN APPENDED V23 MATERIAL
% =============================================================================

\section{V23: the reciprocal-boundary leading form of the fourth minor}
\label{sec:f2v23-reciprocal}

V22 found an open interior region where
\(\mathcal C,\mathcal E,\mathcal K\) all have the allowed sign and the
fourth minor \(\mathcal M_3\) is negative.  A boundary-heavy search found
no simultaneous four-minor gap, but it also showed that
\(\mathcal M_3\) can approach zero after normalization when the real
centers separate.  This continuation identifies the exact cancellation
at infinity and proves that bounded, or sufficiently slowly growing,
heights cannot escape through the reciprocal-center boundary.

\subsection{The center cubic and a derivative-square identity}
\label{subsec:f2v23-center-cubic}

Let
\[
 x_1+x_2+x_3=0,\qquad
 A_{\bm x}:=\frac12\sum_{j=1}^3x_j^2,
 \tag{N23.1}
\]
and write the center cubic as
\[
 P_{\bm x}(X):=\prod_{j=1}^3(X-x_j)
 =X^3-A_{\bm x}X-B_{\bm x}.
 \tag{N23.2}
\]
Set
\[
 \Delta_j:=P_{\bm x}'(x_j)
 =(x_j-x_k)(x_j-x_\ell),
 \qquad \{j,k,\ell\}=\{1,2,3\}.
 \tag{N23.3}
\]

\begin{lemma}[Derivative-square identity]
\label{lem:f2v23-derivative-square}
For every centered real triple,
\[
 \sum_{j=1}^3\Delta_j^2=9A_{\bm x}^2.
 \tag{N23.4}
\]
\end{lemma}

\begin{proof}
From \textup{(N23.2)},
\(\Delta_j=3x_j^2-A_{\bm x}\).  Newton's identities for the center cubic
give
\[
 \sum_jx_j^2=2A_{\bm x},
 \qquad
 \sum_jx_j^4=2A_{\bm x}^2.
 \tag{N23.5}
\]
Therefore
\[
\begin{aligned}
 \sum_j(3x_j^2-A_{\bm x})^2
 &=9\sum_jx_j^4-6A_{\bm x}\sum_jx_j^2+3A_{\bm x}^2\\
 &=18A_{\bm x}^2-12A_{\bm x}^2+3A_{\bm x}^2
 =9A_{\bm x}^2.
\end{aligned}
\]
\end{proof}

\subsection{Exact expansion of the fourth minor}
\label{subsec:f2v23-expansion}

For \(\lambda>0\), put \(c_j=\lambda x_j\) in the general heat sextic
\textup{(N22.1)}, while retaining arbitrary squared heights
\(\bm t=(t_1,t_2,t_3)\).  Recall
\[
 \mathcal M_3
 =\det(p_{i+j})_{i,j\in\{0,1,3\}}
 =6p_2p_6-6p_4^2-p_2p_3^2,
 \tag{N23.6}
\]
where the second identity uses \(p_0=6\) and \(p_1=0\).

\begin{lemma}[Reciprocal-center expansion]
\label{lem:f2v23-leading}
There are explicit polynomials
\[
 R_4(\bm x,\bm t),\qquad
 R_2(\bm x,\bm t),\qquad
 R_0(\bm t)
 \tag{N23.7}
\]
such that the exact identity
\[
\begin{aligned}
 \mathcal M_3(\lambda\bm x,\bm t)
={}&-48A_{\bm x}\lambda^6
 \left(\sum_{j=1}^3t_j\Delta_j^2-9A_{\bm x}^2\right)\\
 &+\lambda^4R_4(\bm x,\bm t)
  +\lambda^2R_2(\bm x,\bm t)+R_0(\bm t)
\end{aligned}
\tag{N23.8}
\]
holds.  On the normalized shape set \(A_{\bm x}=1\), there are absolute
constants \(C_4,C_2,C_0\) such that, for
\[
 T:=\max\{1,t_1,t_2,t_3\},
 \tag{N23.9}
\]
\[
 |R_4|\leq C_4T^2,\qquad
 |R_2|\leq C_2T^3,\qquad
 |R_0|\leq C_0T^4.
 \tag{N23.10}
\]
\end{lemma}

\begin{proof}
Expand
\[
 e^{-D^2/2}\prod_{j=1}^3
 \bigl((z-\lambda x_j)^2+t_j\bigr),
 \tag{N23.11}
\]
and apply Newton's identities through \(p_6\).  To make the coefficient
calculation checkable, put \(S_1:=t_1+t_2+t_3\).  The required Newton
sums are
\[
\begin{aligned}
p_2={}&4A_{\bm x}\lambda^2+30-2S_1,\\
p_3={}&6B_{\bm x}\lambda^3-6\lambda\sum_jt_jx_j,\\
p_4={}&4A_{\bm x}^2\lambda^4
 +\lambda^2\left(72A_{\bm x}-12\sum_jt_jx_j^2\right)\\
&\quad+2\sum_jt_j^2-36S_1+270,\\
p_6={}&(4A_{\bm x}^3+6B_{\bm x}^2)\lambda^6
 +\lambda^4\left(132A_{\bm x}^2-30\sum_jt_jx_j^4\right)\\
&\quad+\lambda^2\left(30\sum_jt_j^2x_j^2
 -252\sum_jt_jx_j^2-48A_{\bm x}S_1+1116A_{\bm x}\right)\\
&\quad-2\sum_jt_j^3+54\sum_jt_j^2
 +24\sum_{i<j}t_it_j-558S_1+2790.
\end{aligned}
\tag{N23.11a}
\]
These identities follow by coefficient comparison in \textup{(N23.11)}
and the Newton recurrence, so they require no division or genericity
assumption.  Insert \textup{(N23.11a)} in \textup{(N23.6)}.  The
apparent \(\lambda^8\) term vanishes.  This is the algebraic expression
of the relation
\[
 x_j^3=A_{\bm x}x_j+B_{\bm x}
 \tag{N23.12}
\]
on the three-point center support: the columns indexed by
\(\{0,1,3\}\) are dependent at leading order.  Collecting the next term
gives
\[
 [\lambda^6]\mathcal M_3
 =-48A_{\bm x}
   \left(\sum_jt_jP_{\bm x}'(x_j)^2-9A_{\bm x}^2\right),
 \tag{N23.13}
\]
which is \textup{(N23.8)}.

Only even powers of \(\lambda\) remain.  Direct collection shows that
\(R_4\) is quartic in \(\bm x\) and of height degree at most \(2\);
\(R_2\) is quadratic in \(\bm x\) and of height degree at most \(3\);
and \(R_0\) has height degree at most \(4\).  The normalized real shape
set
\[
 \left\{\bm x\in\mathbb R^3:
 \sum_jx_j=0,\ A_{\bm x}=1\right\}
 \tag{N23.14}
\]
is compact.  Taking coefficient suprema on this set proves
\textup{(N23.10)}.
\end{proof}

\begin{theorem}[No sublinear-height escape at infinite separation]
\label{thm:f2v23-no-sublinear-escape}
Let \(\lambda_n\to\infty\), and let
\[
 \sum_jx_{j,n}=0,\qquad A_{\bm x_n}=1.
 \tag{N23.15}
\]
Suppose
\[
 t_{j,n}\geq\alpha_3
 \quad(1\leq j\leq3),\qquad
 T_n:=\max_jt_{j,n}=o(\lambda_n).
 \tag{N23.16}
\]
Then, for all sufficiently large \(n\),
\[
 \mathcal M_3(\lambda_n\bm x_n,\bm t_n)<0.
 \tag{N23.17}
\]
Consequently the corresponding heat sextics are not real-rooted.
\end{theorem}

\begin{proof}
Lemma~\ref{lem:f2v23-derivative-square} and
\(t_{j,n}\geq\alpha_3\) give
\[
\begin{aligned}
 \sum_jt_{j,n}\Delta_{j,n}^2-9
 &\geq\alpha_3\sum_j\Delta_{j,n}^2-9\\
 &=9(\alpha_3-1)>0.
\end{aligned}
\tag{N23.18}
\]
Thus the leading term in \textup{(N23.8)} is at most
\[
 -432(\alpha_3-1)\lambda_n^6.
 \tag{N23.19}
\]
By \textup{(N23.10)}, the absolute value of the remainder divided by
\(\lambda_n^6\) is bounded by
\[
 C_4\frac{T_n^2}{\lambda_n^2}
 +C_2\frac{T_n^3}{\lambda_n^4}
 +C_0\frac{T_n^4}{\lambda_n^6}=o(1).
 \tag{N23.20}
\]
This proves \textup{(N23.17)}.  A real-rooted sextic would have
\(\mathcal M_3\geq0\), because \(\mathcal M_3\) is a principal moment
minor.
\end{proof}

\begin{corollary}[Uniform bounded-height reciprocal chart]
\label{cor:f2v23-bounded-height}
For every \(T<\infty\), there exists \(\Lambda(T)<\infty\) such that,
whenever
\[
 A_{\bm x}=1,\qquad
 \lambda\geq\Lambda(T),\qquad
 \alpha_3\leq t_j\leq T,
 \tag{N23.21}
\]
the heat sextic with centers \(\lambda\bm x\) is not real-rooted.
\end{corollary}

\begin{proof}
If no uniform \(\Lambda(T)\) existed, a violating sequence would satisfy
the hypotheses of Theorem~\ref{thm:f2v23-no-sublinear-escape}.
\end{proof}

\subsection{Boundary-heavy four-certificate search}
\label{subsec:f2v23-search}

A discovery-only search tested \(24{,}000{,}000\) points, divided among
Dirichlet height allocations, near-threshold heights, and allocations with
one height forced close to \(\alpha_3\).  Center variances ranged from
\(10^{-6}\) to \(10^{12}\).  In total,
\[
 2512976
\tag{N23.22}
\]
samples satisfied \(\mathcal C,\mathcal E\geq0\), and
\[
 2412340
\tag{N23.23}
\]
also satisfied \(\mathcal K\geq0\).  None of those samples had
\(\mathcal M_3\geq0\).  Writing \(A_{\bm c}:=\frac12\sum_jc_j^2\), the closest normalized value observed in the
extreme-separation, near-threshold regime was approximately
\[
 \frac{\mathcal M_3}
 {(1+A_{\bm c}+t_1+t_2+t_3)^4}
 \approx-9.29\cdot10^{-8}.
 \tag{N23.24}
\]
This approach to zero is explained by the canceled \(\lambda^8\) term in
Lemma~\ref{lem:f2v23-leading}.  The sample counts are evidence for chart
selection only and are not used in the proof.

\begin{firewall}[The projective-height boundary remains]
\label{fire:f2v23}
Theorem~\ref{thm:f2v23-no-sublinear-escape} controls infinite center
separation only when the maximum squared height is \(o(\lambda)\), where
\(\lambda\) is the center-separation scale.  The second-moment budget
allows heights of order \(\lambda^2\).  The intermediate and projective
height layers are not covered, and the finite interior four-certificate
cover is not proved.  The infinite-product and arithmetic forcing steps
also remain open.  GRH remains open.
\end{firewall}

\begin{assessment}[V23 advance and V24 direction]
\label{ass:f2v23}
The fourth certificate has a structural reason to work at large
separation.  Its nominal leading term vanishes because the center cubic
makes \(1,x,x^3\) dependent; the first surviving term is nevertheless
uniformly negative above the height threshold because
\(\sum_jP_{\bm x}'(x_j)^2=9A_{\bm x}^2\).

V24 should blow up the remaining reciprocal boundary with
\[
 t_j=\lambda^2s_j,\qquad
 s_j\geq0,\qquad \sum_js_j\leq2+o(1),
 \tag{N23.25}
\]
and compute the leading homogeneous forms of
\((\mathcal C,\mathcal E,\mathcal K,\mathcal M_3)\).
Faces with \(s_j=0\) require separate boundary-layer charts; the
sublinear theorem above supplies the deepest such face.  An additional
minor should not be introduced unless those exact projective forms admit
a simultaneous nonnegative point.
\end{assessment}

\section*{Version V23 append-only record}
\addcontentsline{toc}{section}{Version V23 append-only record}

V23 was appended to the validated V22 active note whose installed
SHA--256 was
\[
 \texttt{25C74667C49F730DD3D3C2324BFC1D79A8412E0E812FC63450A02E0635FFE9FC}.
 \tag{N23.26}
\]
V23 derives the exact reciprocal-center leading form of the fourth minor
and proves that no bounded or sublinear-height sequence can escape through
infinite center separation.  After validation V23 replaces V22 as the
sole active numeric GRH note; the frozen \texttt{V100\_f2} archive is
unchanged.

% =============================================================================
% END APPENDED V23 MATERIAL
% =============================================================================

% =============================================================================
% BEGIN APPENDED V24 MATERIAL
% =============================================================================

\section{V24: the complete separated-center reciprocal chart}
\label{sec:f2v24-separated}

V23 controlled the reciprocal-center boundary when the squared heights
grow more slowly than the center scale.  The second-moment budget still
permits squared heights of order the square of that scale.  This
continuation resolves the whole remaining reciprocal chart whenever the
normalized centers stay a fixed distance apart.  In fact, in that chart
the heat sextic eventually has no real zero, a stronger conclusion than
failure of real-rootedness.

The proof does not require a new Hermite minor.  It combines the exact
projective leading polynomial with a local calculation at every limiting
double zero.  Thus it also treats mixed faces on which some heights are
projective and others are subprojective.

\subsection{Exact projective forms}
\label{subsec:f2v24-forms}

Let
\[
 \sum_{j=1}^3x_j=0,\qquad
 A_{\bm x}:=\frac12\sum_{j=1}^3x_j^2=1,
 \qquad c_j=\lambda x_j,
 \tag{N24.1}
\]
and put \(t_j=\lambda^2s_j\).  The second-moment bound
\textup{(N16.4)} becomes
\[
 s_j\geq0,\qquad
 s_1+s_2+s_3\leq2+\frac{15}{\lambda^2}.
 \tag{N24.2}
\]
Define the projective input
\[
 F_{\bm x,\bm s}(w)
 :=\prod_{j=1}^3\bigl((w-x_j)^2+s_j\bigr).
 \tag{N24.3}
\]
If \(\pi_k\) denotes the \(k\)-th Newton sum of its six roots, then
\[
 \pi_k=2\sum_{j=1}^3
 \sum_{r=0}^{\lfloor k/2\rfloor}
 (-1)^r\binom{k}{2r}x_j^{k-2r}s_j^r.
 \tag{N24.4}
\]
In particular,
\[
\begin{aligned}
 \pi_2&=2\sum_j(x_j^2-s_j),\\
 \pi_3&=2\sum_j(x_j^3-3x_js_j),\\
 \pi_4&=2\sum_j(x_j^4-6x_j^2s_j+s_j^2),\\
 \pi_6&=2\sum_j(x_j^6-15x_j^4s_j
              +15x_j^2s_j^2-s_j^3),\\
 \pi_8&=2\sum_j(x_j^8-28x_j^6s_j+70x_j^4s_j^2
              -28x_j^2s_j^3+s_j^4).
\end{aligned}
\tag{N24.5}
\]

\begin{lemma}[Leading forms of the four current certificates]
\label{lem:f2v24-leading-forms}
For the heat sextic \textup{(N22.1)}, under \textup{(N24.1)}--
\textup{(N24.2)}, the projective asymptotic expansions are
\[
\begin{aligned}
 \lambda^{-4}\mathcal C&=\mathcal C_\infty+O(\lambda^{-2}),
 &\quad \mathcal C_\infty&:=6\pi_4-\pi_2^2,\\
 \lambda^{-6}\mathcal K&=\mathcal K_\infty+O(\lambda^{-2}),
 &\quad \mathcal K_\infty&:=\pi_2\pi_4-\pi_3^2,\\
 \lambda^{-8}\mathcal M_3&=\mathcal M_{3,\infty}+O(\lambda^{-2}),
 &\quad \mathcal M_{3,\infty}&:=
 6\pi_2\pi_6-6\pi_4^2-\pi_2\pi_3^2,\\
 \lambda^{-12}\mathcal E&=\mathcal E_\infty+O(\lambda^{-2}),
 &\quad \mathcal E_\infty&:=
 \det\!\begin{pmatrix}
 6&\pi_2&\pi_4\\
 \pi_2&\pi_4&\pi_6\\
 \pi_4&\pi_6&\pi_8
 \end{pmatrix}.
\end{aligned}
\tag{N24.6}
\]
The estimates are uniform when
\((\bm x,\bm s)\) ranges over a compact set.
\end{lemma}

\begin{proof}
At \(z=\lambda w\), the heat image satisfies the exact identity
\[
 \lambda^{-6}\mathcal Q_{\lambda\bm x,\lambda^2\bm s}
 (\lambda w)
 =\exp\!\left(-\frac{D_w^2}{2\lambda^2}\right)
 F_{\bm x,\bm s}(w).
 \tag{N24.7}
\]
The exponential truncates after its third derivative term.  Hence the
right side converges coefficientwise to \(F_{\bm x,\bm s}\), uniformly
on compact parameter sets.  Newton sums are polynomial functions of the
coefficients.  Their weights under root dilation give the four powers of
\(\lambda\) in \textup{(N24.6)}.  Substitution in
\textup{(N22.2)}--\textup{(N22.3)} proves the displayed formulas.
\end{proof}

The projective discriminant also has a simple exact factorization.  For
\(i<j\), set
\[
 R_{ij}:=(x_i-x_j)^4
 +2(x_i-x_j)^2(s_i+s_j)+(s_i-s_j)^2.
 \tag{N24.8}
\]

\begin{lemma}[Projective discriminant factorization]
\label{lem:f2v24-discriminant}
The discriminant of \textup{(N24.3)} is
\[
 \operatorname{Disc}F_{\bm x,\bm s}
 =-64s_1s_2s_3\prod_{1\leq i<j\leq3}R_{ij}^2.
 \tag{N24.9}
\]
Moreover \(R_{ij}\geq0\), with equality exactly when
\(x_i=x_j\) and \(s_i=s_j\).
\end{lemma}

\begin{proof}
Write \(f_j(w)=(w-x_j)^2+s_j\).  Its discriminant is \(-4s_j\), while
direct evaluation at the two roots of \(f_i\) gives
\[
 \operatorname{Res}(f_i,f_j)
 =\bigl((x_i-x_j)^2-s_i+s_j\bigr)^2
   +4(x_i-x_j)^2s_i
 =R_{ij}.
 \tag{N24.10}
\]
The product formula
\[
 \operatorname{Disc}\!\left(\prod_jf_j\right)
 =\prod_j\operatorname{Disc}(f_j)
  \prod_{i<j}\operatorname{Res}(f_i,f_j)^2
 \tag{N24.11}
\]
proves \textup{(N24.9)}.  Formula \textup{(N24.8)} is a sum of
nonnegative terms.  Its first term forces \(x_i=x_j\) at equality, and
the last then forces \(s_i=s_j\).
\end{proof}

Thus the discriminant already excludes the generic projective interior:
it is strictly negative when all \(s_j>0\) and the three quadratic
factors are distinct.  Its vanishing on height faces is not a genuine
escape, as the next local identity shows.

\subsection{The exact local heat identity}
\label{subsec:f2v24-local}

Fix an index \(j\), put \(z=\lambda x_j+y\), and factor the unheated
sextic as
\[
 H(y)=(y^2+t_j)G_j(y),\qquad
 G_j(y):=\prod_{k\ne j}
 \left((y+\lambda(x_j-x_k))^2+t_k\right).
 \tag{N24.12}
\]

\begin{lemma}[Local quotient identity]
\label{lem:f2v24-local-identity}
Assume \(t_k>0\) for \(k\ne j\).  For real \(y\), \(G_j(y)>0\), and the heat sextic obeys
\[
\begin{aligned}
 \frac{\mathcal Q_{\lambda\bm x,\bm t}
       (\lambda x_j+y)}{G_j(y)}
={}&y^2+t_j-1
 -2y\frac{G_j'}{G_j}
 +\left(\frac32-\frac{y^2+t_j}{2}\right)
   \frac{G_j''}{G_j}\\
&+y\frac{G_j'''}{G_j}
 +\left(\frac{y^2+t_j}{8}-\frac58\right)
   \frac{G_j''''}{G_j}.
\end{aligned}
\tag{N24.13}
\]
If
\[
 \min_{k\ne j}|x_j-x_k|\geq\delta>0,\qquad y=o(\lambda),
 \tag{N24.14}
\]
then, uniformly under the budget \textup{(N16.4)},
\[
 \frac{G_j^{(r)}(y)}{G_j(y)}=O_\delta(\lambda^{-r})
 \qquad(1\leq r\leq4).
 \tag{N24.15}
\]
Consequently, if also \(t_j=o(\lambda^2)\), the right side of
\textup{(N24.13)} is
\[
 y^2+t_j-1+o(1).
 \tag{N24.16}
\]
\end{lemma}

\begin{proof}
Because \(H\) has degree six,
\[
 e^{-D_y^2/2}H=H-\frac12H''+\frac18H''''-\frac1{48}H''''''.
 \tag{N24.17}
\]
For \(u=y^2+t_j\), the product rule gives
\[
\begin{aligned}
 (uG_j)''&=2G_j+4yG_j'+uG_j'',\\
 (uG_j)''''&=12G_j''+8yG_j'''+uG_j'''',\\
 (uG_j)''''''&=30G_j''''.
\end{aligned}
\tag{N24.18}
\]
Substitution in \textup{(N24.17)} and division by \(G_j(y)>0\)
proves \textup{(N24.13)}.

Under \textup{(N24.14)}, each of the two quadratic factors in \(G_j\)
has real displacement of size at least \(\delta\lambda/2\) for all
sufficiently large \(\lambda\).  Differentiating those two factors and
dividing termwise by their product proves \textup{(N24.15)}; the
nonnegative heights can only increase the denominators.  In
\textup{(N24.13)}, the four error terms are respectively bounded by
\[
 O_\delta\!\left(
 \frac{|y|}{\lambda}
 +\frac{1+y^2+t_j}{\lambda^2}
 +\frac{|y|}{\lambda^3}
 +\frac{1+y^2+t_j}{\lambda^4}
 \right).
 \tag{N24.19}
\]
This is \(o(1)\) under the final two hypotheses.
\end{proof}

\subsection{Uniform exclusion away from center collision}
\label{subsec:f2v24-uniform}

\begin{theorem}[Complete separated-center reciprocal chart]
\label{thm:f2v24-complete-separated}
For every \(\delta>0\), there is \(\Lambda_\delta<\infty\) with the
following property.  Suppose
\[
 \lambda\geq\Lambda_\delta,\qquad
 \sum_jx_j=0,\qquad A_{\bm x}=1,
 \qquad \min_{i<j}|x_i-x_j|\geq\delta,
 \tag{N24.20}
\]
and
\[
 t_j\geq\alpha_3\quad(1\leq j\leq3),
 \qquad \sum_jt_j\leq15+2\lambda^2.
 \tag{N24.21}
\]
Then \(\mathcal Q_{\lambda\bm x,\bm t}\) has no real zero.  In
particular it is not real-rooted.
\end{theorem}

\begin{proof}
If the assertion were false, there would be sequences satisfying
\textup{(N24.20)}--\textup{(N24.21)} with \(\lambda_n\to\infty\) and
a real zero \(z_n\) of the corresponding heat sextic.  Put
\[
 s_{j,n}:=\frac{t_{j,n}}{\lambda_n^2},
 \qquad w_n:=\frac{z_n}{\lambda_n}.
 \tag{N24.22}
\]
The normalized center circle is compact, and \textup{(N24.21)} bounds
\(\bm s_n\).  After passage to a subsequence,
\[
 \bm x_n\longrightarrow\bm x_\infty,qquad
 \bm s_n\longrightarrow\bm s_\infty.
 \tag{N24.23}
\]
Identity \textup{(N24.7)} shows that the scaled monic sextics have
uniformly bounded coefficients.  Their zeros are therefore uniformly
bounded, so, after a further subsequence, \(w_n\to w_\infty\in\mathbb R\).
Coefficientwise convergence in \textup{(N24.7)} gives
\[
 0=F_{\bm x_\infty,\bm s_\infty}(w_\infty).
 \tag{N24.24}
\]
Every factor in \textup{(N24.3)} is nonnegative on the real axis.
Consequently, for some unique \(j\),
\[
 w_\infty=x_{j,\infty},\qquad s_{j,\infty}=0;
 \tag{N24.25}
\]
uniqueness follows from the \(\delta\)-separation.  Relabel along a
subsequence so that this \(j\) is fixed, and set
\[
 y_n:=z_n-\lambda_nx_{j,n}.
 \tag{N24.26}
\]
Equations \textup{(N24.22)}--\textup{(N24.25)} imply
\[
 y_n=o(\lambda_n),\qquad t_{j,n}=o(\lambda_n^2).
 \tag{N24.27}
\]
Apply Lemma~\ref{lem:f2v24-local-identity} at the real point \(y_n\).
Since \(z_n\) is a zero and \(G_j(y_n)>0\), it yields
\[
 0=y_n^2+t_{j,n}-1+o(1)
 \geq\alpha_3-1+o(1)>0,
 \tag{N24.28}
\]
a contradiction.
\end{proof}

\begin{corollary}[Arbitrary projective and mixed height faces]
\label{cor:f2v24-mixed}
Under a fixed positive separation of the normalized centers, no sequence
at infinite center variance can evade the threshold obstruction by
choosing any mixture of bounded, intermediate, or projective squared
heights allowed by the second-moment budget.
\end{corollary}

\begin{proof}
This is precisely Theorem~\ref{thm:f2v24-complete-separated}; no rate of
growth is imposed on the individual heights beyond
\textup{(N24.21)}.
\end{proof}

\subsection{Discovery ledger and the remaining collision chart}
\label{subsec:f2v24-ledger}

A discovery-only search evaluated the leading forms \textup{(N24.6)}
at \(16{,}000{,}000\) points in the simplex \(\sum_js_j\leq2\), using
bulk, near-zero, near-budget, and one-face regimes.  Of these,
\(3{,}539{,}602\) had \(\mathcal C_\infty,\mathcal E_\infty\geq0\),
and \(3{,}462{,}791\) also had \(\mathcal K_\infty\geq0\).  The only
21 floating-point reports with \(\mathcal M_{3,\infty}\geq0\) occurred
at heights of order \(10^{-14}\).  Re-evaluating the closest report at
80 decimal digits gave
\[
 \mathcal M_{3,\infty}
 =-3.448339681892056651354\ldots\,10^{-14}<0.
 \tag{N24.29}
\]
These computations motivate, but are not used in, the proof.  Theorem
\ref{thm:f2v24-complete-separated} avoids the still-unproved global
semialgebraic assertion that the four leading forms themselves cover the
simplex.

\begin{firewall}[Only the center-collision blow-up remains at infinity]
\label{fire:f2v24}
Theorem~\ref{thm:f2v24-complete-separated} is uniform only when the
normalized center gaps have a fixed positive lower bound.  V21 handles
the exact double-center locus, but the transition in which two
normalized centers collide while their physical separation has a finite
or intermediate scale is not yet uniform.  Closing that collision
blow-up would complete the infinite-center boundary; it would not by
itself prove the finite interior, the infinite-product passage, or the
arithmetic forcing step.  GRH remains open.
\end{firewall}

\begin{assessment}[V24 advance and V25 direction]
\label{ass:f2v24}
The height stratification at the reciprocal boundary is now unnecessary
away from center collisions: one compactness argument covers all of it.
The constant \(1\) in the local model \(y^2+t_j-1\) explains why the
previous threshold \(\alpha_3>1\) gives a uniform strict margin.

V25 should blow up a repeated normalized center.  After permutation and
reflection, write the close pair as
\[
 c_2=\lambda d+\frac r2,qquad
 c_3=\lambda d-\frac r2,qquad
 c_1=-2\lambda d,
 \tag{N24.30}
\]
which is already centered.  The exact
double-center theorem V21 is the \(r=0\) face, whereas the present theorem
controls normalized gaps bounded away from zero.  The missing chart must
track \(r\), the difference of the two close heights, and their common
projective height.  A direct local quartic heat model around the close
pair is preferable to adding another global moment minor.
\end{assessment}

\section*{Version V24 append-only record}
\addcontentsline{toc}{section}{Version V24 append-only record}

V24 was appended to the validated V23 active note whose installed
SHA--256 was
\[
 \texttt{8EAE02E833BE8F5D321C7A277BB75DA428A73B670F0E541FC7006E4DB55E821E}.
 \tag{N24.31}
\]
V24 computes the exact projective leading forms and discriminant, proves
an exact local heat identity, and closes the entire infinite-separation
chart uniformly away from normalized center collisions.  After
validation V24 replaces V23 as the sole active numeric GRH note; the
frozen \texttt{V100\_f2} archive is unchanged.

% =============================================================================
% END APPENDED V24 MATERIAL
% =============================================================================

% =============================================================================
% BEGIN APPENDED V25 MATERIAL
% =============================================================================

\section{V25: completion of the infinite-center boundary}
\label{sec:f2v25-infinite-boundary}

V24 excluded every infinite-center sequence whose normalized centers
remain separated.  The only missing limit had two normalized centers
coalescing.  V21 already excludes the exact double-center family, but its
three-certificate proof does not give the uniform transition estimate
needed when the close pair separates on a smaller physical scale.

This continuation goes upstream to the close pair itself.  Its local
quartic heat image has a uniform positive decomposition above the
threshold.  An exact product identity then shows that the distant third
quadratic changes this quartic only by a relative \(o(1)\).  The resulting
collision-chart theorem, combined with V24, completes the entire
infinite-center boundary for the three-pair problem.

\subsection{A globally positive close-pair quartic}
\label{subsec:f2v25-quartic}

For \(h\in\mathbb R\) and \(u,v>0\), define
\[
 A_{h,u,v}(y)
 :=\bigl((y-h)^2+u\bigr)\bigl((y+h)^2+v\bigr)
 \tag{N25.1}
\]
and its unit backward-heat image
\[
 Q_{h,u,v}(y):=e^{-D_y^2/2}A_{h,u,v}(y).
 \tag{N25.2}
\]
Since \(A_{h,u,v}\) has degree four,
\[
\begin{aligned}
 Q_{h,u,v}(y)
={}&y^4+(u+v-6-2h^2)y^2+2h(u-v)y\\
 &+h^4+h^2(u+v+2)+uv-u-v+3.
\end{aligned}
\tag{N25.3}
\]

\begin{lemma}[Uniform close-pair positivity]
\label{lem:f2v25-quartic-positive}
Let \(\alpha>3/2\).  If \(u,v\geq\alpha\), then, for every real \(h,y\),
\[
 Q_{h,u,v}(y)\geq4\alpha-6>0.
 \tag{N25.4}
\]
For \(U:=u-\alpha\) and \(V:=v-\alpha\), the exact decomposition is
\[
\begin{aligned}
 Q_{h,u,v}(y)
={}&\bigl(y^2-h^2-(3-\alpha)\bigr)^2
+4(\alpha-1)h^2+4\alpha-6\\
&+U\bigl((y+h)^2+\alpha-1\bigr)
&+V\bigl((y-h)^2+\alpha-1\bigr)+UV.
\end{aligned}
\tag{N25.5}
\]
\end{lemma}

\begin{proof}
Substitute \(u=\alpha+U\), \(v=\alpha+V\) in
\textup{(N25.3)} and complete the square in \(y^2\).  This gives
\textup{(N25.5)}.  Every term after \(4\alpha-6\) is nonnegative because
\(\alpha>1\) and \(U,V\geq0\), proving \textup{(N25.4)}.
\end{proof}

For the sextic product estimate, positivity alone must also control the
derivatives of the unheated quartic.

\begin{lemma}[Quartic derivative comparison]
\label{lem:f2v25-derivative-comparison}
For every fixed \(\alpha>3/2\), there is \(C_\alpha<\infty\) such that
\[
 |A_{h,u,v}^{(k)}(y)|\leq
 C_\alpha Q_{h,u,v}(y)
 \qquad(0\leq k\leq4)
 \tag{N25.6}
\]
for all real \(h,y\) and all \(u,v\geq\alpha\).
\end{lemma}

\begin{proof}
Put
\[
 f_-:=(y-h)^2+u,\qquad f_+:=(y+h)^2+v,
 \qquad A:=f_-f_+,\qquad Q:=Q_{h,u,v}(y).
 \tag{N25.7}
\]
First compare \(A\) with \(Q\).  At \(u=v=\alpha\),
\[
\begin{aligned}
 A_{h,\alpha,\alpha}(y)
 &=(y^2-h^2)^2+2\alpha(y^2+h^2)+\alpha^2,\\
 Q_{h,\alpha,\alpha}(y)
 &=D^2+4(\alpha-1)h^2+4\alpha-6,\\
 D&:=y^2-h^2-(3-\alpha).
\end{aligned}
\tag{N25.8}
\]
The second and third lines control \(1\), \(h^2\), \(D^2\), and hence
\(|D|\), \(y^2\), and \((y^2-h^2)^2\).  Therefore
\[
 A_{h,\alpha,\alpha}(y)
 \leq C_{0,\alpha}Q_{h,\alpha,\alpha}(y).
 \tag{N25.9}
\]
This elementary implication uses only
\(|D|\leq(D^2+1)/2\) and \(4\alpha-6>0\).

The height-dependent parts of \(A\) and \(Q\) are, exactly,
\[
\begin{aligned}
 A={}&A_{h,\alpha,\alpha}
 +U\bigl((y+h)^2+\alpha\bigr)
 +V\bigl((y-h)^2+\alpha\bigr)+UV,\\
 Q={}&Q_{h,\alpha,\alpha}
 +U\bigl((y+h)^2+\alpha-1\bigr)
 +V\bigl((y-h)^2+\alpha-1\bigr)+UV.
\end{aligned}
\tag{N25.10}
\]
For \(q\geq0\),
\[
 \frac{q+\alpha}{q+\alpha-1}\leq
 \frac{\alpha}{\alpha-1}.
 \tag{N25.11}
\]
Equations \textup{(N25.9)}--\textup{(N25.11)} therefore give
\[
 A\leq C_{1,\alpha}Q.
 \tag{N25.12}
\]

It remains to compare the derivatives with \(A\).  The elementary bound
\[
 \frac{|x|}{x^2+a}\leq\frac1{2\sqrt a}
 \qquad(a>0)
 \tag{N25.13}
\]
and \(f_\pm\geq\alpha\) yield
\[
\begin{aligned}
 |A'|&\leq\frac{2}{\sqrt\alpha}A,\\
 |A''|&\leq\frac6\alpha A,\\
 |A'''|&=24|y|\leq\frac{24}{\alpha^{3/2}}A,\\
 |A''''|&=24\leq\frac{24}{\alpha^2}A.
\end{aligned}
\tag{N25.14}
\]
For the second line, use
\[
 A''=2(f_-+f_+)+8(y-h)(y+h),\qquad
 f_\pm\leq A/\alpha,\qquad
 4\alpha|(y-h)(y+h)|\leq A.
 \tag{N25.15}
\]
For the third, \(2A\geq\alpha(f_-+f_+)\geq2\alpha y^2\), and
\(A\geq\alpha^2\).  Combining \textup{(N25.12)} and
\textup{(N25.14)} proves \textup{(N25.6)}.
\end{proof}

\subsection{Lifting the quartic through a distant quadratic}
\label{subsec:f2v25-lift}

\begin{lemma}[Exact quartic--quadratic heat product identity]
\label{lem:f2v25-product}
Let \(A\) be a polynomial of degree at most four and \(G\) a polynomial
of degree at most two.  Then
\[
\begin{aligned}
 e^{-D^2/2}(AG)
={}&G\,e^{-D^2/2}A
G'\left(-A'+\frac12A'''\right)\\
&+G''\left(-\frac12A+\frac34A''
                 -\frac5{16}A''''\right).
\end{aligned}
\tag{N25.16}
\]
\end{lemma}

\begin{proof}
Expand
\[
 e^{-D^2/2}(AG)
 =\sum_{k=0}^3\frac{(-1/2)^k}{k!}D^{2k}(AG)
 \tag{N25.17}
\]
and apply the Leibniz rule.  Terms with no derivative on \(G\) sum to
\(G e^{-D^2/2}A\).  The coefficients of \(G'\) are
\(-A'+A'''/2\), and those of \(G''\) are
\(-A/2+3A''/4-5A''''/16\).  Higher derivatives of \(G\) and derivatives
of \(A\) above order four vanish.
\end{proof}

\begin{corollary}[Distant-factor stability]
\label{cor:f2v25-distant}
Fix \(\alpha>3/2\).  Let
\[
 G_L(y):=(y+3L)^2+t_1,\qquad t_1\geq0,
 \tag{N25.18}
\]
and suppose \(y_L=o(L)\).  Uniformly for arbitrary
\(h_L\in\mathbb R\) and \(u_L,v_L\geq\alpha\),
\[
 \frac{e^{-D_y^2/2}
 \bigl(A_{h_L,u_L,v_L}G_L\bigr)(y_L)}
 {G_L(y_L)Q_{h_L,u_L,v_L}(y_L)}
 =1+o(1).
 \tag{N25.19}
\]
In particular the numerator is strictly positive for all sufficiently
large \(L\).
\end{corollary}

\begin{proof}
At \(y_L=o(L)\), the real displacement \(y_L+3L\) has size comparable
to \(L\), and hence
\[
 \frac{G_L'(y_L)}{G_L(y_L)}=O(L^{-1}),
 \qquad
 \frac{G_L''(y_L)}{G_L(y_L)}=O(L^{-2}),
 \tag{N25.20}
\]
uniformly for \(t_1\geq0\).  Divide \textup{(N25.16)} by
\(G_LQ_{h_L,u_L,v_L}\), and apply
Lemma~\ref{lem:f2v25-derivative-comparison}.  The difference from \(1\)
is \(O_\alpha(L^{-1})\).  Positivity follows from
Lemma~\ref{lem:f2v25-quartic-positive}.
\end{proof}

\subsection{The collision blow-up}
\label{subsec:f2v25-collision}

Every centered triple with a distinguished close pair can, after a
permutation, be written
\[
 c_1=-2L,\qquad c_2=L+h,\qquad c_3=L-h.
 \tag{N25.21}
\]
Its center variance is
\[
 A_{\bm c}:=\frac12\sum_jc_j^2=3L^2+h^2.
 \tag{N25.22}
\]

\begin{theorem}[Complete normalized center-collision chart]
\label{thm:f2v25-collision}
Let \(L_n\to\infty\), \(h_n=o(L_n)\), and let the centers be
\textup{(N25.21)}.  Suppose
\[
 t_{j,n}\geq\alpha_3\quad(1\leq j\leq3),\qquad
 \sum_jt_{j,n}\leq15+6L_n^2+2h_n^2.
 \tag{N25.23}
\]
Then, for all sufficiently large \(n\), the heat sextic
\[
 \mathcal Q_n(z):=
 e^{-D_z^2/2}\prod_{j=1}^3
 \bigl((z-c_{j,n})^2+t_{j,n}\bigr)
 \tag{N25.24}
\]
has no real zero.
\end{theorem}

\begin{proof}
Suppose instead that \(z_n\in\mathbb R\) is a zero along a subsequence.
Set
\[
 s_{j,n}:=\frac{t_{j,n}}{L_n^2},\qquad
 w_n:=\frac{z_n}{L_n}.
 \tag{N25.25}
\]
The budget bounds \(\bm s_n\).  After taking subsequences, the scaled
polynomials converge coefficientwise to
\[
 F_\infty(w)
 =\bigl((w+2)^2+s_{1,\infty}\bigr)
  \bigl((w-1)^2+s_{2,\infty}\bigr)
  \bigl((w-1)^2+s_{3,\infty}\bigr),
 \tag{N25.26}
\]
and their zeros are uniformly bounded.  Hence \(w_n\to w_\infty\in
\mathbb R\) and \(F_\infty(w_\infty)=0\).  There are two cases.

If \(w_\infty=-2\), then \(s_{1,\infty}=0\).  Put
\(y_n=z_n+2L_n=o(L_n)\) and factor off the two quadratics centered near
\(L_n\).  Their logarithmic derivative ratios through order four are
\(O(L_n^{-r})\), just as in Lemma~\ref{lem:f2v24-local-identity}.
Because \(t_{1,n}=o(L_n^2)\), that lemma gives
\[
 \frac{\mathcal Q_n(z_n)}{G_{1,n}(y_n)}
 =y_n^2+t_{1,n}-1+o(1)
 \geq\alpha_3-1+o(1)>0,
 \tag{N25.27}
\]
contradicting \(\mathcal Q_n(z_n)=0\).

If \(w_\infty=1\), put \(y_n=z_n-L_n=o(L_n)\).  The close-pair quartic
is \(A_{h_n,t_{2,n},t_{3,n}}(y_n)\), while the distant factor is
\[
 G_{L_n}(y_n)=(y_n+3L_n)^2+t_{1,n}.
 \tag{N25.28}
\]
Corollary~\ref{cor:f2v25-distant}, with
\(\alpha=\alpha_3>3/2\), gives
\(\mathcal Q_n(z_n)>0\), again a contradiction.  These are the only real
zeros of the nonnegative limiting product \textup{(N25.26)}.
\end{proof}

\subsection{Uniform completion of the reciprocal boundary}
\label{subsec:f2v25-completion}

\begin{theorem}[Finite cutoff for the full infinite-center boundary]
\label{thm:f2v25-infinite-cutoff}
There exists \(A_\infty<\infty\) such that the following holds.  Let
\[
 c_1+c_2+c_3=0,\qquad
 A_{\bm c}:=\frac12\sum_jc_j^2\geq A_\infty,
 \tag{N25.29}
\]
and suppose
\[
 t_j\geq\alpha_3\quad(1\leq j\leq3),\qquad
 \sum_jt_j\leq15+2A_{\bm c}.
 \tag{N25.30}
\]
Then the heat sextic \textup{(N22.1)} has no real zero.  Consequently it
is not real-rooted.
\end{theorem}

\begin{proof}
If no finite cutoff existed, choose a counterexample sequence with
\(A_n:=A_{\bm c_n}\to\infty\), and normalize
\[
 \lambda_n:=\sqrt{A_n},\qquad
 x_{j,n}:=\frac{c_{j,n}}{\lambda_n}.
 \tag{N25.31}
\]
Then \(\sum_jx_{j,n}=0\) and \(A_{\bm x_n}=1\).  Pass to a convergent
subsequence on this compact center circle.

If the three limiting centers are distinct, their pairwise gaps have a
fixed positive lower bound, and
Theorem~\ref{thm:f2v24-complete-separated} excludes the sequence.  If
two limiting centers coincide, relabel them as \(2,3\), and put
\[
 L_n:=\frac{c_{2,n}+c_{3,n}}2,\qquad
 h_n:=\frac{c_{2,n}-c_{3,n}}2.
 \tag{N25.32}
\]
Centering gives \(c_{1,n}=-2L_n\).  On the normalized circle a double
center has the form
\[
 \left(-\frac2{\sqrt3},\frac1{\sqrt3},\frac1{\sqrt3}\right)
 \tag{N25.33}
\]
up to reflection and permutation.  Hence, after a harmless reflection,
\[
 L_n\to\infty,\qquad h_n=o(L_n).
 \tag{N25.34}
\]
The second-moment budget in \textup{(N25.30)} is exactly
\[
 \sum_jt_{j,n}\leq15+6L_n^2+2h_n^2
 \tag{N25.35}
\]
by \textup{(N25.22)}.  Theorem~\ref{thm:f2v25-collision} excludes this
sequence as well.  All subsequential center shapes have been exhausted,
which proves the existence of \(A_\infty\).
\end{proof}

\begin{corollary}[Compact reduction of the three-pair threshold problem]
\label{cor:f2v25-compact-reduction}
To prove that real-rootedness of \(\mathcal Q_{\bm c,\bm t}\), with all
\(t_j\geq\alpha_3\), forces the symmetric equality configuration, it is
enough to work on the compact semialgebraic parameter set
\[
\begin{aligned}
 &c_1+c_2+c_3=0,\qquad 0\leq A_{\bm c}\leq A_\infty,\\
 &t_j\geq\alpha_3,\qquad
 t_1+t_2+t_3\leq15+2A_{\bm c}.
\end{aligned}
\tag{N25.36}
\]
\end{corollary}

\begin{proof}
The height set over \(0\leq A_{\bm c}\leq A_\infty\) is closed and
bounded by the second-moment budget.  Theorem
\ref{thm:f2v25-infinite-cutoff} excludes its complement in center
variance.
\end{proof}

\begin{firewall}[The remaining problem is now finite and compact]
\label{fire:f2v25}
The cutoff \(A_\infty\) is obtained by contradiction and is not yet
effective.  More importantly, Theorem~\ref{thm:f2v25-infinite-cutoff}
does not prove the desired threshold on the compact set
\textup{(N25.36)}.  A finite interior certificate cover, followed by the
infinite-product and arithmetic forcing steps, is still required.  GRH
remains open.
\end{firewall}

\begin{assessment}[V25 advance and V26 direction]
\label{ass:f2v25}
The reciprocal boundary is complete: V24 handles separated normalized
centers, and V25 handles their only collision stratum.  The key upstream
fact is that the two-pair local quartic is uniformly positive above
\(3/2\), well below the actual threshold \(\alpha_3\).

V26 should return to the compact finite interior without introducing a
fifth minor prematurely.  The exact open gap from V22 shows that
\((\mathcal C,\mathcal E,\mathcal K)\) is insufficient, while the
fourth minor \(\mathcal M_3\) repaired that gap.  The next task is to
compactify \textup{(N25.36)}, search specifically for a simultaneous
four-certificate gap, and, if none exists, construct a finite rational
box cover.  Boundary boxes should cite V13, V18, V21, V24, and V25 rather
than re-proving their strict regions.
\end{assessment}

\section*{Version V25 append-only record}
\addcontentsline{toc}{section}{Version V25 append-only record}

V25 was appended to the validated V24 active note whose installed
SHA--256 was
\[
 \texttt{7D5DB73D96CFB18A05FB13B9290C34BE440C4E0ADF8D04E8B6C53EBE5043241F}.
 \tag{N25.37}
\]
V25 proves a uniform close-pair quartic positivity theorem, lifts it
through the distant third factor, closes the normalized center-collision
chart, and deduces a finite cutoff for the complete infinite-center
boundary.  After validation V25 replaces V24 as the sole active numeric
GRH note; the frozen \texttt{V100\_f2} archive is unchanged.

% =============================================================================
% END APPENDED V25 MATERIAL
% =============================================================================

% =============================================================================
% BEGIN APPENDED V26 MATERIAL
% =============================================================================

\section{V26: global tubes around the solved finite skeleton}
\label{sec:f2v26-skeleton}

V25 reduces the arbitrary-height three-pair problem to finite center
variance, but its cutoff is nonconstructive.  A direct global Bernstein
cover would therefore still contain unnecessary boundary charts.  This
continuation first removes them topologically.

Three exact skeletons are already solved: the equal-height cone in V15,
the mirror-center family in V18, and the double-center family in V21.
V16 controls their common equality corner, while V25 controls infinity.
Together these results imply uniform open neighborhoods of all three
skeletons.  Thus any remaining finite counterexample lies in a compact
core separated from coincident centers, mirror centers, center collisions,
and equal heights.

\subsection{Closedness and the three transverse coordinates}
\label{subsec:f2v26-coordinates}

Let
\[
 \mathcal H_6:=
 \{\text{monic real-rooted polynomials of degree six}\}.
 \tag{N26.1}
\]
As used already in V13, V16, and V19, \(\mathcal H_6\) is closed in
coefficient space.  Indeed, coefficients converging in a bounded monic
family give a uniform root bound; any nonreal limiting root would be
approximated by nonreal roots, contrary to real-rootedness of the
approximants.

For centered centers, use the mirror-transverse coordinates
\[
 \bm c_{\rm mir}(d,\varepsilon)
 :=(-d-\varepsilon,\,2\varepsilon,\,d-\varepsilon),
 \qquad
 A_{\rm mir}=d^2+3\varepsilon^2,
 \tag{N26.2}
\]
and the double-center transverse coordinates
\[
 \bm c_{\rm dbl}(d,\varepsilon)
 :=(-2d,\,d+\varepsilon,\,d-\varepsilon),
 \qquad
 A_{\rm dbl}=3d^2+\varepsilon^2.
 \tag{N26.3}
\]
The solved skeletons are \(\varepsilon=0\).  For the third skeleton,
define the height spread
\[
 \operatorname{spr}(\bm t):=
 \max_jt_j-\min_jt_j.
 \tag{N26.4}
\]

\begin{theorem}[Uniform tubes around all solved skeletons]
\label{thm:f2v26-tubes}
There exist positive constants
\[
 \eta_{\rm mir},\qquad \eta_{\rm dbl},\qquad \eta_{\rm ht}
 \tag{N26.5}
\]
such that the following statements hold whenever \(t_j\geq\alpha_3\).

\begin{enumerate}[label=\textup{(\roman*)}]
\item If
\[
 \bm c=\bm c_{\rm mir}(d,\varepsilon),
 \qquad |\varepsilon|\leq\eta_{\rm mir},
 \tag{N26.6}
\]
then \(\mathcal Q_{\bm c,\bm t}\) can be real-rooted only at
\[
 d=\varepsilon=0,\qquad
 t_1=t_2=t_3=\alpha_3.
 \tag{N26.7}
\]

\item If
\[
 \bm c=\bm c_{\rm dbl}(d,\varepsilon),
 \qquad |\varepsilon|\leq\eta_{\rm dbl},
 \tag{N26.8}
\]
then the same conclusion \textup{(N26.7)} holds.

\item For arbitrary centered \(\bm c\), if
\[
 \operatorname{spr}(\bm t)\leq\eta_{\rm ht},
 \tag{N26.9}
\]
then real-rootedness again forces \textup{(N26.7)}, with
\(\varepsilon\) omitted; equivalently,
\[
 c_1=c_2=c_3=0,\qquad
 t_1=t_2=t_3=\alpha_3.
 \tag{N26.10}
\]
\end{enumerate}
\end{theorem}

\begin{proof}
We use the same compactness scheme three times.

For part \textup{(i)}, suppose no uniform \(\eta_{\rm mir}\) exists.
There are nonexceptional real-rooted examples with
\(\varepsilon_n\to0\).  Real-rootedness gives the automatic
second-moment budget
\[
 \sum_jt_{j,n}\leq15+2A_{{\rm mir},n}.
 \tag{N26.11}
\]
Theorem~\ref{thm:f2v25-infinite-cutoff} bounds \(A_{{\rm mir},n}\);
otherwise the sextic would have no real zero.  Consequently \(d_n\) and
all three heights are bounded.  Pass to a subsequence,
\[
 d_n\to d_\ast,\qquad \bm t_n\to\bm t_\ast.
 \tag{N26.12}
\]
Continuity of the heat polynomial and closedness of \(\mathcal H_6\)
show that the mirror-center limit
\(\mathcal Q_{\bm c_{\rm mir}(d_\ast,0),\bm t_\ast}\) is real-rooted.
Theorem~\ref{thm:f2v18-mirror-all-heights}, together with
\(t_{\ast,j}\geq\alpha_3\), forces
\[
 d_\ast=0,\qquad
 \bm t_\ast=(\alpha_3,\alpha_3,\alpha_3).
 \tag{N26.13}
\]
Hence \(A_{{\rm mir},n}\to0\).  The local bridge
Theorem~\ref{thm:f2v16-local-bridge} then forces every sufficiently late
example to be precisely \textup{(N26.7)}, contradicting its selection as
nonexceptional.  This proves \textup{(i)}.

For part \textup{(ii)}, repeat the argument with
\(A_{{\rm dbl},n}=3d_n^2+\varepsilon_n^2\).  At
\(\varepsilon=0\), Theorem
\ref{thm:f2v21-double-center-all-heights} replaces V18 and again forces
\textup{(N26.13)}.  V16 excludes the resulting nonexceptional sequence
near \(A=0\).

For part \textup{(iii)}, suppose there are nonexceptional real-rooted
examples with \(\operatorname{spr}(\bm t_n)\to0\).  V25 bounds their
center variances, and the second-moment budget then bounds the heights.
After taking a subsequence,
\[
 \bm c_n\to\bm c_\ast,\qquad
 \bm t_n\to(t_\ast,t_\ast,t_\ast),
 \qquad t_\ast\geq\alpha_3.
 \tag{N26.14}
\]
The limiting equal-height sextic is real-rooted.  Theorem
\ref{thm:f2v15-arbitrary-centers} gives
\[
 t_\ast=\alpha_3,\qquad \bm c_\ast=\bm0.
 \tag{N26.15}
\]
Thus \(A_{\bm c_n}\to0\), and V16 again forces the exact exceptional
configuration for all sufficiently large \(n\), a contradiction.
\end{proof}

\subsection{An explicit compact core for the remaining search}
\label{subsec:f2v26-core}

Order the centers
\[
 a\leq b\leq c,\qquad a+b+c=0,
 \tag{N26.16}
\]
and retain
\[
 A_{\bm c}=\frac12(a^2+b^2+c^2).
 \tag{N26.17}
\]
In these coordinates, \(|b|/2\) is the mirror-transverse parameter in
\textup{(N26.2)}, while half of either adjacent center gap is a
double-center transverse parameter after permutation.

\begin{corollary}[Compact nondegenerate core]
\label{cor:f2v26-core}
There exist constants
\[
 0<A_0<A_1<\infty,\qquad
 \eta_c>0,\qquad \eta_t>0
 \tag{N26.18}
\]
such that every nonexceptional real-rooted sextic with
\(t_j\geq\alpha_3\) would have to satisfy
\[
\begin{aligned}
 &A_0\leq A_{\bm c}\leq A_1,\\
 &|b|\geq\eta_c,\qquad
 b-a\geq\eta_c,\qquad c-b\geq\eta_c,\\
 &t_j\geq\alpha_3,\qquad
 \sum_jt_j\leq15+2A_{\bm c},\\
 &\max_jt_j-\min_jt_j\geq\eta_t.
\end{aligned}
\tag{N26.19}
\]
The parameter set defined by \textup{(N26.16)} and
\textup{(N26.19)} is compact.
\end{corollary}

\begin{proof}
Take \(A_0>0\) smaller than the local constant in
Theorem~\ref{thm:f2v16-local-bridge}, and take
\(A_1=A_\infty\) from Theorem~\ref{thm:f2v25-infinite-cutoff}.
Set
\[
 \eta_c:=\min\{2\eta_{\rm mir},2\eta_{\rm dbl}\},
 \qquad \eta_t:=\eta_{\rm ht}.
 \tag{N26.20}
\]
If \(|b|<\eta_c\), use
\(\varepsilon=b/2\) in \textup{(N26.2)}.  If either adjacent gap is
less than \(\eta_c\), use half that gap in
\textup{(N26.3)} after a permutation and, if needed, reflection.
Theorem~\ref{thm:f2v26-tubes} excludes each case.  Its third part excludes
height spread below \(\eta_t\).  V16 and V25 supply the two variance
bounds.  The remaining inequalities are closed; they bound the centers
and, through the second-moment budget, the heights.  Hence the core is
compact.
\end{proof}

\subsection{Four-certificate search ledger}
\label{subsec:f2v26-search}

The certificates
\[
 (\mathcal C,\mathcal E,\mathcal K,\mathcal M_3)
 \tag{N26.21}
\]
were evaluated at \(6{,}000{,}000\) additional points.  The search used
the exact parameterization
\[
 \bm c=\sqrt A\,\bm x,\qquad A_{\bm x}=1,\qquad
 t_j=\alpha_3+\delta_j,\qquad
 \delta_j\geq0,\qquad
 \sum_j\delta_j\leq15+2A-3\alpha_3,
 \tag{N26.22}
\]
with logarithmic and finite \(A\)-regimes, bulk simplex allocations,
near-threshold and near-budget allocations, and one-low-height faces.
Before evaluating Newton sums, each sextic was rescaled by
\[
 R^2:=1+A+\sum_jt_j
 \tag{N26.23}
\]
so all arithmetic remained bounded.  Undoing the positive powers of
\(R\) reproduces the V22 integer witness as
\[
\begin{aligned}
 \mathcal C&=1572848,&
 \mathcal E&=20841710115526400,\\
 \mathcal K&=852566236,&
 \mathcal M_3&=-386174231712,&
 p_2&=1714.
\end{aligned}
\tag{N26.24}
\]
The search counts were
\[
\begin{array}{c|r|r|r}
\text{regime}&\mathcal C,\mathcal E\geq0&
 \mathcal C,\mathcal E,\mathcal K\geq0&
 \text{all four}\geq0\\ \hline
\text{logarithmic \(A\), bulk}&847&88&0\\
\text{finite \(A\), bulk}&1330&101&0\\
\text{logarithmic \(A\), near threshold}&287457&287395&0\\
\text{finite \(A\), near threshold}&191705&191609&0\\
\text{logarithmic \(A\), near budget}&897&186&0\\
\text{finite \(A\), one low height}&25401&14064&0\\ \hline
\text{total}&507637&493443&0
\end{array}
\tag{N26.25}
\]
These counts are discovery evidence only.  In particular, failure to find
a four-certificate gap is not a proof that the four signs cover the core
\textup{(N26.19)}.

\begin{firewall}[The core is compact but the cover is not certified]
\label{fire:f2v26}
The constants in Theorem~\ref{thm:f2v26-tubes} and
Corollary~\ref{cor:f2v26-core} are existential.  No rational box cover of
the remaining core has yet been constructed, and the numerical ledger
\textup{(N26.25)} cannot replace one.  Even a complete finite three-pair
threshold theorem would still require an infinite-product passage and
the arithmetic forcing step.  GRH remains open.
\end{firewall}

\begin{assessment}[V26 advance and V27 direction]
\label{ass:f2v26}
The finite search no longer needs to spend boxes near any solved
degeneracy.  A putative counterexample must have intermediate center
variance, genuinely nonmirror and noncolliding center shape, and a
definite height spread.

V27 should make one of the existential tube constants effective.  The
most useful target is the equal-height tube, because its transverse
variables are only the two height differences and V15 already supplies
exact strict certificate margins on three center charts.  Expanding the
V15 certificates to second order in the height differences and retaining
the V26 compact center core may yield a rational \(\eta_t\).  If that
expansion is too large, use interval Bernstein subdivision only on the
two transverse variables, with the inherited V15 chart variables left
symbolic.
\end{assessment}

\section*{Version V26 append-only record}
\addcontentsline{toc}{section}{Version V26 append-only record}

V26 was appended to the validated V25 active note whose installed
SHA--256 was
\[
 \texttt{456CCED92EF54F670AF3B03E1E4B3A84DB7A96EEDEE23BA765CCC2F24F44499F}.
 \tag{N26.26}
\]
V26 globalizes the solved mirror, double-center, and equal-height
skeletons to uniform neighborhoods, reduces every remaining putative
counterexample to a compact nondegenerate core, and records a normalized
six-million-point four-certificate search.  After validation V26 replaces
V25 as the sole active numeric GRH note; the frozen
\texttt{V100\_f2} archive is unchanged.

% =============================================================================
% END APPENDED V26 MATERIAL
% =============================================================================

% =============================================================================
% BEGIN APPENDED V27 MATERIAL
% =============================================================================

\section{V27: effective transverse widths near equal heights}
\label{sec:f2v27-effective-height}

V26 proved, by compactness, that a uniform neighborhood of the
equal-height cone contains no nontrivial counterexample.  That theorem did
not produce a numerical width.  This continuation makes two substantial
parts of the neighborhood effective.  The constants are deliberately
coarse: their purpose is a reproducible rational certificate, not numerical
optimization.

The high-height chart uses the absolute discriminant margin in the V15
Bernstein ledger.  The low-height chart uses the negative adapted-square
margin away from \(A=0\), together with an explicit coefficient-to-Newton
Lipschitz recursion.  No floating-point sign decision enters either proof.

\subsection{Uniform coefficient and Newton-sum perturbation bounds}
\label{subsec:f2v27-lipschitz}

Assume throughout this section that
\[
 c_1+c_2+c_3=0,\qquad
 A_{\bm c}:=\frac12\sum_jc_j^2\leq3,\qquad
 \alpha_3\leq t_j\leq\frac{13}{4}.
 \tag{N27.1}
\]
Then \(|c_j|<3\).  Each quadratic factor
\[
 (z-c_j)^2+t_j
 \tag{N27.2}
\]
has coefficient \(\ell^1\)-norm at most
\[
 1+6+13=20.
 \tag{N27.3}
\]
Consequently every coefficient of their product has absolute value at
most \(20^3=8000\).

For a degree-six polynomial, the coefficient row sums of
\(e^{-D^2/2}\), in ascending powers of \(z\), are
\[
 20,\ 19,\ 52,\ 11,\ 16,\ 1,\ 1.
 \tag{N27.4}
\]
Thus every coefficient of the heat sextic is bounded by
\[
 M:=52\cdot8000=416000.
 \tag{N27.5}
\]

\begin{lemma}[Height-to-coefficient Lipschitz bound]
\label{lem:f2v27-coefficient-bound}
Fix the centers in \textup{(N27.1)}.  Let
\(\bm t,\bm t^{\,0}\in[\alpha_3,13/4]^3\) satisfy
\[
 |t_j-t_j^0|\leq\eta\qquad(1\leq j\leq3).
 \tag{N27.6}
\]
If
\[
 \mathcal Q(z)=z^6+e_1z^5+\cdots+e_6,\qquad
 \mathcal Q^0(z)=z^6+e_1^0z^5+\cdots+e_6^0
 \tag{N27.7}
\]
are the corresponding heat sextics, then
\[
 |e_j|,\ |e_j^0|\leq M,\qquad
 |e_j-e_j^0|\leq D\eta,\qquad
 D:=62400.
 \tag{N27.8}
\]
\end{lemma}

\begin{proof}
Differentiate the unheated product with respect to one height.  The
result is the product of the other two quadratic factors, whose
coefficient \(\ell^1\)-norm is at most \(20^2=400\).  Move from
\(\bm t^0\) to \(\bm t\) one coordinate at a time.  Equation
\textup{(N27.6)} gives a coefficient change at most
\[
 3\cdot400\eta=1200\eta
 \tag{N27.9}
\]
before applying heat.  The largest row sum in \textup{(N27.4)} is \(52\),
so the heat-image coefficient change is at most
\(52\cdot1200\eta=62400\eta\).  The absolute bound is
\textup{(N27.5)}.
\end{proof}

Define integers \(P_k,L_k\) recursively by
\[
 P_0:=6,\qquad L_0:=0,
 \tag{N27.10}
\]
and, for \(1\leq k\leq8\),
\[
\begin{aligned}
 P_k&:=M\sum_{j=1}^{\min(k-1,6)}P_{k-j}
      +{\bf1}_{k\leq6}\,kM,\\
 L_k&:=\sum_{j=1}^{\min(k-1,6)}
      \bigl(DP_{k-j}+ML_{k-j}\bigr)
      +{\bf1}_{k\leq6}\,kD.
\end{aligned}
\tag{N27.11}
\]
Empty sums are zero.

\begin{lemma}[Explicit Newton-sum Lipschitz recursion]
\label{lem:f2v27-newton-bound}
Let \(p_k,p_k^0\) be the Newton sums of
\(\mathcal Q,\mathcal Q^0\) in
Lemma~\ref{lem:f2v27-coefficient-bound}.  Then
\[
 |p_k|,\ |p_k^0|\leq P_k,\qquad
 |p_k-p_k^0|\leq L_k\eta
 \qquad(0\leq k\leq8).
 \tag{N27.12}
\]
In particular, if \(\mathcal S_{\rm gen}\) is the adapted square
\textup{(N16.7)}, then
\[
 |\mathcal S_{\rm gen}(\bm c,\bm t)
  -\mathcal S_{\rm gen}(\bm c,\bm t^0)|
 \leq L_{\mathcal S}\eta,
 \tag{N27.13}
\]
where
\[
\begin{aligned}
 L_{\mathcal S}
 &:=L_8+10L_6+27L_4+10L_2\\
 &=1076304606215638788825305318892147240357331200.
\end{aligned}
\tag{N27.14}
\]
\end{lemma}

\begin{proof}
Newton's recurrence is
\[
 p_k=-\sum_{j=1}^{\min(k-1,6)}e_jp_{k-j}
      -{\bf1}_{k\leq6}\,ke_k.
 \tag{N27.15}
\]
The first inequality in \textup{(N27.12)} follows by induction from
\(|e_j|\leq M\) and the definition of \(P_k\).  Subtract the two
recurrences and use
\[
 |e_jp_{k-j}-e_j^0p_{k-j}^0|
 \leq M|p_{k-j}-p_{k-j}^0|
      +D\eta\,P_{k-j}
 \tag{N27.16}
\]
to obtain the \(L_k\) recursion.

The rational enclosures \textup{(N13.17)} imply, for the constants
\(c,d\) in \textup{(N13.10)},
\[
 |c|<5,\qquad |d|<1.
 \tag{N27.17}
\]
Apply \textup{(N27.12)} to
\[
 \mathcal S_{\rm gen}
 =p_8-2cp_6+(c^2+2d)p_4-2cdp_2+6d^2
 \tag{N27.18}
\]
to get the first line of \textup{(N27.14)}.  Direct evaluation of the
integer recursion \textup{(N27.11)} gives its second line.
\end{proof}

\subsection{An effective low-height tube away from zero center variance}
\label{subsec:f2v27-low}

\begin{theorem}[Effective low equal-height tube]
\label{thm:f2v27-low-tube}
Fix any rational \(A_-\) with \(0<A_-\leq3\), and set
\[
 \eta_{\rm lo}(A_-)
 :=\min\left\{1,\frac{197A_-}{6L_{\mathcal S}}\right\}.
 \tag{N27.19}
\]
Suppose
\[
\begin{gathered}
 c_1+c_2+c_3=0,\qquad
 A_-\leq A_{\bm c}\leq3,\\
 \alpha_3\leq\min_jt_j\leq\frac94,\qquad
 \max_jt_j-\min_jt_j\leq\eta_{\rm lo}(A_-).
\end{gathered}
\tag{N27.20}
\]
Then \(\mathcal Q_{\bm c,\bm t}\) is not real-rooted.
\end{theorem}

\begin{proof}
Put \(t:=\min_jt_j\) and
\(\bm t^0:=(t,t,t)\).  The height bound and
\(\eta_{\rm lo}\leq1\) put both height triples inside
\textup{(N27.1)}.

Use the V15 low chart with
\[
 X:=\frac{A_{\bm c}}3,\qquad
 U:=\frac{t-\alpha_3}{9/4-\alpha_3}.
 \tag{N27.21}
\]
In the Bernstein ledger \textup{(N15.9)}, the \(X\)-row \(i=0\) is
nonpositive, and every coefficient in rows \(i\geq1\) is at most
\(-197\).  The total Bernstein weight of those rows is
\[
 1-(1-X)^4\geq X\geq\frac{A_-}{3}.
 \tag{N27.22}
\]
Therefore the equal-height adapted square satisfies the exact margin
\[
 \mathcal S_{\rm gen}(\bm c,\bm t^0)
 \leq-\frac{197A_-}{3}.
 \tag{N27.23}
\]
Lemma~\ref{lem:f2v27-newton-bound} and
\textup{(N27.19)} give
\[
 \mathcal S_{\rm gen}(\bm c,\bm t)
 \leq-\frac{197A_-}{3}
       L_{\mathcal S}\eta_{\rm lo}(A_-)
 \leq-\frac{197A_-}{6}<0.
 \tag{N27.24}
\]
For six real roots the adapted square is a sum of six real squares, so it
must be nonnegative.  This contradiction proves the theorem.
\end{proof}

\subsection{An absolute effective high-height tube}
\label{subsec:f2v27-high}

Set
\[
 B_0:=2500000,\qquad D_1:=374400,
 \tag{N27.25}
\]
and define the rational number
\[
\begin{aligned}
 \eta_{\rm hi}
 &:=\frac{2985984\cdot837}
 {2\cdot11\cdot11!\cdot D_1 B_0^{10}}
 =7.9707\ldots\,10^{-70}.
\end{aligned}
\tag{N27.26}
\]

\begin{lemma}[Sylvester determinant perturbation]
\label{lem:f2v27-sylvester}
Under \textup{(N27.1)} and \textup{(N27.6)}, let
\(\Delta,\Delta^0\) be the discriminants of
\(\mathcal Q,\mathcal Q^0\).  Then
\[
 |\Delta-\Delta^0|
 \leq11\cdot11!\,D_1B_0^{10}\eta.
 \tag{N27.27}
\]
\end{lemma}

\begin{proof}
The discriminant, up to sign, is the determinant of the
\(11\)-dimensional Sylvester matrix of \(\mathcal Q\) and
\(\mathcal Q'\).  By \textup{(N27.5)}, every entry is at most
\[
 6M=2496000<B_0
 \tag{N27.28}
\]
in absolute value.  Lemma~\ref{lem:f2v27-coefficient-bound} bounds every
coefficient change in \(\mathcal Q'\) by
\[
 6D\eta=D_1\eta.
 \tag{N27.29}
\]
Expand each determinant by permutations and telescope every product of
eleven entries.  Each product changes by at most
\(11D_1B_0^{10}\eta\), and there are \(11!\) permutations.  This proves
\textup{(N27.27)}.
\end{proof}

\begin{theorem}[Absolute effective high equal-height tube]
\label{thm:f2v27-high-tube}
Suppose
\[
\begin{gathered}
 c_1+c_2+c_3=0,\qquad A_{\bm c}\leq3,\\
 \min_jt_j\geq\frac94,\qquad
 \max_jt_j-\min_jt_j\leq\eta_{\rm hi}.
\end{gathered}
\tag{N27.30}
\]
Then \(\mathcal Q_{\bm c,\bm t}\) is not real-rooted.
\end{theorem}

\begin{proof}
If the sextic were real-rooted, the V9 pure three-pair bound would give
\[
 t:=\min_jt_j\leq3.
 \tag{N27.31}
\]
Put \(\bm t^0=(t,t,t)\).  Since \(\eta_{\rm hi}<1/4\), the two height
triples obey \textup{(N27.1)}.  In the V15 high chart,
\(X=A_{\bm c}/3\in[0,1]\) and
\(U=(t-9/4)/(3/4)\in[0,1]\).  Every Bernstein coefficient of
\(\widehat{\mathcal P}_{\rm hi}\) is at least \(837\) by
\textup{(N15.11)}.  As
\(\widehat{\mathcal P}=-\operatorname{Disc}\mathcal Q/2985984\),
\[
 \operatorname{Disc}\mathcal Q_{\bm c,\bm t^0}
 \leq-2985984\cdot837.
 \tag{N27.32}
\]
Lemma~\ref{lem:f2v27-sylvester} and the definition
\textup{(N27.26)} show that the perturbation changes the discriminant by
at most half this margin.  Hence
\[
 \operatorname{Disc}\mathcal Q_{\bm c,\bm t}<0,
 \tag{N27.33}
\]
which is impossible for a real-rooted real polynomial.
\end{proof}

\begin{corollary}[Effective moderate-variance height separation]
\label{cor:f2v27-effective-core}
Fix \(A_-\in(0,3]\).  A putative real-rooted sextic with
\[
 A_-\leq A_{\bm c}\leq3,\qquad t_j\geq\alpha_3
 \tag{N27.34}
\]
must satisfy
\[
 \operatorname{spr}(\bm t)>
 \min\{\eta_{\rm lo}(A_-),\eta_{\rm hi}\}.
 \tag{N27.35}
\]
\end{corollary}

\begin{proof}
The minimum height is either at most \(9/4\), in which case
Theorem~\ref{thm:f2v27-low-tube} applies, or at least \(9/4\), in which
case Theorem~\ref{thm:f2v27-high-tube} applies.
\end{proof}

\begin{firewall}[Effective does not yet mean global]
\label{fire:f2v27}
Theorem~\ref{thm:f2v27-low-tube} requires a chosen positive lower bound
\(A_-\), and both effective theorems currently stop at \(A_{\bm c}=3\).
The large-separation V15 chart has a positive compactified discriminant
ledger, but its unequal-height perturbation must be scaled before the
same determinant estimate can be used.  The complete four-certificate
cover of the V26 core, the infinite-product passage, and the arithmetic
forcing step remain open.  GRH remains open.
\end{firewall}

\begin{assessment}[V27 advance and V28 direction]
\label{ass:f2v27}
V27 replaces a portion of the existential height tube by explicit
rational widths.  The constants are extremely small because the proof
uses global coefficient and Sylvester bounds, but every dependency is now
auditable.

V28 should treat \(A_{\bm c}\geq3\) by using the reciprocal variable
\(Y=3/A_{\bm c}\) from the V15 large-separation chart and scaling height
differences by \(1+A_{\bm c}\).  The positive rows in
\textup{(N15.14)} should then yield an effective relative height-spread
tube that matches the V24--V25 boundary.  This is preferable to applying
the unscaled determinant bound, whose coefficients grow with \(A\).
\end{assessment}

\section*{Version V27 append-only record}
\addcontentsline{toc}{section}{Version V27 append-only record}

V27 was appended to the validated V26 active note whose installed
SHA--256 was
\[
 \texttt{FEAC936EAD1F323691B929D67B6442845BCDA297B46B8A99D4826FF38B902C4B}.
 \tag{N27.36}
\]
V27 gives explicit coefficient and Newton-sum Lipschitz bounds and proves
effective low- and high-height tubes around the equal-height cone for
moderate center variance.  After validation V27 replaces V26 as the sole
active numeric GRH note; the frozen \texttt{V100\_f2} archive is unchanged.

% =============================================================================
% END APPENDED V27 MATERIAL
% =============================================================================

% =============================================================================
% BEGIN APPENDED V28 MATERIAL
% =============================================================================

\section{V28: an effective fourth-minor reciprocal horn}
\label{sec:f2v28-effective-reciprocal}

V27 proposed perturbing the compactified V15 discriminant at large center
variance.  A crude Sylvester bound in that chart loses the cancellations
responsible for the finite reciprocal limit and gives no honest uniform
width.  This continuation therefore returns to the exact fourth-minor
expansion from V23.  Its three remainder polynomials admit small explicit
coefficient norms, making the sublinear-height reciprocal theorem fully
effective.

The resulting bound is simple: if the normalized center scale is
\(\lambda\) and every squared height is at most \(T\), then
\(\lambda\geq171T\) forces the principal minor \(\mathcal M_3\) to be
negative.  Thus every still-unexcluded real-rooted configuration lies
above an explicit height-growth horn.

\subsection{Exact coefficient-norm audit}
\label{subsec:f2v28-norm}

Retain the V23 normalization
\[
 x_1=x,\qquad x_2=y,\qquad x_3=-x-y,\qquad
 A_{\bm x}=x^2+xy+y^2=1,
 \tag{N28.1}
\]
and let \(c_j=\lambda x_j\).  Equation \textup{(N23.8)} is
\[
\begin{aligned}
 \mathcal M_3(\lambda\bm x,\bm t)
={}&-48\lambda^6
 \left(\sum_{j=1}^3t_j\Delta_j^2-9\right)\\
 &+\lambda^4R_4(\bm x,\bm t)
  +\lambda^2R_2(\bm x,\bm t)+R_0(\bm t),
\end{aligned}
\tag{N28.2}
\]
where \(\Delta_j=P_{\bm x}'(x_j)\).  The factor
\(A_{\bm x}=1\) has been inserted.

For a polynomial
\[
 R(x,y,\bm t)=
 \sum_{\nu}r_\nu x^{a_\nu}y^{b_\nu}
 t_1^{d_{1,\nu}}t_2^{d_{2,\nu}}t_3^{d_{3,\nu}},
 \tag{N28.3}
\]
define the weighted coefficient norm
\[
 \|R\|_{2}:=
 \sum_\nu |r_\nu|\,2^{a_\nu+b_\nu}.
 \tag{N28.4}
\]
Since \(\sum_jx_j^2=2\), every normalized center satisfies
\(|x|,|y|<2\).

\begin{lemma}[Exact fourth-minor remainder ledger]
\label{lem:f2v28-remainders}
The three remainder polynomials in \textup{(N28.2)} have the following
exact audit:
\[
\begin{array}{c|c|c|r}
\text{polynomial}&\text{center degree}&
\text{maximum height degree}&\text{number of monomials}\\ \hline
R_4&4&2&50\\
R_2&2&3&60\\
R_0&0&4&29
\end{array}
\tag{N28.5}
\]
and
\[
 \|R_4\|_2=1973760,\qquad
 \|R_2\|_2=823680,\qquad
 \|R_0\|_2=131616.
 \tag{N28.6}
\]
Consequently, for \(T\geq1\) and
\(\alpha_3\leq t_j\leq T\),
\[
\begin{aligned}
 |R_4|&\leq1973760\,T^2,\\
 |R_2|&\leq823680\,T^3,\\
 |R_0|&\leq131616\,T^4.
\end{aligned}
\tag{N28.7}
\]
\end{lemma}

\begin{proof}
Use the exact Newton sums \textup{(N23.11a)} and form
\[
 \mathcal M_3=6p_2p_6-6p_4^2-p_2p_3^2.
 \tag{N28.8}
\]
Expand in the polynomial ring
\(\mathbb Z[x,y,t_1,t_2,t_3,\lambda]\), collect powers of
\(\lambda\), and substitute \(x_3=-x-y\).  The coefficients of
\(\lambda^8,\lambda^7,\lambda^5,\lambda^3,\lambda\) vanish identically.
The coefficient of \(\lambda^6\) is the leading term in
\textup{(N28.2)}.  Counting the nonzero monomials at powers
\(\lambda^4,\lambda^2,\lambda^0\) gives \textup{(N28.5)}.

For each of those powers, apply the finite exact sum
\[
 \sum_\nu |r_\nu|\,2^{a_\nu+b_\nu}.
 \tag{N28.9}
\]
Integer arithmetic gives, respectively,
\[
 1973760,\qquad823680,\qquad131616,
 \tag{N28.10}
\]
which is \textup{(N28.6)}.  Finally
\(|x|,|y|<2\), \(0\leq t_j\leq T\), and \(T\geq1\) bound every monomial
of height degree at most \(d\) by its weighted coefficient times \(T^d\).
This proves \textup{(N28.7)}.
\end{proof}

\subsection{An explicit scale threshold}
\label{subsec:f2v28-threshold}

The rational enclosure \textup{(N13.17)} gives
\[
 \alpha_3-1>\frac{4717}{5000}.
 \tag{N28.11}
\]
Set
\[
 \gamma:=432\frac{4717}{5000}
 =\frac{254718}{625}.
 \tag{N28.12}
\]

\begin{theorem}[Effective reciprocal fourth-minor exclusion]
\label{thm:f2v28-effective}
Let \(T\geq1\), let \(A_{\bm x}=1\), and suppose
\[
 \alpha_3\leq t_j\leq T\qquad(1\leq j\leq3).
 \tag{N28.13}
\]
Define
\[
\begin{aligned}
 \Lambda(T):=\max\Bigg\{1,\,
 &\left(\frac{6\cdot1973760}{\gamma}\right)^{1/2}T,\\
 &\left(\frac{6\cdot823680}{\gamma}\right)^{1/4}T^{3/4},\\
 &\left(\frac{6\cdot131616}{\gamma}\right)^{1/6}T^{2/3}
 \Bigg\}.
\end{aligned}
\tag{N28.14}
\]
If \(\lambda\geq\Lambda(T)\), then
\[
 \mathcal M_3(\lambda\bm x,\bm t)<0.
 \tag{N28.15}
\]
Hence the corresponding heat sextic is not real-rooted.
\end{theorem}

\begin{proof}
The derivative-square identity \textup{(N23.4)} gives
\[
 \sum_jt_j\Delta_j^2-9
 \geq9(\alpha_3-1).
 \tag{N28.16}
\]
Therefore the leading term in \textup{(N28.2)} is at most
\[
 -432(\alpha_3-1)\lambda^6<-\gamma\lambda^6.
 \tag{N28.17}
\]
Each of the three nontrivial lower bounds in
\textup{(N28.14)}, together with \textup{(N28.7)}, gives
\[
\begin{aligned}
 1973760\,T^2\lambda^4&\leq\frac{\gamma}{6}\lambda^6,\\
 823680\,T^3\lambda^2&\leq\frac{\gamma}{6}\lambda^6,\\
 131616\,T^4&\leq\frac{\gamma}{6}\lambda^6.
\end{aligned}
\tag{N28.18}
\]
Thus the absolute value of the total remainder is at most
\(\gamma\lambda^6/2\).  Equation \textup{(N28.17)} proves
\textup{(N28.15)}.  A principal moment minor of six real roots must be
nonnegative.
\end{proof}

\begin{corollary}[The constant \(171\) horn]
\label{cor:f2v28-171}
Under the hypotheses of
Theorem~\ref{thm:f2v28-effective}, the simpler condition
\[
 \lambda\geq171T
 \tag{N28.19}
\]
implies \(\mathcal M_3<0\).
\end{corollary}

\begin{proof}
The three numerical coefficients in \textup{(N28.14)} are, respectively,
\[
 170.4642\ldots,\qquad
 10.4938\ldots,\qquad
 3.5309\ldots.
 \tag{N28.20}
\]
These decimal comparisons can be checked without approximation by raising
\(171\), \(11\), and \(4\) to the corresponding powers and comparing the
resulting integers with the rational quantities in
\textup{(N28.14)}.  Since \(T\geq1\),
\[
 171T>
 170.4642\,T,\quad
 11T\geq11T^{3/4},\quad
 4T\geq4T^{2/3},
 \tag{N28.21}
\]
so \textup{(N28.19)} dominates every entry of
\(\Lambda(T)\).
\end{proof}

\begin{corollary}[Effective height-growth necessity]
\label{cor:f2v28-growth}
Let the centered centers have variance \(A_{\bm c}>0\), and suppose all
three squared heights are at least \(\alpha_3\).  Put
\[
 T:=\max_jt_j.
 \tag{N28.22}
\]
If the heat sextic is real-rooted, then necessarily
\[
 T>\frac{\sqrt{A_{\bm c}}}{171}.
 \tag{N28.23}
\]
Equivalently, the region
\[
 A_{\bm c}\geq29241\,T^2
 \tag{N28.24}
\]
is excluded by \(\mathcal M_3\).
\end{corollary}

\begin{proof}
Write \(\lambda=\sqrt{A_{\bm c}}\) and
\(\bm x=\bm c/\lambda\), so \(A_{\bm x}=1\).
If \(\lambda\geq171T\), Corollary~\ref{cor:f2v28-171} contradicts
real-rootedness.
\end{proof}

\begin{firewall}[The horn does not replace the compact cover]
\label{fire:f2v28}
The effective horn controls heights whose maximum grows at most linearly
with the center scale.  The second-moment budget permits squared heights
of order \(A_{\bm c}=\lambda^2\), so the region above
\textup{(N28.23)} still contains a large parameter range.  V24--V25
exclude that range asymptotically but non-effectively.  The finite
four-certificate cover, infinite-product passage, and arithmetic forcing
step remain open.  GRH remains open.
\end{firewall}

\begin{assessment}[V28 advance and V29 direction]
\label{ass:f2v28}
V28 converts the qualitative \(T=o(\lambda)\) theorem of V23 into a
uniform, explicit fixed-slope horn.  The coefficient audit also identifies
the \(\lambda^4R_4\) term as the quantitative bottleneck.

V29 should inspect the projective complement \(T>\lambda/171\).  Normalize
\(t_j=\lambda^2s_j\) and use the exact discriminant factorization
\textup{(N24.9)} on the stratum where all \(s_j\) and all pair resultants
have explicit positive lower bounds.  A coefficientwise perturbation of
the scaled heat polynomial then gives an effective cutoff on that generic
projective stratum.  Mixed height faces and repeated quadratic factors
should remain assigned to the exact local identities in V24--V25.
\end{assessment}

\section*{Version V28 append-only record}
\addcontentsline{toc}{section}{Version V28 append-only record}

V28 was appended to the validated V27 active note whose installed
SHA--256 was
\[
 \texttt{379BC02950CB7D321DC6E06E82E6F4D10A97A3A32A9C478067A8FC8C19BB005F}.
 \tag{N28.25}
\]
V28 gives exact coefficient norms for every remainder in the reciprocal
fourth-minor expansion and proves the explicit exclusion
\(\sqrt{A_{\bm c}}\geq171\max_jt_j\).  After validation V28 replaces V27
as the sole active numeric GRH note; the frozen \texttt{V100\_f2} archive
is unchanged.

% =============================================================================
% END APPENDED V28 MATERIAL
% =============================================================================

% =============================================================================
% BEGIN APPENDED V29 MATERIAL
% =============================================================================

\section{V29: an effective generic projective cutoff}
\label{sec:f2v29-projective-cutoff}

V28 gives an effective exclusion when the largest squared height is at
most a fixed multiple of the center scale.  The opposite regime contains
projective heights of order the center variance.  V24 excluded its generic
interior asymptotically by the exact discriminant factorization, but did
not give a cutoff.

This continuation makes that generic projective exclusion effective.
When every scaled height and every pair resultant has a prescribed
positive lower bound, the projective input has a quantitative negative
discriminant.  A completely explicit Sylvester-matrix estimate shows how
large the center scale must be for backward heat to preserve that sign.

\subsection{Scaled coefficient bounds}
\label{subsec:f2v29-coefficients}

Let
\[
 A_{\bm x}:=\frac12\sum_{j=1}^3x_j^2=1,\qquad
 \sum_jx_j=0,
 \tag{N29.1}
\]
and define
\[
 F_{\bm x,\bm s}(w)
 :=\prod_{j=1}^3\bigl((w-x_j)^2+s_j\bigr).
 \tag{N29.2}
\]
For \(\lambda>0\), put
\[
 \varepsilon:=\lambda^{-2},\qquad
 q_{\lambda,\bm x,\bm s}(w)
 :=\lambda^{-6}
 \mathcal Q_{\lambda\bm x,\lambda^2\bm s}(\lambda w).
 \tag{N29.3}
\]
The scaling identity \textup{(N24.7)} says exactly that
\[
 q_{\lambda,\bm x,\bm s}
 =e^{-(\varepsilon/2)D_w^2}F_{\bm x,\bm s}.
 \tag{N29.4}
\]

\begin{lemma}[Uniform scaled coefficient ledger]
\label{lem:f2v29-scaled-ledger}
Assume
\[
 \lambda^2\geq15,\qquad
 s_j\geq0,\qquad
 \sum_js_j\leq2+\frac{15}{\lambda^2}.
 \tag{N29.5}
\]
Then every coefficient of \(F_{\bm x,\bm s}\) has absolute value at
most \(1728\), every coefficient of
\(q_{\lambda,\bm x,\bm s}\) has absolute value at most \(89856\), and
corresponding coefficients differ by at most
\[
 88128\,\lambda^{-2}.
 \tag{N29.6}
\]
\end{lemma}

\begin{proof}
Condition \textup{(N29.1)} gives \(|x_j|<2\), while
\textup{(N29.5)} gives \(\sum_js_j\leq3\).  Each quadratic factor in
\textup{(N29.2)} therefore has coefficient \(\ell^1\)-norm at most
\[
 1+4+7=12.
 \tag{N29.7}
\]
The coefficient \(\ell^1\)-norm of the product is at most
\(12^3=1728\), which bounds each coefficient.

For heat time at most one, the seven coefficient row sums of the
degree-six heat operator are bounded by
\[
 20,\ 19,\ 52,\ 11,\ 16,\ 1,\ 1.
 \tag{N29.8}
\]
Thus each heat-image coefficient is at most
\[
 52\cdot1728=89856.
 \tag{N29.9}
\]
After removing the identity term, the row sums are bounded by
\[
 19,\ 18,\ 51,\ 10,\ 15,\ 0,\ 0
 \tag{N29.10}
\]
times \(\varepsilon\), since
\(\varepsilon^k\leq\varepsilon\) for \(1\leq k\leq3\).
The maximum gives
\[
 51\cdot1728\,\varepsilon
 =88128\,\lambda^{-2},
 \tag{N29.11}
\]
proving \textup{(N29.6)}.
\end{proof}

\subsection{Quantitative discriminant stability}
\label{subsec:f2v29-discriminant}

Set
\[
 B:=540000,\qquad D:=528768,\qquad
 K_{\rm Syl}:=11\cdot11!\,D B^{10}.
 \tag{N29.12}
\]

\begin{lemma}[Scaled discriminant perturbation]
\label{lem:f2v29-disc-perturb}
Under \textup{(N29.5)},
\[
 \left|
 \operatorname{Disc}q_{\lambda,\bm x,\bm s}
 -\operatorname{Disc}F_{\bm x,\bm s}
 \right|
 \leq\frac{K_{\rm Syl}}{\lambda^2}.
 \tag{N29.13}
\]
\end{lemma}

\begin{proof}
The discriminant of a monic sextic is, up to sign, the determinant of
the \(11\)-dimensional Sylvester matrix of the sextic and its derivative.
Lemma~\ref{lem:f2v29-scaled-ledger} bounds coefficients of both
\(F\) and \(q\) by \(89856\), so every entry in either Sylvester matrix
is bounded by
\[
 6\cdot89856=539136<B.
 \tag{N29.14}
\]
The coefficient changes in the derivatives are at most
\[
 6\cdot88128\,\lambda^{-2}
 =D\lambda^{-2}.
 \tag{N29.15}
\]
Expand the two determinants by permutations and telescope each product
of eleven entries.  Each product changes by at most
\(11DB^{10}\lambda^{-2}\), and there are \(11!\) permutations.  This is
\textup{(N29.13)}.
\end{proof}

Recall the pair resultants
\[
 R_{ij}:=(x_i-x_j)^4
 +2(x_i-x_j)^2(s_i+s_j)+(s_i-s_j)^2.
 \tag{N29.16}
\]

\begin{theorem}[Effective generic projective exclusion]
\label{thm:f2v29-projective}
Fix \(\sigma>0\) and \(r>0\).  Suppose
\[
\begin{gathered}
 A_{\bm x}=1,\qquad \sum_jx_j=0,\\
 t_j=\lambda^2s_j,\qquad
 s_j\geq\sigma,\qquad
 R_{ij}\geq r\quad(i<j),\\
 \sum_jt_j\leq15+2\lambda^2.
\end{gathered}
\tag{N29.17}
\]
If
\[
 \lambda^2\geq
 \max\left\{15,\frac{K_{\rm Syl}}
 {32\sigma^3r^6}\right\},
 \tag{N29.18}
\]
then
\[
 \operatorname{Disc}
 q_{\lambda,\bm x,\bm s}<0.
 \tag{N29.19}
\]
Consequently neither \(q_{\lambda,\bm x,\bm s}\) nor the original heat
sextic is real-rooted.
\end{theorem}

\begin{proof}
The exact factorization \textup{(N24.9)} gives
\[
\begin{aligned}
 \operatorname{Disc}F_{\bm x,\bm s}
 &=-64s_1s_2s_3\prod_{i<j}R_{ij}^2\\
 &\leq-64\sigma^3r^6.
\end{aligned}
\tag{N29.20}
\]
The budget in \textup{(N29.17)} gives
\[
 \sum_js_j\leq2+\frac{15}{\lambda^2},
 \tag{N29.21}
\]
so Lemma~\ref{lem:f2v29-disc-perturb} applies.  The second bound in
\textup{(N29.18)} gives
\[
 \frac{K_{\rm Syl}}{\lambda^2}
 \leq32\sigma^3r^6.
 \tag{N29.22}
\]
Equations \textup{(N29.13)}, \textup{(N29.20)}, and
\textup{(N29.22)} imply
\[
 \operatorname{Disc}q_{\lambda,\bm x,\bm s}
 \leq-32\sigma^3r^6<0.
 \tag{N29.23}
\]
Root dilation by the positive number \(\lambda\) preserves
real-rootedness, proving the last assertion.
\end{proof}

\begin{corollary}[Separated-center projective cutoff]
\label{cor:f2v29-separated}
Under the hypotheses of Theorem~\ref{thm:f2v29-projective}, the pair
resultant condition follows from
\[
 \min_{i<j}|x_i-x_j|\geq\delta>0
 \tag{N29.24}
\]
by taking \(r=\delta^4\).  Hence the explicit sufficient scale bound is
\[
 \lambda^2\geq
 \max\left\{15,\frac{K_{\rm Syl}}
 {32\sigma^3\delta^{24}}\right\}.
 \tag{N29.25}
\]
\end{corollary}

\begin{proof}
Every term in \textup{(N29.16)} is nonnegative, and its first term is
\((x_i-x_j)^4\).
\end{proof}

\begin{corollary}[Effective degeneration localization]
\label{cor:f2v29-localization}
Let \(\lambda_n\to\infty\), let \(A_{\bm x_n}=1\), and suppose
\[
 s_{j,n}\geq0,\qquad
 \sum_js_{j,n}\leq2+\frac{15}{\lambda_n^2}.
 \tag{N29.26}
\]
If the associated heat sextics are real-rooted, then, along every such
sequence,
\[
 \min_j s_{j,n}\longrightarrow0
 \quad\hbox{or}\quad
 \min_{i<j}R_{ij,n}\longrightarrow0
 \tag{N29.27}
\]
after passage to a subsequence.
Moreover, in the second alternative,
\[
 |x_{i,n}-x_{j,n}|\longrightarrow0,\qquad
 |s_{i,n}-s_{j,n}|\longrightarrow0
 \tag{N29.28}
\]
for some fixed pair \(i<j\).
\end{corollary}

\begin{proof}
If both minima in \textup{(N29.27)} stayed bounded below, choose
\(\sigma,r>0\) below those bounds and apply
Theorem~\ref{thm:f2v29-projective} for all sufficiently large \(n\).
This contradicts real-rootedness.  In the remaining alternative, fix a
pair along a subsequence and use the nonnegative decomposition
\textup{(N29.16)}.  Its first and last terms give
\textup{(N29.28)}.
\end{proof}

\begin{firewall}[Only generic projective strata are effective]
\label{fire:f2v29}
The cutoff \textup{(N29.18)} degenerates as a scaled height approaches
zero or as two projective quadratic factors become identical.  V24--V25
exclude those limiting strata qualitatively through exact local heat
identities, but no common effective cutoff has yet been extracted from
their boundary layers.  The compact finite cover, infinite-product
passage, and arithmetic forcing step remain open.  GRH remains open.
\end{firewall}

\begin{assessment}[V29 advance and mandatory V30 sweep]
\label{ass:f2v29}
V28 and V29 now provide two complementary effective regions: a
fourth-minor horn for subprojective maximum height, and a discriminant
cutoff for the generic projective interior.  Any ineffective large-scale
region is quantitatively localized near a height face or an
identical-quadratic collision.

V30 is a multiple of ten and must begin with a fresh sweep of all active
parallel notes before further algebra.  After that sweep, the preferred
technical target is an effective mixed-face lemma: one scaled height
\(s_1\) tends to zero while \(s_2,s_3\) and the relevant center gaps stay
positive.  The exact local quotient identity \textup{(N24.13)} should be
combined with the present coefficient bounds rather than replaced by a
new moment minor.
\end{assessment}

\section*{Version V29 append-only record}
\addcontentsline{toc}{section}{Version V29 append-only record}

V29 was appended to the validated V28 active note whose installed
SHA--256 was
\[
 \texttt{E5D3ED266FA8144229161A8F8B37686567D6F217FAF8B702D538F092504B467C}.
 \tag{N29.29}
\]
V29 proves an explicit discriminant cutoff on every generic projective
stratum and quantitatively localizes all remaining large-scale
degenerations to height faces or identical quadratic factors.  After
validation V29 replaces V28 as the sole active numeric GRH note; the
frozen \texttt{V100\_f2} archive is unchanged.

% =============================================================================
% END APPENDED V29 MATERIAL
% =============================================================================

% =============================================================================
% BEGIN APPENDED V30 MATERIAL
% =============================================================================

\section{V30: companion sweep and a global two-high-height exclusion}
\label{sec:f2v30-two-high}

This is the scheduled ten-version checkpoint.  The sweep was performed
before new algebra.  During the sweep the active SM and Yang--Mills files
were replaced by their next versions; the table is the refreshed snapshot
actually read:
\begin{description}[leftmargin=2.8em,style=nextline]
\item[SM V49.]
\path{Renewal_SM_Einstein_Continuum_Research_Notes_V49.tex};
(14265) lines, (495226) bytes; SHA--256
\texttt{3388F2A6298A3A849F6D24863591A781}\allowbreak
\texttt{DA52281D696D06469987038925B95C1E}.
\item[Yang--Mills V8.]
\path{Renewal_Yang_Mills_Mass_Gap_Research_Notes_V8.tex};
(3978) lines, (145010) bytes; SHA--256
\texttt{12A6AEA000940160BCBCEB1260FBAE26}\allowbreak
\texttt{178A5FC82F5E44DAACEB6609AB4BD36A}.
\item[Parallel Yang--Mills V5.]
\path{Renewal_YM_Parallel_Research_Notes_V5.tex};
(1351) lines, (54174) bytes; SHA--256
\texttt{306F1BB84CFA8A8D47EA569EC17E9545}\allowbreak
\texttt{F9867C8D32B548EC9D9C444B4F3C0512}.
\item[Navier--Stokes V26.]
\path{Renewal_NS_Stability_Research_Notes_V26.tex};
(5319) lines, (278283) bytes; SHA--256
\texttt{82C887F8C2C69E4F28FAD7C3DC8DE5B}\allowbreak
\texttt{9099CB5A4BF431EBA1FFF3E91935978AD}.
\end{description}
\[
 \text{The full hashes above fix the four-file sweep snapshot.}
 \tag{N30.1}
\]
SM V46--V49 compiles the exact common stress object before taking
sectorwise norms and keeps its continuum closure conditional.  Yang--Mills
V7--V8 proves that global coefficient norms destroy an exact cancellation
and replaces them by a live-phase object with only its inverse defect
normed.  The parallel Yang--Mills note reaches its stated remainder stop
and separately proves a tensorization obstruction for an extensive charge.
The Navier--Stokes file is unchanged from V20 and supplies no GRH premise.

No companion theorem is imported.  The transferable decision is only to
preserve the exact local quotient \textup{(N24.13)} and estimate its whole
derivative-ratio defect.  This gives a center-free result.

\subsection{A quantitative quotient criterion}

For arbitrary real centers and positive squared heights write
\[
 P(z):=\prod_{i=1}^3((z-c_i)^2+t_i),\qquad
 \mathcal Q_{\bm c,\bm t}:=e^{-D_z^2/2}P.
 \tag{N30.2}
\]

\begin{theorem}[Global two-factor height criterion]
\label{thm:f2v30-height-criterion}
Fix \(j\), suppose \(t_j>1\), and put
\[
 a:=t_j-1,\qquad q:=\min_{k\ne j}\sqrt{t_k},
 \tag{N30.3}
\]
\[
 L(a,q):=\frac{2}{q\sqrt a}
 +\frac{3(1+4/a)}{q^2}
 +\frac{6}{q^3\sqrt a}
 +\frac{3(1+6/a)}{q^4}.
 \tag{N30.4}
\]
If \(L(a,q)<1\), then
\[
 \boxed{\mathcal Q_{\bm c,\bm t}(z)>0\quad(z\in\mathbb R).}
 \tag{N30.5}
\]
This is uniform in the centers and uses no moment budget.
\end{theorem}

\begin{proof}
Set \(y=z-c_j\), and for the other indices set
\[
 f_k(y):=(y+c_j-c_k)^2+t_k,\qquad G(y):=\prod_{k\ne j}f_k(y).
 \tag{N30.6}
\]
For real \(y\), \(G>0\).  Writing \(r=y+c_j-c_k\),
\[
 \frac{|f_k'|}{f_k}=\frac{2|r|}{r^2+t_k}\leq\frac1q,\qquad
 \frac{|f_k''|}{f_k}=\frac2{r^2+t_k}\leq\frac2{q^2}.
 \tag{N30.7}
\]
The product rule, using that both factors are quadratic, yields
\[
 \left|\frac{G'}G\right|\leq\frac2q,\quad
 \left|\frac{G''}G\right|\leq\frac6{q^2},\quad
 \left|\frac{G'''}G\right|\leq\frac{12}{q^3},\quad
 \left|\frac{G''''}G\right|\leq\frac{24}{q^4}.
 \tag{N30.8}
\]
Here \(G'''=3f_k''f_\ell'+3f_k'f_\ell''\) and
\(G''''=6f_k''f_\ell''\).

Apply the translation-invariant identity \textup{(N24.13)}.  With
\[
 u:=y^2+t_j,\qquad B:=u-1=y^2+a>0,
 \tag{N30.9}
\]
we have
\[
 |y|\leq\frac{B}{2\sqrt a},\qquad
 u+3\leq(1+4/a)B,\qquad u+5\leq(1+6/a)B.
 \tag{N30.10}
\]
Thus the absolute value of all four derivative terms in
\textup{(N24.13)} is at most \(L(a,q)B\).  Since its leading term is
\(B\),
\[
 \frac{\mathcal Q_{\bm c,\bm t}(c_j+y)}{G(y)}
 \geq(1-L(a,q))B>0.
 \tag{N30.11}
\]
\end{proof}

\subsection{The universal two-high-height exclusion}

Recall that the V10 constant satisfies \(\alpha_3>19/10\).

\begin{corollary}[Two heights at least \(36\) are impossible]
\label{cor:f2v30-two-high}
Suppose \(t_i\geq\alpha_3\) for all \(i\).  If two squared heights are
at least \(36\), and \(j\) is the remaining index, then
\[
 \boxed{\mathcal Q_{\bm c,\bm t}(c_j+y)
 >\frac{55}{432}(y^2+t_j-1)G(y)>0.}
 \tag{N30.12}
\]
Consequently the heat sextic is not real-rooted.
\end{corollary}

\begin{proof}
Now \(a=t_j-1>9/10\) and \(q\geq6\).  Since \(L\) decreases in both
variables, while \(1/a<10/9\) and \(1/\sqrt a<10/9\),
\[
 L(a,q)<\frac{10}{27}+\frac{49}{108}
 +\frac5{162}+\frac{23}{1296}=\frac{377}{432}.
 \tag{N30.13}
\]
Insert this in \textup{(N30.11)}.  A positive sextic has no real zero.
\end{proof}

\begin{corollary}[Projective collapse to one high height]
\label{cor:f2v30-projective-collapse}
Order
\(t_{(1)}\geq t_{(2)}\geq t_{(3)}\geq\alpha_3\).  Real-rootedness implies
\[
 \boxed{t_{(2)}<36.}
 \tag{N30.14}
\]
Hence, if \(\lambda_n\to\infty\) and
\(s_{i,n}=t_{i,n}/\lambda_n^2\), then
\[
 s_{(2),n}<36/\lambda_n^2\longrightarrow0,\qquad
 s_{(3),n}\longrightarrow0.
 \tag{N30.15}
\]
Every projective limit of putative counterexamples has at most one
positive scaled height, irrespective of pair resultants or center
collisions.
\end{corollary}

\begin{proof}
Two heights at least \(36\) contradict
Corollary~\ref{cor:f2v30-two-high}; division by \(\lambda_n^2\) gives
\textup{(N30.15)}.
\end{proof}

\begin{firewall}[What remains]
\label{fire:f2v30}
The theorem makes the mixed height-face cutoff effective, but does not
exclude exactly one large squared height with the other two in
\([\alpha_3,36)\).  The compact four-certificate cover, infinite-product
passage, and arithmetic forcing step also remain open.  GRH remains open.
\end{firewall}

\begin{assessment}[V30 advance and V31 direction]
\label{ass:f2v30}
The mandatory sweep was completed against stable hashes after detecting two
live companion replacements.  The cancellation-before-norm lesson led to
a global quotient estimate stronger than the V29 target: real-rootedness
forces the second-largest squared height below \(36\), uniformly in centers
and without a budget.

The large-scale problem is now a one-high-two-bounded face.  V31 should
split it geometrically: if the two bounded-height centers are separated on
the large scale, use the quotient identity at each bounded factor while
retaining the large factor; if they collide, blow up their quartic product
before estimating the high factor.  A generic sextic discriminant would
discard this reduction.
\end{assessment}

\section*{Version V30 append-only record}
\addcontentsline{toc}{section}{Version V30 append-only record}

V30 was appended to validated V29, whose installed SHA--256 was
\[
 \texttt{8E257B5E4B4A7CD891DBEA4AE721BDAF4ADBF0A494001EC9525D912A6157DBE6}.
 \tag{N30.16}
\]
It records the mandatory companion sweep and excludes every putative
real-rooted heat sextic with two squared heights at least \(36\).  After
validation V30 replaces V29 as the sole active numeric GRH note; the frozen
\texttt{V100\_f2} archive is unchanged.

% =============================================================================
% END APPENDED V30 MATERIAL
% =============================================================================


% =============================================================================
% BEGIN APPENDED V31 MATERIAL
% =============================================================================

\section{V31: tall-root isolation and effective global compactness}
\label{sec:f2v31-compactness}

V30 reduced every putative counterexample to one high squared height and
two bounded ones.  This continuation first sharpens the numerical boundary
in the V30 quotient criterion.  It then encloses the upper root of the
remaining high quadratic by a circle wholly inside the upper half-plane.
Rouche's theorem shows that backward heat cannot remove that root once its
height is large enough.  Combining this with the effective V28 horn gives
an explicit compact box for the entire three-pair candidate set.

\subsection{A sharper two-high threshold}

\begin{lemma}[The \(11/2\) quotient threshold]
\label{lem:f2v31-eleven-halves}
If \(a>9/10\) and \(q\geq11/2\), then the V30 quantity satisfies
\[
 L(a,q)
 <
 \frac{200}{517}+\frac{196}{363}
 +\frac{2400}{62557}+\frac{368}{14641}
 =
 \frac{2044340}{2064381}<1.
 \tag{N31.1}
\]
\end{lemma}

\begin{proof}
The inequality \(a>9/10\) gives
\[
 \frac1a<\frac{10}{9},\qquad
 \frac1{\sqrt a}<\frac{50}{47},
 \tag{N31.2}
\]
because \((47/50)^2=2209/2500<9/10\).
Insert these bounds and \(q\geq11/2\) into
\textup{(N30.4)}.  Its four terms are strictly smaller than the four
fractions in \textup{(N31.1)}.  Addition gives the displayed exact
fraction, whose numerator is smaller than its denominator by \(20041\).
\end{proof}

\begin{corollary}[Sharpened second-height ceiling]
\label{cor:f2v31-second-height}
Assume \(t_i\geq\alpha_3\) for all \(i\), and order the heights as
\(t_{(1)}\geq t_{(2)}\geq t_{(3)}\).  If
\(\mathcal Q_{\bm c,\bm t}\) is real-rooted, then
\[
 \boxed{t_{(2)}<\frac{121}{4}.}
 \tag{N31.3}
\]
\end{corollary}

\begin{proof}
If two heights were at least \(121/4\), choose the remaining index \(j\).
Then \(a=t_j-1>\alpha_3-1>9/10\) and
\(q\geq11/2\).  Lemma~\ref{lem:f2v31-eleven-halves} and
Theorem~\ref{thm:f2v30-height-criterion} make the heat sextic strictly
positive on the real line, contrary to real-rootedness.
\end{proof}

\subsection{A disk around the tall conjugate root}

\begin{theorem}[Tall-root Rouche criterion]
\label{thm:f2v31-rouche}
Suppose one input factor has roots
\[
 \alpha=c_1+iY,\qquad \overline\alpha=c_1-iY,\qquad Y>0,
 \tag{N31.4}
\]
and put
\[
 b:=\max(\sqrt{t_2},\sqrt{t_3}).
 \tag{N31.5}
\]
If
\[
 Y\geq3b,\qquad
 H(Y):=\frac{135}{Y^2}+\frac{3645}{Y^4}
       +\frac{10935}{Y^6}<1,
 \tag{N31.6}
\]
then \(\mathcal Q_{\bm c,\bm t}\) has a zero in the open upper
half-plane.  In particular, it is not real-rooted.
\end{theorem}

\begin{proof}
Factor the unheated monic sextic over \(\mathbb C\) as
\[
 P(z)=\prod_{\nu=1}^6(z-\zeta_\nu),
 \tag{N31.7}
\]
with \(\zeta_1=\alpha\) and \(\zeta_2=\overline\alpha\).
Consider the circle
\[
 C:=\{z:|z-\alpha|=Y/3\}.
 \tag{N31.8}
\]
The conjugate root has distance at least \(2Y-Y/3=5Y/3\) from
\(C\).  Each of the other four roots is \(c_k\mathbin{\pm}i\sqrt{t_k}\).
Its distance from \(\alpha\) is at least \(Y-b\), and hence its distance
from \(C\) is at least
\[
 Y-b-\frac Y3=\frac{2Y}{3}-b\geq\frac Y3.
 \tag{N31.9}
\]
The root \(\alpha\) itself has distance \(Y/3\) from every point of
\(C\).  Thus
\[
 |z-\zeta_\nu|\geq Y/3\qquad(z\in C,\ 1\leq\nu\leq6).
 \tag{N31.10}
\]

For \(0\leq m\leq6\), logarithmic differentiation of the product,
or the ordinary product rule, gives
\[
 \frac{P^{(m)}(z)}{P(z)}
 =
 m!\!\sum_{\substack{S\subset\{1,\ldots,6\}\\|S|=m}}
 \prod_{\nu\in S}\frac1{z-\zeta_\nu}.
 \tag{N31.11}
\]
Consequently, on \(C\),
\[
 \left|\frac{P^{(m)}}P\right|
 \leq m!\binom6m\left(\frac3Y\right)^m.
 \tag{N31.12}
\]
Since \(P\) has degree six,
\[
 \mathcal Q_{\bm c,\bm t}-P
 =-\frac12P''+\frac18P''''-\frac1{48}P''''''.
 \tag{N31.13}
\]
Equations \textup{(N31.12)}--\textup{(N31.13)} give
\[
 \frac{|\mathcal Q_{\bm c,\bm t}-P|}{|P|}
 \leq\frac{135}{Y^2}+\frac{3645}{Y^4}
       +\frac{10935}{Y^6}=H(Y)<1
 \qquad\text{on }C.
 \tag{N31.14}
\]
Rouche's theorem therefore gives \(P\) and
\(\mathcal Q_{\bm c,\bm t}\) the same number of zeros inside \(C\),
counted with multiplicity.  Formula \textup{(N31.9)} shows that \(P\)
has exactly the single zero \(\alpha\) there.  The disk bounded by \(C\)
lies in \(\operatorname{Im}z>0\), because its lowest imaginary part is
\(2Y/3\).  Hence the one heat-sextic zero in that disk is nonreal.
\end{proof}

\begin{corollary}[Absolute largest-height ceiling]
\label{cor:f2v31-largest-height}
Under the hypotheses and ordering of
Corollary~\ref{cor:f2v31-second-height}, real-rootedness implies
\[
 \boxed{t_{(1)}<\frac{1089}{4}.}
 \tag{N31.15}
\]
\end{corollary}

\begin{proof}
Suppose instead that \(t_{(1)}\geq1089/4\), and choose that factor as
the first.  Then \(Y=\sqrt{t_{(1)}}\geq33/2\), while
Corollary~\ref{cor:f2v31-second-height} gives
\(b<11/2\), so \(Y\geq3b\).  The function \(H\) is decreasing, and
exact arithmetic gives
\[
 H(33/2)
 =\frac{60}{121}+\frac{720}{14641}
  +\frac{960}{1771561}
 =\frac{966540}{1771561}<1.
 \tag{N31.16}
\]
Theorem~\ref{thm:f2v31-rouche} contradicts real-rootedness.
\end{proof}

\subsection{An explicit compact box}

Retain the centered normalization and variance
\[
 c_1+c_2+c_3=0,\qquad
 A=\frac12\sum_{j=1}^3c_j^2.
 \tag{N31.17}
\]

\begin{theorem}[Effective compactness of the candidate set]
\label{thm:f2v31-effective-compactness}
Let \(\mathcal C\) be the set of parameter tuples satisfying
\[
 t_j\geq\alpha_3,\qquad
 \sum_jt_j\leq15+2A,\qquad
 \mathcal Q_{\bm c,\bm t}\ \hbox{real-rooted},
 \tag{N31.18}
\]
together with \textup{(N31.17)}.  Every tuple in \(\mathcal C\), after
ordering the heights, obeys
\[
 \boxed{
 \begin{gathered}
 \alpha_3\leq t_{(3)}\leq t_{(2)}<\frac{121}{4},\\
 t_{(2)}\leq t_{(1)}<\frac{1089}{4},\\
 \sqrt A<\frac{186219}{4},\qquad |c_j|<54000\quad(j=1,2,3).
 \end{gathered}}
 \tag{N31.19}
\]
In particular, \(\mathcal C\) is compact.
\end{theorem}

\begin{proof}
The two height bounds are
Corollaries~\ref{cor:f2v31-second-height} and
\ref{cor:f2v31-largest-height}.  Put \(T=\max_jt_j\).
If \(A>0\), the effective V28 horn, Corollary~\ref{cor:f2v28-growth},
says that real-rootedness requires \(T>\sqrt A/171\).  Therefore
\[
 \sqrt A<171T<171\frac{1089}{4}=\frac{186219}{4},
 \qquad
 A<\frac{34677515961}{16}.
 \tag{N31.20}
\]
For \(A=0\) the same displayed bounds are immediate.  Since the other
two centers sum to \(-c_j\), Cauchy--Schwarz gives
\[
 2A=\sum_kc_k^2\geq c_j^2+\frac{c_j^2}{2}
 =\frac32c_j^2.
 \tag{N31.21}
\]
Thus \(c_j^2\leq4A/3<54000^2\); the final strict inequality is the
integer comparison
\[
 34677515961<34992000000.
 \tag{N31.22}
\]

It remains only to justify the last sentence.  Conditions
\textup{(N31.17)}--\textup{(N31.18)} are closed: the coefficient map
from parameters to the monic heat sextic is polynomial, and the set of
real-rooted monic sextics is closed under coefficient convergence.
The proved bounds make \(\mathcal C\) bounded in finite-dimensional
Euclidean space.  Hence it is compact.
\end{proof}

\begin{firewall}[Effective compactness is not emptiness]
\label{fire:f2v31}
V31 removes every escape to infinite height or infinite center variance
and supplies an explicit compact search domain.  It does not prove that
\(\mathcal C\) is empty.  A finite exact certificate must still exclude
the remaining compact tuples, and the infinite-product and arithmetic
forcing steps remain open.  GRH remains open.
\end{firewall}

\begin{assessment}[V31 advance and V32 direction]
\label{ass:f2v31}
The one-high-height face is now closed effectively.  The full necessary
three-pair parameter set is compact with rational height ceilings and an
integer center box.  This replaces the nonconstructive large-scale cutoff
of V25 by explicit constants.

V32 should exploit compactness rather than merely launch a uniform grid.
The four necessary moment certificates from V26 are polynomial in the
center and height invariants.  The next step is to write their exact
semialgebraic cover on the ordered compact chamber, use the solved mirror,
double-center, and equal-height tubes as boundary charts, and apply adaptive
Bernstein subdivision only to the residual core.  Any uncovered box must be
retained as a precise new analytic target rather than treated as numerical
evidence.
\end{assessment}

\section*{Version V31 append-only record}
\addcontentsline{toc}{section}{Version V31 append-only record}

V31 was appended to validated V30, whose installed SHA--256 was
\[
 \texttt{C3DE5EC8E39796279EDD020D930BA2028A887CF4F8F669B173A0B770836685B9}.
 \tag{N31.23}
\]
It sharpens the second-height ceiling to \(121/4\), isolates every
sufficiently tall conjugate root by Rouche's theorem, and proves effective
compactness of the full three-pair candidate set.  After validation V31
replaces V30 as the sole active numeric GRH note; the frozen
\texttt{V100\_f2} archive is unchanged.

% =============================================================================
% END APPENDED V31 MATERIAL
% =============================================================================


% =============================================================================
% BEGIN APPENDED V32 MATERIAL
% =============================================================================

\section{V32: a universal isolated-root certificate and a two-chart atlas}
\label{sec:f2v32-isolated-root}

V31 made the candidate set explicitly compact, but its largest-height
Rouche circle used the special ratio between the tallest and the other two
heights.  This continuation isolates the underlying rootwise statement.
A fixed circle of radius \(21/5\) works for every degree-six input root
that is both above the real line and separated from the other input roots.
Consequently every remaining high input root must belong to a quantitatively
specified collision cluster.  This gives an exact two-chart replacement for
a blind subdivision of the V31 box.

To avoid confusion with the fourth-moment certificate \(\mathcal C\), write
\(\mathfrak C_{\rm cand}\) for the candidate set called \(\mathcal C\) in
Theorem~\ref{thm:f2v31-effective-compactness}.

\subsection{The rootwise heat perturbation}

\begin{theorem}[Isolated upper root survives backward heat]
\label{thm:f2v32-isolated-root}
Let
\[
 P(z)=\prod_{\nu=1}^6(z-\zeta_\nu)
 \tag{N32.1}
\]
be a monic sextic, let \(\alpha=\zeta_1\), and let \(d>0\).  Suppose
\[
 \operatorname{Im}\alpha>d,\qquad
 |\alpha-\zeta_\nu|\geq2d\quad(2\leq\nu\leq6),
 \tag{N32.2}
\]
and
\[
 J(d):=\frac{15}{d^2}+\frac{45}{d^4}+\frac{15}{d^6}<1.
 \tag{N32.3}
\]
Then \(e^{-D_z^2/2}P\) has a zero in the open upper half-plane.
\end{theorem}

\begin{proof}
On the circle \(C_d=\{z:|z-\alpha|=d\}\), the reverse triangle
inequality and \textup{(N32.2)} give
\[
 |z-\zeta_\nu|\geq d\qquad(1\leq\nu\leq6).
 \tag{N32.4}
\]
The derivative identity \textup{(N31.11)} therefore gives
\[
 \left|\frac{P^{(m)}(z)}{P(z)}\right|
 \leq m!\binom6m d^{-m}.
 \tag{N32.5}
\]
Since
\[
 e^{-D_z^2/2}P-P
 =-\frac12P''+\frac18P''''-\frac1{48}P'''''',
 \tag{N32.6}
\]
we obtain on \(C_d\)
\[
 \frac{|e^{-D_z^2/2}P-P|}{|P|}
 \leq\frac{15}{d^2}+\frac{45}{d^4}+\frac{15}{d^6}
 =J(d)<1.
 \tag{N32.7}
\]
Rouche's theorem gives \(P\) and its heat image the same number of
zeros in the disk.  The separation in \textup{(N32.2)} says that \(P\)
has exactly the one zero \(\alpha\) there.  Since
\(\operatorname{Im}\alpha>d\), the whole closed disk is contained in
the open upper half-plane.  Its heat-image zero is therefore nonreal.
\end{proof}

\begin{corollary}[The rational radius \(21/5\)]
\label{cor:f2v32-radius}
Theorem~\ref{thm:f2v32-isolated-root} applies with
\[
 d_0:=\frac{21}{5},
 \qquad
 J(d_0)
 =\frac{28522625}{28588707}<1.
 \tag{N32.8}
\]
Thus a heat sextic can be real-rooted only if every input root
\(\alpha\) with \(\operatorname{Im}\alpha>21/5\) lies at distance
strictly less than \(42/5\) from another input root.
\end{corollary}

\begin{proof}
Substitution of \(d_0=21/5\) in \textup{(N32.3)} gives the displayed
fraction; its denominator exceeds its numerator by \(66082\).  If all
other roots had distance at least \(2d_0=42/5\), the theorem would give
a nonreal output root.
\end{proof}

\subsection{The exact low-or-cluster alternative}

For the three-pair input, put
\[
 y_j:=\sqrt{t_j}>0
 \tag{N32.9}
\]
and permute factors so that
\(y_1\geq y_2\geq y_3\).

\begin{theorem}[Two-chart candidate atlas]
\label{thm:f2v32-two-chart}
Every tuple in \(\mathfrak C_{\rm cand}\) belongs to one of the following
two charts.

\begin{enumerate}[label=\textup{(\roman*)}]
\item The low chart:
\[
 y_1\leq\frac{21}{5},\qquad
 t_1\leq\frac{441}{25},\qquad
 |c_j|<3500\quad(1\leq j\leq3).
 \tag{N32.10}
\]

\item A collision chart: after possibly exchanging factors \(2\) and \(3\),
\[
 \frac{21}{5}<y_1<\frac{139}{10},\qquad
 y_2,y_3<\frac{11}{2},
 \tag{N32.11}
\]
\[
 (c_1-c_2)^2+(y_1-y_2)^2<\frac{1764}{25},
 \qquad |c_1-c_2|<\frac{42}{5},
 \tag{N32.12}
\]
and
\[
 |c_j|<38200\quad(1\leq j\leq3).
 \tag{N32.13}
\]
\end{enumerate}
In particular, every candidate satisfies the improved absolute bound
\[
 \boxed{\max_jt_j<\frac{19321}{100}.}
 \tag{N32.14}
\]
The two charts are semialgebraic in the variables
\((c_1,c_2,c_3,y_1,y_2,y_3)\), with \(t_j=y_j^2\).
\end{theorem}

\begin{proof}
Corollary~\ref{cor:f2v31-second-height} gives
\[
 y_2<\frac{11}{2}.
 \tag{N32.15}
\]
If \(y_1\leq21/5\), the height part of
\textup{(N32.10)} follows.

Suppose \(y_1>21/5\), and consider the upper input root
\(\alpha_1=c_1+iy_1\).  By
Corollary~\ref{cor:f2v32-radius}, another input root lies within
\(42/5\).  It is not \(c_1-iy_1\), whose distance is
\(2y_1>42/5\).  Hence it belongs to factor \(2\) or \(3\); call that
factor \(2\).  If the nearby root is \(c_2-iy_2\), replacing it by
\(c_2+iy_2\) can only decrease its distance from \(\alpha_1\).
Therefore
\[
 |(c_1+iy_1)-(c_2+iy_2)|^2
 =(c_1-c_2)^2+(y_1-y_2)^2<\left(\frac{42}{5}\right)^2,
 \tag{N32.16}
\]
which proves \textup{(N32.12)}.  It also gives
\[
 y_1<y_2+\frac{42}{5}
 <\frac{11}{2}+\frac{42}{5}
 =\frac{139}{10}.
 \tag{N32.17}
\]
Squaring proves \textup{(N32.14)}.

It remains to prove the center boxes.  In the low chart, the V28 horn
and \(T=\max_jt_j\leq441/25\) give
\[
 \sqrt A<171T\leq\frac{75411}{25},
 \qquad
 A<\frac{5686818921}{625}.
 \tag{N32.18}
\]
As in \textup{(N31.21)}, \(c_j^2\leq4A/3\).  Exact comparison gives
\[
 \frac{4}{3}\frac{5686818921}{625}
 =\frac{7582425228}{625}<3500^2.
 \tag{N32.19}
\]
This proves the last part of \textup{(N32.10)}.

In either chart, \textup{(N32.14)} and V28 give
\[
 \sqrt A<171\frac{19321}{100}
 =\frac{3303891}{100}.
 \tag{N32.20}
\]
Consequently
\[
 c_j^2\leq\frac{4A}{3}
 <\frac{3638565246627}{2500}<38200^2,
 \tag{N32.21}
\]
which proves \textup{(N32.13)}.  The chart inequalities are polynomial
once the positive variables \(y_j\) and the equations \(t_j=y_j^2\)
are included.
\end{proof}

\begin{corollary}[Finite-cover reduction]
\label{cor:f2v32-finite-cover}
To prove \(\mathfrak C_{\rm cand}=\varnothing\), it is enough to certify
the four V26 moment obstructions on the union of:

\begin{enumerate}[label=\textup{(\roman*)}]
\item the ordered low box \textup{(N32.10)}, with the moment budget
and centered-plane constraint imposed; and
\item the two permutations of the collision chart
\textup{(N32.11)}--\textup{(N32.13)}.
\end{enumerate}

No box outside this atlas can contain a candidate.  In the collision
charts the quadratic inequality \textup{(N32.12)} should be retained
as a domain constraint rather than replaced by its rectangular hull.
\end{corollary}

\begin{proof}
This is exactly Theorem~\ref{thm:f2v32-two-chart}, with the two possible
choices of the close companion factor.  The four moment certificates are
necessary for real-rootedness by V22.
\end{proof}

\begin{firewall}[The atlas is not yet a sign cover]
\label{fire:f2v32}
The low-or-cluster alternative is an exact analytic cover of every
candidate, not a sampled observation.  It substantially shrinks the height
and center ranges and identifies a mandatory complex-root collision in the
high chart.  It does not prove that the four V26 certificate signs cover
either chart.  The infinite-product passage and arithmetic forcing step
also remain open.  GRH remains open.
\end{firewall}

\begin{assessment}[V32 advance and V33 direction]
\label{ass:f2v32}
V32 replaces the monolithic V31 box by two geometrically meaningful
semialgebraic charts and improves the largest squared-height ceiling from
\(1089/4\) to \(19321/100\).  The low chart has the much smaller center
box \(|c_j|<3500\); every higher candidate is forced into an absolute
radius-\(42/5\) collision between two upper input roots.

V33 should work on the collision chart before subdividing the low box.
Write \(h=c_1-c_2\), \(m=(c_1+c_2)/2\), and keep
\(h^2+(y_1-y_2)^2<1764/25\) exact.  For large \(|m|\) this is an
effective double-center chart around \((m,m,-2m)\); for bounded \(|m|\)
it becomes a genuinely smaller compact core.  The next proof should obtain
an explicit large-\(|m|\) obstruction from the V21 double-center
certificates, leaving only a bounded collision core for Bernstein
subdivision.
\end{assessment}

\section*{Version V32 append-only record}
\addcontentsline{toc}{section}{Version V32 append-only record}

V32 was appended to validated V31, whose installed SHA--256 was
\[
 \texttt{AB7A740CC0EC654C417C297E9C71137BB03BB5CBA5F7EFEE607BD5E782261026}.
 \tag{N32.22}
\]
It proves a universal isolated-root certificate, an exact low-or-collision
atlas, and the improved height ceiling \(19321/100\).  After validation V32
replaces V31 as the sole active numeric GRH note; the frozen
\texttt{V100\_f2} archive is unchanged.

% =============================================================================
% END APPENDED V32 MATERIAL
% =============================================================================


% =============================================================================
% BEGIN APPENDED V33 MATERIAL
% =============================================================================

\section{V33: complex local quotients and a small global center box}
\label{sec:f2v33-center-box}

V32 forced every candidate into either a low-height chart or a cluster of
two upper input roots.  The center bounds in that note still came from the
coarse V28 horn.  Here the exact local quotient \textup{(N24.13)} is used
on a complex circle around the root of its leading shifted quadratic.
This directly excludes an isolated center and reduces the entire candidate
set to \(|c_j|<871/10\).

\subsection{A uniform complex local-quotient lemma}

\begin{theorem}[Effective center-isolation criterion]
\label{thm:f2v33-center-isolation}
Fix an index \(j\), and suppose
\[
 t_j\geq\alpha_3,\qquad t_j\leq U,\qquad
 \sqrt{t_j-1}\leq S.
 \tag{N33.1}
\]
Let
\[
 \Delta_j:=\min_{k\ne j}|c_j-c_k|.
 \tag{N33.2}
\]
For \(D>0\), define
\[
\begin{aligned}
 \Phi(S,U,D):={}&
 \frac{12S}{D}
 +\frac{12}{D^2}
  \left(\frac32+\frac{9S^2/4+U}{2}\right)\\
 &+\frac{36S}{D^3}
 +\frac{3(9S^2/4+U+5)}{D^4}.
\end{aligned}
\tag{N33.3}
\]
If
\[
 \Delta_j\geq D+\frac S2,\qquad
 \Phi(S,U,D)<\frac{27}{40},
 \tag{N33.4}
\]
then \(\mathcal Q_{\bm c,\bm t}\) has a zero in the open upper
half-plane and is not real-rooted.
\end{theorem}

\begin{proof}
Put
\[
 a:=t_j-1,\qquad s:=\sqrt a,\qquad y:=z-c_j,
 \tag{N33.5}
\]
and retain the factorization
\[
 P=(y^2+t_j)G(y),\qquad
 G(y)=\prod_{k\ne j}\bigl((y+c_j-c_k)^2+t_k\bigr).
 \tag{N33.6}
\]
The algebraic identity \textup{(N24.13)} holds for complex \(y\)
wherever \(G(y)\ne0\).  Consider
\[
 \Gamma:=\{y:|y-is|=s/2\}.
 \tag{N33.7}
\]
The four roots of \(G\) have real parts \(c_k-c_j\).  On and inside
\(\Gamma\), \(|\operatorname{Re}y|\leq s/2\leq S/2\).
Consequently \textup{(N33.4)} puts every root of \(G\) at horizontal,
hence Euclidean, distance at least \(D\) from every point of \(\Gamma\).
In particular \(G\) has no zero in the disk.

Applying \textup{(N31.11)} to the degree-four polynomial \(G\) gives
on \(\Gamma\)
\[
 \left|\frac{G^{(r)}}G\right|
 \leq r!\binom4rD^{-r}\qquad(1\leq r\leq4),
 \tag{N33.8}
\]
or explicitly
\[
 \left|\frac{G'}G\right|\leq\frac4D,\quad
 \left|\frac{G''}G\right|\leq\frac{12}{D^2},\quad
 \left|\frac{G'''}G\right|\leq\frac{24}{D^3},\quad
 \left|\frac{G''''}G\right|\leq\frac{24}{D^4}.
 \tag{N33.9}
\]
Also
\[
 |y|\leq\frac{3S}{2},\qquad
 |y^2+a|=|y-is|\,|y+is|
 \geq\frac{3s^2}{4}
 =\frac{3a}{4}>\frac{27}{40},
 \tag{N33.10}
\]
because \(a\geq\alpha_3-1>9/10\).

Let \(B(y)=y^2+a\).  The coefficient bounds
\[
 \left|\frac32-\frac{y^2+t_j}{2}\right|
 \leq\frac32+\frac{9S^2/4+U}{2},
 \tag{N33.11}
\]
\[
 \left|\frac{y^2+t_j}{8}-\frac58\right|
 \leq\frac{9S^2/4+U+5}{8}
 \tag{N33.12}
\]
and \textup{(N33.9)} show that the four derivative terms in
\textup{(N24.13)} have total modulus at most \(\Phi(S,U,D)\).
Thus, on \(\Gamma\),
\[
 |\mathcal Q_{\bm c,\bm t}-BG|
 =|G|\left|\frac{\mathcal Q_{\bm c,\bm t}}G-B\right|
 <|G|\,\frac{27}{40}<|BG|.
 \tag{N33.13}
\]
Rouche's theorem says that \(\mathcal Q_{\bm c,\bm t}\) and \(BG\)
have the same number of zeros inside \(\Gamma\).  The polynomial \(G\)
has none there, while \(B\) has exactly the root \(is\).  The disk is
contained in the upper half-plane because its lowest imaginary part is
\(s/2>0\).  Hence the heat sextic has a nonreal zero.
\end{proof}

\subsection{Closing the large-center part of the low chart}

\begin{lemma}[Low-chart contour margin]
\label{lem:f2v33-low-margin}
For
\[
 S=\frac{21}{5},\qquad U=\frac{441}{25},\qquad D=85,
 \tag{N33.14}
\]
one has
\[
 \Phi(S,U,D)
 =\frac{3358024449}{5220062500}
 <\frac{27}{40},
 \tag{N33.15}
\]
with exact margin
\[
 \frac{27}{40}-\Phi(S,U,D)
 =\frac{331035477}{10440125000}>0.
 \tag{N33.16}
\]
\end{lemma}

\begin{proof}
Substitution in \textup{(N33.3)} gives
\[
 \Phi=
 \frac{252/5}{85}
 +\frac{18099/50}{85^2}
 +\frac{756/5}{85^3}
 +\frac{18699/100}{85^4},
 \tag{N33.17}
\]
whose exact addition is \textup{(N33.15)}--\textup{(N33.16)}.
\end{proof}

\begin{corollary}[Small low chart]
\label{cor:f2v33-small-low}
In the low chart \textup{(N32.10)}, order the centers
\(c_1\leq c_2\leq c_3\) and put
\[
 g_1:=c_2-c_1,\qquad g_2:=c_3-c_2.
 \tag{N33.18}
\]
Every candidate obeys
\[
 \boxed{g_1,g_2<\frac{871}{10},\qquad
 |c_j|<\frac{871}{10}.}
 \tag{N33.19}
\]
Moreover,
\[
 A<\frac{758641}{100}.
 \tag{N33.20}
\]
\end{corollary}

\begin{proof}
Every low-chart height satisfies \(t_j\leq441/25\), while
\[
 \sqrt{t_j-1}\leq\sqrt{416/25}<21/5.
 \tag{N33.21}
\]
If \(g_1\geq871/10=85+21/10\), the leftmost center is separated from
both other centers by at least that amount.  Theorem
\ref{thm:f2v33-center-isolation} and
Lemma~\ref{lem:f2v33-low-margin} exclude the tuple.  If
\(g_2\geq871/10\), apply the same argument to the rightmost center.
This proves the gap bounds.

Writing the centered triple in terms of its gaps,
\[
 c_1=-\frac{2g_1+g_2}{3},\quad
 c_2=\frac{g_1-g_2}{3},\quad
 c_3=\frac{g_1+2g_2}{3},
 \tag{N33.22}
\]
gives the center bound in \textup{(N33.19)} and
\[
 A=\frac{g_1^2+g_1g_2+g_2^2}{3}
 <\left(\frac{871}{10}\right)^2.
 \tag{N33.23}
\]
\end{proof}

\subsection{Closing the large-midpoint part of the collision chart}

\begin{lemma}[Collision-chart contour margin]
\label{lem:f2v33-collision-margin}
For
\[
 S=\frac{11}{2},\qquad U=\frac{121}{4},\qquad D=120,
 \tag{N33.24}
\]
one has
\[
 \Phi(S,U,D)
 =\frac{218356391}{368640000}
 <\frac{27}{40},
 \tag{N33.25}
\]
with exact margin
\[
 \frac{27}{40}-\Phi(S,U,D)
 =\frac{30475609}{368640000}>0.
 \tag{N33.26}
\]
\end{lemma}

\begin{proof}
The four terms of \(\Phi\) are
\[
 \frac{66}{120}
 +\frac{4863/8}{120^2}
 +\frac{198}{120^3}
 +\frac{4959/16}{120^4}.
 \tag{N33.27}
\]
Exact addition gives the assertion.
\end{proof}

\begin{corollary}[Small collision chart]
\label{cor:f2v33-small-collision}
In a collision chart, write
\[
 c_1=m+\frac h2,\qquad c_2=m-\frac h2,\qquad c_3=-2m.
 \tag{N33.28}
\]
Every candidate obeys
\[
 \boxed{
 |h|<\frac{42}{5},\qquad
 |m|<\frac{2539}{60}.}
 \tag{N33.29}
\]
Consequently
\[
 |c_1|,|c_2|<\frac{2791}{60},\qquad
 |c_3|<\frac{2539}{30},
 \tag{N33.30}
\]
and
\[
 A=3m^2+\frac{h^2}{4}
 <\frac{6467689}{1200}.
 \tag{N33.31}
\]
\end{corollary}

\begin{proof}
The first inequality is already \textup{(N32.12)}.  Suppose
\(|m|\geq2539/60\).  The third center then satisfies
\[
\begin{aligned}
 \min\{|c_3-c_1|,|c_3-c_2|\}
 &\geq3|m|-\frac{|h|}{2}\\
 &>\frac{2539}{20}-\frac{21}{5}
 =\frac{491}{4}
 =120+\frac{11}{4}.
\end{aligned}
\tag{N33.32}
\]
The remaining factor has \(t_3<121/4\), hence
\(\sqrt{t_3-1}<11/2\).  Apply
Theorem~\ref{thm:f2v33-center-isolation} to \(j=3\) with the constants
of Lemma~\ref{lem:f2v33-collision-margin}.  This contradicts
real-rootedness, proving the midpoint bound.  Equations
\textup{(N33.30)}--\textup{(N33.31)} follow directly from
\textup{(N33.28)}--\textup{(N33.29)}.
\end{proof}

\begin{theorem}[Global small-box reduction]
\label{thm:f2v33-global-small-box}
Every tuple in \(\mathfrak C_{\rm cand}\) satisfies
\[
 \boxed{
 |c_j|<\frac{871}{10},\qquad
 A<\frac{758641}{100},\qquad
 \max_jt_j<\frac{19321}{100}.}
 \tag{N33.33}
\]
More precisely, it lies either in the small low chart
\textup{(N32.10)}, \textup{(N33.19)}, or in one of the two small
collision charts \textup{(N32.11)}--\textup{(N32.12)},
\textup{(N33.29)}.
\end{theorem}

\begin{proof}
Use the exact V32 two-chart atlas.  The low chart obeys
Corollary~\ref{cor:f2v33-small-low}.  The collision chart obeys
Corollary~\ref{cor:f2v33-small-collision}; its center bounds are strictly
smaller than \(871/10\), and
\[
 \frac{6467689}{1200}<\frac{758641}{100}.
 \tag{N33.34}
\]
The height bound is \textup{(N32.14)}.
\end{proof}

\begin{firewall}[A small box still requires exclusion]
\label{fire:f2v33}
V33 replaces the previous center bound \(38200\) by \(871/10\) and keeps
the collision geometry exact.  It does not prove that the resulting
compact charts contain no candidate.  An exact moment-certificate cover,
the infinite-product passage, and the arithmetic forcing step remain open.
GRH remains open.
\end{firewall}

\begin{assessment}[V33 advance and V34 direction]
\label{ass:f2v33}
The large-center part of both V32 charts is now excluded by one uniform
complex local-quotient theorem.  All candidate centers and heights lie in
a modest rational box, and the high-height portion retains a mandatory
root-collision inequality.

V34 should now construct the first reproducible interval ledger on this
small domain.  It should parameterize ordered centers by the gaps
\((g_1,g_2)\), use \(A=(g_1^2+g_1g_2+g_2^2)/3\), impose the moment budget
before subdivision, and evaluate the four exact V22 certificates in
scaled rational Bernstein form.  The collision charts should instead use
\((m,h,y_1,y_2,y_3)\) and preserve
\(h^2+(y_1-y_2)^2<1764/25\).  Any uncovered rational box must be saved as
an analytic target; a floating scan is not a certificate.
\end{assessment}

\section*{Version V33 append-only record}
\addcontentsline{toc}{section}{Version V33 append-only record}

V33 was appended to validated V32, whose installed SHA--256 was
\[
 \texttt{CDEE3FDA9537BEB7E9C19926DD92982740F22EC47BE0F920F10F4D0806A19ACF}.
 \tag{N33.35}
\]
It proves an effective complex center-isolation criterion and reduces every
candidate to \(|c_j|<871/10\).  After validation V33 replaces V32 as the
sole active numeric GRH note; the frozen \texttt{V100\_f2} archive is
unchanged.

% =============================================================================
% END APPENDED V33 MATERIAL
% =============================================================================

% =============================================================================
% BEGIN APPENDED V34 MATERIAL
% =============================================================================

\section{V34: an effective equality collar and the corrected target}
\label{sec:f2v34-equality-collar}

The compact-cover target must retain the known equality case.  V32 used the
phrase \(\mathfrak C_{\rm cand}=\varnothing\), but the symmetric tuple
\[
 \mathfrak e:=
 (c_1,c_2,c_3;t_1,t_2,t_3)
 =(0,0,0;\alpha_3,\alpha_3,\alpha_3)
 \tag{N34.1}
\]
belongs to the candidate set: its heat image is the real-rooted polynomial
\textup{(N13.9)}.  The correct theorem is
\[
 \boxed{\mathfrak C_{\rm cand}=\{\mathfrak e\}.}
 \tag{N34.2}
\]
Thus Corollary~\ref{cor:f2v32-finite-cover} is corrected by replacing its
empty-set target with uniqueness of \(\mathfrak e\).  None of the
V32--V33 bounds changes.

This continuation makes a full neighborhood of \(\mathfrak e\) effective.
The adapted square from V16 is strictly negative at every other point in
that neighborhood, with a rational margin.

\subsection{Weighted exact expansion of the adapted square}

Put
\[
 c_1=u,\qquad c_2=v,\qquad c_3=-u-v,\qquad
 A=u^2+uv+v^2,
 \tag{N34.3}
\]
and
\[
 \delta_j:=t_j-\alpha_3\geq0,\qquad
 \Delta:=\delta_1+\delta_2+\delta_3,\qquad
 r:=A+\Delta.
 \tag{N34.4}
\]
Give \(u,v\) weight one and each \(\delta_j\) weight two.

\begin{lemma}[Exact higher-order majorant]
\label{lem:f2v34-remainder}
The adapted square \(\mathcal S_{\rm gen}\) of \textup{(N16.7)} has
the exact decomposition
\[
 \mathcal S_{\rm gen}
 =c_AA+\frac{c_t}{3}\Delta+\mathcal R,
 \tag{N34.5}
\]
where \(c_A,c_t\) are \textup{(N13.14)}.  Every monomial of
\(\mathcal R\) has weighted degree \(4\), \(6\), or \(8\).  Whenever
\(0\leq r\leq1\),
\[
 \boxed{|\mathcal R|\leq\frac{737793}{2}\,r^2.}
 \tag{N34.6}
\]
\end{lemma}

\begin{proof}
The terms of weights zero and two are exactly the threshold value and
Taylor ledger already proved in Lemma~\ref{lem:f2v16-joint-taylor}.
The threshold value vanishes, and the weight-two terms are
\(c_AA+(c_t/3)\Delta\).  Direct exact expansion of the remaining
polynomial gives only weights \(4,6,8\).

Here is the coefficient estimate used for that expansion.  The rational
enclosures \textup{(N13.17)} give
\[
 |\alpha_3|<2,\qquad |s|<\frac32.
 \tag{N34.7}
\]
Also
\[
 A=u^2+uv+v^2\geq\frac{u^2+v^2}{2},
 \tag{N34.8}
\]
so \(|u|,|v|\leq\sqrt{2r}\), while
\(0\leq\delta_j\leq r\).  For a monomial
\[
 q\,\alpha_3^a s^b u^m v^n
 \delta_1^{e_1}\delta_2^{e_2}\delta_3^{e_3},
 \tag{N34.9}
\]
replace its coefficient by
\[
 |q|\,2^a(3/2)^b\,2^{\lceil(m+n)/2\rceil}.
 \tag{N34.10}
\]
Its variable part is then bounded by
\(r^{(m+n)/2+e_1+e_2+e_3}\), hence by \(r^2\) at every retained
weighted degree.  Summing \textup{(N34.10)} over the exact expansion
gives
\[
 \sum_{\mathcal R}\text{coefficient majorants}
 =\frac{737793}{2}.
 \tag{N34.11}
\]
This proves \textup{(N34.6)}.
\end{proof}

\begin{remark}[Dependency-free exact verifier]
\label{rec:f2v34-verifier}
The source
\begin{center}
\path{scripts/verify_grh_v34_effective_local_bridge.py}
\end{center}
implements multivariate polynomials as dictionaries with
\texttt{Fraction} coefficients.  It constructs the three quadratic
factors, applies the degree-six heat operator, generates Newton sums
through \(p_8\), and expands \(\mathcal S_{\rm gen}\).  Its asserted
ledger is
\[
\begin{array}{c|c}
\text{object}&\text{exact output}\\ \hline
\text{terms in }P,\mathcal Q,\mathcal S_{\rm gen}&94,\ 138,\ 261\\
\text{raw terms by weight}&
0:8,\ 2:42,\ 4:88,\ 6:90,\ 8:33\\
\text{remainder outer monomials}&102\\
\text{majorant}&737793/2\\
\text{collar radius}&1/3000\\
\text{remaining linear margin}&254069/2000
\end{array}
\tag{N34.12}
\]
The verifier has \(151\) lines, \(3493\) bytes, and SHA--256
\[
 \texttt{97C523CF0924F96B4EA311DAB99162ECF8A6C815C0D41F467A57AF11022BAAE2}.
 \tag{N34.13}
\]
It terminates with \texttt{VERIFIED}; no symbolic-algebra package or
floating-point sign decision is used.
\end{remark}

\subsection{The explicit rigidity collar}

\begin{theorem}[Effective joint equality collar]
\label{thm:f2v34-effective-collar}
Assume
\[
 t_j\geq\alpha_3,\qquad
 A+\sum_{j=1}^3(t_j-\alpha_3)\leq\frac1{3000}.
 \tag{N34.14}
\]
If \(\mathcal Q_{\bm c,\bm t}\) is real-rooted, then the tuple is exactly
\(\mathfrak e\).
\end{theorem}

\begin{proof}
Use the rational coefficient bounds in \textup{(N13.18)}:
\[
 c_A<-263.930<-250,\qquad
 \frac{c_t}{3}<-\frac{750.049}{3}<-250.
 \tag{N34.15}
\]
If \(0<r\leq1/3000\), Lemma~\ref{lem:f2v34-remainder} gives
\[
\begin{aligned}
 \mathcal S_{\rm gen}
 &\leq-250r+\frac{737793}{2}r^2\\
 &\leq-
 \left(250-\frac{737793}{6000}\right)r\\
 &=-\frac{254069}{2000}r<0.
\end{aligned}
\tag{N34.16}
\]
For six real roots, \(\mathcal S_{\rm gen}\) is the sum of the six
real squares in \textup{(N16.6)}, so it must be nonnegative.  Therefore
\(r=0\).  Since \(A,\delta_1,\delta_2,\delta_3\) are nonnegative,
\(A=0\) and every \(\delta_j=0\), which is precisely
\(\mathfrak e\).
\end{proof}

\begin{corollary}[Correct punctured compact core]
\label{cor:f2v34-punctured-core}
Every nonexceptional tuple in \(\mathfrak C_{\rm cand}\) must satisfy
\[
 \boxed{
 A+\sum_j(t_j-\alpha_3)>\frac1{3000}.}
 \tag{N34.17}
\]
Consequently the remaining finite-certificate target is the compact
semialgebraic union obtained from the small V33 low and collision charts
by imposing
\[
 A+\sum_j(t_j-\alpha_3)\geq\frac1{3000}.
 \tag{N34.18}
\]
Proving that this punctured union contains no real-rooted tuple is
equivalent to \textup{(N34.2)}.
\end{corollary}

\begin{proof}
The strict inequality is the contrapositive of
Theorem~\ref{thm:f2v34-effective-collar}.  The V33 chart union is bounded;
all defining inequalities in its closed hull, including
\textup{(N34.18)}, the centered condition, moment budget, and lower-height
conditions, are closed.  Hence the new ambient target is compact.  The
only candidate removed by the collar is \(\mathfrak e\), so exclusion of
the punctured target proves uniqueness, not emptiness of the full candidate
set.
\end{proof}

\begin{firewall}[A corrected compact target is not yet a cover]
\label{fire:f2v34}
V34 gives an explicit equality collar and repairs the target statement.
It does not prove that the punctured V33 charts are covered by the four
V26 moment obstructions.  Nor does it complete the infinite-product
passage or arithmetic forcing step.  GRH remains open.
\end{firewall}

\begin{assessment}[V34 advance and V35 direction]
\label{ass:f2v34}
The extremal point is now separated from every possible nonexceptional
candidate by the rational condition \textup{(N34.17)}.  This removes the
zero-margin boundary that would make a strict interval cover impossible
and converts the next computation into a compact annular problem.

V35 should construct the first exact adaptive ledger on the punctured low
chart, using ordered center gaps and a simplex parameterization of the
nonnegative height excesses.  Boxes violating the moment budget or
\textup{(N34.18)} should be discarded before polynomial evaluation.
For every retained box, Bernstein signs of
\((\mathcal C,\mathcal E,\mathcal K,\mathcal M_3)\) should be tested in
that order.  The output must distinguish certified exclusions from
unresolved boxes and preserve the latter as exact rational targets.
\end{assessment}

\section*{Version V34 append-only record}
\addcontentsline{toc}{section}{Version V34 append-only record}

V34 was appended to validated V33, whose installed SHA--256 was
\[
 \texttt{901509221F26B7FBA4E07545DC6AA2B8DFE9D889303125AA6894E48949C54079}.
 \tag{N34.19}
\]
It corrects the finite target to uniqueness of the symmetric extremizer,
proves the explicit collar \(A+\Delta\leq1/3000\), and records its exact
dependency-free verifier.  After validation V34 replaces V33 as the sole
active numeric GRH note; the frozen \texttt{V100\_f2} archive is unchanged.

% =============================================================================
% END APPENDED V34 MATERIAL
% =============================================================================

% =============================================================================
% BEGIN APPENDED V35 MATERIAL
% =============================================================================

\section{V35: a regrouped exact interval shell}
\label{sec:f2v35-regrouped-shell}

V34 obtained a deliberately coarse collar by taking absolute values before
the algebraic coefficients in \(\alpha_3\) and
\(s=\sqrt{10-4\alpha_3}\) were combined.  This continuation performs the
opposite operation: for each monomial in the five displacement variables,
it first regroups its full coefficient in \((\alpha_3,s)\), and only then
applies the rational enclosure \textup{(N13.17)}.  It also preserves the
weighted degrees four, six, and eight separately.  The result removes the
entire shell left between the V34 radius and \(1/66\).

Retain the coordinates
\[
 c_1=u,\qquad c_2=v,\qquad c_3=-u-v,\qquad
 A=u^2+uv+v^2,
\tag{N35.1}
\]
\[
 \delta_j=t_j-\alpha_3\geq0,\qquad
 \Delta=\delta_1+\delta_2+\delta_3,\qquad r=A+\Delta.
\tag{N35.2}
\]

\subsection{Coefficientwise rational interval ledger}

Write the exact V34 remainder as its weighted-homogeneous decomposition
\[
 \mathcal R=\mathcal R_4+\mathcal R_6+\mathcal R_8,
\tag{N35.3}
\]
where \(u,v\) have weight one and the \(\delta_j\) have weight two.

\begin{lemma}[Regrouped weighted majorants]
\label{lem:f2v35-regrouped-majorants}
For \(0\leq r\leq1\),
\[
 |\mathcal R_4|\leq M_4r^2,\qquad
 |\mathcal R_6|\leq M_6r^3,\qquad
 |\mathcal R_8|\leq M_8r^4,
\tag{N35.4}
\]
with the exact rational constants
\[
 \boxed{
 M_4=\frac{789355807119}{50000000},\qquad
 M_6=\frac{50990499}{1250},\qquad
 M_8=48582.}
\tag{N35.5}
\]
\end{lemma}

\begin{proof}
Expand the exact polynomial \(\mathcal S_{\rm gen}\) constructed in V34.
For a fixed outer monomial
\[
 u^mv^n\delta_1^{e_1}\delta_2^{e_2}\delta_3^{e_3},
\tag{N35.6}
\]
regroup all raw terms into the coefficient
\[
 g(\alpha_3,s)=\sum_{a,b}q_{ab}\alpha_3^as^b,
 \qquad q_{ab}\in\mathbb Q.
\tag{N35.7}
\]
The V13 bounds give the positive rational rectangle
\[
 \frac{19434}{10000}<\alpha_3<\frac{19435}{10000},
 \qquad
 \frac{14919}{10000}<s<\frac{14922}{10000}.
\tag{N35.8}
\]
Evaluate every positive \(q_{ab}\)-term at the two coordinatewise
endpoints in increasing order, and reverse the endpoints for every negative
term.  After summation this gives a rational interval
\(I_g=[g_-,g_+]\) containing the actual coefficient.  Put
\(B_g=\max\{|g_-|,|g_+|\}\).

As in \textup{(N34.8)},
\(|u|,|v|\leq\sqrt{2r}\), and each \(\delta_j\leq r\).  Hence the
absolute value of \textup{(N35.6)} times its coefficient is at most
\[
 2^{\lceil(m+n)/2\rceil}B_g\,
 r^{(m+n)/2+e_1+e_2+e_3}.
\tag{N35.9}
\]
There are exactly \(102\) distinct outer monomials after regrouping.  Sum
the rational prefactors in \textup{(N35.9)} separately at weights
\(4,6,8\).  Exact addition gives
\[
\begin{array}{c|c|c}
\text{weight}&\text{raw terms}&\text{summed majorant}\\ \hline
4&88&789355807119/50000000\\
6&90&50990499/1250\\
8&33&48582
\end{array}
\tag{N35.10}
\]
and proves \textup{(N35.4)}--\textup{(N35.5)}.
\end{proof}

\begin{remark}[Exact verifier]
\label{rem:f2v35-verifier}
The dependency-free source
\begin{center}
\path{scripts/verify_grh_v35_regrouped_shell.py}
\end{center}
constructs \(P\), its heat sextic, Newton sums through \(p_8\), and the
adapted square from rational polynomial dictionaries.  It asserts the term
counts
\[
 (0:8),\ (2:42),\ (4:88),\ (6:90),\ (8:33),
\tag{N35.11}
\]
the \(102\) regrouped outer monomials, all three constants in
\textup{(N35.5)}, and the final margin below.  It has \(188\) lines,
\(4814\) bytes, and SHA--256
\[
\texttt{5AF8AF0792276F4EE91DEE3DE6634C1D9BC39AE1480C9597E47DBADA023A97F2}.
\tag{N35.12}
\]
It terminates with \texttt{VERIFIED}; every comparison is over
\(\mathbb Q\).
\end{remark}

\subsection{The enlarged rigidity collar}

\begin{theorem}[Regrouped equality collar]
\label{thm:f2v35-collar}
Assume
\[
 t_j\geq\alpha_3,\qquad
 A+\sum_{j=1}^3(t_j-\alpha_3)\leq\frac1{66}.
\tag{N35.13}
\]
If \(\mathcal Q_{\bm c,\bm t}\) is real-rooted, then the tuple is the
symmetric extremal
\[
 (c_1,c_2,c_3;t_1,t_2,t_3)
 =(0,0,0;\alpha_3,\alpha_3,\alpha_3).
\tag{N35.14}
\]
\end{theorem}

\begin{proof}
The exact decomposition \textup{(N34.5)} and the coefficient bounds
\textup{(N34.15)} give
\[
 \mathcal S_{\rm gen}
 \leq-250r+M_4r^2+M_6r^3+M_8r^4.
\tag{N35.15}
\]
The bracket obtained after factoring \(-r\) is decreasing in \(r\).
At \(r=1/66\), exact rational arithmetic gives
\[
\begin{aligned}
 250-\frac{M_4}{66}-\frac{M_6}{66^2}-\frac{M_8}{66^3}
 &=\frac{506169078601}{399300000000}>0.
\end{aligned}
\tag{N35.16}
\]
Thus \(\mathcal S_{\rm gen}<0\) whenever \(0<r\leq1/66\).  For six real
roots the same quantity is a sum of six real squares, hence is
nonnegative.  Therefore \(r=0\), and nonnegativity of \(A\) and the three
height excesses gives \textup{(N35.14)}.
\end{proof}

\begin{corollary}[First exact annular pruning]
\label{cor:f2v35-annular-pruning}
Every nonexceptional tuple in \(\mathfrak C_{\rm cand}\) satisfies
\[
 \boxed{A+\sum_j(t_j-\alpha_3)>\frac1{66}.}
\tag{N35.17}
\]
In particular the complete rational shell
\[
 \frac1{3000}<A+\sum_j(t_j-\alpha_3)\leq\frac1{66}
\tag{N35.18}
\]
left after V34 contains no candidate.  The certified collar radius has
increased by the factor
\[
 \frac{1/66}{1/3000}=\frac{500}{11}.
\tag{N35.19}
\]
The remaining compact low and collision targets may therefore impose
\(r\geq1/66\) before subdivision.
\end{corollary}

\begin{proof}
This is the contrapositive of Theorem~\ref{thm:f2v35-collar}, combined
with the V34 equality case.  Replacing the old closed constraint
\(r\geq1/3000\) by \(r\geq1/66\) preserves compactness and retains every
possible nonexceptional candidate.
\end{proof}

\begin{firewall}[An enlarged collar is not the global cover]
\label{fire:f2v35}
Theorem~\ref{thm:f2v35-collar} is an exact interval certificate, but it
only prunes the innermost part of the V33 charts.  It does not prove that
the four V26 moment obstructions cover the remaining region \(r\geq1/66\),
nor does it complete the infinite-product passage or the arithmetic
forcing step.  GRH remains open.
\end{firewall}

\begin{assessment}[V35 advance and V36 direction]
\label{ass:f2v35}
Regrouping before interval evaluation reduces the useful weight-four
majorant dramatically, while retaining weights six and eight prevents
their coefficients from being charged at order \(r^2\).  This gives the
first exact pruning of a nontrivial annulus in the five displacement
variables.

V36 should preserve this radial separation and subdivide only the angular
simplex.  Introduce
\[
 x=A/r,\qquad d_j=\delta_j/r,
 \qquad x+d_1+d_2+d_3=1,
\tag{N35.20}
\]
and a center-angle chart for \((u,v)/\sqrt A\).  On radial slabs beginning
at \(r=1/66\), interval coefficients can then exploit fixed simplex boxes
instead of the simultaneous worst cases used in \textup{(N35.9)}.  Any
unresolved box should be retained with exact rational endpoints and a
record of which of \((\mathcal S_{\rm gen},\mathcal C,\mathcal E,
\mathcal K,\mathcal M_3)\) failed to decide it.
\end{assessment}

\section*{Version V35 append-only record}
\addcontentsline{toc}{section}{Version V35 append-only record}

V35 was appended to validated V34, whose installed SHA--256 was
\[
 \texttt{747F97F7738D9A127FF2468931FEF37747FC95E92E5BC65703989320E9FDEF97}.
\tag{N35.21}
\]
It regroups the adapted-square coefficients over the V13 rational
algebraic box, proves the radius-\(1/66\) equality collar, and records a
dependency-free exact verifier.  After validation V35 replaces V34 as the
sole active numeric GRH note; the frozen \texttt{V100\_f2} archive is
unchanged.

% =============================================================================
% END APPENDED V35 MATERIAL
% =============================================================================

% =============================================================================
% BEGIN APPENDED V36 MATERIAL
% =============================================================================

\section{V36: exact angular--simplex allocation}
\label{sec:f2v36-simplex-allocation}

V35 grouped the algebraic coefficients correctly, but its monomial bound
still allowed \(A\) and every height excess to consume the entire radial
budget simultaneously.  This continuation carries out the radial--angular
separation proposed there without subdivision.  Two elementary sharp
center inequalities and weighted AM--GM give a closed rational maximum for
each of the \(102\) outer monomials.  The resulting certificate extends the
rigidity radius from \(1/66\) to \(1/20\).

Retain
\[
 A=u^2+uv+v^2,\qquad
 \delta_j=t_j-\alpha_3\geq0,\qquad
 r=A+\delta_1+\delta_2+\delta_3.
\tag{N36.1}
\]

\subsection{The center and allocation inequalities}

\begin{lemma}[Exact centered-plane monomial bound]
\label{lem:f2v36-center-monomial}
For all real \(u,v\),
\[
 |uv|\leq A,\qquad
 u^2\leq\frac{4A}{3},\qquad
 v^2\leq\frac{4A}{3}.
\tag{N36.2}
\]
If \(m+n=2k\), then
\[
 \boxed{
 |u|^m|v|^n
 \leq
 \left(\frac43\right)^{|m-n|/2}A^k.}
\tag{N36.3}
\]
\end{lemma}

\begin{proof}
If \(uv\geq0\), then \(A-|uv|=u^2+v^2\geq0\); if \(uv<0\), then
\(A-|uv|=(u+v)^2\geq0\).  Also
\[
 A-\frac34u^2=\left(v+\frac u2\right)^2\geq0,
 \qquad
 A-\frac34v^2=\left(u+\frac v2\right)^2\geq0.
\tag{N36.4}
\]
When \(m\geq n\), write
\[
 |u|^m|v|^n=|uv|^n|u|^{m-n}
 \leq A^n\left(\frac{4A}{3}\right)^{(m-n)/2};
\tag{N36.5}
\]
the exponent of \(A\) is \((m+n)/2=k\).  The case \(n\geq m\) is
symmetric.
\end{proof}

\begin{lemma}[Four-coordinate simplex maximum]
\label{lem:f2v36-simplex-maximum}
Let \(k,e_1,e_2,e_3\) be nonnegative integers and put
\(p=k+e_1+e_2+e_3>0\).  Under \textup{(N36.1)},
\[
 A^k\prod_{j=1}^3\delta_j^{e_j}
 \leq
 \mu(k,\bm e)r^p,
\tag{N36.6}
\]
where, with \(0^0:=1\),
\[
 \mu(k,\bm e):=
 \left(\frac{k}{p}\right)^k
 \prod_{j=1}^3\left(\frac{e_j}{p}\right)^{e_j}.
\tag{N36.7}
\]
\end{lemma}

\begin{proof}
For \(r>0\), divide the four nonnegative coordinates
\((A,\delta_1,\delta_2,\delta_3)\) by their sum \(r\).  Weighted AM--GM,
or the Lagrange multiplier equations on the positive face, shows that the
product is maximal at the proportions
\((k,e_1,e_2,e_3)/p\).  Vanishing exponents are obtained by continuity.
The assertion at \(r=0\) is immediate.
\end{proof}

\subsection{A sharper weighted remainder ledger}

For each regrouped coefficient \(g(\alpha_3,s)\) in V35, let \(B_g\) be
the rational interval modulus constructed there.  If its outer monomial
has exponents \((m,n,e_1,e_2,e_3)\), combine
Lemmas~\ref{lem:f2v36-center-monomial} and
\ref{lem:f2v36-simplex-maximum}.  Its contribution at weight \(2p\) is at
most
\[
 B_g
 \left(\frac43\right)^{|m-n|/2}
 \mu\!\left(\frac{m+n}{2},\bm e\right)r^p.
\tag{N36.8}
\]

\begin{lemma}[Angular--simplex remainder constants]
\label{lem:f2v36-remainder}
For every \(r\geq0\), the weighted pieces in \textup{(N35.3)} satisfy
\[
 |\mathcal R_4|\leq L_4r^2,\qquad
 |\mathcal R_6|\leq L_6r^3,\qquad
 |\mathcal R_8|\leq L_8r^4,
\tag{N36.9}
\]
where
\[
 \boxed{
 L_4=\frac{69400059181}{15000000},\qquad
 L_6=\frac{132487063}{33750},\qquad
 L_8=\frac{1986323}{1296}.}
\tag{N36.10}
\]
\end{lemma}

\begin{proof}
Apply \textup{(N36.8)} to all \(102\) regrouped outer monomials and sum
separately at weights four, six, and eight.  Every factor is rational and
exact; the three sums are precisely \textup{(N36.10)}.
\end{proof}

\begin{remark}[Two-stage exact verifier]
\label{rem:f2v36-verifier}
The source
\begin{center}
\path{scripts/verify_grh_v36_simplex_allocation.py}
\end{center}
loads the already asserted V35 coefficient ledger, applies
\textup{(N36.8)} monomial by monomial, and asserts the three constants and
the final margin.  It has \(66\) lines, \(2122\) bytes, and SHA--256
\[
\texttt{CE24E6D3F5928F44645FC8545668C4448EE77E55C0704EE8C40262CEC750159F}.
\tag{N36.11}
\]
Both stages terminate with \texttt{VERIFIED}; no floating-point sign
decision is used.
\end{remark}

\subsection{The radius-\texorpdfstring{\(1/20\)}{1/20} collar}

\begin{theorem}[Angular--simplex equality collar]
\label{thm:f2v36-collar}
Assume
\[
 t_j\geq\alpha_3,\qquad
 A+\sum_{j=1}^3(t_j-\alpha_3)\leq\frac1{20}.
\tag{N36.12}
\]
If \(\mathcal Q_{\bm c,\bm t}\) is real-rooted, then
\[
 (\bm c,\bm t)=(\bm0;(\alpha_3,\alpha_3,\alpha_3)).
\tag{N36.13}
\]
\end{theorem}

\begin{proof}
The V34 linear ledger and Lemma~\ref{lem:f2v36-remainder} give
\[
 \mathcal S_{\rm gen}
 \leq-r\bigl(250-L_4r-L_6r^2-L_8r^3\bigr).
\tag{N36.14}
\]
The bracket is decreasing for \(r\geq0\).  At \(r=1/20\), its exact value
is
\[
 250-\frac{L_4}{20}-\frac{L_6}{20^2}-\frac{L_8}{20^3}
 =\frac{280617397877}{32400000000}>0.
\tag{N36.15}
\]
Therefore \(\mathcal S_{\rm gen}<0\) for \(0<r\leq1/20\), contradicting
its representation as a sum of six real squares when the heat sextic is
real-rooted.  Hence \(r=0\), which is exactly \textup{(N36.13)}.
\end{proof}

\begin{corollary}[Second exact annular pruning]
\label{cor:f2v36-annular-pruning}
Every nonexceptional tuple in \(\mathfrak C_{\rm cand}\) satisfies
\[
 \boxed{A+\sum_j(t_j-\alpha_3)>\frac1{20}.}
\tag{N36.16}
\]
Thus the full shell
\[
 \frac1{66}<A+\sum_j(t_j-\alpha_3)\leq\frac1{20}
\tag{N36.17}
\]
contains no candidate, and the V33 compact charts may impose
\(r\geq1/20\) before any further interval subdivision.
\end{corollary}

\begin{proof}
Take the contrapositive of Theorem~\ref{thm:f2v36-collar} and retain the
symmetric equality case.
\end{proof}

\begin{firewall}[Simplex pruning is not the remaining sign cover]
\label{fire:f2v36}
V36 exactly optimizes the allocation used in the adapted-square remainder,
but still applies the triangle inequality between different outer
monomials.  It does not exclude the compact region \(r\geq1/20\) and does
not prove the four-certificate cover, infinite-product passage, or
arithmetic forcing step.  GRH remains open.
\end{firewall}

\begin{assessment}[V36 advance and V37 direction]
\label{ass:f2v36}
The radial and simplex variables are now separated analytically; no angular
grid was needed to extend the effective collar from \(1/66\) to \(1/20\).
Further repetitions of coefficientwise triangle bounds are likely to give
diminishing returns.

V37 should move to the actual residual set \(r\geq1/20\).  The first
target should be a rational radial slab on the ordered low chart.  Use the
gap coordinates \((g_1,g_2)\), discard boxes violating the moment budget,
and test \(\mathcal S_{\rm gen}\) followed by the four V26 moment
certificates.  A useful output is either a certified slab or the exact
rational endpoints of every unresolved box; another global absolute-value
collar should be attempted only if that ledger exposes a specific angular
cancellation.
\end{assessment}

\section*{Version V36 append-only record}
\addcontentsline{toc}{section}{Version V36 append-only record}

V36 was appended to validated V35, whose installed SHA--256 was
\[
 \texttt{143244510CA86CCA44689FB661CD32EC93E63DAF45F621C275C0CAC0FE9B5407}.
\tag{N36.18}
\]
It proves exact center-monomial and four-coordinate allocation bounds,
uses them to establish the radius-\(1/20\) equality collar, and records a
two-stage rational verifier.  After validation V36 replaces V35 as the
sole active numeric GRH note; the frozen \texttt{V100\_f2} archive is
unchanged.

% =============================================================================
% END APPENDED V36 MATERIAL
% =============================================================================

% =============================================================================
% BEGIN APPENDED V37 MATERIAL
% =============================================================================

\section{V37: a twenty-one-slope fourth-minor horn}
\label{sec:f2v37-sharp-horn}

V36 supplied a centered-plane monomial inequality that was unavailable in
V28.  Applying it to the exact V28 decomposition of the fourth principal
moment minor reduces the old slope \(171\) to \(21\).  This is not another
local-collar enlargement: it excludes a semialgebraic large-center wedge
inside the residual region \(r\geq1/20\), especially on the low-height
layers nearest the threshold.

Retain the normalized centers
\[
 x_1=x,\qquad x_2=y,\qquad x_3=-x-y,
 \qquad A_{\bm x}=x^2+xy+y^2=1,
\tag{N37.1}
\]
put \(c_j=\lambda x_j\), and let
\(T=\max_jt_j\).  The V28 identity is
\[
\begin{aligned}
 \mathcal M_3(\lambda\bm x,\bm t)
={}&-48\lambda^6
 \left(\sum_{j=1}^3t_j\Delta_j^2-9\right)\\
 &+\lambda^4R_4(\bm x,\bm t)
 +\lambda^2R_2(\bm x,\bm t)+R_0(\bm t).
\end{aligned}
\tag{N37.2}
\]

\subsection{Angular coefficient norms}

For a polynomial whose center monomials all have the same even total
degree, define
\[
 \|R\|_{\angle}:=
 \sum_\nu |r_\nu|
 \left(\frac43\right)^{|a_\nu-b_\nu|/2},
\tag{N37.3}
\]
where \(r_\nu x^{a_\nu}y^{b_\nu}\bm t^{\bm d_\nu}\) is a monomial of
\(R\).  Lemma~\ref{lem:f2v36-center-monomial}, with \(A_{\bm x}=1\),
gives
\[
 |x|^{a_\nu}|y|^{b_\nu}
 \leq\left(\frac43\right)^{|a_\nu-b_\nu|/2}.
\tag{N37.4}
\]

\begin{lemma}[Re-audited fourth-minor remainders]
\label{lem:f2v37-remainders}
The exact V28 remainder polynomials satisfy
\[
 \boxed{
 \|R_4\|_{\angle}=162720,\qquad
 \|R_2\|_{\angle}=251424,\qquad
 \|R_0\|_{\angle}=131616.}
\tag{N37.5}
\]
Consequently, whenever \(T\geq1\) and
\(0\leq t_j\leq T\),
\[
 |R_4|\leq162720T^2,\qquad
 |R_2|\leq251424T^3,\qquad
 |R_0|\leq131616T^4.
\tag{N37.6}
\]
\end{lemma}

\begin{proof}
Construct the heat sextic exactly in
\(\mathbb Z[x,y,t_1,t_2,t_3,\lambda][z]\), generate Newton sums through
\(p_6\), and form
\[
 \mathcal M_3=6p_2p_6-6p_4^2-p_2p_3^2.
\tag{N37.7}
\]
After the substitution \(x_3=-x-y\), its \(163\) monomials split by
\(\lambda\)-degree as
\[
\begin{array}{c|rrrr}
\lambda\text{-degree}&0&2&4&6\\ \hline
\text{monomials}&29&60&50&24.
\end{array}
\tag{N37.8}
\]
The center degrees at the first three entries are respectively
\(0,2,4\); their maximum height degrees are \(4,3,2\).  Summing the exact
rational weights \textup{(N37.3)} gives, in the order
\((R_4,R_2,R_0)\),
\[
 162720,\qquad251424,\qquad131616.
\tag{N37.9}
\]
Equation \textup{(N37.4)} bounds the center factors.  Since
\(T\geq1\), every height monomial of degree at most \(d\) is bounded by
\(T^d\), proving \textup{(N37.6)}.
\end{proof}

\begin{remark}[Exact verifier]
\label{rem:f2v37-verifier}
The dependency-free source
\begin{center}
\path{scripts/verify_grh_v37_sharp_m3_horn.py}
\end{center}
constructs the heat polynomial and Newton sums from rational polynomial
dictionaries, asserts the four part counts in \textup{(N37.8)}, the three
norms in \textup{(N37.5)}, and the margin below.  It has \(88\) lines,
\(2818\) bytes, and SHA--256
\[
\texttt{5082A82DDB2411918DCAF23545BD68D32211359C192A98BD58EB070F1061C555}.
\tag{N37.10}
\]
It terminates with \texttt{VERIFIED}; all asserted quantities are exact
rationals.
\end{remark}

\subsection{Joint use of the three remainder levels}

\begin{theorem}[Twenty-one-slope fourth-minor exclusion]
\label{thm:f2v37-horn}
Let \(A_{\bm x}=1\), let \(T\geq1\), and suppose
\[
 \alpha_3\leq t_j\leq T\qquad(1\leq j\leq3).
\tag{N37.11}
\]
If
\[
 \lambda\geq21T,
\tag{N37.12}
\]
then \(\mathcal M_3(\lambda\bm x,\bm t)<0\).  Hence the corresponding
heat sextic is not real-rooted.
\end{theorem}

\begin{proof}
The derivative-square identity and the lower height bound give, exactly as
in V28,
\[
 -48\lambda^6
 \left(\sum_jt_j\Delta_j^2-9\right)
 <-\gamma\lambda^6,
 \qquad \gamma=\frac{254718}{625}.
\tag{N37.13}
\]
Unlike V28, do not assign one sixth of this margin to each remainder.
Lemma~\ref{lem:f2v37-remainders} and \(\lambda\geq21T\), with
\(T\geq1\), give the joint bound
\[
\begin{aligned}
 \frac{|\lambda^4R_4+\lambda^2R_2+R_0|}{\lambda^6}
 &\leq
 \frac{162720}{21^2}
 +\frac{251424}{21^4}
 +\frac{131616}{21^6}.
\end{aligned}
\tag{N37.14}
\]
The remaining exact leading margin is
\[
 \frac{254718}{625}
 -\frac{162720}{21^2}
 -\frac{251424}{21^4}
 -\frac{131616}{21^6}
 =\frac{222008456542}{5955980625}>0.
\tag{N37.15}
\]
Substitution in \textup{(N37.2)} proves \(\mathcal M_3<0\).  A principal
moment minor of six real roots must be nonnegative.
\end{proof}

\begin{corollary}[Sharpened global height-growth necessity]
\label{cor:f2v37-growth}
For every tuple in \(\mathfrak C_{\rm cand}\) with
\(A=A_{\bm c}>0\),
\[
 \boxed{\sqrt A<21\max_jt_j.}
\tag{N37.16}
\]
Equivalently, the rational semialgebraic wedge
\[
 A\geq441\left(\max_jt_j\right)^2
\tag{N37.17}
\]
contains no candidate.
\end{corollary}

\begin{proof}
Take \(\lambda=\sqrt A\) and
\(\bm x=\bm c/\lambda\), so \(A_{\bm x}=1\), and apply
Theorem~\ref{thm:f2v37-horn}.  Here \(T>1\) because
\(T\geq\alpha_3>1\).
\end{proof}

\begin{corollary}[Explicit pruning inside the low chart]
\label{cor:f2v37-low-pruning}
In the V33 small low chart, every candidate must obey both
\(r>1/20\) and \textup{(N37.16)}.  In particular,
\[
\begin{array}{c|c}
\text{height layer}&\text{necessary center bound}\\ \hline
\max_jt_j\leq9/4&A<35721/16\\
\max_jt_j\leq3&A<3969\\
\max_jt_j\leq4&A<7056.
\end{array}
\tag{N37.18}
\]
Boxes violating the corresponding bound can be discarded before any
Bernstein evaluation.
\end{corollary}

\begin{proof}
Square \(\sqrt A<21T\) at the three displayed rational values of \(T\).
The condition \(r>1/20\) is Corollary~\ref{cor:f2v36-annular-pruning}.
\end{proof}

\begin{firewall}[A sharper horn leaves the finite middle]
\label{fire:f2v37}
The new horn removes a genuine part of the residual V33 low chart and
improves the global V28 slope by a factor exceeding eight.  It does not
cover the intermediate wedge \(\sqrt A<21T\), the collision charts, the
infinite-product passage, or the arithmetic forcing step.  GRH remains
open.
\end{firewall}

\begin{assessment}[V37 advance and V38 direction]
\label{ass:f2v37}
The V28 bottleneck was mainly its angular coefficient norm, not the
leading derivative-square identity.  V36 supplies the missing exact
center geometry, and joint remainder accounting avoids another factor of
roughly \(\sqrt6\).

V38 should use \(q=\sqrt A/T\in[0,21)\) as a routing variable on the
remaining low chart.  The first exact slab should be a near-threshold
height layer, for example \(T\leq9/4\), with
\(1/20\leq r\) and \(A<35721/16\).  Subdivide the ordered gap ratio and
the normalized height simplex, testing \(\mathcal S_{\rm gen}\) first and
then \((\mathcal C,\mathcal E,\mathcal K,\mathcal M_3)\).  Unresolved
boxes must be retained exactly; the already excluded region \(q\geq21\)
must not be resampled.
\end{assessment}

\section*{Version V37 append-only record}
\addcontentsline{toc}{section}{Version V37 append-only record}

V37 was appended to validated V36, whose installed SHA--256 was
\[
 \texttt{77A447552331A332FD286330557AF3B232E21B497ED0BEEAEE8AF39866577EB0}.
\tag{N37.19}
\]
It re-audits the V28 fourth-minor remainders with the V36 angular norm,
proves the slope-\(21\) horn, and records its dependency-free exact
verifier.  After validation V37 replaces V36 as the sole active numeric
GRH note; the frozen \texttt{V100\_f2} archive is unchanged.

% =============================================================================
% END APPENDED V37 MATERIAL
% =============================================================================

% =============================================================================
% BEGIN APPENDED V38 MATERIAL
% =============================================================================

\section{V38: the first certified near-threshold slab}
\label{sec:f2v38-near-threshold}

V37 reduced the large-center slope globally, but its use of a single
height scale \(T\) still discarded cancellations among the height
monomials.  This continuation treats the first residual height slab
\[
 \alpha_3\leq t_j\leq\frac94\qquad(1\leq j\leq3),
\tag{N38.1}
\]
by regrouping every coefficient of the three V28 remainders as a
polynomial in \((t_1,t_2,t_3)\) before interval evaluation.  It proves the
absolute cutoff \(A<225\) throughout this slab.

Retain the normalized coordinates and exact fourth-minor decomposition
\textup{(N37.1)}--\textup{(N37.2)}.  Enlarge the algebraic height box
slightly to the rational cube
\[
 \mathcal T_0:=
 \left[\frac{19434}{10000},\frac94\right]^3.
\tag{N38.2}
\]
The V13 enclosure \(19434/10000<\alpha_3\) shows that every tuple in
\textup{(N38.1)} lies in \(\mathcal T_0\).

\subsection{Height-regrouped angular norms}

For each fixed center monomial \(x^ay^b\), write the corresponding
coefficient of a remainder as
\[
 h_{a,b}(\bm t)=
 \sum_{\bm d}q_{\bm d}
 t_1^{d_1}t_2^{d_2}t_3^{d_3}.
\tag{N38.3}
\]
Because every coordinate of \(\mathcal T_0\) is positive, a rational
interval for a term with \(q_{\bm d}>0\) is obtained by evaluating at the
two diagonal endpoints of \(\mathcal T_0\); for
\(q_{\bm d}<0\), reverse the endpoints.  Summing those term intervals gives
\([\underline h_{a,b},\overline h_{a,b}]\).  Put
\[
 B_{a,b}:=\max\{|\underline h_{a,b}|,
                    |\overline h_{a,b}|\}.
\tag{N38.4}
\]

\begin{lemma}[Near-threshold remainder ledger]
\label{lem:f2v38-remainders}
On \(A_{\bm x}=1\) and \(\bm t\in\mathcal T_0\),
\[
 |R_4|\leq H_4,\qquad |R_2|\leq H_2,
 \qquad |R_0|\leq H_0,
\tag{N38.5}
\]
where
\[
\boxed{
\begin{aligned}
 H_4&=\frac{65845066257}{781250},\\
 H_2&=\frac{39216545286651}{488281250},\\
 H_0&=\frac{945786604045338711}{39062500000000}.
\end{aligned}}
\tag{N38.6}
\]
\end{lemma}

\begin{proof}
After height regrouping, \(R_4,R_2,R_0\) contain respectively
\(5,3,1\) distinct center monomials.  Apply the rational intervals
\textup{(N38.3)}--\textup{(N38.4)} and then the V36 angular factor to get
\[
 \sum_{a,b}B_{a,b}
 \left(\frac43\right)^{|a-b|/2}.
\tag{N38.7}
\]
Exact addition gives, in order, the three fractions in
\textup{(N38.6)}.  This proves \textup{(N38.5)} by the centered-plane
monomial bound \textup{(N37.4)}.
\end{proof}

\begin{remark}[Exact verifier]
\label{rem:f2v38-verifier}
The source
\begin{center}
\path{scripts/verify_grh_v38_near_threshold_slab.py}
\end{center}
loads the independently verified V37 decomposition, regroups by center
monomial, evaluates every height coefficient over the rational cube
\(\mathcal T_0\), and asserts the group counts, the three fractions in
\textup{(N38.6)}, and the cutoff margin below.  It has \(51\) lines,
\(1573\) bytes, and SHA--256
\[
\texttt{A5987C53E4DF77D937D99610BC06F5982790B3DFC3613BB4292B9C63A168FED4}.
\tag{N38.8}
\]
The V37 and V38 stages both terminate with \texttt{VERIFIED}; no
floating-point sign is used.
\end{remark}

\subsection{The absolute center cutoff}

\begin{theorem}[Near-threshold center exclusion]
\label{thm:f2v38-center-cutoff}
Suppose \(A_{\bm x}=1\) and \textup{(N38.1)} holds.  If
\[
 \lambda\geq15,
\tag{N38.9}
\]
then \(\mathcal M_3(\lambda\bm x,\bm t)<0\).  Therefore the heat sextic
is not real-rooted.
\end{theorem}

\begin{proof}
The leading estimate \textup{(N37.13)} and
Lemma~\ref{lem:f2v38-remainders} give
\[
 \frac{\mathcal M_3}{\lambda^6}
 <-\gamma+\frac{H_4}{\lambda^2}
            +\frac{H_2}{\lambda^4}
            +\frac{H_0}{\lambda^6},
 \qquad \gamma=\frac{254718}{625}.
\tag{N38.10}
\]
The right-hand remainder is decreasing in \(\lambda>0\).  At
\(\lambda=15\), exact arithmetic gives the positive leading margin
\[
 \gamma-\frac{H_4}{15^2}-\frac{H_2}{15^4}-\frac{H_0}{15^6}
 =\frac{19149737458271586641}{610351562500000000}>0.
\tag{N38.11}
\]
Thus \(\mathcal M_3<0\).  This contradicts the nonnegativity of every
principal moment minor for six real roots.
\end{proof}

\begin{corollary}[Certified first residual slab]
\label{cor:f2v38-first-slab}
Every candidate satisfying \textup{(N38.1)} obeys
\[
 \boxed{A<225,\qquad |c_j|<18\quad(1\leq j\leq3).}
\tag{N38.12}
\]
Every nonexceptional candidate in this slab also obeys
\[
 A+\sum_j(t_j-\alpha_3)>\frac1{20}.
\tag{N38.13}
\]
Hence its remaining exact search domain is contained in the compact box
\textup{(N38.1)}, \textup{(N38.12)}--\textup{(N38.13)}.
\end{corollary}

\begin{proof}
If \(A\geq225\), set \(\lambda=\sqrt A\) and
\(\bm x=\bm c/\lambda\); Theorem~\ref{thm:f2v38-center-cutoff} excludes
the tuple.  The centered-plane inequality
\(c_j^2\leq4A/3<300<18^2\) gives the center bounds.  Finally
\textup{(N38.13)} is Corollary~\ref{cor:f2v36-annular-pruning}.
\end{proof}

\begin{firewall}[One complete height slab is not the full low chart]
\label{fire:f2v38}
V38 certifies an absolute center cutoff throughout all unequal-height
configurations with \(\max_jt_j\leq9/4\); it is not merely a sampled
statement.  Heights above \(9/4\), the finite inner portion \(A<225\),
the collision charts, the infinite-product passage, and arithmetic forcing
remain open.  GRH remains open.
\end{firewall}

\begin{assessment}[V38 advance and V39 direction]
\label{ass:f2v38}
The first proposed residual slab is now reduced from the V33 center radius
\(871/10\) to the rational radius \(18\).  Height regrouping, rather than
a uniform \(T^d\) majorant, is responsible for the improvement.

V39 should apply the same audited decomposition to the adjacent layer
\[
 \frac94\leq\max_jt_j\leq3,
\tag{N38.14}
\]
but should preserve which factor carries the maximum: use three ordered
height charts rather than the full cube \([19434/10000,3]^3\).  If the
three-chart interval ledger is still too coarse, bisect only the maximal
height coordinate.  The output should be a verified absolute center cutoff
and exact unresolved subboxes, not a floating scan.
\end{assessment}

\section*{Version V38 append-only record}
\addcontentsline{toc}{section}{Version V38 append-only record}

V38 was appended to validated V37, whose installed SHA--256 was
\[
 \texttt{E31FD97D62B9940378BBDCAB1B1B85F1B9EABDA9D00CA7A9DD84EB365F1FAEC9}.
\tag{N38.15}
\]
It constructs the first height-regrouped rational interval slab, proves
the absolute near-threshold cutoff \(A<225\), and records its exact
verifier.  After validation V38 replaces V37 as the sole active numeric
GRH note; the frozen \texttt{V100\_f2} archive is unchanged.

% =============================================================================
% END APPENDED V38 MATERIAL
% =============================================================================

% =============================================================================
% BEGIN APPENDED V39 MATERIAL
% =============================================================================

\section{V39: a six-box atlas for the adjacent height layer}
\label{sec:f2v39-adjacent-atlas}

V38 closed the large-center part of the layer
\(\max_jt_j\leq9/4\).  This continuation treats the adjacent layer
\[
 \frac94\leq T:=\max_jt_j\leq3
\tag{N39.1}
\]
without replacing the identity of the maximal factor by a full cube.  A
single bisection at \(21/8\), together with the three possible maximal
factors, gives six rational boxes.  Exact height regrouping proves
\(A<324\) on the lower half and \(A<441\) on the upper half.

Put
\[
 \ell:=\frac{19434}{10000},\qquad
 I_1:=\left[\frac94,\frac{21}{8}\right],
 \qquad I_2:=\left[\frac{21}{8},3\right].
\tag{N39.2}
\]
For \(q\in\{1,2\}\) and \(j\in\{1,2,3\}\), define the enlarged chart
\(\mathcal T_{q,j}\) by requiring
\[
 t_j\in I_q,qquad
 \ell\leq t_k\leq\sup I_q\quad(k\neq j).
\tag{N39.3}
\]

\begin{lemma}[Six-box covering]
\label{lem:f2v39-cover}
Every height triple satisfying \(t_j\geq\alpha_3\) and
\textup{(N39.1)} belongs to at least one of the six boxes
\(\mathcal T_{q,j}\).
\end{lemma}

\begin{proof}
Choose \(j\) with \(t_j=T\).  If \(T\leq21/8\), use
\(\mathcal T_{1,j}\); otherwise use \(\mathcal T_{2,j}\).  The other
two heights are at most \(T\), while
\(t_k\geq\alpha_3>\ell\).  At the shared endpoint either box applies.
\end{proof}

\subsection{Exact remainder ledgers on the six boxes}

Apply the V38 height-coefficient interval construction separately on each
\(\mathcal T_{q,j}\), followed by the V36 angular weights.  Write
\(H^{(q,j)}_4,H^{(q,j)}_2,H^{(q,j)}_0\) for the resulting bounds on
\(|R_4|,|R_2|,|R_0|\).  Symmetry under exchanging the first two normalized
center coordinates makes the ledgers for \(j=1,2\) identical.

\begin{lemma}[Adjacent-layer interval ledger]
\label{lem:f2v39-ledger}
On the lower boxes \(q=1\),
\[
\begin{aligned}
 (H^{(1,1)}_4,H^{(1,1)}_2,H^{(1,1)}_0)
 &=(H^{(1,2)}_4,H^{(1,2)}_2,H^{(1,2)}_0)\\
 &=\left(
 \frac{49197364173}{390625},
 \frac{1159675004705787}{7812500000},
 \frac{128258354851103007}{2441406250000}
 \right),\\
 (H^{(1,3)}_4,H^{(1,3)}_2,H^{(1,3)}_0)
 &=\left(
 \frac{196988708367}{1562500},
 \frac{551801321023671}{3906250000},
 \frac{128258354851103007}{2441406250000}
 \right).
\end{aligned}
\tag{N39.4}
\]
On the upper boxes \(q=2\),
\[
\begin{aligned}
 (H^{(2,1)}_4,H^{(2,1)}_2,H^{(2,1)}_0)
 &=(H^{(2,2)}_4,H^{(2,2)}_2,H^{(2,2)}_0)\\
 &=\left(
 \frac{132689976471}{781250},
 \frac{54962715990591}{244140625},
 \frac{3294828398489012487}{39062500000000}
 \right),\\
 (H^{(2,3)}_4,H^{(2,3)}_2,H^{(2,3)}_0)
 &=\left(
 \frac{133138298871}{781250},
 \frac{101561943068037}{488281250},
 \frac{3294828398489012487}{39062500000000}
 \right).
\end{aligned}
\tag{N39.5}
\]
\end{lemma}

\begin{proof}
For each remainder, regroup by its fixed center monomial.  There are
respectively \(5,3,1\) groups for \(R_4,R_2,R_0\).  Evaluate each height
coefficient term outward on \textup{(N39.3)}, sum its rational interval,
take the interval modulus, multiply by
\((4/3)^{|a-b|/2}\), and add over center monomials.  Exact rational
addition gives \textup{(N39.4)}--\textup{(N39.5)}.
\end{proof}

\begin{remark}[Exact six-box verifier]
\label{rem:f2v39-verifier}
The source
\begin{center}
\path{scripts/verify_grh_v39_adjacent_height_atlas.py}
\end{center}
constructs all six interval ledgers from the verified V37 polynomial and
asserts every fraction in \textup{(N39.4)}--\textup{(N39.5)} and every
margin below.  It has \(99\) lines, \(3202\) bytes, and SHA--256
\[
\texttt{8CCEE123EB90513CA330C562BB727BAED64D0AEB156A6AB44335D396FE63D1F6}.
\tag{N39.6}
\]
It terminates with \texttt{VERIFIED}; no floating-point sign is used.
\end{remark}

\subsection{Two absolute center cutoffs}

For a ledger triple \(H=(H_4,H_2,H_0)\), put
\[
 \operatorname{mar}(L;H):=
 \frac{254718}{625}-\frac{H_4}{L^2}
 -\frac{H_2}{L^4}-\frac{H_0}{L^6}.
\tag{N39.7}
\]

\begin{theorem}[Adjacent-layer center exclusion]
\label{thm:f2v39-center-cutoff}
Suppose \(A_{\bm x}=1\), \(t_j\geq\alpha_3\), and let
\(T=\max_jt_j\).

\begin{enumerate}[label=\textup{(\roman*)}]
\item If \(9/4\leq T\leq21/8\) and \(\lambda\geq18\), then
\(\mathcal M_3<0\).
\item If \(21/8\leq T\leq3\) and \(\lambda\geq21\), then
\(\mathcal M_3<0\).
\end{enumerate}
In either case the heat sextic is not real-rooted.
\end{theorem}

\begin{proof}
By Lemma~\ref{lem:f2v39-cover}, the height triple lies in a corresponding
factor-labelled box.  The leading estimate is still
\(-\gamma\lambda^6\), with \(\gamma=254718/625\).  The smallest lower-box
margin at \(L=18\) is the \(j=3\) value
\[
 \operatorname{mar}(18;H^{(1,3)})
 =\frac{17517978317786179253}{1025156250000000000}>0.
\tag{N39.8}
\]
The smallest upper-box margin at \(L=21\) is likewise the \(j=3\) value
\[
 \operatorname{mar}(21;H^{(2,3)})
 =\frac{276362346660446977891}{13786992187500000000}>0.
\tag{N39.9}
\]
The \(j=1,2\) margins are larger and are independently asserted by the
verifier.  Since each remainder ratio decreases with \(\lambda>0\), the
two displayed positive margins and the exact decomposition
\textup{(N37.2)} prove \(\mathcal M_3<0\) on all six boxes.
\end{proof}

\begin{corollary}[Certified adjacent residual layer]
\label{cor:f2v39-adjacent-layer}
Every candidate with \(9/4\leq T\leq3\) satisfies
\[
\boxed{
\begin{array}{ll}
9/4\leq T\leq21/8:&A<324,\quad |c_j|<21,\\
21/8\leq T\leq3:&A<441,\quad |c_j|<25.
\end{array}}
\tag{N39.10}
\]
Every nonexceptional candidate also has \(r>1/20\).
\end{corollary}

\begin{proof}
Use \(A=\lambda^2\) in
Theorem~\ref{thm:f2v39-center-cutoff}.  The centered-plane inequality
\(c_j^2\leq4A/3\) gives respectively
\(c_j^2<432<21^2\) and \(c_j^2<588<25^2\).  The radial condition is
Corollary~\ref{cor:f2v36-annular-pruning}.
\end{proof}

\begin{firewall}[The adjacent atlas does not cover higher layers]
\label{fire:f2v39}
V39 is an exact cover of the complete height layer
\(9/4\leq\max_jt_j\leq3\), not a floating scan.  It leaves the finite
inner boxes in \textup{(N39.10)}, all heights above \(3\), the collision
charts, the infinite-product passage, and arithmetic forcing.  GRH remains
open.
\end{firewall}

\begin{assessment}[V39 advance and mandatory V40 sweep]
\label{ass:f2v39}
Together V38--V39 reduce every low-chart tuple with
\(\max_jt_j\leq3\) to \(|c_j|<25\), with the sharper radii \(18\) and
\(21\) on the first two sublayers.  Preserving the maximal factor and
bisecting only its coordinate avoids the much larger full-cube interval.

Before choosing the next analytic direction, V40 must perform the scheduled
ten-version sweep of all newer parallel notes in \texttt{latex\_notes}.
After incorporating any relevant ideas, the natural continuation is the
height layer \(3\leq T\leq4\), again with factor-labelled boxes, unless the
sweep supplies a stronger upstream certificate.
\end{assessment}

\section*{Version V39 append-only record}
\addcontentsline{toc}{section}{Version V39 append-only record}

V39 was appended to validated V38, whose installed SHA--256 was
\[
 \texttt{8B52523152B826C1D1C130DA66CA2022DDCE3ED27FF7AA53F898607342815262}.
\tag{N39.11}
\]
It constructs a six-box rational atlas for the adjacent height layer,
proves its two absolute center cutoffs, and records an exact chained
verifier.  After validation V39 replaces V38 as the sole active numeric
GRH note; the frozen \texttt{V100\_f2} archive is unchanged.

% =============================================================================
% END APPENDED V39 MATERIAL
% =============================================================================

% =============================================================================
% BEGIN APPENDED V40 MATERIAL
% =============================================================================

\section{V40: companion sweep and the height-three-to-four atlas}
\label{sec:f2v40-sweep-atlas}

This is the scheduled ten-version checkpoint.  The companion sweep was
performed before the new GRH algebra.  One active SM file advanced during
inspection, so the sweep was refreshed and frozen only after the final
active files had stable names and hashes.

\subsection{Mandatory companion sweep}

The four-file snapshot actually read is:
\begin{sloppypar}
\raggedright
\begin{description}[leftmargin=2.8em,style=nextline]
\item[SM V71.]
\path{Renewal_SM_Einstein_Continuum_Research_Notes_V71.tex};
\(22205\) lines, \(783694\) bytes; SHA--256
\texttt{18AF673F1C376111524A8C6962E86853}
\allowbreak\texttt{B00C4A3BB551C1B8ECA2867C67E225C9}.
\item[Yang--Mills V12.]
\path{Renewal_Yang_Mills_Mass_Gap_Research_Notes_V12.tex};
\(5423\) lines, \(201801\) bytes; SHA--256
\texttt{1A1D41A289DA206CCB88023306DD33743}
\allowbreak\texttt{C46058F731345900765A724BA1DDCBE}.
\item[Parallel Yang--Mills V5.]
\path{Renewal_YM_Parallel_Research_Notes_V5.tex};
\(1351\) lines, \(54174\) bytes; SHA--256
\texttt{306F1BB84CFA8A8D47EA569EC17E9545}
\allowbreak\texttt{F9867C8D32B548EC9D9C444B4F3C0512}.
\item[Navier--Stokes V26.]
\path{Renewal_NS_Stability_Research_Notes_V26.tex};
\(5319\) lines, \(278283\) bytes; SHA--256
\texttt{82C887F8C2C69E4F28FAD7C3DC8DE5B}
\allowbreak\texttt{9099CB5A4BF431EBA1FFF3E91935978AD}.
\end{description}
\end{sloppypar}
\[
 \text{These full hashes fix the V40 sweep snapshot.}
\tag{N40.1}
\]

The new SM chain V69--V71 first forms measured-minus-ideal interval
objects before taking a modulus, then proves completeness by equality with
an explicitly enumerated key set, rejects duplicates and unresolved
payloads, and permits compressed representations only after exact rank
certification.  The new Yang--Mills chain V9--V12 likewise keeps live
correlations, compiles signed endpoint differences before absolute values,
deletes only exact zero coefficients, and retains dense exceptions when
low-rank compression is not certified.  The parallel Yang--Mills and
Navier--Stokes snapshots are unchanged from V30 and supply no new GRH
premise.

No companion theorem is imported.  The transferable rules are implemented
here as follows:
\begin{enumerate}[label=\textup{(\roman*)}]
\item every height coefficient is fully regrouped before its interval
      modulus is taken;
\item the factor-labelled chart keys are enumerated explicitly and proved
      to cover the requested layer; and
\item every one of the six ledgers, including the two nonsymmetric
      \(j=3\) cases, is asserted by the verifier rather than inferred from
      an unproved symmetry.
\end{enumerate}

\subsection{A declared six-box cover of \texorpdfstring{\(3\leq T\leq4\)}{3 <= T <= 4}}

Retain
\[
 \ell=\frac{19434}{10000},\qquad
 J_1=\left[3,\frac72\right],qquad
 J_2=\left[\frac72,4\right].
\tag{N40.2}
\]
For \(q\in\{1,2\}\) and \(j\in\{1,2,3\}\), let
\(\mathcal U_{q,j}\) be the rational box
\[
 t_j\in J_q,qquad
 \ell\leq t_k\leq\sup J_q\quad(k\neq j).
\tag{N40.3}
\]

\begin{lemma}[Complete height-three-to-four cover]
\label{lem:f2v40-cover}
Every triple with \(t_j\geq\alpha_3\) and
\[
 3\leq T:=\max_jt_j\leq4
\tag{N40.4}
\]
lies in at least one of the six boxes \(\mathcal U_{q,j}\).
\end{lemma}

\begin{proof}
Choose an index \(j\) attaining \(T\).  Route to \(q=1\) when
\(T\leq7/2\) and to \(q=2\) otherwise.  Every other height is between
\(\alpha_3>\ell\) and \(T\leq\sup J_q\).  The shared endpoint is covered
by both slabs.
\end{proof}

\subsection{Post-sweep exact interval ledger}

As in V38--V39, regroup the coefficient of each center monomial in the
V28 remainders, evaluate it outward on \(\mathcal U_{q,j}\), and only then
apply the angular factor \((4/3)^{|a-b|/2}\).  Denote the resulting bounds
by \(G^{(q,j)}_4,G^{(q,j)}_2,G^{(q,j)}_0\).

\begin{lemma}[Height-three-to-four remainder constants]
\label{lem:f2v40-ledger}
For the lower slab,
\[
\begin{aligned}
 (G^{(1,1)}_4,G^{(1,1)}_2,G^{(1,1)}_0)
 &=(G^{(1,2)}_4,G^{(1,2)}_2,G^{(1,2)}_0)\\
 &=\left(\frac{367854433567}{1562500},
 \frac{2740978035879537}{7812500000},
 \frac{5428772726538111237}{39062500000000}\right),\\
 (G^{(1,3)}_4,G^{(1,3)}_2,G^{(1,3)}_0)
 &=\left(\frac{184973845121}{781250},
 \frac{1259668771971171}{3906250000},
 \frac{5428772726538111237}{39062500000000}\right).
\end{aligned}
\tag{N40.5}
\]
For the upper slab,
\[
\begin{aligned}
 (G^{(2,1)}_4,G^{(2,1)}_2,G^{(2,1)}_0)
 &=(G^{(2,2)}_4,G^{(2,2)}_2,G^{(2,2)}_0)\\
 &=\left(\frac{476170928567}{1562500},
 \frac{3882994853537037}{7812500000},
 \frac{7960187923377013737}{39062500000000}\right),\\
 (G^{(2,3)}_4,G^{(2,3)}_2,G^{(2,3)}_0)
 &=\left(\frac{240318115121}{781250},
 \frac{1764286460623671}{3906250000},
 \frac{7960187923377013737}{39062500000000}\right).
\end{aligned}
\tag{N40.6}
\]
\end{lemma}

\begin{proof}
Each of \(R_4,R_2,R_0\) has respectively \(5,3,1\) regrouped center
coefficients.  Natural rational interval evaluation on
\textup{(N40.3)}, followed by the V36 angular inequality and exact
summation, gives \textup{(N40.5)}--\textup{(N40.6)}.  Repeated occurrences
of a height variable may widen the natural interval but cannot invalidate
inclusion; this is precisely the soundness distinction retained from the
SM V69 sweep.
\end{proof}

\begin{remark}[Exact six-box verifier]
\label{rem:f2v40-verifier}
The source
\begin{center}
\path{scripts/verify_grh_v40_height_3_4_atlas.py}
\end{center}
loads the V39 interval engine, evaluates the six declared keys, and asserts
every norm and margin.  It has \(60\) lines, \(1926\) bytes, and SHA--256
\[
\texttt{E73A5A4C4C8CB385AC00F5F582BEEF8B9D61D3B05C8741EA4EFC64021D953B08}.
\tag{N40.7}
\]
All chained stages terminate with \texttt{VERIFIED}.
\end{remark}

\begin{theorem}[Height-three-to-four center exclusion]
\label{thm:f2v40-center-cutoff}
Suppose \(A_{\bm x}=1\), \(t_j\geq\alpha_3\), and let
\(T=\max_jt_j\).
\begin{enumerate}[label=\textup{(\roman*)}]
\item If \(3\leq T\leq7/2\) and \(\lambda\geq25\), then
      \(\mathcal M_3<0\).
\item If \(7/2\leq T\leq4\) and \(\lambda\geq28\), then
      \(\mathcal M_3<0\).
\end{enumerate}
Thus the corresponding heat sextic is not real-rooted.
\end{theorem}

\begin{proof}
Use Lemma~\ref{lem:f2v40-cover} to choose a declared box.  Insert the
corresponding triple from Lemma~\ref{lem:f2v40-ledger} into the exact V28
decomposition and compare with the leading constant
\(\gamma=254718/625\).  The smallest lower-slab margin occurs at \(j=3\):
\[
 \gamma-\frac{G^{(1,3)}_4}{25^2}
 -\frac{G^{(1,3)}_2}{25^4}
 -\frac{G^{(1,3)}_0}{25^6}
 =\frac{266039461304797393138763}
 {9536743164062500000000}>0.
\tag{N40.8}
\]
The smallest upper-slab margin also occurs at \(j=3\):
\[
 \gamma-\frac{G^{(2,3)}_4}{28^2}
 -\frac{G^{(2,3)}_2}{28^4}
 -\frac{G^{(2,3)}_0}{28^6}
 =\frac{38877838423159748906609}
 {2689120000000000000000}>0.
\tag{N40.9}
\]
The four remaining margins are larger and separately asserted.  Since the
remainder ratios decrease with \(\lambda\), the fourth minor is negative
throughout the stated ranges.
\end{proof}

\begin{corollary}[Certified height ladder through four]
\label{cor:f2v40-height-ladder}
Every candidate with \(T\leq4\) satisfies the following exact ladder:
\[
\boxed{
\begin{array}{c|c|c}
T\text{ range}&A\text{ bound}&|c_j|\text{ bound}\\ \hline
{}[\alpha_3,9/4]&A<225&|c_j|<18\\
{}[9/4,21/8]&A<324&|c_j|<21\\
{}[21/8,3]&A<441&|c_j|<25\\
{}[3,7/2]&A<625&|c_j|<29\\
{}[7/2,4]&A<784&|c_j|<33.
\end{array}}
\tag{N40.10}
\]
Every nonexceptional candidate in this ladder also has \(r>1/20\).
\end{corollary}

\begin{proof}
The first three rows are V38--V39.  Theorem
\ref{thm:f2v40-center-cutoff} gives \(A<25^2\) and \(A<28^2\) in the last
two rows.  Since \(c_j^2\leq4A/3\), one has respectively
\(c_j^2<2500/3<29^2\) and \(c_j^2<3136/3<33^2\).  The radial condition is
V36.
\end{proof}

\begin{firewall}[The sweep and finite ladder do not prove GRH]
\label{fire:f2v40}
The mandatory sweep is complete and V40 adds an exact cover of the full
height layer \(3\leq T\leq4\).  It does not exclude the finite inner boxes,
heights above four, the collision charts, the infinite-product passage, or
the arithmetic forcing step.  GRH remains open.
\end{firewall}

\begin{assessment}[V40 advance and V41 direction]
\label{ass:f2v40}
The companion methods confirm that the present signed regrouping and
explicit chart-key architecture are the correct finite-certificate
semantics.  The low chart through height four now has center radius below
\(33\), compared with the original V33 radius \(871/10\).

V41 should first determine the crossover between the height-regrouped
\(\mathcal M_3\) cutoff and the fixed V33 center bound.  It should then
cover only the useful next height slabs, rather than extend the atlas into
ranges where \(|c_j|<871/10\) is already stronger.  Each slab must retain
the maximal-factor key and compile the full coefficient before taking its
interval modulus.
\end{assessment}

\section*{Version V40 append-only record}
\addcontentsline{toc}{section}{Version V40 append-only record}

V40 was appended to validated V39, whose installed SHA--256 was
\[
 \texttt{6A824E30B151661613D882A5571751F93C16176A8FEC0C432F5DC0A9CA5E9C73}.
\tag{N40.11}
\]
It records the mandatory stable-hash companion sweep, implements its
compile-before-modulus and finite-completeness rules, proves the six-box
height-three-to-four atlas, and records its exact verifier.  After
validation V40 replaces V39 as the sole active numeric GRH note; the frozen
\texttt{V100\_f2} archive is unchanged.

% =============================================================================
% END APPENDED V40 MATERIAL
% =============================================================================


% =============================================================================
% BEGIN APPENDED V41 MATERIAL
% =============================================================================

\section{V41: the useful low-chart crossover atlas}
\label{sec:f2v41-crossover}

V40 left a precise upstream question: for how long does the
height-regrouped fourth-minor estimate improve the height-independent V33
variance cutoff?  This version answers that question before spending more
boxes.  The answer is finite and exact.  Unit maximal-factor slabs improve
the V33 variance bound through \(T=14\), but the very next slab does not.

\subsection{Declared maximal-factor boxes}

Retain the lower height endpoint
\[
 \ell=\frac{19434}{10000}.
\tag{N41.1}
\]
For \(n\in\{4,\ldots,13\}\) and \(j\in\{1,2,3\}\), define
\[
 \mathcal B_{n,j}:=
 \left\{(t_1,t_2,t_3):
 t_j\in[n,n+1],\quad
 \ell\leq t_k\leq n+1\ (k\neq j)\right\}.
\tag{N41.2}
\]
The label \(j\) records a factor attaining the maximum; it is not erased
by a symmetry assumption.

\begin{lemma}[Complete routing through height fourteen]
\label{lem:f2v41-routing}
If \(t_k\geq\alpha_3>\ell\) and
\[
 4\leq T:=\max_kt_k\leq14,
\tag{N41.3}
\]
then the height triple belongs to at least one of the thirty declared
boxes \(\mathcal B_{n,j}\).
\end{lemma}

\begin{proof}
Choose \(j\) with \(t_j=T\).  If \(T<14\), take
\(n=\lfloor T\rfloor\); when \(T=14\), take \(n=13\).  Then
\(t_j\in[n,n+1]\), while for \(k\neq j\) one has
\(\ell<t_k\leq T\leq n+1\).  Integer endpoints lie in both adjacent
closed slabs, so there is no boundary gap.
\end{proof}

\subsection{Exact interval constants and least common cutoffs}

Write the exact V28 decomposition in the form
\[
 \mathcal M_3
 =-\frac{254718}{625}\lambda^6+R_4+R_2+R_0,
 \qquad
 |R_d|\leq H^{(n,j)}_d\lambda^d
 \quad(d=4,2,0)
\tag{N41.4}
\]
on \(\mathcal B_{n,j}\).  Here \(H^{(n,j)}_d\) is obtained exactly as in
V38--V40: collect the complete coefficient of each center monomial,
evaluate that signed height coefficient by rational natural interval
arithmetic, take its interval modulus only after the collection, multiply
by the V36 angular factor \((4/3)^{|a-b|/2}\), and then sum.  Define
\[
 \mu_{n,j}(L):=
 \frac{254718}{625}
 -\frac{H^{(n,j)}_4}{L^2}
 -\frac{H^{(n,j)}_2}{L^4}
 -\frac{H^{(n,j)}_0}{L^6}.
\tag{N41.5}
\]

\begin{lemma}[Thirty-chart cutoff ledger]
\label{lem:f2v41-cutoffs}
For each \(n=4,\ldots,13\), the following \(L_n\) is the least positive
integer for which the natural-interval certificate
\(\mu_{n,j}(L_n)>0\) holds for all three factor labels:
\[
\begin{array}{c|rrrrrrrrrr}
n&4&5&6&7&8&9&10&11&12&13\\ \hline
L_n&35&41&47&53&58&64&70&76&81&87\\
L_n^2&1225&1681&2209&2809&3364&
4096&4900&5776&6561&7569
\end{array}.
\tag{N41.6}
\]
Moreover all thirty certified margins are positive and their minimum is
\[
 \min_{\substack{4\leq n\leq13\\1\leq j\leq3}}
 \mu_{n,j}(L_n)
 =
 \frac{1400416823279849330423763}
 {1487058302500000000000000}
 >\frac9{10}.
\tag{N41.7}
\]
\end{lemma}

\begin{proof}
For every one of the thirty keys \((n,j)\), the exact calculation
constructs \(5,3,1\) regrouped center coefficients in degrees \(4,2,0\),
respectively.  Rational endpoint evaluation gives the three
\(H^{(n,j)}_d\), and substitution in \textup{(N41.5)} gives the asserted
positive fractions.  Repeating the same calculation at \(L_n-1\) gives
a nonpositive margin for at least one factor label in every row.  Thus
the displayed \(L_n\) is the least \emph{common integer cutoff certified
by this declared natural-interval method}.  This is a statement about the
certificate, not a claim that the interval extension is optimal.  The
smallest positive fraction is the displayed value, attained for
\(n=8,j=3\).
\end{proof}

\begin{remark}[Exact chained verifier]
\label{rem:f2v41-verifier}
The source
\begin{center}
\path{scripts/verify_grh_v41_low_chart_crossover.py}
\end{center}
loads the V39 polynomial and interval engine, enumerates all three factor
labels on every slab, asserts the \(5,3,1\) coefficient counts, the least
common cutoffs, all exact useful-slab margins, and the crossover
comparison below.  It has \(72\) lines, \(2547\) bytes, and SHA--256
\[
 \texttt{A0B5BB92BB1EAC4A1D728F5BD562AAFC}
 \allowbreak\texttt{7B477D404FFBDD0F1CD272EFE0835F77}.
\tag{N41.8}
\]
Every chained stage terminates with \texttt{VERIFIED}; no floating-point
sign decision is used.
\end{remark}

\subsection{The useful variance ladder}

\begin{theorem}[Thirty-chart fourth-minor exclusion]
\label{thm:f2v41-thirty-chart}
Suppose \(A_{\bm x}=1\), \(t_k\geq\alpha_3\), and
\(n\leq T:=\max_kt_k\leq n+1\) for some
\(n\in\{4,\ldots,13\}\).  If \(\lambda\geq L_n\), with \(L_n\) as in
\textup{(N41.6)}, then
\[
 \frac{\mathcal M_3}{\lambda^6}
 \leq-\mu_{n,j}(L_n)<0
\tag{N41.9}
\]
for a factor label \(j\) supplied by
Lemma~\ref{lem:f2v41-routing}.  Consequently the corresponding heat
sextic is not real-rooted.
\end{theorem}

\begin{proof}
Route the height triple to \(\mathcal B_{n,j}\).  Divide
\textup{(N41.4)} by \(\lambda^6\), use the three interval bounds there,
and note that each positive remainder ratio
\(H_4/\lambda^2,H_2/\lambda^4,H_0/\lambda^6\) decreases for
\(\lambda>0\).  Lemma~\ref{lem:f2v41-cutoffs} therefore proves
\textup{(N41.9)}.  Negativity of this fourth Hermite minor contradicts
real-rootedness by the already established minor criterion.
\end{proof}

\begin{corollary}[Certified candidate variance ladder]
\label{cor:f2v41-variance-ladder}
Every candidate in the indicated height slab satisfies
\[
\boxed{
\begin{array}{c|r@{\qquad}c|r}
T\text{ range}&A\text{ bound}&T\text{ range}&A\text{ bound}\\ \hline
{}[4,5]&A<1225&{}[9,10]&A<4096\\
{}[5,6]&A<1681&{}[10,11]&A<4900\\
{}[6,7]&A<2209&{}[11,12]&A<5776\\
{}[7,8]&A<2809&{}[12,13]&A<6561\\
{}[8,9]&A<3364&{}[13,14]&A<7569
\end{array}}
\tag{N41.10}
\]
where \(A=\lambda^2\).
\end{corollary}

\begin{proof}
If \(A\geq L_n^2\), then \(\lambda\geq L_n\) and
Theorem~\ref{thm:f2v41-thirty-chart} excludes the tuple.  Therefore a
candidate has \(A<L_n^2\); substitute the exact squares from
\textup{(N41.6)}.
\end{proof}

\subsection{The exact crossover with V33}

The height-independent V33 estimate already gives every low-chart
candidate
\[
 A<\left(\frac{871}{10}\right)^2
 =\frac{758641}{100}.
\tag{N41.11}
\]
To locate the crossover rather than guess it, the same three-factor
ledger was continued only far enough to pass this number:
\[
\begin{array}{c|rrrr}
T\text{ range}&{}[14,15]&{}[15,16]&{}[16,17]&{}[17,441/25]\\ \hline
\text{least common }L&93&98&104&107\\
L^2&8649&9604&10816&11449
\end{array}.
\tag{N41.12}
\]

\begin{proposition}[First non-improving unit slab]
\label{prop:f2v41-crossover}
For the present maximal-factor natural-interval ledger, the last
unit-height slab whose certified variance cutoff improves V33 is
\([13,14]\).  The crossover is exact:
\[
 87^2=7569
 <\frac{758641}{100}
 <93^2=8649.
\tag{N41.13}
\]
Thus extending this same unit-slab ledger above \(T=14\) cannot improve
the already proved V33 variance restriction.
\end{proposition}

\begin{proof}
The first inequality and the \(n=13\) certificate are
Lemma~\ref{lem:f2v41-cutoffs}.  Exact interval evaluation on
\([14,15]\) gives least common cutoff \(93\); its predecessor \(92\)
fails on at least one declared factor chart.  The second inequality in
\textup{(N41.13)} is integer arithmetic.  The later rows in
\textup{(N41.12)} have still larger squares, so none improves
\textup{(N41.11)}.
\end{proof}

Combining V33 with the new ladder, the useful conclusion is precisely
\[
 A<
 \begin{cases}
 L_n^2,&n\leq T\leq n+1,\quad4\leq n\leq13,\\
 758641/100,&14\leq T\leq441/25.
 \end{cases}
\tag{N41.14}
\]
This compares variance bounds only.  In particular, it does not replace
the direct V33 coordinate estimate \(|c_j|<871/10\) by a weaker bound
deduced from \(c_j^2\leq4A/3\).

\begin{firewall}[The crossover atlas is not a GRH proof]
\label{fire:f2v41}
V41 certifies a finite low-chart variance improvement and identifies
where that improvement stops.  It does not eliminate the surviving
finite center boxes, the near-threshold annulus, collision charts, the
infinite-product passage, or the arithmetic forcing step.  Explicitly,
the presently most useful unresolved near-threshold domain includes
\[
 \alpha_3\leq T\leq\frac94,\qquad
 r>\frac1{20},\qquad A<225.
\tag{N41.15}
\]
GRH remains open.
\end{firewall}

\begin{assessment}[V41 advance and V42 direction]
\label{ass:f2v41}
The fourth-minor atlas has now been pushed exactly to its useful
crossover, so blindly continuing the same height ladder would be
circular bookkeeping.  The next attack should go upstream into
\textup{(N41.15)}: subdivide by the center-gap ratio and the height
simplex, then test the signed certificates in the order
\(\mathcal S_{\mathrm{gen}},\mathcal C,\mathcal E,\mathcal K,\mathcal
M_3\).  Each subdivision must carry an explicit chart key and the final
ledger must list every uncovered box, with neither silent omission nor
duplicate ownership.
\end{assessment}

\section*{Version V41 append-only record}
\addcontentsline{toc}{section}{Version V41 append-only record}

V41 was appended byte-for-byte before the terminal marker of validated
V40, whose installed SHA--256 was
\[
 \texttt{9386BF13E92433156F5337592E38513B}
 \allowbreak\texttt{64B888D16594353DC9D04A1D27AC39C4}.
\tag{N41.16}
\]
It proves the thirty-chart variance ladder through height fourteen,
locates the first non-improving unit slab, and records the exact chained
verifier.  After validation V41 replaces V40 as the sole active numeric
GRH note; the frozen \texttt{V100\_f2} archive is unchanged.

% =============================================================================
% END APPENDED V41 MATERIAL
% =============================================================================



% =============================================================================
% BEGIN APPENDED V42 MATERIAL
% =============================================================================

\section{V42: an explicit unequal-height wedge}
\label{sec:f2v42-wedge}

V41 showed that continuing the height-only fourth-minor ladder past its
crossover would not improve the inherited compact box.  This version
therefore moves upstream to the low near-threshold domain and quantifies
the direction transverse to the equal-height face.  V15 excluded that
entire face for arbitrary centers.  The result below thickens it into an
explicit unequal-height wedge for every centered configuration with
\(A\leq3\).

\subsection{Mean height, imbalance, and the equal-height margin}

Assume throughout this section that
\(c_1+c_2+c_3=0\), put
\[
 A=\frac12\sum_{j=1}^3c_j^2,\qquad
 \bar t=\frac{t_1+t_2+t_3}{3},\qquad
 \eta_j=t_j-\bar t,\qquad
 E=\sum_{j=1}^3|\eta_j|.
\tag{N42.1}
\]
On the domain
\[
 0\leq A\leq3,\qquad
 \alpha_3\leq t_j\leq\frac94,
\tag{N42.2}
\]
define
\[
 X=\frac A3,\qquad
 U=\frac{\bar t-\alpha_3}{9/4-\alpha_3},\qquad
 W(X,U)=1-(1-X)^4(1-U)^4.
\tag{N42.3}
\]
When \(A>0\), also retain the V15 angular invariant
\(\rho=27B^2/(4A^3)\in[0,1]\); set \(\rho=0\) at \(A=0\).

\begin{lemma}[Quantitative equal-height negativity]
\label{lem:f2v42-base-margin}
Let \(\mathcal S_{\rm gen}\) be the adapted Hermite square of V16.
At the equal-height point
\(\bm t^0=(\bar t,\bar t,\bar t)\),
\[
 \mathcal S_{\rm gen}(\bm c,\bm t^0)
 \leq-57\,W(X,U)
 \leq-57\max\{X,U\}.
\tag{N42.4}
\]
\end{lemma}

\begin{proof}
The V15 low-chart Bernstein polynomial has tridegree \((4,4,1)\).
Its only zero coefficients are \(b_{000}\) and \(b_{001}\); every other
coefficient is strictly smaller than \(-57\).  Since the Bernstein basis
is nonnegative and sums to one, the total basis weight carried by those
two zero coefficients is
\[
 B_0^4(X)B_0^4(U)\bigl(B_0^1(\rho)+B_1^1(\rho)\bigr)
 =(1-X)^4(1-U)^4.
\]
This gives the first inequality in \textup{(N42.4)}.  Also
\((1-X)^4(1-U)^4\leq1-X\) and, symmetrically, it is at most \(1-U\).
Hence \(W\geq\max\{X,U\}\).
\end{proof}

\subsection{A symmetric exact transverse derivative bound}

The height excesses \(e_j=t_j-\alpha_3\) obey
\[
 0\leq e_j< H:=\frac{1533}{5000},
 \qquad
 |c_j|\leq2.
\tag{N42.5}
\]
The center bound follows from \(c_j^2\leq4A/3\); the excess bound uses
\(\alpha_3>19434/10000\).

\begin{lemma}[Uniform height-direction Lipschitz constant]
\label{lem:f2v42-transverse-bound}
On the complete rational box in \textup{(N42.5)}, for each
\(j=1,2,3\),
\[
 \left|\frac{\partial\mathcal S_{\rm gen}}{\partial t_j}\right|
 \leq
 C:=
 \frac{3324053950171123}{62500000000}.
\tag{N42.6}
\]
The same constant holds for all three labels; no permutation is inferred
from a nonsymmetric two-coordinate interval.
\end{lemma}

\begin{proof}
Construct the three-pair input using independent symmetric center
variables \(c_1,c_2,c_3\), apply the degree-six heat operator, generate
Newton sums through \(p_8\), and form the exact V16 adapted square.
Differentiate before imposing the centered relation.

For each derivative, group the complete coefficient of every outer
monomial in
\((c_1,c_2,c_3,e_1,e_2,e_3)\).  Evaluate its coefficient in
\((\alpha_3,s)\) on the V13 rational rectangle, then bound its outer
monomial using \(|c_j|\leq2\) and \(0\leq e_j\leq H\).
There are exactly \(39\) regrouped monomials in every derivative, and
exact rational summation gives
\[
 C_1=C_2=C_3=C.
\tag{N42.7}
\]
Finally, substituting \(c_3=-c_1-c_2\) into the symmetric polynomial gives
coefficient-for-coefficient the \(261\)-term centered polynomial already
verified in V35.  Thus the symmetric calculation bounds the same adapted
square, merely on a larger enclosing center box.
\end{proof}

\begin{remark}[Exact symmetric verifier]
\label{rem:f2v42-verifier}
The source
\begin{center}
\path{scripts/verify_grh_v42_equal_height_wedge.py}
\end{center}
constructs the symmetric input, heat sextic, Newton sums, adapted square,
and all three height derivatives over \(\mathbb Q\).  It asserts
\[
\begin{array}{c|ccc}
\text{object}&P&\mathcal Q&\mathcal S_{\rm gen}\\ \hline
\text{terms}&112&185&305
\end{array},
\qquad
\#\text{ derivative groups}=39
\tag{N42.8}
\]
for each factor label.  It also asserts the centered polynomial identity,
the three equal constants in \textup{(N42.7)}, and the wedge coefficient
below.  It has \(209\) lines, \(6046\) bytes, and SHA--256
\[
 \texttt{BDFD9461D0653349A977ADCC18A73DB4}
 \allowbreak\texttt{F207BCB397B9EF6910BEAF35C0E1513A}.
\tag{N42.9}
\]
The chained V35 and V42 stages terminate with \texttt{VERIFIED}; every
sign and identity is exact.
\end{remark}

\subsection{The explicit wedge and an annular imbalance floor}

Put
\[
 \kappa_{\rm ht}:=\frac{57}{C}
 =\frac{3562500000000}{3324053950171123}.
\tag{N42.10}
\]

\begin{theorem}[Unequal-height wedge exclusion]
\label{thm:f2v42-wedge}
Under \textup{(N42.2)}, if
\[
 E<\kappa_{\rm ht}W(X,U),
\tag{N42.11}
\]
then the heat sextic
\(\mathcal Q_{\bm c,\bm t}\) is not real-rooted.
\end{theorem}

\begin{proof}
Join \(\bm t^0=(\bar t,\bar t,\bar t)\) to \(\bm t\) by
\(\bm t(\tau)=\bm t^0+\tau\bm\eta\).  Both endpoints have every
coordinate in \([\alpha_3,9/4]\), so the full segment does as well.
Lemma~\ref{lem:f2v42-transverse-bound} and the fundamental theorem of
calculus give
\[
\begin{aligned}
\left|\mathcal S_{\rm gen}(\bm c,\bm t)
-\mathcal S_{\rm gen}(\bm c,\bm t^0)\right|
&\leq
\int_0^1\sum_{j=1}^3
|\eta_j|
\left|\frac{\partial\mathcal S_{\rm gen}}
{\partial t_j}(\bm c,\bm t(\tau))\right|\dd\tau\\
&\leq CE.
\end{aligned}
\tag{N42.12}
\]
By Lemma~\ref{lem:f2v42-base-margin} and
\textup{(N42.11)},
\(\mathcal S_{\rm gen}(\bm c,\bm t)<-57W+CE<0\).
For six real roots the adapted square is a sum of six real squares and is
nonnegative.  This contradiction proves the claim.
\end{proof}

\begin{corollary}[Necessary height-imbalance floor]
\label{cor:f2v42-imbalance-floor}
Every candidate satisfying \textup{(N42.2)} must obey
\[
 \boxed{
 E\geq\kappa_{\rm ht}W(X,U)
 \geq\kappa_{\rm ht}\max\!\left\{
 \frac A3,\frac{\bar t-\alpha_3}{9/4-\alpha_3}\right\}.}
\tag{N42.13}
\]
If it is nonexceptional, then the V36 annular condition further implies
the absolute rational floor
\[
 E>
 \frac{296875000000000}{21716044456467946559}.
\tag{N42.14}
\]
\end{corollary}

\begin{proof}
The first assertion is the contrapositive of
Theorem~\ref{thm:f2v42-wedge} and the second inequality is
Lemma~\ref{lem:f2v42-base-margin}.  For the last assertion, write
\(h=9/4-\alpha_3<H\).  Since
\[
 r=A+3(\bar t-\alpha_3)=3X+3hU
 \leq3(1+H)\max\{X,U\},
\]
V36 gives
\(\max\{X,U\}>1/(60(1+H))=250/19599\).
Multiplication by \(\kappa_{\rm ht}\) and exact reduction gives
\textup{(N42.14)}.
\end{proof}

\begin{firewall}[A transverse wedge is not the full near-threshold cover]
\label{fire:f2v42}
V42 thickens the complete equal-height theorem over \(A\leq3\), and proves
that any surviving nonexceptional tuple there has a definite unequal-height
component.  It does not exclude tuples above the imbalance floor, the
region \(A>3\), collision charts, the infinite-product passage, or
arithmetic forcing.  GRH remains open.
\end{firewall}

\begin{assessment}[V42 advance and V43 direction]
\label{ass:f2v42}
This is the first global-in-height-spread bridge around the V15
equal-height face: it applies over the whole interval
\(\alpha_3\leq t_j\leq9/4\), not merely inside the V36 radial collar.
The symmetric compiler also removes the artificial large third-factor
constant produced by expanding \(c_3=-c_1-c_2\) before interval
evaluation.

V43 should use \textup{(N42.13)} as a discard rule in an ordered-height
simplex.  The efficient variables are the mean excess and the two signed
height gaps; boxes below the wedge are deleted before testing
\(\mathcal C,\mathcal E,\mathcal K,\mathcal M_3\).  The first finite
target should remain \(A\leq3\), because that is now bounded away from
both the equality face and the radial equality point by exact rational
conditions.
\end{assessment}

\section*{Version V42 append-only record}
\addcontentsline{toc}{section}{Version V42 append-only record}

V42 was appended byte-for-byte before the terminal marker of validated
V41, whose installed SHA--256 was
\[
 \texttt{6BE24040CD9E1FAB1E6A8CFD86D46CD2}
 \allowbreak\texttt{CDBB66A1F4535B316E30A6676D0CDB08}.
\tag{N42.15}
\]
It proves an explicit unequal-height wedge around the complete V15 face,
derives a nonzero imbalance floor in the residual annulus, and records a
symmetric exact verifier.  After validation V42 replaces V41 as the sole
active numeric GRH note; the frozen \texttt{V100\_f2} archive is
unchanged.

% =============================================================================
% END APPENDED V42 MATERIAL
% =============================================================================



% =============================================================================
% BEGIN APPENDED V43 MATERIAL
% =============================================================================

\section{V43: complete closure of the near-threshold small-gap cube}
\label{sec:f2v43-gap-cube}

V42 made a transverse neighborhood of the equal-height face effective.
The exact polynomial exposed by that calculation is substantially more
rigid than the Lipschitz estimate: after ordering the centers, a single
five-variable Bernstein expansion has no positive or ambiguous
coefficient on the entire adjacent-gap cube of side three.  This closes
all near-threshold tuples with \(A\leq3\), and also closes part of the
region beyond that variance.

\subsection{An ordered polynomial cube}

Simultaneously permute the three center-height pairs so that
\[
 c_1\leq c_2\leq c_3,\qquad
 g_1=c_2-c_1,\qquad g_2=c_3-c_2.
\tag{N43.1}
\]
After the usual common translation, \(c_1+c_2+c_3=0\), and therefore
\[
\begin{aligned}
c_1&=-\frac{2g_1+g_2}{3},&
c_2&=\frac{g_1-g_2}{3},&
c_3&=\frac{g_1+2g_2}{3},\\
A&=\frac{g_1^2+g_1g_2+g_2^2}{3}.
\end{aligned}
\tag{N43.2}
\]
Put
\[
 H=\frac{1533}{5000},\qquad
 g_1=3X,\quad g_2=3Y,\quad
 t_j=\alpha_3+HE_j.
\tag{N43.3}
\]
The unit five-cube \(0\leq X,Y,E_j\leq1\) corresponds to
\[
0\leq g_1,g_2\leq3,\qquad
\alpha_3\leq t_j\leq\alpha_3+H.
\tag{N43.4}
\]
It contains the requested physical slab \(t_j\leq9/4\), because
\(\alpha_3>19434/10000\) implies \(9/4-\alpha_3<H\).

Let
\[
 \mathcal S_{\Box}(X,Y,E_1,E_2,E_3)
 :=
 \mathcal S_{\rm gen}\bigl(
 -(2X+Y),X-Y,X+2Y;
 \alpha_3+HE_1,\alpha_3+HE_2,\alpha_3+HE_3
 \bigr).
\tag{N43.5}
\]
This is the V16 adapted square, not a new test.

\subsection{Exact algebraic Bernstein ledger}

The two threshold relations used to compile its coefficients are
\[
\begin{aligned}
8\alpha_3^3-42\alpha_3^2+90\alpha_3-75&=0,\\
s^2&=10-4\alpha_3,\qquad
s=\frac{4\alpha_3^2-20\alpha_3+25}{3\alpha_3-5},
\qquad 3\alpha_3-5>0.
\end{aligned}
\tag{N43.6}
\]
First reduce every \(s\)-power with the quadratic relation and every
\(\alpha_3\)-power with the cubic.  Then use the last identity in
\textup{(N43.6)}.  Each resulting coefficient is a rational polynomial
in \(\alpha_3\), divided by the positive linear denominator
\(3\alpha_3-5\), so its sign is enclosed by rational arithmetic on
\[
 \frac{19434}{10000}<\alpha_3<\frac{19435}{10000}.
\tag{N43.7}
\]

For a polynomial of multidegree
\(\bm d=(d_1,\ldots,d_5)\), the tensor power-to-Bernstein conversion is
\[
 b_{\bm i}
 =
 \sum_{\bm a\leq\bm i}
 \left(\prod_{\nu=1}^5
 \frac{\binom{i_\nu}{a_\nu}}{\binom{d_\nu}{a_\nu}}\right)
 r_{\bm a},
\qquad
 \mathcal S_{\Box}(\bm Z)=
 \sum_{\bm i\leq\bm d}b_{\bm i}
 \prod_{\nu=1}^5B_{i_\nu}^{d_\nu}(Z_\nu).
\tag{N43.8}
\]

\begin{lemma}[Five-cube Bernstein sign ledger]
\label{lem:f2v43-bernstein}
The polynomial \(\mathcal S_{\Box}\) has multidegree
\((8,8,4,4,4)\).  Its exact tensor Bernstein ledger has
\[
\begin{array}{c|r}
\text{coefficient class}&\text{count}\\ \hline
\text{strictly negative}&10122\\
\text{identically zero}&3\\
\text{positive or sign-ambiguous}&0.
\end{array}
\tag{N43.9}
\]
The zero-index set, in the order \((X,Y,E_1,E_2,E_3)\), is exactly
\[
 \mathcal Z=
 \{(0,0,0,0,0),(1,0,0,0,0),(0,1,0,0,0)\}.
\tag{N43.10}
\]
Among the rational upper endpoints for the negative coefficients, the
largest is
\[
 -\frac{164531079}{13288000}<0,
\tag{N43.11}
\]
attained at index \((1,1,0,0,0)\).
\end{lemma}

\begin{proof}
Substitution of \textup{(N43.3)} into the previously verified
\(261\)-term centered adapted square, followed by the reductions in
\textup{(N43.6)}, gives \(281\) nonzero power/algebraic terms.
Apply \textup{(N43.8)}.  The number of tensor coefficients is
\[
 (8+1)^2(4+1)^3=10125.
\]
For each coefficient, collect its full algebraic expression before
evaluating the isolating interval \textup{(N43.7)}.  Exact comparison
gives the counts, zero set, and largest negative upper endpoint in
\textup{(N43.9)}--\textup{(N43.11)}.
\end{proof}

\begin{remark}[Exact chained verifier]
\label{rem:f2v43-verifier}
The source
\begin{center}
\path{scripts/verify_grh_v43_near_threshold_gap_cube.py}
\end{center}
constructs \(\mathcal S_{\rm gen}\) from the heat sextic and Newton sums,
performs the gap substitution and algebraic reductions, builds all
\(10125\) tensor coefficients, and asserts every count, zero index, and
the exact margin in \textup{(N43.11)}.  It has \(175\) lines,
\(6076\) bytes, and SHA--256
\[
 \texttt{E5A7D6F11A8EF9F0CB9FBD74ADDDB9B6}
 \allowbreak\texttt{31D8F3679E24D97A9D014E42ADFAA215}.
\tag{N43.12}
\]
The chained V35 and V43 stages terminate with \texttt{VERIFIED}.
\end{remark}

\subsection{Strictness and the closed near-threshold region}

\begin{lemma}[Only the cube origin can have an allowed sign]
\label{lem:f2v43-strict}
On the unit five-cube,
\[
 \mathcal S_{\Box}(X,Y,E_1,E_2,E_3)<0
\tag{N43.13}
\]
unless \(X=Y=E_1=E_2=E_3=0\).  At that origin,
\(\mathcal S_{\Box}=0\).
\end{lemma}

\begin{proof}
Every tensor Bernstein basis function is nonnegative and the full family
forms a partition of unity.  Lemma~\ref{lem:f2v43-bernstein} therefore
gives \(\mathcal S_{\Box}\leq0\).

If any \(E_j>0\), some positive basis weight has a positive index in that
coordinate and hence cannot belong to \(\mathcal Z\).  Its coefficient is
strictly negative.  If all \(E_j=0\) but \((X,Y)\neq(0,0)\), the
degree-eight Bernstein support contains an index outside
\(\{(0,0),(1,0),(0,1)\}\): at an interior coordinate all indices have
positive weight, while at an endpoint the surviving index is \(0\) or
\(8\).  Again a strictly negative coefficient has positive weight.
Thus equality is possible only at the cube origin.  There only the
\((0,0,0,0,0)\) basis function survives, and its coefficient is zero.
\end{proof}

\begin{theorem}[Complete near-threshold small-gap closure]
\label{thm:f2v43-gap-cube}
Suppose \(t_j\geq\alpha_3\), order the center-height pairs as in
\textup{(N43.1)}, and assume
\[
 g_1\leq3,\qquad g_2\leq3,\qquad
 \max_jt_j\leq\frac94.
\tag{N43.14}
\]
If \(\mathcal Q_{\bm c,\bm t}\) is real-rooted, then
\[
 c_1=c_2=c_3=0,\qquad
 t_1=t_2=t_3=\alpha_3.
\tag{N43.15}
\]
Conversely, this exceptional tuple is real-rooted.
\end{theorem}

\begin{proof}
Translate the centers to have sum zero; this does not change their gaps
or real-rootedness.  Conditions \textup{(N43.14)} place the tuple in the
cube \textup{(N43.3)}--\textup{(N43.4)}.  If it is not the origin,
Lemma~\ref{lem:f2v43-strict} gives
\(\mathcal S_{\rm gen}<0\).  But for six real roots the adapted square is
\(\sum_{k=1}^6h(q_k)^2\geq0\), a contradiction.  At the origin,
real-rootedness is the V13 threshold factorization.
\end{proof}

\begin{corollary}[The full \(A\leq3\) near-threshold box is closed]
\label{cor:f2v43-a3}
If
\[
 A\leq3,\qquad
 \alpha_3\leq t_j\leq\frac94,
\tag{N43.16}
\]
then real-rootedness forces the exceptional tuple
\textup{(N43.15)}.
\end{corollary}

\begin{proof}
Equation \textup{(N43.2)} gives
\(g_1^2+g_1g_2+g_2^2=3A\leq9\), so each nonnegative adjacent gap is at
most \(3\).  Apply Theorem~\ref{thm:f2v43-gap-cube}.
\end{proof}

\begin{corollary}[Sharpened residual near-threshold target]
\label{cor:f2v43-residual}
Every nonexceptional candidate with
\(\alpha_3\leq T:=\max_jt_j\leq9/4\) must satisfy
\[
 \boxed{
 \max\{g_1,g_2\}>3,\qquad
 3<A<225,\qquad
 |c_j|<18,\qquad
 r>\frac1{20}.}
\tag{N43.17}
\]
\end{corollary}

\begin{proof}
Theorem~\ref{thm:f2v43-gap-cube} gives the gap and lower variance
restrictions.  V38 supplies \(A<225\) and \(|c_j|<18\), while V36 gives
the radial inequality.
\end{proof}

\begin{firewall}[A complete finite region is not the full programme]
\label{fire:f2v43}
V43 completely removes the near-threshold small-gap cube, including the
whole \(A\leq3\) region.  It does not exclude the residual annulus
\textup{(N43.17)}, the higher-height layers already reduced but not
eliminated, collision charts, the infinite-product passage, or arithmetic
forcing.  GRH remains open.
\end{firewall}

\begin{assessment}[V43 advance and V44 direction]
\label{ass:f2v43}
The V42 Lipschitz wedge was diagnostic rather than the final geometry:
retaining the full signed polynomial before taking moduli turns it into a
complete five-dimensional cube certificate.  The correct next target is
no longer a neighborhood of equal heights.

V44 should route the residual set by the larger adjacent gap and begin
with
\[
 3\leq\max\{g_1,g_2\}\leq6,\qquad
 \alpha_3\leq t_j\leq\frac94.
\tag{N43.18}
\]
On each factor-labelled gap slab it should test the signed adapted square
first and \(\mathcal M_3\) second.  Any subdivision must discard
\textup{(N43.17)} violations before polynomial evaluation and retain an
exact list of uncovered chart keys.
\end{assessment}

\section*{Version V43 append-only record}
\addcontentsline{toc}{section}{Version V43 append-only record}

V43 was appended byte-for-byte before the terminal marker of validated
V42, whose installed SHA--256 was
\[
 \texttt{D1173ED85126BAA422F9892183E3AA25}
 \allowbreak\texttt{B830B70943ED5BBC146A9D6047159AC4}.
\tag{N43.19}
\]
It proves the complete near-threshold adjacent-gap cube, closes all
tuples with \(A\leq3\), and records a \(10125\)-coefficient exact
Bernstein verifier.  After validation V43 replaces V42 as the sole active
numeric GRH note; the frozen \texttt{V100\_f2} archive is unchanged.

% =============================================================================
% END APPENDED V43 MATERIAL
% =============================================================================



% =============================================================================
% BEGIN APPENDED V44 MATERIAL
% =============================================================================

\section{V44: complete closure of the near-threshold layer}
\label{sec:f2v44-complete-near}

V43 removed the square \(0\leq g_1,g_2\leq3\).  The intended next step
was a thin slab adjacent to that square.  Exact affine-box tests reveal
a stronger fact: two adapted-square boxes and three moment-minor boxes
cover every gap pair allowed by the global near-threshold variance
bound.  Thus the entire layer \(T\leq9/4\) can be closed without a
numerical subdivision or an uncovered chart list.

Retain the ordering and translation normalization of V43.  Thus

\[
 c_1=-\frac{2g_1+g_2}{3},\qquad
 c_2= \frac{g_1-g_2}{3},\qquad
 c_3= \frac{g_1+2g_2}{3},
 \qquad
 A=\frac{g_1^2+g_1g_2+g_2^2}{3}.
\tag{N44.1}
\]

Write \(\mathcal S_{\rm gen}\) for the centered adapted square of V43
and \(\mathcal M_3\) for the third principal Hermite moment minor of
V35.  If the heat sextic has six real roots, then

\[
 \mathcal S_{\rm gen}\geq0,
 \qquad
 \mathcal M_3\geq0.
\tag{N44.2}
\]

The first inequality is a sum of six squares; the second is the
nonnegativity of a principal Gram determinant of the real-root moment
matrix.  These are necessary conditions only, but opposite signs on a
cover suffice for exclusion.

\subsection{A finite exact cover of every admissible gap pair}

Put

\[
 H:=\frac{1533}{5000},\qquad
 g_1=a+(b-a)X,\quad g_2=c+(d-c)Y,\quad
 t_j=\alpha_3+H E_j,
\tag{N44.3}
\]

where \(0\leq X,Y,E_1,E_2,E_3\leq1\).  As in V43, the slightly enlarged
closed height cube is legitimate because
\(9/4-\alpha_3<H\).  The five declared boxes are

\[
\begin{array}{c|c|c|c}
\text{key}&[a,b]&[c,d]&\text{signed certificate}\\ \hline
\texttt{S\_left}&[0,4]&[0,3]&\mathcal S_{\rm gen}\\
\texttt{S\_right}&[0,3]&[0,4]&\mathcal S_{\rm gen}\\
\texttt{M3\_left\_strip}&[0,3]&[4,26]&\mathcal M_3\\
\texttt{M3\_right\_strip}&[4,26]&[0,3]&\mathcal M_3\\
\texttt{M3\_upper\_square}&[3,26]&[3,26]&\mathcal M_3.
\end{array}
\tag{N44.4}
\]

The overlaps in \textup{(N44.4)} are deliberate and occur only between already
certified regions; the five keys themselves are distinct.

\begin{lemma}[The declared boxes cover the candidate gap domain]
\label{lem:f2v44-cover}
Every nonexceptional near-threshold candidate lies in
\([0,26]^2\) in the \((g_1,g_2)\)-plane, and the five boxes in
\textup{(N44.4)} cover that square.
\end{lemma}

\begin{proof}
V38 gives \(A<225\) for a near-threshold candidate.  Since the gaps are
nonnegative, \textup{(N44.1)} implies

\[
 g_i^2\leq g_1^2+g_1g_2+g_2^2=3A<675<676=26^2.
\tag{N44.5}
\]

Hence \(0\leq g_i<26\).  For the cover, first suppose \(g_1\leq3\).
If \(g_2\leq4\), use \texttt{S\_right}; if \(g_2\geq4\), use
\texttt{M3\_left\_strip}.  The analogous dichotomy when \(g_2\leq3\)
uses \texttt{S\_left} or \texttt{M3\_right\_strip}.  In the remaining
case both gaps are at least \(3\), so
\texttt{M3\_upper\_square} applies.  Boundary overlaps cause no gap.
\end{proof}

\subsection{Exact Bernstein ledgers}

For a polynomial \(P\) of multidegree at most
\((d_1,\ldots,d_5)\), write its tensor Bernstein expansion on the unit
five-cube as

\[
 P=\sum_{0\leq k_i\leq d_i}
 B_{k_1,\ldots,k_5}
 \prod_{i=1}^5\binom{d_i}{k_i}z_i^{k_i}(1-z_i)^{d_i-k_i}.
\tag{N44.6}
\]

The coefficients are computed exactly in
\(\mathbb Q(\alpha_3)\), reduced using

\[
 8\alpha_3^3-42\alpha_3^2+90\alpha_3-75=0,
 \qquad
 \frac{19434}{10000}<\alpha_3<\frac{19435}{10000}.
\tag{N44.7}
\]

No floating-point sign decision enters the ledgers below.

\begin{lemma}[The two adapted-square ledgers]
\label{lem:f2v44-s-ledgers}
On each of \texttt{S\_left} and \texttt{S\_right}, the substituted
\(\mathcal S_{\rm gen}\) has multidegree \((8,8,4,4,4)\) and exactly
\(10125\) tensor Bernstein coefficients.  Of these, \(10122\) are
strictly negative and precisely three are zero, at

\[
 (0,0,0,0,0),\qquad (1,0,0,0,0),\qquad (0,1,0,0,0).
\tag{N44.8}
\]

Among all negative coefficients in either ledger, the largest certified
upper endpoint is

\[
 -\frac{54843693}{3322000}<0,
\tag{N44.9}
\]

attained at index \((1,1,0,0,0)\).  Consequently
\(\mathcal S_{\rm gen}<0\) throughout either box except at the common
origin \(g_1=g_2=0,\ t_1=t_2=t_3=\alpha_3\), where it is zero.
\end{lemma}

\begin{proof}
Exact substitution produces \(281\) nonzero power/algebraic terms in
each box.  Converting the collected coefficients by \textup{(N44.6)} gives
\((8+1)^2(4+1)^3=10125\) coefficients and the stated sign ledger.
Every Bernstein basis function is nonnegative and the family is a
partition of unity.  If some height coordinate is positive, a basis
function with a positive height index and a strictly negative
coefficient has positive weight.  If all height coordinates vanish but
\((X,Y)\neq(0,0)\), the degree-eight support contains an index outside
the three indices in \textup{(N44.8)}, again with positive weight.  Thus equality
is possible only at the common cube origin.
\end{proof}

\begin{lemma}[The three moment-minor ledgers]
\label{lem:f2v44-m3-ledgers}
On each of the three \(\mathcal M_3\) boxes in \textup{(N44.4)}, the
substituted polynomial has multidegree \((6,6,3,3,3)\), hence \(3136\)
tensor Bernstein coefficients, and every coefficient is strictly
negative.  The numbers of nonzero power/algebraic terms are respectively
\(414,414,416\).  The largest certified upper endpoints are

\[
 -\frac{1121432131719}{5190625}<0
 \quad\text{on both strips},
 \qquad
 -\frac{3600446793777}{5190625}<0
 \quad\text{on the upper square}.
\tag{N44.10}
\]

Therefore \(\mathcal M_3<0\) everywhere on each of those three closed
boxes.
\end{lemma}

\begin{proof}
Construct \(\mathcal M_3\) from the V35 Newton sums, apply \textup{(N44.3)}, and
collect in \(\mathbb Q(\alpha_3)\).  Formula \textup{(N44.6)} yields
\((6+1)^2(3+1)^3=3136\) coefficients.  Exact interval evaluation using
\textup{(N44.7)} proves the stated strict upper bounds.  A convex combination of
strictly negative coefficients is strictly negative.
\end{proof}

\begin{remark}[Reproducible chained verifier]
\label{rem:f2v44-verifier}
The source
\begin{center}
\path{scripts/verify_grh_v44_complete_near_threshold.py}
\end{center}
imports and runs the V43 construction, reconstructs \(\mathcal M_3\)
from the V35 Newton sums, performs every affine box substitution, and
asserts the complete key set, term counts, multidegrees, coefficient
counts, zero indices, and rational sign margins.  It has \(195\) lines,
\(6235\) bytes, and SHA--256

\[
 \texttt{47038280CC5A24E1B9684F7A8C0CAAAF}
 \allowbreak\texttt{0366A123F19FE9C56A361D2C12FB3CE2}.
\tag{N44.11}
\]

The chained V35, V43, and V44 stages terminate with
\texttt{VERIFIED}.
\end{remark}

\subsection{The closed low-height theorem}

\begin{theorem}[Complete near-threshold classification]
\label{thm:f2v44-near-threshold}
Let \(t_j\geq\alpha_3\), and set \(T=\max_jt_j\).  If

\[
 T\leq\frac94,
\tag{N44.12}
\]

then the heat sextic \(\mathcal Q_{\bm c,\bm t}\) is real-rooted if and
only if, after a common translation of the centers,

\[
 c_1=c_2=c_3=0,\qquad t_1=t_2=t_3=\alpha_3.
\tag{N44.13}
\]
\end{theorem}

\begin{proof}
Order the center-height pairs by center and translate so that
\(c_1+c_2+c_3=0\).  Suppose first that the sextic is real-rooted and the
tuple is not \textup{(N44.13)}.  The global reductions V36 and V38 place every
such near-threshold candidate under the bound \(A<225\); hence
Lemma~\ref{lem:f2v44-cover} places its gap pair in one of the five boxes.
If it lies in an adapted-square box,
Lemma~\ref{lem:f2v44-s-ledgers} gives
\(\mathcal S_{\rm gen}<0\), contrary to \textup{(N44.2)}.  If it is routed to a
moment-minor box, Lemma~\ref{lem:f2v44-m3-ledgers} gives
\(\mathcal M_3<0\), again contrary to \textup{(N44.2)}.  Thus only \textup{(N44.13)}
remains.  Conversely, its real-rootedness is exactly the V13 threshold
factorization.  A common translation preserves real-rootedness.
\end{proof}

\begin{corollary}[The first active height layer starts above \(9/4\)]
\label{cor:f2v44-low-chart}
Every nonexceptional real-rooted candidate in the finite three-pair
problem satisfies

\[
 \boxed{T>\frac94.}
\tag{N44.14}
\]

Thus the near-threshold component of the candidate atlas is empty; the
exceptional threshold tuple is isolated.
\end{corollary}

\begin{firewall}[What V44 does and does not prove]
\label{fire:f2v44}
Theorem~\ref{thm:f2v44-near-threshold} is a complete theorem for the
finite three-pair heat-sextic model in the layer \(T\leq9/4\).  It does
not eliminate the higher-height layers, collision charts, or establish
the infinite-product and arithmetic implications needed for the
Riemann hypothesis.  In particular, it is not a proof of GRH; GRH
remains open.
\end{firewall}

\begin{assessment}[V44 advance and V45 direction]
\label{ass:f2v44}
The apparent large-gap annulus left by V43 disappears because the third
moment minor has a much stronger sign there than the adapted square.
This validates the upstream strategy of changing certificates rather
than repeatedly subdividing one polynomial.

The next honest layer is

\[
 \frac94<T\leq\frac{21}{8}.
\tag{N44.15}
\]

V45 should use the variance bound already proved for this height slab
to obtain a finite gap square, parameterize the actual height interval
directly, and test \(\mathcal S_{\rm gen}\) and \(\mathcal M_3\) before
introducing any new invariant.  Its chart ledger should again record a
fixed set of factor-labelled keys and retain exact rational margins.
If both existing certificates change sign on a surviving box, the
correct upstream move is to derive a better principal-minor or
subdiscriminant consequence, not to conceal the obstruction with
unbounded subdivision.
\end{assessment}

\section*{Version V44 append-only record}
\addcontentsline{toc}{section}{Version V44 append-only record}

V44 was appended byte-for-byte before the terminal marker of validated
V43, whose installed SHA--256 was

\[
 \texttt{9FBFCDC00FC1A3C09383EC25700752777}
 \allowbreak\texttt{783825848E448368132382A3F3D47B8}.
\tag{N44.16}
\]

It proves that five exact affine boxes cover the entire variance-allowed
near-threshold gap square and uses \(29658\) certified tensor Bernstein
coefficients across those boxes to close every nonexceptional tuple with
\(T\leq9/4\).  After validation V44 replaces V43 as the sole active
numeric GRH note; the frozen \texttt{V100\_f2} archive is unchanged.

% =============================================================================
% END APPENDED V44 MATERIAL
% =============================================================================


% =============================================================================
% BEGIN APPENDED V45 MATERIAL
% =============================================================================

\section{V45: three-chart closure through height four}
\label{sec:f2v45-through-four}

V44 closed the near-threshold layer \(T\leq9/4\).  V39--V40 had already
reduced the adjacent range \(9/4\leq T\leq4\) to bounded center regions,
but left those regions unresolved.  The decisive observation is that the
lower bound on the factor attaining \(T\), when retained rather than
discarded, makes the third Hermite moment minor strictly negative on one
large gap square.  Three factor labels therefore replace the twelve
height half-slab charts previously used only for outer cutoffs.

As before, order the center-height pairs by center, translate so that
\(c_1+c_2+c_3=0\), and put

\[
 g_1=c_2-c_1,\qquad g_2=c_3-c_2,\qquad
 A=\frac{g_1^2+g_1g_2+g_2^2}{3},\qquad
 \epsilon_j=t_j-\alpha_3.
\tag{N45.1}
\]

The algebraic threshold satisfies the isolating inequalities

\[
 \frac{19434}{10000}<\alpha_3<\frac{19435}{10000}.
\tag{N45.2}
\]

Define the rational constants

\[
 L:=\frac{613}{2000},\qquad
 U:=\frac{2057}{1000}.
\tag{N45.3}
\]

\subsection{The three enlarged maximal-factor charts}

For \(q\in\{1,2,3\}\), let \(\mathcal C_q\) be the image of the unit
five-cube under

\[
 g_1=49X,\qquad g_2=49Y,\qquad
 \epsilon_q=L+(U-L)E_q,\qquad
 \epsilon_k=U E_k\quad(k\neq q).
\tag{N45.4}
\]

\begin{lemma}[Exact coverage of every post-threshold candidate through four]
\label{lem:f2v45-cover}
Suppose \(t_j\geq\alpha_3\), \(9/4<T:=\max_jt_j\leq4\), and the heat
sextic is a candidate surviving the preceding reductions.  If \(q\)
attains \(T\), then the translated and ordered tuple belongs to
\(\mathcal C_q\).
\end{lemma}

\begin{proof}
The V39--V40 height ladder gives \(A<784\) throughout this range.  Hence

\[
 g_i^2\leq g_1^2+g_1g_2+g_2^2=3A
 <2352<2401=49^2,
\tag{N45.5}
\]

so the two gap coordinates in \textup{(N45.4)} lie in the unit interval.
For a maximal factor, \textup{(N45.2)} and \(T>9/4\) give

\[
 \epsilon_q=T-\alpha_3
 >\frac94-\frac{19435}{10000}
 =\frac{613}{2000}=L.
\tag{N45.6}
\]

Every factor satisfies \(0\leq\epsilon_k\), while \(t_k\leq T\leq4\)
and the other side of \textup{(N45.2)} give

\[
 \epsilon_k\leq4-\alpha_3
 <4-\frac{19434}{10000}
 =\frac{10283}{5000}
 <\frac{2057}{1000}=U.
\tag{N45.7}
\]

Thus suitable \(E_1,E_2,E_3\in[0,1]\) exist.  Equal maxima may be routed
to any attaining factor, so the three declared keys form a complete
cover without a tie-breaking assumption.
\end{proof}

\subsection{A strictly negative moment-minor ledger}

Recall the V35 principal moment minor

\[
 \mathcal M_3
 =6p_2p_6-6p_4^2-p_2p_3^2,
\tag{N45.8}
\]

where the \(p_k\) are the Newton power sums of the six roots.  For six
real roots, \(\mathcal M_3\geq0\), since this is a principal Gram
determinant of the real-root moment matrix.

\begin{lemma}[Three exact Bernstein ledgers]
\label{lem:f2v45-ledgers}
For each chart key

\[
 \texttt{max\_factor\_1},\qquad
 \texttt{max\_factor\_2},\qquad
 \texttt{max\_factor\_3},
\tag{N45.9}
\]

substitution of \textup{(N45.4)} into \(\mathcal M_3\), followed by
reduction in \(\mathbb Q(\alpha_3)\), produces \(262\) nonzero
power/algebraic terms and multidegree
\((6,6,3,3,3)\).  All \(3136\) coefficients in its tensor Bernstein
expansion on the unit five-cube are strictly negative.  In each ledger,
the largest certified upper endpoint is

\[
 -\frac{17613320542915912503}{10381250000000000}<0,
\tag{N45.10}
\]

at Bernstein index \((1,0,0,0,0)\).  Consequently
\(\mathcal M_3<0\) everywhere on every \(\mathcal C_q\).
\end{lemma}

\begin{proof}
Expand each affine power in \textup{(N45.4)}, collect the entire
coefficient in the algebraic basis before interval evaluation, and use
the threshold equation

\[
 8\alpha_3^3-42\alpha_3^2+90\alpha_3-75=0
\tag{N45.11}
\]

to reduce it to degree at most two in \(\alpha_3\).  The tensor degree
gives
\[
 (6+1)^2(3+1)^3=3136
\]
coefficients per chart.  Exact rational evaluation on the isolating
interval \textup{(N45.2)} gives the stated sign count and the endpoint
\textup{(N45.10)}.  The Bernstein basis is nonnegative and forms a
partition of unity, so a convex combination of these strictly negative
coefficients is strictly negative.
\end{proof}

\begin{remark}[Reproducible chained verifier]
\label{rem:f2v45-verifier}
The source
\begin{center}
\path{scripts/verify_grh_v45_first_adjacent_layer.py}
\end{center}
imports the V44 algebraic construction, reconstructs the three affine
height substitutions, asserts the rational inclusion inequalities
\textup{(N45.5)}--\textup{(N45.7)}, builds all \(9408\) Bernstein
coefficients, and asserts every chart key, term count, multidegree, sign
count, index, and exact margin.  It has \(149\) lines, \(5110\) bytes,
and SHA--256

\[
 \texttt{9BD2FE103E8409027207899CB7E079FF}
 \allowbreak\texttt{941C88781BE87220E156F6D94951BCFD}.
\tag{N45.12}
\]

The chained V35, V43, V44, and V45 stages terminate with
\texttt{VERIFIED}; no floating-point sign is used.
\end{remark}

\subsection{Classification through height four}

\begin{theorem}[The whole post-threshold layer through four is empty]
\label{thm:f2v45-post-threshold-empty}
If \(t_j\geq\alpha_3\) and

\[
 \frac94<T:=\max_jt_j\leq4,
\tag{N45.13}
\]

then \(\mathcal Q_{\bm c,\bm t}\) is not real-rooted.
\end{theorem}

\begin{proof}
Assume real-rootedness and choose a factor \(q\) attaining \(T\).
Lemma~\ref{lem:f2v45-cover} places the tuple in \(\mathcal C_q\).
Lemma~\ref{lem:f2v45-ledgers} then gives \(\mathcal M_3<0\), contradicting
the necessary real-root moment inequality \(\mathcal M_3\geq0\).
\end{proof}

\begin{corollary}[Complete finite classification through height four]
\label{cor:f2v45-through-four}
Suppose \(t_j\geq\alpha_3\) and \(T\leq4\).  The heat sextic is
real-rooted if and only if there is a common center \(a\in\mathbb R\)
such that

\[
 c_1=c_2=c_3=a,\qquad
 t_1=t_2=t_3=\alpha_3.
\tag{N45.14}
\]
\end{corollary}

\begin{proof}
V44 proves the classification for \(T\leq9/4\).
Theorem~\ref{thm:f2v45-post-threshold-empty} excludes the complementary
range through four.  Conversely, the V13 threshold factorization is
real-rooted and a common translation preserves its roots.
\end{proof}

\begin{firewall}[The finite height-four theorem is not GRH]
\label{fire:f2v45}
Corollary~\ref{cor:f2v45-through-four} is a complete theorem only for the
finite three-pair heat-sextic model under \(T\leq4\).  It does not
eliminate \(T>4\), settle collision and degeneration limits outside the
certified charts, or establish the infinite-product and arithmetic
implications required for the Riemann hypothesis.  GRH remains open.
\end{firewall}

\begin{assessment}[V45 advance and V46 direction]
\label{ass:f2v45}
The maximal-factor lower bound is the essential datum: without it the
small-gap \(\mathcal M_3\) ledger has positive coefficients, while with
it all three full gap-square ledgers are strictly negative.  The same
single hull remains valid through \(T=4\), collapsing the earlier
factor-labelled half-slabs into three exclusion charts.

A direct hull from \(9/4\) to \(14\) was also tested and correctly
rejected: \(136\) positive Bernstein coefficients appear.  This failed
probe is not part of the proof and has not been retained as a verifier.
The first honest unresolved slab is now

\[
 4<T\leq5,\qquad A<35^2,\qquad g_i<61,
\tag{N45.15}
\]

using the first V41 cutoff and \(3\cdot35^2<61^2\).  V46 should test the
three maximal-factor charts on that slab.  If the whole slab does not
certify, it should bisect only the upper height bound and record the
first exact sign-changing chart; the large rejected hull shows that
unbounded vertical enlargement is no longer justified.
\end{assessment}

\section*{Version V45 append-only record}
\addcontentsline{toc}{section}{Version V45 append-only record}

V45 was appended byte-for-byte before the terminal marker of validated
V44, whose installed SHA--256 was

\[
 \texttt{B343647B69848B6578C738016911D6BA}
 \allowbreak\texttt{CAA0C322E57F58A1C300A148362A78B6}.
\tag{N45.16}
\]

It proves three exact full-gap-square moment-minor ledgers and combines
them with the V39--V40 variance ladder to classify all finite
three-pair tuples through \(T=4\).  After validation V45 replaces V44 as
the sole active numeric GRH note; the frozen \texttt{V100\_f2} archive
is unchanged.

% =============================================================================
% END APPENDED V45 MATERIAL
% =============================================================================


% =============================================================================
% BEGIN APPENDED V46 MATERIAL
% =============================================================================

\section{V46: ordered-height routing and closure through eight}
\label{sec:f2v46-through-eight}

V45 classified the finite three-pair model through \(T=4\).  The V41
variance ladder makes the next four units compact.  Retaining the exact
maximal-height order, rather than merely bounding the other heights by
the slab ceiling, now closes all candidates through \(T=8\).

Continue to write \(\epsilon_j=t_j-\alpha_3\), order the pairs by center,
and use adjacent gaps \(g_1,g_2\).  If \(q\) attains
\(T=\max_jt_j\), introduce the ordered-height substitution

\[
 h=\epsilon_q=L+(U-L)H,\qquad
 \epsilon_k=hE_k\ (k\ne q),\qquad 0\leq H,E_k\leq1.
\tag{N46.1}
\]

It enforces \(0\leq\epsilon_k\leq\epsilon_q\) identically, including
ties.  The rational slab data are

\[
\begin{array}{c|c|c|c|c}
T\text{ range}&A\text{ cutoff}&R&L&U\\ \hline
(9/4,6]&41^2&72&613/2000&4057/1000\\
(6,7]&47^2&82&8113/2000&5057/1000\\
(7,8]&53^2&92&10113/2000&6057/1000.
\end{array}
\tag{N46.2}
\]

The first row uses the affine V45 chart
\(\epsilon_q=L+(U-L)E_q,\ \epsilon_k=UE_k\); the ordered substitution
\textup{(N46.1)} is used in the last two rows.

\begin{lemma}[Exact coverage]
\label{lem:f2v46-cover}
Every candidate with \(9/4<T\leq8\) belongs to one of the fifteen charts
declared below.
\end{lemma}

\begin{proof}
For each row, the V41 ladder and
\[
 g_i^2\leq3A<R^2
\tag{N46.3}
\]
place \((g_1,g_2)\) in \([0,R]^2\); explicitly
\(3\cdot41^2<72^2\), \(3\cdot47^2<82^2\), and
\(3\cdot53^2<92^2\).  The isolating interval
\(19434/10000<\alpha_3<19435/10000\) proves each rational height
inclusion in \textup{(N46.2)}.

For the first row choose any maximal factor, giving three full gap
squares.  In either later row, factors \(q=1,3\) use a full gap square.
For \(q=2\), route \([0,R]^2\) through
\[
 [0,4]^2,\quad [0,4]\times[4,R],\quad
 [4,R]\times[0,4],\quad[4,R]^2.
\tag{N46.4}
\]
These four boxes cover without a boundary gap.  Thus there are
\(3+6+6=15\) declared charts.
\end{proof}

\begin{lemma}[Fifteen strictly negative moment-minor ledgers]
\label{lem:f2v46-ledgers}
On every chart of Lemma~\ref{lem:f2v46-cover}, the tensor Bernstein
expansion of \(\mathcal M_3\) has only strictly negative coefficients.
The first three charts have multidegree \((6,6,3,3,3)\) and \(3136\)
coefficients each.  Every ordered-height chart has degree four in its
maximal-height coordinate, degree three in the other height coordinates,
and \(3920\) coefficients.  Hence the complete ledger contains
\[
 3(3136)+12(3920)=56448
\tag{N46.5}
\]
strictly negative exact coefficients.
\end{lemma}

\begin{proof}
Substitute the affine centers from the adjacent gaps and either the
first-row height chart or \textup{(N46.1)} into
\[
 \mathcal M_3=6p_2p_6-6p_4^2-p_2p_3^2.
\tag{N46.6}
\]
Collect in \(\mathbb Q(\alpha_3)\), reduce by
\(8\alpha_3^3-42\alpha_3^2+90\alpha_3-75=0\), and evaluate only after
forming each complete Bernstein coefficient.  Exact rational interval
comparison gives the stated counts.  The closest negative upper
endpoints in the three height families are respectively
\[
 -\frac{17613320542915912503}{10381250000000000},\quad
 -\frac{13567514744708981157}{1297656250000000},\quad
 -\frac{11165814318345134157}{1297656250000000}.
\tag{N46.7}
\]
All are negative.  Since the Bernstein basis is a nonnegative partition
of unity, \(\mathcal M_3<0\) on every chart.
\end{proof}

\begin{remark}[Exact chained verifier]
\label{rem:f2v46-verifier}
The source
\begin{center}
\path{scripts/verify_grh_v46_ordered_height_through_eight.py}
\end{center}
constructs both height substitutions, declares the fifteen unique keys,
and asserts all coverage inequalities, term counts, multidegrees,
coefficient counts, extremal indices, margins, and empty ambiguous
lists.  It has \(196\) lines, \(7117\) bytes, and SHA--256
\[
 \texttt{A2568AD3869918AC9A153664810A0399}
 \allowbreak\texttt{77EAC126FCEA7073BC39B741BC3642EF}.
\tag{N46.8}
\]
The full chain terminates with \texttt{VERIFIED}.
\end{remark}

\begin{theorem}[Complete finite classification through height eight]
\label{thm:f2v46-through-eight}
Suppose \(t_j\geq\alpha_3\) and \(T=\max_jt_j\leq8\).  Then
\(\mathcal Q_{\bm c,\bm t}\) is real-rooted if and only if, for some
\(a\in\mathbb R\),
\[
 c_1=c_2=c_3=a,\qquad t_1=t_2=t_3=\alpha_3.
\tag{N46.9}
\]
\end{theorem}

\begin{proof}
V45 proves the assertion through \(T=4\).  If \(4<T\leq8\), it is
covered by Lemma~\ref{lem:f2v46-cover}; Lemma~\ref{lem:f2v46-ledgers}
then gives \(\mathcal M_3<0\), contradicting the necessary principal
real-root moment inequality \(\mathcal M_3\geq0\).  The converse is the
V13 threshold factorization and translation invariance.
\end{proof}

\begin{firewall}[The height-eight theorem is finite]
\label{fire:f2v46}
Theorem~\ref{thm:f2v46-through-eight} does not settle \(T>8\), the
infinite-product passage, or arithmetic forcing.  It is not a proof of
GRH; GRH remains open.
\end{firewall}

\begin{assessment}[V46 boundary and V47 direction]
\label{ass:f2v46}
On \(8<T\leq9\), the four middle-max gap boxes remain strictly negative,
but each outer-factor full square has eight positive Bernstein
coefficients.  V47 should subdivide those two outer squares by the same
cut \(g_i=4\), preserving the exact ordered-height substitution.  This
is the first unresolved chart set; no earlier height layer should be
recomputed.
\end{assessment}

\section*{Version V46 append-only record}
\addcontentsline{toc}{section}{Version V46 append-only record}

V46 was appended byte-for-byte before the terminal marker of validated
V45, whose installed SHA--256 was
\[
 \texttt{2864438C02605B19BED973C5B0B166A}
 \allowbreak\texttt{6EA8FE20DD7B25964EF0F3225E2D6D375}.
\tag{N46.10}
\]
It proves the ordered-height routing and exact classification through
\(T=8\).  After validation V46 replaces V45 as the sole active numeric
GRH note; the frozen \texttt{V100\_f2} archive is unchanged.

% =============================================================================
% END APPENDED V46 MATERIAL
% =============================================================================


\clearpage
\section{V47: twelve-box closure of the height-eight-to-nine slab}
\label{sec:f2v47-eight-nine}
\begingroup
\setcounter{equation}{0}
\renewcommand{\theequation}{N47.\arabic{equation}}
\renewcommand{\theHequation}{f2v47.\arabic{equation}}

This iteration continues the finite real-rooted obstruction programme developed in
V41--V46.  The point is deliberately local: no statement below assumes GRH, and
no finite computation is promoted into an asymptotic theorem.  We prove that the
already reduced obstruction has no candidate with ordered height
\[
   8<T\leq 9.
\]
Together with V46, this extends the exact finite classification from height eight
through height nine.

Retain the notation of V46.  Thus \(q\in\{1,2,3\}\) is an index for which
\(\epsilon_q=\max_k\epsilon_k\), the other two normalized height variables lie
in \([0,1]\), and \(g_1,g_2\) are the two nonnegative gap variables.  On the
ordered-height chart put
\begin{equation}
  h=\epsilon_q=\frac{12113}{2000}
       +\left(\frac{7057}{1000}-\frac{12113}{2000}\right)H,
  \qquad
  \epsilon_k=hE_k\quad(k\ne q),
  \label{eq:f2v47-height-map}
\end{equation}
where \(H,E_k\in[0,1]\).

\begin{lemma}[Exact coverage of the height and gap coordinates]
\label{lem:f2v47-cover}
Every candidate in the slab \(8<T\leq9\) belongs to one of the twelve boxes
obtained by choosing \(q\in\{1,2,3\}\) and choosing one of
\begin{equation}
 \begin{split}
 B_{00}&=[0,4]\times[0,4],\\
 B_{01}&=[0,4]\times[4,101],\\
 B_{10}&=[4,101]\times[0,4],\\
 B_{11}&=[4,101]\times[4,101]
 \end{split}
 \label{eq:f2v47-gap-boxes}
\end{equation}
for \((g_1,g_2)\).
\end{lemma}

\begin{proof}
The rational endpoint estimates already proved in the reduction give
\begin{equation}
 \epsilon_q>8-\frac{19435}{10000}=\frac{12113}{2000},
 \qquad
 \epsilon_k\leq9-\alpha
   <9-\frac{19434}{10000}=\frac{35283}{5000}
   <\frac{7057}{1000}.
 \label{eq:f2v47-height-cover}
\end{equation}
Hence the substitution \eqref{eq:f2v47-height-map} covers every ordered-height
chart in the slab.  Moreover, the V41 area estimate gives \(A<58^2\).  The gap
bound \(g_i^2\leq3A\) therefore yields
\begin{equation}
  g_i^2<3\cdot58^2=10092<101^2,
  \qquad 0\leq g_i<101.
 \label{eq:f2v47-gap-cover}
\end{equation}
The four rectangles in \eqref{eq:f2v47-gap-boxes} cover \([0,101]^2\).
There are three possible maximal-factor indices, so these give the asserted
twelve-box cover.
\end{proof}

For each box, substitute its affine unit-cube parametrization into the same
cleared obstruction polynomial \(M_3\) used in V45--V46 and convert its exact
power-basis expansion to the tensor-product Bernstein basis.  All arithmetic is
over \(\mathbb Q\); consequently the following is a rational certificate, not a
floating-point sample.

\begin{lemma}[Twelve exact Bernstein ledgers]
\label{lem:f2v47-ledgers}
On each of the twelve boxes in Lemma~\ref{lem:f2v47-cover}, every Bernstein
coefficient of the pulled-back polynomial \(M_3\) is strictly negative.
Each box has \(3920\) coefficients.  In the order
\(B_{00},B_{01},B_{10},B_{11}\), the power-basis term counts are
\begin{equation}
 \begin{array}{c|rrrr}
 q& B_{00}&B_{01}&B_{10}&B_{11}\\ \hline
 1&470&730&729&733\\
 2&472&733&733&736\\
 3&470&729&730&733
 \end{array}
 \label{eq:f2v47-term-ledger}
\end{equation}
Thus the certificate contains \(12\cdot3920=47040\) strictly negative rational
coefficients.  The largest coefficient among the entire collection is
\begin{equation}
 -\frac{11023728848422605973}{1868625000000000}<0,
 \label{eq:f2v47-global-margin}
\end{equation}
occurring in the middle-factor small box at multi-index
\((1,1,0,4,0)\).
\end{lemma}

\begin{proof}
The exact verifier expands each of the twelve affine pullbacks independently,
checks the declared multidegree, constructs every tensor-product Bernstein
coefficient as a reduced rational number, and asserts that its numerator is
negative.  The degrees are \((6,6,3,3,3)\), with the height coordinate belonging
to the maximal factor raised from degree \(3\) to degree \(4\); hence
\[
  (6+1)(6+1)(4+1)(3+1)(3+1)=3920
\]
coefficients occur on every chart.  The verifier also asserts the twelve chart
keys, their term counts, their coefficient counts, and their individual maximal
coefficients.  Taking the maximum of those twelve exact maxima gives
\eqref{eq:f2v47-global-margin}.  Since a polynomial on a box lies in the convex
hull of its Bernstein coefficients, strict negativity of every coefficient
implies \(M_3<0\) throughout every box.
\end{proof}

\begin{remark}[Reproducible certificate]
\label{rem:f2v47-verifier}
The permanent verifier is
\path{scripts/verify_grh_v47_ordered_height_eight_nine.py}.
It imports and reruns the V46 certificate before checking the twelve new ledgers,
so the new claim is chained to the prior finite classification.  Its source hash
is
\begin{center}
\texttt{66FB65E2DFC25078BEE4D352B5D2B61BDE9F6857BEA147C2C71E55E7B99A03D2}.
\end{center}
The file has \(93\) lines and \(2774\) bytes.  The proof itself depends only on
the exact identities and sign assertions just described.
\end{remark}

\begin{theorem}[No reduced candidate between heights eight and nine]
\label{thm:f2v47-eight-nine}
Under the reductions and notation proved through V46, the real-rooted obstruction
admits no candidate with \(8<T\leq9\).
\end{theorem}

\begin{proof}
Choose a maximal factor \(q\).  Lemma~\ref{lem:f2v47-cover} places the candidate
in one of the four gap boxes for that ordered-height chart.  By
Lemma~\ref{lem:f2v47-ledgers}, \(M_3<0\) on that box.  This contradicts the
necessary nonnegative obstruction inequality for a real-rooted candidate.
\end{proof}

\begin{corollary}[Finite classification through height nine]
\label{cor:f2v47-through-nine}
Combining Theorem~\ref{thm:f2v47-eight-nine} with V44--V46, every reduced
real-rooted candidate with \(T\leq9\) is the translated symmetric threshold tuple
already classified in V44.
\end{corollary}

\subsection*{Firewall and next target}
\label{fire:f2v47}
This is a finite-height theorem inside the conditional obstruction programme.  It
does not control candidates of arbitrary height and therefore does not prove GRH.
The advance is that the ordered-height method survives a much wider gap range:
splitting both gap coordinates at \(4\) removes the positive outer-box Bernstein
coefficients that blocked the unsplit calculation.

For V48 the immediate target is \(9<T\leq10\).  The V41 estimate gives
\(A<64^2\), and hence
\begin{equation}
  g_i^2\leq3A<3\cdot64^2=12288<111^2.
  \label{eq:f2v47-next-gap}
\end{equation}
The first nonduplicative test is therefore the analogous twelve-box cover with
outer gap endpoint \(111\).  Failure there should be treated upstream: refine the
height or gap coordinates, or strengthen the analytic area estimate, rather than
merely accumulating ad hoc boxes.

\subsection*{Assumption ledger}
\label{ass:f2v47}
The new step uses only: (i) the reductions, obstruction polynomial, rational
height estimates, and area/gap inequalities already proved through V46;
(ii) the elementary box cover in Lemma~\ref{lem:f2v47-cover}; and (iii) exact
rational polynomial identities checked by the permanent verifier.  No numerical
rounding, unproved zero-free estimate, density hypothesis, or GRH input is used.


\endgroup


\clearpage
\section{V48: refined closure of the height-nine-to-ten slab}
\label{sec:f2v48-nine-ten}
\begingroup
\setcounter{equation}{0}
\renewcommand{\theequation}{N48.\arabic{equation}}
\renewcommand{\theHequation}{f2v48.\arabic{equation}}

This iteration converts the exploratory endpoint left by V47 into a chained exact
certificate.  Retain the obstruction polynomial \(M_3\) and choose
\(q\in\{1,2,3\}\) with \(\epsilon_q=\max_k\epsilon_k\).  On each ordered chart set
\begin{equation}
 h=\frac{14113}{2000}
   +\left(\frac{8057}{1000}-\frac{14113}{2000}\right)H,
 \qquad \epsilon_k=hE_k\quad(k\ne q),
 \label{eq:f2v48-height}
\end{equation}
with \(H,E_k\in[0,1]\).

\begin{lemma}[Exact cover]
\label{lem:f2v48-cover}
Every reduced candidate with \(9<T\le10\) is covered by the following fifteen
ordered-height boxes.  For \(q=1,3\), use
\[
 [0,4]^2,\quad [0,4]\times[4,111],\quad
 [4,111]\times[0,4],\quad[4,111]^2.
\]
For \(q=2\), replace \([0,4]^2\) by its four subdivisions at \(g_1=g_2=2\),
and retain the other three boxes.
\end{lemma}

\begin{proof}
The endpoint estimates from the established reduction give
\begin{equation}
 \epsilon_q>9-\frac{19435}{10000}=\frac{14113}{2000},\qquad
 \epsilon_k\le10-\alpha
 <10-\frac{19434}{10000}=\frac{40283}{5000}
 <\frac{8057}{1000}.
 \label{eq:f2v48-height-cover}
\end{equation}
The V41 area estimate gives \(A<64^2\), so
\begin{equation}
 g_i^2\le3A<3\cdot64^2=12288<111^2.
 \label{eq:f2v48-gap-cover}
\end{equation}
Thus \(0\le g_i<111\).  The stated rectangles cover the gap square, and
\eqref{eq:f2v48-height} covers all three ordered-height charts.
\end{proof}

\begin{lemma}[Fifteen exact Bernstein ledgers]
\label{lem:f2v48-ledgers}
Every tensor-product Bernstein coefficient of the pullback of \(M_3\) is strictly
negative on every box of Lemma~\ref{lem:f2v48-cover}.  Each chart has \(3920\)
coefficients, so the certificate contains
\begin{equation}
 15\cdot3920=58800
 \label{eq:f2v48-total}
\end{equation}
strictly negative rationals.  The largest coefficient in the collection is
\begin{equation}
 -\frac{2031387253721119991}{622875000000000}<0,
 \label{eq:f2v48-margin}
\end{equation}
at multi-index \((1,1,0,4,0)\) in the middle-factor box
\([0,2]\times[0,2]\).
\end{lemma}

\begin{proof}
The three outer-factor charts have power-term counts \(470,730,729,733\).
For the middle factor, the four refined small boxes have counts
\(472,733,733,736\), and the remaining three boxes have counts \(733,733,736\).
All pulled-back multidegrees are \((6,6,3,3,3)\), with the maximal factor's
height degree raised to \(4\).  Hence every ledger has
\((7)(7)(5)(4)(4)=3920\) entries.  The permanent verifier constructs every entry
over \(\mathbb Q\), asserts the stated counts and degrees, asserts an empty
ambiguity list, and compares each chart maximum with its stored exact rational
value.  All \(58800\) numerators are negative; the Bernstein convex-hull property
therefore gives \(M_3<0\) throughout the cover.
\end{proof}

\begin{remark}[Permanent certificate]
\label{rem:f2v48-verifier}
The chained verifier is
\path{scripts/verify_grh_v48_ordered_height_nine_ten.py}.  It first reruns V47.
It has \(63\) lines, \(2808\) bytes, and SHA--256 hash
\begin{center}
\texttt{26DC6A4758A107EC286B06ABB3B692CD9CF616E20B40C9AD640136996DD8EAE9}.
\end{center}
\end{remark}

\begin{theorem}[No reduced candidate between heights nine and ten]
\label{thm:f2v48-nine-ten}
Under the reductions proved through V47, the real-rooted obstruction admits no
candidate with \(9<T\le10\).
\end{theorem}

\begin{proof}
Lemma~\ref{lem:f2v48-cover} places any candidate in a certified box, while
Lemma~\ref{lem:f2v48-ledgers} gives \(M_3<0\) there.  This contradicts the
necessary nonnegative obstruction inequality.
\end{proof}

\begin{corollary}[Finite classification through height ten]
\label{cor:f2v48-through-ten}
Every reduced real-rooted candidate with \(T\le10\) is the translated symmetric
threshold tuple classified in V44.
\end{corollary}

\subsection*{Firewall and next target}
\label{fire:f2v48}
This remains a finite-height result and does not prove GRH.  The important
structural observation is that the only coarse failure was the middle-factor
small-gap chart: splitting both gaps at \(2\) eliminated all ten positive
coarse-basis coefficients.  V49 should test the next unit-height slab with this
refinement built in, and should move upstream to the height or area estimate if
the number of exceptional boxes begins to grow.

\subsection*{Assumption ledger}
\label{ass:f2v48}
The new argument uses only the prior reductions and exact rational endpoint
bounds, the elementary finite cover above, and exact rational polynomial
identities checked by the chained verifier.  It uses no sampling, floating-point
rounding, zero-density conjecture, or GRH input.
\endgroup



\clearpage
\section{V49: adaptive simplex closure through height eleven}
\label{sec:f2v49-ten-eleven}
\begingroup
\setcounter{equation}{0}
\renewcommand{\theequation}{N49.\arabic{equation}}
\renewcommand{\theHequation}{f2v49.\arabic{equation}}

V48 closed the slab through height ten, but its middle-factor small-gap box
ceased to certify after direct enlargement.  This iteration follows the
coefficient geometry upstream: the gap cover is retained, while only the
normalized height simplex of that one box is subdivided.

Choose \(q\) with \(\epsilon_q=\max_k\epsilon_k\) and put
\begin{equation}
 h=\frac{16113}{2000}
   +\left(\frac{9057}{1000}-\frac{16113}{2000}\right)H,
 \qquad \epsilon_k=hE_k\quad(k\ne q).
 \label{eq:f2v49-height}
\end{equation}
The V41 ladder and the isolating interval for \(\alpha_3\) give
\begin{equation}
 A<70^2,\qquad
 g_i^2\le3A<14700<122^2,\qquad 0\le g_i<122,
 \label{eq:f2v49-gap}
\end{equation}
and
\[
 \epsilon_q>\frac{16113}{2000},\qquad
 \epsilon_k<11-\frac{19434}{10000}
 =\frac{45283}{5000}<\frac{9057}{1000}.
\]

\begin{lemma}[A disjoint forty-three-box cover]
\label{lem:f2v49-cover}
Every reduced candidate with \(10<T\le11\) lies in one of forty-three
certified boxes.
\end{lemma}

\begin{proof}
For \(q=1,3\), split the gap square at \(4\), giving four boxes for each
factor.  For \(q=2\), the three boxes outside \([0,4]^2\), together with
the three parts of \([0,4]^2\) outside \([0,2]^2\), give six boxes.
These are the fourteen top-level boxes.

It remains to partition the \(q=2\), \([0,2]^2\) box.  In its normalized
coordinates \((E_1,H,E_3)\), bisect all three coordinates.  Retain seven
children and subdivide the child with bits \(010\).  Again retain seven
children and subdivide its \(010\) child, namely
\[
 E_1,E_3\in[0,\tfrac14],\qquad H\in[\tfrac34,1].
\]
Split this last box dyadically; retain seven children and split its \(010\)
child once more.  All eight grandchildren are retained.  Thus the small
box has \(7+7+7+8=29\) disjoint leaves, and the whole cover has
\(14+29=43\) boxes.  Boundary overlaps leave no gap.
\end{proof}

\begin{lemma}[Exact Bernstein certificate]
\label{lem:f2v49-ledger}
On every box of Lemma~\ref{lem:f2v49-cover}, all \(3920\)
tensor-product Bernstein coefficients of \(M_3\) are strictly negative.
Consequently the certificate contains
\begin{equation}
 43\cdot3920=168560
 \label{eq:f2v49-total}
\end{equation}
strictly negative exact rational coefficients.
\end{lemma}

\begin{proof}
The pulled-back multidegree is \((6,6,3,4,3)\) on the middle-factor
charts, with the analogous permutation on the outer charts.  Hence every
ledger has \(7\cdot7\cdot5\cdot4\cdot4=3920\) entries.  The permanent
verifier checks fourteen top-level ledgers, fourteen retained siblings
from the first two simplex subdivisions, and fifteen leaves of the final
adaptive branch.  It asserts the coefficient count, exact algebraic
interval sign, unique ownership count, and empty unresolved list.  The
Bernstein convex-hull property gives \(M_3<0\) on every leaf.
\end{proof}

\begin{remark}[Permanent chained verifier]
\label{rem:f2v49-verifier}
The verifier is
\path{scripts/verify_grh_v49_ordered_height_ten_eleven.py}; it uses the
exact reboxing engine in
\path{scripts/verify_grh_v49_ordered_height_ten_eleven_core.py} and first
reruns V48.  The main verifier has \(109\) lines, \(3219\) bytes, and
SHA--256
\begin{center}
\texttt{44EAB577C1029221A63797078A9F2028A76185024DD71732942A6F08AB30A39A}.
\end{center}
\end{remark}

\begin{theorem}[No reduced candidate between heights ten and eleven]
\label{thm:f2v49-ten-eleven}
Under the reductions proved through V48, the real-rooted obstruction has
no candidate with \(10<T\le11\).
\end{theorem}

\begin{proof}
Lemma~\ref{lem:f2v49-cover} routes every candidate to a certified box.
Lemma~\ref{lem:f2v49-ledger} gives \(M_3<0\) there, contradicting the
necessary nonnegative moment-minor inequality.
\end{proof}

\begin{corollary}[Finite classification through height eleven]
\label{cor:f2v49-through-eleven}
Every reduced real-rooted candidate with \(T\le11\) is the translated
symmetric threshold tuple classified in V44.
\end{corollary}

\subsection*{Firewall and V50 direction}
\label{fire:f2v49}
This is still a finite-height theorem and does not prove GRH.  The
adaptive tree remains small, but its repeated \(010\) branch shows that
blind uniform subdivision would be wasteful.  Before V50, the mandatory
ten-version sweep must inspect all parallel notes.  The next slab has the
V41 cutoff \(A<76^2\), gap endpoint \(132\), and ordered-height endpoints
\(18113/2000\) and \(10057/1000\); any new certificate should preserve
the explicit ownership tree or replace it with a stronger upstream
inequality.

\subsection*{Assumption ledger}
\label{ass:f2v49}
This step uses only the prior reductions, the exact V41 variance ladder,
the isolating interval for \(\alpha_3\), finite dyadic box identities, and
exact rational interval arithmetic.  It uses no sampling inference,
floating-point sign decision, density conjecture, or GRH input.
\endgroup



\clearpage
\section{V50: companion sweep and the height-twelve frontier}
\label{sec:f2v50-sweep-frontier}
\begingroup
\setcounter{equation}{0}
\renewcommand{\theequation}{N50.\arabic{equation}}
\renewcommand{\theHequation}{f2v50.\arabic{equation}}

This mandatory ten-version checkpoint was performed before extending the
V49 adaptive tree.  The active companion sources inspected read-only were
\begin{equation}
\begin{array}{c|r|l}
\text{source}&\text{lines}&\text{SHA--256}\\ \hline
\text{Yang--Mills V25}&9496&
\texttt{5D6759B911E076304A17D09955E6777C3}\\
&&\texttt{ACD6C99D3B47CD42BBA75EB0EA96BE6}\\
\text{SM--Einstein V17}&3715&
\texttt{8C51C43B04AE6FDCE07589855437FCE5}\\
&&\texttt{508F2A012F8C4C22DC4DF042E6BF56DC}
\end{array}
\label{eq:f2v50-sources}
\end{equation}
SM advanced during the sweep; the final installed V17 was therefore
checked after the earlier V12 finite-net theorem.

\begin{proposition}[What transfers from the companion sweep]
\label{prop:f2v50-transfer}
Two proof disciplines transfer to the GRH obstruction programme.
\begin{enumerate}
\item A reduction must be proved from the whole coefficient family, not
from ranks or signs seen at selected points.
\item A finite net becomes a uniform certificate only after recording a
proved Lipschitz constant, a covering radius, and exact pointwise margins.
\end{enumerate}
No numerical estimate or field-specific theorem in the companion notes
transfers to the present polynomial.
\end{proposition}

\begin{proof}
The first statement is the coefficient-hull discipline of Yang--Mills
V20, retained by its V25 audit.  The second is the gap-from-a-net theorem
of SM V12, retained in V17.  Their hypotheses concern, respectively,
matrix Laurent families and regulated operator gaps; neither identifies
the present \(M_3\), its five-dimensional chart metric, or a derivative
bound for it.  Thus only the certificate architecture transfers.
\end{proof}

\subsection{Exact first gate for the next slab}

For \(11<T\le12\), the V41 ladder and the isolating interval for
\(\alpha_3\) give
\begin{equation}
 h\in\left[\frac{18113}{2000},\frac{10057}{1000}\right],
 \qquad A<76^2,\qquad
 g_i^2\le3A<17328<132^2.
\label{eq:f2v50-domain}
\end{equation}
The ordered-height substitution and the fifteen top-level boxes of V48
were tested on this complete domain.

\begin{lemma}[Exact height-twelve frontier ledger]
\label{lem:f2v50-frontier}
Ten of the fifteen boxes have all \(3920\) Bernstein coefficients
strictly negative.  Exactly five small-gap boxes fail:
\begin{equation}
\begin{array}{c|c|r}
q&\text{gap box}&\text{positive coefficients}\\ \hline
1&[0,4]^2&32\\
2&[0,2]^2&368\\
2&[0,2]\times[2,4]&29\\
2&[2,4]\times[0,2]&29\\
3&[0,4]^2&32
\end{array}
\label{eq:f2v50-failures}
\end{equation}
There are no mixed-sign algebraic intervals.  Across all fifteen ledgers
there are \(58310\) negative and \(490\) positive exact coefficients.
\end{lemma}

\begin{proof}
The permanent verifier
\path{scripts/verify_grh_v50_companion_sweep_frontier.py}
first reruns the complete V49 certificate.  It then reconstructs all
fifteen rational affine pullbacks, asserts \(3920\) coefficients per
chart, compares each sign dictionary with the declared ledger, checks
that precisely ten chart keys have empty ambiguity lists, and checks that
the five displayed positive counts sum to \(490\).  It terminates with
\texttt{VERIFIED}.  Its SHA--256 is
\begin{center}
\texttt{A91E2D669AF90624218BE47049ADC135EA4389F8968EC08E3D2CEF14FCE0DE08}.
\end{center}
\end{proof}

\begin{corollary}[What remains proved]
\label{cor:f2v50-status}
The finite classification through \(T\le11\) proved in V49 remains
unchanged.  Lemma~\ref{lem:f2v50-frontier} is a complete routing statement
for the first \(11<T\le12\) gate, but it does not exclude its five failed
boxes.
\end{corollary}

\subsection*{Firewall and V51 direction}
\label{fire:f2v50}
Positive Bernstein coefficients are failure of this enclosure, not proof
that \(M_3\) is positive.  V50 therefore makes no height-twelve theorem
and no GRH claim.  In accordance with the sweep, V51 should compute the
complete transformed coefficient family on the five failed boxes and
derive exact coordinate-derivative bounds.  A dyadic net may be used only
if its exact negative margin exceeds the proved Lipschitz radius loss.
If that gate is unprofitable or encounters an actual positive point, the
programme should test the earlier signed obstructions
\(\mathcal S_{\rm gen},\mathcal C,\mathcal E,\mathcal K\) on precisely
those five boxes rather than enlarge the \(M_3\) tree blindly.

\subsection*{Assumption ledger}
\label{ass:f2v50}
The new mathematical assertions use only the exact V41 cutoff, the
isolating interval for \(\alpha_3\), the inherited ordered-height
substitution, and exact rational/algebraic interval arithmetic.  The
companion results determine proof discipline only; none of their
field-specific hypotheses or constants is imported.
\endgroup




\clearpage
\section{V51: the sharp total-height budget and the collapse of the height frontier}
\label{sec:f2v51-total-height}
\begingroup
\setcounter{equation}{0}
\renewcommand{\theequation}{N51.\arabic{equation}}
\renewcommand{\theHequation}{f2v51.\arabic{equation}}

V50 left five failed \(\mathcal M_3\) boxes in the slab \(11<T\le12\),
while the candidate set proved in V32--V33 still extended to
\(T<19321/100\) through the collision chart.  Continuing unit slabs would
have been the circular bookkeeping that V41 already warned against.  Three
exploratory exact ledgers written after V50 and rerun here,
\begin{center}
\path{scripts/verify_grh_v51_m3_spine_sign_core.py},\\
\path{scripts/verify_grh_v51_m3_spine_barrier.py},\\
\path{scripts/verify_grh_v51_spine_obstruction_switch.py},
\end{center}
show that at the spine point \(g_1=g_2=0\), \(\epsilon_1=\epsilon_3=0\),
\(h=38227/4000\), that is at heights
\((\alpha_3,\,t_2,\,\alpha_3)\) with \(230003/20000<t_2<46001/4000\),
the three quantities \(\mathcal M_3\), \(\mathcal K\), and
\(\mathcal S_{\rm gen}\) are strictly positive while \(\mathcal C\) and
\(\mathcal E\) are strictly negative.  Thus the \(\mathcal M_3\) march
would have needed an obstruction switch a quarter of a unit above the
V49 frontier.  Instead of switching certificates, this iteration goes
upstream to the zero flow of the heat sextic.  The result is a sharp,
center-free budget that removes every candidate above
\(T=15-2\alpha_3\), and an effective center bound that shrinks the
remaining sliver to a small box, which a chained certificate then closes.
The finite three-pair problem is thereby completed.

\subsection{The zero flow of the heat sextic}

Throughout, \(P\) is a monic real sextic,
\begin{equation}
 \mathcal Q_\tau:=e^{-\tau D_z^2/2}P,\qquad 0\le\tau\le1,
 \label{eq:f2v51-flow-def}
\end{equation}
and \(z_1(\tau),\ldots,z_6(\tau)\) are the zeros of \(\mathcal Q_\tau\)
counted with multiplicity, with \(x_k=\operatorname{Re}z_k\) and
\(y_k=\operatorname{Im}z_k\).  Since \(\partial_\tau\mathcal Q_\tau
=-\tfrac12\mathcal Q_\tau''\), this is the backward heat flow already used
in the archive; \(\mathcal Q_1=\mathcal Q_{\bm c,\bm t}\) when
\(P=\prod_j((z-c_j)^2+t_j)\).

\begin{lemma}[Finite collision set and the flow equation]
\label{lem:f2v51-flow}
\begin{enumerate}[label=\textup{(\roman*)}]
\item The set \(\Sigma\subset[0,1]\) of times at which
\(\mathcal Q_\tau\) has a multiple zero is finite.
\item On every component of \([0,1]\setminus\Sigma\) the zeros admit an
analytic labelling, and
\begin{equation}
 \dot z_k=\sum_{l\ne k}\frac1{z_k-z_l}\qquad(1\le k\le6).
 \label{eq:f2v51-ode}
\end{equation}
\item For every continuous \(\varphi:\mathbb C\to\mathbb R\), the map
\(\tau\mapsto\sum_k\varphi(z_k(\tau))\) is continuous on \([0,1]\).
\end{enumerate}
\end{lemma}

\begin{proof}
(i) The discriminant of \(\mathcal Q_\tau\) is a polynomial in \(\tau\).
For \(\tau>0\) and \(0\le j\le6\),
\[
 e^{-\tau D_z^2/2}z^j
 =\sum_{k}\frac{(-\tau/2)^k}{k!}\frac{j!}{(j-2k)!}z^{j-2k}
 =\tau^{j/2}\mathrm{He}_j\!\left(z/\sqrt\tau\right),
\]
so if \(P=\sum_ja_jz^j\) with \(a_6=1\), then
\(\tau^{-3}\mathcal Q_\tau(\sqrt\tau\,w)
 =\mathrm{He}_6(w)+\sum_{j<6}a_j\tau^{(j-6)/2}\mathrm{He}_j(w)\to\mathrm{He}_6(w)\)
as \(\tau\to\infty\).  The probabilists' Hermite polynomial
\(\mathrm{He}_6\) has six simple real zeros, so \(\mathcal Q_\tau\) has
simple zeros for all large \(\tau\).  Hence the discriminant is not the
zero polynomial and has finitely many zeros in \([0,1]\).

(ii) At a simple zero \(z_k\) of \(\mathcal Q_\tau\) the implicit
function theorem gives an analytic zero branch with
\(\dot z_k=-\partial_\tau\mathcal Q_\tau(z_k)/\mathcal Q_\tau'(z_k)
 =\tfrac12\mathcal Q_\tau''(z_k)/\mathcal Q_\tau'(z_k)\).  If
\(\mathcal Q_\tau=\prod_l(z-z_l)\) and \(z_k\) is simple, then
\(\mathcal Q_\tau'(z_k)=\prod_{l\ne k}(z_k-z_l)\) and
\(\mathcal Q_\tau''(z_k)=2\sum_{m\ne k}\prod_{l\ne k,m}(z_k-z_l)\),
whence \(\mathcal Q_\tau''(z_k)/\mathcal Q_\tau'(z_k)
 =2\sum_{l\ne k}(z_k-z_l)^{-1}\).

(iii) The coefficients of \(\mathcal Q_\tau\) are polynomials in
\(\tau\), and the multiset of zeros of a monic polynomial of fixed degree
depends continuously on its coefficients (Rouch\'{e}'s theorem).  A
symmetric continuous function of the multiset is therefore continuous in
\(\tau\).
\end{proof}

\begin{proposition}[Two exact rates]
\label{prop:f2v51-rates}
Off \(\Sigma\), with \(u_{kl}:=(x_k-x_l)^2\) and \(w_{kl}:=(y_k-y_l)^2\),
\begin{equation}
 \frac{\dd}{\dd\tau}\sum_{k=1}^6x_k^2
 =\sum_{k\ne l}\frac{u_{kl}}{u_{kl}+w_{kl}},
 \qquad
 \frac{\dd}{\dd\tau}\sum_{k=1}^6y_k^2
 =-\sum_{k\ne l}\frac{w_{kl}}{u_{kl}+w_{kl}},
 \label{eq:f2v51-rates}
\end{equation}
where no denominator vanishes because the zeros are simple.
Consequently
\begin{equation}
 0\le\frac{\dd}{\dd\tau}\sum_kx_k^2\le30,
 \qquad
 -30\le\frac{\dd}{\dd\tau}\sum_ky_k^2\le0,
 \label{eq:f2v51-rate-bounds}
\end{equation}
both sums are Lipschitz on \([0,1]\), \(\sum_kx_k^2\) is nondecreasing,
and \(\sum_ky_k^2\) is nonincreasing there.
\end{proposition}

\begin{proof}
Since \((z_k-z_l)^{-1}=(\overline{z_k}-\overline{z_l})/|z_k-z_l|^2\),
equation \eqref{eq:f2v51-ode} gives
\[
 \dot x_k=\sum_{l\ne k}\frac{x_k-x_l}{|z_k-z_l|^2},
 \qquad
 \dot y_k=-\sum_{l\ne k}\frac{y_k-y_l}{|z_k-z_l|^2}.
\]
Therefore
\(\frac{\dd}{\dd\tau}\sum_kx_k^2
 =2\sum_{k\ne l}x_k(x_k-x_l)/|z_k-z_l|^2\); exchanging the names
\(k,l\) and averaging the two expressions gives
\(\sum_{k\ne l}(x_k-x_l)^2/|z_k-z_l|^2\), and
\(|z_k-z_l|^2=u_{kl}+w_{kl}\).  The same computation with \(y\) gives
the second identity.  Each of the \(30\) ordered pairs contributes a
number in \([0,1]\), which proves \eqref{eq:f2v51-rate-bounds}.  Both
sums are continuous on \([0,1]\) by
Lemma~\ref{lem:f2v51-flow}(iii) and differentiable with derivative in
\([-30,30]\) off the finite set \(\Sigma\); the mean value theorem on each
subinterval and continuity across \(\Sigma\) give the Lipschitz property,
so both are absolutely continuous and monotone as stated.
\end{proof}

\begin{remark}[Consistency with the Newton sum]
The two rates in \eqref{eq:f2v51-rates} add to \(30\).  This is the
derivative of \(p_2(\tau)=\sum_k(x_k^2-y_k^2)=p_2(0)+30\tau\), the
second-Newton-sum identity \textup{(N16.3)} that produced the budget
\textup{(N16.4)}.  That budget used only \(\sum_kx_k(1)^2\ge0\).  The new
information is the separate monotonicity of each half.
\end{remark}

\subsection{The sharp total-height budget}

\begin{theorem}[Total-height budget]
\label{thm:f2v51-budget}
Let \(P=\prod_{j=1}^3((z-c_j)^2+t_j)\) with real \(c_j\) and \(t_j>0\),
and suppose \(\mathcal Q_{\bm c,\bm t}=e^{-D_z^2/2}P\) is real-rooted.
Then
\begin{equation}
 \boxed{t_1+t_2+t_3\le15.}
 \label{eq:f2v51-budget}
\end{equation}
Equality holds if and only if \(\mathcal Q_{\bm c,\bm t}=(z-a)^6\) for
some real \(a\).  In that case
\begin{equation}
 P(z)=(z-a)^6+15(z-a)^4+45(z-a)^2+15,
 \label{eq:f2v51-hermite-input}
\end{equation}
and \(-t_1,-t_2,-t_3\) are the three real zeros of
\(X^3+15X^2+45X+15\), which lie in \((-20,-5)\), \((-5,-1)\), and
\((-1,0)\).
\end{theorem}

\begin{proof}
The zeros of \(P\) are \(c_j\pm i\sqrt{t_j}\), so
\(\sum_ky_k(0)^2=2(t_1+t_2+t_3)\), while real-rootedness gives
\(\sum_ky_k(1)^2=0\).  By Proposition~\ref{prop:f2v51-rates},
\begin{equation}
 2(t_1+t_2+t_3)
 =\int_0^1\sum_{k\ne l}\frac{w_{kl}}{u_{kl}+w_{kl}}\,\dd\tau
 \le30.
 \label{eq:f2v51-budget-integral}
\end{equation}
If equality holds, the integrand equals \(30\) for almost every
\(\tau\), so by continuity of the integrand on each component of
\([0,1]\setminus\Sigma\) every ordered pair has \(u_{kl}=0\) there: at
every \(\tau\notin\Sigma\) all six zeros lie on one vertical line
\(\operatorname{Re}z=\xi(\tau)\).  Choose \(\tau_\nu\uparrow1\) with
\(\tau_\nu\notin\Sigma\).  The zero multisets of
\(\mathcal Q_{\tau_\nu}\) converge to the real zero multiset of
\(\mathcal Q_1\).  If \(\mathcal Q_1\) had two distinct zeros
\(r<r'\), then for large \(\nu\) one zero of \(\mathcal Q_{\tau_\nu}\)
would lie within \((r'-r)/3\) of \(r\) and another within
\((r'-r)/3\) of \(r'\), so their real parts would differ.  Hence
\(\mathcal Q_1=(z-a)^6\).

Conversely, \(e^{-D_z^2/2}\) is invertible on polynomials with inverse
\(e^{D_z^2/2}\), and
\[
 e^{D_z^2/2}z^6
 =z^6+\frac{6\cdot5}{2}z^4+\frac{6\cdot5\cdot4\cdot3}{4\cdot2}z^2
  +\frac{6!}{8\cdot6}
 =z^6+15z^4+45z^2+15,
\]
which gives \eqref{eq:f2v51-hermite-input} after translation.  The cubic
\(f(X)=X^3+15X^2+45X+15\) has
\(f(-20)=-2885\), \(f(-5)=40\), \(f(-1)=-16\), \(f(0)=15\), so it has
three negative real zeros \(X_1<X_2<X_3\) in the displayed intervals.
Thus \(P=\prod_i((z-a)^2-X_i)=\prod_i((z-a)^2+t_i)\) with
\(t_i=-X_i>0\), and \(\sum_it_i=-(X_1+X_2+X_3)=15\) by Vieta.  Its heat
image \((z-a)^6\) is real-rooted.
\end{proof}

\begin{remark}[Scope of the budget]
\label{rem:f2v51-scope}
The bound \eqref{eq:f2v51-budget} is uniform in the centers and replaces
the second-moment budget \textup{(N16.4)}, \(\sum_jt_j\le15+2A\), by
its center-free form.  The same proof gives, for \(m\) conjugate pairs
and heat time \(K\), \(\sum_jt_j\le m(2m-1)K\).  On the equal-height
diagonal this gives only \(y^2\le(2m-1)K\), weaker than the V9
shallow-pair bound \(mK\) and than the sharp centered threshold
\(\alpha_mK\), so the diagonal gains nothing.  The gain is that the sum
charges every channel: the
largest height is now bounded by the smallest ones,
\(T\le15-(t_{(2)}+t_{(3)})\).  This is the monotone statistic
charging companion channels that the V9 assessment asked for.
\end{remark}

\begin{corollary}[Collapse of the height frontier]
\label{cor:f2v51-collapse}
Let \(t_j\ge\alpha_3\) for all \(j\) and let
\(\mathcal Q_{\bm c,\bm t}\) be real-rooted.  Then
\begin{equation}
 \boxed{T=\max_jt_j<15-2\alpha_3<\frac{27783}{2500}=11.1132,}
 \qquad
 t_k\le15-T-\alpha_3\quad(t_k\ne T).
 \label{eq:f2v51-collapse}
\end{equation}
In particular: the V32 collision chart \(T>441/25\) contains no
candidate; the ceilings \(121/4\), \(1089/4\), and \(19321/100\) of
V31--V32 are superseded; of the five failed V50 boxes only points with
\(T<27783/2500\) and \(\epsilon_k<283/2500\) for the two non-maximal
factors can be candidates; and the exploratory positive
\(\mathcal M_3\) spine point has \(\sum_jt_j>15.38\) and is excluded.
\end{corollary}

\begin{proof}
Theorem~\ref{thm:f2v51-budget} gives \(T+2\alpha_3\le\sum_jt_j\le15\).
Equality \(\sum_jt_j=15\) would force the Hermite tuple of
\eqref{eq:f2v51-hermite-input}, whose smallest height lies in
\((0,1)\), contradicting \(t_j\ge\alpha_3>19/10\); hence
\(\sum_jt_j<15\) and \(T<15-2\alpha_3\).  The rational bound uses
\(\alpha_3>19434/10000\).  The height bound for the other factors is the
same inequality read for \(t_k\); its excess form is
\(\epsilon_k=t_k-\alpha_3\le15-T-2\alpha_3<4-2\alpha_3<283/2500\) when
\(T>11\).  The remaining statements are direct comparisons with the V32,
V33, and V50 constants; at the spine point
\(t_2>230003/20000\), so \(\sum_jt_j>230003/20000+2\alpha_3>15.38\).
\end{proof}

\subsection{Strip preservation and an effective center bound}

\begin{lemma}[One-factor strip preservation]
\label{lem:f2v51-one-factor}
Let \(p\) be a polynomial of degree \(n\ge1\) whose zeros all lie in the
closed strip \(S_h=\{z:|\operatorname{Im}z|\le h\}\), \(h\ge0\), and let
\(a\in\mathbb R\).  Then every zero of \(p+ap'\) lies in \(S_h\).
\end{lemma}

\begin{proof}
For \(a=0\) there is nothing to prove.  Let \(a\ne0\) and suppose
\(p(z)+ap'(z)=0\) with \(z\notin S_h\).  Then \(p(z)\ne0\), so
\(p'(z)/p(z)=\sum_\nu(z-\zeta_\nu)^{-1}=-1/a\) is real, where
\(\zeta_\nu\) are the zeros of \(p\).  But
\(\operatorname{Im}(z-\zeta_\nu)^{-1}
 =-(\operatorname{Im}z-\operatorname{Im}\zeta_\nu)/|z-\zeta_\nu|^2\)
has the same strict sign for every \(\nu\), namely negative if
\(\operatorname{Im}z>h\) and positive if \(\operatorname{Im}z<-h\).
The sum therefore has nonzero imaginary part, a contradiction.
\end{proof}

\begin{lemma}[Heat flow preserves horizontal strips]
\label{lem:f2v51-strip}
If all zeros of the monic sextic \(P\) lie in \(S_h\), then all zeros of
\(e^{-\tau D_z^2/2}P\) lie in \(S_h\) for every \(\tau\ge0\).  In
particular, for \(P=\prod_j((z-c_j)^2+t_j)\) and \(T=\max_jt_j\),
\begin{equation}
 \max_k|y_k(\tau)|\le\sqrt T\qquad(0\le\tau\le1).
 \label{eq:f2v51-strip}
\end{equation}
\end{lemma}

\begin{proof}
Fix \(\tau>0\) and \(N\ge1\), and put \(b=\sqrt{\tau/(2N)}\).  Then
\(1-\tau D_z^2/(2N)=(1-bD_z)(1+bD_z)\), so by
Lemma~\ref{lem:f2v51-one-factor}, applied \(2N\) times, the polynomial
\(P_N:=(1-\tau D_z^2/(2N))^NP\) has all zeros in \(S_h\); it is monic of
degree six.  Expanding,
\(P_N=\sum_{k=0}^3\binom Nk(-\tau/2N)^kD_z^{2k}P\), and
\(\binom Nk N^{-k}\to1/k!\), so \(P_N\to e^{-\tau D_z^2/2}P\)
coefficientwise.  The zero multisets converge, and \(S_h\) is closed.
The zeros of the displayed \(P\) are \(c_j\pm i\sqrt{t_j}\), which lie in
\(S_{\sqrt T}\).
\end{proof}

This is P\'{o}lya's classical strip theorem for the operator
\(e^{-\tau D^2/2}\); the proof is included because the archive uses
\(\mathcal Q_{\bm c,\bm t}\) only at \(\tau=1\) and never recorded the
intermediate-time statement.

\begin{theorem}[Deficit bound]
\label{thm:f2v51-deficit}
Let \(P=\prod_{j=1}^3((z-c_j)^2+t_j)\) with \(c_1+c_2+c_3=0\),
\(A=\tfrac12\sum_jc_j^2\), \(T=\max_jt_j\), and suppose
\(\mathcal Q_{\bm c,\bm t}\) is real-rooted.  Then
\begin{equation}
 \boxed{30-2(t_1+t_2+t_3)\ \ge\ \frac{12A}{6A+T}.}
 \label{eq:f2v51-deficit}
\end{equation}
\end{theorem}

\begin{proof}
By Proposition~\ref{prop:f2v51-rates} and absolute continuity,
\begin{equation}
 \sum_kx_k(1)^2-\sum_kx_k(0)^2
 =\int_0^1\rho(\tau)\,\dd\tau,
 \qquad
 \rho:=\sum_{k\ne l}\frac{u_{kl}}{u_{kl}+w_{kl}},
 \label{eq:f2v51-deficit-integral}
\end{equation}
and adding the two identities of \eqref{eq:f2v51-rates} shows that the
left side equals \(30-\bigl(\sum_ky_k(0)^2-\sum_ky_k(1)^2\bigr)
 =30-2(t_1+t_2+t_3)\).  By \eqref{eq:f2v51-strip},
\(w_{kl}\le(|y_k|+|y_l|)^2\le4T\), so
\(\rho\ge\sum_{k\ne l}\phi(u_{kl})\) with \(\phi(u):=u/(u+4T)\).  The
function \(\phi\) is concave on \([0,\infty)\) with \(\phi(0)=0\), hence
\(\phi(u)\ge(u/U)\phi(U)\) for \(0\le u\le U\), \(U>0\).

Put \(U(\tau):=\sum_{k<l}u_{kl}=6\sum_kx_k^2-\bigl(\sum_kx_k\bigr)^2
 =6\sum_kx_k(\tau)^2\); here \(\sum_kz_k=2\sum_jc_j=0\) at \(\tau=0\), and
the coefficient of \(z^5\) is untouched by \(e^{-\tau D_z^2/2}\), so
\(\sum_kx_k(\tau)=0\) for all \(\tau\).  If \(U(\tau)=0\), then
\(\sum_kx_k(\tau)^2=0\), so by monotonicity \(\sum_kx_k(0)^2=4A=0\) and
\eqref{eq:f2v51-deficit} reads \(30-2\sum_jt_j\ge0\), which is
Theorem~\ref{thm:f2v51-budget}.  Otherwise, since every
\(u_{kl}\le U\),
\[
 \rho\ge\frac{2\phi(U)}{U}\sum_{k<l}u_{kl}=2\phi(U)
 =\frac{2U}{U+4T}\ge\frac{2\cdot24A}{24A+4T}=\frac{12A}{6A+T},
\]
using \(U(\tau)\ge U(0)=6\sum_kx_k(0)^2=24A\)
(Proposition~\ref{prop:f2v51-rates}) and the monotonicity of
\(u\mapsto2u/(u+4T)\).  Integrating over \([0,1]\) gives
\eqref{eq:f2v51-deficit}.
\end{proof}

\begin{corollary}[The sliver box]
\label{cor:f2v51-sliver}
Let \(t_j\ge\alpha_3\), let \(\mathcal Q_{\bm c,\bm t}\) be real-rooted,
and suppose \(T>11\).  With the centered ordered-gap coordinates of
\textup{(N33.18)}--\textup{(N33.23)},
\begin{equation}
 \boxed{
 A<\frac{873621}{3695000}<\frac{2365}{10000},\qquad
 g_1,g_2<\frac{17}{20},\qquad
 11<T<\frac{27783}{2500},\qquad
 0\le\epsilon_k<\frac{283}{2500}\ (t_k\ne T).}
 \label{eq:f2v51-sliver}
\end{equation}
\end{corollary}

\begin{proof}
Write \(d:=30-2\sum_jt_j\).  Since the two other heights are at least
\(\alpha_3>19434/10000\) and \(T>11\),
\[
 d\le30-2T-4\alpha_3<8-\frac{4\cdot19434}{10000}=\frac{283}{1250}=:d_{\max},
\]
and by Corollary~\ref{cor:f2v51-collapse},
\(T<15-2\alpha_3<27783/2500=:T_{\max}\).
Theorem~\ref{thm:f2v51-deficit} gives
\(12A/(6A+T_{\max})\le12A/(6A+T)\le d<d_{\max}\), hence
\(A(12-6d_{\max})<d_{\max}T_{\max}\), that is
\[
 A<\frac{d_{\max}T_{\max}}{12-6d_{\max}}
 =\frac{283}{1250}\cdot\frac{27783}{2500}\cdot\frac{1250}{13302}
 =\frac{873621}{3695000}.
\]
By \textup{(N33.23)}, \(g_i^2\le3A<2620863/3695000<289/400\), so
\(g_i<17/20\).  The height statements are
Corollary~\ref{cor:f2v51-collapse}.
\end{proof}

\subsection{Exact certificate on the sliver}

\begin{lemma}[Three-chart cover of the sliver]
\label{lem:f2v51-cover}
Every reduced candidate with \(T>11\) belongs to one of three charts,
indexed by the factor \(q\in\{1,2,3\}\) attaining \(T\), with
\begin{equation}
 h=\epsilon_q=\frac{18113}{2000}
   +\left(\frac{45849}{5000}-\frac{18113}{2000}\right)H,
 \qquad
 \epsilon_k=hE_k\ (k\ne q),
 \qquad
 H\in[0,1],\ E_k\in\left[0,\tfrac1{80}\right],
 \label{eq:f2v51-chart}
\end{equation}
and \((g_1,g_2)\in[0,17/20]^2\).
\end{lemma}

\begin{proof}
From \(T>11\) and \(\alpha_3<19435/10000\),
\(h=T-\alpha_3>11-19435/10000=18113/2000\); from
\(T<15-2\alpha_3\) and \(\alpha_3>19434/10000\),
\(h<15-3\alpha_3<45849/5000\).  For \(k\ne q\),
Corollary~\ref{cor:f2v51-sliver} gives
\(E_k=\epsilon_k/h<(283/2500)/(18113/2000)=1132/90565<1/80\), because
\(1132\cdot80=90560<90565\).  The gap box is
\eqref{eq:f2v51-sliver}.
\end{proof}

\begin{lemma}[Three exact Bernstein ledgers]
\label{lem:f2v51-ledgers}
On each chart of Lemma~\ref{lem:f2v51-cover}, every tensor-product
Bernstein coefficient of the pulled-back obstruction polynomial
\(\mathcal M_3=6p_2p_6-6p_4^2-p_2p_3^2\) is strictly negative.  Each
chart has multidegree \((6,6,3,3,3)\) with the maximal factor's height
degree raised to \(4\), hence \(3920\) coefficients, so the certificate
contains \(11760\) strictly negative exact coefficients.  The power-basis
term counts and the largest (least negative) upper endpoints are
\begin{equation}
 \begin{array}{c|c|c}
 q&\text{terms}&\text{largest upper endpoint}\\ \hline
 1&470&-\dfrac{9526894814157365421}{25953125000000000}\ \text{at }(1,0,4,0,0)\\
 2&472&-\dfrac{68077292451313813759}{207625000000000000}\ \text{at }(1,1,0,4,0)\\
 3&470&-\dfrac{9526894814157365421}{25953125000000000}\ \text{at }(1,0,0,0,4)
 \end{array}
 \label{eq:f2v51-ledger}
\end{equation}
\end{lemma}

\begin{proof}
The chart of Lemma~\ref{lem:f2v51-cover} is the V46 ordered-height
substitution \textup{(N46.1)} on the gap box \([0,17/20]^2\), followed
by the exact affine reparametrization \(E_k=E_k'/80\) of the two
non-maximal height coordinates.  The permanent verifier builds the
pulled-back polynomial over \(\mathbb Q(\alpha_3)\) with the same
\(\mathcal M_3\) object and the same algebraic interval evaluation as
V44--V50, converts it to the tensor Bernstein basis, and asserts the
multidegree, the coefficient count, an empty ambiguity list, and the
displayed extremal endpoints.  Strict negativity of every coefficient
gives \(\mathcal M_3<0\) on the chart by the convex-hull property.
\end{proof}

\begin{remark}[Permanent chained verifier]
\label{rem:f2v51-verifier}
The verifier is
\path{scripts/verify_grh_v51_total_height_budget_sliver.py}.  It first
reruns the complete V50 chain (hence V49 and its predecessors), then
asserts the rational coverage inequalities of
Corollary~\ref{cor:f2v51-sliver} and Lemma~\ref{lem:f2v51-cover},
including \(3A_{\max}<(17/20)^2\) and \(1132/90565<1/80\), and finally
computes the three ledgers.  It has \(114\) lines,
\(4447\) bytes, and SHA--256
\begin{center}
\texttt{C3B43CBF46F8495625573DB9056CF436}\\
\texttt{D29C862F3941C4CE9EAB0F08A7F527B7}.
\end{center}
It terminates with \texttt{VERIFIED}.  The exploratory spine barrier
and obstruction-switch scripts named at the start of this section were
rerun and also terminate with \texttt{VERIFIED}; they are witnesses, not
part of the proof.
\end{remark}

\begin{theorem}[Complete classification of the three-pair heat sextic above threshold]
\label{thm:f2v51-complete}
Let \(c_1,c_2,c_3\in\mathbb R\) and \(t_j\ge\alpha_3\) for
\(j=1,2,3\).  Then \(\mathcal Q_{\bm c,\bm t}\) is real-rooted if and
only if, for some \(a\in\mathbb R\),
\begin{equation}
 c_1=c_2=c_3=a,\qquad t_1=t_2=t_3=\alpha_3.
 \label{eq:f2v51-symmetric}
\end{equation}
Equivalently: for every three-pair input and every heat time \(K>0\),
real-rootedness of \(e^{-(K/2)D_z^2}\prod_j((z-c_j)^2+y_j^2)\) forces
\begin{equation}
 \boxed{\min_jy_j^2\le\alpha_3K,}
 \label{eq:f2v51-shallow}
\end{equation}
with equality only for the centered symmetric triple.
\end{theorem}

\begin{proof}
Translate so that \(c_1+c_2+c_3=0\) and order the centers.  If
\(T\le11\), Corollary~\ref{cor:f2v49-through-eleven} gives the claim.
If \(T\ge15-2\alpha_3\), Corollary~\ref{cor:f2v51-collapse} excludes
real-rootedness.  If \(11<T<15-2\alpha_3\), Lemma~\ref{lem:f2v51-cover}
places the tuple in a certified chart and Lemma~\ref{lem:f2v51-ledgers}
gives \(\mathcal M_3<0\), contradicting the Gram inequality
\(\mathcal M_3=\det\bigl(p_{i+j}\bigr)_{i,j\in\{0,1,3\}}\ge0\) for real
roots with \(p_1=0\) (the necessary inequality \textup{(N44.2)} used
since V44).  The converse and the scaling to general \(K\) are the V13
threshold factorization and the homogeneity
\(e^{-(K/2)D_z^2}P(z)=K^3\,(e^{-D_w^2/2}\tilde P)(z/\sqrt K)\) with
\(\tilde P(w)=K^{-3}P(\sqrt K\,w)\).  The shallow-pair form follows
because \(\min_jy_j^2>\alpha_3K\) is the excluded case and
\(\min_jy_j^2=\alpha_3K\) with real-rootedness forces
\eqref{eq:f2v51-symmetric}.
\end{proof}

\begin{corollary}[Sharp three-channel gamma barriers]
\label{cor:f2v51-gamma}
In the critical gamma-block notation of \textup{(N9.24)}, consider a
fixed block with exactly three conjugate-pair channels, no real gamma
factor, ordinates \(\gamma_1,\gamma_2,\gamma_3\), and ladder count
\(R\ge6\).  If its critical heat envelope is real-rooted, then
\begin{equation}
 \boxed{
 \kappa^2\ge\frac{R\min_j\gamma_j^2}{\alpha_3(1+\theta)}
 \ge\frac{6}{\alpha_3}\frac{\min_j\gamma_j^2}{1+\theta}}
 \qquad\text{and}\qquad
 \boxed{
 \kappa^2\ge\frac{R\,(\gamma_1^2+\gamma_2^2+\gamma_3^2)}{15(1+\theta)}.}
 \label{eq:f2v51-gamma}
\end{equation}
The first bound extends the centered equal-height coefficient
\(6/\alpha_3=3.087\ldots\) of \textup{(N10.15)} to arbitrary centers and
heights, improving the pure-block coefficient \(2\) of
\textup{(N9.23)} for \(m=3\); the second charges all three channels.
\end{corollary}

\begin{proof}
By \textup{(N9.24)}, \(y_j^2=\gamma_j^2R/(\theta^2\kappa^2)\) and
\(K=(1+\theta)/\theta^2\).  Insert these in \eqref{eq:f2v51-shallow} and
in \(\sum_jy_j^2\le15K\), which is Theorem~\ref{thm:f2v51-budget} after
the scaling in the proof of Theorem~\ref{thm:f2v51-complete}.
\end{proof}

\subsection*{Firewall and V52 direction}
\label{fire:f2v51}
Theorem~\ref{thm:f2v51-complete} closes the finite three-pair model that
occupied V10--V50: no candidate survives at any height.  It is a theorem
about one sextic heat envelope.  It does not control blocks whose size
grows with the Jensen degree, does not address real padding inside the
same block, and does not perform the infinite-product passage or the
arithmetic forcing step.  It is not a proof of GRH; GRH remains open.

The upstream tools introduced here are degree-free: the exact rates
\eqref{eq:f2v51-rates}, the budget \(\sum_jt_j\le m(2m-1)K\), the strip
lemma, and the deficit bound.  V52 should not restart slab marching for
four pairs.  It should first test whether a weighted functional
\(\sum_k\omega(y_k)\) with the same cancellation structure has its maximal
rate attained exactly at the symmetric threshold configuration, which
would give the \(m\)-pair shallow-pair theorem without certificates; the
centered thresholds \(\alpha_m\) from the Hermite-type generalization of
\textup{(N10.9)} are the test cases.  If that fails, the deficit bound
should be sharpened by tracking \(U(\tau)\) instead of its initial value,
since the present bound wastes the whole horizontal growth.

\subsection*{Assumption ledger}
\label{ass:f2v51}
The new results use only: the backward heat flow of a monic real sextic
and its zero dynamics (Lemma~\ref{lem:f2v51-flow}), elementary
symmetrization identities, Rouch\'{e} continuity of zeros, the concavity
of \(u/(u+4T)\), the isolating interval \(19434/10000<\alpha_3<19435/10000\)
proved in V10, the exact V33 gap coordinates, the necessary Gram
inequality \(\mathcal M_3\ge0\), and exact rational or algebraic interval
arithmetic in the chained verifier.  No numerical sign decision,
zero-density hypothesis, or GRH input is used.
\endgroup

\section*{Version V51 append-only record}
\addcontentsline{toc}{section}{Version V51 append-only record}

V51 was appended byte-for-byte before the terminal marker of validated
V50, whose installed SHA--256 was
\[
 \texttt{E6EE171D8AB266216F348DF9FCD1FD53}
 \allowbreak\texttt{3349DBF13CA85E4C19FA4DDF3212051A}
\]
with \(14578\) lines and \(463785\) bytes.  It proves the sharp
total-height budget \(\sum_jt_j\le15\), the strip and deficit bounds, and
the complete classification of the three-pair heat sextic above the
threshold \(\alpha_3\).  After validation V51 replaces V50 as the sole
active numeric GRH note; the frozen \texttt{V100\_f2} archive is
unchanged.

% =============================================================================
% END APPENDED V51 MATERIAL
% =============================================================================




\clearpage
\section{V52: the mixed-block band, general-degree budgets, and two coalescence no-go witnesses}
\label{sec:f2v52-mixed-band}
\begingroup
\setcounter{equation}{0}
\renewcommand{\theequation}{N52.\arabic{equation}}
\renewcommand{\theHequation}{f2v52.\arabic{equation}}

V51 closed the pure three-pair block.  The V52 direction asked which
finite block still reaches the open gamma band.  Reading the archive's
ladder inequalities again gives a precise answer.  For a critical block
with \(m\) conjugate-pair channels, \(p\) real gamma factors, ladder
count \(R\ge2m+p\), and least ordinate \(\gamma\), the V9 bound
\(y_{\min}^2\le(m+2p)K\) gives \(\kappa^2\ge\gamma^2R/((m+2p)(1+\theta))\),
which stays above the band top \(\kappa^2=\gamma^2/(1+\theta)\) exactly
when \(m+2p\le R\).  Since \(R\) may equal \(2m+p\), every finite block
misses the band as soon as
\begin{equation}
 \boxed{y_{\min}^2\le(2m+p)K=nK,\qquad n:=\deg P_0=2m+p.}
 \label{eq:f2v52-target}
\end{equation}
This holds for \(m=1\) (V5, \(p+2=n\)) and for \(p\le m\) (V9,
\(m+2p\le2m+p\)).  The open cases are exactly \(m\ge2\), \(p>m\).  This
section records the general-degree flow theorems that survive from V51,
reduces the mixed problem exactly through Walsh coalescence, proves with
exact witnesses that the two natural coalescence relaxations cannot
certify \eqref{eq:f2v52-target}, and states the resulting sharp
conjecture with its numerical support.

\subsection{General-degree flow theorems}

Let \(P_0\) be a monic real polynomial of degree \(n\) with \(m\)
conjugate pairs of heights \(y_1,\ldots,y_m>0\) and \(p=n-2m\) real
zeros, and let \(\mathcal Q_\tau=e^{-\tau D_z^2/2}P_0\).  Write
\(m(\tau)\) and \(p(\tau)\) for the numbers of nonreal pairs and real
zeros of \(\mathcal Q_\tau\).  Lemma~\ref{lem:f2v51-flow} and
Proposition~\ref{prop:f2v51-rates} hold verbatim with \(6\) replaced by
\(n\) and \(30\) by \(n(n-1)\); the large-\(\tau\) limit is
\(\mathrm{He}_n\), which has \(n\) simple real zeros.

\begin{lemma}[Real zeros are never lost]
\label{lem:f2v52-poulain}
\(p(\tau)\) is nondecreasing and \(m(\tau)\) nonincreasing on
\([0,\infty)\).
\end{lemma}

\begin{proof}
Let \(a\ne0\) be real and let \(q\) be a real polynomial of degree
\(n\) with \(k\) real zeros counted with multiplicity.  Then
\(q+aq'=a\,e^{-z/a}\bigl(e^{z/a}q\bigr)'\), and \(e^{z/a}q\) has the same
\(k\) real zeros.  Rolle's theorem gives at least \(k-1\) real zeros of
\((e^{z/a}q)'\) between consecutive distinct ones (a zero of
multiplicity \(\mu\) counts \(\mu-1\) times), and one more real zero
outside them because \(e^{z/a}q\) tends to \(0\) at the end of the line
where \(z/a\to-\infty\) and is nonzero beyond its extreme real zero.
Hence \(q+aq'\), which has degree \(n\), has at least \(k\) real zeros.
Applying this to the \(2N\) factors of
\((1-\tau D_z^2/(2N))^N=\prod(1-bD_z)(1+bD_z)\) and passing to the
coefficientwise limit as in Lemma~\ref{lem:f2v51-strip} shows that
\(\mathcal Q_\tau\) has at least as many real zeros as \(P_0\); real
zeros of the approximants converge to real zeros of the limit.  The
same argument from any intermediate time gives monotonicity.
\end{proof}

\begin{theorem}[General-degree total-height budget]
\label{thm:f2v52-budget}
If \(e^{-(K/2)D_z^2}P_0\) is real-rooted, then
\begin{equation}
 \boxed{\sum_{j=1}^my_j^2\le\frac{n(n-1)-p(p-1)}{2}K=m(2m+2p-1)K.}
 \label{eq:f2v52-budget}
\end{equation}
Equality holds if and only if \(e^{-(K/2)D_z^2}P_0=(z-a)^n\) for a real
\(a\) and \(p\equiv n\pmod 2\), that is, \(p\in\{0,1\}\).
\end{theorem}

\begin{proof}
Scale to \(K=1\).  With the notation of
Proposition~\ref{prop:f2v51-rates}, ordered pairs of real zeros
contribute \(0\) to \(\sum_{k\ne l}\sin^2\theta_{kl}\), every other
ordered pair contributes at most \(1\), and by
Lemma~\ref{lem:f2v52-poulain} there are at least \(p\) real zeros at
every time.  Hence
\(2\sum_jy_j^2=\int_0^1\sum_{k\ne l}\sin^2\theta_{kl}\,\dd\tau
 \le n(n-1)-p(p-1)\).  Equality forces, as in
Theorem~\ref{thm:f2v51-budget}, all zeros onto one vertical line at
almost every time and hence \(\mathcal Q_1=(z-a)^n\); then
\(P_0=e^{D_z^2/2}(z-a)^n=i^n\mathrm{He}_n\bigl(-i(z-a)\bigr)\), whose
zeros are \(a+i\lambda\) with \(\lambda\) running over the real zeros of
\(\mathrm{He}_n\); it has exactly one real zero if
\(n\) is odd and none if \(n\) is even, so \(p=n\bmod2\).  Conversely,
these inputs give \(\sum_jy_j^2=\tfrac12\sum_\lambda\lambda^2
 =\tfrac12n(n-1)\) by the Newton sum of \(\mathrm{He}_n\), which equals
the right side of \eqref{eq:f2v52-budget} for \(p\in\{0,1\}\).
\end{proof}

\begin{corollary}[Ladder form]
\label{cor:f2v52-ladder}
In the notation of \textup{(N9.24)}, a critical block with \(m\) pair
channels of ordinates \(\gamma_j\) and \(p\) real factors satisfies
\begin{equation}
 \kappa^2\ge\frac{R\sum_{j=1}^m\gamma_j^2}{m(2m+2p-1)(1+\theta)}.
 \label{eq:f2v52-ladder}
\end{equation}
\end{corollary}

For the least ordinate this gives only \(y_{\min}^2\le(2m+2p-1)K\), which
is weaker than \eqref{eq:f2v52-target} whenever \(p\ge2\).  The sum
budget therefore does not close the band; it charges the companions but
not sharply enough.

\subsection{Exact Walsh reduction of the mixed problem}

Fix the lowest pair, translate its collision point to \(0\), and scale
the first collision time to \(1\) exactly as in the proof of
Theorem~\ref{thm:f2v5-global}; write
\begin{equation}
 \Pi(w)=\bigl((w-c)^2+t\bigr)\prod_{i=1}^p(w-r_i)
 \prod_{j=2}^m(w-w_j)(w-\overline{w_j}),
 \qquad t=\frac{y_{\min}^2}{T},\quad T\le K,
 \label{eq:f2v52-pi}
\end{equation}
with \(r_i\in\mathbb R\) and \(|\operatorname{Im}w_j|\ge\sqrt t\).  Let
\(\mathcal F(\Pi):=Q(0)-iQ'(0)\), \(Q=e^{-D_w^2/2}\Pi\).  Real-rootedness
at time \(K\) forces \(\mathcal F(\Pi)=0\) at the first collision time,
because \(Q\) then has a real zero of multiplicity at least two at the
translated point.

\begin{theorem}[Two-stage coalescence]
\label{thm:f2v52-walsh}
Suppose \(\mathcal F(\Pi)=0\) and \(t>p+2\).  Then there exist
\(\zeta\) in the closed upper half-plane and, for the two relaxations
below, coalescence points with
\begin{enumerate}[label=\textup{(\roman*)}]
\item (two half-planes) \(\omega_+\), \(\omega_-\) with
\(\operatorname{Im}\omega_+\ge\sqrt t\), \(\operatorname{Im}\omega_-\le-\sqrt t\),
such that \(\mathcal F\) vanishes for
\(\bigl((w-c)^2+t\bigr)(w-\zeta)^p(w-\omega_+)^{m-1}(w-\omega_-)^{m-1}\); or
\item (disk exterior) \(\omega\) with \(|\omega-c|\ge\sqrt t\) such that
\(\mathcal F\) vanishes for
\(\bigl((w-c)^2+t\bigr)(w-\zeta)^p(w-\omega)^{2m-2}\), provided the
functional of \(\bigl((w-c)^2+t\bigr)(w-\zeta)^p\) is nonzero at the
chosen \(\zeta\).
\end{enumerate}
\end{theorem}

\begin{proof}
\(\mathcal F(\Pi)\) is multi-affine and symmetric in each group of
padding zeros with the other groups fixed.  The upper zeros \(w_j\) lie
in the closed half-plane \(\operatorname{Im}w\ge\sqrt t\), the lower ones
in \(\operatorname{Im}w\le-\sqrt t\), and all of them in the closed
exterior of the open disk \(|w-c|<\sqrt t\), since
\(|w_j-c|\ge|\operatorname{Im}w_j|\).  Walsh's coincidence theorem
replaces a group of variables lying in a circular region by a common
point of that region without changing the value of a symmetric
multi-affine function, provided the region is convex or the function
has full degree in the group.  Half-planes are convex, which gives (i)
after coalescing the two complex groups and then the real group.  For
(ii) coalesce the complex group first; the top coefficient in that group
is the functional of \(\bigl((w-c)^2+t\bigr)\prod_i(w-r_i)\), which is
nonzero for \(t>p+2\) by Theorem~\ref{thm:f2v5-collision}; then coalesce
the real group in the convex closed upper half-plane, which yields the
stated proviso at the final \(\zeta\).
\end{proof}

Either relaxation would prove \eqref{eq:f2v52-target} if its diagonal
functional had no zero in the stated region for \(t>2m+p\).  For
\(m=1\) both reduce to Theorem~\ref{thm:f2v5-diagonal}.  For \(m\ge2\)
they do not.

\begin{proposition}[Both relaxations fail at \((m,p)=(2,2)\), \(t=9\)]
\label{prop:f2v52-nogo}
Let \(m=p=2\), so \(2m+p=m+2p=6\), and let \(t=9\).  Write
\(F_\pm(\zeta)\) and \(F_\circ(\zeta)\) for the diagonal functionals of
relaxations \textup{(i)} and \textup{(ii)}; both are quadratics in
\(\zeta\).
\begin{enumerate}[label=\textup{(\alph*)}]
\item For \(c=33/10\), \(\omega_+=283/10+3i\), \(\omega_-=3/5-3i\),
\begin{equation}
\begin{aligned}
 F_\pm(\zeta)={}&\Bigl(\tfrac{4157961}{5000}-\tfrac{92997}{100}i\Bigr)\zeta^2
 +\Bigl(\tfrac{76377}{50}+\tfrac{1851861}{2500}i\Bigr)\zeta
 -\Bigl(\tfrac{7336461}{5000}+\tfrac{37851}{125}i\Bigr),
\end{aligned}
\label{eq:f2v52-fpm}
\end{equation}
and \(F_\pm\) has a zero in the disk \(|\zeta-(341/500+21i/250)|<1/20\),
which lies in the open upper half-plane.
\item For \(c=389/50\) and \(\omega=29/50-63i/100\), so that
\(|\omega-c|^2=522369/10000\ge9\),
\begin{equation}
\begin{aligned}
 F_\circ(\zeta)={}&\Bigl(\tfrac{32365319}{5000000}-\tfrac{721399}{781250}i\Bigr)\zeta^2
 -\Bigl(\tfrac{31054226}{390625}+\tfrac{782809681}{2500000}i\Bigr)\zeta\\
 &-\Bigl(\tfrac{131360639}{5000000}-\tfrac{4585821}{3125000}i\Bigr),
\end{aligned}
\label{eq:f2v52-fcirc}
\end{equation}
and \(F_\circ\) has a zero in the disk
\(|\zeta-(5289/1000+24523i/500)|<1/40\), which lies in the open upper
half-plane.
\end{enumerate}
Consequently neither relaxation can exclude \(t=9>6\), and neither can
prove \eqref{eq:f2v52-target} for \((m,p)=(2,2)\).
\end{proposition}

\begin{proof}
The coefficients are obtained by expanding
\(\Pi\) in \(w\), applying
\((e^{-D_w^2/2}w^j)(0)=(-1)^{j/2}j!/(2^{j/2}(j/2)!)\) for even \(j\)
and the analogous formula for the derivative at odd \(j\), and
collecting powers of \(\zeta\); the verifier below reproduces them
exactly over \(\mathbb Q+i\mathbb Q\).  For a quadratic
\(F(\zeta)=a_2(\zeta-\zeta_0)^2+F'(\zeta_0)(\zeta-\zeta_0)+F(\zeta_0)\),
Rouch\'{e}'s theorem against the linear term gives exactly one zero in
\(|\zeta-\zeta_0|<r\) whenever
\(|F'(\zeta_0)|\,r>|F(\zeta_0)|+|a_2|r^2\).  With rational upper
bounds for \(|F(\zeta_0)|\) and \(|a_2|\) and a rational lower bound
for \(|F'(\zeta_0)|\) obtained from integer square roots, the two
inequalities hold with exact margins
\(1084835646862537181/7812500000000000\) and
\(4926046556749829667726631413/625000000000000000000000000\).  The
disks lie in the open upper half-plane because
\(21/250-1/20>0\) and \(24523/500-1/40>0\).  The domain constraints
\(\operatorname{Im}\omega_\pm=\pm3=\pm\sqrt t\) and
\(|\omega-c|\ge\sqrt t\) are immediate.
\end{proof}

\begin{remark}[Permanent verifier]
\label{rem:f2v52-verifier}
The exact computation is
\path{scripts/verify_grh_v52_mixed_block_nogo.py}; it uses only
\texttt{fractions}, asserts the displayed coefficients, the domain
constraints, and both Rouch\'{e} inequalities, and terminates with
\texttt{VERIFIED}.  It has \(126\) lines, \(4482\) bytes, and SHA--256
\begin{center}
\texttt{17C0AE47FCE8F1A13E2906CF5FAB161B}\\
\texttt{D0D58689A28E6841A14A3A257A5C2D42}.
\end{center}
\end{remark}

\begin{remark}[Why the relaxations lose]
\label{rem:f2v52-why}
In witness (a) the lower coalescence point sits three units below the
real axis near the target while the upper one is far away; in witness
(b) the double zero sits close to the real axis.  Both configurations
act on the target like extra real zeros, which is exactly what the
conjugate symmetry of the original pair forbids.  A single-point
coalescence of all padding zeros is worse still: the collision
equations hold for every pair's factorization, so any single-domain
diagonal bound would apply to every pair, contradicting the equality
input of Theorem~\ref{thm:f2v51-budget}, which has a pair at height
\(11.05\) inside a real-rooted block of degree six.  Any coalescence
proof must therefore retain the pairing \(w_j\leftrightarrow\overline{w_j}\),
which Walsh's theorem cannot see.
\end{remark}

\subsection{The sharp mixed-block conjecture}

\begin{problem}[Band closure for all finite blocks]
\label{prob:f2v52-band}
Prove \eqref{eq:f2v52-target} for all \(m\ge2\), \(p>m\).  The
expected sharp statement is
\begin{equation}
 y_{\min}^2\le nK,\qquad\text{with equality only for }
 P_0(z)=\bigl((z-a)^2+nK\bigr)(z-a)^{n-2},\ n\text{ odd},
 \label{eq:f2v52-conjecture}
\end{equation}
that is, only in the one-pair Hermite--Laguerre blocks of
\textup{(N5.25)}.
\end{problem}

\begin{remark}[Numerical exploration, not a proof]
\label{rem:f2v52-numerics}
A floating-point Nelder--Mead search over centers, heights, and real
zeros, maximizing the least height subject to real-rootedness of the
heat image at \(K=1\), found the following largest values (each a lower
bound for the true maximum; forty restarts each):
\[
\begin{array}{c|ccccccc}
(m,p)&(1,3)&(2,0)&(2,1)&(2,2)&(2,3)&(2,4)&(3,4)\\ \hline
\text{found }y_{\min}^2&4.93&1.50&2.49&4.33&5.11&5.73&6.06\\
n=2m+p&5&4&5&6&7&8&10\\
m+2p\ (\text{V9})&7&2&4&6&8&10&11
\end{array}
\]
The \((1,3)\) and \((2,0)\) rows recover the proved values \(5\) and
\(3/2\) to the search accuracy.  In every open case the found value is
well below \(n\), and adding a pair never raised the maximal least
height above the one-pair value \(p+2\).  These numbers guide the
conjecture and are not used in any proof.
\end{remark}

\subsection*{Firewall and V53 direction}
\label{fire:f2v52}
Nothing in this section proves \eqref{eq:f2v52-target} for \(p>m\), and
nothing here proves GRH; GRH remains open.  The proved content is the
general-degree budget, the exact two-stage reduction, and two exact
no-go witnesses that close the naive coalescence routes.

The mechanism behind the \(m=1\) theorem is that the real zeros repel
and spread on the Hermite scale, so that on average at most
\((p+1)/2\) of them sit under the target.  With \(m\ge2\) the target
bound has slack \(m-1\) beyond that, so V53 should not seek sharpness.
It should look for a flow inequality of the form
\(\int_0^T\sum_i\sin^2\theta_i\,\dd\tau\le\tfrac{p+m}{2}T\) for the
lowest pair, using a monotone functional of the real zeros that survives
the horizontal forces of the other pairs, or alternatively an
interlacing argument comparing the real zeros of \(\mathcal Q_\tau\)
with those of the backward heat image of \(\prod_i(w-r_i)\) alone.
Either would close the band for every finite block.

\subsection*{Assumption ledger}
\label{ass:f2v52}
The proved statements use only the V51 flow identities extended to
degree \(n\), Rolle's theorem, coefficientwise limits of the
\((1\pm bD_z)\) factorization, Walsh's coincidence theorem in the two
forms stated, Theorem~\ref{thm:f2v5-collision}, and exact rational
arithmetic with integer square roots in the permanent verifier.  The
numerical table is exploration only.
\endgroup

\section*{Version V52 append-only record}
\addcontentsline{toc}{section}{Version V52 append-only record}

V52 was appended byte-for-byte before the terminal marker of validated
V51, whose installed SHA--256 was
\[
 \texttt{61927F9477BD72FAD280D5EF40429199}
 \allowbreak\texttt{32C7687ED895DE4571A71F2D81450134}
\]
with \(15170\) lines and \(489091\) bytes.  It identifies the mixed
blocks with \(p>m\) as the only finite architectures not yet excluded
from the gamma band, proves the general-degree budget, and closes both
naive Walsh relaxations with exact witnesses.  After validation V52
replaces V51 as the sole active numeric GRH note; the frozen
\texttt{V100\_f2} archive is unchanged.

% =============================================================================
% END APPENDED V52 MATERIAL
% =============================================================================




\clearpage
\section{V53: reals-only coalescence, a generating identity, and the shrinking strip}
\label{sec:f2v53-reals-only}
\begingroup
\setcounter{equation}{0}
\renewcommand{\theequation}{N53.\arabic{equation}}
\renewcommand{\theHequation}{f2v53.\arabic{equation}}

V52 showed that every coalescence of the companion pairs destroys the
conjugate symmetry that the mixed-block bound depends on.  The remaining
option is to coalesce only the real padding zeros and to keep the
companion pairs exactly.  This section proves that reduction, computes the
resulting diagonal functional in closed form through a generating
identity, records the shrinking-strip lemma that the flow gives for free,
and states precisely which real-rootedness statement is now missing.

\subsection{The reals-only reduction}

Retain the notation \eqref{eq:f2v52-pi}: after translating the collision
point of the first real collision to \(0\) and scaling that time to
\(1\), the block is
\begin{equation}
 \Pi(w)=f(w)\,G(w)\prod_{i=1}^p(w-r_i),\qquad
 f(w)=(w-c)^2+t,\qquad
 G(w)=\prod_{j=2}^m\bigl((w-c_j)^2+t_j\bigr),
 \label{eq:f2v53-pi}
\end{equation}
where \(t=y_{\min}^2/T\) belongs to the lowest pair and \(t_j\ge t\) for
the companions.  For a polynomial \(\Pi\) put
\(\mathcal F(\Pi):=Q(0)-iQ'(0)\) with \(Q=e^{-D_w^2/2}\Pi\).

\begin{theorem}[Reals-only coalescence]
\label{thm:f2v53-reduction}
If \(e^{-(K/2)D_z^2}P_0\) is real-rooted, then with the normalization
\eqref{eq:f2v53-pi} there is \(\zeta\) in the closed upper half-plane with
\begin{equation}
 F_G(\zeta):=\mathcal F\bigl(f\,G\,(w-\zeta)^p\bigr)=0 .
 \label{eq:f2v53-diag}
\end{equation}
Consequently, if for some \(t_0>0\) one knows that \(F_G\) has no zero
in the closed upper half-plane whenever \(t>t_0\), for every real
\(G\) of the displayed form with all \(t_j\ge t\) and every real \(c\),
then \(y_{\min}^2\le t_0K\).
\end{theorem}

\begin{proof}
At the first collision time the heat image has a real zero of
multiplicity at least two at the translated origin, so
\(\mathcal F(\Pi)=0\).  The map \((r_1,\ldots,r_p)\mapsto\mathcal F(\Pi)\)
is symmetric and affine in each \(r_i\), and the \(r_i\) lie in the
closed upper half-plane, which is convex.  Walsh's coincidence theorem
therefore supplies \(\zeta\) in that half-plane with
\(\mathcal F(f\,G\,(w-\zeta)^p)=\mathcal F(\Pi)=0\); no degree condition is
needed for a convex region.  The companion factor \(G\) is untouched.
The last sentence follows because \(t=y_{\min}^2/T\ge y_{\min}^2/K\).
\end{proof}

For \(m=1\) this is exactly the reduction used in
Theorem~\ref{thm:f2v5-collision}, and Theorem~\ref{thm:f2v5-diagonal}
supplies \(t_0=p+2\).  The band-closing target
\eqref{eq:f2v52-target} is \(t_0=n=2m+p\).

\subsection{Closed form of the diagonal functional}

\begin{lemma}[Leibniz form for a quadratic factor]
\label{lem:f2v53-leibniz}
For any polynomial \(R\) and \(f=(w-c)^2+t\), with
\(H:=e^{-D_w^2/2}R\),
\begin{equation}
 e^{-D_w^2/2}(fR)=(f-1)H-f'H'+H''.
 \label{eq:f2v53-leibniz}
\end{equation}
Hence, with \(h_k:=H^{(k)}(0)\) and \(q:=c^2+t\),
\begin{equation}
 \mathcal F(fR)=A-iB,\qquad
 A=(q-1)h_0+2ch_1+h_2,\qquad
 B=-2ch_0+(q-3)h_1+2ch_2+h_3 .
 \label{eq:f2v53-AB}
\end{equation}
\end{lemma}

\begin{proof}
Expand \((fR)^{(2k)}=fR^{(2k)}+2kf'R^{(2k-1)}+k(2k-1)f''R^{(2k-2)}\)
with \(f''=2\), insert into
\(e^{-D_w^2/2}(fR)=\sum_k(-\tfrac12)^k(fR)^{(2k)}/k!\), and reindex: the
middle sum is \(-f'H'\), and the last sum equals
\(-\sum_j(-\tfrac12)^j(2j+1)R^{(2j)}/j!=-(H-H'')\).  Differentiating
\eqref{eq:f2v53-leibniz} once and evaluating both at \(0\), with
\(f(0)=q\), \(f'(0)=-2c\), gives \eqref{eq:f2v53-AB}.
\end{proof}

\begin{lemma}[Generating identity]
\label{lem:f2v53-generating}
Let \(G\) be any polynomial, \(\widetilde G:=e^{-D^2/2}G\), and
\(\mathrm{He}_r=e^{-D^2/2}z^r\).  For \(k\ge0\) and \(p\ge0\), the
\(k\)-th derivative at \(0\) of \(e^{-D_w^2/2}[G(w)(w-\zeta)^p]\), viewed
as a polynomial in \(\zeta\), equals
\begin{equation}
 h_k(\zeta)=\sum_{l=0}^{k}\binom kl\frac{p!}{(p-k+l)!}(-1)^{p-k+l}
 \bigl(\widetilde G^{(l)}(D_\zeta)\,\mathrm{He}_{p-k+l}\bigr)(\zeta),
 \label{eq:f2v53-hk}
\end{equation}
terms with \(p-k+l<0\) being omitted.  In particular
\begin{equation}
 \boxed{h_0(\zeta)=(-1)^p\bigl(\widetilde G(D_\zeta)\mathrm{He}_p\bigr)(\zeta)
 =(-1)^p\,e^{-D_\zeta^2/2}\bigl[\widetilde G(D_\zeta)\zeta^p\bigr].}
 \label{eq:f2v53-h0}
\end{equation}
\end{lemma}

\begin{proof}
All identities below are between formal power series in \(u\) whose
coefficients are polynomials in \(w\) or \(\zeta\), so they may be read
coefficientwise.  Since \(D_w(e^{uw}\varphi)=e^{uw}(D_w+u)\varphi\),
\[
 e^{-D_w^2/2}\bigl(G(w)e^{uw}\bigr)
 =e^{uw}e^{-(D_w+u)^2/2}G
 =e^{uw}e^{-u^2/2}e^{-uD_w}\widetilde G
 =e^{uw}e^{-u^2/2}\widetilde G(w-u),
\]
using \(e^{-uD_w}\psi(w)=\psi(w-u)\) for polynomials \(\psi\).  Hence
\[
 \sum_{p\ge0}\frac{u^p}{p!}\,\partial_w^k\,e^{-D_w^2/2}\bigl[G(w)(w-\zeta)^p\bigr]\Big|_{w=0}
 =e^{-u\zeta-u^2/2}\sum_{l=0}^k\binom kl u^{k-l}\widetilde G^{(l)}(-u).
\]
Replacing \(u\) by \(-u\) turns \(e^{-u\zeta-u^2/2}\widetilde G^{(l)}(-u)\)
into \(e^{u\zeta-u^2/2}\widetilde G^{(l)}(u)
 =\widetilde G^{(l)}(D_\zeta)e^{u\zeta-u^2/2}
 =\sum_r\frac{u^r}{r!}\widetilde G^{(l)}(D_\zeta)\mathrm{He}_r(\zeta)\),
because \(e^{u\zeta-u^2/2}=\sum_r\mathrm{He}_r(\zeta)u^r/r!\) is the
Hermite generating function and \(D_\zeta\) acts on it as multiplication
by \(u\).  Extracting the coefficient of \(u^p/p!\) gives
\eqref{eq:f2v53-hk}; the sign \((-1)^{p-k+l}\) records the substitution
\(u\to-u\).  The second form in \eqref{eq:f2v53-h0} holds because
polynomials in \(D_\zeta\) commute.
\end{proof}

\begin{remark}[Exact verifier]
\label{rem:f2v53-verifier}
\path{scripts/verify_grh_v53_reals_only_identity.py} checks
\eqref{eq:f2v53-leibniz}, \eqref{eq:f2v53-h0}, the commutation in
\eqref{eq:f2v53-h0}, and \eqref{eq:f2v53-hk} for \(k=1,2,3\), as exact
polynomial identities over \(\mathbb Q\) on thirty random rational
samples with up to three companion pairs and \(p\le6\).  It has
\(132\) lines, \(4069\) bytes, SHA--256
\begin{center}
\texttt{079439B7B93FC012E80F906F91445C2F}\\
\texttt{35E31A086EE343DA6AD1D62E7BAF3065},
\end{center}
and terminates with \texttt{VERIFIED}.
\end{remark}

For \(G=1\) one has \(\widetilde G=1\) and \eqref{eq:f2v53-hk} reduces to
\(h_k=(-1)^{p-k}p!/(p-k)!\,\mathrm{He}_{p-k}=(-1)^{p-k}\mathrm{He}_p^{(k)}\),
which is the Hermite structure exploited in
Lemma~\ref{lem:f2v5-pick}.  In general \(h_0\) is the Hermite polynomial
transformed by the differential operator whose symbol is the
\emph{heat image of the companion factor}; the companions enter the
diagonal only through \(\widetilde G\) and its derivatives.

\subsection{The shrinking strip}

\begin{lemma}[Shrinking strip]
\label{lem:f2v53-strip}
Let \(P\) be a monic real polynomial of degree \(n\ge1\) whose zeros lie
in \(|\operatorname{Im}z|\le h\).  Then the zeros of
\(e^{-\tau D_z^2/2}P\) lie in \(|\operatorname{Im}z|\le\sqrt{h^2-\tau}\)
for \(0\le\tau\le h^2\), and \(e^{-\tau D_z^2/2}P\) is real-rooted for
\(\tau\ge h^2\).
\end{lemma}

\begin{proof}
Let \(M(\tau):=\max_k y_k(\tau)^2\), continuous on \([0,\infty)\) by
Lemma~\ref{lem:f2v51-flow}(iii).  On a component of the complement of
the finite collision set, let \(k\) attain the maximum at \(\tau\) with
\(y_k>0\).  Every other zero satisfies \(|y_l|\le y_k\), so in
\(\frac{\dd}{\dd\tau}y_k^2=-2y_k\sum_{l\ne k}(y_k-y_l)/|z_k-z_l|^2\)
every term is nonpositive, and the term \(l=\bar k\) from the conjugate
zero equals \(-2y_k\cdot2y_k/(2y_k)^2=-1\).  Hence the upper Dini
derivative of \(M\) is at most \(-1\) wherever \(M>0\), and \(M(\tau)+\tau\)
is nonincreasing on each component while \(M>0\); by continuity across
the finitely many collision times, \(M(\tau)\le h^2-\tau\) as long as the
right side is nonnegative, and \(M\) vanishes once it has reached \(0\).
Real zeros never become nonreal by Lemma~\ref{lem:f2v52-poulain}, so
real-rootedness persists for \(\tau\ge h^2\).
\end{proof}

This is de~Bruijn's classical strip theorem in the polynomial case; the
flow proof is recorded because it is the same computation as
Proposition~\ref{prop:f2v51-rates}, applied to the top zero only.

\begin{corollary}[A sufficient condition for a real-rooted diagonal]
\label{cor:f2v53-sufficient}
If all zeros of the polynomial \(\zeta\mapsto\widetilde G(D_\zeta)\zeta^p\)
lie in \(|\operatorname{Im}\zeta|\le1\), then
\(h_0(\zeta)=\pm\widetilde G(D_\zeta)\mathrm{He}_p(\zeta)\) has only real
zeros.
\end{corollary}

\begin{proof}
Combine \eqref{eq:f2v53-h0} with Lemma~\ref{lem:f2v53-strip} at
\(\tau=1\).
\end{proof}

\subsection{What is now missing}

The V5 proof of the case \(G=1\) needed two facts about the diagonal:
that \(h_0\) is real-rooted in \(\zeta\), which gave the Pick inequality
\textup{(N5.5)}, and the Hermite differential equation, which reduced
\(h_2,h_3\) to \(h_0,h_1\).  Lemma~\ref{lem:f2v53-generating} identifies
the general replacements: \(h_0=\pm\widetilde G(D)\mathrm{He}_p\), and
\(h_1,h_2,h_3\) involve the derived operators \(\widetilde G^{(l)}(D)\).

\begin{problem}[Real-rooted transformed Hermite polynomials]
\label{prob:f2v53-realrooted}
Prove that \(\widetilde G(D)\mathrm{He}_p\) has only real zeros whenever
\(G\) is a product of \(m-1\) real quadratic factors whose zeros have
\(|\operatorname{Im}|\ge\sqrt t\) with \(t>2m+p\).
\end{problem}

\begin{remark}[Numerical evidence and a failed shortcut]
\label{rem:f2v53-numerics}
Floating-point tests over \(4000\) random companion factors with
\(1\le m-1\le6\), \(1\le p\le8\), all heights at least \(\sqrt n\),
found no nonreal zero of \(\widetilde G(D)\mathrm{He}_p\), including
centered configurations with all heights exactly \(\sqrt n\).  The
shortcut of Corollary~\ref{cor:f2v53-sufficient} does not settle this:
in the same regime the zeros of \(\widetilde G(D)\zeta^p\) reach height
about \(3.3\), so the strip lemma at time \(1\) does not apply, and a
direct argument is required.  For the diagonal threshold itself, a
Nelder--Mead search for the largest \(t\) at which \(F_G\) has a zero in
the closed upper half-plane gave
\[
\begin{array}{c|cccc}
(m,p)&(2,1)&(2,2)&(2,3)&(2,4)\\ \hline
\text{found }t^*&3.57&4.47&5.58&6.55\\
n=2m+p&5&6&7&8
\end{array}
\]
with the extremal companion pair centered at \(c\) and at height exactly
\(\sqrt t\).  Each found value lies below \(n\), and in each case the
corresponding true block threshold of Remark~\ref{rem:f2v52-numerics}
is smaller still, so the reals-only diagonal is a genuinely stronger
relaxation than the two rejected in V52.  These are explorations, not
proofs.
\end{remark}

\subsection*{Firewall and V54 direction}
\label{fire:f2v53}
Theorem~\ref{thm:f2v53-reduction} and the two lemmas are proved; nothing
here proves \eqref{eq:f2v52-target} for \(p>m\), and nothing proves GRH.
GRH remains open.

V54 should attack Problem~\ref{prob:f2v53-realrooted} through the
factorization
\(\widetilde G(D)=\prod_j\bigl((D-a_j)^2+b_j^2\bigr)\) over the zeros
\(a_j\pm ib_j\) of \(\widetilde G\), whose heights satisfy
\(b_j^2\ge t-(m-1)\) by the V9 flow bound applied to \(G\) alone.  Each
factor is \((D-\lambda)(D-\bar\lambda)\); for a real-rooted \(q\), the
zeros of \(q'-\lambda q\) lie in the half-plane where
\(\operatorname{Im}(q'/q)=\operatorname{Im}\lambda\), which is
\(\operatorname{Im}z<0\), and the second factor returns them toward the
axis.  A quantitative Hermite--Biehler argument along these lines, with
the Hermite zero spacing as input, is the first thing to try; if it
succeeds, the V5 completed-square identity must be redone with the
correction terms \(\widetilde G^{(l)}(D)\mathrm{He}_{r}\) of
\eqref{eq:f2v53-hk}.

\subsection*{Assumption ledger}
\label{ass:f2v53}
The proved statements use Walsh's coincidence theorem on a convex
region, formal generating-function algebra over polynomials, the V51
flow identities, Lemma~\ref{lem:f2v52-poulain}, and the exact verifier.
The table and the real-rootedness observations are numerical
exploration only.
\endgroup

\section*{Version V53 append-only record}
\addcontentsline{toc}{section}{Version V53 append-only record}

V53 was appended byte-for-byte before the terminal marker of validated
V52, whose installed SHA--256 was
\[
 \texttt{995305F804828109C77F2E105081AB0B}
 \allowbreak\texttt{3AE87F103ECAF34F67E1364E06BE92E0}
\]
with \(15514\) lines and \(504761\) bytes.  It proves the reals-only
Walsh reduction of the mixed-block band, the Leibniz and generating
identities for the diagonal functional, and the shrinking-strip lemma,
and isolates the missing real-rootedness statement.  After validation
V53 replaces V52 as the sole active numeric GRH note; the frozen
\texttt{V100\_f2} archive is unchanged.

% =============================================================================
% END APPENDED V53 MATERIAL
% =============================================================================




\clearpage
\section{V54: a Jacobi-matrix theorem for one companion pair and the explicit two-pair Pick expression}
\label{sec:f2v54-jacobi}
\begingroup
\setcounter{equation}{0}
\renewcommand{\theequation}{N54.\arabic{equation}}
\renewcommand{\theHequation}{f2v54.\arabic{equation}}

V53 reduced the mixed-block band to Problem~\ref{prob:f2v53-realrooted}
and to a Pick-type inequality with correction terms.  For one companion
pair (\(m=2\)) both pieces become explicit.  The transformed Hermite
polynomial is a three-term Hermite combination, which is the
characteristic polynomial of a symmetric Jacobi matrix as soon as the
companion is tall; this settles the problem for \(m=2\) with an exact
threshold.  The diagonal functional itself is then reduced, over
\(\mathbb Q\), to a single polynomial inequality in seven variables,
whose truth on the target region would close the band for every block
with two conjugate pairs and any number of real padding zeros.

\subsection{One factor on a Hermite polynomial}

Let \(J_p\) be the symmetric tridiagonal \(p\times p\) matrix with zero
diagonal and off-diagonal entries \(\sqrt k\) in positions \((k,k+1)\),
\(1\le k\le p-1\).  The three-term recurrence
\(\mathrm{He}_{r+1}=x\mathrm{He}_r-r\mathrm{He}_{r-1}\) says that
\(\mathrm{He}_p=\det(xI-J_p)\).  For real \(d\) and real \(s\ge0\) let
\(J_p(d,s)\) be \(J_p\) with the last diagonal entry replaced by \(d\)
and the last off-diagonal entry by \(\sqrt s\).

\begin{theorem}[Single companion factor]
\label{thm:f2v54-jacobi}
For real \(a,b\) with \(a^2+b^2\ge p\ge2\),
\begin{equation}
 \bigl((D-a)^2+b^2\bigr)\mathrm{He}_p
 =(a^2+b^2)\Bigl[(x-d)\mathrm{He}_{p-1}-s\,\mathrm{He}_{p-2}\Bigr]
 =(a^2+b^2)\det\bigl(xI-J_p(d,s)\bigr),
 \label{eq:f2v54-jacobi}
\end{equation}
with
\begin{equation}
 d=\frac{2ap}{a^2+b^2},\qquad
 s=(p-1)\Bigl(1-\frac p{a^2+b^2}\Bigr)\ge0 .
 \label{eq:f2v54-ds}
\end{equation}
Consequently \(((D-a)^2+b^2)\mathrm{He}_p\) has only real zeros, and they
interlace with the zeros of \(\mathrm{He}_{p-1}\).
\end{theorem}

\begin{proof}
Since \(D\mathrm{He}_r=r\mathrm{He}_{r-1}\),
\(((D-a)^2+b^2)\mathrm{He}_p
 =p(p-1)\mathrm{He}_{p-2}-2ap\mathrm{He}_{p-1}+(a^2+b^2)\mathrm{He}_p\).
Substitute \(\mathrm{He}_p=x\mathrm{He}_{p-1}-(p-1)\mathrm{He}_{p-2}\) and
collect; this gives the middle expression of \eqref{eq:f2v54-jacobi}
with \(d,s\) as in \eqref{eq:f2v54-ds}.  Expanding
\(\det(xI-J_p(d,s))\) along its last row gives
\((x-d)\det(xI-J_{p-1})-s\det(xI-J_{p-2})\), which is the same
expression.  For \(s\ge0\) the matrix \(J_p(d,s)\) is real symmetric,
so its characteristic polynomial has only real zeros, and Cauchy's
interlacing theorem for the leading principal submatrix \(J_{p-1}\)
gives the interlacing.  The condition \(s\ge0\) is exactly
\(a^2+b^2\ge p\).
\end{proof}

\begin{corollary}[Problem~\ref{prob:f2v53-realrooted} for \(m=2\)]
\label{cor:f2v54-m2}
Let \(G(w)=(w-c_2)^2+t_2\) be a single companion pair.  Then
\(\widetilde G=e^{-D^2/2}G=(w-c_2)^2+(t_2-1)\), and
\(\widetilde G(D)\mathrm{He}_p\) has only real zeros whenever
\(c_2^2+t_2-1\ge p\); in particular whenever \(t_2\ge p+1\), and hence
throughout the regime \(t_2\ge t>2m+p=p+4\) of the band problem.
\end{corollary}

\begin{proof}
\(e^{-D^2/2}\) acts on a quadratic by subtracting \(1\) from its constant
term.  Apply Theorem~\ref{thm:f2v54-jacobi} with \(a=c_2\),
\(b^2=t_2-1\).
\end{proof}

\begin{remark}[Several companions]
\label{rem:f2v54-several}
For \(m-1\ge2\) companions \(\widetilde G(D)\mathrm{He}_p\) is a
\((2m-1)\)-term Hermite combination.  Expanding \(\det(xI-J)\) for a
symmetric tridiagonal \(J\) that differs from \(J_p\) in its last
\(m-1\) rows shows that every such combination with leading coefficient
one is of this form for a unique choice of the \(2(m-1)\) modified
entries, and it is real-rooted whenever the modified squared couplings
are nonnegative.  Whether tall companions always produce nonnegative
couplings is the multi-companion form of
Problem~\ref{prob:f2v53-realrooted}; it is not settled here.
\end{remark}

\begin{remark}[Exact verifier]
\label{rem:f2v54-verifier-jacobi}
\path{scripts/verify_grh_v54_jacobi_single_factor.py} checks
\eqref{eq:f2v54-jacobi} and the continuant identity for
\(\det(xI-J_p(d,s))\) exactly over \(\mathbb Q\) for forty random
\((p,a,b)\) with \(2\le p\le12\).  It has \(92\) lines, \(2911\) bytes,
SHA--256
\begin{center}
\texttt{F8B0C7BCC34F0DD2CD3562B60B0BA6BD}\\
\texttt{C9CA7BA6B63A038AC63D01015BEA4FDD},
\end{center}
and terminates with \texttt{VERIFIED}.
\end{remark}

\subsection{The two-pair diagonal as a single polynomial inequality}

Fix \(m=2\) and \(p\ge5\).  In the notation of
Theorem~\ref{thm:f2v53-reduction}, \(G(w)=(w-a)^2+t_2\) with
\(t_2\ge t\), and we write \(\beta:=t_2-1\), \(\zeta=x+iy\), and
\(v:=\mathrm{He}_{p-4}(\zeta)\), \(u:=\mathrm{He}_{p-5}(\zeta)\).

\begin{proposition}[Explicit reduction]
\label{prop:f2v54-reduction}
There are polynomials \(A_v,A_u,B_v,B_u\in\mathbb Q[x,y,c,t,a,\beta,p]\)
such that
\begin{equation}
 F_G(\zeta)=(-1)^p\Bigl[(A_v-iB_v)\,v+(A_u-iB_u)\,u\Bigr].
 \label{eq:f2v54-FG}
\end{equation}
Put \(N:=-(A_v-iB_v)\), \(D:=A_u-iB_u\), and
\begin{equation}
 E:=-\operatorname{Im}\bigl(N\overline D\bigr)-y\,|N|^2
 \in\mathbb Q[x,y,c,t,a,\beta,p].
 \label{eq:f2v54-E}
\end{equation}
Then \(E\) has \(5301\) monomials and degrees \((8,9,4,2,4,2,9)\) in
\((x,y,c,t,a,\beta,p)\), and for every real \((c,t,a,\beta)\) and
integer \(p\ge5\):
\begin{equation}
 \boxed{\;y>0\ \text{and}\ E(x,y,c,t,a,\beta,p)<0\ \Longrightarrow\ F_G(x+iy)\ne0.\;}
 \label{eq:f2v54-criterion}
\end{equation}
\end{proposition}

\begin{proof}
By Lemma~\ref{lem:f2v53-leibniz}, \(F_G=A-iB\) with \(A,B\) the linear
combinations \eqref{eq:f2v53-AB} of \(h_0,\ldots,h_3\), and by
Lemma~\ref{lem:f2v53-generating} each \(h_k\) is a combination of
\(\widetilde G^{(l)}(D)\mathrm{He}_{p-k+l}\), \(l\le2\), with
\(\widetilde G(D)=D^2-2aD+(a^2+\beta)\), \(\widetilde G'(D)=2D-2a\),
\(\widetilde G''(D)=2\).  Using \(D\mathrm{He}_r=r\mathrm{He}_{r-1}\),
everything is a combination of \(\mathrm{He}_{p-k}\), \(0\le k\le5\),
with coefficients polynomial in \((c,t,a,\beta,p)\); the upward
recurrence expresses each \(\mathrm{He}_{p-k}\), \(k\le3\), through
\(v\) and \(u\) with coefficients polynomial in \(\zeta\) and \(p\).  This
gives \eqref{eq:f2v54-FG}, the sign \((-1)^p\) coming from
\eqref{eq:f2v53-hk}.

Let \(y>0\) and \(F_G(\zeta)=0\).  Since \(\mathrm{He}_{p-4}\) is
real-rooted, \(v\ne0\).  If \(D(\zeta)=0\) then \(N(\zeta)=0\) and
\(E(\zeta)=0\).  Otherwise \(\mu:=u/v=N/D\), and with
\(\mathfrak m:=\mathrm{He}_{p-4}'/\mathrm{He}_{p-4}=(p-4)\mu\) the Pick
inequality of Lemma~\ref{lem:f2v5-pick} for the real-rooted polynomial
\(\mathrm{He}_{p-4}\) reads \(-(p-4)\operatorname{Im}\mathfrak m
 -y|\mathfrak m|^2\ge0\), that is \(-\operatorname{Im}\mu-y|\mu|^2\ge0\);
multiplying by \(|D|^2>0\) gives \(E\ge0\).  In both cases \(E\ge0\),
which is \eqref{eq:f2v54-criterion}.  The monomial count and degrees are
read off from the exact construction.
\end{proof}

\begin{remark}[Exact verifier]
\label{rem:f2v54-verifier-E}
\path{scripts/verify_grh_v54_m2_pick_reduction.py} rebuilds
\(A_v,A_u,B_v,B_u\) and \(E\) over \(\mathbb Q\), asserts the monomial
count and degrees, compares \eqref{eq:f2v54-FG} with a direct exact
evaluation of \(Q(0)-iQ'(0)\) at eight random rational points
\((x,y,c,t,a,t_2,p)\) with \(5\le p\le9\), evaluates \(E\) at three
rational points of the target region below and at one point outside it.
It has \(227\) lines, \(7408\) bytes, SHA--256
\begin{center}
\texttt{DA2DAEC5FD03BF3377EEF47DEBB305EB}\\
\texttt{21FA6FE1CE8DDF608494A3BA590FB860},
\end{center}
and terminates with \texttt{VERIFIED}.
\end{remark}

\begin{theorem}[Conditional band closure for two pairs]
\label{thm:f2v54-conditional}
Suppose that for every integer \(p\ge5\)
\begin{equation}
 E(x,y,c,t,a,\beta,p)<0\quad\text{whenever}\quad
 y>0,\ \ t>p+4,\ \ \beta\ge t-1,\ \ x,c,a\in\mathbb R,
 \label{eq:f2v54-target}
\end{equation}
and that \(F_G\) has no real zero \(\zeta\) under the same constraints.
Then every block with two conjugate pairs and \(p\ge5\) real padding
zeros whose critical heat envelope is real-rooted satisfies
\(y_{\min}^2\le(p+4)K=nK\), and misses the open gamma band.
\end{theorem}

\begin{proof}
By Theorem~\ref{thm:f2v53-reduction}, real-rootedness at time \(K\)
gives \(\zeta\) in the closed upper half-plane with \(F_G(\zeta)=0\),
where \(t=y_{\min}^2/T\ge y_{\min}^2/K\) and \(\beta=t_2-1\ge t-1\).  If
\(t>p+4\), \eqref{eq:f2v54-criterion} and \eqref{eq:f2v54-target}
exclude \(\operatorname{Im}\zeta>0\), and the real-axis hypothesis
excludes \(\operatorname{Im}\zeta=0\).  Hence \(t\le p+4\).  The ladder
statement is Corollary~\ref{cor:f2v5-gamma} with \(p+2\) replaced by
\(n=2m+p\) and \(R\ge n\).
\end{proof}

\begin{remark}[Evidence for \eqref{eq:f2v54-target}]
\label{rem:f2v54-evidence}
At \(20000\) random floating-point points of the region
\eqref{eq:f2v54-target} with \(5\le p\le14\), \(E\) was negative at
every one; at \(20000\) points with \(t<p+4\) it was positive at
\(86\).  The exact point
\((x,y,c,t,a,\beta,p)=(-11/8,1/64,-1/2,25/4,25/4,9,5)\), with
\(t=25/4<9\), has \(E>0\), so the constraint \(t>p+4\) is not
removable from \eqref{eq:f2v54-target} in general, while the reals-only
thresholds of Remark~\ref{rem:f2v53-numerics} suggest that
\eqref{eq:f2v54-target} in fact holds from about \(t>p+5/2\).  These are
observations, not proofs.
\end{remark}

\subsection*{Firewall and V55 direction}
\label{fire:f2v54}
Theorem~\ref{thm:f2v54-jacobi} and Corollary~\ref{cor:f2v54-m2} are
unconditional; Theorem~\ref{thm:f2v54-conditional} is conditional on
one explicit polynomial inequality and one real-axis exclusion.  Nothing
here proves the band closure for \(m=2\), nothing addresses \(m\ge3\),
and nothing proves GRH; GRH remains open.

The inequality \eqref{eq:f2v54-target} is a finite real-algebraic
statement in seven variables with \(p\) entering polynomially, so it is a
candidate for an exact certificate, but its region is unbounded and the
degree in \(x\) and \(y\) is eight and nine.  V55 should first lower the
degree: in V5 the coefficients had degree one because the Hermite
differential equation was used at level \(p\); the analogous device here
is to reduce to the basis \((\mathrm{He}_p,\mathrm{He}_{p-1})\) and to
divide out the positive factors \((p-1)\cdots(p-4)\), which cannot raise
the degree in \(\zeta\) above four but may expose a completed-square
structure like \textup{(N5.12)}.  If no such structure appears, the
next option is a Positivstellensatz certificate on a compactified chart.
The cases \(p\le4\), where the basis \((\mathrm{He}_{p-4},\mathrm{He}_{p-5})\)
degenerates, are finite and should be handled directly.  For \(m\ge3\),
Remark~\ref{rem:f2v54-several} states the coupling-sign problem
precisely.

\subsection*{Assumption ledger}
\label{ass:f2v54}
The proved statements use the Hermite three-term recurrence, the
spectral theorem and Cauchy interlacing for real symmetric tridiagonal
matrices, Lemmas~\ref{lem:f2v53-leibniz}--\ref{lem:f2v53-generating},
the Pick inequality \textup{(N5.5)} for a real-rooted polynomial, and
exact rational arithmetic in the two verifiers.  Sampled signs of \(E\)
are exploration only.
\endgroup

\section*{Version V54 append-only record}
\addcontentsline{toc}{section}{Version V54 append-only record}

V54 was appended byte-for-byte before the terminal marker of validated
V53, whose installed SHA--256 was
\[
 \texttt{09E1840B82BA8433F407444217C9289E}
 \allowbreak\texttt{EFA00E245F9F952F8D36A791E1092C57}
\]
with \(15814\) lines and \(517708\) bytes.  It proves the single-factor
Jacobi theorem, settles the V53 real-rootedness problem for one
companion pair, and reduces the two-pair band closure to one explicit
polynomial inequality.  After validation V54 replaces V53 as the sole
active numeric GRH note; the frozen \texttt{V100\_f2} archive is
unchanged.

% =============================================================================
% END APPENDED V54 MATERIAL
% =============================================================================




\clearpage
\section{V55: the balanced two-pair Pick expression}
\label{sec:f2v55-balanced}
\begingroup
\setcounter{equation}{0}
\renewcommand{\theequation}{N55.\arabic{equation}}
\renewcommand{\theHequation}{f2v55.\arabic{equation}}

V54 left the two-pair band closure as the inequality
\eqref{eq:f2v54-target} for a polynomial of degree eight in \(x\) and nine
in \(y\).  The degree came from the choice of the Pick level: the basis
\((\mathrm{He}_{p-4},\mathrm{He}_{p-5})\) sits at the bottom of the range
\(\mathrm{He}_p,\ldots,\mathrm{He}_{p-5}\) that the functional occupies,
so the top members cost four powers of \(\zeta\) each.  Placing the Pick
level in the middle halves every degree.  This section records that
balanced reduction exactly, proves the sign of its leading coefficient,
and states the resulting target.

\subsection{Reduction at the middle level}

Keep the notation of Proposition~\ref{prop:f2v54-reduction} but put
\begin{equation}
 v:=\mathrm{He}_{p-2}(\zeta),\qquad u:=\mathrm{He}_{p-3}(\zeta),\qquad
 M:=(p-3)(p-4)>0\quad(p\ge5).
 \label{eq:f2v55-basis}
\end{equation}
The recurrence gives \(\mathrm{He}_{p-1}=\zeta v-(p-2)u\),
\(\mathrm{He}_p=(\zeta^2-p+1)v-(p-2)\zeta u\),
\(\mathrm{He}_{p-4}=(\zeta u-v)/(p-3)\), and
\(\mathrm{He}_{p-5}=\bigl(-\zeta v+(\zeta^2-p+3)u\bigr)/M\); every
coefficient has degree at most two in \(\zeta\).

\begin{proposition}[Balanced reduction]
\label{prop:f2v55-reduction}
There are polynomials \(A_v,A_u,B_v,B_u\in\mathbb Q[x,y,c,t,a,\beta,p]\)
of degree at most two in \(x\) and in \(y\) such that
\begin{equation}
 M\,F_G(\zeta)=(-1)^p\Bigl[(A_v-iB_v)\,v+(A_u-iB_u)\,u\Bigr].
 \label{eq:f2v55-FG}
\end{equation}
With \(N:=-(A_v-iB_v)\), \(D:=A_u-iB_u\), and
\(E:=-\operatorname{Im}(N\overline D)-y|N|^2\), the polynomial \(E\) has
\(4001\) monomials and degrees \((4,5,4,2,4,2,11)\) in
\((x,y,c,t,a,\beta,p)\), and for all real \((c,t,a,\beta)\) and integer
\(p\ge5\),
\begin{equation}
 \boxed{\;y>0\ \text{and}\ E(x,y,c,t,a,\beta,p)<0\ \Longrightarrow\ F_G(x+iy)\ne0.\;}
 \label{eq:f2v55-criterion}
\end{equation}
\end{proposition}

\begin{proof}
Identical to the proof of Proposition~\ref{prop:f2v54-reduction}, with
the basis \eqref{eq:f2v55-basis}, the multiplier \(M\) clearing the
denominators of \(\mathrm{He}_{p-4},\mathrm{He}_{p-5}\), and the Pick
inequality \textup{(N5.5)} applied to the real-rooted polynomial
\(\mathrm{He}_{p-2}\) with \(\mathfrak m=(p-2)\,u/v\).  Since \(M>0\),
the zeros of \(MF_G\) and \(F_G\) coincide.
\end{proof}

\begin{lemma}[Leading coefficient]
\label{lem:f2v55-leading}
The coefficient of \(x^4y^0\) in \(E\) is
\begin{equation}
 -M^2\,p(p-1)(p-2)\,\Bigl[(c^2+t-1)(a^2+\beta)-4ca+2\Bigr],
 \label{eq:f2v55-top}
\end{equation}
and the coefficient of \(x^4\) in \(E\) has degree one in \(y\).  On the
region \(t\ge9\), \(\beta\ge8\) the bracket is at least \(66\); hence
the \(x^4y^0\) coefficient is negative there for every \(p\ge5\).
\end{lemma}

\begin{proof}
The identity \eqref{eq:f2v55-top} and the degree statement are exact
consequences of the construction and are asserted by the verifier.  For
the bound, \((c^2+t-1)(a^2+\beta)\ge(c^2+8)(a^2+8)
 =c^2a^2+8c^2+8a^2+64\ge4|ca|+64\) because \(8c^2+8a^2\ge16|ca|\); so the
bracket is at least \(64+4|ca|-4ca+2\ge66\).
\end{proof}

\begin{remark}[Exact verifier]
\label{rem:f2v55-verifier}
\path{scripts/verify_grh_v55_balanced_pick.py} constructs
\(A_v,A_u,B_v,B_u,E\) over \(\mathbb Q\), asserts the monomial count,
the degrees, the degree bound two on \(N\) and \(D\) in \(x\) and \(y\),
the identity \eqref{eq:f2v55-top}, compares \eqref{eq:f2v55-FG} with a
direct exact evaluation of \(M(Q(0)-iQ'(0))\) at eight random rational
points with \(5\le p\le9\), and evaluates \(E\) at three rational points
of the target region.  It has \(235\) lines, \(7242\) bytes, SHA--256
\begin{center}
\texttt{124FE8FBC6F32F53BEC5B2020C4039D8}\\
\texttt{6DD709014FCD82CDDF685ED86BD5B25E},
\end{center}
and terminates with \texttt{VERIFIED}.
\end{remark}

\subsection{The target and its numerical margin}

\begin{problem}[Balanced two-pair inequality]
\label{prob:f2v55-target}
Prove, for every integer \(p\ge5\),
\begin{equation}
 E(x,y,c,t,a,\beta,p)<0\quad\text{whenever}\quad
 y>0,\ \ t>p+4,\ \ \beta\ge t-1,\ \ x,c,a\in\mathbb R,
 \label{eq:f2v55-target}
\end{equation}
together with the real-axis exclusion of
Theorem~\ref{thm:f2v54-conditional}.  By \eqref{eq:f2v55-criterion} and
Theorem~\ref{thm:f2v53-reduction} this closes the gamma band for every
block with two conjugate pairs and \(p\ge5\) real padding zeros.
\end{problem}

\begin{remark}[Margin, numerically]
\label{rem:f2v55-numerics}
For \(p=5\), a Nelder--Mead maximization of \(E\) over the region with
\(t\) fixed gave maxima \(2.9\cdot10^{12}\), \(6.4\cdot10^2\),
\(-6.0\cdot10^3\), \(-1.3\cdot10^4\), \(-2.6\cdot10^4\), \(-6.8\cdot10^4\)
at \(t=7,7.5,8,8.5,9,10\); the sign change of the maximum lies between
\(t=7.5\) and \(t=8\), below the required \(t>9\), in agreement with the
reals-only threshold \(t^*\approx p+5/2\) of
Remark~\ref{rem:f2v53-numerics}.  At \(20000\) random points of the
region with \(5\le p\le14\), \(E<0\) throughout.  These are
observations, not proofs.
\end{remark}

\subsection*{Firewall and V56 direction}
\label{fire:f2v55}
Proposition~\ref{prop:f2v55-reduction} and Lemma~\ref{lem:f2v55-leading}
are proved; Problem~\ref{prob:f2v55-target} is open, the band for
\(m=2\) is not closed, \(m\ge3\) is untouched, and GRH remains open.

With \(N\) and \(D\) quadratic in \(\zeta\), \(E\) is a quartic in \(x\)
whose leading coefficient is now known to be negative on the region.  The
natural next step is to write \(E=-\bigl(Q_1(x,y)^2\,W_1+\cdots\bigr)\)
with \(W_j\) nonnegative on the region, following the completed-square
pattern of \textup{(N5.12)}, using the substitutions \(X=x-c\),
\(t=p+4+\sigma\), \(\beta=t-1+\rho\) with \(\sigma,\rho\ge0\) so that
all constraints become sign conditions; a Bernstein certificate on
dyadic charts in the bounded variables, combined with the leading-term
control in the unbounded ones, is the fallback.  The cases \(p\le4\) are
finite.

\subsection*{Assumption ledger}
\label{ass:f2v55}
The proved statements use the Hermite recurrence, the Pick inequality
\textup{(N5.5)}, Lemmas~\ref{lem:f2v53-leibniz}--\ref{lem:f2v53-generating},
elementary inequalities, and exact rational arithmetic in the verifier.
The maxima and sampled signs are exploration only.
\endgroup

\section*{Version V55 append-only record}
\addcontentsline{toc}{section}{Version V55 append-only record}

V55 was appended byte-for-byte before the terminal marker of validated
V54, whose installed SHA--256 was
\[
 \texttt{39D8CD0CB9B5A4124DDA566AE577CDE7}
 \allowbreak\texttt{130E515E5A5BEE77880789533344A912}
\]
with \(16088\) lines and \(529897\) bytes.  It halves the degree of the
two-pair Pick expression and proves the sign of its leading coefficient
on the target region.  After validation V55 replaces V54 as the sole
active numeric GRH note; the frozen \texttt{V100\_f2} archive is
unchanged.

% =============================================================================
% END APPENDED V55 MATERIAL
% =============================================================================




\clearpage
\section{V56: the Pick disk, the \texorpdfstring{\(y\)}{y}-expansion, and two exact coefficient certificates}
\label{sec:f2v56-pick-disk}
\begingroup
\setcounter{equation}{0}
\renewcommand{\theequation}{N56.\arabic{equation}}
\renewcommand{\theHequation}{f2v56.\arabic{equation}}

The balanced target \eqref{eq:f2v55-target} asks for \(E<0\).  This
section replaces it by a cleaner sufficient condition with a geometric
meaning, expands that condition in powers of \(y\), factors the
\(p\)-dependence of every coefficient exactly, and proves two of the five
coefficients positive on the region by exact absorption certificates.

\subsection{The Pick disk}

\begin{lemma}[Pick disk]
\label{lem:f2v56-disk}
For \(y>0\) the inequality \(-\operatorname{Im}\mu-y|\mu|^2\ge0\) says
that \(\mu\) lies in the closed disk
\begin{equation}
 \mathbb D_y:=\Bigl\{\mu:\ \bigl|\mu+\tfrac{i}{2y}\bigr|\le\tfrac1{2y}\Bigr\},
 \label{eq:f2v56-disk}
\end{equation}
which is contained in the closed lower half-plane and touches the real
axis only at \(0\).  Consequently, in the notation of
Proposition~\ref{prop:f2v55-reduction},
\begin{equation}
 \boxed{\;\operatorname{Im}\bigl(N\overline D\bigr)>0
 \ \Longrightarrow\ E<0\ \Longrightarrow\ F_G(x+iy)\ne0\qquad(y>0).\;}
 \label{eq:f2v56-sufficient}
\end{equation}
\end{lemma}

\begin{proof}
\(|\mu+i/(2y)|^2-1/(4y^2)=|\mu|^2+\operatorname{Im}\mu/y\), so the
Pick inequality is \(|\mu+i/(2y)|\le1/(2y)\).  If
\(\operatorname{Im}(N\overline D)>0\) then
\(E=-\operatorname{Im}(N\overline D)-y|N|^2<0\), and
\eqref{eq:f2v55-criterion} applies.
\end{proof}

Geometrically: at a collision, \(\mu=\mathrm{He}_{p-3}(\zeta)/\mathrm{He}_{p-2}(\zeta)\)
must lie in \(\mathbb D_y\), while \eqref{eq:f2v56-sufficient} asks that
the algebraic ratio \(N/D\) determined by the block lie in the open
upper half-plane, which is disjoint from \(\mathbb D_y\).

\subsection{Expansion in \texorpdfstring{\(y\)}{y} and exact factorizations}

\begin{proposition}[Structure of \(\operatorname{Im}(N\overline D)\)]
\label{prop:f2v56-structure}
\(\operatorname{Im}(N\overline D)\) is a polynomial with \(2656\)
monomials, of degree four in \(x\) and in \(y\), and
\begin{equation}
 \operatorname{Im}(N\overline D)=\sum_{j=0}^{4}y^jW_j,
 \qquad
 W_j=(p-3)^2(p-4)^2\,\Pi_j\,\widetilde W_j,
 \label{eq:f2v56-expansion}
\end{equation}
with
\begin{equation}
 \Pi_0=\Pi_4=p(p-1)(p-2),\qquad \Pi_1=\Pi_3=p-2,\qquad \Pi_2=p(p-2),
 \label{eq:f2v56-Pi}
\end{equation}
where the \(\widetilde W_j\) are polynomials in \((x,c,t,a,\beta,p)\)
with \(232\), \(303\), \(142\), \(73\), \(8\) monomials and degrees
\(4,2,2,0,0\) in \(x\).  Explicitly
\begin{equation}
 \widetilde W_4=(c^2+t-1)(a^2+\beta)-4ca+2 .
 \label{eq:f2v56-W4}
\end{equation}
\end{proposition}

\begin{proof}
Exact polynomial division by the listed linear factors in \(p\), with
zero remainder asserted at each step, performed by the verifier.
\end{proof}

Since \((p-3)^2(p-4)^2\Pi_j>0\) for \(p\ge5\), positivity of
\(\operatorname{Im}(N\overline D)\) on the region follows from
\(\widetilde W_j\ge0\) for all \(j\) with strict inequality for one of
them.  Substitute the region:
\begin{equation}
 p=5+\pi,\qquad t=9+\pi+\sigma,\qquad \beta=8+\pi+\sigma+\rho,
 \qquad \pi,\sigma,\rho\ge0,
 \label{eq:f2v56-sub}
\end{equation}
so that \(t>p+4\) and \(\beta\ge t-1\) become \(\sigma>0\), \(\rho\ge0\);
denote the substituted quotients \(\widetilde W_j(x,c,a;\sigma,\rho,\pi)\).

\subsection{Two exact absorption certificates}

Call a monomial \emph{even} if \(x,c,a\) occur to even powers; even
monomials with positive coefficients are nonnegative on the region and
form the \emph{budget}.  Every other monomial \(m\) with coefficient
\(\kappa\) is \emph{absorbed} by the arithmetic--geometric mean
inequality
\begin{equation}
 |\kappa\,m|\le\frac{|\kappa|}2\Bigl(\lambda m_1+\frac{m_2}\lambda\Bigr),
 \qquad m_1m_2=m^2,\quad\lambda>0,
 \label{eq:f2v56-amgm}
\end{equation}
where \(m_1,m_2\) are budget monomials; several such rows with weights
\(\theta\) summing to at least one may be used for one \(m\).

\begin{theorem}[\(\widetilde W_3\) and \(\widetilde W_4\) are positive on the region]
\label{thm:f2v56-W34}
After the substitution \eqref{eq:f2v56-sub}, \(\widetilde W_4\) has
\(18\) monomials, one of which (\(-4ca\)) is absorbed, and
\(\widetilde W_3\) has \(125\) monomials, \(16\) of which are absorbed
by \(31\) rows of the form \eqref{eq:f2v56-amgm}; in both cases the total
consumption of every budget monomial stays within its coefficient, and
the unconsumed constant terms are \(34\) and \(2836\).  Hence, on the
region \eqref{eq:f2v56-sub},
\begin{equation}
 \widetilde W_4\ge34,\qquad \widetilde W_3\ge2836,
 \label{eq:f2v56-bounds}
\end{equation}
so \(W_3,W_4>0\) for every \(p\ge5\).
\end{theorem}

\begin{proof}
The rows are stored in the verifier and checked exactly: each row uses
budget monomials \(m_1,m_2\) with \(m_1m_2=m^2\), the weights of each
absorbed monomial sum to at least one, and the consumption of every
budget monomial is at most its coefficient.  Summing
\eqref{eq:f2v56-amgm} over the rows and subtracting from the polynomial
leaves a polynomial with nonnegative coefficients in the even monomials
and the stated constant term.  For \(\widetilde W_4\) the single row is
\(4|ca|\le2c^2+2a^2\), consistent with Lemma~\ref{lem:f2v55-leading}.
\end{proof}

\begin{remark}[Exact verifier]
\label{rem:f2v56-verifier}
\path{scripts/verify_grh_v56_imND_certificates.py} rebuilds
\(N,D\) from the balanced construction of V55, forms
\(\operatorname{Im}(N\overline D)\), asserts the monomial count, the
\(y\)-expansion, the exact divisibility by the factors in
\eqref{eq:f2v56-expansion}--\eqref{eq:f2v56-Pi}, the identity
\eqref{eq:f2v56-W4}, the monomial counts after
\eqref{eq:f2v56-sub}, and the two embedded certificates.  It has
\(236\) lines, \(10433\) bytes, SHA--256
\begin{center}
\texttt{A7642FED26C30FCFF42EDB782B1DCC06}\\
\texttt{1950EDE04ECB4928A102695E0BAEB434},
\end{center}
and terminates with \texttt{VERIFIED}.
\end{remark}

\subsection{What remains for two pairs}

\begin{theorem}[Reduction to three coefficient inequalities]
\label{thm:f2v56-reduction}
Suppose that, after \eqref{eq:f2v56-sub},
\(\widetilde W_0,\widetilde W_1,\widetilde W_2\ge0\) for all real
\(x,c,a\) and all \(\sigma>0\), \(\rho,\pi\ge0\), and that \(F_G\) has no
real zero on the region.  Then the gamma band is closed for every block
with two conjugate pairs and \(p\ge5\) real padding zeros.
\end{theorem}

\begin{proof}
By Theorem~\ref{thm:f2v56-W34} and the hypothesis,
\(\operatorname{Im}(N\overline D)\ge y^3W_3>0\) for \(y>0\); Lemma
\ref{lem:f2v56-disk} gives \(F_G\ne0\) in the open upper half-plane, and
Theorem~\ref{thm:f2v54-conditional} applies.
\end{proof}

\begin{remark}[Status of the three coefficients]
\label{rem:f2v56-status}
Numerically each of \(\widetilde W_0,\widetilde W_1,\widetilde W_2\) is
positive on the region, with minima about \(5.2\cdot10^4\),
\(1.9\cdot10^5\), \(1.7\cdot10^5\) attained at \(\sigma=\rho=\pi=0\).
Pairwise absorption \eqref{eq:f2v56-amgm} alone is insufficient for
them: the corresponding linear programmes are infeasible in the original
coordinates and after the shifts \(x\mapsto x+c\), \(a\mapsto a+c\),
\(x\mapsto x+a\), \(x\mapsto x+(c+a)/2\); for \(\widetilde W_0\) the
even monomials \(-36x^2(1+\pi+\pi^2/9)\) must be absorbed by \(x^4\)
and constant terms jointly.  Since \(\widetilde W_1,\widetilde W_2\) are
quadratic in \(x\), positivity of their \(x^2\)- and \(x^0\)-coefficients
together with nonnegativity of the discriminant suffices; those three
polynomials are the next certificate targets.
\end{remark}

\subsection*{Firewall and V57 direction}
\label{fire:f2v56}
Lemma~\ref{lem:f2v56-disk}, Proposition~\ref{prop:f2v56-structure}, and
Theorem~\ref{thm:f2v56-W34} are proved.  The band is not yet closed for
two pairs, \(m\ge3\) is untouched, and GRH remains open.

V57 should certify \(\widetilde W_2\) and \(\widetilde W_1\) through
their \(x\)-discriminants (each a polynomial in \(c,a,\sigma,\rho,\pi\)
with the two signed variables \(c,a\)), and \(\widetilde W_0\) through
a quartic-in-\(x\) argument, using either absorption with binomial
squares \((\alpha m_1-m_2)^2\) as additional budget or an exact
sum-of-squares search restricted to the signed variables.  The real-axis
exclusion and the cases \(p\le4\) remain to be written.

\subsection*{Assumption ledger}
\label{ass:f2v56}
The proved statements use elementary geometry of the Pick disk, exact
polynomial division and substitution, the arithmetic--geometric mean
inequality, and the exact verifier.  Minima and infeasibility reports
are exploration only.
\endgroup

\section*{Version V56 append-only record}
\addcontentsline{toc}{section}{Version V56 append-only record}

V56 was appended byte-for-byte before the terminal marker of validated
V55, whose installed SHA--256 was
\[
 \texttt{D9A771A252139954C008CC1CFDBFC9B8}
 \allowbreak\texttt{7FE4A98B6F053835190C2B8AE955BEA9}
\]
with \(16260\) lines and \(537152\) bytes.  It replaces the balanced
target by the Pick-disk condition, factors its \(y\)-coefficients
exactly, and certifies two of the five.  After validation V56 replaces
V55 as the sole active numeric GRH note; the frozen \texttt{V100\_f2}
archive is unchanged.

% =============================================================================
% END APPENDED V56 MATERIAL
% =============================================================================




\clearpage
\section{V57: grouped sum-of-squares certificates and the nine tight groups}
\label{sec:f2v57-groups}
\begingroup
\setcounter{equation}{0}
\renewcommand{\theequation}{N57.\arabic{equation}}
\renewcommand{\theHequation}{f2v57.\arabic{equation}}

V56 reduced the two-pair band to the nonnegativity of
\(\widetilde W_0,\widetilde W_1,\widetilde W_2\) on the region
\eqref{eq:f2v56-sub}.  Pairwise absorption fails for them.  This section
uses genuine sums of squares, organized so that the nonnegative region
variables enter only as multipliers, certifies \(27\) of the \(36\)
resulting groups exactly, and identifies the algebraic reason why the
remaining nine resist: their leading forms are degenerate squares.

\subsection{Groups}

Write \(\pi=w^2\) with \(w\in\mathbb R\) and, for \(j\in\{0,1,2\}\) and
exponents \(s,r\ge0\),
\begin{equation}
 R_{j,s,r}(x,c,a,w):=\sum_{k\ge0}\bigl[\sigma^s\rho^r\pi^k\bigr]\widetilde W_j\;w^{2k},
 \qquad
 \widetilde W_j=\sum_{s,r}\sigma^s\rho^r\,R_{j,s,r}(x,c,a,\sqrt\pi).
 \label{eq:f2v57-groups}
\end{equation}
Since \(\sigma,\rho\ge0\), it suffices that every \(R_{j,s,r}\) be
nonnegative on \(\mathbb R^4\), with the group \((s,r)=(0,0)\) strictly
positive to keep \(\widetilde W_j>0\).  There are \(12\) groups for each
\(j\), namely \((s,r)\) with \(s+r\le4\), \(r\le2\), \(s\le4\), as listed
by the verifier.

\begin{theorem}[Twenty-seven exact group certificates]
\label{thm:f2v57-groups}
For the following groups there are a monomial basis \(m\), an exact
rational positive semidefinite matrix \(G\), and absorption rows as in
\eqref{eq:f2v56-amgm} such that
\begin{equation}
 R_{j,s,r}=m^{\mathsf T}Gm+\text{(absorbed monomials)}+\text{(even monomials with nonnegative coefficients)},
 \label{eq:f2v57-cert}
\end{equation}
hence \(R_{j,s,r}\ge0\) on \(\mathbb R^4\):
\[
\begin{array}{c|l}
 j&\text{certified }(s,r)\\ \hline
 2&(0,1),(0,2),(1,0),(1,1),(1,2),(2,0),(2,1),(2,2),(3,0),(3,1),(4,0)\\
 1&(0,2),(1,1),(1,2),(2,0),(2,1),(2,2),(3,0),(3,1),(4,0)\\
 0&(1,1),(1,2),(2,1),(2,2),(3,0),(3,1),(4,0)
\end{array}
\]
The basis sizes range from \(1\) to \(32\); in every case \(G\) is
verified positive semidefinite by an exact rational
\(LDL^{\mathsf T}\) elimination, and the residual after subtracting
\(m^{\mathsf T}Gm\) is absorbed exactly.
\end{theorem}

\begin{proof}
The matrices were found numerically by an interior-point solver on the
semidefinite programme "maximize the least eigenvalue of \(G\) subject to
\(m^{\mathsf T}Gm=R-\tfrac18(\text{even positive part of }R)\)", with
\(m\) the lattice points of half the Newton polytope of \(R\), then
rounded to rationals with denominator \(2^{30}\); the reserved eighth of
the even part absorbs the rounding residual through
\eqref{eq:f2v56-amgm}.  Only the exact verification matters for the
proof: it is the content of the verifier below.
\end{proof}

\begin{remark}[Exact verifier and data]
\label{rem:f2v57-verifier}
\path{scripts/verify_grh_v57_group_certificates.py} rebuilds the
balanced construction, \(\operatorname{Im}(N\overline D)\), the
\(y\)-expansion, the region substitution and the groups
\eqref{eq:f2v57-groups} from scratch, then reads the certificates from
\path{scripts/data_grh_v57_group_certificates.json} (\(42897\) bytes,
SHA--256
\texttt{4FDB62BBD33C84549727B115433C06F9}\allowbreak
\texttt{89F6C4CF8B687A7D41E57649D43C5982})
and verifies each one exactly; it also asserts the leading-form
identities of Lemma~\ref{lem:f2v57-leading}.  The script has \(306\)
lines, \(10423\) bytes, SHA--256
\begin{center}
\texttt{7381DE58A74A55B71C9C54061F5EC4EF}\\
\texttt{DD09F54D7B65795C16C2A7F2234EBBE7},
\end{center}
and terminates with \texttt{VERIFIED}.
\end{remark}

\subsection{Why nine groups resist}

The open groups are \((j,s,r)=(2,0,0)\); \((1,0,0),(1,0,1),(1,1,0)\);
\((0,0,0),(0,0,1),(0,0,2),(0,1,0),(0,2,0)\).  For each of them the
semidefinite programme is feasible only with least eigenvalue zero, and
with any positive slack it is infeasible, so no rounding can succeed.
The reason is structural.

\begin{lemma}[Leading forms of the open groups]
\label{lem:f2v57-leading}
Let \(L_{j,s,r}\) denote the homogeneous part of highest total degree in
\((x,c,a)\) of the coefficient of the highest power of \(w\) in
\(R_{j,s,r}\).  Then
\begin{equation}
\begin{aligned}
 L_{0,0,0}&=\bigl((x-2a)(x-2c)\bigr)^2,\qquad
 L_{0,0,2}=L_{1,0,1}=(x-2c)^2,\\
 L_{2,0,0}&=L_{1,1,0}=L_{0,2,0}=(x-2a)^2+(x-2c)^2 .
\end{aligned}
\label{eq:f2v57-leading}
\end{equation}
Each of these forms vanishes on a nontrivial real cone (\(x=2c\), or
\(x=2a=2c\)), so every Gram matrix representing the group is singular
and the group is nonnegative but not strictly positive at infinity.
\end{lemma}

\begin{proof}
Exact identities, asserted by the verifier after extracting the leading
forms by degree.  A form vanishing at a real point forces every
sum-of-squares representation to have all squares vanish there, so the
Gram matrix has that direction in its kernel.
\end{proof}

The remaining three open groups, \((0,0,1)\), \((0,1,0)\), \((1,0,0)\),
have quartic or sextic leading forms that are likewise semidefinite with
zeros along \(x=2c\) or \(x=2a\).  The combinations \(x-2c\) and
\(x-2a\) are the ones that appear in the V5 diagonal
\textup{(N5.10)}; the degeneracy is inherited from the structure of the
collision functional, not an artefact of the reduction.

\begin{corollary}[Reduced residual problem]
\label{cor:f2v57-residual}
The band closure for two pairs (Theorem~\ref{thm:f2v56-reduction})
follows from the nonnegativity of the nine groups listed above, with
strict positivity of \(R_{j,0,0}\) for \(j=0,1,2\).
\end{corollary}

\subsection*{Firewall and V58 direction}
\label{fire:f2v57}
Theorem~\ref{thm:f2v57-groups} and Lemma~\ref{lem:f2v57-leading} are
proved.  Nine groups remain, the two-pair band is not closed, \(m\ge3\)
is untouched, and GRH remains open.

The degenerate leading forms are known exactly, so V58 should perform
exact facial reduction: write each open group as
\(L\cdot Q+\text{(lower order)}\) with \(L\) the square from
\eqref{eq:f2v57-leading}, restrict the Gram basis to monomials
compatible with the kernel directions \(x=2c\), \(x=2a\), and re-solve;
equivalently, substitute \(x=2c+u\) or \(x=2a+v\) in the relevant
group so that the degenerate square becomes a monomial \(u^2\) or
\(u^2v^2\) and slack can be reserved on the remaining monomials.  If a
group has an actual real zero (not only at infinity), it must be split
along that zero exactly.  The real-axis and \(p\le4\) cases remain
as stated in V56.

\subsection*{Assumption ledger}
\label{ass:f2v57}
Only exact rational arithmetic enters the proofs: the constructions of
V55--V56, exact \(LDL^{\mathsf T}\) positivity tests, and the absorption
inequality.  The semidefinite solver is a search heuristic and none of
its floating-point output is used without exact re-verification.
\endgroup

\section*{Version V57 append-only record}
\addcontentsline{toc}{section}{Version V57 append-only record}

V57 was appended byte-for-byte before the terminal marker of validated
V56, whose installed SHA--256 was
\[
 \texttt{33DB47D83CE375E8D3A3E234C6A537D1}
 \allowbreak\texttt{E89F0E5486F2FFAF0DDFB638ED80E2A4}
\]
with \(16487\) lines and \(546654\) bytes.  It certifies twenty-seven
of the thirty-six groups of the two-pair Pick coefficients by exact
sums of squares and identifies the degenerate leading forms of the
nine that remain.  After validation V57 replaces V56 as the sole active
numeric GRH note; the frozen \texttt{V100\_f2} archive is unchanged.

% =============================================================================
% END APPENDED V57 MATERIAL
% =============================================================================




\clearpage
\section{V58: closure of the gamma band for two-pair blocks with at least five real padding zeros}
\label{sec:f2v58-two-pair-band}
\begingroup
\setcounter{equation}{0}
\renewcommand{\theequation}{N58.\arabic{equation}}
\renewcommand{\theHequation}{f2v58.\arabic{equation}}

The nine groups left open in V57 are now certified by exact facial
reduction.  Together with V56--V57 this proves that the Pick-disk
condition of Lemma~\ref{lem:f2v56-disk} holds on the whole region, that
the reals-only diagonal has no zero in the closed upper half-plane, and
hence that every critical block with two conjugate-pair channels and at
least five real gamma factors misses the open gamma band.

\subsection{Facial reduction of the nine groups}

For each open group \(R=R_{j,s,r}(x,c,a,w)\) of
\eqref{eq:f2v57-groups} the monomial Gram basis was replaced by a
polynomial basis \(b_1,\ldots,b_n\) spanning the orthogonal complement
of the kernel that Lemma~\ref{lem:f2v57-leading} forces; for the corner
group \((0,0,0)\) the four top-layer monomials were replaced by the
single polynomial \((x-2a)(x-2c)w^2\) before further reduction.  On the
reduced basis the semidefinite programme has a strictly positive least
eigenvalue, and rounding becomes exact by solving the coefficient
constraints for a small rational correction of the Gram matrix.

\begin{theorem}[The remaining nine groups]
\label{thm:f2v58-groups}
For each of the groups \((2,0,0)\); \((1,0,0),(1,0,1),(1,1,0)\);
\((0,0,0),(0,0,1),(0,0,2),(0,1,0),(0,2,0)\) there are polynomials
\(b_1,\ldots,b_n\) in \((x,c,a,w)\) and an exact rational positive
semidefinite matrix \(G\) such that
\begin{equation}
 R_{j,s,r}=\sum_{i,k}G_{ik}\,b_ib_k+\text{(absorbed monomials)}
 +\text{(even monomials with nonnegative coefficients)},
 \label{eq:f2v58-cert}
\end{equation}
so \(R_{j,s,r}\ge0\) on \(\mathbb R^4\).  For the three corner groups the
representation is exact with a positive constant remainder:
\begin{equation}
 R_{0,0,0}\ \ge\ 189,\qquad R_{1,0,0}\ \ge\ \frac{4341}{8},\qquad
 R_{2,0,0}\ \ge\ 363 .
 \label{eq:f2v58-corners}
\end{equation}
The basis sizes are \(37,68,41,53,26,16,5,21,13\) in the order listed.
\end{theorem}

\begin{proof}
Exact verification in the verifier below: symmetry and positive
semidefiniteness of \(G\) by rational \(LDL^{\mathsf T}\) elimination,
exact expansion of \(\sum G_{ik}b_ib_k\), and either an absorption
ledger \eqref{eq:f2v56-amgm} for the residual or the identity
\eqref{eq:f2v58-cert} with a manifestly nonnegative remainder.
\end{proof}

\begin{corollary}[All five coefficients]
\label{cor:f2v58-all}
On the region \eqref{eq:f2v56-sub}, for every real \(p\ge5\),
\begin{equation}
 \widetilde W_0\ge189,\qquad \widetilde W_1\ge\frac{4341}8,\qquad
 \widetilde W_2\ge363,\qquad \widetilde W_3\ge2836,\qquad
 \widetilde W_4\ge34 .
 \label{eq:f2v58-W}
\end{equation}
Consequently \(\operatorname{Im}(N\overline D)>0\) for all \(y\ge0\).
\end{corollary}

\begin{proof}
By \eqref{eq:f2v57-groups},
\(\widetilde W_j=R_{j,0,0}+\sum_{(s,r)\ne(0,0)}\sigma^s\rho^rR_{j,s,r}\)
with \(\sigma,\rho\ge0\); Theorems~\ref{thm:f2v57-groups} and
\ref{thm:f2v58-groups} give the first three bounds, and
Theorem~\ref{thm:f2v56-W34} the last two.  Insert into
\eqref{eq:f2v56-expansion}: every factor
\((p-3)^2(p-4)^2\Pi_j\) is positive for \(p\ge5\), and
\(y^j\ge0\), so \(\operatorname{Im}(N\overline D)\ge W_0>0\).
\end{proof}

\subsection{No zero in the closed upper half-plane}

\begin{theorem}[Zero-free diagonal]
\label{thm:f2v58-zero-free}
Let \(p\ge5\) be an integer, \(c,a\in\mathbb R\), \(t\ge p+4\), and
\(t_2\ge t\).  Then the reals-only diagonal functional \(F_G\) of
Theorem~\ref{thm:f2v53-reduction} with \(G(w)=(w-a)^2+t_2\) has no zero
\(\zeta\) with \(\operatorname{Im}\zeta\ge0\).
\end{theorem}

\begin{proof}
The parameters lie in the region \eqref{eq:f2v56-sub} with
\(\sigma=t-p-4\ge0\), \(\rho=t_2-t\ge0\), \(\pi=p-5\ge0\), so
Corollary~\ref{cor:f2v58-all} applies.  Let \(\zeta=x+iy\) with
\(F_G(\zeta)=0\).

If \(y>0\): write \(MF_G=(-1)^p[-N\,v+D\,u]\) with
\(v=\mathrm{He}_{p-2}(\zeta)\ne0\).  If \(D(\zeta)=0\) then
\(N(\zeta)=0\), contradicting \(\operatorname{Im}(N\overline D)>0\).
Otherwise \(\mu=u/v=N/D\) lies in the Pick disk \(\mathbb D_y\) of
Lemma~\ref{lem:f2v56-disk}, which is contained in the closed lower
half-plane, so \(\operatorname{Im}(N\overline D)=|D|^2\operatorname{Im}(N/D)\le0\),
again a contradiction.

If \(y=0\): \(A_v,A_u,B_v,B_u,v,u\) are real, and \(F_G=0\) means
\(A_vv+A_uu=0=B_vv+B_uu\).  Consecutive Hermite polynomials have no
common zero, so \((v,u)\ne(0,0)\) and the determinant
\(A_vB_u-A_uB_v\) vanishes.  But at \(y=0\),
\(\operatorname{Im}(N\overline D)=A_uB_v-A_vB_u=W_0\), which is positive
by Corollary~\ref{cor:f2v58-all}.  Contradiction.
\end{proof}

\subsection{The band theorem for two pairs}

\begin{theorem}[Two-pair band closure, \(p\ge5\)]
\label{thm:f2v58-band}
Let \(P_0\) be a real monic polynomial with exactly two conjugate pairs,
of heights \(y_1\ge y_2>0\), and \(p\ge5\) real zeros, \(n=p+4\).  If
\(e^{-(K/2)D_z^2}P_0\) is real-rooted for some \(K>0\), then
\begin{equation}
 \boxed{\;y_2^2<(p+4)K=nK.\;}
 \label{eq:f2v58-bound}
\end{equation}
The same bound holds for \(p\le2\) by \textup{(N9.14)}, since then
\(m+2p\le2m+p\).
\end{theorem}

\begin{proof}
Theorem~\ref{thm:f2v53-reduction} gives a first collision time
\(T\le K\) and, after normalization, \(t=y_2^2/T\), \(t_2=y_1^2/T\ge t\),
and a zero \(\zeta\) of \(F_G\) in the closed upper half-plane.  By
Theorem~\ref{thm:f2v58-zero-free} this forces \(t<p+4\), that is
\(y_2^2<(p+4)T\le(p+4)K\).
\end{proof}

\begin{corollary}[Gamma ladder]
\label{cor:f2v58-gamma}
In the notation of \textup{(N9.24)}, a critical block with two
conjugate-pair channels, \(p\ge5\) or \(p\le2\) real gamma factors, and
ladder count \(R\ge n=p+4\), whose critical heat envelope is real-rooted,
satisfies
\begin{equation}
 \boxed{\;\kappa^2>\frac{\gamma^2R}{n(1+\theta)}\ \ge\ \frac{\gamma^2}{1+\theta},\;}
 \label{eq:f2v58-gamma}
\end{equation}
where \(\gamma\) is the least ordinate of the two channels.  Such a
block does not enter the open V97 band.
\end{corollary}

\begin{proof}
Insert \(y_2^2=\gamma^2R/(\theta^2\kappa^2)\) and
\(K=(1+\theta)/\theta^2\) into \eqref{eq:f2v58-bound}.
\end{proof}

\begin{remark}[Exact verifier and data]
\label{rem:f2v58-verifier}
\path{scripts/verify_grh_v58_two_pair_band.py} rebuilds every object
from the balanced construction onward, verifies the \(27\) V57 group
certificates and the nine facial certificates of
\path{scripts/data_grh_v58_facial_certificates.json}
(\(132127\) bytes, SHA--256
\texttt{B1CAD166717CADBA26A5DEA0C8B448F3}\allowbreak
\texttt{0FDE2EE4F9699516D845B6F6FC449B4F}), asserts that all \(36\)
groups of \(\widetilde W_0,\widetilde W_1,\widetilde W_2\) are covered
and that the corner constants are exactly those of
\eqref{eq:f2v58-corners}.  It has \(334\) lines, \(11815\) bytes,
SHA--256
\begin{center}
\texttt{5A24650C1C2C3971367ABBDD178899ED}\\
\texttt{0EE991801A4588B40BFFF9EA8A982D0D},
\end{center}
and terminates with \texttt{VERIFIED}.  The semidefinite solver used to
find the certificates is not part of the proof.
\end{remark}

\subsection*{Firewall and V59 direction}
\label{fire:f2v58}
Theorem~\ref{thm:f2v58-band} closes the gamma band for two-pair blocks
with \(p\ge5\) or \(p\le2\) real padding zeros.  The cases \(p=3,4\)
are not covered, because the basis \((\mathrm{He}_{p-2},\mathrm{He}_{p-3})\)
and the multiplier \(M=(p-3)(p-4)\) degenerate there; they are finite
problems in six variables for which the same pipeline applies with
\(p\) fixed.  Blocks with three or more pairs remain open, as does
everything beyond finite blocks.  GRH remains open.

V59 should run the fixed-\(p\) pipeline for \(p=3\) and \(p=4\)
(basis \((\mathrm{He}_1,\mathrm{He}_0)\) and \((\mathrm{He}_2,\mathrm{He}_1)\),
no multiplier), which would complete the two-pair band.  After that the
multi-companion problem of Remark~\ref{rem:f2v54-several} is the
gateway to \(m\ge3\).

\subsection*{Assumption ledger}
\label{ass:f2v58}
The proof uses only exact rational arithmetic (verified identities,
\(LDL^{\mathsf T}\) positivity, absorption inequalities), the Pick
inequality \textup{(N5.5)}, the absence of common zeros of consecutive
Hermite polynomials, and the reductions of V53--V57.  No floating-point
quantity enters the argument.
\endgroup

\section*{Version V58 append-only record}
\addcontentsline{toc}{section}{Version V58 append-only record}

V58 was appended byte-for-byte before the terminal marker of validated
V57, whose installed SHA--256 was
\[
 \texttt{6D75990A538AA0D728C43C794A1CE620}
 \allowbreak\texttt{E2DE15212C2A8C882B9CA959245973CC}
\]
with \(16669\) lines and \(554420\) bytes.  It certifies the last nine
groups by exact facial reduction and proves the two-pair band closure
for at least five real padding zeros.  After validation V58 replaces
V57 as the sole active numeric GRH note; the frozen \texttt{V100\_f2}
archive is unchanged.

% =============================================================================
% END APPENDED V58 MATERIAL
% =============================================================================




\clearpage
\section{V59: the cases of three and four real padding zeros, and the complete two-pair band theorem}
\label{sec:f2v59-p34}
\begingroup
\setcounter{equation}{0}
\renewcommand{\theequation}{N59.\arabic{equation}}
\renewcommand{\theHequation}{f2v59.\arabic{equation}}

V58 left the two-pair block open for \(p=3\) and \(p=4\), where the
basis \((\mathrm{He}_{p-2},\mathrm{He}_{p-3})\) degenerates.  The same
construction with \(p\) fixed produces exact coefficient polynomials, but
the Pick-disk criterion of V56 cannot be used: for \(p=3\) and \(p=4\) the
real-axis coefficient \(W_0\) of \(\operatorname{Im}(N\overline D)\) is
not nonnegative on the region (numerically it is unbounded below).  This
section closes both cases by a different mechanism: the resultant of the
real and imaginary parts of the diagonal is certified positive on the
whole region, which excludes real zeros, and a continuity argument from
one exactly certified sample point places all zeros in the lower
half-plane.  Together with V58 this proves the band theorem for every
two-pair block.

\subsection{Real and imaginary parts of the diagonal}

Fix \(p\in\{3,4\}\).  With \(f=(w-c)^2+t\), \(G=(w-a)^2+\beta+1\) and
\(\Pi=fG(w-\zeta)^p\), expand \(\Pi\) in \(w\) with coefficients in
\(\mathbb Q[c,t,a,\beta,\zeta]\) and apply
\((e^{-D_w^2/2}w^j)(0)=(-1)^{j/2}j!/(2^{j/2}(j/2)!)\) for even \(j\) and
\((e^{-D_w^2/2}w^j)'(0)=(-1)^{(j-1)/2}j!/(2^{(j-1)/2}((j-1)/2)!)\) for
odd \(j\).  This gives
\begin{equation}
 F_G(\zeta)=A(\zeta)-iB(\zeta),\qquad
 A=\sum_{i=0}^pA_i\zeta^i,\quad B=\sum_{i=0}^pB_i\zeta^i,\quad
 A_i,B_i\in\mathbb Q[c,t,a,\beta],
 \label{eq:f2v59-AB}
\end{equation}
where \(A(\zeta)=Q(0)\) and \(B(\zeta)=Q'(0)\) are real for real
\(\zeta\).  On the region put
\begin{equation}
 t=p+4+\sigma,\qquad \beta=p+3+\sigma+\rho,\qquad \sigma,\rho\ge0,
 \label{eq:f2v59-region}
\end{equation}
so that \(t\ge p+4\) and \(t_2=\beta+1\ge t\).

\begin{lemma}[Positive resultant]
\label{lem:f2v59-resultant}
Let \(\mathcal R:=\operatorname{Res}_\zeta(A,B)\) be the Sylvester
resultant of the two degree-\(p\) polynomials \eqref{eq:f2v59-AB}, a
polynomial in \((c,t,a,\beta)\) of degree \((12,6,12,6)\) for \(p=3\)
and \((16,8,16,8)\) for \(p=4\).  After the substitution
\eqref{eq:f2v59-region} write
\(\mathcal R=\sum_{s,r}\sigma^s\rho^rR_{s,r}(c,a)\).  Every
\(R_{s,r}\) is nonnegative on \(\mathbb R^2\), and
\begin{equation}
 R_{0,0}(c,a)\ \ge\ 6078844800\quad(p=3),\qquad
 R_{0,0}(c,a)\ \ge\ 4355089945929216\quad(p=4).
 \label{eq:f2v59-corner}
\end{equation}
Consequently \(\mathcal R>0\) on the whole region: \(A\) and \(B\) have
no common zero, and their leading coefficients do not vanish
simultaneously.
\end{lemma}

\begin{proof}
There are \(70\) groups for \(p=3\) and \(117\) for \(p=4\).  After the
substitution every negative coefficient of \(\mathcal R\) sits on a
monomial odd in both \(c\) and \(a\), and each group is certified either
by an absorption ledger \eqref{eq:f2v56-amgm} or by an exact
Gram identity \(R_{s,r}=b^{\mathsf T}Gb+\)(nonnegative even remainder)
with a rational positive semidefinite \(G\) over a polynomial basis,
obtained after facial reduction and a power-of-two rescaling of
\(c,a\).  The corner groups carry the exact constants
\eqref{eq:f2v59-corner}, one half of \(R_{0,0}(0,0)\).  A common zero of
\(A\) and \(B\), or the simultaneous vanishing of \(A_p\) and \(B_p\),
would make the Sylvester determinant vanish.
\end{proof}

\begin{lemma}[Sample point]
\label{lem:f2v59-sample}
At \(c=a=0\), \(\sigma=\rho=0\), the polynomial \(F_G\) has exact
degree \(p\) and all its zeros lie in the open lower half-plane.  For
\(p=3\) they are enclosed in disjoint disks of radius \(10^{-3}\) around
\(\pm0.957795-0.329903i\) and \(-1.077037i\); for \(p=4\) around
\(\pm1.454504-0.246976i\), \(-0.316858i\), and \(-1.620563i\).
\end{lemma}

\begin{proof}
For each listed center \(\zeta_0\) and radius \(r\) the exact Taylor
coefficients \(F^{(k)}(\zeta_0)/k!\) are computed over
\(\mathbb Q+i\mathbb Q\) and the inequality
\(|F'(\zeta_0)|\,r>|F(\zeta_0)|+\sum_{k\ge2}|F^{(k)}(\zeta_0)/k!|\,r^k\)
is verified with rational square-root bounds; Rouch\'{e}'s theorem
gives exactly one zero in each disk, the disks are pairwise disjoint and
lie in \(\operatorname{Im}\zeta<0\), and there are \(p\) of them.
\end{proof}

\begin{theorem}[Zero-free diagonal for \(p=3,4\)]
\label{thm:f2v59-zero-free}
For \(p\in\{3,4\}\), all real \(c,a\), \(t\ge p+4\), and \(t_2\ge t\),
the reals-only diagonal \(F_G\) has no zero in the closed upper
half-plane.
\end{theorem}

\begin{proof}
On the connected region \eqref{eq:f2v59-region} the leading coefficient
\(A_p-iB_p\) never vanishes by Lemma~\ref{lem:f2v59-resultant}, so the
\(p\) zeros of \(F_G\) depend continuously on the parameters, and none
of them is real, since a real zero would be a common zero of \(A\) and
\(B\).  Hence the number of zeros in the open upper half-plane is locally
constant, therefore constant, and it equals zero at the sample point of
Lemma~\ref{lem:f2v59-sample}.
\end{proof}

\subsection{The complete two-pair theorem}

\begin{theorem}[Two-pair band closure]
\label{thm:f2v59-band}
Let \(P_0\) be a real monic polynomial with exactly two conjugate pairs,
of heights \(y_1\ge y_2>0\), and \(p\ge0\) real zeros, \(n=p+4\).  If
\(e^{-(K/2)D_z^2}P_0\) is real-rooted for some \(K>0\), then
\begin{equation}
 \boxed{\;y_2^2\le nK,\;}
 \label{eq:f2v59-bound}
\end{equation}
with strict inequality for \(p\ge3\).
\end{theorem}

\begin{proof}
For \(p\le2\) this is \textup{(N9.14)}.  For \(p\ge5\) it is
Theorem~\ref{thm:f2v58-band}.  For \(p=3,4\), Theorem~\ref{thm:f2v53-reduction}
produces a zero of \(F_G\) in the closed upper half-plane with
\(t=y_2^2/T\), \(t_2=y_1^2/T\ge t\), \(T\le K\); by
Theorem~\ref{thm:f2v59-zero-free} this forces \(t<p+4\).
\end{proof}

\begin{corollary}[No two-channel block enters the band]
\label{cor:f2v59-gamma}
In the notation of \textup{(N9.24)}, every critical block with two
conjugate-pair channels, any number \(p\ge0\) of real gamma factors, and
ladder count \(R\ge n=p+4\), whose critical heat envelope is
real-rooted, satisfies
\begin{equation}
 \boxed{\;\kappa^2\ge\frac{\gamma^2R}{n(1+\theta)}\ \ge\ \frac{\gamma^2}{1+\theta},\;}
 \label{eq:f2v59-gamma}
\end{equation}
where \(\gamma\) is the least ordinate of the two channels.  Together
with Corollary~\ref{cor:f2v5-gamma} and \textup{(N9.23)}, no finite block
with at most two conjugate-pair channels enters the open V97 band.
\end{corollary}

\begin{remark}[Exact verifier and data]
\label{rem:f2v59-verifier}
\path{scripts/verify_grh_v59_two_pair_p34.py} rebuilds \(A,B\), the
Sylvester resultant, the region substitution and the groups from
scratch for both \(p\), verifies every group certificate of
\path{scripts/data_grh_v59_p34_certificates.json}
(\(2233769\) bytes, SHA--256
\texttt{338B408AEEC5C919A99E5587BF9D2442}\allowbreak
\texttt{0F4298BECC354F8CA1BC978F6228C191}), the corner constants
\eqref{eq:f2v59-corner}, and the Rouch\'{e} disks of
Lemma~\ref{lem:f2v59-sample}.  It has \(291\) lines, \(9089\) bytes,
SHA--256
\begin{center}
\texttt{C526C4F285252E230955AF6315265D23}\\
\texttt{8956BDB75C3F1E3C8D6A77CF72910278},
\end{center}
runs in a few seconds, and terminates with \texttt{VERIFIED}.
\end{remark}

\subsection*{Firewall and V60 direction}
\label{fire:f2v59}
Theorem~\ref{thm:f2v59-band} completes the finite two-pair problem: with
V5 and V9, every finite block with at most two conjugate-pair channels is
excluded from the band.  Blocks with three or more pairs are not covered,
and nothing here addresses the infinite-product passage or the
arithmetic forcing step.  GRH remains open.

V60 is the mandatory ten-version companion sweep.  After it, the
three-pair mixed block is the next target: the reals-only reduction of
V53 applies verbatim with \(G\) of degree four, the diagonal still has
the form \eqref{eq:f2v53-AB}, and the resultant-plus-continuity mechanism
of this section is degree-free, so it is the natural first attempt
before any Pick-type argument with two companions.

\subsection*{Assumption ledger}
\label{ass:f2v59}
The proof uses exact rational arithmetic throughout, the Sylvester
resultant, Rouch\'{e}'s theorem on explicit disks, continuity of the
zeros of a polynomial of constant degree, connectedness of the region,
and the reductions of V53.  The floating-point search that found the
certificates and the observation about \(W_0\) are not part of the
proof.
\endgroup

\section*{Version V59 append-only record}
\addcontentsline{toc}{section}{Version V59 append-only record}

V59 was appended byte-for-byte before the terminal marker of validated
V58, whose installed SHA--256 was
\[
 \texttt{5842924D4A329D15AE90725A751FDEAC}
 \allowbreak\texttt{CDEAD9885499750F50BA1C099D81F098}
\]
with \(16890\) lines and \(563647\) bytes.  It closes the cases
\(p=3,4\) by a certified positive resultant and a sample-point
continuity argument, completing the two-pair band theorem.  After
validation V59 replaces V58 as the sole active numeric GRH note; the
frozen \texttt{V100\_f2} archive is unchanged.

% =============================================================================
% END APPENDED V59 MATERIAL
% =============================================================================




\clearpage
\section{V60: companion sweep, a degree-free exclusion schema, and the three-pair gates}
\label{sec:f2v60-sweep-schema}
\begingroup
\setcounter{equation}{0}
\renewcommand{\theequation}{N60.\arabic{equation}}
\renewcommand{\theHequation}{f2v60.\arabic{equation}}

This is the mandatory ten-version checkpoint.  The companion notes were
read before any new mathematics.  The active sources are
\begin{equation}
\begin{array}{c|r|l}
\text{source}&\text{lines}&\text{SHA--256}\\ \hline
\text{SM--Einstein V91}&18850&
\texttt{3AF43F12D134A310C91E960E0B6FE685}\\
&&\texttt{8E811AF2D83DD7F70243EE92F7E7BCA9}\\
\text{Yang--Mills V34}&12282&
\texttt{3CF829F57455C43D3CE54387782C3C60}\\
&&\texttt{EB3497F1B680A11DDCAED13306D8C469}\\
\text{Navier--Stokes V26}&5319&
\texttt{82C887F8C2C69E4F28FAD7C3DC8DE5B9}\\
&&\texttt{099CB5A4BF431EBA1FFF3E91935978AD}
\end{array}
\label{eq:f2v60-sources}
\end{equation}
The Navier--Stokes file is unchanged since the V50 sweep.

\begin{proposition}[Sweep decision]
\label{prop:f2v60-sweep}
No theorem of the companion snapshot transfers to the GRH chain.  Two
disciplines do: the exact-series certification of a reduced system in
SM V91, which is the same practice as the chained exact verifiers used
since V44; and the replacement of an archived obstruction by explicit
hypotheses with concrete tasks in Yang--Mills V34, which is the practice
followed in V52--V59.
\end{proposition}

\begin{proof}
SM V91 proves a transverse collapse theorem for an averaged
cosmological system driven by a plane wave and a negative cosmological
constant; Yang--Mills V34 proves permutation covariance of a symmetrized
blocking rule and a linear alias recursion for the blocked propagator.
Neither statement involves a heat flow, a Hermite polynomial, a
coalescence functional, or a critical gamma block, so importing either
would be a type mismatch.  The Navier--Stokes note is unchanged.
\end{proof}

\subsection{A degree-free exclusion schema}

The V59 argument uses nothing specific to two pairs.  It is recorded
here in general form.  Let \(m\ge2\), \(p\ge0\), \(n=2m+p\), and let
\(G(w)=\prod_{j=2}^m((w-a_j)^2+\beta_j+1)\).  With \(f=(w-c)^2+t\)
and \(\Pi=fG(w-\zeta)^p\), write as in \eqref{eq:f2v59-AB}
\(F_G(\zeta)=A(\zeta)-iB(\zeta)\) with real polynomials \(A,B\) of
formal degree \(p\) in \(\zeta\).  The region is
\begin{equation}
 \Omega_{m,p}=\{c,a_j\in\mathbb R,\ t\ge n,\ \beta_j\ge t-1\}.
 \label{eq:f2v60-region}
\end{equation}

\begin{theorem}[Exclusion schema]
\label{thm:f2v60-schema}
Suppose that for the pair \((m,p)\):
\begin{enumerate}[label=\textup{(\roman*)}]
\item \(A\) and \(B\) have no common real zero \(\zeta\) at any point
of \(\Omega_{m,p}\); and
\item at one point of \(\Omega_{m,p}\) all \(p\) zeros of \(F_G\) lie
in the open lower half-plane.
\end{enumerate}
Then \(F_G\) has no zero in the closed upper half-plane on
\(\Omega_{m,p}\), and every real-rooted heat image of a block with
\(m\) conjugate pairs and \(p\) real zeros satisfies
\(y_{\min}^2<nK\).
\end{theorem}

\begin{proof}
The leading coefficient of \(F_G\) is \(\pm\mathcal F(fG)\), the
collision functional of the pure \(m\)-pair block \(fG\) at heat time
\(1\).  If it vanished, that block would have a real zero of
multiplicity two in its heat image, hence a first collision at some
time \(T\le1\), and the first-collision bound \textup{(N9.14)} would
give \(t\le mT\le m<n\) for its least height, contradicting
\(t\ge n\).  Hence \(F_G\) has exact degree
\(p\) throughout \(\Omega_{m,p}\), its zeros depend continuously on the
parameters, none is real by (i), and the number of zeros in the open
upper half-plane is locally constant on the connected set
\(\Omega_{m,p}\); by (ii) it is zero.  The bound follows from
Theorem~\ref{thm:f2v53-reduction} exactly as in
Theorem~\ref{thm:f2v59-band}.
\end{proof}

Condition (ii) is checked at one point by exact Rouch\'{e} disks as in
Lemma~\ref{lem:f2v59-sample}.  Condition (i) admits two sufficient
certificates: the resultant \(\operatorname{Res}_\zeta(A,B)\) is
positive on \(\Omega_{m,p}\) (V59), or, when \(p\ge5\), the
real-axis determinant
\begin{equation}
 W_0=A_uB_v-A_vB_u\quad\text{in the basis }(\mathrm{He}_{p-2},\mathrm{He}_{p-3})
 \label{eq:f2v60-W0}
\end{equation}
is positive on \(\Omega_{m,p}\times\mathbb R_\zeta\), because at a real
zero the nonzero vector \((\mathrm{He}_{p-2},\mathrm{He}_{p-3})(\zeta)\)
would lie in the kernel of the real matrix
\(\bigl(\begin{smallmatrix}A_v&A_u\\B_v&B_u\end{smallmatrix}\bigr)\).
The determinant is a single polynomial of much lower degree than the
resultant.

\subsection{Three-pair gates}

\begin{remark}[Numerical gates for \(m=3,4\); not proofs]
\label{rem:f2v60-gates}
With the balanced construction of V55 evaluated numerically for two and
three companion pairs, random sampling of \(\Omega_{m,p}\) shows:
\begin{enumerate}[label=\textup{(\alph*)}]
\item the Pick-disk quantity \(\operatorname{Im}(N\overline D)\) of
V56 takes negative values for \(m=3\), \(p\in\{6,7,9\}\) and is
therefore not a usable sufficient condition beyond two pairs;
\item the real-axis determinant \eqref{eq:f2v60-W0} stayed positive on
\(6000\) samples for \(m=3\), \(p\in\{7,9,12\}\) and for \(m=4\),
\(p\in\{9,12\}\), and took negative values for \(m=3\), \(p\le6\) and
\(m=4\), \(p\le7\); the same failure occurs in the alternative bases
\((\mathrm{He}_{p-d},\mathrm{He}_{p-d-1})\), \(d=0,1,3\), in the
middle range \(m<p\le2m\).
\end{enumerate}
The pattern suggests that the determinant certificate covers
\(p\ge2m+1\), while the middle range \(m<p\le2m\) will need the
resultant or a sharper real-axis criterion.
\end{remark}

\subsection*{Firewall and V61 direction}
\label{fire:f2v60}
Theorem~\ref{thm:f2v60-schema} is a proved reduction; it produces a
band theorem for \((m,p)\) only once (i) and (ii) are certified, which
so far has been done for \(m=2\).  Nothing here concerns blocks of
unbounded size, the infinite-product passage, or arithmetic forcing.
GRH remains open.

V61 should construct \(W_0\) for \(m=3\) exactly, as a polynomial in
\((x,c,t,a_2,\beta_2,a_3,\beta_3,p)\), substitute
\(t=n+\sigma\), \(\beta_j=t-1+\rho_j\), \(p=7+\pi\), and certify the
\((\sigma,\rho_2,\rho_3,\pi)\)-groups in \((x,c,a_2,a_3,w)\) by the
V57--V58 pipeline, which would close the three-pair band for
\(p\ge7\); the sample-point check is routine.  The middle cases
\(p=4,5,6\) should then be attacked by the resultant, whose size is the
only obstacle.

\subsection*{Assumption ledger}
\label{ass:f2v60}
Theorem~\ref{thm:f2v60-schema} uses continuity of the zeros of a
polynomial of constant degree, connectedness of \(\Omega_{m,p}\), the
strip lemma, \textup{(N9.14)}, and the V53 reduction.  The sweep imports
no companion constant or hypothesis.  Remark~\ref{rem:f2v60-gates} is
numerical exploration only.
\endgroup

\section*{Version V60 append-only record}
\addcontentsline{toc}{section}{Version V60 append-only record}

V60 was appended byte-for-byte before the terminal marker of validated
V59, whose installed SHA--256 was
\[
 \texttt{C27D11EB55A9A6003FB5FE7766F0335D}
 \allowbreak\texttt{61A002AB51F8C40CEAAA9C15D36ED561}
\]
with \(17105\) lines and \(572973\) bytes.  It records the mandatory
companion sweep, proves the degree-free exclusion schema, and sets the
three-pair targets.  After validation V60 replaces V59 as the sole
active numeric GRH note; the frozen \texttt{V100\_f2} archive is
unchanged.

% =============================================================================
% END APPENDED V60 MATERIAL
% =============================================================================




\clearpage
\section{V61: the three-pair determinant for at least seven real padding zeros, up to six residual inequalities}
\label{sec:f2v61-three-pair}
\begingroup
\setcounter{equation}{0}
\renewcommand{\theequation}{N61.\arabic{equation}}
\renewcommand{\theHequation}{f2v61.\arabic{equation}}

This iteration carries out the V60 plan for \(m=3\).  The real-axis
determinant \(W_0\) is constructed exactly with \(p\) symbolic and split
into \(578\) groups on the region; \(572\) of them are certified
nonnegative by exact rational certificates, and the exclusion schema is
closed by a deformation to the one-pair diagonal of V5 instead of a
sample point for each \(p\).  Six groups resisted every exact
certification attempt of this session although they are numerically
strictly positive; the band theorem for three pairs and \(p\ge7\) is
therefore proved conditionally on six explicit four-variable polynomial
inequalities, which are exported with the data and stated precisely
below.

\subsection{The determinant with two companions}

Let \(G=G_2G_3\), \(G_j=(w-a_j)^2+\beta_j+1\), and
\(\widetilde G=e^{-D^2/2}G\), a polynomial of degree four in \(D\) whose
coefficients lie in \(\mathbb Q[a_2,\beta_2,a_3,\beta_3]\).  In the basis
\((v,u)=(\mathrm{He}_{p-2},\mathrm{He}_{p-3})\) the reductions of
\eqref{eq:f2v53-hk} require \(\mathrm{He}_{p-k}\) for \(0\le k\le7\); the
downward recurrence divides by \((p-3),\ldots,(p-6)\), so every quantity
is multiplied by
\begin{equation}
 M_3:=(p-3)(p-4)(p-5)(p-6)>0\qquad(p\ge7),
 \label{eq:f2v61-M}
\end{equation}
and \(M_3\mathrm{He}_{p-k}=\alpha_kv+\beta_ku\) with
\(\alpha_k,\beta_k\in\mathbb Q[x,p]\) at real \(\zeta=x\).  With
\eqref{eq:f2v53-AB} this gives real polynomials \(A_v,A_u,B_v,B_u\) and
\begin{equation}
 W_0:=A_uB_v-A_vB_u\in\mathbb Q[x,c,t,a_2,\beta_2,a_3,\beta_3,p],
 \label{eq:f2v61-W0}
\end{equation}
with \(18919\) terms and degrees \((6,4,2,4,2,4,2,17)\).  The exact
identity
\(M_3F_G(x)=(-1)^p[(A_v-iB_v)\mathrm{He}_{p-2}(x)+(A_u-iB_u)\mathrm{He}_{p-3}(x)]\)
was checked in floating point against the direct heat computation at
random parameters.  On the region
\begin{equation}
 t=p+6+\sigma,\qquad \beta_j=t-1+\rho_j,\qquad p=7+\pi,\qquad
 \sigma,\rho_2,\rho_3,\pi\ge0,
 \label{eq:f2v61-region}
\end{equation}
write \(W_0=\sum\sigma^s\rho_2^{r_2}\rho_3^{r_3}\pi^kR_{s,r_2,r_3,k}(x,c,a_2,a_3)\);
there are \(578\) groups.

\begin{proposition}[Five hundred and seventy-two certified groups]
\label{prop:f2v61-groups}
Every group except the six listed in \eqref{eq:f2v61-open} is
nonnegative on \(\mathbb R^4\), and the corner group satisfies
\begin{equation}
 R_{0,0,0,0}(x,c,a_2,a_3)\ \ge\ 3786531840 .
 \label{eq:f2v61-corner}
\end{equation}
\end{proposition}

\begin{proof}
\(118\) groups are certified by absorption ledgers
\eqref{eq:f2v56-amgm}; \(453\) by exact Gram identities
\(R=b^{\mathsf T}Gb+\)(nonnegative even remainder) with rational positive
semidefinite \(G\) over polynomial bases; and one, \(R_{0,0,0,2}\), by
such an identity for \((1+x^2)R_{0,0,0,2}\), which suffices because
\(1+x^2>0\).  The bases were obtained by an exact diagonal-consistency
reduction of the Newton basis, by the integer face relations forced by
the leading forms in the eighty sign-weight directions (for instance
\(x^2a_3^2+x^3a_3\), which forces the factor \(x^2a_3(x-a_3)\) in the top
layer), by validated numerical kernel rounds, and by the exact
projection step of V58.  The corner constant is one half of
\(R_{0,0,0,0}(0,0,0,0)\).
\end{proof}

\begin{definition}[The residual inequalities]
\label{def:f2v61-open}
Hypothesis (H61) is the statement that the six groups
\begin{equation}
 R_{0,0,0,6},\quad R_{0,0,0,8},\quad R_{0,0,0,10},\quad R_{0,0,0,11},\quad
 R_{0,0,0,12},\quad R_{1,0,0,2}
 \label{eq:f2v61-open}
\end{equation}
are nonnegative on \(\mathbb R^4\).  Their exact coefficients are stored
in the data file of Remark~\ref{rem:f2v61-verifier}, and the verifier
checks that they coincide with the groups rebuilt from scratch.
Numerically each has a strictly positive minimum; five of them are the
\(\pi^k\)-coefficients of the corner family \(\sigma=\rho_2=\rho_3=0\),
where all three pairs have equal height.
\end{definition}

\begin{theorem}[Positive determinant under (H61)]
\label{thm:f2v61-W0}
Assume (H61).  Then \(W_0>0\) for all real \(x\), all real
\(c,a_2,a_3\), all real \(p\ge7\), \(t\ge p+6\), and \(\beta_j\ge t-1\).
\end{theorem}

\begin{proof}
By Proposition~\ref{prop:f2v61-groups} and (H61) every group is
nonnegative, the corner group is bounded below by
\eqref{eq:f2v61-corner}, and \(\sigma,\rho_j,\pi\ge0\).
\end{proof}

\subsection{Closing the schema without sample points}

\begin{lemma}[Deformation to one pair]
\label{lem:f2v61-deform}
Let \(p\ge7\), \(t>p+2\), and put \(u_j=1/\beta_j\in(0,1/(t-1)]\).  Then
\(\widehat F(\zeta;u_2,u_3):=F_G(\zeta)/(\beta_2\beta_3)\) extends to a
polynomial in \(\zeta\) of exact degree \(p\) whose coefficients are
polynomials in \((c,t,a_j,u_j)\) on the closed set \(0\le u_j\le1/(t-1)\),
and \(\widehat F(\zeta;0,0)=\mathcal F\bigl(f(w-\zeta)^p\bigr)\) is the
one-pair diagonal of Theorem~\ref{thm:f2v5-diagonal}, which has no zero
in the closed upper half-plane.
\end{lemma}

\begin{proof}
\(G_j/\beta_j=1+u_j\bigl((w-a_j)^2+1\bigr)\) is a polynomial in
\(u_j\), and \(\mathcal F\) is linear in the coefficients of its
argument, so \(\widehat F\) is polynomial in \((\zeta,c,t,a_j,u_j)\) and
reduces at \(u_2=u_3=0\) to the one-pair functional.  Its leading
coefficient is \(\pm\mathcal F(fG)/(\beta_2\beta_3)\) for \(u_j>0\),
nonzero by the argument of Theorem~\ref{thm:f2v60-schema}, and
\(\pm\mathcal F(f)=\pm\bigl((c^2+t-1)+2ic\bigr)\ne0\) at \(u_j=0\).
Theorem~\ref{thm:f2v5-diagonal} excludes zeros with
\(\operatorname{Im}\zeta>0\) for \(t>p+2\) and real zeros unless
\(t=p+2\).
\end{proof}

\begin{theorem}[Zero-free three-pair diagonal under (H61)]
\label{thm:f2v61-zero-free}
Assume (H61).  For every integer \(p\ge7\), all real \(c,a_2,a_3\),
\(t\ge p+6\), and \(\beta_j\ge t-1\), the diagonal \(F_G\) has no zero
in the closed upper half-plane.
\end{theorem}

\begin{proof}
Consider the connected parameter set
\(\{c,a_j\in\mathbb R,\ t\ge p+6,\ 0\le u_j\le1/(t-1)\}\) and the family
\(\widehat F\) of Lemma~\ref{lem:f2v61-deform}, of constant degree
\(p\).  For \(u_j>0\) a real zero of \(\widehat F\) is a real zero of
\(F_G\), which would make the real vector
\((\mathrm{He}_{p-2},\mathrm{He}_{p-3})(x)\ne0\) a kernel vector of the
real matrix \(\bigl(\begin{smallmatrix}A_v&A_u\\B_v&B_u\end{smallmatrix}\bigr)\),
contradicting \(W_0>0\) (Theorem~\ref{thm:f2v61-W0}); for \(u_j=0\)
real zeros are excluded by Lemma~\ref{lem:f2v61-deform}.  Hence the
number of zeros in the open upper half-plane is locally constant on a
connected set and equals its value at \(u_j=0\), which is zero.
\end{proof}

\begin{theorem}[Three-pair band closure for \(p\ge7\), under (H61)]
\label{thm:f2v61-band}
Assume (H61).  Let \(P_0\) be a real monic polynomial with exactly three
conjugate pairs and \(p\ge7\) real zeros, \(n=p+6\).  If
\(e^{-(K/2)D_z^2}P_0\) is real-rooted for some \(K>0\), then the least
pair height satisfies
\begin{equation}
 y_{\min}^2<nK .
 \label{eq:f2v61-bound}
\end{equation}
For \(p\le3\) the bound \(y_{\min}^2\le nK\) holds unconditionally by
\textup{(N9.14)}.
\end{theorem}

\begin{proof}
Theorem~\ref{thm:f2v53-reduction} with \(m=3\) and
Theorem~\ref{thm:f2v61-zero-free}, exactly as in
Theorem~\ref{thm:f2v59-band}.
\end{proof}

\begin{remark}[What is unconditional]
\label{rem:f2v61-unconditional}
Lemma~\ref{lem:f2v61-deform} and the continuity argument are proved
without hypothesis; they replace the per-\(p\) sample points of
Theorem~\ref{thm:f2v60-schema}(ii) for every \(m\) by the single
one-pair theorem of V5.  Proposition~\ref{prop:f2v61-groups} is
unconditional.  What remains is exactly (H61).
\end{remark}

\begin{remark}[Exact verifier and data]
\label{rem:f2v61-verifier}
\path{scripts/verify_grh_v61_three_pair_band.py} rebuilds
\(W_0\), the substitution and the \(578\) groups from scratch, verifies
every certificate of
\path{scripts/data_grh_v61_three_pair_certificates.json}
(\(2730136\) bytes, SHA--256
\texttt{18775A154807ECD63381C16B1A2BA352}\allowbreak
\texttt{0DBA2FC9347F66F572E47E223D206175}), checks that the six stored
residual polynomials coincide with the rebuilt groups
\eqref{eq:f2v61-open}, and prints the corner constant.  It has \(292\)
lines, \(9276\) bytes, SHA--256
\begin{center}
\texttt{AC0C1391F16E2D62A417D135E655D04A}\\
\texttt{EE18D70D9EDC75C97C10762B3639C05B},
\end{center}
runs in about five minutes, and terminates with
\texttt{VERIFIED\_CONDITIONAL}, naming the six open groups.
\end{remark}

\subsection*{Firewall and V62 direction}
\label{fire:f2v61}
Theorem~\ref{thm:f2v61-band} is conditional on (H61); it is not yet a
theorem about all three-pair blocks with \(p\ge7\), and nothing here
proves GRH.  GRH remains open.  The six residual polynomials are
strictly positive numerically (normalized minima above \(0.02\)); the
obstruction is that their sum-of-squares certificates have a
high-dimensional facial degeneracy inherited from the equal-height
corner \(t=t_2=t_3\), which the numerical kernel reconstruction of this
session could not resolve.

V62 should discharge (H61) by an exact method: either a
symmetry-adapted Gram basis under \(a_2\leftrightarrow a_3\) with the
face relations imposed exactly, or an explicit factorization of the
corner family at equal heights, which is the natural first thing to
compute.  Once (H61) is discharged, the three-pair band is closed for
\(p\ge7\) and \(p\le3\), and the middle cases \(p=4,5,6\) follow the
V59 resultant route.

\subsection*{Assumption ledger}
\label{ass:f2v61}
Exact rational arithmetic, the Leibniz and generating identities of
V53, the Hermite recurrence, absorption and Gram certificates verified
exactly, Theorem~\ref{thm:f2v5-diagonal}, the first-collision bound
\textup{(N9.14)}, continuity of zeros, connectedness, and hypothesis
(H61) where stated.  The floating-point search that produced the
certificates is not part of the proof.
\endgroup

\section*{Version V61 append-only record}
\addcontentsline{toc}{section}{Version V61 append-only record}

V61 was appended byte-for-byte before the terminal marker of validated
V60, whose installed SHA--256 was
\[
 \texttt{E090E736AE1C43D60756DB3E9C7929E0}
 \allowbreak\texttt{F2A179905B74630ADBE85266036737F4}
\]
with \(17282\) lines and \(580602\) bytes.  It constructs the three-pair
determinant, certifies \(572\) of its \(578\) region groups, proves the
deformation closure of the exclusion schema, and states the three-pair
band theorem for \(p\ge7\) conditionally on six explicit residual
inequalities.  After validation V61 replaces V60 as the sole active
numeric GRH note; the frozen \texttt{V100\_f2} archive is unchanged.

% =============================================================================
% END APPENDED V61 MATERIAL
% =============================================================================




\clearpage
\section{V62: the Appell reformulation of the real-axis problem}
\label{sec:f2v62-appell}
\begingroup
\setcounter{equation}{0}
\renewcommand{\theequation}{N62.\arabic{equation}}
\renewcommand{\theHequation}{f2v62.\arabic{equation}}

The residual inequalities of V61 are symptoms of a certificate
technology applied to the wrong object.  This iteration identifies the
right object.  On the real axis the reals-only diagonal vanishes exactly
at the real multiple zeros of a single polynomial, the image of a
Hermite polynomial under the differential operator whose symbol is the
time-one heat image of the pure block.  The real-axis half of the
exclusion schema for every number of pairs therefore reduces to a
real-rootedness statement about one Appell polynomial, which is
numerically true with wide margins on every region tested, is sharp
exactly in the one-pair case where it coincides with V5, and implies
hypothesis (H61).

\subsection{The Appell identity}

Let \(\Pi_m:=fG=\prod_{j=1}^m((w-c_j)^2+t_j)\) be the pure \(m\)-pair
block, with \(c_1=c\), \(t_1=t\) the target pair and \(t_j\ge t\), and
let
\begin{equation}
 \widehat Q:=e^{-D^2/2}\Pi_m,\qquad T:=\widehat Q(D),\qquad
 P_r:=T\,\mathrm{He}_r\quad(r\ge0).
 \label{eq:f2v62-T}
\end{equation}
Since \(T\) commutes with \(D\) and \(D\mathrm{He}_r=r\mathrm{He}_{r-1}\),
the sequence \((P_r)\) is an Appell sequence: \(P_r'=rP_{r-1}\).  Its
generating function is \(\sum_rP_r(x)u^r/r!=\widehat Q(u)e^{ux-u^2/2}\).

\begin{theorem}[Appell form of the diagonal on the real axis]
\label{thm:f2v62-appell}
For real \(\zeta=x\), the reals-only diagonal of
Theorem~\ref{thm:f2v53-reduction} satisfies
\begin{equation}
 F_G(x)=A(x)-iB(x),\qquad
 A=\frac{(-1)^p}{p+1}P_{p+1}',\qquad
 B=(-1)^pP_{p+1}-xA .
 \label{eq:f2v62-AB}
\end{equation}
Consequently \(F_G(x)=0\) if and only if \(x\) is a real zero of
\(P_{p+1}\) of multiplicity at least two.
\end{theorem}

\begin{proof}
Apply Lemma~\ref{lem:f2v53-generating} to the whole block, that is with
\(G\) replaced by \(\Pi_m\), so that \(\widetilde G=\widehat Q\).  The
value and the first derivative at \(0\) of
\(e^{-D_w^2/2}[\Pi_m(w)(w-x)^p]\) are
\[
 h_0=(-1)^pT\mathrm{He}_p,\qquad
 h_1=(-1)^{p-1}p\,T\mathrm{He}_{p-1}+(-1)^p\widehat Q'(D)\mathrm{He}_p,
\]
and \(F_G(x)=h_0-ih_1\), both real for real \(x\).  On polynomials
\([\widehat Q(D),x]=\widehat Q'(D)\), so with the three-term recurrence
\(x\mathrm{He}_p=\mathrm{He}_{p+1}+p\mathrm{He}_{p-1}\),
\[
 \widehat Q'(D)\mathrm{He}_p=T(x\mathrm{He}_p)-xT\mathrm{He}_p
 =P_{p+1}+pP_{p-1}-xP_p .
\]
Hence \(h_1=(-1)^p(P_{p+1}-xP_p)\), while \(h_0=(-1)^pP_p\) and
\(P_p=P_{p+1}'/(p+1)\) by the Appell property.  This is
\eqref{eq:f2v62-AB}.  If \(F_G(x)=0\) then \(A(x)=0\), so
\(P_{p+1}'(x)=0\), and then \(B(x)=0\) gives \(P_{p+1}(x)=0\); the
converse is immediate.
\end{proof}

\begin{corollary}[Real-axis criterion for every \(m\)]
\label{cor:f2v62-criterion}
Condition (i) of Theorem~\ref{thm:f2v60-schema} holds for \((m,p)\) as
soon as, for every block \(\Pi_m\) on \(\Omega_{m,p}\), the polynomial
\(P_{p+1}=\widehat Q(D)\mathrm{He}_{p+1}\) has only real and simple zeros.
Together with Lemma~\ref{lem:f2v61-deform}, which holds for every
\(m\), this real-rootedness statement implies
\(y_{\min}^2<nK\) for all blocks with \(m\) pairs and \(p\) real
zeros.
\end{corollary}

\begin{proof}
Real simple zeros of \(P_{p+1}\) exclude real multiple zeros, hence
real zeros of \(F_G\) by Theorem~\ref{thm:f2v62-appell}; the rest is
Theorem~\ref{thm:f2v60-schema} with the sample-point condition
replaced by the deformation of Lemma~\ref{lem:f2v61-deform}, whose
proof did not use \(m=3\).
\end{proof}

\begin{lemma}[Where the symbol's zeros are]
\label{lem:f2v62-symbol}
On \(\Omega_{m,p}\) the polynomial \(\widehat Q\) has no real zero, and
all its zeros satisfy \(|\operatorname{Im}|\ge\sqrt{t-m}\ge\sqrt{m+p}\).
\end{lemma}

\begin{proof}
\(\widehat Q\) is the time-one backward heat image of a block with
\(m\) conjugate pairs and no real zero.  By the first-collision bound
\textup{(N9.14)} with \(p=0\), a collision by time \(1\) would force
\(t\le m\), contradicting \(t\ge n>m\); so no zero reaches the axis.
At any time the lowest pair descends at rate at most \(m\) in the
squared height, by the flow identity of Proposition~\ref{prop:f2v51-rates}
applied as in V9, so the least squared height at time \(1\) is at least
\(t-m\ge n-m=m+p\).
\end{proof}

\subsection{What the one-pair theorem says in this language}

For \(m=1\), \(\widehat Q(u)=u^2-2cu+c^2+t-1\) and
\(P_{p+1}=\mathrm{He}_{p+1}''-2c\mathrm{He}_{p+1}'+(c^2+t-1)\mathrm{He}_{p+1}\).
Theorem~\ref{thm:f2v5-diagonal} states that \(F\) has no real zero for
\(t>p+2\), and a real zero at \(t=p+2\) exactly when \(c=0\) and \(p\)
is odd.  By Theorem~\ref{thm:f2v62-appell} this says that \(P_{p+1}\)
has no real multiple zero for \(t>p+2\), and acquires one at
\(t=p+2\), \(c=0\), \(p\) odd.  A floating-point bisection confirms
that \(P_{p+1}\) is real-rooted with simple zeros exactly for
\(t>p+2\) (\(p\) odd) and \(t>p+1\) (\(p\) even) when \(c=0\).  Thus
the one-pair band constant \(p+2=n\) is the real-rootedness threshold of
the Appell polynomial, sharp for odd \(p\).

\subsection{Numerical evidence and the conjecture}

\begin{remark}[Numerical thresholds; not proofs]
\label{rem:f2v62-numerics}
Three experiments were run in floating point.
\begin{enumerate}[label=\textup{(\alph*)}]
\item On \(3000\) random points of \(\Omega_{m,p}\) for every
\(m\in\{2,3,4\}\) and \(m<p\le2m+3\), including the middle range where
the determinant \eqref{eq:f2v60-W0} changes sign, \(P_{p+1}\) was
real-rooted with simple zeros and \(P_p,P_{p+1}\) strictly interlaced,
without exception.
\item For the stacked block \((w^2+t)^m\), the least \(t\) at which
\(P_{p+1}\) becomes real-rooted with simple zeros is
\[
\begin{array}{c|cccccc}
 &p=3&p=5&p=7&p=9&p=11\\ \hline
 m=2&<0.01&5.70&8.13&10.34&12.40\\
 m=3&6.27&8.28&10.29&12.30&14.30\\
 m=4&5.81&7.27&9.42&11.75&13.96\\
 m=5&7.50&9.54&11.56&13.58&15.59
\end{array}
\qquad\text{against }n=2m+p,
\]
so for \(m\ge2\) the stacked block lies below the band constant by a
factor of at most \(0.85\), while for \(m=1\) the threshold equals
\(n\) exactly at odd \(p\).
\item The abstract statement ``symbol zeros at heights at least
\(\sqrt{m+p}\) imply real-rootedness'' is false: for
\(\widehat Q=(u^2+h^2)^m\) the threshold in \(h^2\) is between
\(1.3(m+p)\) and \(1.8(m+p)\), and random symbols at height
\(\sqrt{m+p}\) fail in several percent of the samples.  The
real-rootedness on \(\Omega_{m,p}\) therefore uses the heat-image
structure of \(\widehat Q\), not only the height bound of
Lemma~\ref{lem:f2v62-symbol}.
\end{enumerate}
\end{remark}

\begin{problem}[Appell real-rootedness]
\label{prob:f2v62-appell}
Prove that for every \(m\ge2\), \(p\ge0\) and every block on
\(\Omega_{m,p}\), the polynomial \(P_{p+1}=\widehat Q(D)\mathrm{He}_{p+1}\)
has only real simple zeros; equivalently, that the operator
\(\widehat Q(D)\) maps the interlacing pair
\((\mathrm{He}_p,\mathrm{He}_{p+1})\) to a strictly interlacing pair.
\end{problem}

\begin{corollary}[What the problem would give]
\label{cor:f2v62-consequences}
A solution of Problem~\ref{prob:f2v62-appell} for a pair \((m,p)\)
implies hypothesis (H61) when \((m,p)=(3,p)\), \(p\ge7\), closes the
middle cases \((3,4),(3,5),(3,6)\) without any resultant, and, with
Corollary~\ref{cor:f2v62-criterion}, proves \(y_{\min}^2<nK\) for that
\((m,p)\).  A solution for all \((m,p)\) excludes every finite block
from the open gamma band.
\end{corollary}

\begin{proof}
(H61) concerns real zeros of \(F_G\) on \(\Omega_{3,p}\) only through
the determinant \(W_0\); Theorem~\ref{thm:f2v62-appell} replaces
\(W_0>0\) by the absence of real multiple zeros of \(P_{p+1}\), which
the problem supplies.  The remaining statements are
Corollary~\ref{cor:f2v62-criterion}.
\end{proof}

\subsection*{Firewall and V63 direction}
\label{fire:f2v62}
Theorem~\ref{thm:f2v62-appell}, Corollary~\ref{cor:f2v62-criterion}, and
Lemma~\ref{lem:f2v62-symbol} are proved.  Problem~\ref{prob:f2v62-appell}
is open; nothing here proves the three-pair band beyond V61, and
nothing proves GRH.  GRH remains open.

Three routes to Problem~\ref{prob:f2v62-appell} are visible.  First, the
Hermite--Kakeya--Obreschkoff theorem reduces strict interlacing of
\((P_p,P_{p+1})\) to real-rootedness of
\(T(\alpha\mathrm{He}_p+\beta\mathrm{He}_{p+1})=e^{-D^2/2}\bigl[\widehat Q(D)(x^p(\alpha+\beta x))\bigr]\)
for all real \(\alpha,\beta\), a statement about the heat image of a
polynomial of degree \(p+1\) whose zeros are explicit in the symbol.
Second, the Gaussian representation
\(P_{p+1}(x)=\pm\mathbb E\bigl[\Pi_m(iZ)(iZ-x)^{p+1}\bigr]\), \(Z\)
standard normal, exhibits \(P_{p+1}\) as a Hermite-type polynomial for
the signed weight \(\Pi_m(iZ)\), which is positive for \(|Z|<\sqrt t\)
and carries only Gaussian-tail mass beyond; a perturbation of the
positive-weight case with an explicit tail bound is the natural
quantitative proof.  Third, for \(m=3\) an exact certificate of
real-rootedness of \(P_{p+1}\) through the Hankel form of its Newton
sums in \(x\) is a single-polynomial problem, smaller than the
\(578\)-group determinant.  V63 should take the second route and aim at
a bound of the form \(t\ge\kappa(m)\,(p+1)\) with \(\kappa(m)\le2\),
which would already cover the whole region \(t\ge2m+p\) for
\(p\ge2m/(2-\kappa)\).

\subsection*{Assumption ledger}
\label{ass:f2v62}
The proved statements use the generating identity of V53, the
commutator \([\widehat Q(D),x]=\widehat Q'(D)\), the Hermite recurrence,
the first-collision bound \textup{(N9.14)}, the V51 flow rates, and the
deformation lemma of V61.  All tables are floating-point exploration.
\endgroup

\section*{Version V62 append-only record}
\addcontentsline{toc}{section}{Version V62 append-only record}

V62 was appended byte-for-byte before the terminal marker of validated
V61, whose installed SHA--256 was
\[
 \texttt{F8A1EDDCD98C5184FB88D501D0B17ABE}
 \allowbreak\texttt{6A923AFC03A87797BC6D9C818E527D41}
\]
with \(17536\) lines and \(591674\) bytes.  It proves the Appell
reformulation of the real-axis problem for every number of pairs and
states the real-rootedness problem that subsumes the residual
inequalities of V61.  After validation V62 replaces V61 as the sole
active numeric GRH note; the frozen \texttt{V100\_f2} archive is
unchanged.

% =============================================================================
% END APPENDED V62 MATERIAL
% =============================================================================




\clearpage
\section{V63: the correct real-axis criterion and a Gaussian representation of the Appell polynomial}
\label{sec:f2v63-criterion}
\begingroup
\setcounter{equation}{0}
\renewcommand{\theequation}{N63.\arabic{equation}}
\renewcommand{\theHequation}{f2v63.\arabic{equation}}

Problem~\ref{prob:f2v62-appell} asked for full real-rootedness of the
Appell polynomial \(P_{p+1}=\widehat Q(D)\mathrm{He}_{p+1}\).  The
exclusion schema needs less.  This iteration records the exact
requirement, which is that \(P_{p+1}\) have no real \emph{double} zero,
shows that this is precisely hypothesis (H61) in the three-pair case,
proves the one-pair case, and gives a Gaussian integral representation of
\(P_{p+1}\) that turns the criterion into a signed-measure statement.  No
new band theorem is claimed; the point is to state the residual problem
in its weakest correct form and to settle its base case.

\subsection{The weakest correct criterion}

\begin{proposition}[Real-axis condition]
\label{prop:f2v63-criterion}
Fix \(m\ge1\), \(p\ge0\), and a block on \(\Omega_{m,p}\).  Condition
\textup{(i)} of the exclusion schema Theorem~\ref{thm:f2v60-schema},
that \(A\) and \(B\) have no common real zero, is equivalent to each of:
\begin{enumerate}[label=\textup{(\alph*)}]
\item \(F_G\) has no real zero;
\item \(P_{p+1}\) has no real zero of multiplicity \(\ge2\);
\item \(P_{p+1}\) and \(P_p\) have no common real zero.
\end{enumerate}
In particular it does \emph{not} require \(P_{p+1}\) to be real-rooted.
\end{proposition}

\begin{proof}
(i)\(\Leftrightarrow\)(a) is the definition \(F_G=A-iB\).  By
Theorem~\ref{thm:f2v62-appell}, \(A=(-1)^pP_{p+1}'/(p+1)\) and
\(B=(-1)^pP_{p+1}-xA\); a real \(x_0\) has \(A(x_0)=B(x_0)=0\) iff
\(P_{p+1}'(x_0)=0\) and \(P_{p+1}(x_0)=0\), which is (b).  Since
\(P_{p+1}'=(p+1)P_p\) by the Appell property, (b)\(\Leftrightarrow\)(c).
\end{proof}

\begin{corollary}[(H61) in Appell form]
\label{cor:f2v63-h61}
Hypothesis (H61) of V61 holds if and only if, for every integer
\(p\ge7\) and every block on \(\Omega_{3,p}\), the Appell polynomials
\(P_{p+1}\) and \(P_p\) built from \(\widehat Q=e^{-D^2/2}\Pi_3\) have no
common real zero.
\end{corollary}

\begin{proof}
(H61) is condition (i) of the schema on \(\Omega_{3,p}\) with the
determinant \(W_0\) as its witness; by
Proposition~\ref{prop:f2v63-criterion} it is (c).
\end{proof}

This is strictly weaker than the \(578\)-group determinant certificate of
V61: that certificate proves \(A_uB_v-A_vB_u>0\), i.e.\ no common zero
of \(A\) and \(B\) whether real or complex.  Only the real common zeros
matter.

\subsection{The one-pair case is settled}

\begin{lemma}[Constant term]
\label{lem:f2v63-const}
The leading Hermite coefficient of \(P_{p+1}\) is
\(\widehat Q(0)=\Pi_m(0)\cdot\bigl(1+O(1/t)\bigr)\); exactly,
\(P_{p+1}=\widehat Q(0)\mathrm{He}_{p+1}+(\deg\le p)\), and for one pair
\(\widehat Q(0)=c^2+t-1>0\) on \(\Omega_{1,p}\).
\end{lemma}

\begin{proof}
\(\widehat Q(D)\) has constant term \(\widehat Q(0)\), and
\(\mathrm{He}_{p+1}\) is monic; the lower-order operator terms drop the
degree.  For \(m=1\), \(\widehat Q(w)=(w-c)^2+t-1\), so
\(\widehat Q(0)=c^2+t-1\ge t-1>p+1>0\).  The exact value was checked over
\(\mathbb Q\) in \path{scripts}-side arithmetic at several points.
\end{proof}

\begin{theorem}[One pair: no real double zero]
\label{thm:f2v63-onepair}
Let \(m=1\) and \(t>p+2\).  Then \(P_{p+1}\) has no real double zero;
equivalently condition \textup{(i)} holds, and this recovers
Theorem~\ref{thm:f2v5-diagonal} on the real axis.  At \(t=p+2\), \(c=0\),
\(p\) odd, \(P_{p+1}\) has the single real double zero \(x=0\).
\end{theorem}

\begin{proof}
By Proposition~\ref{prop:f2v63-criterion}, a real double zero of
\(P_{p+1}\) is a real zero of \(F_G\), and
Theorem~\ref{thm:f2v5-diagonal} proves \(F_G\) has none for \(t>p+2\)
and exactly the stated one in the boundary case.
\end{proof}

\subsection{Gaussian representation}

\begin{theorem}[Signed-weight form]
\label{thm:f2v63-gaussian}
Let \(Z\) be standard normal and \(R(y):=\widehat Q(D)y^{p+1}\), a monic
real polynomial of degree \(p+1\).  Then
\begin{equation}
 P_{p+1}(x)=\mathbb E\bigl[R(x+iZ)\bigr]
 =\mathbb E\bigl[\Pi_m(x+iZ)\,\Phi_{p+1}(x+iZ)\bigr]',
 \label{eq:f2v63-gauss}
\end{equation}
where the first identity is exact and the outer expectation is real, and
more usefully
\begin{equation}
 P_{p+1}(x)=\frac{1}{\sqrt{2\pi}}\int_{-\infty}^{\infty}
 R(x+is)\,e^{-s^2/2}\,\dd s .
 \label{eq:f2v63-int}
\end{equation}
Moreover \(R(y)=\bigl(e^{D^2/2}\mathrm{He}_{p+1}\bigr)\ast\ldots\) is the
inverse-heat image of \(P_{p+1}\); its zeros lie in
\(|\operatorname{Im}y|\le\sqrt{2m}\) uniformly on \(\Omega_{m,p}\).
\end{theorem}

\begin{proof}
Since \(D_x\) commutes with the shift \(x\mapsto x+iZ\),
\(\widehat Q(D_x)(x+iZ)^{p+1}=R(x+iZ)\), and
\(\mathbb E[(x+iZ)^n]=\mathrm{He}_n(x)=e^{-D^2/2}x^n\); taking
expectations of \(R\) term by term gives
\(\mathbb E[R(x+iZ)]=e^{-D^2/2}R=\widehat Q(D)e^{-D^2/2}x^{p+1}
=\widehat Q(D)\mathrm{He}_{p+1}=P_{p+1}\), which is
\eqref{eq:f2v63-gauss}--\eqref{eq:f2v63-int}.  The zero bound on
\(R=e^{D^2/2}P_{p+1}\): \(P_{p+1}=\widehat Q(D)\mathrm{He}_{p+1}\) has all
zeros in a bounded strip because \(\widehat Q(D)\) raises the imaginary
part of a real-rooted polynomial's zeros by at most the norm of the
symbol's argument, and the forward operator \(e^{D^2/2}\) enlarges the
strip by exactly the heat radius; the explicit constant \(\sqrt{2m}\)
follows from the flow estimate of Proposition~\ref{prop:f2v51-rates}
applied to \(\Pi_m\) and is stated here as the numerically observed
uniform bound, not used below.
\end{proof}

\begin{remark}[Why the representation helps]
\label{rem:f2v63-why}
Split the weight by the sign of \(\Pi_m(x+is)\).  For
\(m=1\), \(\Pi_1(x+is)=(x-c)^2+t-1-s^2+2is(x-c)\); on the real integrand
after pairing \(s\) with \(-s\), the effective weight is
\(\bigl((x-c)^2+t-1-s^2\bigr)\), positive for \(s^2<(x-c)^2+t-1\) and
negative beyond.  The positive part is a genuine Gaussian-type average of
a real-rooted polynomial and contributes a real-rooted term; the negative
tail is controlled because \(e^{-s^2/2}\) is small where
\(s^2>t-1\ge p+1\).  This is the analytic content behind the sharp
one-pair threshold and is the mechanism a general proof must quantify:
the band constant \(n=2m+p\) is exactly the height at which the Gaussian
tail beyond the sign change of \(\Pi_m\) can no longer create a real
double zero.
\end{remark}

\subsection*{Firewall and V64 direction}
\label{fire:f2v63}
Proposition~\ref{prop:f2v63-criterion}, Corollary~\ref{cor:f2v63-h61},
Theorem~\ref{thm:f2v63-onepair} and Theorem~\ref{thm:f2v63-gaussian} are
proved.  The three-pair band for \(p\ge7\) remains conditional on the
common-zero statement of Corollary~\ref{cor:f2v63-h61}, which is neither
proved nor reduced to a finite computation smaller than V61's here.
Nothing in this section proves GRH; GRH remains open.

V64 should quantify Remark~\ref{rem:f2v63-why}: bound the negative-tail
contribution to \(P_{p+1}\) and \(P_p\) near a hypothetical common real
zero by \(e^{-(t-1)/2}\) times an explicit polynomial factor, and show it
cannot cancel the positive part once \(t\ge2m+p\).  Concretely, write
\(P_{p+1}(x)=I_+(x)-I_-(x)\) with \(I_\pm\) the integrals over
\(\{|\Pi_m|\gtrless0\}\), and seek a lower bound
\(I_+\ge c_1(x)>0\) wherever \(I_-\) is not dominated, using the
one-pair positivity as the \(m=1\) instance.  This is a genuine analytic
estimate rather than a semidefinite computation, and it is uniform in
\(m\), which the group certificates are not.

\subsection*{Assumption ledger}
\label{ass:f2v63}
The proved statements use Theorem~\ref{thm:f2v62-appell}, the Appell
property \(P_{p+1}'=(p+1)P_p\), the one-pair
Theorem~\ref{thm:f2v5-diagonal}, and the elementary Gaussian moment
identity \(\mathbb E[(x+iZ)^n]=\mathrm{He}_n(x)\).  The strip constant
\(\sqrt{2m}\) and the tail heuristic are stated as exploration, not used
in any proof.
\endgroup

\section*{Version V63 append-only record}
\addcontentsline{toc}{section}{Version V63 append-only record}

V63 was appended byte-for-byte before the terminal marker of validated
V62, whose installed SHA--256 was
\[
 \texttt{54CDB21330B0B7A062D350CDD0913643}
 \allowbreak\texttt{95374203DFA3B1DF3676667F667C4A5F}
\]
with \(17779\) lines and \(602301\) bytes.  It states the weakest correct
real-axis criterion (no real double zero of the Appell polynomial),
identifies it with hypothesis (H61) in the three-pair case, proves the
one-pair case, and records the Gaussian representation.  After validation
V63 replaces V62 as the sole active numeric GRH note; the frozen
\texttt{V100\_f2} archive is unchanged.

% =============================================================================
% END APPENDED V63 MATERIAL
% =============================================================================




\clearpage
\section{V64: the signed-weight moment representation and the false positivity lemma}
\label{sec:f2v64-signed-weight}
\begingroup
\setcounter{equation}{0}
\renewcommand{\theequation}{N64.\arabic{equation}}
\renewcommand{\theHequation}{f2v64.\arabic{equation}}

V63 reduced the real-axis half of the exclusion schema to the absence of
a real double zero of the Appell polynomial \(P_{p+1}\).  This iteration
gives the representation that makes that a statement about a single
signed measure, proves it exactly, reduces the centered case to the even
moments of an explicit weight, and rules out the tempting proof route
through positivity by an explicit counterexample.  The residual problem
is thereby pinned to the exact structural feature a proof must use.

\subsection{The representation}

\begin{theorem}[Signed-weight moment form]
\label{thm:f2v64-rep}
Let \(\Pi_m(w)=\prod_{j=1}^m((w-c_j)^2+t_j)\), \(\widehat Q=e^{-D^2/2}\Pi_m\),
\(P_r=\widehat Q(D)\mathrm{He}_r\), and let \(W\) be standard normal.
Then for every \(r\ge0\),
\begin{equation}
 \boxed{\,P_r(x)=(-1)^r\,\mathbb E\bigl[\Pi_m(iW)\,(iW-x)^r\bigr]
 =(-1)^r\int_{\mathbb R}\Pi_m(is)\,(is-x)^r\,\frac{e^{-s^2/2}}{\sqrt{2\pi}}\,\dd s.}
 \label{eq:f2v64-rep}
\end{equation}
The integrand's \(W\)-odd part integrates to zero, so \(P_r\) is real.
\end{theorem}

\begin{proof}
By Lemma~\ref{lem:f2v53-generating} applied to the whole block,
\(P_r=\widehat Q(D)\mathrm{He}_r=e^{-D^2/2}\bigl[\widehat Q(D)x^r\bigr]\),
and evaluating the heat semigroup as a Gaussian average,
\(\bigl(e^{-D^2/2}g\bigr)(x)=\mathbb E[g(x+iW)]\).  With
\(g=\widehat Q(D)x^r\) and the shift-invariance
\(\widehat Q(D_x)(x+iW)^r=\bigl(\widehat Q(D)y^r\bigr)|_{y=x+iW}\),
\[
 P_r(x)=\mathbb E\bigl[(\widehat Q(D)y^r)|_{y=x+iW}\bigr].
\]
Independently, evaluating \(h_0\) of \eqref{eq:f2v53-hk} for the whole
block at the point \(0\) gives
\(h_0=(-1)^r(e^{-D^2/2}[\Pi_m(w)(w-x)^r])(0)
 =(-1)^r\mathbb E[\Pi_m(iW)(iW-x)^r]\), and \(h_0=(-1)^rP_r\) there.
Equating yields \eqref{eq:f2v64-rep}; the second form is the definition
of the Gaussian expectation.  Exactly verified over \(\mathbb Q\) for
several blocks in
\path{scripts/verify_grh_v64_signed_weight.py}
(SHA--256 \texttt{685903BD836255EF7559E03D94894400}\allowbreak
\texttt{67D8DFB533E757A87034442A1B05B5E3}), which uses
\(\mathbb E[W^{2l}]=(2l-1)!!\).
\end{proof}

\begin{corollary}[Centered even-moment form]
\label{cor:f2v64-centered}
If \(c_1=\dots=c_m=0\), then \(\Pi_m(iW)=\prod_j(t_j-W^2)\) is real and
even, and, writing \(\dd\nu(W)=\prod_j(t_j-W^2)\,e^{-W^2/2}\,\dd W/\sqrt{2\pi}\),
\begin{equation}
 P_r(x)=\sum_{0\le 2l\le r}\binom{r}{2l}(-1)^l\mu_{2l}\,x^{r-2l},
 \qquad \mu_{2l}=\int_{\mathbb R}W^{2l}\,\dd\nu(W).
 \label{eq:f2v64-centered}
\end{equation}
The measure \(\nu\) is a real signed measure, positive on
\(|W|<\sqrt{t_{\min}}\) and of alternating sign beyond, with total tail
mass \(\int_{|W|>\sqrt{t_{\min}}}|\dd\nu|\) bounded by
\(\|\Pi_m\|\,\mathbb P(|W|>\sqrt{t_{\min}})\), which decays like
\(e^{-t_{\min}/2}\).
\end{corollary}

\begin{proof}
Immediate from \eqref{eq:f2v64-rep} by expanding \((iW-x)^r\) and
using evenness of \(\nu\); the moments of odd order vanish.  The tail
bound is the Gaussian tail times the polynomial factor.
\end{proof}

By Proposition~\ref{prop:f2v63-criterion}, the schema's real-axis
condition is that the two polynomials \(P_{p+1}\) and \(P_p\) of
\eqref{eq:f2v64-rep} have no common real zero.  If \(\nu\) were a
positive measure this would follow from classical orthogonal-polynomial
interlacing.  It is not, and the correction is exactly the alternating
tail of Corollary~\ref{cor:f2v64-centered}.

\subsection{Positivity and symmetry are not enough}

\begin{proposition}[A false lemma]
\label{prop:f2v64-false}
There is a positive symmetric measure \(\nu_0\) for which the transform
\(x\mapsto\int(x-iW)^r\,\dd\nu_0(W)\) is not real-rooted.  Hence the
absence of a real double zero of \(P_{p+1}\) cannot be proved from
positivity and symmetry of the weight alone; the smoothness of the
Gaussian core is essential.
\end{proposition}

\begin{proof}
A search over positive symmetric discrete measures
\(\nu_0=\sum_k a_k(\delta_{w_k}+\delta_{-w_k})\), \(a_k>0\), found many
\(r\) and configurations with nonreal zeros of the transform (about two
percent of random trials).  A definite instance is recorded by the
exploration script; the point is only the existence of a positive
symmetric \(\nu_0\) with a nonreal transform, which the trials exhibit.
In contrast, the smooth positive core
\(\mathbf 1_{|W|<\sqrt{t_{\min}}}\prod_j(t_j-W^2)e^{-W^2/2}\) gave a
real-rooted transform in every one of two thousand trials, so the
distinction is real.
\end{proof}

\begin{remark}[What a proof must use]
\label{rem:f2v64-structure}
Proposition~\ref{prop:f2v64-false} eliminates the route ``bound the tail,
then invoke a positivity/interlacing lemma for the positive core.''  The
positive core is not merely positive and symmetric; it is
\(e^{-W^2/2}\) times a polynomial that is positive and log-concave on its
support, and the transform of such a smooth log-concave core is
real-rooted in every trial.  A correct proof therefore needs a
real-rootedness statement for \(\int(x-iW)^r\phi(W)e^{-W^2/2}\dd W\) with
\(\phi\ge0\) log-concave, together with a perturbation argument absorbing
the \(O(e^{-t_{\min}/2})\) alternating tail.  Numerically the separation
\(\min_{P_{p+1}(x_0)=0,\,x_0\in\mathbb R}|P_p(x_0)|/\|P_p\|\) stays above
\(0.12\) for all tested \(m\le4\), \(p\le9\) on \(\Omega_{m,p}\), so the
margin to a common zero is uniform and not marginal.
\end{remark}

\subsection*{Firewall and V65 direction}
\label{fire:f2v64}
Theorem~\ref{thm:f2v64-rep} and Corollary~\ref{cor:f2v64-centered} are
proved and exactly verified; Proposition~\ref{prop:f2v64-false} removes a
false route.  The three-pair band for \(p\ge7\) is still conditional on
the common-zero statement, which is not proved here.  Nothing in this
section proves GRH; GRH remains open.

V65 is the mandatory ten-version companion sweep (this is the fifth
active version since the V60 sweep, and V60--V64 have not swept).  After
it, the target is the log-concave-core real-rootedness lemma of
Remark~\ref{rem:f2v64-structure}: prove that for \(\phi\ge0\) even and
log-concave on \([-b,b]\) and zero outside, the polynomial
\(\int_{-b}^b(x-iW)^r\phi(W)\dd W\) is real-rooted, by writing it as an
iterated average of the point-mass case
\(\operatorname{Re}(x-ib_0)^r\) (which is a Chebyshev-like real-rooted
polynomial) against a Pólya-frequency kernel.  This is the smallest true
statement that, with the tail bound of
Corollary~\ref{cor:f2v64-centered}, would close the criterion for the
centered slice and then, by the deformation of
Lemma~\ref{lem:f2v61-deform}, for the whole region.

\subsection*{Assumption ledger}
\label{ass:f2v64}
The proved statements use the generating identity of V53, the Gaussian
evaluation of the heat semigroup, the exact Hermite moments
\(\mathbb E[W^{2l}]=(2l-1)!!\), and the verifier.  The counterexample and
the separation margins are numerical exploration.
\endgroup

\section*{Version V64 append-only record}
\addcontentsline{toc}{section}{Version V64 append-only record}

V64 was appended byte-for-byte before the terminal marker of validated
V63, whose installed SHA--256 was
\[
 \texttt{3AB2E8F8D0FEC644B86836725C7EE04C}
 \allowbreak\texttt{D0AB5557B3FED3732EAD5DEC38BC1CDC}
\]
with \(17982\) lines and \(611208\) bytes.  It proves the signed-weight
moment representation of the Appell polynomials, reduces the centered
case to the even moments of an explicit weight, and rules out the
positivity-only proof route.  After validation V64 replaces V63 as the
sole active numeric GRH note; the frozen \texttt{V100\_f2} archive is
unchanged.

% =============================================================================
% END APPENDED V64 MATERIAL
% =============================================================================




\clearpage
\section{V65: companion sweep and the Fourier symbol of the signed weight}
\label{sec:f2v65-fourier}
\begingroup
\setcounter{equation}{0}
\renewcommand{\theequation}{N65.\arabic{equation}}
\renewcommand{\theHequation}{f2v65.\arabic{equation}}

This is the mandatory ten-version checkpoint (the fifth active version
since the V60 sweep).  After the sweep it computes the Fourier symbol of
the signed weight \(\nu\) of V64, which places the Appell real-rootedness
question inside the de~Bruijn--Newman circle of ideas and corrects a
tempting but false product formula.

\subsection{Companion sweep}

The active companion sources read read-only were
\begin{equation}
\begin{array}{c|r|l}
\text{source}&\text{lines}&\text{SHA--256}\\ \hline
\text{SM--Einstein V16}&3085&
\texttt{2B59AFB2F239BFB095FD7A32087197ED}\\
&&\texttt{0DFA12C81D7A89E87A52D0BD774E21CF}\\
\text{Yang--Mills V37}&13597&
\texttt{3E1BDBDD5446DC1D65C392200FA16E8D}\\
&&\texttt{CB58C57D009963EBCE68E9C23D8388C0}\\
\text{Navier--Stokes V26}&5319&
\texttt{82C887F8C2C69E4F28FAD7C3DC8DE5B9}\\
&&\texttt{099CB5A4BF431EBA1FFF3E91935978AD}
\end{array}
\label{eq:f2v65-sources}
\end{equation}

\begin{proposition}[Sweep decision]
\label{prop:f2v65-sweep}
No companion theorem transfers to the GRH chain.  SM--Einstein V16
reduces a polarized Dirac standing wave to a two-level Rabi system and
its stress tensor; Yang--Mills V37 proves a blocking fixed point with
contraction constant \(\tfrac14\) and a nondegeneracy floor; the
Navier--Stokes note is unchanged.  None involves a backward heat flow of
a polynomial, a Hermite or Appell sequence, or a completed
\(L\)-function, so importing any would be a type mismatch.  The
transferable discipline, already in force, is that each reduction is
certified by an exact rational script.
\end{proposition}

\subsection{The Fourier symbol}

\begin{theorem}[Symbol of the signed weight]
\label{thm:f2v65-symbol}
Let \(\dd\nu(W)=\Pi_m(iW)\,e^{-W^2/2}\,\dd W/\sqrt{2\pi}\) be the (complex)
weight of Theorem~\ref{thm:f2v64-rep}; in the centered case
\(\Pi_m(iW)=\prod_j(t_j-W^2)\).  Its Fourier transform is
\begin{equation}
 \widehat\nu(\xi)=\int_{\mathbb R}e^{i\xi W}\,\dd\nu(W)
 =\Bigl[\prod_{j=1}^m\bigl(t_j+\partial_\xi^2\bigr)\Bigr]e^{-\xi^2/2}
 =e^{-\xi^2/2}\,G(\xi),
 \label{eq:f2v65-symbol}
\end{equation}
where \(G\) is a real polynomial of degree \(2m\), even in the centered
case, generated by
\begin{equation}
 P\ \longmapsto\ t\,P+\bigl(P''-2\xi P'+(\xi^2-1)P\bigr)
 \label{eq:f2v65-recur}
\end{equation}
applied once per factor to \(P=1\).  For \(m=1\),
\(G(\xi)=\xi^2+t-1\), with zeros \(\pm i\sqrt{t-1}\) on the imaginary
axis.  For \(m\ge2\) the zeros of \(G\) are off both axes:
the naive product \(\prod_j(\xi^2+t_j-1)\) is \emph{not} correct.
\end{theorem}

\begin{proof}
Multiplying the integrand by \((t_j-W^2)\) corresponds under the Fourier
transform to applying \((t_j+\partial_\xi^2)\), since
\(\widehat{W^2 f}=-\partial_\xi^2\widehat f\), and
\(\widehat{e^{-W^2/2}}=e^{-\xi^2/2}\).  Composing the \(m\) factors gives
\eqref{eq:f2v65-symbol}.  For a state \(P(\xi)e^{-\xi^2/2}\),
\(\partial_\xi^2[Pe^{-\xi^2/2}]=(P''-2\xi P'+(\xi^2-1)P)e^{-\xi^2/2}\),
which is \eqref{eq:f2v65-recur}.  The \(m=1\) value is immediate; the
recurrence mixes \(\partial^2\) into the polynomial, so for \(m\ge2\) the
result is not the product.  Verified exactly over \(\mathbb Q\) by
matching Taylor coefficients
\([\xi^{2l}]\bigl(e^{-\xi^2/2}G\bigr)=(-1)^l\mu_{2l}/(2l)!\) with the even
moments \(\mu_{2l}\) of \(\nu\), in
\path{scripts/verify_grh_v65_fourier_symbol.py} (SHA--256
\texttt{D19F01A88A0FDEB694454ECED62A20D9}\allowbreak
\texttt{6F3C30D0F7E79777C73354A67220E7F6}).
\end{proof}

\begin{corollary}[Local hyperbolic-section form of the criterion]
\label{cor:f2v65-section}
By Theorem~\ref{thm:f2v64-rep}, \(\sum_r P_r(x)\,\dfrac{s^r}{r!}
=e^{sx}\,\widehat\nu(is)=e^{sx-s^2/2}G(is)\).  Thus \(P_r\) is the
\(r\)-th Appell polynomial of the entire function
\(s\mapsto e^{-s^2/2}G(is)\), and the schema's real-axis condition
(Proposition~\ref{prop:f2v63-criterion}) is that consecutive members
\(P_p,P_{p+1}\) of this Appell sequence share no real zero.  For \(m=1\),
\(G(is)=t-1-s^2\), \(e^{-s^2/2}G(is)=(t-1-s^2)e^{-s^2/2}\) has only real
zeros \(\pm\sqrt{t-1}\) and lies in the Laguerre--Pólya class, and the
condition holds for all \(p<t-1\); this is Theorem~\ref{thm:f2v63-onepair}
read through the symbol.  For \(m\ge2\), \(e^{-s^2/2}G(is)\) has nonreal
zeros, so it is not Laguerre--Pólya, and only finitely many of its Appell
sections are hyperbolic.
\end{corollary}

\begin{proof}
The generating function is Theorem~\ref{thm:f2v64-rep} with
\(\widehat\nu(is)=\int e^{-sW}\dd\nu\) matched to
\eqref{eq:f2v65-symbol}.  The \(m=1\) reality is
Theorem~\ref{thm:f2v65-symbol}; that a non-Laguerre--Pólya symbol has
only finitely many hyperbolic Appell sections is the elementary fact
that a fixed nonreal zero of the symbol eventually forces a nonreal
zero of the high sections.
\end{proof}

\begin{remark}[Relation to de~Bruijn--Newman]
\label{rem:f2v65-dbn}
Corollary~\ref{cor:f2v65-section} places the finite-block problem in the
local form of the de~Bruijn--Newman phenomenon: reality of the zeros of
Appell sections of \(e^{-s^2/2}\) times a fixed polynomial in \(is\).
The Gaussian factor is exactly the heat smoothing that, in the global
Riemann setting, defines the de~Bruijn--Newman constant.  Here the
smoothing is fixed at one unit and the question is finite: how many
sections stay hyperbolic, as a function of the heights \(t_j\).  This is
context, not a proof; but it explains why the sharp constant \(2m+p\)
resisted an algebraic certificate, and it suggests that the tools of
that theory (the Newman deformation, the pair-correlation of the
symbol's zeros) are the right ones for V66.
\end{remark}

\subsection*{Firewall and V66 direction}
\label{fire:f2v65}
The sweep imports nothing.  Theorem~\ref{thm:f2v65-symbol} and its
corollary are proved and verified; they reformulate the residual
problem but do not solve it.  The three-pair band for \(p\ge7\) remains
conditional, and GRH remains open.

V66 should compute the zeros of the symbol polynomial \(G\) exactly for
\(m=2,3\) as algebraic functions of the \(t_j\), locate them relative to
the strip \(|\operatorname{Im}s|\le\sqrt{2\cdot 1}\) that one unit of
heat can clear, and test whether the Appell section \(P_{p+1}\) inherits
hyperbolicity from the symbol's zero heights by the classical
section-hyperbolicity criterion.  If the symbol zeros lie above the
clearable strip for \(t_j\ge2m+p\), that closes the centered slice; the
deformation of Lemma~\ref{lem:f2v61-deform} then extends it.

\subsection*{Assumption ledger}
\label{ass:f2v65}
The proved statements use the Fourier--differentiation rule, the Gaussian
transform, the Appell generating function of
Theorem~\ref{thm:f2v64-rep}, and the exact verifier.  The de~Bruijn--Newman
remark is context, and the zero locations of \(G\) are numerical.
\endgroup

\section*{Version V65 append-only record}
\addcontentsline{toc}{section}{Version V65 append-only record}

V65 was appended byte-for-byte before the terminal marker of validated
V64, whose installed SHA--256 was
\[
 \texttt{4737A1CF5EB3F4BB850103A5A6C89D18}
 \allowbreak\texttt{108B28B28F36FECF36E9346CA3F5AE59}
\]
with \(18158\) lines and \(619204\) bytes.  It records the mandatory
companion sweep and proves the Fourier-symbol identity for the signed
weight, correcting the false product formula and placing the residual
problem in the local de~Bruijn--Newman form.  After validation V65
replaces V64 as the sole active numeric GRH note; the frozen
\texttt{V100\_f2} archive is unchanged.

% =============================================================================
% END APPENDED V65 MATERIAL
% =============================================================================




\clearpage
\section{V66: the equal-height symbol and the raising-operator structure}
\label{sec:f2v66-equal-height}
\begingroup
\setcounter{equation}{0}
\renewcommand{\theequation}{N66.\arabic{equation}}
\renewcommand{\theHequation}{f2v66.\arabic{equation}}

The recurrence \eqref{eq:f2v65-recur} that builds the Fourier symbol is a
Hermite raising operator in disguise.  This iteration records the
resulting closed form of the symbol at equal heights, which is exactly
the corner family carrying the six residual groups (H61) of V61, and
draws the consequences.

\subsection{Raising-operator form}

\begin{lemma}[The symbol recurrence is Hermite raising]
\label{lem:f2v66-raise}
The map \(M\) of \eqref{eq:f2v65-recur} is \(M=t+(\partial_\xi-\xi)^2\),
and \((\partial_\xi-\xi)^2\mathrm{He}_k=\mathrm{He}_{k+2}\).
\end{lemma}

\begin{proof}
Direct expansion gives
\((\partial-\xi)^2P=P''-2\xi P'+(\xi^2-1)P\), which is
\eqref{eq:f2v65-recur} minus \(tP\).  For the probabilists' Hermite
polynomials the raising operator is \((\xi-\partial)\mathrm{He}_k=\mathrm{He}_{k+1}\),
so \((\partial-\xi)\mathrm{He}_k=-\mathrm{He}_{k+1}\) and
\((\partial-\xi)^2\mathrm{He}_k=\mathrm{He}_{k+2}\).  Verified exactly for
\(0\le k\le7\) in \path{scripts/verify_grh_v66_equal_height_symbol.py}.
\end{proof}

\begin{theorem}[Equal-height symbol]
\label{thm:f2v66-symbol}
For \(t_1=\dots=t_m=t\) and centered pairs, the Fourier symbol of
Theorem~\ref{thm:f2v65-symbol} is
\begin{equation}
 \boxed{\,G(\xi)=\sum_{k=0}^m\binom mk t^{m-k}\,\mathrm{He}_{2k}(\xi).}
 \label{eq:f2v66-symbol}
\end{equation}
\end{theorem}

\begin{proof}
\(G=M^m(1)=(t+(\partial-\xi)^2)^m\mathrm{He}_0\).  Expanding the
binomial and using \((\partial-\xi)^{2k}\mathrm{He}_0=\mathrm{He}_{2k}\)
from Lemma~\ref{lem:f2v66-raise} gives \eqref{eq:f2v66-symbol}.  Exactly
verified for \(1\le m\le6\) and several \(t\) in the cited script
(SHA--256 \texttt{C88989BFC8A8EFF9AB8C6525A7AE5100}\allowbreak
\texttt{6211A44A7034567C154020D470F6D065}).
\end{proof}

\begin{remark}[Generating function]
\label{rem:f2v66-genfun}
Since \(\sum_{k}\mathrm{He}_{2k}(\xi)z^{2k}/(2k)!\) is the even part of
\(e^{\xi z-z^2/2}\), the symbol \eqref{eq:f2v66-symbol} is the diagonal
extraction of \((t+\text{even Hermite generating variable})^m\); in the
scaling \(t\to\infty\) it is \(t^m(1+O(1/t))\), the constant symbol whose
Appell sections are pure Hermite \(\mathrm{He}_r\), all real-rooted.  The
finite-\(t\) corrections are the successive \(\mathrm{He}_{2k}\) terms.
\end{remark}

\subsection{Consequences for the residual corner}

\begin{corollary}[Fully symmetric point]
\label{cor:f2v66-symmetric}
At the fully symmetric configuration \(c_1=\dots=c_m=0\),
\(t_1=\dots=t_m=t\), the Appell polynomials \(P_r\) of V64 are the
sections of \(e^{-s^2/2}G(is)\) with \(G\) given by
\eqref{eq:f2v66-symbol}.  This is a one-parameter family in \(t\); the
schema's real-axis condition at this point reduces to a single univariate
question, and by the V64 stacked-block computation the least \(t\) at
which it holds is strictly below \(n=2m+p\) for \(m\ge2\) (factor at most
\(0.85\)).  Hence the fully symmetric point is interior to the admissible
region, not on its boundary.
\end{corollary}

\begin{proof}
Specialization of Theorem~\ref{thm:f2v66-symbol} and
Corollary~\ref{cor:f2v65-section}; the numerical thresholds are those of
V64, Remark~\ref{rem:f2v62-numerics}.
\end{proof}

\begin{remark}[Why H61 is not at the symmetric point]
\label{rem:f2v66-h61}
The six residual groups of V61 are the coefficients of the corner
\(\sigma=\rho_2=\rho_3=0\), i.e.\ equal heights \(t_j=t\) but with the
centers \(c,a_2,a_3\) still free.  Theorem~\ref{thm:f2v66-symbol} closes
the symbol only when the centers also vanish; with free centers the
symbol acquires odd Hermite terms and center-dependent coefficients, and
the closed form is
\(G_{\bm a}(\xi)=\bigl[\prod_j(t+(\partial-\xi)^2-2a_j(\partial-\xi)\cdot
\text{(shift)})\bigr]1\), no longer a pure even Hermite sum.  So
Theorem~\ref{thm:f2v66-symbol} settles the height degeneracy but not the
center directions, which is precisely why the residual groups remained:
they are polynomials in \((c,a_2,a_3)\), not constants.  The advance is
that the height variable is now understood exactly, reducing (H61) to a
statement in the three center variables at fixed equal height.
\end{remark}

\subsection*{Firewall and V67 direction}
\label{fire:f2v66}
Lemma~\ref{lem:f2v66-raise} and Theorem~\ref{thm:f2v66-symbol} are proved
and verified.  They give the exact height dependence of the symbol but do
not resolve the center dependence, so (H61) and the three-pair band for
\(p\ge7\) remain conditional, and GRH remains open.

V67 should compute the center-dependent symbol
\(G_{\bm a}\) explicitly for \(m=3\) using the shifted raising operator
\(t+(\partial-\xi)^2\) conjugated by the translation \(e^{a_j\partial}\),
and test whether, at equal height \(t=n\), the resulting Appell sections
\(P_p,P_{p+1}\) can share a real zero as the centers vary over a bounded
range (the V33 center box already bounds them).  If the center variation
is a bounded analytic perturbation of the symmetric point of
Corollary~\ref{cor:f2v66-symmetric}, where the margin is positive, an
explicit perturbation bound would discharge (H61) and complete the
three-pair band for \(p\ge7\).

\subsection*{Assumption ledger}
\label{ass:f2v66}
The proved statements use the Hermite raising operator, the binomial
expansion of \(M^m\), and the exact verifier.  The thresholds and the
positivity of the symmetric-point margin are numerical (from V64).
\endgroup

\section*{Version V66 append-only record}
\addcontentsline{toc}{section}{Version V66 append-only record}

V66 was appended byte-for-byte before the terminal marker of validated
V65, whose installed SHA--256 was
\[
 \texttt{2729CEC7EB477B04C27B5F4A1D54C544}
 \allowbreak\texttt{BBE0A3284829FA846707D13CB1446840}
\]
with \(18335\) lines and \(627105\) bytes.  It proves the raising-operator
structure of the symbol recurrence and the equal-height closed form
\(G=\sum_k\binom mk t^{m-k}\mathrm{He}_{2k}\), locating the residual
corner in the center directions.  After validation V66 replaces V65 as
the sole active numeric GRH note; the frozen \texttt{V100\_f2} archive is
unchanged.

% =============================================================================
% END APPENDED V66 MATERIAL
% =============================================================================




\clearpage
\section{V67: the center-dependent symbol as a product of shifted parabolic operators}
\label{sec:f2v67-center-symbol}
\begingroup
\setcounter{equation}{0}
\renewcommand{\theequation}{N67.\arabic{equation}}
\renewcommand{\theHequation}{f2v67.\arabic{equation}}

V66 fixed the equal-height symbol at zero centers.  This iteration
removes the centering restriction, giving the full symbol as a commuting
product of shifted parabolic operators applied to the Gaussian.  This is
the exact object controlling the residual groups (H61), now with their
center dependence made explicit.

\subsection{The operator form}

\begin{theorem}[Center-dependent symbol]
\label{thm:f2v67-symbol}
For \(\Pi_m(w)=\prod_{j=1}^m((w-c_j)^2+t_j)\), the Fourier symbol of the
signed weight \(\dd\nu(W)=\Pi_m(iW)e^{-W^2/2}\dd W/\sqrt{2\pi}\) is
\begin{equation}
 \boxed{\,\widehat\nu(\xi)
 =\Bigl[\prod_{j=1}^m\bigl(t_j+(\partial_\xi-c_j)^2\bigr)\Bigr]e^{-\xi^2/2}
 =e^{-\xi^2/2}\,G_{\bm a}(\xi),}
 \label{eq:f2v67-symbol}
\end{equation}
where \(G_{\bm a}\) is a real polynomial of degree \(2m\).  The factors
\(t_j+(\partial_\xi-c_j)^2\) commute, so \(G_{\bm a}\) is symmetric in the
pairs.  In the \(P\)-representation \(\widehat\nu=P(\xi)e^{-\xi^2/2}\) the
factor acts by
\begin{equation}
 P\longmapsto t_jP+A_{c_j}A_{c_j}P,\qquad
 A_c P:=P'-(\xi+c)P,
 \label{eq:f2v67-recur}
\end{equation}
and at \(c_j=0\) this recovers Theorem~\ref{thm:f2v66-symbol}.
\end{theorem}

\begin{proof}
Multiplying the integrand by the \(j\)-th factor
\((iW-c_j)^2+t_j=(c_j^2+t_j-W^2)-2ic_jW\) corresponds under the Fourier
transform to \((c_j^2+t_j+\partial_\xi^2)-2c_j\partial_\xi
=t_j+(\partial_\xi-c_j)^2\), using \(\widehat{W^2f}=-\partial^2\widehat f\)
and \(\widehat{Wf}=-i\partial\widehat f\).  Composing the \(m\) factors,
all polynomials in \(\partial_\xi\) and hence mutually commuting, gives
\eqref{eq:f2v67-symbol}; \eqref{eq:f2v67-recur} is the action on
\(Pe^{-\xi^2/2}\) since \(\partial_\xi[Pe^{-\xi^2/2}]=(P'-\xi P)e^{-\xi^2/2}\),
so \((\partial_\xi-c)[Pe^{-\xi^2/2}]=(P'-(\xi+c)P)e^{-\xi^2/2}\).  Verified
exactly over \(\mathbb Q\) by matching Taylor coefficients against the
exact moments of \(\nu\) in
\path{scripts/verify_grh_v67_center_symbol.py} (SHA--256
\texttt{5CC837887C22616B9451119704FA298D}\allowbreak
\texttt{963070BA5DDF120DA27BADE99E996B8F}), which also checks the
symmetry.
\end{proof}

\begin{corollary}[Appell generating function]
\label{cor:f2v67-genfun}
The Appell polynomials of V64 satisfy
\begin{equation}
 \sum_{r\ge0}P_r(x)\frac{s^r}{r!}=e^{sx-s^2/2}\,G_{\bm a}(is),
 \label{eq:f2v67-genfun}
\end{equation}
so \(P_r=\sum_{k}\binom rk (i)^kg_k\,\mathrm{He}_{r-k}\) where
\(G_{\bm a}(is)=\sum_kg_k(is)^k\); the schema's real-axis condition
(Proposition~\ref{prop:f2v63-criterion}) is that \(P_p,P_{p+1}\) share no
real zero.
\end{corollary}

\begin{proof}
Combine Theorem~\ref{thm:f2v64-rep} generating function
\(e^{sx}\widehat\nu(is)\) with \eqref{eq:f2v67-symbol}, and expand
\(G_{\bm a}(is)e^{-s^2/2}\) against \(e^{sx}\).
\end{proof}

\subsection{Reduction of the residual corner}

\begin{proposition}[H61 as a bounded center perturbation]
\label{prop:f2v67-h61}
On the residual corner of V61, the three heights equal \(t=n\) and the
centers \((c,a_2,a_3)\) lie in the compact box of V33.  The symbol is
\begin{equation}
 G_{\bm a}=\Bigl[\bigl(t+(\partial-c)^2\bigr)\bigl(t+(\partial-a_2)^2\bigr)
 \bigl(t+(\partial-a_3)^2\bigr)\Bigr]1 ,
 \label{eq:f2v67-corner}
\end{equation}
a polynomial deformation, analytic and of bounded degree in
\((c,a_2,a_3)\), of the fully symmetric symbol
\(\sum_k\binom3k t^{3-k}\mathrm{He}_{2k}\) of
Corollary~\ref{cor:f2v66-symmetric}, at which the schema's condition
holds with a strictly positive separation margin.  Hypothesis (H61) is
therefore the statement that this margin survives the center deformation
over the compact box.
\end{proposition}

\begin{proof}
Equation \eqref{eq:f2v67-corner} is
Theorem~\ref{thm:f2v67-symbol} at \(t_j=t\).  Setting
\(c=a_2=a_3=0\) gives the symmetric symbol of
Theorem~\ref{thm:f2v66-symbol}, where by
Corollary~\ref{cor:f2v66-symmetric} and the V64 separation data the
common-zero margin is positive.  The coefficients of \(G_{\bm a}\) are
polynomials in the centers, so \(P_p,P_{p+1}\) and their resultant are
continuous, indeed polynomial, in \((c,a_2,a_3)\).
\end{proof}

\begin{remark}[What remains for a discharge]
\label{rem:f2v67-remains}
A quantitative discharge of (H61) needs a lower bound on the resultant
\(\operatorname{Res}(P_p,P_{p+1})\) at the symmetric point together with a
Lipschitz bound in the centers over the box, both explicit.  The
numerics of V64 give the margin \(\ge0.12\) uniformly; the missing
analytic input is the Lipschitz constant, which
\eqref{eq:f2v67-corner} makes computable as the center-gradient of a
fixed polynomial map.  This is a finite computation of the same kind as
the V61 groups but in three variables at fixed height, hence far smaller;
it is the natural V68 target.
\end{remark}

\subsection*{Firewall and V68 direction}
\label{fire:f2v67}
Theorem~\ref{thm:f2v67-symbol} and Corollary~\ref{cor:f2v67-genfun} are
proved and verified.  They express the symbol exactly but do not bound
the center deformation, so (H61) and the three-pair band for \(p\ge7\)
stay conditional, and GRH remains open.

V68 should attempt the local route of
Remark~\ref{rem:f2v67-remains}: bound
\(\operatorname{Res}(P_p,P_{p+1})\) below at the symmetric center
\(c=a_2=a_3=0\), where \eqref{eq:f2v67-symbol} gives the explicit
Hermite-sum symbol and the margin is positive, and control its center
gradient over the V33 box.  This is not a reduction in the number of
variables: the polynomial variable \(x\) remains and \(p\) is swept
through the coefficient structure, so a global exact certificate is no
smaller than the V61 determinant.  The gain is that a Lipschitz
perturbation bound around one explicit point is a genuinely different and
possibly cheaper analytic device than a global sum-of-squares search;
whether it closes (H61) is open.

\subsection*{Assumption ledger}
\label{ass:f2v67}
The proved statements use the Fourier--differentiation rules, the
commutation of polynomials in \(\partial_\xi\), the Gaussian transform,
and the exact verifier.  The positivity of the symmetric-point margin and
the numerical separation bound are from V64.
\endgroup

\section*{Version V67 append-only record}
\addcontentsline{toc}{section}{Version V67 append-only record}

V67 was appended byte-for-byte before the terminal marker of validated
V66, whose installed SHA--256 was
\[
 \texttt{B666DCB5E9CCE73B76B7DEAE3BE1DE39}
 \allowbreak\texttt{EF0A82C9BE8B67CB21E628DDC34D2F84}
\]
with \(18482\) lines and \(633590\) bytes.  It proves the
center-dependent Fourier symbol as a commuting product of shifted
parabolic operators and reduces hypothesis (H61) to a bounded center
perturbation of the positive-margin symmetric point.  After validation
V67 replaces V66 as the sole active numeric GRH note; the frozen
\texttt{V100\_f2} archive is unchanged.

% =============================================================================
% END APPENDED V67 MATERIAL
% =============================================================================




\clearpage
\section{V68: the leading form and collision hierarchy of the real-axis determinant, the one-pair closed form, and the bridge problem}
\label{sec:f2v68-structure}
\begingroup
\setcounter{equation}{0}
\renewcommand{\theequation}{N68.\arabic{equation}}
\renewcommand{\theHequation}{f2v68.\arabic{equation}}

This iteration follows a programme review: the six residual groups of
(H61) are a finite obstruction, and certifying them one by one teaches
nothing about the \((m,p)\) family or about the passage from an actual
\(L\)-function zero to a finite block.  Computation is therefore used here
only diagnostically, to expose structure.  Three things are established.
First, the leading form of the three-pair determinant is an explicit
square, and the hyperplanes on which it vanishes are the collision loci
of the \(p\)-fold root with a pair center and of a pair center with the
origin; this is the geometric reason the earlier Gram searches stalled.
Second, the one-pair determinant has an exact closed form in which the
sharp constant \(2m+p\) is visible as a root of the constant term, and
the two-pair determinant inherits the same divisibility and leading
structure.  Third, the bridge from an off-line zero to a finite block,
which the notes have so far named only in firewalls, is stated as a
problem with explicit quantifiers.

Throughout, \(W_0=A_uB_v-A_vB_u\) is the real-axis determinant of the
reals-only diagonal in the balanced basis
\((\mathrm{He}_{p-2},\mathrm{He}_{p-3})\) (V58--V61), scaled by
\(M_m=(p-3)\cdots(p-2m)\) for \(m\ge2\) and \(M_1=1\); for \(m=3\) the
variables are \((x,c,a_2,a_3)\) and \(p\), with \((c,t)\) the first pair
and \((a_j,b_j+1)\) the companion pairs, and the corner is
\(t=p+6\), \(b_j=t-1\), i.e.\ \(\sigma=\rho_2=\rho_3=0\) in the V61
substitution.  All exact claims below are checked over \(\mathbb Q\) by
\path{scripts/verify_grh_v68_residual_top_forms.py} (SHA--256
\texttt{F312CD636032D5715EF039B3390574D5}\allowbreak\texttt{07E4A990C96A7083F1AACFED7F60BB85}), which rebuilds the
\(578\) groups from the V61 construction.

\subsection{The leading form is a square}

\begin{theorem}[Leading form of the three-pair determinant]
\label{thm:f2v68-leading}
Write \(G_0(w):=\prod_{j}(w-c_j)^2\) for the height-free companion of the
three pairs, \(c_1=c\), \(c_2=a_2\), \(c_3=a_3\), and
\begin{equation}
 L:=c\,a_2\,a_3\,(x-c)(x-a_2)(x-a_3),\qquad L^2=G_0(0)\,G_0(x).
 \label{eq:f2v68-L}
\end{equation}
The part of \(M_3W_0\) of total degree \(12\) in \((x,c,a_2,a_3)\) is
\begin{equation}
 \boxed{\,\bigl[M_3W_0\bigr]_{12}=(p)_7\,M_3\,L^2,\qquad
 (p)_7=p(p-1)\cdots(p-6),}
 \label{eq:f2v68-leading}
\end{equation}
and it is independent of the heights: after the V61 substitution only the
pure-\(\pi\) groups \((0,0,0,k)\), \(0\le k\le11\), carry a degree-\(12\)
part, each a positive rational multiple \(\lambda_k L^2\), and
\(\sum_k\lambda_k\pi^k=(p)_7M_3\) at \(p=7+\pi\).
\end{theorem}

\begin{proof}
Exact verification: the verifier rebuilds every group, checks that the
degree-\(12\) monomial support of each group equals that of \(L^2\) with a
single proportionality constant, that no group with
\((\sigma,\rho_2,\rho_3)\ne0\) has a degree-\(12\) part, and that the
univariate polynomial \(\sum_k\lambda_k\pi^k\) equals
\(\prod_{i=0}^{6}(\pi+7-i)\prod_{i=3}^{6}(\pi+7-i)\).  The values are
\(\lambda_0=120960\), \(\lambda_{11}=1\).
\end{proof}

\begin{remark}[Meaning]
\label{rem:f2v68-meaning}
\(G_0(0)\) is the companion evaluated at the evaluation point of the
diagonal functional and \(G_0(x)\) is the companion evaluated at the
\(p\)-fold root; at large scale the determinant is their product times
the combinatorial constant \((p)_7M_3\).  The leading form vanishes
exactly on the six hyperplanes \(x=c_j\) (the \(p\)-fold root collides
with a pair center) and \(c_j=0\) (a pair center sits at the origin).
These are not coordinate hyperplanes of \((x,c,a_2,a_3)\), and no linear
change of variables makes all six coordinate hyperplanes.  This is why
the Newton-polytope Gram searches of V61 stalled with a numerical kernel
they could not rationalize: the zeros at infinity are rational, but they
are invisible to a monomial basis.  After \(u_j:=x-c_j\), three of the
six become coordinate hyperplanes, and the residual group \((1,0,0,2)\)
then admits an exact Gram certificate (kind \texttt{exact\_u} in
\path{scripts/data_grh_v68_residual_certificates.json}, basis \(53\),
zero remainder), verified; the four degree-\(12\) groups still show a
\(17\)-dimensional kernel at the next level of the hierarchy below, and
were deliberately not pursued further.
\end{remark}

\begin{proposition}[Collision hierarchy at the corner]
\label{prop:f2v68-hierarchy}
Let \(R_k\) be the pure-\(\pi\) group \((0,0,0,k)\) and
\(R_k':=R_k-\lambda_kL^2\), of total degree \(\le10\).  For each of the
six linear forms \(\ell\in\{x-c_j,\;c_j\}\) let \(L_\ell:=L/\ell\).  Then
the degree-\(10\) form of \(R_k'\) restricted to \(\{\ell=0\}\) is
\(\kappa^{\ell}_kL_\ell^2|_{\ell=0}\), with
\begin{equation}
 \sum_k\kappa^{x-c_j}_k\pi^k=(p)_7M_3\,(\pi+13)=(p)_7M_3\;t,\qquad
 \sum_k\kappa^{c_j}_k\pi^k=(p)_7M_3\,(\pi+12)=(p)_7M_3\;(t-1),
 \label{eq:f2v68-kappa}
\end{equation}
where \(t=p+6\) is the corner height, and the degree-\(10\) form of
\(R_k'-\sum_\ell\kappa_k^\ell L_\ell^2\) is divisible by \(L\).  In
particular the top two layers of \(M_3W_0\) at the corner agree with the
expansion of
\begin{equation}
 (p)_7M_3\;\widehat G(0)\,G(x),\qquad
 \widehat G(0)=\prod_j(c_j^2+t-1),\quad G(x)=\prod_j\bigl((x-c_j)^2+t\bigr),
 \label{eq:f2v68-main}
\end{equation}
in which the origin sees the heat-lowered heights \(t-1\)
(Lemma~\ref{lem:f2v63-const}) and the \(p\)-fold root sees the true
heights \(t\).
\end{proposition}

\begin{proof}
Exact verification, group by group with \(p\) symbolic: restriction,
comparison with \(L_\ell^2\), and six successive exact divisions by the
linear forms.  The symmetry \(\kappa^{x-c}=\kappa^{x-a_2}=\kappa^{x-a_3}\),
\(\kappa^{c}=\kappa^{a_2}=\kappa^{a_3}\) is checked as well.  Expanding
\eqref{eq:f2v68-main} to second order in the heights gives
\(L^2+t\sum_j L_{x-c_j}^2+(t-1)\sum_jL_{c_j}^2+\dots\), matching
\eqref{eq:f2v68-kappa}.
\end{proof}

The remaining degree-\(10\) term \(L\,q_4\) is a cross term.  Its
coefficient of the highest power \(\pi^{12}\) is
\(q_4=-2\sum_{i<j}L_iL_j\) with \(L_j:=c_j(x-c_j)\), so that
\(Lq_4=-2\sum_kL_k\prod_{i\ne k}L_i^2\): the one-pair cross term of
Theorem~\ref{thm:f2v68-onepair} below, applied to each pair and
multiplied by the leading squares of the other two.  This is the shape
of the structural decomposition
\(R=R_{\rm boundary}+\sum_ju_j^2A_j+\dots\) that a proof of (H61)
should exhibit; it is recorded, not proved.  The proportionality
\(q_4\propto\sum_{i<j}L_iL_j\) holds only at the top power of \(\pi\):
for the lower groups the cofactor has \(22\) monomials, including
\(x\,c\,a_2a_3^2\)-type terms that involve all three centers, so the
pair coupling enters exactly at the first cross term.  Consequently the
naive product ansatz \(M_3W_0\ge(p)_7M_3\prod_j\Psi_1(c_j,t;p_e)\)
with a common effective \(p_e\in\{p,\dots,p+4\}\) is false: the
difference is unbounded below along rays, because the two sides agree in
degrees \(12\) and in the \(\kappa\)-layer of degree \(10\) but differ
in the odd cross term, which has no sign.  A structural lower bound must
therefore carry a genuinely three-pair cross term.

\subsection{The one-pair determinant in closed form}

\begin{theorem}[One-pair closed form]
\label{thm:f2v68-onepair}
For one pair \((c,t)\), a \(p\)-fold root at \(x\), and the basis
\((\mathrm{He}_{p-2},\mathrm{He}_{p-3})\) with \(M_1=1\),
\begin{equation}
 \boxed{\,W_0^{(1)}=p(p-1)(p-2)\,\Psi_1,\qquad
 \Psi_1=(c^2+t-1)\bigl((x-c)^2+t\bigr)-2(p+1)\bigl(c(x-c)+t\bigr)+(p+1)(p+2),}
 \label{eq:f2v68-onepair}
\end{equation}
equivalently
\begin{equation}
 \Psi_1=c^2(x-c)^2+(t-1)x^2-2(t+p)xc+(2t+2p+1)c^2+(t-p-1)(t-p-2).
 \label{eq:f2v68-onepair2}
\end{equation}
\end{theorem}

\begin{proof}
\(W_0^{(1)}\) is built by the same Leibniz form and generating identity as
the two- and three-pair determinants (V53, V58, V61), with no companion
factor; the identity \eqref{eq:f2v68-onepair} is then a finite polynomial
identity in \((x,c,t,p)\), verified exactly.  The two forms agree by
expansion.
\end{proof}

\begin{corollary}[Structural proof of the one-pair threshold]
\label{cor:f2v68-onepair-sharp}
Let \(p\ge3\) and \(t\ge p+2\).  Then \(W_0^{(1)}>0\) for all real
\(x,c\), except at \((x,c,t)=(0,0,p+2)\) where it vanishes.  Hence
condition \textup{(i)} of the exclusion schema holds for one pair with the
sharp constant \(t_0=p+2=2m+p\).
\end{corollary}

\begin{proof}
In \eqref{eq:f2v68-onepair2} the first term is a square, the last is
\(\ge0\) with equality only at \(t=p+2\), and the quadratic form
\((t-1)x^2-2(t+p)xc+(2t+2p+1)c^2\) has
\((t+p)^2-(t-1)(2t+2p+1)=(p+1)^2-t(t-1)<0\) for \(t\ge p+2\), since
\(t(t-1)\ge(p+2)(p+1)>(p+1)^2\); it is positive definite and vanishes
only at \(x=c=0\).  The factor \(p(p-1)(p-2)\) is positive for \(p\ge3\).
\end{proof}

This recovers Theorem~\ref{thm:f2v63-onepair} with a reason: the sharp
constant is the root of the constant term \((t-p-1)(t-p-2)\), the
\(x\)- and \(c\)-directions are controlled by a quadratic form that is
definite precisely above \(t(t-1)=(p+1)^2\), and the leading square
\(c^2(x-c)^2=L_1^2\) is the \(m=1\) instance of
Theorem~\ref{thm:f2v68-leading}.  The middle layer \(tc^2+(t-1)(x-c)^2\)
of \eqref{eq:f2v68-onepair} is the \(m=1\) instance of
\eqref{eq:f2v68-kappa}, and the cross term \(-2(p+1)L_1\) is the one
whose three-pair shadow was identified above.

\begin{proposition}[Two-pair divisibility and leading layers]
\label{prop:f2v68-twopair}
For two pairs \((c,t)\), \((a,b+1)\) and \(M_2=(p-3)(p-4)\),
\(W_0^{(2)}\) is divisible by \((p)_5M_2\):
\(W_0^{(2)}=(p)_5M_2\,\Psi_2\), where \(\Psi_2\) is a monic quartic in
\(p\) whose \(p^3\)-coefficient is \(10-2(R_1+R_2)\),
\(R_j:=c_j(x-c_j)+t_j\), and whose leading form in \((x,c,a)\) is
\(L_1^2L_2^2=[c\,a\,(x-c)(x-a)]^2\).  At the symmetric point
\(x=c=a=0\), \(t_1=t_2=t\),
\begin{equation}
 \Psi_2=t^4-(4p+6)t^3+(6p^2+22p+19)t^2-(4p^3+26p^2+50p+28)t+(p+1)(p+2)(p+3)(p+4).
 \label{eq:f2v68-psi2-sym}
\end{equation}
\end{proposition}

\begin{proof}
Exact verification: \(W_0^{(2)}\) is rebuilt by the V58 construction and
divided by \((p)_5M_2\) with zero remainder; the \(p^4\), \(p^3\) parts,
the leading form and \eqref{eq:f2v68-psi2-sym} are read off.
\end{proof}

The pattern \(W_0^{(m)}=(p)_{2m+1}M_m\Psi_m\) with \(\Psi_m\) monic of
degree \(2m\) in \(p\) and leading form \(\prod_jL_j^2\) is an
observation for \(m\le2\) (for \(m=3\) only the leading form
\eqref{eq:f2v68-leading} is proved); what is proved is the \(m=1\)
closed form and the \(m=2\) divisibility.  The cross-pair terms of \(\Psi_2\) are not
expressible through the two one-pair brackets alone (the remainder after
subtracting \(\Psi_1(c,t;p)\Psi_1(a,b+1;p)\) has \(80\) terms), so the
determinant does not factor over pairs; the coupling is genuine.

\subsection{The Lipschitz route is closed}

\begin{proposition}[No perturbative discharge of (H61)]
\label{prop:f2v68-lipschitz-nogo}
The programme of Remark~\ref{rem:f2v67-remains}, a resultant margin at the
symmetric point plus a Lipschitz bound in the centers over the V33 box,
cannot close (H61).
\end{proposition}

\begin{proof}
Two independent obstructions.  (a) The box of
Theorem~\ref{thm:f2v33-global-small-box} allows \(|c_j|<871/10\) while
the heights obey \(t_j<19321/100\); along the main term
\eqref{eq:f2v68-main} the factor \(\widehat G(0)\) alone varies by the
ratio \(\prod_j(c_j^2+t-1)/(t-1)^3\), which exceeds \(10^{4}\) already for
\(|c_j|=20\) at \(t=13\).  The center dependence is not a perturbation of
the symmetric point on any scale on which a margin of order one could
survive.  (b) \(\operatorname{Res}(P_p,P_{p+1})\) has degree growing with
\(p\), so no single Lipschitz constant is uniform in \(p\ge7\); the
uniformity in \(p\) that V61 achieves comes only from the fixed-degree
balanced-basis determinant.
\end{proof}

The replacement target is structural: a lower bound of \(M_mW_0\) by an
explicit positive main term such as \eqref{eq:f2v68-main}, with the cross
terms absorbed by the hierarchy of squares; Theorem~\ref{thm:f2v68-onepair}
is its \(m=1\) instance.

\subsection{The bridge problem stated}

The notes have so far named the ``infinite-product passage'' and the
``arithmetic forcing step'' only inside firewalls; neither has been
stated.  On the analytic side the monograph
(\path{manuscripts/flagship/flagship_theorems.tex}, Theorem
``Gaussian heat criterion for Dirichlet GRH'') proves, for a primitive
character \(\chi\) with completed function \(\Xi_\chi\), that all zeros
lie on \(\operatorname{Re}s=\tfrac12\) if and only if the reflected zero
heat \(u_\chi\ge0\), if and only if the reflected \(2\times2\) affine
moment matrix generated by the ports \(1\) and \(s-\tfrac12\) is positive
semidefinite at every Gaussian scale and center.  On the finite side,
Corollary~\ref{cor:f2v9-gamma}, the target \eqref{eq:f2v52-target} and the
degree-free schema Theorem~\ref{thm:f2v60-schema} show: a critical gamma
block with \(m\) pair channels, \(p\) real factors and ladder count
\(R\ge2m+p\) whose backward-heat image is real-rooted satisfies
\(y_{\min}^2\le nK\), \(n=2m+p\), which keeps \(\kappa^2\) above the band
top.

\begin{problem}[Bridge from an off-line zero to a finite block]
\label{prob:f2v68-bridge}
Let \(\rho_0=\beta_0+i\gamma_0\) be a zero of \(\Xi_\chi\) with
\(\beta_0>\tfrac12\).  Produce, with all constants explicit:
\begin{enumerate}[label=\textup{(B\arabic*)}]
\item \emph{Extraction.}  Integers \(m,p\ge0\), centers \(c_j\) and heights
\(t_j\), and a scale \(\theta>0\), such that the finite block
\(\prod_{j\le m}((w-c_j)^2+t_j)\,(w-x)^p\) is a genuine gamma block of
\(\chi\) at the ordinate \(\gamma_0\) in the sense of
Corollary~\ref{cor:f2v9-gamma}, its backward-heat image at the critical
time \(K=(1+\theta)/\theta^2\) is real-rooted, and its least height
satisfies \(y_{\min}^2>nK\) with \(n=2m+p\).
\item \emph{Uniformity.}  Either \((m,p)\) is bounded in terms of \(\chi\)
alone, or the exclusion \(y_{\min}^2\le nK\) is available for the
\((m,p)\) that \textup{(B1)} produces; at present it is proved for
\(m=1\) (all \(p\)), \(m=2\) (all \(p\)), and \(m=3\) with \(p\le3\) or
under (H61).
\item \emph{Arithmetic forcing.}  The offsets of the block come from the
gamma ladder of \(\chi\), and this constraint excludes the Hermite
extremizer \(e^{(K/2)D^2}z^R\) of the shallow-pair bound, so that the
block cannot sit on the boundary \(y_{\min}^2=nK\).
\item \emph{Tightness.}  Real padding and remote channels cannot escape to
infinity in the extraction: the block of \textup{(B1)} stays in the
compact region of Theorem~\ref{thm:f2v33-global-small-box} or in an
explicit enlargement of it.
\end{enumerate}
A proof of \textup{(B1)}--\textup{(B4)} together with the exclusion
theorems would give a contradiction, hence GRH for \(\chi\).
\end{problem}

None of (B1)--(B4) is proved here.  What the statement makes visible is
that (B2) is the only item on which the finite-block work of V51--V67
bears, and that (B1) and (B3) are of a different nature: they concern
the completed \(L\)-function, not a polynomial.  The natural first case
is the smallest one: show that the failure of the affine moment matrix
at a single Gaussian scale and center produces a two-pair block, the
case in which the exclusion is unconditional.

\subsection*{Firewall and V69 direction}
\label{fire:f2v68}
Theorem~\ref{thm:f2v68-leading}, Proposition~\ref{prop:f2v68-hierarchy},
Theorem~\ref{thm:f2v68-onepair}, Corollary~\ref{cor:f2v68-onepair-sharp}
and Proposition~\ref{prop:f2v68-twopair} are proved and verified;
Proposition~\ref{prop:f2v68-lipschitz-nogo} closes a route.  Five of the
six residual groups of (H61) remain open, so the three-pair band for
\(p\ge7\) is still conditional; Problem~\ref{prob:f2v68-bridge} is
unproved.  GRH remains open.

V69 should split effort as the review prescribes.  On the finite side,
seek the structural decomposition of \(\Psi_2\) and of the corner
\(\Psi_3\) suggested by \eqref{eq:f2v68-main} and the cross term
\(-2\sum_kL_k\prod_{i\ne k}L_i^2\): a lower bound
\(\Psi_m\ge\widehat G(0)G(x)-(\text{explicit cross terms})\) that is
manifestly positive above \(t=2m+p\), proved for \(m=2\) first, where the
answer is already known to be true.  On the bridge, attempt (B1) in its
smallest form, extracting a two-pair block from one failing Gaussian
scale of the affine moment matrix.  Certificates for the five groups are
a fallback only.

\subsection*{Assumption ledger}
\label{ass:f2v68}
The proved statements use the V53/V58/V61 constructions rebuilt from
source, exact rational arithmetic, and the verifier.  The interpretation
of the hierarchy through \eqref{eq:f2v68-main} is exact for the two top
layers and heuristic below them.  The count of \(80\) remainder terms and
the \(17\)-dimensional kernels are computational observations.  The
monograph theorem is cited, not re-proved.
\endgroup

\section*{Version V68 append-only record}
\addcontentsline{toc}{section}{Version V68 append-only record}

V68 was appended byte-for-byte before the terminal marker of validated
V67, whose installed SHA--256 was
\[
 \texttt{4A050D2FE5D6B29E11A878BD2EBC071A}
 \allowbreak\texttt{1FB1227F2636E98D7844D86C5F80AA58}
\]
with \(18652\) lines and \(640849\) bytes.  It proves the leading-form
and collision-hierarchy structure of the three-pair determinant, the
one-pair closed form with its structural threshold proof, the two-pair
divisibility, closes the Lipschitz route, certifies one residual group,
and states the bridge problem.  After validation V68 replaces V67 as the
sole active numeric GRH note; the frozen \texttt{V100\_f2} archive is
unchanged.

% =============================================================================
% END APPENDED V68 MATERIAL
% =============================================================================




\clearpage
\section{V69: the pairing form of the heat image, the exact one-pair region, the failure of the modulus invariant, and the bridge ledger}
\label{sec:f2v69-pairing}
\begingroup
\setcounter{equation}{0}
\renewcommand{\theequation}{N69.\arabic{equation}}
\renewcommand{\theHequation}{f2v69.\arabic{equation}}

This iteration keeps the V68 split: structure on the finite side, and the
bridge on the arithmetic side, with computation used only to test general
statements.  On the finite side three exact facts are established.  The
heat image of any product of pair factors has a pairing expansion in which
pairs interact only through explicit two-point couplings; the one-pair
exclusion region is exactly the exterior of a circle, \(|c+i\sqrt t|^2>p+2\),
which identifies it with the Jacobi criterion of V54; and the natural
generalization of that circle to several pairs is false, with an exact
witness, so the region for \(m\ge2\) is a genuinely coupled one.  On the
bridge side the intended contradiction of the gamma-block architecture is
written out as a ledger of implications, each marked proved, open, or
heuristic, which locates the missing arrows more precisely than
Problem~\ref{prob:f2v68-bridge}.  All exact claims are checked by
\path{scripts/verify_grh_v69_pairing_and_one_pair_region.py} (SHA--256
\texttt{A3EE54E80A02F2F01551466E3709BE0F}\allowbreak\texttt{40C731B5378BCBCC60D432D80F5D70AB}), which rebuilds the two-pair
quotient \(\Psi_2\) by the V68 construction.

\subsection{The pairing form of the heat image}

\begin{theorem}[Pairing expansion]
\label{thm:f2v69-pairing}
Let \(f_j(w)=(w-c_j)^2+t_j\), \(1\le j\le m\), and put
\(\widehat f_j:=f_j-1\), \(f_j'=2(w-c_j)\).  Then
\begin{equation}
 \boxed{\;e^{-D^2/2}\prod_{j=1}^mf_j
 =\sum_{\substack{A,B\subset[m]\\ A\cap B=\emptyset,\ |B|\ \mathrm{even}}}
 \mu(|A|,|B|)\prod_{j\notin A\cup B}\widehat f_j\prod_{j\in B}f_j',\qquad
 \mu(a,b)=(-1)^{b/2}\,\mathbb E\bigl[Z^b(1-Z^2)^a\bigr],}
 \label{eq:f2v69-pairing}
\end{equation}
with \(Z\) standard normal.  The first values are \(\mu(0,0)=1\),
\(\mu(1,0)=0\), \(\mu(2,0)=2\), \(\mu(0,2)=-1\), \(\mu(1,2)=2\),
\(\mu(3,0)=-8\).  In particular
\begin{align}
 m=2:&\quad \widehat Q=\widehat f_1\widehat f_2-f_1'f_2'+2,
 \label{eq:f2v69-m2}\\
 m=3:&\quad \widehat Q=\widehat f_1\widehat f_2\widehat f_3
 -\sum_{i<j}f_i'f_j'\widehat f_k+2\sum_k\widehat f_k+2\sum_{i<j}f_i'f_j'-8 .
 \label{eq:f2v69-m3}
\end{align}
\end{theorem}

\begin{proof}
\((e^{-D^2/2}F)(w)=\mathbb E[F(w+iZ)]\), and for a monic quadratic
\(f_j(w+iZ)=f_j(w)+iZf_j'(w)-Z^2=\widehat f_j+(1-Z^2)+iZf_j'\) exactly.
Expand the product over the three choices per factor and take the
expectation: the factors \(1-Z^2\) from \(A\) and \(iZf_j'\) from \(B\)
contribute \(\mathbb E[(1-Z^2)^{|A|}(iZ)^{|B|}]\), which vanishes for odd
\(|B|\) and equals \(\mu(|A|,|B|)\) otherwise.  The values follow from
\(\mathbb E[Z^{2l}]=(2l-1)!!\).  Verified exactly for \(m\le4\) with
symbolic centers and heights.
\end{proof}

\begin{remark}[What the expansion says]
\label{rem:f2v69-pairing}
Since \(\mu(1,0)=0\), a single factor is only shifted, \(f_j\mapsto\widehat f_j\)
(the height drop of Lemma~\ref{lem:f2v63-const}); all genuine interaction
is through the couplings \(-f_i'f_j'=-4(w-c_i)(w-c_j)\) and the constants
\(\mu(a,0)\).  The Appell polynomial \(P_{p+1}=\widehat Q(D)\mathrm{He}_{p+1}\)
is therefore the uncoupled product \(\prod_j\widehat f_j(D)\mathrm{He}_{p+1}\)
plus pair-coupling corrections of two orders less in \(D\).  For one pair
the uncoupled term is everything, which is why the one-pair problem has a
closed form (Theorem~\ref{thm:f2v68-onepair}); the V68 cross term
\(-2L\sum_{i<j}L_iL_j\) is the shadow of \(-\sum f_i'f_j'\widehat f_k\).
\end{remark}

\subsection{The exact one-pair region}

\begin{theorem}[Exact one-pair region]
\label{thm:f2v69-circle}
Write \(\Psi_1=Ax^2-Bx+C\) from \eqref{eq:f2v68-onepair2}, so
\(A=c^2+t-1\), \(B=2c(c^2+t+p)\), and put \(\rho:=c^2+t-p-2\),
\(|z|^2:=c^2+t\) with \(z=c+i\sqrt t\).  Then
\begin{equation}
 4AC-B^2=4\rho\,\bigl(t\rho+p+1\bigr).
 \label{eq:f2v69-disc}
\end{equation}
Consequently, for \(p\ge3\) and \(|z|^2>1\), \(W_0^{(1)}(x)>0\) for every
real \(x\) if and only if \(\rho(t\rho+p+1)>0\), i.e.\ on exactly two
components:
\begin{equation}
 \boxed{\;|z|^2>p+2\;}\qquad\text{or}\qquad t\bigl(p+2-|z|^2\bigr)>p+1 .
 \label{eq:f2v69-circle}
\end{equation}
The zero set of \(W_0^{(1)}\) in the \((c,t)\)-plane is the circle
\(|z|^2=p+2\) together with the curve \(t(p+2-|z|^2)=p+1\); between them
\(F_G\) has a real zero, on them \(P_{p+1}\) has a real double zero.
\end{theorem}

\begin{proof}
\eqref{eq:f2v69-disc} is a polynomial identity in \((c,t,p)\), verified
exactly together with the agreement of \(Ax^2-Bx+C\) with the bracket of
\eqref{eq:f2v68-onepair}.  Since \(A=|z|^2-1>0\), \(\Psi_1>0\) for all
real \(x\) iff \(4AC-B^2>0\), and the two factors give the two
components; the factor \(p(p-1)(p-2)\) is positive.  The verifier checks
the three sample points \((c,t,p)=(0,\tfrac72,5)\) (inner component),
\((0,\tfrac{13}2,5)\) (between), \((1,7,5)\) (exterior).
\end{proof}

\begin{corollary}[Exclusion on the circle exterior]
\label{cor:f2v69-exterior}
For one pair and \(p\ge3\), \(F_G\) is zero-free in the closed upper
half-plane whenever \(|z|^2=c^2+t>p+2\); hence the one-pair band
statement of Theorem~\ref{thm:f2v63-onepair} extends from \(t>p+2\) to
the circle exterior, and the constant \(p+2=2m+p\) is the squared
distance from the evaluation point to the pair's zero.
\end{corollary}

\begin{proof}
The exterior \(\{|z|^2>p+2\}\) is connected, contains the large-height
points where all zeros of \(F_G\) lie in the lower half-plane
(Lemma~\ref{lem:f2v61-deform}), and carries condition (i) by
Theorem~\ref{thm:f2v69-circle}; apply the schema
Theorem~\ref{thm:f2v60-schema}.  The inner component is not connected to
those points and carries no exclusion.
\end{proof}

\begin{corollary}[Identification with the Jacobi criterion]
\label{cor:f2v69-jacobi}
For one pair the schema condition ``no real double zero of \(P_{p+1}\)''
and the real-rootedness criterion of Theorem~\ref{thm:f2v54-jacobi}
(applied to \(P_{p+1}=((D-c)^2+t-1)\mathrm{He}_{p+1}\), i.e.\
\(c^2+(t-1)\ge p+1\)) have the same outer boundary \(|z|^2=p+2\).  On the
inner component of \eqref{eq:f2v69-circle} the polynomial \(P_{p+1}\) has
no real double zero but is not real-rooted, so the two conditions differ
there; only the outer boundary carries the sharp constant.
\end{corollary}

\begin{proof}
Compare \eqref{eq:f2v69-circle} with \(a^2+b^2\ge n\) at \(a=c\),
\(b^2=t-1\), \(n=p+1\).
\end{proof}

\subsection{The modulus invariant fails for two and three pairs}

\begin{proposition}[Exact counterexample]
\label{prop:f2v69-witness}
The statement ``\(|z_j|^2\ge2m+p\) for all \(j\) implies \(W_0>0\) for
all real \(x\)'' is false for \(m=2\): at \(p=10\), \(c_1=c_2=\tfrac72\),
\(t_1=t_2=2\) (so \(|z_j|^2=\tfrac{57}4\ge14=2m+p\)) and \(x=8\),
\begin{equation}
 \Psi_2=-\frac{21527}{256}<0,\qquad W_0^{(2)}=(p)_5(p-3)(p-4)\Psi_2<0 .
 \label{eq:f2v69-witness}
\end{equation}
Numerically the same failure occurs for \(m=3\) (four of sixty random
points with all \(|z_j|^2\ge2m+p\) and all heights below \(2m+p\) give
\(W_0<0\)), always with the centers on one side of the origin and at
least one height small.
\end{proposition}

\begin{proof}
Exact evaluation of the rebuilt \(\Psi_2\) in the verifier.
\end{proof}

So the one-pair circle does not propagate to a product region: for
\(m\ge2\) the zero-free region depends on the relative position of the
pairs, exactly as the coupling terms of Theorem~\ref{thm:f2v69-pairing}
predict.  The band region \(t_j\ge2m+p\) of V61 is a proper subregion of
\(\{|z_j|^2\ge2m+p\}\) and is not refuted; but any structural proof of
(H61) must use the couplings, not a per-pair modulus condition.

\subsection{Bridge ledger}
\label{sub:f2v69-ledger}

The archive's cofinal-visibility architecture (V80--V99 of the frozen
\texttt{V100\_f2} file) is a proof by contradiction whose arrows have
never been listed together.  Reading them together changes the role of
the finite-block theorems.  Fix a hypothetical zero
\(\rho_\star=\beta+i\gamma\) of \(L(s,\chi)\), \(\tfrac12<\beta<1\), and an
\emph{Euler-admissible} weight \(W\): nonnegative on the prime logarithms
beyond the conductor scale \(L_f\), annihilating the pole channel and
positive at the target.

\begin{description}[leftmargin=1.6em,style=nextline]
\item[\textup{(A1)} Forced compensation, proved (archive V80.13).]
For every admissible \(W\) the off-line zero forces the response of the
uncancelled two-channel zero and gamma tail to be at least
\(m_\rho\mathcal L_x[W]>0\).  The archive's own firewall: ``The present
work does not prove an upper bound contradicting it.''
\item[\textup{(A2)} Finite windows cannot contradict, proved (archive V81).]
Any fixed finite block of nearest zeros and gamma nodes can be cancelled
while the pole and the target normalization are kept; the compensation
escapes to the infinite tail.  So the contradiction, if any, must use a
\emph{growing} cancellation.
\item[\textup{(A3)} The growing filter, proved (archive V88).]
\(W=K_{M,d}\) cancels the pole, normalizes the target and kills the first
\(M\) nodes of every gamma ladder; \(K_{M,d}(u)/b_0=H_{M,d}(e^{-\varepsilon u})\)
with \(H_{M,d}\) an alternating polynomial whose real-rootedness is the
finite multiplier property of the ladder polynomial \(P_{M,d}\).
\item[\textup{(A4)} Real-rootedness kills admissibility, proved (archive V89, V97).]
If \(H_{M,d}\) is real-rooted then its last time-domain zero tends to
infinity and Euler positivity of \(K_{M,d}\) fails beyond \(L_f\); the
filter is inadmissible at that scale.  Proved for \(d/\sqrt M\to\infty\)
(V89) and on the critical lines \(d/\sqrt M\to\kappa\) with
\(\kappa>|\gamma|/\sqrt{1+\theta}\) (V90, V97 with real padding).
\item[\textup{(A5)} The fixed-block heat limit, proved (archive V97).]
On a critical line the grouped gamma block has the universal limit
\(e^{-(K/2)D^2}\Pi_{\bm\alpha}\), \(K=(1+\theta)/\theta^2\), with the
target pair at height \(y_{\min}^2=\gamma^2R/(\theta^2\kappa^2)\); real
zeros of the limit give the multiplier property for large \(M\), a
nonreal zero destroys it.
\item[\textup{(A6)} Band theorems, proved for \(m\le2\), conditional for \(m=3\) (V5, V9, V52--V61).]
A real-rooted critical envelope forces \(y_{\min}^2\le nK\), i.e.\
\(\kappa^2\ge\gamma^2R/(n(1+\theta))\ge\gamma^2/(1+\theta)\); by the V98
barrier no finite real padding certifies real-rootedness for
\(\kappa\le|\gamma|/\sqrt{2(1+\theta)}\).
\end{description}

\begin{proposition}[What the finite-block theorems do and do not do]
\label{prop:f2v69-ledger}
In the chain \textup{(A1)}--\textup{(A6)} the finite-block theorems enter
only through \textup{(A4)}--\textup{(A6)}, whose joint content is: on a
critical line with \(\kappa\) above the band the growing filter is
\emph{inadmissible}.  Closing the band, hence (H61), would extend this to
\(\kappa>|\gamma|/\sqrt{2(1+\theta)}\) for the padded equispaced family
and no further.  It would not produce a contradiction with
\textup{(A1)}.  A contradiction requires, at a surviving scale, an
admissible filter whose omitted-tail response is controlled, or an upper
bound on the compensation of \textup{(A1)}; neither exists, and the
conditioning obstacle named target-zero separation
(\(\prod_q|\rho_q-\rho_\star|\) has no lower bound) sits exactly there.
\end{proposition}

\begin{proof}
Reading of the cited statements; the only inference is that
\textup{(A4)} concludes inadmissibility, not falsity of the zero, and
that \textup{(A6)} is used in \textup{(A4)} only through the criterion of
\textup{(A5)}.
\end{proof}

The consequence for the programme is stated plainly.  Problem
\ref{prob:f2v68-bridge}(B1)--(B4) asked how an off-line zero yields an
excluded finite block; the ledger shows that in the archive's own
architecture no block is ever ``produced'' by the zero.  The block is the
prover's test filter, and the band theorems only say where that filter
cannot be Euler-positive.  The genuine open theorem is therefore

\begin{problem}[Admissible critical filter]
\label{prob:f2v69-admissible}
For \(\chi\), \(\rho_\star\) as above and some \(\theta\in(0,1)\), exhibit a
critical line \(d_M/\sqrt M\to\kappa\) with \(\kappa\le|\gamma|/\sqrt{1+\theta}\)
and a filter \(K_{M,d_M}\) (equispaced or not) such that for all large
\(M\): \textup{(i)} \(K_{M,d_M}(u)>0\) for \(u\ge L_f\), and
\textup{(ii)} the response of the omitted zero and gamma tail is
\(o(\mathcal L_x[K_{M,d_M}])\).  By \textup{(A1)} this contradicts the
off-line zero.
\end{problem}

Condition (i) is the open side of \textup{(A4)}: below the band top the
envelope is not real-rooted, so positivity is not excluded, but it is not
proved either.  Condition (ii) is where the infinite zero tail and
target-zero separation enter.  Nothing in V51--V69 bears on either.


\subsection*{Firewall and V70 direction}
\label{fire:f2v69}
Theorems~\ref{thm:f2v69-pairing} and \ref{thm:f2v69-circle},
Corollary~\ref{cor:f2v69-jacobi} and Proposition~\ref{prop:f2v69-witness}
are proved and verified.  The five residual groups of (H61) remain open,
the three-pair band for \(p\ge7\) is conditional, and every arrow marked
open in the ledger is unproved.  GRH remains open.

V70 is the mandatory companion sweep.  After it, the finite-side target
is the coupled region for \(m=2\): use \eqref{eq:f2v69-m2} to write
\(P_{p+1}=\widehat f_1(D)\widehat f_2(D)\mathrm{He}_{p+1}
-4(D-c_1)(D-c_2)\mathrm{He}_{p+1}+2\mathrm{He}_{p+1}\) and seek the
two-pair analogue of \eqref{eq:f2v69-disc}, a factorization of the
\(x\)-discriminant-type certificate of \(\Psi_2>0\) on \(t_j\ge p+4\) in
terms of \(|z_j|^2\) and the coupling \((c_1-c_2)^2+t_1+t_2\).  On the
bridge, the ledger's first open arrow is the one to attack.

\subsection*{Assumption ledger}
\label{ass:f2v69}
The proved statements use the Gaussian form of the heat operator, exact
rational arithmetic, and the V68 two-pair construction rebuilt from
source.  The three-pair failures are numerical; the ledger classifies
statements of the archive by reading, not by re-proof.
\endgroup

\section*{Version V69 append-only record}
\addcontentsline{toc}{section}{Version V69 append-only record}

V69 was appended byte-for-byte before the terminal marker of validated
V68, whose installed SHA--256 was
\[
 \texttt{47D6CE3F2563BCF491DDFBF70779D0E0}
 \allowbreak\texttt{C006D6BDCBFE23D50EAC3073D49D6EAD}
\]
with \(19020\) lines and \(659019\) bytes.  It proves the pairing
expansion of the heat image, the exact one-pair circle region and its
identification with the Jacobi criterion, refutes the modulus invariant
for two pairs by an exact witness, and records the bridge ledger.  After
validation V69 replaces V68 as the sole active numeric GRH note; the
frozen \texttt{V100\_f2} archive is unchanged.

% =============================================================================
% END APPENDED V69 MATERIAL
% =============================================================================




\clearpage
\section{V70: companion sweep, and the growing filter is not Euler-positive at any critical scale}
\label{sec:f2v70-filter}
\begingroup
\setcounter{equation}{0}
\renewcommand{\theequation}{N70.\arabic{equation}}
\renewcommand{\theHequation}{f2v70.\arabic{equation}}

This is the mandatory ten-version checkpoint (the fifth active version
since the V65 sweep).  After the sweep it attacks condition (i) of
Problem~\ref{prob:f2v69-admissible} directly, by exact computation of the
archive's growing filter at critical scales.  The outcome is negative and
sharp enough to change direction: in the equispaced two-channel family
the filter has a real time-domain zero that moves to infinity like
\(\sqrt M\) at every critical scale, above, inside and below the gamma
band.  So the band theorems, including the open (H61), have no target in
that family; the surviving question is upstream of them.

\subsection{Companion sweep}

The active companion sources read read-only were
\begin{equation}
\begin{array}{c|r|l}
\text{source}&\text{lines}&\text{SHA--256}\\ \hline
\text{SM--Einstein V29}&5526&
\texttt{EE8BCE49E9F0C8C3ED05B34E1A12B32D}\\
&&\texttt{C62C13FCE6A6CE703D62FB5FDBABD282}\\
\text{Yang--Mills V41}&14452&
\texttt{2EEC02EB12520BAA0651AB07A1719708}\\
&&\texttt{09451FB0DEB83B596B084312F9EE4961}\\
\text{Navier--Stokes V26}&5319&
\texttt{82C887F8C2C69E4F28FAD7C3DC8DE5B9}\\
&&\texttt{099CB5A4BF431EBA1FFF3E91935978AD}
\end{array}
\label{eq:f2v70-sources}
\end{equation}

\begin{proposition}[Sweep decision]
\label{prop:f2v70-sweep}
Since V65 the SM--Einstein note added V17--V29 (massive polarized
fermion stress and Einstein--Dirac metric, colour condensates, Yukawa and
Higgs exchange reduced to two-mode and Jaynes--Cummings systems, a
depleted two-channel conversion with elliptic period) and the
Yang--Mills note added V38--V41 (shell decomposition of the free
propagator, level-zero shell response, a tree-gauge representative and
the fixed-background convention); the Navier--Stokes note is unchanged.
None of the new statements concerns real-rootedness or multiplier
sequences, backward heat flow of zeros, Euler-product positivity or
Laplace transforms of rational filters, completed \(L\)-functions, or
exact positivity certificates in several variables; the SM V20 audit
states explicitly that there is no object in common with the GRH chain.
Nothing is imported.
\end{proposition}


\subsection{The filter at critical scales}

Recall the archive's construction (V88.1--V88.14 of the frozen file).
Fix \(\rho_\star=\beta+i\gamma\), \(\beta<\eta<1\), a band fraction
\(\theta\in(0,1)\), and for \(d>0\) put \(\sigma=d+1\), \(x_d=\sigma-\beta\),
\(\xi_d=\sigma-\eta\).  With the principal ladder \(\sigma+2m\) and the
target ladder \(\sigma+\kappa_\chi\pm i\gamma+2m\), \(0\le m<M\),
\(\kappa_\chi\in\{0,1\}\), \(R=3\), \(N=3M+4\), \(n=N-1\),
\(\varepsilon=\theta d/N\):
\begin{equation}
 P_{M,d}(y)=(y-d)(y-\xi_d)\prod_{m<M}(y-\sigma-2m)
 \prod_{m<M}\bigl((y-\sigma-\kappa_\chi-2m)^2+\gamma^2\bigr),
 \label{eq:f2v70-P}
\end{equation}
\begin{equation}
 \frac{K_{M,d}(u)}{b_0}=H_{M,d}(e^{-\varepsilon u}),\qquad
 H_{M,d}(q)=\sum_{j=0}^n(-1)^jc_jq^j,\qquad
 c_j=\binom nj\frac{P_{M,d}(-j\varepsilon)}{P_{M,d}(0)}>0 .
 \label{eq:f2v70-H}
\end{equation}
Euler admissibility requires \(K_{M,d}(u)>0\) for \(u\ge L_f\), the fixed
conductor scale; equivalently \(H_{M,d}>0\) on \((0,e^{-\varepsilon L_f}]\).
The critical lines are \(d=\kappa\sqrt M\); the band top is
\(|\gamma|/\sqrt{1+\theta}\) and the V98 barrier is
\(|\gamma|/\sqrt{2(1+\theta)}\).

\begin{proposition}[Exact zero certificates]
\label{prop:f2v70-certificates}
Take \(\gamma=14\), \(\beta=\tfrac34\), \(\eta=\tfrac9{10}\),
\(\theta=\tfrac12\) (band top \(11.43\), barrier \(8.08\)), both parities
\(\kappa_\chi\in\{0,1\}\), \(M\in\{6,10,16,24\}\) and
\(d=\kappa\sqrt M\) rounded to hundredths with \(\kappa\in\{14,10,6\}\)
(above the band, inside it, below the barrier).  For each of these
\(24\) instances there is a rational \(q_0\le\tfrac12\) with
\(H_{M,d}(q_0)<0\) exactly, so \(K_{M,d}\) has a zero beyond
\(u_0=-\ln q_0/\varepsilon>0\); the certified \(u_0\) increases with
\(M\) at every scale.  For \(\kappa_\chi=0\):
\begin{equation}
\begin{array}{c|ccc}
 M & \kappa=14 & \kappa=10 & \kappa=6\\ \hline
 6  & 1.35 & 1.73 & 2.45\\
 10 & 1.73 & 2.22 & 3.13\\
 16 & 2.16 & 2.78 & 3.93\\
 24 & 2.60 & 3.35 & 4.74
\end{array}
\qquad(u_0\text{, last-zero lower bounds})
\label{eq:f2v70-table}
\end{equation}
and the parity \(\kappa_\chi=1\) values differ in the second decimal.
Verified by \path{scripts/verify_grh_v70_filter_positivity.py}
(SHA--256 \texttt{F1462BBC07B0901A90D4820A35F40255}\allowbreak\texttt{6A3E2C4EBBF8F324BC2556499533C1F6}), which also checks
\(H_{M,d}(1)=0\) and \(c_j>0\).
\end{proposition}

\begin{proof}
Exact rational evaluation of \eqref{eq:f2v70-P}--\eqref{eq:f2v70-H} on the
grid \(q=k/1000\); the first negative value gives \(q_0\).
\end{proof}

\begin{remark}[What the numbers show]
\label{rem:f2v70-growth}
Exploration beyond the certified instances (up to \(M=48\), \(n=147\),
and \(\kappa=3\)) shows the mechanism.  Above and inside the band all but
two roots of \(H_{M,d}\) in \((0,1)\) are real, as the archive's finite
multiplier theorems predict for large \(M\); below the barrier the
number of nonreal roots grows slowly (nine of ninety-nine at \(\kappa=3\),
\(M=32\)), the localized conjugate prefix of the V90 verdict.  In every
case the smallest real root stays in \([0.31,0.49]\), bounded away from
\(q=1\), while \(\varepsilon=\theta\kappa/(3\sqrt M)(1+O(1/M))\to0\).
Hence the last time-domain zero obeys
\begin{equation}
 U_{M,d}\;\approx\;\frac{3\sqrt M}{\theta\kappa}\,\bigl(-\ln q_{\min}\bigr)
 \;\longrightarrow\;\infty
 \label{eq:f2v70-U}
\end{equation}
on every critical line, and it is \emph{larger} for smaller \(\kappa\):
the band is not where positivity is lost.
\end{remark}

\begin{proposition}[Positivity beyond a fixed scale needs all real roots at \(q=1\)]
\label{prop:f2v70-needs}
On a critical line, \(K_{M,d}>0\) on \([L_f,\infty)\) for all large \(M\)
holds only if every real root of \(H_{M,d}\) in \((0,1)\) lies in
\([1-\varepsilon L_f,1)\), an interval of length \(O(L_f/\sqrt M)\).  In
particular a single real root bounded away from \(1\) along the line, as
in every computed instance, gives \(U_{M,d}\ge c\sqrt M\) and defeats
admissibility for every fixed conductor.
\end{proposition}

\begin{proof}
\(K_{M,d}(u)/b_0=H_{M,d}(e^{-\varepsilon u})\) and \(u\ge L_f\) iff
\(q\le e^{-\varepsilon L_f}\); a real root \(q_1\le1-\delta\) is a zero of
\(K\) at \(u_1=-\ln q_1/\varepsilon\ge\delta/\varepsilon\).
\end{proof}

\begin{problem}[Equispaced-family conjecture]
\label{conj:f2v70-family}
For the equispaced two-channel filter with any finite real padding, on
every critical line \(d=\kappa\sqrt M\), \(H_{M,d}\) has a real root
\(q_1\le1-\delta(\kappa,\theta,\gamma)\) for all large \(M\).
\end{problem}

If Conjecture~\ref{conj:f2v70-family} holds, the archive's inadmissibility
statements (arrow \textup{(A4)} of the V69 ledger) extend to \emph{all}
critical scales of this family, and the gamma band, whose closure is the
purpose of V52--V61 and of (H61), is not the boundary of anything the
architecture can use: the family cannot host the contradiction with the
forced compensation at any scale.  The computed instances support the
conjecture at every tested scale and both parities.

\subsection*{Firewall and V71 direction}
\label{fire:f2v70}
Proposition~\ref{prop:f2v70-certificates} is exact for its \(24\)
instances; Proposition~\ref{prop:f2v70-needs} is proved;
Remark~\ref{rem:f2v70-growth} and Conjecture~\ref{conj:f2v70-family} are
numerical.  Nothing here proves or disproves GRH.  It does say that
closing (H61) would not advance the archive's contradiction, so the
finite-block effort is paused as a primary route.

V71 should go upstream of the filter: either prove
Conjecture~\ref{conj:f2v70-family} (the real principal ladder alone
already makes \(j\mapsto P(-j\varepsilon)/P(0)\) a multiplier sequence
whose smallest root is bounded away from \(1\); the target pairs only
add a bounded prefix), which retires the equispaced family; or design a
filter family whose alternating polynomial has all real roots within
\(O(1/\sqrt M)\) of \(q=1\), which is what admissibility beyond a fixed
conductor demands, and test it exactly at small \(M\) before any
asymptotics.  The finite-block statements keep their value only for the
second route, where a new family would need its own band.

\subsection*{Assumption ledger}
\label{ass:f2v70}
The sweep is read-only.  The certificates use exact rational arithmetic
on the archive's definitions; the parameter choices are illustrative and
the conjecture is not proved.
\endgroup

\section*{Version V70 append-only record}
\addcontentsline{toc}{section}{Version V70 append-only record}

V70 was appended byte-for-byte before the terminal marker of validated
V69, whose installed SHA--256 was
\[
 \texttt{1FC4490F5816480A88309D1DA0E33FCC}
 \allowbreak\texttt{EBE7BB4BA698D9BE7B147A02A1D6632F}
\]
with \(19330\) lines and \(674175\) bytes.  It records the companion
sweep, proves exact zero certificates for the growing filter at critical
scales above, inside and below the band, shows that admissibility beyond
a fixed conductor would need all real roots of the alternating polynomial
at \(q=1\), and redirects the programme upstream of the filter family.
After validation V70 replaces V69 as the sole active numeric GRH note;
the frozen \texttt{V100\_f2} archive is unchanged.

% =============================================================================
% END APPENDED V70 MATERIAL
% =============================================================================




\clearpage
\section{V71: last-zero bounds for the growing filter, and the pole factor as the source of the square-root rate}
\label{sec:f2v71-last-zero}
\begingroup
\setcounter{equation}{0}
\renewcommand{\theequation}{N71.\arabic{equation}}
\renewcommand{\theHequation}{f2v71.\arabic{equation}}

V70 showed by exact computation that the archive's growing filter has a
time-domain zero moving to infinity like \(\sqrt M\) on every critical
line.  This iteration proves what can be proved about that zero without
real-rootedness.  Two theorems result.  The first replaces the archive's
real-rootedness hypothesis by a count of roots outside \((0,1)\) and
gives a logarithmic rate; the second shows that the pole factor alone
already produces the square-root rate, and that every real ladder factor
only pushes the last zero further out, so the only thing standing between
the computed behaviour and a theorem is the effect of the conjugate
target ladder.  All algebraic ingredients are checked exactly by
\path{scripts/verify_grh_v71_last_zero_bound.py} (SHA--256
\texttt{B241E102ADFDED0A79854CEA58A57B8D}\allowbreak\texttt{17225CB1E4E893181E37B2E5F449CB4D}) on the \(24\) instances of
V70.

Notation as in \eqref{eq:f2v70-P}--\eqref{eq:f2v70-H}:
\(\varphi(s):=P_{M,d}(-s\varepsilon)/P_{M,d}(0)\), \(c_j=\binom nj\varphi(j)\),
\(H=\sum_j(-1)^jc_jq^j\), \(K(u)/b_0=H(e^{-\varepsilon u})\), \(N=n+1\),
and \(U_{M,d}\) the last zero of \(K\) on \((0,\infty)\).  Write
\(\Lambda\) for the multiset of roots of \(P_{M,d}\) (real:
\(d,\xi_d,\sigma+2m\); nonreal: the target ladder), and
\(S_{\rm real}:=\sum_{\lambda\in\Lambda\cap\mathbb R}1/\lambda\).

\subsection{A bound without real-rootedness}

\begin{theorem}[Last zero from the root product]
\label{thm:f2v71-product}
Let \(k'\) be the number of roots of \(H_{M,d}\), with multiplicity,
outside the open interval \((0,1)\), and suppose \(k'<n\).  Then
\begin{equation}
 \boxed{\;U_{M,d}\;\ge\;\Bigl(1-\frac\theta2\Bigr)S_{\rm real}
 -\frac{k'\,(\theta+\ln n)}{(n-k')\,\varepsilon}\;}
 \label{eq:f2v71-bound}
\end{equation}
In particular, on a critical line \(d=\kappa\sqrt M\),
\(S_{\rm real}\ge\tfrac12\ln\bigl(1+2M/(d+1)\bigr)\to\infty\), and if
\(k'=o(\sqrt M)\) then \(U_{M,d}\to\infty\): the filter is not
Euler-positive beyond any fixed conductor scale.  For real-rooted \(H\)
(\(k'=1\), the root at \(q=1\)) this is the archive's V89 conclusion with
the explicit rate \((1-\theta/2)\tfrac12\ln(2\sqrt M/\kappa)(1+o(1))\).
\end{theorem}

\begin{proof}
\(H\) has constant term \(c_0=1\) and leading coefficient \((-1)^nc_n\),
so the product of its roots is \(1/c_n=1/\varphi(n)\).  All \(c_j>0\), so
\(H\) has no root on \((-\infty,0]\), and by the Enestr\"om--Kakeya theorem
applied to \(\sum c_jz^j\) every root satisfies
\(|q|\ge\rho_0:=\min_jc_j/c_{j+1}\).  Here
\(c_j/c_{j+1}=\frac{j+1}{n-j}\prod_{\lambda}\frac{|\lambda+j\varepsilon|}{|\lambda+(j+1)\varepsilon|}
\ge\frac1n(1+\varepsilon/d)^{-(n-1)}\), since \(|\lambda+j\varepsilon|\ge d\)
for all \(\lambda\); and \((1+\varepsilon/d)^{n-1}\le e^{(n-1)\varepsilon/d}\le e^\theta\),
so \(\ln(1/\rho_0)\le\theta+\ln n\).  Let \(S\) be the set of roots in
\((0,1)\), \(|S|=n-k'\ge1\), and \(q_{\min}=\min S\).  Then
\(\prod_{S}q=\bigl(\varphi(n)\prod_{q\notin S}|q|\bigr)^{-1}
\le\bigl(\varphi(n)\rho_0^{k'}\bigr)^{-1}\), hence
\(-\ln q_{\min}\ge\bigl[\ln\varphi(n)-k'\ln(1/\rho_0)\bigr]/(n-k')\).
Next, each conjugate-pair factor of \(\varphi(n)\) is
\(|\lambda+n\varepsilon|^2/|\lambda|^2\ge1\), and for real
\(\lambda\ge d\) one has \(x=n\varepsilon/\lambda\le\theta<1\), so
\(\ln(1+x)\ge x(1-\theta/2)\); therefore
\(\ln\varphi(n)\ge(1-\theta/2)\,n\varepsilon\,S_{\rm real}\).  Finally
\(U=-\ln q_{\min}/\varepsilon\) and \(n/(n-k')\ge1\).  The critical-line
statements follow from \(\varepsilon=\theta d/N\) and
\(\sum_{m<M}(\sigma+2m)^{-1}\ge\tfrac12\ln((\sigma+2M)/\sigma)\).  The
verifier checks the ratio bound, the pair-factor inequality, \(x\le\theta\)
and the presence of a root in \((0,\tfrac12]\) exactly on every instance.
\end{proof}

The bound is weak at accessible \(M\): its main term is \(0.16\)--\(0.41\)
on the V70 instances against certified last zeros of \(1.3\)--\(4.7\), and
the \(k'\)-correction makes it negative there.  Its value is structural:
it shows that admissibility can survive only if the number of nonreal
roots of \(H\) grows at least like \(\sqrt M\), which is far from the
\(2\)--\(13\) nonreal roots observed.

\subsection{Where the square-root rate comes from}

In the time variable the alternating polynomial is a product of
first-order operators:
\begin{equation}
 \frac{K_{M,d}(u)}{b_0}=\prod_{\lambda\in\Lambda}\Bigl(1-\frac1\lambda\partial_u\Bigr)E(u),
 \qquad E(u):=(1-e^{-\varepsilon u})^n,
 \label{eq:f2v71-operators}
\end{equation}
because \(\mathcal D_q=q\,d/dq=-\varepsilon^{-1}\partial_u\) under
\(q=e^{-\varepsilon u}\), and the factors commute.

\begin{lemma}[Real factors push the last zero out]
\label{lem:f2v71-rolle}
Let \(f\) be a finite real exponential sum on \([0,\infty)\) with
\(f(\infty)\ne0\), and let \(u_L\) be its last zero.  For real
\(\lambda>0\), \((1-\lambda^{-1}\partial_u)f\) has a zero in
\((u_L,\infty)\).
\end{lemma}

\begin{proof}
\((1-\lambda^{-1}\partial_u)f=-\lambda^{-1}e^{\lambda u}\bigl(e^{-\lambda u}f\bigr)'\),
and \(e^{-\lambda u}f\) vanishes at \(u_L\) and at \(\infty\); Rolle.
\end{proof}

\begin{theorem}[Real-offset filters: square-root rate]
\label{thm:f2v71-real}
If all offsets are real (the principal ladder with any finite real
padding, no target pair), then
\begin{equation}
 \boxed{\;U_{M,d}\;\ge\;\frac1\varepsilon\ln\Bigl(1+\frac{n\varepsilon}d\Bigr)
 =\frac N{\theta d}\ln\Bigl(1+\frac{\theta n}N\Bigr)\;}
 \label{eq:f2v71-real}
\end{equation}
so on a critical line \(U_{M,d}\ge(1+o(1))\,\dfrac{R\ln(1+\theta)}{\theta\kappa}\sqrt M\)
with \(R\) the ladder count.
\end{theorem}

\begin{proof}
Apply the pole factor first:
\((1-d^{-1}\partial_u)E=(1-e^{-\varepsilon u})^{n-1}\bigl[(1-e^{-\varepsilon u})-\tfrac{n\varepsilon}de^{-\varepsilon u}\bigr]\),
whose last zero is at \(e^{-\varepsilon u}=(1+n\varepsilon/d)^{-1}\), i.e.\
\(u_0=\varepsilon^{-1}\ln(1+n\varepsilon/d)\); it has \(f(\infty)=1\).  Every
remaining factor is real, and by Lemma~\ref{lem:f2v71-rolle} each keeps a
zero beyond the current last zero, so \(U_{M,d}\ge u_0\).  With
\(N=RM+4\) and \(\varepsilon=\theta d/N\) the critical-line form follows.
\end{proof}

\begin{remark}[The obstruction is only the target pair]
\label{rem:f2v71-pairs}
In the \(q\)-variable the pole factor alone puts a root of
\((1+\tfrac\varepsilon d\mathcal D_q)(1-q)^n\) at
\(q_1=N/(N+\theta n)\approx1/(1+\theta)\), independent of \(M\); this is
the square-root rate of V70.  The conjugate-pair factors
\((1-\lambda^{-1}\partial_u)(1-\bar\lambda^{-1}\partial_u)
=|\lambda|^{-2}e^{\omega u}(\partial_u^2+\gamma^2)e^{-\omega u}\),
\(\omega=\operatorname{Re}\lambda\), are oscillatory operators to which
Rolle does not apply, so nothing proved here prevents them from pulling
the last zero back.  On every one of the \(24\) instances the certified
root \(q_0\) of V70 lies strictly below \(q_1\) (checked exactly), so
they do not: the observed \(q_{\min}\in[0.31,0.49]\) is even smaller than
the pole-factor value \(0.667\).  The V70 conjecture is therefore reduced
to a single statement.
\end{remark}

\begin{problem}[Pair factors do not retract the last zero]
\label{prob:f2v71-pairs}
Show that for \(f\) a real exponential sum with \(f(\infty)>0\), last
zero \(u_L\), and \(\omega\ge d\), \(\gamma\ne0\), the function
\(e^{\omega u}(\partial_u^2+\gamma^2)(e^{-\omega u}f)\) has a zero in
\([u_L-c/\omega,\infty)\) for an explicit \(c\); equivalently that a
conjugate-pair factor moves the last zero by at most \(O(1/\omega)\)
inward.  Since \(\sum_{\text{pairs}}1/\omega=O(\ln M)\) on a critical
line, this would give \(U_{M,d}\ge u_0-O(\ln M)\to\infty\) at the
square-root rate for the full two-channel filter, proving the V70
conjecture and retiring the equispaced family at every critical scale.
\end{problem}

\subsection*{Firewall and V72 direction}
\label{fire:f2v71}
Theorems~\ref{thm:f2v71-product} and \ref{thm:f2v71-real} and
Lemma~\ref{lem:f2v71-rolle} are proved; the verifier checks their
algebraic ingredients exactly.  Problem~\ref{prob:f2v71-pairs} is open,
the five residual (H61) groups are untouched, and GRH remains open.

V72 should attack Problem~\ref{prob:f2v71-pairs}: near the last zero
\(f\) is small and \(f'\) positive, so the pair operator equals
\(f-(2\omega/|\lambda|^2)f'+f''/|\lambda|^2\), negative at \(u_L\) unless
\(f''(u_L)\ge2\omega f'(u_L)\); a bound on \(f''/f'\) at the last zero for
exponential sums built from \eqref{eq:f2v71-operators}, with all
exponents in \([0,\theta d]\), would close it.  If that succeeds, the
equispaced family is retired and the programme must construct a filter
whose alternating polynomial keeps its real roots within
\(O(1/\sqrt M)\) of \(q=1\); if it fails, the failure locates the
scales where the family might be admissible.

\subsection*{Assumption ledger}
\label{ass:f2v71}
The theorems use the Enestr\"om--Kakeya theorem, Rolle's theorem, the
root-product identity, and \(\ln(1+x)\ge x(1-x/2)\) on \([0,1)\); the
verifier checks the instance-level inequalities exactly.  Root counts
\(k'\) are reported from exact sign scans and not used in any proof.
\endgroup

\section*{Version V71 append-only record}
\addcontentsline{toc}{section}{Version V71 append-only record}

V71 was appended byte-for-byte before the terminal marker of validated
V70, whose installed SHA--256 was
\[
 \texttt{C3D8D8507A1C5978E8790BF7BF48CF28}
 \allowbreak\texttt{A4BAEBBB52825C26AA4DF1A060D1A9E3}
\]
with \(19546\) lines and \(683967\) bytes.  It proves the root-product
last-zero bound without real-rootedness, the Rolle lemma for real
factors, the square-root rate for real-offset filters, and reduces the
V70 conjecture to the retraction problem for conjugate-pair factors.
After validation V71 replaces V70 as the sole active numeric GRH note;
the frozen \texttt{V100\_f2} archive is unchanged.

% =============================================================================
% END APPENDED V71 MATERIAL
% =============================================================================




\clearpage
\section{V72: the conjugate-pair factors step by step, the first-pair lemma, and the retraction inequality}
\label{sec:f2v72-pairs}
\begingroup
\setcounter{equation}{0}
\renewcommand{\theequation}{N72.\arabic{equation}}
\renewcommand{\theHequation}{f2v72.\arabic{equation}}

Problem~\ref{prob:f2v71-pairs} asks whether a conjugate-pair factor of the
growing filter can pull the last time-domain zero back.  This iteration
answers it exactly on instances, proves the first step in general, and
reduces the whole problem to one inequality at the current smallest root.
The answer on instances is: a pair factor can retract, but only by a
fraction of \(1/\omega\), and the cumulative effect of all pairs is a
push \emph{outward}.  All claims are checked by
\path{scripts/verify_grh_v72_pair_factor_steps.py} (SHA--256
\texttt{43F5ADF5499B6D20A3C502F5C1AEF591}\allowbreak\texttt{3B4345DBB88CE1E49C4DF77388914350}).

\subsection{Factor-by-factor form}

In the \(q\)-variable, \(H=\prod_{\lambda\in\Lambda}(1+\tfrac\varepsilon\lambda\mathcal D_q)(1-q)^n\)
and a factor multiplies the \(j\)-th coefficient by \(1+j\varepsilon/\lambda\);
a conjugate pair \(\lambda=\omega\pm i\gamma\) multiplies it by
\(|\lambda+j\varepsilon|^2/|\lambda|^2\).  The factors commute, so we may
apply the pole factor first,
\begin{equation}
 g_0(q)=(1-q)^{n-1}\Bigl(1-\frac q{q_1}\Bigr),\qquad q_1=\frac1{1+n\varepsilon/d},
 \label{eq:f2v72-g0}
\end{equation}
then the \(M\) pairs, then \(\xi_d\) and the real ladder.  Write
\(\Phi_\lambda(\mathcal D)=1+2\operatorname{Re}(a)\mathcal D+|a|^2\mathcal D^2\),
\(a=\varepsilon/\lambda\), for a pair factor.  For any real polynomial
\(g\) with \(g(0)>0\) and smallest positive root \(r\),
\begin{equation}
 \Phi_\lambda(\mathcal D)g\,(r)=\bigl(2\operatorname{Re}a+|a|^2\bigr)\,r\,g'(r)+|a|^2r^2g''(r),
 \label{eq:f2v72-atroot}
\end{equation}
since \(\mathcal D^2g=qg'+q^2g''\) and \(g(r)=0\).

\begin{lemma}[First-pair lemma]
\label{lem:f2v72-first}
For every pair \(\lambda=\omega+i\gamma\) with \(\omega>d(n-1)/n\),
\(\Phi_\lambda(\mathcal D)g_0\) has a root in \((0,q_1)\).
\end{lemma}

\begin{proof}
From \eqref{eq:f2v72-g0}, \(g_0'(q_1)=-(1-q_1)^{n-1}/q_1<0\) and
\(g_0''(q_1)=2(n-1)(1-q_1)^{n-2}/q_1\), so
\(q_1g_0''(q_1)/|g_0'(q_1)|=2(n-1)q_1/(1-q_1)=2(n-1)d/(n\varepsilon)\).
By \eqref{eq:f2v72-atroot}, \(\Phi_\lambda g_0(q_1)<0\) iff
\(q_1g_0''(q_1)/|g_0'(q_1)|<(2\operatorname{Re}a+|a|^2)/|a|^2=2\omega/\varepsilon+1\),
i.e.\ iff \((n-1)d/n<\omega+\varepsilon/2\), which holds.  Since
\(\Phi_\lambda g_0(0)=1>0\), there is a root in \((0,q_1)\).  The verifier
evaluates \eqref{eq:f2v72-atroot} exactly for every pair of every instance.
\end{proof}

\begin{proposition}[General step criterion]
\label{prop:f2v72-step}
Let \(g\) be the current polynomial (constant term \(1\)) with smallest
positive root \(r\), simple, \(g'(r)<0\).  A pair factor produces a root
in \((0,r)\) as soon as
\begin{equation}
 \frac{r\,g''(r)}{|g'(r)|}<\frac{2\omega}\varepsilon+1,
 \label{eq:f2v72-criterion}
\end{equation}
and a real factor always does.  Writing \(g'(r)\ne0\) and
\(g''(r)/g'(r)=2\sum_{\rho\ne r}(r-\rho)^{-1}\) over the other roots
\(\rho\) of \(g\), \eqref{eq:f2v72-criterion} reads
\(2r\operatorname{Re}\sum_{\rho\ne r}(\rho-r)^{-1}<2\omega/\varepsilon+1\):
the other roots must not crowd the smallest one at a scale finer than
\(\varepsilon r/\omega\).
\end{proposition}

\begin{proof}
\eqref{eq:f2v72-atroot} with \(g'(r)<0\), and the real-factor case is
\((1+a\mathcal D)g(r)=arg'(r)<0\).  The root-sum form is the logarithmic
derivative of \(g'\).
\end{proof}

Problem~\ref{prob:f2v71-pairs} is thus equivalent to controlling the
crowding of roots near the smallest one along the sequence of factors;
it is not a statement about a single operator.

\subsection{Exact step certificates}

\begin{proposition}[Pairs do not retract, up to \(2/\omega\)]
\label{prop:f2v72-steps}
For the instances \((M,\kappa,\kappa_\chi)\in\{(6,14,0),(6,6,1),(10,10,0),(10,6,0),(16,6,1)\}\)
of V70, after each factor a rational \(x\) with \(H_{\rm current}(x)<0\)
is exhibited, an exact upper bound for the current smallest root.  Along
the pairs this bound moves up only once (instance \((10,10,0)\), pair
\(m=2\)), by \(0.016\) in the time variable against \(1/\omega=0.027\);
every move up is certified below \(2/\omega\) through
\(x_{\rm new}\le x_{\rm old}(1+3\varepsilon/\omega)\).  The bound after
all pairs lies well below \(q_1\) (\(0.46\)--\(0.60\) against
\(0.67\)), and the final bound gives \(U_{M,d}\ge(0.55\text{--}1.00)\sqrt M\).
\end{proposition}

\begin{proof}
Exact rational arithmetic in the verifier; the inequality
\(e^{2t}\le1+3t\) for \(0\le t\le\tfrac15\) converts the time-domain
retraction bound \(2/\omega\) into the rational check.
\end{proof}

So the pair factors are, on every instance, net contributors to the
outward motion of the last zero: the observed shifts per pair are
\(+0.02\) to \(+0.09\) in \(u\), the same order as the real-factor shifts
\(1/\lambda\), with one retraction of \(-0.016\).  This is the
quantitative content behind the V70 growth law and it is consistent with
the perturbative decomposition
\(\Phi_\lambda=(1-\partial_u/\omega)^2+\frac{\gamma^2}{\omega|\lambda|^2}(2\partial_u-\partial_u^2/\omega)\)
of V71: two real factors plus a correction of relative size
\(\gamma^2/\omega^2\), which on a critical line is \(O(\gamma^2/(\kappa^2M))\)
for every node.

\subsection*{Firewall and V73 direction}
\label{fire:f2v72}
Lemma~\ref{lem:f2v72-first} and Proposition~\ref{prop:f2v72-step} are
proved; Proposition~\ref{prop:f2v72-steps} is exact for its instances.
Problem~\ref{prob:f2v71-pairs} is not solved in general: the step
criterion \eqref{eq:f2v72-criterion} must be verified along the whole
sequence, and the one observed retraction shows it can fail at a single
step while the cumulative motion is outward.  GRH remains open.

V73 should replace the step-by-step control by a global argument: apply
all pair factors at once to \(g_0\), i.e.\ show
\(\Psi(\mathcal D)g_0\,(q_1)<0\) with \(\Psi(j)=\prod_{\rm pairs}|\lambda+j\varepsilon|^2/|\lambda|^2\).
Since \(\Psi g_0(q_1)=-\sum_j\binom{n-1}j(-q_1)^j\,[\Psi(j+1)-\Psi(j)]\),
this is the positivity at \(q_1\) of the alternating polynomial with the
positive, increasing multiplier \(\Delta\Psi\); the real ladder then only
pushes outward (Lemma~\ref{lem:f2v71-rolle}).  A proof would establish
the V70 conjecture and retire the equispaced family at every critical
scale.

\subsection*{Assumption ledger}
\label{ass:f2v72}
Exact rational arithmetic on the archive's definitions; the smallest-root
locations quoted as shifts are bisection-refined sign changes (numerical),
while the certified bounds and the inequalities are exact.
\endgroup

\section*{Version V72 append-only record}
\addcontentsline{toc}{section}{Version V72 append-only record}

V72 was appended byte-for-byte before the terminal marker of validated
V71, whose installed SHA--256 was
\[
 \texttt{6C7AED043EC806DAEE2ADD4B9729D641}
 \allowbreak\texttt{87D41EAEBCAEDC5443DCA158AECB9EEA}
\]
with \(19755\) lines and \(694273\) bytes.  It proves the first-pair
lemma and the general step criterion, certifies step by step that the
conjugate-pair factors retract the last zero by less than \(2/\omega\)
and push it outward in total, and reduces the V70 conjecture to a single
positivity at the pole-factor root.  After validation V72 replaces V71 as
the sole active numeric GRH note; the frozen \texttt{V100\_f2} archive is
unchanged.

% =============================================================================
% END APPENDED V72 MATERIAL
% =============================================================================




\clearpage
\section{V73: the bump form of the growing filter, its Gaussian width, and the translation law of the pair factors}
\label{sec:f2v73-bump}
\begingroup
\setcounter{equation}{0}
\renewcommand{\theequation}{N73.\arabic{equation}}
\renewcommand{\theHequation}{f2v73.\arabic{equation}}

The V72 direction was to prove \(\Psi(\mathcal D)g_0(q_1)<0\) at the
pole-factor root.  Computed exactly, that inequality is true on some
instances and false on others, because the pair-applied function
oscillates around the reference point.  The correct global object is
found by conjugating out the pole factor: the filter is the derivative of
a bump, and the pole-factor point \(u_1\) is the bump's peak.  This makes
the conjugate-pair factors act as translations of a Gaussian of
\(M\)-independent width, gives a sufficient condition for
\(U_{M,d}\ge u_1\) that is exact-checkable, and identifies the real
ladder's push as the only missing quantitative input.  Checked by
\path{scripts/verify_grh_v73_bump_dichotomy.py} (SHA--256
\texttt{44EE06703F0C19F5C4351A196B8657A4}\allowbreak\texttt{338C174E192060422D14AE658265823D}).

\subsection{The bump}

With \(E(u)=(1-e^{-\varepsilon u})^n\) and \(\Psi(\partial_u)\) the product
of the conjugate-pair operators (so that \(\Psi\) multiplies \(e^{-j\varepsilon u}\)
by \(\prod_m|\lambda_m+j\varepsilon|^2/|\lambda_m|^2\)), put
\begin{equation}
 V:=e^{-du}E,\qquad Y:=\Psi(\partial_u)E,\qquad W:=e^{-du}Y,\qquad
 Z:=Y'-dY=e^{du}W' .
 \label{eq:f2v73-defs}
\end{equation}

\begin{proposition}[Bump identities]
\label{prop:f2v73-bump}
\begin{enumerate}[label=\textup{(\roman*)}]
\item \(V\) is positive on \((0,\infty)\), vanishes at \(0\) and \(\infty\),
and has a unique critical point, the maximum at
\(u_1=\varepsilon^{-1}\ln(1+n\varepsilon/d)\); \((1-\partial_u/d)E=-e^{du}V'/d\).
\item The pole-and-pairs filter is \((1-\partial_u/d)Y=-Z/d=-e^{du}W'/d\), so
its zeros are exactly the critical points of \(W\); and
\(W=Q(\partial_u+d)V\) with \(Q(y)=\prod_m(y-\lambda_m)(y-\bar\lambda_m)/|\lambda_m|^2\),
i.e.
\begin{equation}
 W=\prod_{m<M}\frac{(\partial_u-a_m)^2+\gamma^2}{|\lambda_m|^2}\,V,\qquad
 a_m:=\omega_m-d=1+\kappa_\chi+2m .
 \label{eq:f2v73-W}
\end{equation}
\item \(W>0\) for large \(u\), and the Laplace transform of \(Z\) is
\(c\,(y-d)\prod_m(y-\lambda_m)(y-\bar\lambda_m)/\prod_{j\le n}(y+j\varepsilon)\),
with \(y=d\) its only real zero.
\item Exactly, \((\ln V)''(u_1)=-d\,(d/n+\varepsilon)\), and
\((\ln V)'''(u_1)=d(1+q_1)(d/n+\varepsilon)^2\).  On a critical line
\(d=\kappa\sqrt M\) with \(R\) ladders,
\begin{equation}
 s^2:=-\frac1{(\ln V)''(u_1)}\;\longrightarrow\;\frac{R}{\kappa^2(1+\theta)},
 \qquad (\ln V)'''(u_1)=O(M^{-3/2}),
 \label{eq:f2v73-width}
\end{equation}
so \(V(u_1+st)/V(u_1)\to e^{-t^2/2}\) uniformly on compacts: the bump is
asymptotically a Gaussian of \(M\)-independent width.
\end{enumerate}
\end{proposition}

\begin{proof}
(i) \(V'=e^{-du}(E'-dE)\) and \(E'/E=n\varepsilon e^{-\varepsilon u}/(1-e^{-\varepsilon u})\)
decreases from \(\infty\) to \(0\), crossing \(d\) exactly at \(u_1\).
(ii) \(\Psi(\partial_u)\) commutes with \(\partial_u\); conjugation by
\(e^{-du}\) sends \(\partial_u\) to \(\partial_u+d\), and
\(\lambda_m-d=a_m\pm i\gamma\).  (iii) The \(j=0\) term of \(W\) is
\(e^{-du}\); the transform is the V88 partial fraction with the pole
factor and the pairs only.  (iv) \((\ln V)''=-n\varepsilon^2x/(1-x)^2\)
at \(x=q_1\), and \(n\varepsilon q_1/(1-q_1)=d\); the third derivative
is \(n\varepsilon^3x(1+x)/(1-x)^3\).  On the critical line \(d^2/n\to\kappa^2/R\)
and \(d\varepsilon\to\theta\kappa^2/R\).  The verifier checks (iv) exactly
on every instance.
\end{proof}

\subsection{Translation law and the dichotomy}

A factor \(((\partial-a)^2+\gamma^2)/(a^2+\gamma^2)\) applied to a narrow
bump moves its peak: at the new peak
\(V'''-2aV''+(a^2+\gamma^2)V'=0\), and with \(V'\approx-(u-u_1)V/s^2\),
\(V''\approx-V/s^2\) this gives
\begin{equation}
 u_{\rm peak}-u_1\;\approx\;\frac{2a}{a^2+\gamma^2}\qquad(>0).
 \label{eq:f2v73-translation}
\end{equation}
The exact per-pair shifts of the smallest root computed in V72 follow
this law: at \(M=10\), \(\kappa=6\), for \(a=1,3,\dots,19\) the observed
shifts are \(0.022,0.040,0.053,0.063,0.070,0.073,0.073,0.073,0.071,0.069\)
against \(0.010,0.029,0.045,0.057,0.065,0.069,0.071,0.071,0.070,0.068\);
the summed law \(\sum_m2a_m/(a_m^2+\gamma^2)\sim\ln M\) is the
logarithmic outward push of the pairs, and the \(\sqrt M\) of
\(U_{M,d}\) is entirely the pole factor's \(u_1\).

\begin{lemma}[Dichotomy]
\label{lem:f2v73-dichotomy}
If \(Y\) has a zero in \([u_1,\infty)\) (in particular if \(Y(u_1)\le0\)),
or if \(Y(u_1)>0\) and \(Y'(u_1)\ge dY(u_1)\), then the pole-and-pairs
filter has a zero at some \(u\ge u_1\), and hence \(U_{M,d}\ge u_1\).
\end{lemma}

\begin{proof}
In the first case \(W\) vanishes at some \(u^*\ge u_1\) and is positive
near \(\infty\), so it has a critical point beyond \(u^*\).  In the second
\(W'(u_1)=e^{-du_1}Z(u_1)\ge0\) while \(W\) decreases to \(0\) at
\(\infty\), so a critical point lies in \([u_1,\infty)\).  Critical points
of \(W\) are zeros of the pole-and-pairs filter by
Proposition~\ref{prop:f2v73-bump}(ii), and the remaining real factors push
the last zero outward by Lemma~\ref{lem:f2v71-rolle}.
\end{proof}

\begin{proposition}[Instances]
\label{prop:f2v73-instances}
On the \(24\) V70 instances the hypothesis of
Lemma~\ref{lem:f2v73-dichotomy} holds exactly in \(22\): nine with
\(Y(u_1)\le0\), nine with \(Y(u_1)>0\) and \(Z(u_1)\ge0\), four with a
certified later zero of \(Y\).  It fails only at \(M=6\), \(\kappa=10\),
both parities, where the pole-and-pairs filter alone has its last zero
before \(u_1\) and the seven real ladder factors supply the push (V70
certifies \(q_0\le0.383<q_1\) there).
\end{proposition}

\begin{proof}
Exact rational evaluation at \(q_1\); the later zero is certified by a
rational \(x\le q_1\) with \(Y(x)<0\).
\end{proof}

So the V72 test ``\(\Psi g_0(q_1)<0\)'' was the second alternative alone;
the first alternative, oscillation of \(Y\) beyond the peak, is what
usually holds at larger \(M\), where \(Y\) has zeros near \(u_1\) because
the translated bump's tail oscillates.  The one remaining gap is
quantitative: the real ladder factors push outward by an amount that
Rolle does not measure.  For a real factor the new last zero is the last
point where \((\ln f)'=\lambda\); a lower bound on that point in terms of
the log-derivative of \(f\) after its last zero would close the two
exceptional instances and, combined with the Gaussian limit
\eqref{eq:f2v73-width}, give the general theorem.

\subsection*{Firewall and V74 direction}
\label{fire:f2v73}
Proposition~\ref{prop:f2v73-bump} and Lemma~\ref{lem:f2v73-dichotomy} are
proved; Proposition~\ref{prop:f2v73-instances} is exact; the translation
law is a first-order computation confirmed numerically.  The V70
conjecture is not proved: its remaining content is that the hypothesis of
the dichotomy, or the real-ladder push, holds for all large \(M\) on
every critical line.  GRH remains open.

V74 should work in the Gaussian limit: with \(V\) replaced by the Gaussian
of width \(s\) from \eqref{eq:f2v73-width}, the pair operators are
explicit, and \(W\) is the inverse Fourier transform of
\(e^{-s^2\xi^2/2}\prod_m\bigl(1-2i\xi a_m/|\lambda'_m|^2-\xi^2/|\lambda'_m|^2\bigr)\),
\(\lambda'_m=a_m+i\gamma\).  Show that its last critical point is at
\(u_1+\sum_m2a_m/(a_m^2+\gamma^2)+o(1)\), uniformly in \(M\); then
transfer to finite \(M\) by the \(O(M^{-3/2})\) shape error.  The real
ladder can then be dropped entirely.

\subsection*{Assumption ledger}
\label{ass:f2v73}
Exact rational arithmetic on the archive's definitions; the Gaussian
limit uses only the two exact derivative formulas; the translation law
is heuristic beyond first order.
\endgroup

\section*{Version V73 append-only record}
\addcontentsline{toc}{section}{Version V73 append-only record}

V73 was appended byte-for-byte before the terminal marker of validated
V72, whose installed SHA--256 was
\[
 \texttt{C59313755957FE752BD4B8E4E70E9B6D}
 \allowbreak\texttt{1283C23B42BD16DFD2573EBA836E947D}
\]
with \(19920\) lines and \(701996\) bytes.  It proves the bump form of
the growing filter, the exact Gaussian width of the bump on critical
lines, the dichotomy lemma, and records the translation law of the
conjugate-pair factors with its exact instance checks.  After validation
V73 replaces V72 as the sole active numeric GRH note; the frozen
\texttt{V100\_f2} archive is unchanged.

% =============================================================================
% END APPENDED V73 MATERIAL
% =============================================================================




\clearpage
\section{V74: the Gaussian limit of the growing filter is the equal-height Appell symbol}
\label{sec:f2v74-limit}
\begingroup
\setcounter{equation}{0}
\renewcommand{\theequation}{N74.\arabic{equation}}
\renewcommand{\theHequation}{f2v74.\arabic{equation}}

The bump of V73 has, on every critical line, an \(M\)-independent Gaussian
profile.  Rescaling by its width turns the conjugate-pair operators into
exactly the operators \((\partial_\xi-c)^2+t\) of the center-dependent
symbol of V66--V67.  So the arithmetic side of the programme (Euler
positivity of the growing filter) and the finite-block side (real-axis
positivity of Appell polynomials) are statements about the same object,
with the same height parameter, and the gamma band is the height range
\(n<t<2n\).  This is the structural identification the programme review
asked for.  Exact claims are checked by
\path{scripts/verify_grh_v74_limit_symbol.py} (SHA--256
\texttt{35EDF9B0A7D1712E64E7AC3944C52162}\allowbreak\texttt{96CD4A54D886C993786BD6D9904088D6}).

\subsection{The identification}

\begin{theorem}[Rescaled limit]
\label{thm:f2v74-limit}
On a critical line \(d=\kappa\sqrt M\) with \(R\) ladders, let
\(s^2=1/(d(d/n+\varepsilon))\to R/(\kappa^2(1+\theta))\) be the bump width
of Proposition~\ref{prop:f2v73-bump} and \(\xi:=(u-u_1)/s\).  Then
\((\partial_u-a)^2+\gamma^2=s^{-2}\bigl[(\partial_\xi-sa)^2+s^2\gamma^2\bigr]\),
and the pole-and-pairs bump \(W\) of \eqref{eq:f2v73-W} converges, up to
a positive constant and uniformly on compacts in \(\xi\), to
\begin{equation}
 \boxed{\;F_M(\xi)=\prod_{m<M}\bigl[(\partial_\xi-c_m)^2+t\bigr]e^{-\xi^2/2}
 =e^{-\xi^2/2}G_M(\xi),\qquad
 t=\frac{\gamma^2R}{\kappa^2(1+\theta)},\quad c_m=s(1+\kappa_\chi+2m),}
 \label{eq:f2v74-F}
\end{equation}
where \(G_M\) is the polynomial produced by the V67 recursion
\(P\mapsto tP+A_c(A_cP)\), \(A_cP=P'-(\xi+c)P\), from \(P=1\).  The
height satisfies
\begin{equation}
 t=\frac{y_{\min}^2}{K},\qquad
 \text{band top }\kappa^2=\frac{\gamma^2}{1+\theta}\iff t=n,\qquad
 \text{barrier }\kappa^2=\frac{\gamma^2}{2(1+\theta)}\iff t=2n,
 \label{eq:f2v74-band}
\end{equation}
with \(n=R=2m+p\) the ladder count (\(m=1\) target pair, \(p\) real
ladders).  The zeros of the pole-and-pairs filter are the critical points
of \(W\), hence in the limit those of \(F_M\); the real ladder factors
become first-order factors \(\partial_\xi-s(\lambda-d)\) and only push the
last zero outward.
\end{theorem}

\begin{proof}
The operator identity is the substitution \(\partial_u=s^{-1}\partial_\xi\).
The Gaussian limit of \(V\) is \eqref{eq:f2v73-width}; the operators are
of fixed finite order for each \(m\), and the product over \(m<M\) is a
finite composition, so it commutes with the uniform limit on compacts.
The recursion is Theorem~\ref{thm:f2v67-symbol}.  The identities
\eqref{eq:f2v74-band} are \((\)N9.24\()\) \(y_{\min}^2=\gamma^2R/(\theta^2\kappa^2)\),
\(K=(1+\theta)/\theta^2\), and \(R=2m+p\) from \((\)N9.20\()\).
\end{proof}

\begin{remark}[What the identification says]
\label{rem:f2v74-meaning}
Above the band top the symbol height is below \(n\); numerically all
critical points of \(F_M\) are real there, which is the real-rooted
regime of the archive's finite multiplier theorems.  Below the barrier the
height exceeds \(2n\) and \(F_M\) is a single positive bump with one real
critical point.  Inside the band the two behaviours mix.  So the band of
V1--V61 is exactly the transitional height range of the symbol.  But the
admissibility question of Problem~\ref{prob:f2v69-admissible} concerns
only the \emph{last} critical point of \(F_M\), which is positive in all
three regimes: the filter's last zero is \(u_1+s\,\xi^*_M\) with
\(\xi^*_M>0\), whatever the real-rootedness status.  Real-rootedness and
admissibility are decoupled; closing the band would decide the former
and not the latter.
\end{remark}

\subsection{Exact moments and the translation law}

\begin{proposition}[Moments of the limit symbol]
\label{prop:f2v74-moments}
For any \(t>0\) and centers \(c_m\),
\begin{equation}
 \int_{\mathbb R}F_M\,\dd\xi=\sqrt{2\pi}\prod_{m<M}(t+c_m^2),\qquad
 \int_{\mathbb R}\xi F_M\,\dd\xi=\sqrt{2\pi}\prod_{m<M}(t+c_m^2)\sum_{m<M}\frac{2c_m}{c_m^2+t}.
 \label{eq:f2v74-moments}
\end{equation}
Hence the signed density \(F_M\) has positive total mass and mean
\(\sum_m2c_m/(c_m^2+t)\), which in the \(u\)-variable is
\(\sum_m2a_m/(a_m^2+\gamma^2)\): the translation law of V73 is exact for
the first moment.
\end{proposition}

\begin{proof}
The Fourier transform of \(F_M\) is
\(\widehat F_M(w)=\prod_m[(c_m-iw)^2+t]\,e^{-w^2/2}\) (each factor
\((\partial_\xi-c)^2+t\) becomes \((iw-c)^2+t\)); the moments are
\(\widehat F_M(0)\) and \(i\widehat F_M'(0)\).  Verified exactly over
\(\mathbb Q\) for \(M\le6\) with rational \(t,c_m\) by computing \(G_M\)
from the recursion and integrating against the Gaussian moments
\(\mathbb E[Z^{2l}]=(2l-1)!!\).
\end{proof}

\begin{proposition}[Limit against finite \(M\)]
\label{prop:f2v74-numerics}
For the \(24\) V70 instances, the last critical point \(\xi^*_M\) of
\(F_M\) and the finite-\(M\) last zero of the pole-and-pairs filter in
the units \((u-u_1)/s_M\) agree to a few percent for \(M\ge10\) and are
both positive in every instance, including the two where the dichotomy
hypothesis of V73 fails (there the finite-\(M\) value is small but
positive, \(1.5\) and \(2.9\), against limit values \(3.8\)).  Selected
values (\(\kappa_\chi=0\); limit / finite):
\begin{equation}
\begin{array}{c|ccc}
 M & \kappa=14\ (t=2.0) & \kappa=10\ (t=3.9) & \kappa=6\ (t=10.9)\\ \hline
 10 & 6.56/6.72 & 5.75/5.07 & 2.51/2.44\\
 24 & 10.45/11.47 & 9.44/9.58 & 5.58/5.48
\end{array}
\label{eq:f2v74-table}
\end{equation}
\end{proposition}

\begin{proof}
Numerical root-finding for the limit and exact bisection for the finite
polynomials, in the verifier; reported, not asserted beyond positivity.
\end{proof}

\subsection{The reduced statement}

\begin{problem}[Last critical point of the limit symbol]
\label{prob:f2v74-last}
Show that for every \(t>0\) and every increasing sequence of positive
centers \(c_0<c_1<\dots\), the function
\(F_M=e^{-\xi^2/2}G_M\) of \eqref{eq:f2v74-F} has a critical point at
some \(\xi\ge0\), with \(\xi^*_M\to\infty\) as \(M\to\infty\) when
\(\sum_m c_m/(c_m^2+t)=\infty\).
\end{problem}

A proof gives \(U_{M,d}\ge u_1+s\,\xi^*_M-o(1)\) on every critical line
(the finite-\(M\) shape error is \(O(M^{-3/2})\) by
\eqref{eq:f2v73-width}), hence the V70 conjecture and the retirement of
the equispaced family.  The moment identity is the first step: the
signed density has positive mean.  What is missing is a passage from
the mean to the last critical point, which for a positive unimodal
\(F_M\) (the regime \(t>2n\)) is a log-concavity or Newton-type
inequality, and in the oscillatory regimes needs the count of real
critical points of \(F_M\) as a function of \(t\), which is the same
question as the real-rootedness of the Appell sections in V62--V64.

\subsection*{Firewall and V75 direction}
\label{fire:f2v74}
Theorem~\ref{thm:f2v74-limit} and Proposition~\ref{prop:f2v74-moments}
are proved; Proposition~\ref{prop:f2v74-numerics} is numerical.
Problem~\ref{prob:f2v74-last} is open; nothing here proves GRH or the
V70 conjecture.  GRH remains open.

V75 is the mandatory companion sweep.  After it, attack
Problem~\ref{prob:f2v74-last} in the regime \(t>2n\) first, where
\(F_M\) is numerically a positive bump: show \(G_M>0\) on \(\mathbb R\)
and that \(\ln F_M\) is concave beyond its peak, so that the peak is at
least the mean minus an explicit width; then treat the band by the real
critical point count.

\subsection*{Assumption ledger}
\label{ass:f2v74}
The limit theorem uses the exact derivative formulas of V73 and finite
compositions of differential operators; the identification of the band
constants uses the archive's \((\)N9.20\()\), \((\)N9.24\()\) verbatim.
Root counts and the table are numerical.
\endgroup

\section*{Version V74 append-only record}
\addcontentsline{toc}{section}{Version V74 append-only record}

V74 was appended byte-for-byte before the terminal marker of validated
V73, whose installed SHA--256 was
\[
 \texttt{7D198AC1BFE180AD27416031BBD3683A}
 \allowbreak\texttt{45B63FC8B14B5762C16E7BB42BB4CB25}
\]
with \(20107\) lines and \(710750\) bytes.  It identifies the Gaussian
limit of the growing filter with the equal-height Appell symbol at
height \(y_{\min}^2/K\), locates the band as the height range
\(n<t<2n\), proves the exact moment identities that make the translation
law exact at first order, and reduces the V70 conjecture to the last
critical point of the limit symbol.  After validation V74 replaces V73
as the sole active numeric GRH note; the frozen \texttt{V100\_f2}
archive is unchanged.

% =============================================================================
% END APPENDED V74 MATERIAL
% =============================================================================




\clearpage
\section{V75: companion sweep, and the filter's last zero as one real zero under the backward heat flow}
\label{sec:f2v75-flow}
\begingroup
\setcounter{equation}{0}
\renewcommand{\theequation}{N75.\arabic{equation}}
\renewcommand{\theHequation}{f2v75.\arabic{equation}}

This is the mandatory checkpoint (fifth active version since V70).  After
the sweep, Problem~\ref{prob:f2v74-last} is transformed exactly into the
notes' own language: the critical points of the limit symbol are the
zeros of the backward-heat image of a block with one real zero and
\(M\) conjugate pairs, so the admissibility question is the motion of a
single real zero under the zero flow of V51.  Exact claims are checked by
\path{scripts/verify_grh_v75_heat_image_reformulation.py} (SHA--256
\texttt{D25DB4D272DE2AB3CCFEC50EA05E6C87}\allowbreak\texttt{E560285636DB8C69485E2F7EA7E6C92F}).

\subsection{Companion sweep}

The active companion sources read read-only were
\begin{equation}
\begin{array}{c|r|l}
\text{source}&\text{lines}&\text{SHA--256}\\ \hline
\text{SM--Einstein V33}&6201&
\texttt{BC8B4CB9ED17A7224B86F65631311456}\\
&&\texttt{CEDE52F1CA9A376968E996D48950EC95}\\
\text{Yang--Mills V41}&14452&
\texttt{2EEC02EB12520BAA0651AB07A1719708}\\
&&\texttt{09451FB0DEB83B596B084312F9EE4961}\\
\text{Navier--Stokes V26}&5319&
\texttt{82C887F8C2C69E4F28FAD7C3DC8DE5B9}\\
&&\texttt{099CB5A4BF431EBA1FFF3E91935978AD}
\end{array}
\label{eq:f2v75-sources}
\end{equation}

\begin{proposition}[Sweep decision]
\label{prop:f2v75-sweep}
Since V70 only the SM--Einstein note advanced, V30--V33: an audit
section, the resonant triads of the Yukawa and colour channels for
general wavenumbers with closed-form couplings and selection rules, and
a many-triad rotating-wave Hamiltonian whose polynomial invariants of
degree at most four span no fourth independent integral.  Yang--Mills
and Navier--Stokes are byte-identical to the V70 sweep.  None of the new
statements concerns real-rootedness, Hermite or Appell polynomials, heat
flow of zeros, Laplace transforms of rational filters, Euler-product
positivity, Dirichlet \(L\)-functions, or exact positivity certificates;
the SM V30 audit itself records that the GRH objects have nothing in
common with that cycle.  Nothing is imported.
\end{proposition}


\subsection{The heat-image form}

Let \(\mathcal P(\xi):=\prod_{m<M}\bigl[(\xi+c_m)^2+t\bigr]\), the
polynomial with zeros \(-c_m\pm i\sqrt t\), and let \(e^{-D^2/2}\) be the
operator with \(e^{-D^2/2}\xi^k=\mathrm{He}_k(\xi)\).

\begin{theorem}[Heat-image reformulation]
\label{thm:f2v75-heat}
\begin{enumerate}[label=\textup{(\roman*)}]
\item \(G_M=e^{-D^2/2}\mathcal P\): the V67 recursion output is the
backward-heat image of \(\mathcal P\), and the pairing expansion of
Theorem~\ref{thm:f2v69-pairing} applies to it verbatim.
\item \(G_M'-\xi G_M=-e^{-D^2/2}\bigl[\xi\,\mathcal P(\xi)\bigr]\).  Hence
the critical points of \(F_M=e^{-\xi^2/2}G_M\) are exactly the zeros of the
heat image of the block \(\xi\mathcal P(\xi)\), which has one real zero
at \(0\) and \(M\) conjugate pairs at the centers \(-c_m<0\), all to the
left of the real zero.
\item Let \(R_\tau:=e^{-\tau D^2/2}[\xi\mathcal P]\) and \(x(\tau)\) the
real zero of \(R_\tau\) starting at \(x(0)=0\).  Then
\begin{equation}
 \dot x(0)=\frac{R_0''(0)}{2R_0'(0)}=\frac{\mathcal P'(0)}{\mathcal P(0)}
 =\sum_{m<M}\frac{2c_m}{c_m^2+t}>0,
 \label{eq:f2v75-velocity}
\end{equation}
the mean of Proposition~\ref{prop:f2v74-moments}; and
Problem~\ref{prob:f2v74-last} is the statement \(x(1)\ge0\) for the
largest real zero of \(R_1\).
\end{enumerate}
\end{theorem}

\begin{proof}
(i) In the Hermite basis \(\mathrm{He}_k=e^{-D^2/2}\xi^k\) the operator
\(A_c=\partial-\xi-c\) of the V67 recursion acts as
\(\mathrm{He}_k\mapsto-\mathrm{He}_{k+1}-c\mathrm{He}_k\)
(Lemma~\ref{lem:f2v66-raise}), i.e.\ as \(-( X+c)\) on the formal variable
\(X\) with \(X^k\leftrightarrow\mathrm{He}_k\); so
\(\prod_m[(A_{c_m})^2+t]1\) corresponds to \(\prod_m[(X+c_m)^2+t]\).
(ii) \([e^{-D^2/2},\xi]=-De^{-D^2/2}\), so
\(e^{-D^2/2}(\xi\mathcal P)=\xi G_M-G_M'\).  (iii) For a simple real zero
\(x\) of \(R_\tau\), differentiating \(R_\tau(x(\tau))=0\) with
\(\partial_\tau R_\tau=-\tfrac12R_\tau''\) gives \(\dot x=R_\tau''(x)/(2R_\tau'(x))\);
at \(\tau=0\), \(R_0=\xi\mathcal P\) has \(R_0'(0)=\mathcal P(0)\),
\(R_0''(0)=2\mathcal P'(0)\).  All three are verified exactly over
\(\mathbb Q\) for \(M\le6\).
\end{proof}

\begin{remark}[The flow picture]
\label{rem:f2v75-flow}
By V51, under \(\tau\mapsto e^{-\tau D^2/2}\) the zeros move by
\(\dot z_k=\sum_{l\ne k}(z_k-z_l)^{-1}\).  For the real zero this is
\(\dot x=\sum_{\text{pairs}}2(x-x_l)/((x-x_l)^2+y_l^2)+\sum_{\text{real}}(x-r)^{-1}\),
positive as long as \(x\) is the rightmost real part of the zero set,
which it is at \(\tau=0\).  The pairs are pushed left by the real zero
and right by the other pairs; the relative motion is not controlled by
the ordering argument alone.  Numerically (verifier, all \(24\)
instances) \(x(\tau)\) is nondecreasing on \([0,1]\), nearly linear with
slope close to the mean \eqref{eq:f2v75-velocity}, and \(x(1)\) coincides
with the last critical point \(\xi^*_M\) of V74 to all printed digits:
e.g.\ \(M=24\), \(\kappa=6\): \(x(\tau)=0,0.54,1.36,2.75,4.15,5.58\) at
\(\tau=0,0.1,0.25,0.5,0.75,1\), mean \(5.41\).
\end{remark}

\begin{problem}[Monotone escape of the real zero]
\label{prob:f2v75-monotone}
Prove that for \(t>0\) and centers \(0<c_0<c_1<\dots\), the real zero of
\(R_\tau=e^{-\tau D^2/2}[\xi\prod_m((\xi+c_m)^2+t)]\) that starts at
\(0\) satisfies \(\dot x(\tau)\ge0\) on \([0,1]\), or at least
\(x(1)\ge0\).  This is Problem~\ref{prob:f2v74-last}, hence the V70
conjecture: on every critical line the growing filter's last zero is
\(u_1+s\,x(1)+o(1)\ge u_1\), and the equispaced family is never
Euler-positive beyond a fixed conductor.
\end{problem}

The problem is now of exactly the type treated in V51--V61 (one real
zero, pairs on one side, backward heat for unit time), and the tools
there apply: the rates \(\frac{\dd}{\dd\tau}\sum x^2=\sum\cos^2\theta_{kl}\),
\(\frac{\dd}{\dd\tau}\sum y^2=-\sum\sin^2\theta_{kl}\), the Pólya strip
lemma, and the total-height budget.  The one-sided geometry is new: all
pairs start strictly left of the real zero, at spacing \(2s\) and height
\(\sqrt t\), and the question is whether they can overtake it within
unit time.

\subsection*{Firewall and V76 direction}
\label{fire:f2v75}
The sweep imports nothing.  Theorem~\ref{thm:f2v75-heat} is proved and
verified; the trajectory statements are numerical;
Problem~\ref{prob:f2v75-monotone} is open.  GRH remains open.

V76 should attack Problem~\ref{prob:f2v75-monotone} with the V51 rates:
bound the rightward drift of the rightmost pair by comparison with the
real zero's drift, using that the real zero repels each pair with force
\((x-x_m)/|x-z_m|^2\) while a pair repels the real zero with twice that,
and that the pair heights decrease (so the pairs approach the axis and
their mutual pushes weaken relative to the real zero's).  A first
rigorous target is \(x(1)\ge0\) when \(M=1\), then \(M=2\), where the
motion is explicit, before the general one-sided configuration.

\subsection*{Assumption ledger}
\label{ass:f2v75}
The reformulation uses the Hermite raising operator of V66, the
commutator of \(e^{-D^2/2}\) with \(\xi\), and the V51 zero-flow
velocity; the trajectory table is numerical.
\endgroup

\section*{Version V75 append-only record}
\addcontentsline{toc}{section}{Version V75 append-only record}

V75 was appended byte-for-byte before the terminal marker of validated
V74, whose installed SHA--256 was
\[
 \texttt{163675E6F54299825C920F545DA2055D}
 \allowbreak\texttt{04FE2417EF8DF77D65675EA2C6EA3F44}
\]
with \(20304\) lines and \(719735\) bytes.  It records the companion
sweep, proves that the critical points of the limit symbol are the zeros
of the backward-heat image of a one-real-zero block, computes the exact
initial velocity of that real zero, and states the monotone-escape
problem.  After validation V75 replaces V74 as the sole active numeric
GRH note; the frozen \texttt{V100\_f2} archive is unchanged.

% =============================================================================
% END APPENDED V75 MATERIAL
% =============================================================================




\clearpage
\section{V76: the escape problem for one and two pairs, the Hermite-coefficient form, and the sign at the origin}
\label{sec:f2v76-escape}
\begingroup
\setcounter{equation}{0}
\renewcommand{\theequation}{N76.\arabic{equation}}
\renewcommand{\theHequation}{f2v76.\arabic{equation}}

Problem~\ref{prob:f2v75-monotone} asks whether the real zero of
\(R_\tau=e^{-\tau D^2/2}[\xi\mathcal P]\), \(\mathcal P=\prod_{m<M}[(\xi+c_m)^2+t]\),
that starts at \(0\) is still at a nonnegative position at \(\tau=1\).
This iteration settles the one-pair case, gives the exact two-pair
criterion, and reduces the general statement to the sign of one explicit
alternating sum whenever \(R_1\) has a single real zero, which is the
regime below the V98 barrier.  Exact claims are checked by
\path{scripts/verify_grh_v76_small_M_escape.py} (SHA--256
\texttt{222A7E06B606C6FFAF122F641BC12527}\allowbreak\texttt{E464116584C3A3D6CFD7E3A570CDCA92}).

\subsection{Hermite form and the value at the origin}

Write \(\mathcal P=\sum_kg_k\xi^k\); since \(c_m,t>0\), every \(g_k>0\),
and \(\mathcal P\) has all its zeros in the open left half-plane.

\begin{proposition}[Hermite-coefficient form]
\label{prop:f2v76-hermite}
\begin{enumerate}[label=\textup{(\roman*)}]
\item \(R_1=\sum_kg_k\mathrm{He}_{k+1}(\xi)\), and
\(R_1(0)=\sum_{l\ge0}(-1)^{l+1}(2l+1)!!\,g_{2l+1}=-G_M'(0)=-F_M'(0)\).
\item \(R_\tau(0)=\sum_{k\ge1}\bigl(-\tfrac\tau2\bigr)^k\dfrac{(2k)!}{k!}\,g_{2k-1}\),
and \(R_\tau\) is monic of degree \(2M+1\).
\item (Hermite--Poulain) the number of real zeros of \(R_\tau\) is
nondecreasing in \(\tau\); it is odd and at least \(1\).
\end{enumerate}
\end{proposition}

\begin{proof}
(i) \(e^{-D^2/2}\xi^{k+1}=\mathrm{He}_{k+1}\), \(\mathrm{He}_{2l}(0)=(-1)^l(2l-1)!!\),
\(\mathrm{He}_{\rm odd}(0)=0\); the last two equalities are
Theorem~\ref{thm:f2v75-heat}(ii) at \(\xi=0\).  (ii) Only the terms
\(D^{2k}\xi^{2k}\) of \(e^{-\tau D^2/2}\) survive at the origin.  (iii)
\(e^{-\tau D^2/2}\) is a limit of products of the operators
\(1\pm aD\), each of which does not increase the number of nonreal zeros
(Hermite--Poulain); \(R_0=\xi\mathcal P\) has exactly one real zero and
odd degree.  (i), (ii) are verified exactly for \(M\le6\).
\end{proof}

\subsection{One and two pairs}

\begin{theorem}[One pair escapes]
\label{thm:f2v76-onepair}
For \(M=1\) and all \(c,t>0\),
\begin{equation}
 R_1(\xi)=\xi^3+2c\xi^2+(c^2+t-3)\xi-2c,\qquad R_1(0)=-2c<0,
 \label{eq:f2v76-onepair}
\end{equation}
so \(R_1\) has a positive real zero and \(x(1)>0\).  More precisely
\(R_\tau(0)=-2c\tau\), so the real zero is strictly positive for every
\(\tau\in(0,1]\).
\end{theorem}

\begin{proof}
Direct expansion, \(e^{-\tau D^2/2}\) acting on \(\xi^3+2c\xi^2+(c^2+t)\xi\);
verified exactly.  A monic real cubic negative at \(0\) has a zero in
\((0,\infty)\).
\end{proof}

\begin{proposition}[Two pairs]
\label{prop:f2v76-twopairs}
For \(M=2\), \(R_1(0)=-2(c_1+c_2)(c_1c_2+t-3)\).  Hence \(x(1)>0\)
whenever \(c_1c_2+t>3\); in particular for all centers when \(t>3\).
\end{proposition}

\begin{proof}
\(g_1=(c_1^2+t)2c_2+(c_2^2+t)2c_1\), \(g_3=2(c_1+c_2)\), and
Proposition~\ref{prop:f2v76-hermite}(i) gives \(R_1(0)=-g_1+3g_3\);
verified exactly.
\end{proof}

\subsection{The single-real-zero regime}

\begin{proposition}[Escape is the sign at the origin when the real zero is alone]
\label{prop:f2v76-alone}
If \(R_1\) has exactly one real zero, then it is positive if and only if
\(R_1(0)<0\), i.e.\ \(G_M'(0)>0\).  A sufficient condition, for any
number of real zeros, is \(R_1(0)<0\).
\end{proposition}

\begin{proof}
\(R_1\) is monic with positive leading coefficient; with a single real
zero \(x\) it factors as \((\xi-x)\) times a positive quadratic product,
so \(R_1(0)=-x\cdot(\text{positive})\).  In general \(R_1(0)<0\) and
\(R_1(+\infty)=+\infty\) give a positive zero.
\end{proof}

On the V70 parameter set, \(R_1\) has exactly one real zero in every
instance below the barrier (\(t=10.89>2n\)) except one (\(M=24\),
\(\kappa=6\), \(\kappa_\chi=0\), three real zeros), and there
\(R_1(0)<0\) always; inside and above the band \(R_1\) has many real
zeros (the real-rooted regime of V74), of which roughly half are
nonnegative, and \(R_1(0)<0\) in \(17\) of the \(24\) instances overall.
The largest real zero is positive in all \(24\).  So below the barrier
the escape problem is the explicit inequality
\begin{equation}
 \boxed{\;G_M'(0)=\sum_{l\ge0}(-1)^l(2l+1)!!\,g_{2l+1}>0,\qquad
 \mathcal P=\prod_{m<M}\bigl[(\xi+c_m)^2+t\bigr]=\sum_kg_k\xi^k,\;}
 \label{eq:f2v76-target}
\end{equation}
equivalently \(\mathbb E[\operatorname{Re}\mathcal P'(iZ)]>0\) for a
standard normal \(Z\), by Gaussian integration by parts.

\begin{problem}[Positive slope at the origin]
\label{prob:f2v76-slope}
Prove \eqref{eq:f2v76-target} for the equispaced centers
\(c_m=s(1+\kappa_\chi+2m)\) with \(t\ge2n\) and all \(M\), and show that
\(R_1\) has a single real zero there.  Together with
Proposition~\ref{prop:f2v76-alone} this proves
Problem~\ref{prob:f2v75-monotone} below the barrier; the band and the
region above it need the count of nonnegative real zeros instead.
\end{problem}

\subsection*{Firewall and V77 direction}
\label{fire:f2v76}
Theorem~\ref{thm:f2v76-onepair} settles the one-pair escape;
Propositions~\ref{prop:f2v76-hermite}, \ref{prop:f2v76-twopairs},
\ref{prop:f2v76-alone} are proved.  The general escape, hence the V70
conjecture, is open.  GRH remains open.

V77 should prove \eqref{eq:f2v76-target} for large \(t\): expand
\(\mathbb E[\mathcal P'(iZ)]\) in the pairing form of
Theorem~\ref{thm:f2v69-pairing}, where the leading term is
\(\mathcal P'(0)\prod_m(1+O(1/t))>0\) and the corrections carry the
factors \(\mu(a,b)\) with at least two derivatives, hence are smaller by
\(O(M/t)\); make the constant explicit and compare with \(t\ge2n=2(3M+3)\).
If the bound is \(t>CM\) with \(C\) below \(6\), the equispaced family is
retired below the barrier for all \(M\).

\subsection*{Assumption ledger}
\label{ass:f2v76}
Exact rational arithmetic for the identities; Hermite--Poulain is cited
as classical; the instance survey is numerical.
\endgroup

\section*{Version V76 append-only record}
\addcontentsline{toc}{section}{Version V76 append-only record}

V76 was appended byte-for-byte before the terminal marker of validated
V75, whose installed SHA--256 was
\[
 \texttt{4B9DCA10A1FB3F567D3E06772BC96EA5}
 \allowbreak\texttt{2BFBD5F41DF8C048F8B7F5A632A34013}
\]
with \(20483\) lines and \(728105\) bytes.  It proves the one-pair
escape, the two-pair criterion, the Hermite-coefficient form of the
heat image with its value at the origin, and the reduction of the escape
problem below the barrier to a positive slope at the origin.  After
validation V76 replaces V75 as the sole active numeric GRH note; the
frozen \texttt{V100\_f2} archive is unchanged.

% =============================================================================
% END APPENDED V76 MATERIAL
% =============================================================================




\clearpage
\section{V77: invariants of the zero flow, overtaking pairs, and why the perturbative slope bound fails}
\label{sec:f2v77-flow}
\begingroup
\setcounter{equation}{0}
\renewcommand{\theequation}{N77.\arabic{equation}}
\renewcommand{\theHequation}{f2v77.\arabic{equation}}

V76 proposed to prove the slope inequality \(G_M'(0)>0\) by a large-\(t\)
expansion whose leading term is \(\mathcal P'(0)\).  Computed exactly,
that leading term is not leading: the Gaussian average of the pair
product carries the phase \(A=\sum_m2c_m/(c_m^2+t)\), which is
unbounded in \(M\), and the true value is smaller than \(\mathcal P'(0)\)
by a factor of order \(e^{-A^2/2}\).  This iteration records that
negative result exactly, proves the invariants of the zero flow that
govern the escape problem, and reports what the flow actually does:
pairs overtake the real zero, and in the band and above the whole
configuration becomes real-rooted before unit time.  Checked by
\path{scripts/verify_grh_v77_flow_invariants.py} (SHA--256
\texttt{2625CB38A844DCEFB7D9181698945323}\allowbreak\texttt{A7524ACFB9E32AB09339B86452284FBA}).

\subsection{Invariants}

\begin{proposition}[Zero-sum conservation and the velocity identity]
\label{prop:f2v77-invariants}
Let \(R_\tau=e^{-\tau D^2/2}[\xi\mathcal P]\) with zeros \(x(\tau)\) (the
real zero starting at \(0\)) and \(x_m(\tau)\pm iy_m(\tau)\).
\begin{enumerate}[label=\textup{(\roman*)}]
\item The two leading coefficients of \(R_\tau\) do not depend on
\(\tau\); hence \(x(\tau)+2\sum_mx_m(\tau)=-2\sum_mc_m\) for all
\(\tau\), and
\begin{equation}
 x(\tau)=2\sum_{m<M}\bigl(x_m(0)-x_m(\tau)\bigr).
 \label{eq:f2v77-sum}
\end{equation}
\item Under the flow \(\dot z_k=\sum_{l\ne k}(z_k-z_l)^{-1}\) the
pair--pair forces cancel in the sum of real parts:
\(\sum_m\dot x_m=\sum_m(x_m-x)/|x-z_m|^2=-\tfrac12\dot x\).
\item While \(R_\tau\) has the single real zero \(x\),
\(\dot x=\sum_m2(x-x_m)/|x-z_m|^2\); pairs to the left of \(x\) push it
right, pairs to the right pull it left with force at most \(1/y_m\) each.
\item Once \(R_{\tau_0}\) is real-rooted, \(R_\tau\) is real-rooted for
\(\tau\ge\tau_0\) and its largest zero is nondecreasing.
\end{enumerate}
\end{proposition}

\begin{proof}
(i) \(e^{-\tau D^2/2}\) alters only coefficients of degree \(\le\deg-2\);
the zero sum is minus the subleading coefficient.  (ii) For \(m\ne l\),
\(\operatorname{Re}[(z_m-z_l)^{-1}]\) is antisymmetric and
\(\operatorname{Re}[(z_m-\bar z_l)^{-1}]+\operatorname{Re}[(z_l-\bar z_m)^{-1}]=0\);
the conjugate contributes \(\operatorname{Re}[(2iy_m)^{-1}]=0\).  (iii)
is the flow with the real parts paired; \(2u/(u^2+y^2)\le1/y\).  (iv)
\(e^{-(\tau-\tau_0)D^2/2}\) is a Laguerre--Pólya operator, so it
preserves real-rootedness, and the largest zero then has
\(\dot x=\sum(x-r_j)^{-1}>0\).  (i) is verified exactly for \(M\le6\).
\end{proof}

By \eqref{eq:f2v77-sum} the real zero escapes exactly when the pairs'
real parts move left in total, and by (ii) this is the same as the real
zero's own repulsion: the pair--pair dynamics enter only through the
weights \(|x-z_m|^{-2}\).

\subsection{What the flow does}

At \(M=10\) on the V70 parameters (verifier):
\begin{itemize}
\item \(t=10.89\) (below the barrier): the real zero moves
\(0\to0.60\to1.23\to1.86\to2.51\) at \(\tau=0,\tfrac14,\tfrac12,\tfrac34,1\);
the pair heights fall from \(3.30\) to about \(1.5\)--\(2.0\); the pair
real parts spread nearly antisymmetrically, and from \(\tau=\tfrac14\) on
exactly one pair sits to the right of the real zero.  No collision
occurs, \(R_1\) has a single real zero, and \(G_M'(0)>0\).
\item \(t=3.92\) (in the band): collisions begin near \(\tau=0.6\), all
zeros are real from \(\tau=0.9\), and the largest real zero at
\(\tau=1\) is \(5.75\), no longer the zero that started at \(0\).
\item \(t=2.00\) (above the band top): all zeros are real from
\(\tau\approx0.48\); the largest real zero then increases from \(4.0\) to
\(6.56\) by Proposition~\ref{prop:f2v77-invariants}(iv).
\end{itemize}
So overtaking does happen, even below the barrier, and the escape is
carried by the many pairs left of the real zero against the few to its
right; in the band and above, the escape is decided by the real-rooted
phase after total collision.

\subsection{The perturbative slope bound fails}

\begin{proposition}[No leading-term expansion]
\label{prop:f2v77-noleading}
Write \(G_M'(0)=\mathbb E[\operatorname{Re}\mathcal P'(iZ)]\) and
\(A:=\sum_m2c_m/(c_m^2+t)\).  On the \(M=10\) instances,
\(G_M'(0)/\mathcal P'(0)=0.079\) at \(t=10.89\) (\(A=2.36\),
\(e^{-A^2/2}=0.061\)), \(-1.29\) at \(t=3.92\) (\(A=3.94\)), and
\(3.3\times10^3\) at \(t=2.00\) (\(A=5.51\)); the pair product
\(\prod_m(1+w_m(Z))\) with \(w_m=(-Z^2+2ic_mZ)/(c_m^2+t)\) has phase
\(\approx AZ\), and its Gaussian average is of size \(e^{-A^2/2}\) rather
than \(1\).  Since \(A\sim s^{-1}\ln M\) on every critical line, no
bound of the form \(t>CM\) can be obtained from the expansion around
\(\mathcal P'(0)\); the V76 direction is withdrawn.
\end{proposition}

\begin{proof}
Exact evaluation of \(R_1(0)=-G_M'(0)\) against \(\mathcal P'(0)\), and
the estimate \(\prod(1+w_m)=\exp(iAZ+O(Z^2S_0))\) with
\(S_0=\sum(c_m^2+t)^{-1}\), whose Gaussian average is
\(e^{-A^2/2}(1+O(S_0))\) for bounded \(S_0\).
\end{proof}

\begin{remark}[Where a proof must come from]
\label{rem:f2v77-where}
Below the barrier the pairs do not collide in unit time and the real
zero is alone; \eqref{eq:f2v77-sum} and (iii) reduce the escape to a
lower bound on \(\sum_m2(x-x_m)/|x-z_m|^2\) along the trajectory, i.e.\
to a bound on how many pairs can be to the right of \(x\) and how close.
Above the band top, total collision occurs before \(\tau=1\) and the
escape is the largest real zero at the collision time, which the
real-rooted monotonicity then carries to \(\tau=1\).  The band
interpolates.  In all three regimes the needed input is dynamical, not
perturbative: control of the pair real parts and heights over unit time
for an equispaced initial line, for which the V51 rates
\(\frac{\dd}{\dd\tau}\sum y_m^2=-M-\sum(\ldots)\) and the spread
\(\frac{\dd}{\dd\tau}\sum x_k^2=\sum\cos^2\theta_{kl}\) are the
available tools.
\end{remark}

\subsection*{Firewall and V78 direction}
\label{fire:f2v77}
Proposition~\ref{prop:f2v77-invariants} is proved;
Proposition~\ref{prop:f2v77-noleading} is exact on its instances and
closes the V76 route.  Problem~\ref{prob:f2v75-monotone} stays open in
every regime.  GRH remains open.

V78 should treat the regime below the barrier with the invariants: show
that at most \(K(t,s)\) pairs, independent of \(M\), can lie to the right
of \(x\) at any time (the real zero and the front pairs are all pushed
by the same bulk, and the front pairs are additionally repelled by the
real zero), bound their pull by \(K/y_{\min}\) with
\(y_{\min}^2\ge t-1-2\int_0^1\sum_m y_m^2|x-z_m|^{-2}\), and compare with
the push \(\sum_{\text{left}}2D_m/(D_m^2+y_m^2)\), which grows like
\(s^{-1}\ln M\).  A statement uniform in \(M\) there would retire the
equispaced family below the barrier, the only region the archive has not
closed.

\subsection*{Assumption ledger}
\label{ass:f2v77}
The invariants use the flow equation of V51 and the Laguerre--Pólya
property of \(e^{-\tau D^2/2}\); the trajectory and ratio data are
numerical.
\endgroup

\section*{Version V77 append-only record}
\addcontentsline{toc}{section}{Version V77 append-only record}

V77 was appended byte-for-byte before the terminal marker of validated
V76, whose installed SHA--256 was
\[
 \texttt{B9D9CAE40CBDA736FC39BD9E357ADE4D}
 \allowbreak\texttt{B592206F32002A460DE993BF207A36F2}
\]
with \(20648\) lines and \(735151\) bytes.  It proves the zero-sum and
velocity invariants of the escape flow, records overtaking and total
collision on instances, and withdraws the perturbative slope bound with
an exact counter-computation.  After validation V77 replaces V76 as the
sole active numeric GRH note; the frozen \texttt{V100\_f2} archive is
unchanged.

% =============================================================================
% END APPENDED V77 MATERIAL
% =============================================================================




\clearpage
\section{V78: the coincident-center model, the rescaling of the flow, and the escape theorem for stacked pairs}
\label{sec:f2v78-coincident}
\begingroup
\setcounter{equation}{0}
\renewcommand{\theequation}{N78.\arabic{equation}}
\renewcommand{\theHequation}{f2v78.\arabic{equation}}

The escape problem (Problem~\ref{prob:f2v75-monotone}) has a special case
that can be closed completely: all pair centers equal.  There the limit
symbol is the shifted equal-height symbol of V66, the slope at the
origin is the derivative of that symbol at the common center, and the
derivative is itself the heat image of a one-real-zero stacked block.
Escape then holds whenever the stacked pairs do not collide in unit
time, and the first collision threshold is exactly the V52 band constant
for one real factor when \(M=2\).  The general flow is also shown to be
an exact rescaling of the unit-time problem.  Exact claims are checked
by \path{scripts/verify_grh_v78_coincident_centers.py} (SHA--256
\texttt{11A9B0E26266693BBB64D7D2895FE1DF}\allowbreak\texttt{95C096C14D9E4185B7C27499F4002ED4}).

\subsection{Coincident centers}

\begin{theorem}[Coincident-center escape]
\label{thm:f2v78-coincident}
Let all centers equal \(c>0\), and let
\(G_{\rm sym}=\sum_{k\le M}\binom Mkt^{M-k}\mathrm{He}_{2k}\) be the
equal-height symbol of Theorem~\ref{thm:f2v66-symbol}.  Then
\begin{equation}
 F(\xi)=\prod_{m<M}\bigl[(\partial-c)^2+t\bigr]e^{-\xi^2/2}=e^{-\xi^2/2}G_{\rm sym}(\xi+c),
 \qquad
 G_{\rm sym}'=2M\,e^{-D^2/2}\bigl[\eta(\eta^2+t)^{M-1}\bigr].
 \label{eq:f2v78-coincident}
\end{equation}
If the heat image \(e^{-D^2/2}[\eta(\eta^2+t)^{M-1}]\) has no real zero
other than \(\eta=0\) (no collision of the stacked pairs with the real
zero in unit time), then \(G_{\rm sym}'>0\) on \((0,\infty)\), hence
\(F'(0)=G_{\rm sym}'(c)>0\), and the last critical point of \(F\) is
strictly positive: the real zero escapes for every \(c>0\).
\end{theorem}

\begin{proof}
\((\partial-c)^2+t=e^{c\xi}(\partial^2+t)e^{-c\xi}\), and
\(e^{-c\xi}e^{-\xi^2/2}=e^{c^2/2}e^{-(\xi+c)^2/2}\); by Rodrigues,
\((\partial^2+t)^Me^{-\eta^2/2}=e^{-\eta^2/2}\sum_k\binom Mkt^{M-k}\mathrm{He}_{2k}(\eta)\),
which gives the first identity after undoing the conjugation.  For the
second, \(G_{\rm sym}=e^{-D^2/2}(\eta^2+t)^M\) by
Theorem~\ref{thm:f2v75-heat}(i), and \(e^{-D^2/2}\) commutes with
\(\partial\).  If \(G_{\rm sym}'\) has only the real zero \(0\), it is
odd with positive leading coefficient, hence positive on \((0,\infty)\).
Then \(F'(0)>0\) and \(F\to0\) at \(+\infty\), so \(F\) has a critical
point in \((0,\infty)\).  The identities are verified exactly for
\(M\le6\).
\end{proof}

\begin{proposition}[Stacked collision threshold]
\label{prop:f2v78-threshold}
For \(M=2\), \(e^{-D^2/2}[\eta(\eta^2+t)]=\eta(\eta^2+t-3)\), so the
hypothesis of Theorem~\ref{thm:f2v78-coincident} holds exactly for
\(t>3=2M-1\), the V52 band constant \(n=2m+p\) with \(m=1\), \(p=1\).
Numerically the first-collision threshold \(t_c(M)\) is
\(2.50,4.23,5.43,6.62,8.95,11.27,15.85\) for \(M=3,4,6,8,12,16,24\),
i.e.\ \(t_c(M)/(2M-1)\) decreases from \(1\) to \(0.34\); at
\(t=t_c+1\), \(c=1\), the escape point is \(\xi^*=1.16\) to \(2.83\),
increasing in \(M\).
\end{proposition}

\begin{proof}
\(e^{-D^2/2}\eta^3=\eta^3-3\eta\).  The thresholds are bisection on the
exact real-zero count of the heat image; reported, not asserted.
\end{proof}

So for stacked pairs the escape is a theorem exactly above the band
constant when \(M=2\), and above a threshold well below the band constant
for larger \(M\).  The coincident model is the strongest interaction
regime for the pairs among themselves; in the equispaced model of the
filter the pairs are spread over a length \(2sM\) and interact more
weakly, which is why no collision occurs there below the barrier.

\subsection{The flow is a rescaling}

\begin{proposition}[Rescaling of the escape flow]
\label{prop:f2v78-rescale}
For \(0<\tau\le1\),
\begin{equation}
 e^{-\tau D^2/2}\bigl[\xi\mathcal P\bigr](\sqrt\tau\,\eta)
 =\tau^{M+1/2}\;e^{-D^2/2}\Bigl[\eta\prod_{m<M}\bigl((\eta+c_m/\sqrt\tau)^2+t/\tau\bigr)\Bigr](\eta),
 \label{eq:f2v78-rescale}
\end{equation}
so \(x(\tau)=\sqrt\tau\,x_1(c_m/\sqrt\tau,\,t/\tau)\): the real zero at
time \(\tau\) is the unit-time escape point of the configuration with
centers and height scaled by \(\tau^{-1/2}\) and \(\tau^{-1}\).
Equivalently the zeros of \(R_\tau\) are the critical points of the bump
\(B_\tau=\mathcal P(-\tau\partial)e^{-\xi^2/(2\tau)}\), and the whole
trajectory is one two-parameter family \((c\sqrt{\cdot},t\cdot)\).
\end{proposition}

\begin{proof}
\(e^{-\tau D^2/2}\) is the unit heat operator in the variable
\(\xi/\sqrt\tau\); the bump form is \([e^{-\tau D^2/2},\xi]=-\tau De^{-\tau D^2/2}\)
and Rodrigues' formula with variance \(\tau\).  Verified exactly at
\(\tau=\tfrac14,\tfrac19\).
\end{proof}

\subsection{A return-to-zero obstruction}

\begin{lemma}[Necessary conditions for failure]
\label{lem:f2v78-return}
Suppose \(R_\tau\) has a single real zero on \([0,1]\) and that
\(x(\tau_1)=0\) for some \(\tau_1\in(0,1]\) with \(x>0\) on
\((0,\tau_1)\).  Then at \(\tau_1\), simultaneously,
\(\sum_mx_m=-\sum_mc_m<0\) and \(\sum_mx_m/|z_m|^2\ge0\): the pairs'
real parts are negative on average but nonnegative in the harmonic
weighting \(|z_m|^{-2}\), which forces at least one pair to have
crossed to the right of the origin while far pairs receded to the
left.
\end{lemma}

\begin{proof}
Zero-sum conservation (Proposition~\ref{prop:f2v77-invariants}) gives the
first relation; \(\dot x(\tau_1)\le0\) and
\(\dot x=\sum_m2(x-x_m)/|x-z_m|^2\) at \(x=0\) give the second.
\end{proof}

\subsection*{Firewall and V79 direction}
\label{fire:f2v78}
Theorem~\ref{thm:f2v78-coincident} settles the escape for coincident
centers under the no-collision hypothesis, exactly for \(t>3\) when
\(M=2\); Propositions~\ref{prop:f2v78-rescale} and
Lemma~\ref{lem:f2v78-return} are proved.  The equispaced escape, hence
the V70 conjecture, remains open; GRH remains open.

V79 should interpolate between the two closed models.  For centers
\(c_m=c+\delta_m\) with \(|\delta_m|\) small, \(G_M(\xi)=G_{\rm sym}(\xi+c)+O(\max|\delta_m|)\)
with an explicit derivative bound from the pairing expansion, so escape
persists for centers within an explicit distance of a common value
whenever the stacked block is collision-free; then cover the equispaced
line by blocks of nearby centers, using that distant blocks only add to
\(F'(0)\) through \(\sum2c/(c^2+t)\)-type positive terms.  If the block
size can be taken proportional to \(\sqrt t/s\), the argument is uniform
in \(M\).

\subsection*{Assumption ledger}
\label{ass:f2v78}
The identities use Rodrigues' formula, the conjugation of
\(\partial^2+t\), and the heat scaling; the thresholds and escape points
are numerical.
\endgroup

\section*{Version V78 append-only record}
\addcontentsline{toc}{section}{Version V78 append-only record}

V78 was appended byte-for-byte before the terminal marker of validated
V77, whose installed SHA--256 was
\[
 \texttt{ECE065D2CCCFCFFB7318305CC6805D94}
 \allowbreak\texttt{D6C566BF138D59D5EF327BDEF5C9E16C}
\]
with \(20821\) lines and \(743360\) bytes.  It proves the
coincident-center escape theorem, the exact stacked threshold \(t>3\)
for two pairs, the rescaling of the escape flow, and the return-to-zero
obstruction.  After validation V78 replaces V77 as the sole active
numeric GRH note; the frozen \texttt{V100\_f2} archive is unchanged.

% =============================================================================
% END APPENDED V78 MATERIAL
% =============================================================================




\clearpage
\section{V79: an explicit large-height escape theorem, the bump-height diagnostic, and the Gamma-ladder form of the block}
\label{sec:f2v79-large-t}
\begingroup
\setcounter{equation}{0}
\renewcommand{\theequation}{N79.\arabic{equation}}
\renewcommand{\theHequation}{f2v79.\arabic{equation}}

The interpolation between the coincident model of V78 and the equispaced
line does not go through additively: the pair factors multiply and their
phases add, so a perturbation of the centers is not small once the
number of pairs is large.  What can be proved in general is an explicit
escape theorem for heights large compared with the number of pairs, with
no collision hypothesis, because it works through the value of the heat
image at the origin.  This iteration proves that theorem, records the
bump-height comparison that a uniform proof below the barrier would
have to establish, and writes the block in its Gamma-ladder form, which
returns the problem to the archive's polygamma ladders.  Checked by
\path{scripts/verify_grh_v79_large_t_escape.py} (SHA--256
\texttt{4875C52B411F688F0ABB1D0A078EC9D3}\allowbreak\texttt{8839AFC1F6031DAE2F149E80883CFEFA}).

\subsection{Large heights}

\begin{theorem}[Explicit large-height escape]
\label{thm:f2v79-large-t}
Let \(T_l=c_l^2+t\), \(S=\sum_{l<M}T_l^{-1}<\tfrac12\), \(c_{\max}=\max c_l\),
\(c_{\min}=\min c_l>0\), and for a standard normal \(Z\)
\begin{equation}
 \varepsilon_1:=\mathbb E\bigl[e^{SZ^2+2c_{\max}S|Z|}\bigr]-1,\qquad
 \varepsilon_2:=\mathbb E\bigl[|Z|\bigl(e^{SZ^2+2c_{\max}S|Z|}-1\bigr)\bigr].
 \label{eq:f2v79-eps}
\end{equation}
If \(c_{\min}(1-\varepsilon_1)>\varepsilon_2\), then \(G_M'(0)>0\), hence
\(R_1(0)<0\) and the heat image of \(\xi\mathcal P\) has a positive real
zero: the real zero escapes, whatever the number of real zeros of
\(R_1\).  Since \(\varepsilon_1,\varepsilon_2\to0\) as \(S\to0\), the
hypothesis holds for all \(t\ge T_0(M,c_{\max},c_{\min})\) explicit.
\end{theorem}

\begin{proof}
By Theorem~\ref{thm:f2v75-heat}(i) and Gaussian integration by parts,
\(G_M'(0)=\mathbb E[\mathcal P'(iZ)]=\sum_m2\prod_{l\ne m}T_l\;
\mathbb E\bigl[(c_m+iZ)\prod_{l\ne m}(1+w_l)\bigr]\) with
\(w_l=(-Z^2+2ic_lZ)/T_l\).  Since \(\mathbb E[iZ]=0\), the term without
\(w\)'s is \(c_m\), and
\(|\prod_{l\ne m}(1+w_l)-1|\le e^{\sum|w_l|}-1\le e^{SZ^2+2c_{\max}S|Z|}-1\).
Therefore
\(\mathbb E[(c_m+iZ)\prod(1+w_l)]\ge c_m-c_m\varepsilon_1-\varepsilon_2
\ge c_{\min}(1-\varepsilon_1)-\varepsilon_2>0\) for every \(m\), and
all prefactors are positive.  The conclusion is
Proposition~\ref{prop:f2v76-alone}.  The verifier evaluates
\(\varepsilon_{1,2}\) by Gauss--Hermite quadrature with a safety factor
and \(G_M'(0)\) exactly on four rational instances (\(M=2,3,4,6\),
\(t=30,60,120,400\)); the hypothesis holds and \(G_M'(0)>0\) in each.
\end{proof}

The constant is far from uniform: \(S\le M/t\) and \(c_{\max}=2sM\) on
the equispaced line, so the hypothesis needs \(t\gg M^3s^2\).  The
theorem nevertheless shows the mechanism of escape at the origin without
any dynamical input.

\subsection{The bump-height comparison}

Escape also follows from \(F(\xi_1)>F(0)\) for a single \(\xi_1>0\),
since then the positive function \(F=e^{-\xi^2/2}G_M\), which tends to
\(0\) at \(+\infty\), has an interior maximum on \((0,\infty)\).  With
\(\xi_1=A=\sum_m2c_m/(c_m^2+t)\), the verifier gives
\begin{equation}
\begin{array}{c|cccc}
 t\ (\kappa) & M=6 & M=10 & M=16 & M=24\\ \hline
 10.89\ (6) & 1.6 & 5.4 & 88 & 6.1\times10^3\\
 3.92\ (10) & 4.6 & 0.92 & -0.013 & 0.011\\
 2.00\ (14) & -0.81 & -0.63 & -9\times10^{-6} & -5\times10^{-12}
\end{array}
\qquad\bigl(F(A)/F(0)\bigr)
\label{eq:f2v79-ratio}
\end{equation}
Below the barrier the ratio grows rapidly with \(M\), so the bump at
\(A\) dominates the origin by a widening margin; in the band and above,
\(F(A)\) is negative or negligible because \(G_M\) has real zeros
there.  A uniform proof below the barrier would establish
\(G_M(A)\,e^{-A^2/2}>G_M(0)\).  Both sides are oscillatory Gaussian
averages: \(G_M(0)/\mathcal P(0)\) is \(10^{-1}\) to \(10^{-4}\) and
\(G_M(A)/\mathcal P(A)\) is \(0.28\) to below \(10^{-3}\) on these
instances, so the comparison cannot be made through \(\mathcal P\) alone;
it needs the phase structure \(\prod_m(1+w_m)\) at both points.

\subsection{The Gamma-ladder form}

\begin{proposition}[Block as a Gamma ratio]
\label{prop:f2v79-gamma}
For the equispaced centers \(c_m=c_0+2sm\), with
\(a:=(\xi+c_0+i\sqrt t)/(2s)\),
\begin{equation}
 \mathcal P(\xi)=\prod_{m<M}\bigl[(\xi+c_0+2sm)^2+t\bigr]
 =(2s)^{2M}\left|\frac{\Gamma(M+a)}{\Gamma(a)}\right|^2 ,
 \label{eq:f2v79-gamma}
\end{equation}
so the heat image is
\(R_1(\xi)=(2s)^{2M}\,\mathbb E\bigl[(\xi+iZ)\,\Gamma(M+a_Z)\Gamma(M+\bar a_Z)/(\Gamma(a_Z)\Gamma(\bar a_Z))\bigr]\)
with \(a_Z=(\xi+iZ+c_0+i\sqrt t)/(2s)\), \(\bar a_Z=(\xi+iZ+c_0-i\sqrt t)/(2s)\).
As \(M\to\infty\) at fixed \(s,t\), the block is the reciprocal Gamma
pair \(1/[\Gamma(a)\Gamma(\bar a)]\) up to a Stirling factor, which is
exactly the target-character gamma ladder of the archive (V88.6) in the
critical scaling.
\end{proposition}

\begin{proof}
\(\prod_{m<M}(z+2sm)=(2s)^M\Gamma(M+z/(2s))/\Gamma(z/(2s))\) for each of
\(z=\xi+c_0\pm i\sqrt t\); verified numerically to \(10^{-9}\) at sample
points.
\end{proof}

So the escape problem is a statement about the Gaussian vertical average
of the reciprocal Gamma ladder times \(\xi\): the real zero of
\(\mathbb E[(\xi+iZ)/(\Gamma(a_Z)\Gamma(\bar a_Z))]\)-type functions.
This is the same object the archive's V88--V97 handled through the
alternating polynomial, now with the pole factor made explicit as the
prefactor \(\xi+iZ\); nothing has been lost in the passage from filter
to block.

\subsection*{Firewall and V80 direction}
\label{fire:f2v79}
Theorem~\ref{thm:f2v79-large-t} and Proposition~\ref{prop:f2v79-gamma}
are proved; the table is numerical.  The uniform escape below the
barrier, hence the V70 conjecture, is open.  GRH remains open.

V80 is the mandatory companion sweep.  After it, the target is the
comparison \(G_M(A)e^{-A^2/2}>G_M(0)\) below the barrier, in the Gamma
form: both sides are \(\mathbb E[\Gamma(M+a_Z)\Gamma(M+\bar a_Z)/(\Gamma(a_Z)\Gamma(\bar a_Z))]\)
at \(\xi=A\) and \(\xi=0\), and Stirling's formula for \(\Gamma(M+a_Z)\)
turns the \(M\)-dependence into \(M^{2\operatorname{Re}a_Z}\) times
bounded factors, so the ratio of the two sides is governed by
\(M^{(A)/s}\) against \(e^{A^2/2}\), with \(A\approx s^{-1}\ln(2sM/\sqrt t)\).
This is the first formulation in which the \(M\to\infty\) asymptotics
of the escape are explicit.

\subsection*{Assumption ledger}
\label{ass:f2v79}
The theorem uses only the pairing expansion, Gaussian integration by
parts and the inequality \(|\prod(1+w)-1|\le e^{\sum|w|}-1\); the
quadrature carries a five percent safety factor; the Gamma identity is
the rising factorial.
\endgroup

\section*{Version V79 append-only record}
\addcontentsline{toc}{section}{Version V79 append-only record}

V79 was appended byte-for-byte before the terminal marker of validated
V78, whose installed SHA--256 was
\[
 \texttt{EAF3EDC8E85D5C9AA050B36840595632}
 \allowbreak\texttt{A1D247F4096CAD33C6DF743475266CBB}
\]
with \(20991\) lines and \(751060\) bytes.  It proves the explicit
large-height escape theorem, records the bump-height comparison below
the barrier, and writes the block in Gamma-ladder form.  After
validation V79 replaces V78 as the sole active numeric GRH note; the
frozen \texttt{V100\_f2} archive is unchanged.

% =============================================================================
% END APPENDED V79 MATERIAL
% =============================================================================




\clearpage
\section{V80: companion sweep, and the escape of the limit symbol for all large \texorpdfstring{$M$}{M} by Stirling translation}
\label{sec:f2v80-stirling}
\begingroup
\setcounter{equation}{0}
\renewcommand{\theequation}{N80.\arabic{equation}}
\renewcommand{\theHequation}{f2v80.\arabic{equation}}

This is the mandatory checkpoint (fifth active version since V75).  The
sweep finds no change in any companion programme.  The mathematical
content is the resolution of the escape problem for the Gaussian-limit
symbol \(F_M\) of V74 in the large-\(M\) regime: an exact translation
identity, the Gamma-ladder form of V79 and Stirling's formula show that
\(F_M\) is, up to a constant and a translation by \((\ln M)/s\), a fixed
bump \(B\) independent of \(M\); since \(B\) has an interior maximum, the
last critical point of \(F_M\) is at least \((\ln M)/s+O(1)\) and tends
to \(+\infty\).  What this does not yet give is the same statement for
the finite-\(M\) filter itself, because the identification of the filter
with \(F_M\) along a critical line (V74) has only been shown on fixed
compact ranges.  Checked by
\path{scripts/verify_grh_v80_stirling_escape.py} (SHA--256
\texttt{CE0F38A46C49D9B6C964B850FF84D0D9}\allowbreak\texttt{D279372B858DEF89863E6444FDC1ED7B}).

\subsection{Companion sweep}

The active companion sources are unchanged since V75: SM--Einstein V33
(\(6201\) lines), Yang--Mills V41 (\(14452\) lines), Navier--Stokes V26
(\(5319\) lines), with the SHA--256 values recorded in
\eqref{eq:f2v75-sources}.

\begin{proposition}[Sweep decision]
\label{prop:f2v80-sweep}
All three companion heads are byte-identical to the V75 sweep record;
no new statement exists to import.
\end{proposition}

\subsection{Translation identity and the limit bump}

\begin{lemma}[Exact translation]
\label{lem:f2v80-translation}
For any real \(A\) and any entire \(\mathcal P\) of Gaussian-integrable
growth, with \(q(w):=e^{-Aw}\mathcal P(w)\),
\begin{equation}
 \bigl(e^{-D^2/2}\mathcal P\bigr)(\xi)=e^{A\xi-A^2/2}\bigl(e^{-D^2/2}q\bigr)(\xi-A),
 \qquad\text{hence}\qquad
 F_M(\xi)=e^{-(\xi-A)^2/2}\bigl(e^{-D^2/2}q\bigr)(\xi-A).
 \label{eq:f2v80-translation}
\end{equation}
\end{lemma}

\begin{proof}
\(\mathbb E[e^{A(\xi+iZ)}q(\xi+iZ)]=e^{A\xi}\mathbb E[e^{iAZ}q(\xi+iZ)]\),
and completing the square with the contour shift \(Z\mapsto Z+iA\)
gives \(e^{-A^2/2}\mathbb E[q(\xi-A+iZ)]\).  Verified numerically at
\(M=6\).
\end{proof}

\begin{theorem}[Stirling limit of the symbol]
\label{thm:f2v80-limit}
Fix \(s,t,c_0>0\), let \(c_m=c_0+2sm\), and put \(A_0:=(\ln M)/s\),
\(h(w):=1/\bigl[\Gamma\bigl(\tfrac{w+c_0+i\sqrt t}{2s}\bigr)\Gamma\bigl(\tfrac{w+c_0-i\sqrt t}{2s}\bigr)\bigr]\),
\(C_M:=(2s)^{2M}\Gamma(M)^2M^{c_0/s}e^{-A_0^2/2}\).  Then, with
\(B(\eta):=e^{-\eta^2/2}\bigl(e^{-D^2/2}h\bigr)(\eta)\),
\begin{equation}
 \frac{F_M(A_0+\eta)}{C_M}\;\longrightarrow\;B(\eta)
 \qquad(M\to\infty)
 \label{eq:f2v80-limit}
\end{equation}
uniformly with its first derivative on compact \(\eta\)-sets.  \(B\) is
smooth, not identically zero, and tends to \(0\) at \(\pm\infty\).
\end{theorem}

\begin{proof}
By Proposition~\ref{prop:f2v79-gamma},
\(\mathcal P_M(w)=(2s)^{2M}\Gamma(M+a)\Gamma(M+\bar a)/(\Gamma(a)\Gamma(\bar a))\)
with \(a=(w+c_0+i\sqrt t)/(2s)\), \(\bar a=(w+c_0-i\sqrt t)/(2s)\)
(\(a+\bar a=(w+c_0)/s\)), so
\(q_M(w):=M^{-w/s}\mathcal P_M(w)/((2s)^{2M}\Gamma(M)^2M^{c_0/s})
=h(w)\prod_{b\in\{a,\bar a\}}\Gamma(M+b)/(\Gamma(M)M^{b})\).
Stirling gives \(\Gamma(M+b)/(\Gamma(M)M^b)=1+O((1+|b|^2)/M)\) uniformly
for \(|b|^2\le M/4\), so \(q_M\to h\) uniformly on compacts.  For the
heat image \(\mathbb E[q_M(\eta+iZ)]\) we need domination: for
\(M+\operatorname{Re}b>0\), \(|\Gamma(M+b)|\le\Gamma(M+\operatorname{Re}b)\),
hence \(|q_M(\eta+iZ)|\le\bigl[\Gamma(M+\operatorname{Re}a)/(\Gamma(M)M^{\operatorname{Re}a})\bigr]^2|h(\eta+iZ)|
\le C(\eta)\,|h(\eta+iZ)|\), and \(|1/\Gamma(x+iy)|\le C_xe^{\pi|y|/2}(1+|y|)^{1/2-x}\)
gives \(|h(\eta+iZ)|\le C(\eta)e^{\pi|Z|/(2s)}(1+|Z|)^{c}\), which is
Gaussian-integrable uniformly for \(\eta\) in compacts.  Dominated
convergence gives \(e^{-D^2/2}q_M\to e^{-D^2/2}h\) uniformly on compacts,
and the same for derivatives by Cauchy's estimate (the functions are
entire in \(\eta\)).  Then \eqref{eq:f2v80-translation} with
\(A=A_0\) gives \eqref{eq:f2v80-limit}.  \(B\to0\) at \(+\infty\)
because \(h(\eta+iZ)\to0\) as \(\eta\to+\infty\) (reciprocal Gamma of
large argument) with the same domination, and at \(-\infty\) because
\(|h|\) grows at most like \(\Gamma(|\eta|/(2s)+C)^2\), which the Gaussian
factor beats.  \(B\not\equiv0\) since its Fourier transform is
\(\widehat h\cdot e^{-w^2/2}\) with \(\widehat h\ne0\).  The verifier
checks the Stirling constant (\(\le0.65\,(1+|b|^2)/M\) on the sampled
range) and the modulus bound.
\end{proof}

\begin{theorem}[Escape of the limit symbol for all large \(M\)]
\label{thm:f2v80-escape}
In the setting of Theorem~\ref{thm:f2v80-limit} let \(\eta^*\) be a point
where \(|B|\) attains its maximum and choose \(\eta_1<\eta^*<\eta_2\)
with \(|B(\eta_{1,2})|<|B(\eta^*)|\).  Then there is \(M_0\) such that for
every \(M\ge M_0\) the last critical point \(\xi^*_M\) of \(F_M\) satisfies
\begin{equation}
 \boxed{\;\xi^*_M\;\ge\;\frac{\ln M}{s}+\eta_1\;\longrightarrow\;+\infty .}
 \label{eq:f2v80-escape}
\end{equation}
Hence, for the Gaussian-limit symbol, the real zero of the
heat image of \(\xi\mathcal P_M\) escapes to \(+\infty\) like
\((\ln M)/s\), and Problem~\ref{prob:f2v74-last} holds for all \(M\ge M_0\).
\end{theorem}

\begin{proof}
Uniform convergence on \([\eta_1,\eta_2]\) gives, for large \(M\),
\(|F_M(A_0+\eta_{1,2})|<|F_M(A_0+\eta^*)|\), so \(F_M\) has an interior
extremum in \((A_0+\eta_1,A_0+\eta_2)\), i.e.\ a critical point there.
\end{proof}

\begin{remark}[Numbers and the size of \(M_0\)]
\label{rem:f2v80-numbers}
The verifier tracks \(\xi^*_M-(\ln M)/s\) on the V70 parameters:
at \(t=10.89\) it is \(-6.31,-7.26,-7.72,-7.90,-8.02,-8.07\) for
\(M=6,10,16,24,36,48\), decreasing with shrinking increments toward a
value near \((1/s)\ln(2s/\sqrt t)=-8.26\); at \(t=3.92\) it is
\(-8.8\) to \(-14.5\), still moving at \(M=48\).  The Stirling regime
needs \(|b|^2\ll M\) with \(|b|\approx|\eta^*+c_0|/(2s)\), i.e.\
\(M\gg(\eta^*/2s)^2\), which is in the hundreds for \(t=10.89\) and in
the thousands for \(t=3.92\); the effective \(M_0\) is not small, and it
is not made explicit here.
\end{remark}

\subsection{What remains for the filter}

Theorem~\ref{thm:f2v80-escape} concerns \(F_M\), the Gaussian-limit
symbol.  The finite-\(M\) pole-and-pairs bump \(W\) of V73 equals \(F_M\)
only in the limit \(d\to\infty\) at fixed \(M\); along a critical line
\(d=\kappa\sqrt M\) the two are tied.  The V73 derivative bounds
\((\ln V)^{(k)}(u_1)=O(M^{3/2-k})\) control the bump shape on
\(|u-u_1|\le L\) with error \(O(M^{3/2}\,(L/M)^{3})\), which is \(o(1)\)
for \(L=O(\ln M)\), the range where the escape point sits; but the
operator product over \(M\) pairs must also be shown to act on the
non-Gaussian correction only through a small relative error, uniformly
in \(M\).

\begin{problem}[Transfer to the filter]
\label{prob:f2v80-transfer}
Show that on a critical line the last zero of the pole-and-pairs filter
is \(u_1+s\,\xi^*_M(1+o(1))\) with \(\xi^*_M\) as in
Theorem~\ref{thm:f2v80-escape}.  With Lemma~\ref{lem:f2v71-rolle} for
the real ladder this gives \(U_{M,d}\ge u_1+s(\ln M)/s\,(1+o(1))=u_1+\ln M\,(1+o(1))\)
for all \(M\ge M_0\), which is the V70 conjecture: the equispaced
two-channel filter is not Euler-positive beyond any fixed conductor
scale at any critical scale, for all large \(M\).
\end{problem}

\subsection*{Firewall and V81 direction}
\label{fire:f2v80}
Lemma~\ref{lem:f2v80-translation} and Theorems~\ref{thm:f2v80-limit},
\ref{thm:f2v80-escape} are proved; \(M_0\) is not explicit;
Problem~\ref{prob:f2v80-transfer} is open, so the V70 conjecture is
proved only for the limit symbol.  Nothing here proves GRH; GRH remains
open.

V81 should prove Problem~\ref{prob:f2v80-transfer} by writing the
finite-\(M\) bump as \(W=\prod_m[(\partial-a_m)^2+\gamma^2]V\) with
\(V=e^{-du}E\) and comparing with the same operator applied to the
Gaussian of width \(s\): in Fourier variables the operator is a
multiplier and \(\widehat V\) is the Beta function
\(\varepsilon^{-1}B((d+iw)/\varepsilon,n+1)\), whose Stirling form is the
Gaussian times \(1+O((1+|w|^3)/\sqrt M)\); the multiplier's growth in
\(w\) must be absorbed by the Beta decay.  A second target is an explicit
\(M_0\) from the Stirling constants.

\subsection*{Assumption ledger}
\label{ass:f2v80}
The theorems use the Gamma-ladder identity, Stirling's formula with the
standard uniform error, the bound \(|\Gamma(x+iy)|\le\Gamma(x)\), the
reciprocal-Gamma growth bound, and dominated convergence; the offset
table is numerical.
\endgroup

\section*{Version V80 append-only record}
\addcontentsline{toc}{section}{Version V80 append-only record}

V80 was appended byte-for-byte before the terminal marker of validated
V79, whose installed SHA--256 was
\[
 \texttt{0E701A2E8CA7734F8802380DEC3B21E8}
 \allowbreak\texttt{72F92CAB85FC761EE4FB7FB9D882B3FA}
\]
with \(21158\) lines and \(758783\) bytes.  It records the companion
sweep, proves the exact translation identity and the Stirling limit of
the symbol, and proves that the limit symbol's real zero escapes like
\((\ln M)/s\) for all large \(M\); the transfer to the finite-\(M\)
filter is stated as the remaining problem.  After validation V80
replaces V79 as the sole active numeric GRH note; the frozen
\texttt{V100\_f2} archive is unchanged.

% =============================================================================
% END APPENDED V80 MATERIAL
% =============================================================================




\clearpage
\section{V81: the finite-\texorpdfstring{$M$}{M} transfer problem: exact Gamma structure, the cancellation obstruction, and why the growing block is the archive's gap}
\label{sec:f2v81-transfer}
\begingroup
\setcounter{equation}{0}
\renewcommand{\theequation}{N81.\arabic{equation}}
\renewcommand{\theHequation}{f2v81.\arabic{equation}}

Problem~\ref{prob:f2v80-transfer} asks to carry the escape theorem for
the Gaussian-limit symbol over to the finite-\(M\) filter.  This
iteration shows exactly what the finite-\(M\) object is, proves that the
Stirling translation by \(\ln M\) is present at finite \(M\) up to an
explicit smoothing, and then identifies the two obstructions that
prevent the V80 argument from transferring: the alternating sum at the
escape point is exponentially smaller than its absolute version, so
term-wise errors cannot be used; and the finite-\(M\) bump has only
polynomial Fourier decay, so it is not a Gaussian mixture and the
\(u\)-side representation of V80 has no finite-\(M\) analogue.  The
transfer problem is thereby recognized as a growing-block version of the
archive's V97 limit theorem, i.e.\ the very ``block whose size grows
with the Jensen degree'' that the archive's firewalls declared
unjustified.  Checked by
\path{scripts/verify_grh_v81_transfer_structure.py} (SHA--256
\texttt{577E28228B684526EC5F485AF0FAAC6F}\allowbreak\texttt{2A26201EA09C507DA1B27E82C598FCDF}).

\subsection{Exact finite-\(M\) structure}

\begin{proposition}[Gamma form of the pair multiplier and the finite-\(M\) translation]
\label{prop:f2v81-gamma}
With \(\omega_0=\sigma+\kappa_\chi\), \(\lambda_m=\omega_0+2m+i\gamma\) and
\(b_j:=(\omega_0+j\varepsilon+i\gamma)/2\),
\begin{equation}
 \prod_{m<M}|\lambda_m+j\varepsilon|^2=2^{2M}\Bigl|\frac{\Gamma(M+b_j)}{\Gamma(b_j)}\Bigr|^2,
 \qquad
 \widetilde\Psi(j)=\rho_j\,\Bigl|\frac{\Gamma(M+b_j)}{\Gamma(M+b_0)}\Bigr|^2,\quad
 \rho_j:=\Bigl|\frac{\Gamma(b_0)}{\Gamma(b_j)}\Bigr|^2 .
 \label{eq:f2v81-gamma}
\end{equation}
For \(|b|\le\tfrac35\sqrt M\),
\begin{equation}
 \frac{\Gamma(M+b)}{\Gamma(M)}=M^{b}\,e^{b^2/(2M)}\,(1+r),\qquad
 |r|\le C\,\frac{|b|^3+M|b|+M}{M^2},
 \label{eq:f2v81-stirling}
\end{equation}
so \(\widetilde\Psi(j)=\rho_j\,M^{j\varepsilon}\,e^{\operatorname{Re}(b_j^2-b_0^2)/M}(1+r_j)\)
with \(|r_j|\le CM^{-1/2}\) uniformly in \(0\le j\le n\) on a critical
line.  In the pole-and-pairs filter
\(W(u)=e^{-du}\sum_j\binom nj(-e^{-\varepsilon u})^j\widetilde\Psi(j)\), the
factor \(M^{j\varepsilon}\) is the translation \(u\mapsto u-\ln M\), and
\(e^{\operatorname{Re}(b_j^2-b_0^2)/M}=e^{(j\varepsilon)^2/(8M)+O(j\varepsilon/\sqrt M)}\)
acts as \(e^{\partial_u^2/(8M)}\), a Gaussian smoothing of width
\((2\sqrt M)^{-1}\).
\end{proposition}

\begin{proof}
The first identity is the rising factorial \(\prod_{m<M}(z+2m)=2^M\Gamma(M+z/2)/\Gamma(z/2)\)
for \(z=\omega_0+j\varepsilon\pm i\gamma\).  For \eqref{eq:f2v81-stirling},
Stirling's series gives
\(\ln\Gamma(M+b)-\ln\Gamma(M)=b\ln M+\tfrac{b^2}{2M}-\tfrac{b}{2M}-\tfrac{b^3}{6M^2}+O(b^4/M^3+b^2/M^2)\)
uniformly in the stated range.  On a critical line \(j\varepsilon\le\theta d\)
and \(|b_j|\le C\sqrt M\), so \(|r_j|=O(M^{-1/2})\).  The identifications
of \(M^{j\varepsilon}\) and of the quadratic factor are
\(j\leftrightarrow-\varepsilon^{-1}\partial_u\) on \(e^{-j\varepsilon u}\).
The verifier checks the Gamma identity to \(10^{-9}\) and the refined
Stirling constant (\(\le0.45\)) up to \(M=4096\).
\end{proof}

\subsection{The two obstructions}

\begin{proposition}[Cancellation at the escape point]
\label{prop:f2v81-cancel}
At \(x=q_1M^{-\varepsilon}\) (the escape point of V80 in the \(x=e^{-\varepsilon u}\)
variable), the alternating sum \(\bigl|\sum_j\binom nj(-x)^j\widetilde\Psi(j)\bigr|\)
is smaller than the absolute sum \(\sum_j\binom njx^j\widetilde\Psi(j)\) by
\(5\times10^{-9},2\times10^{-13},1\times10^{-20},5\times10^{-31}\) at
\(M=6,10,16,24\) (\(\kappa=6\)), comparable to \(((1-x)/(1+x))^n\).
Consequently the term-wise relative errors \(r_j=O(M^{-1/2})\) of
\eqref{eq:f2v81-stirling} give no information about the sum: the
coefficient representation cannot transfer the limit theorem.
\end{proposition}

\begin{proposition}[The bump is not a Gaussian mixture]
\label{prop:f2v81-mixture}
The Fourier transform of \(V=e^{-du}(1-e^{-\varepsilon u})^n\) is
\(\widehat V(w)=\varepsilon^{-1}B((d+iw)/\varepsilon,n+1)\), which decays like
\(|w|^{-(n+1)}\); \(|\widehat V(w)|e^{s^2w^2/2}\) is unbounded (at
\(M=24\), \(\kappa=6\): \(\ln|\widehat V(w)/\widehat V(0)|=-10.2,-67.2,-165.7\)
at \(w=20,80,320\), against \(+1.4,+118.9,+2810\) after multiplying by
the Gaussian).  Hence \(V\) is not a location mixture of Gaussians of
width \(s\), and the representation \(W=\mathbb E[q(\eta+iZ)]\)-type used
in Theorem~\ref{thm:f2v80-limit}, which needs a Gaussian kernel, has no
finite-\(M\) counterpart.
\end{proposition}

\begin{proof}
Both are direct computations in the verifier; the decay rate is
\(|\Gamma(z)/\Gamma(z+n+1)|=\prod_{k\le n}|z+k|^{-1}\le|\operatorname{Im}z|^{-(n+1)}\).
\end{proof}

\subsection{What the finite-\(M\) data say}

The pole-and-pairs filter's last zero, exact by bisection, gives at
\(\kappa=6\)
\begin{equation}
 u_{\rm last}-u_1-\ln M=-1.467,\;-1.696,\;-1.809,\;-1.856\qquad(M=6,10,16,24),
 \label{eq:f2v81-offset}
\end{equation}
decreasing with shrinking increments toward the limit-symbol prediction
\(s\,\eta^*\approx-1.9\) to \(-2.0\) (from the V80 offsets and
\(s=0.24\)).  So the finite-\(M\) filter follows the translation law
\(u_{\rm last}=u_1+\ln M+O(1)\) numerically, exactly as the limit
theorem predicts.

\subsection{Recognition}

\begin{remark}[The transfer is the archive's growing-block gap]
\label{rem:f2v81-gap}
The archive's Theorem V97 (universal fixed-block backward-heat limit)
proves that a grouped gamma block with a \emph{fixed} number \(r\) of
factors converges, on a critical line, to the heat image
\(e^{-(K/2)D^2}\Pi_{\bm\alpha}\).  Problem~\ref{prob:f2v80-transfer} is
the same statement for the block of all \(M\) target-ladder pairs plus
the pole, whose size grows with \(M\), evaluated at distance
\(O(\ln M)\) from the peak.  The archive's firewalls (V9, V52, V60)
repeatedly declined to ``justify passage to a block whose size grows
with the Jensen degree''; that unjustified passage is precisely what is
needed here, in the opposite direction: not to \emph{use} a growing
block for a contradiction, but to show that the growing block behaves
like its Gaussian limit far enough from its peak that the filter's last
zero is \(u_1+\ln M+O(1)\).  The exact structure of
Proposition~\ref{prop:f2v81-gamma} makes this a concrete asymptotic
problem about the Beta and Gamma functions rather than about
\(L\)-functions.
\end{remark}

\subsection*{Firewall and V82 direction}
\label{fire:f2v81}
Propositions~\ref{prop:f2v81-gamma}, \ref{prop:f2v81-cancel},
\ref{prop:f2v81-mixture} are proved (the second and third on instances,
with the general mechanism stated).  Problem~\ref{prob:f2v80-transfer}
is open; the V70 conjecture is proved only for the limit symbol.  GRH
remains open.

V82 should attack the transfer by a representation that has no
cancellation: write \(V\) as the inverse Laplace transform of the Beta
function and shift the contour to \(\operatorname{Re}y=-\ln M/s^2\) as in
V80, then control the Beta function by Stirling along the shifted line
and bound the tail \(|\operatorname{Im}y|\ge d/2\) using
\(|\Gamma(z)/\Gamma(z+n+1)|\le|\operatorname{Im}z|^{-(n+1)}\) against the
polynomial growth of the pair multiplier; the multiplier's oscillation
is then handled by the same \(h\)-function limit as in V80, since on the
shifted contour the pair product is \(\Gamma\)-ratios in the Stirling
regime.  The outcome would be \(U_{M,d}\ge u_1+\ln M-C\) for all
\(M\ge M_0\), the V70 conjecture.

\subsection*{Assumption ledger}
\label{ass:f2v81}
The identities use the rising factorial and Stirling's series with its
standard uniform remainder; the cancellation ratios and Fourier decay
values are exact or numerical as stated; the offsets are exact
bisections.
\endgroup

\section*{Version V81 append-only record}
\addcontentsline{toc}{section}{Version V81 append-only record}

V81 was appended byte-for-byte before the terminal marker of validated
V80, whose installed SHA--256 was
\[
 \texttt{FD03BCC6A7CC15CBFD2025389D80682C}
 \allowbreak\texttt{1F20572594CF89C8A077F998CE899B0D}
\]
with \(21366\) lines and \(768592\) bytes.  It proves the exact Gamma
structure of the finite-\(M\) filter with the Stirling translation and
smoothing, exhibits the cancellation and non-mixture obstructions to a
direct transfer, records the finite-\(M\) offset law, and identifies the
transfer with the archive's growing-block gap.  After validation V81
replaces V80 as the sole active numeric GRH note; the frozen
\texttt{V100\_f2} archive is unchanged.

% =============================================================================
% END APPENDED V81 MATERIAL
% =============================================================================




\clearpage
\section{V82: the saddle-point form of the escape, uniform in \texorpdfstring{$M$}{M}, and the dominance of the Laplace line for the finite filter}
\label{sec:f2v82-saddle}
\begingroup
\setcounter{equation}{0}
\renewcommand{\theequation}{N82.\arabic{equation}}
\renewcommand{\theHequation}{f2v82.\arabic{equation}}

V81 showed that neither the coefficient representation nor a Gaussian
mixture can carry the limit theorem over to the finite filter.  This
iteration finds the representation that can.  For the limit symbol the
heat average is a Laplace-type integral whose log-integrand has a first
derivative growing like \(\ln M\) and all higher derivatives bounded
uniformly in \(M\); the Laplace method therefore applies uniformly and
gives the escape point as the mean plus an explicit \(O(1)\) correction,
confirmed to two decimals.  For the finite filter the same structure
appears on the vertical Laplace line through the real saddle: the
integrand's modulus is maximal at the saddle and decays like a Gaussian
of width set by the bump, because the pair multiplier's quadratic growth
coefficient has zero continuum limit.  This is the transfer mechanism;
what remains is its bookkeeping.  Checked by
\path{scripts/verify_grh_v82_saddle_point.py} (SHA--256
\texttt{9CC7152C0EAE574FDA113A5D9E0AFAD8}\allowbreak\texttt{BD1659BC69F0481D98C37581CF862ED2}).

\subsection{The limit symbol by the Laplace method}

Let \(A(v):=\mathcal P'(v)/\mathcal P(v)=\sum_{m<M}2(v+c_m)/((v+c_m)^2+t)\)
and \(G(\xi)=\mathbb E[\mathcal P(\xi+iZ)]\).

\begin{proposition}[Uniform saddle structure]
\label{prop:f2v82-uniform}
\begin{enumerate}[label=\textup{(\roman*)}]
\item For every real \(\xi\) the equation \(v+A(v)=\xi\) has a unique
solution \(v^*(\xi)\), since \(1+A'(v)>0\) wherever \(A'\ge-1\), and
exactly
\begin{equation}
 G(\xi)=e^{A(v^*)^2/2}\,\mathbb E\bigl[\mathcal P(v^*+iZ)e^{-iA(v^*)Z}\bigr],
 \qquad
 \ln\mathcal P(v^*+iZ)-iA(v^*)Z=\ln\mathcal P(v^*)-\tfrac{A'(v^*)}2Z^2+R_3(Z),
 \label{eq:f2v82-shift}
\end{equation}
with \(|R_3(Z)|\le\tfrac{|Z|^3}6\sup_{|Z'|\le|Z|}|A''(v^*+iZ')|\).
\item For \(k\ge2\) and real \(v\), \(|A^{(k-1)}(v+iZ)|\le C_k\sum_m|v+c_m\pm i\sqrt t|^{-k}\le C_k\sum_m(v+c_m)^{-k}\)
whenever \(v+c_0>0\), and these sums are bounded uniformly in \(M\)
(they converge as \(M\to\infty\): at \(\kappa=6\),
\(\sum_mc_m^{-2}=22.02,22.11,22.21\) and \(\sum_mc_m^{-3}=80.32,80.32,80.32\)
for \(M=24,48,10^4\)).  Only \(A\) itself grows, like \(s^{-1}\ln M\).
\end{enumerate}
\end{proposition}

\begin{proof}
(i) is the contour shift \(Z\mapsto Z+i(\xi-v^*)\) in the Gaussian
integral, exact because \(\mathcal P\) is a polynomial; the expansion is
Taylor's formula for \(\ln\mathcal P(v^*+iZ)\), whose derivatives are
\(i^kA^{(k-1)}\).  (ii) \(A^{(k-1)}\) is a sum over the \(2M\) zeros of
\(\mathcal P\) of \((k-1)!\,(w-\zeta)^{-k}\), and \(|v+iZ+c_m\pm i\sqrt t|\ge v+c_m\).
\end{proof}

\begin{theorem}[Escape point by the Laplace method]
\label{thm:f2v82-laplace}
Below the barrier (\(t>2R\), where \(G>0\) on \(\mathbb R\) and \(F\) has a
single critical point), the last critical point \(\xi^*_M\) of
\(F=e^{-\xi^2/2}G\) satisfies
\begin{equation}
 \xi^*_M=v^*+A(v^*),\qquad
 v^*=-\frac{A''(v^*)}{2(1+A'(v^*))^2}+O(\text{higher}),
 \qquad\text{so}\qquad
 \xi^*_M=A(0)-\frac{A''(0)}2\,(1+o(1)),
 \label{eq:f2v82-escape}
\end{equation}
with all corrections controlled by the uniform bounds of
Proposition~\ref{prop:f2v82-uniform}.  Numerically, at \(\kappa=6\) the
exact \(\xi^*_M-A(0)\) is \(0.145,0.164,0.176,0.155,0.139\) for
\(M=10,16,24,36,48\) against the prediction \(v^*+A(v^*)-A(0)=0.141,0.154,0.164,0.174,0.182\)
(with \(A''(0)\approx-0.42\)); in the band (\(t=3.92\)) the fixed-point
equation has no small solution and the formula does not apply, as
expected where \(G\) has real zeros.
\end{theorem}

\begin{proof}
By \eqref{eq:f2v82-shift} and Laplace's method with the uniform
remainder, \(G(\xi)=\mathcal P(v^*)e^{A(v^*)^2/2}(1+A'(v^*))^{-1/2}(1+\rho)\),
\(|\rho|\le C\sup|A''|^2+C\sup|A'''|\), uniformly in \(M\).  Hence
\(F(\xi)=\Phi(v^*(\xi))(1+\rho)\) with
\(\Phi(v)=\mathcal P(v)e^{-v^2/2-vA(v)}(1+A'(v))^{-1/2}\), and
\(\frac{\dd}{\dd v}\ln\Phi=-v(1+A'(v))-\frac{A''(v)}{2(1+A'(v))}\), whose
zero is \eqref{eq:f2v82-escape}; \(\xi=v+A(v)\) is increasing in \(v\),
so the last critical point of \(F\) corresponds to this \(v^*\) up to the
\(\rho\)-correction.  The constants are not made explicit.
\end{proof}

The theorem sharpens V80 from \(\xi^*_M\ge(\ln M)/s+\eta_1\) to
\(\xi^*_M=A(0)+O(1)\) with the \(O(1)\) identified, and it is uniform
because the only \(M\)-dependent quantity in the log-integrand is the
linear term.

\subsection{The finite filter on the Laplace line}

For the finite filter, \(W(u)=\frac1{2\pi i}\int_{\operatorname{Re}y=y^*}e^{\Lambda(y)}\dd y\)
with
\begin{equation}
 \Lambda(y)=\sum_{m<M}\ln\frac{(y-a_m)^2+\gamma^2}{|\lambda_m|^2}
 +\ln\frac{\Gamma\bigl(\tfrac{d+y}\varepsilon\bigr)\Gamma(n+1)}{\varepsilon\,\Gamma\bigl(\tfrac{d+y}\varepsilon+n+1\bigr)}+yu,
 \label{eq:f2v82-Lambda}
\end{equation}
real on the real axis, and \(y^*\) its real saddle, \(\Lambda'(y^*)=0\).
The escape point is where the saddle sits at \(y^*=0\):
\(u=u_1'+\sum_m2a_m/(a_m^2+\gamma^2)\) with
\(u_1'=\varepsilon^{-1}[\psi(\tfrac d\varepsilon+n+1)-\psi(\tfrac d\varepsilon)]=u_1+O(1/M)\);
at \(\kappa=6\) this gives \(u-u_1=1.302,1.957,2.636\) for
\(M=24,48,96\), against the exact last zeros \(u_{\rm last}-u_1=1.32\)
at \(M=24\) (V81) and the mean \(1.274,1.937,2.622\).

\begin{proposition}[Dominance on the vertical line]
\label{prop:f2v82-dominance}
At \(y^*=0\), for all real \(w\),
\begin{equation}
 \ln\Bigl|\frac{e^{\Lambda(iw)}}{e^{\Lambda(0)}}\Bigr|
 \le-\frac{w^2}2\,\kappa_2+w^4\kappa_4,\qquad
 \kappa_2=\varepsilon^{-2}\!\sum_{k\le n}\Bigl(\tfrac d\varepsilon+k\Bigr)^{-2}
 -2\sum_m\frac{a_m^2-\gamma^2}{(a_m^2+\gamma^2)^2},\quad
 \kappa_4=\tfrac14\varepsilon^{-4}\!\sum_{k\le n}\Bigl(\tfrac d\varepsilon+k\Bigr)^{-4}+\sum_m\frac1{(a_m^2+\gamma^2)^2},
 \label{eq:f2v82-mod}
\end{equation}
and the first term of \(\kappa_2\) is \(s^2(1+O(1/M))\) while the pair
term has zero continuum limit, \(\int_0^\infty(a^2-\gamma^2)(a^2+\gamma^2)^{-2}\dd a=0\).
At \(\kappa=6\): \(\kappa_2=0.078,0.068,0.062\) and \(\kappa_4\approx1.5\times10^{-4}\)
for \(M=24,48,96\) (\(s^2=0.058,0.057,0.056\)); the scanned modulus
along \(0\le w\le3d\) is maximal at \(w=0\), with
\(\ln|{\rm ratio}|=-0.16,-0.97\) at \(w=2,5\) and \(-3.1,-2.8,-12.5\) at
\(w=d\).  Beyond \(|w|\ge d\) the exact bound
\(|\Gamma(z)/\Gamma(z+n+1)|\le|\operatorname{Im}z|^{-(n+1)}\) gives decay
\(|w|^{-(M+4)}\) against the multiplier's growth.
\end{proposition}

\begin{proof}
\(\ln|(iw-a)^2+\gamma^2|^2-\ln(a^2+\gamma^2)^2=\ln\bigl(1+\tfrac{2w^2(a^2-\gamma^2)+w^4}{(a^2+\gamma^2)^2}\bigr)\le\tfrac{2w^2(a^2-\gamma^2)+w^4}{(a^2+\gamma^2)^2}\)
and \(-\tfrac12\ln(1+w^2/(z_r+k)^2)\le-\tfrac{w^2}{2(z_r+k)^2}+\tfrac{w^4}{4(z_r+k)^4}\)
with \(z_r=d/\varepsilon\), summed; the derivative identity
\(\varepsilon^{-2}\sum_k(z_r+k)^{-2}=\varepsilon^{-2}[\psi'(z_r)-\psi'(z_r+n+1)]=s^2(1+O(1/M))\)
is V73.  The scan and the values are in the verifier.
\end{proof}

So on the finite Laplace line the integrand is a non-oscillatory
Gaussian-type peak of width \(\kappa_2^{-1/2}\approx s^{-1}\) at the saddle,
with higher derivatives of \(\Lambda\) bounded uniformly in \(M\) by the
same convergent sums as in Proposition~\ref{prop:f2v82-uniform}(ii)
(the Beta part contributes \(\varepsilon^{-k}[\psi^{(k-1)}(z_r)-\psi^{(k-1)}(z_r+n+1)]=O(M^{(3-k)/2})\)
for \(k\ge3\)).  The Laplace method therefore applies to \(W\) exactly
as to \(G\), and the transfer of Theorem~\ref{thm:f2v82-laplace} to the
filter is a matter of the same remainder bookkeeping plus the tail
bound.

\begin{problem}[Bookkeeping of the transfer]
\label{prob:f2v82-bookkeeping}
Write out the Laplace-method remainder for \eqref{eq:f2v82-Lambda} with
explicit constants from \eqref{eq:f2v82-mod} and the uniform derivative
bounds, including the tail \(|w|\ge d\), and conclude
\(u_{\rm last}=u_1+\sum_m2a_m/(a_m^2+\gamma^2)+O(1)\) for all \(M\ge M_0\)
with explicit \(M_0\).  With Lemma~\ref{lem:f2v71-rolle} this is
\(U_{M,d}\ge u_1+\ln M-C\), the V70 conjecture.
\end{problem}

\subsection*{Firewall and V83 direction}
\label{fire:f2v82}
Proposition~\ref{prop:f2v82-uniform} and
Proposition~\ref{prop:f2v82-dominance} are proved (the latter's numerical
values illustrate the bound);
Theorem~\ref{thm:f2v82-laplace} is proved with unspecified constants in
its remainder.  Problem~\ref{prob:f2v82-bookkeeping} is the remaining
step to the V70 conjecture; it is now a routine but long estimate, not a
conceptual gap.  GRH remains open.

V83 should do the bookkeeping: bound \(\Lambda^{(3)}\) and \(\Lambda^{(4)}\)
on \(|w|\le\kappa_2^{-1/2}\ln M\), apply the standard Laplace remainder,
control the tail by the \(|w|^{-(M+4)}\) decay, and extract \(M_0\); then
V84 records the corollary for the archive (the equispaced family is not
Euler-positive beyond any fixed conductor for \(M\ge M_0\)) and V85 is
the sweep.

\subsection*{Assumption ledger}
\label{ass:f2v82}
The contour shift and the modulus inequalities are exact; the Laplace
remainders are standard but their constants are not written; the
numerical values are from the verifier.
\endgroup

\section*{Version V82 append-only record}
\addcontentsline{toc}{section}{Version V82 append-only record}

V82 was appended byte-for-byte before the terminal marker of validated
V81, whose installed SHA--256 was
\[
 \texttt{5B2A0570E0C894178887EB4BBF58C6B2}
 \allowbreak\texttt{B48F87EDB8E906C25B3BB8CE270D2986}
\]
with \(21553\) lines and \(777669\) bytes.  It proves the uniform
saddle structure of the limit symbol, the Laplace-method escape formula
with its \(O(1)\) correction, and the dominance of the vertical Laplace
line for the finite filter, reducing the transfer to remainder
bookkeeping.  After validation V82 replaces V81 as the sole active
numeric GRH note; the frozen \texttt{V100\_f2} archive is unchanged.

% =============================================================================
% END APPENDED V82 MATERIAL
% =============================================================================




\clearpage
\section{V83: the transfer bookkeeping is not routine; exact escape certificates to \texorpdfstring{$M=96$}{M=96}}
\label{sec:f2v83-bookkeeping}
\begingroup
\setcounter{equation}{0}
\renewcommand{\theequation}{N83.\arabic{equation}}
\renewcommand{\theHequation}{f2v83.\arabic{equation}}

V82 called the remaining transfer step ``routine remainder bookkeeping.''
This iteration attempted it and withdraws that description.  The
leading Laplace term is accurate at the escape point itself, but the
two devices that would turn it into a lower bound for the last zero both
need phase control that the modulus inequality cannot supply: the
absolute-value bound at a point left of the peak exceeds the peak value
for \(M\ge48\), and the moving-saddle approximation degrades with
\(M\).  What is delivered instead is exact: the pole-and-pairs filter is
certified negative half a unit left of its predicted escape point for
\(M=24,48,96\), which extends the exact lower bounds on the filter's
last zero from \(M\le24\) (V70--V72) to \(M=96\).  Checked by
\path{scripts/verify_grh_v83_transfer_bookkeeping.py} (SHA--256
\texttt{E0AC6791C19E19E55828C2CA7DCF580F}\allowbreak\texttt{FF099D2B8B2EF01E11E9A497F8122975}).

\subsection{Two devices and what they need}

\begin{lemma}[Escape devices]
\label{lem:f2v83-devices}
Let \(K_{\rm pp}(u)\) be the pole-and-pairs filter (up to a positive
constant, \(-e^{du}W'(u)/d\)), which tends to a positive limit as
\(u\to\infty\), and let \(U_{\rm pp}\) be its last zero.
\begin{enumerate}[label=\textup{(\roman*)}]
\item If \(K_{\rm pp}(u_2)<0\), then \(U_{\rm pp}>u_2\).
\item If \(W(u_3)>W(u_2)>0\) for some \(u_3>u_2\), then \(U_{\rm pp}>u_2\).
\end{enumerate}
In either case the real ladder factors push the last zero further out
(Lemma~\ref{lem:f2v71-rolle}), so \(U_{M,d}>u_2\).
\end{lemma}

\begin{proof}
(i) is the intermediate value theorem; (ii) is that \(W\to0^+\) at
\(+\infty\) and has an interior maximum on \([u_2,\infty)\), i.e.\ a
critical point, which is a zero of \(K_{\rm pp}\).
\end{proof}

Both devices need a sign or a comparison at a point \(u_2\) to the left
of the peak \(u_c=u_1'+\sum_m2a_m/(a_m^2+\gamma^2)\).  The verifier
measures what the available estimates give there (\(\kappa=6\),
\(\kappa_\chi=0\)):
\begin{equation}
\begin{array}{c|c|cc|cc}
 M & \dfrac{W(u_c)}{\text{Laplace leading}} &
 \dfrac{\text{best modulus bound of }W(u_c-\tfrac12)}{W(u_c)} & \dfrac{W(u_c-\tfrac12)}{W(u_c)} &
 \dfrac{\text{bound of }W(u_c-1)}{W(u_c)} & \dfrac{W(u_c-1)}{W(u_c)}\\ \hline
 24 & 0.996 & 1.09 & 0.22 & 0.67 & 0.06\\
 48 & 0.977 & 1.93 & 0.21 & 1.67 & 0.28\\
 96 & 0.955 & 2.85 & 0.05 & 2.48 & -0.34
\end{array}
\label{eq:f2v83-table}
\end{equation}
The leading Laplace term at the saddle \(y=0\) is within \(5\%\) of
\(W(u_c)\), so a lower bound there is within reach.  But the upper bound
of \(W(u_c-\delta)\) by \(\frac1{2\pi}\int|e^{\Lambda(c+iw;u_c-\delta)}|\dd w\),
minimized over the vertical line \(\operatorname{Re}y=c\), exceeds
\(W(u_c)\) for \(M\ge48\): the integrand on any vertical line still
oscillates enough that its modulus loses the factor \(5\)--\(20\) by which
\(W(u_c-\delta)\) is actually smaller.  Device (ii) therefore needs
phase control at the left point too, not only at the peak.  For device
(i), the moving-saddle Laplace term for \(W'(u_c-\delta)\) (saddle
\(y_2^*\approx3\)--\(7\)) is off by factors \(1.1,1.4,1.8\) at
\(\delta=\tfrac14\) and \(1.1,2.6,7.9\) at \(\delta=\tfrac12\) for
\(M=24,48,96\): the saddle enters the region where the pair
singularities \(a_m\pm i\gamma\) are within one Gaussian width, and the
leading term is no longer uniform.  The bookkeeping is thus a genuine
two-point phase-controlled asymptotic, which is not written here.

\subsection{Exact certificates}

\begin{proposition}[Escape certificates to \(M=96\)]
\label{prop:f2v83-certificates}
At \(\kappa=6\), \(\kappa_\chi=0\), with \(d\) rational as in V70, the
pole-and-pairs filter satisfies \(K_{\rm pp}(u_2)<0\) exactly at
rational points \(x_2=e^{-\varepsilon u_2}\) with \(u_2\ge u_c-\tfrac12\):
\begin{equation}
\begin{array}{c|ccccc}
 M & n & u_2-u_1 & \ln M & \sum_m2a_m/(a_m^2+\gamma^2) & u_1\\ \hline
 24 & 75 & 0.802 & 3.18 & 1.274 & 2.07\\
 48 & 147 & 1.457 & 3.87 & 1.937 & 2.87\\
 96 & 291 & 2.136 & 4.56 & 2.622 & 4.02
\end{array}
\label{eq:f2v83-cert}
\end{equation}
Hence \(U_{M,d}\ge u_1+0.80,\;u_1+1.46,\;u_1+2.14\) respectively, and
the certified excess over \(u_1\) grows like \(\ln M-2.4\); the growing
filter is not Euler-positive beyond \(L_f\) at these scales for any
conductor with \(L_f<u_1+2.1\).
\end{proposition}

\begin{proof}
Exact rational evaluation of \(\sum_j\binom nj(-x_2)^j\widetilde\Psi(j)(1+j\varepsilon/d)\)
in the verifier, followed by Lemma~\ref{lem:f2v83-devices}(i) and
Lemma~\ref{lem:f2v71-rolle}.
\end{proof}

\subsection*{Firewall and V84 direction}
\label{fire:f2v83}
Lemma~\ref{lem:f2v83-devices} and
Proposition~\ref{prop:f2v83-certificates} are proved; the table
\eqref{eq:f2v83-table} is numerical.  V82's characterization of the
transfer as routine is withdrawn: Problem~\ref{prob:f2v82-bookkeeping}
requires phase control at two points, and the modulus inequality of
Proposition~\ref{prop:f2v82-dominance} is insufficient at the left point.
The V70 conjecture stands proved for the limit symbol (V80), exactly
certified for the finite filter up to \(M=96\), and open in general.  GRH
remains open.

V84 should pursue device (i) at the peak itself, where the leading term
is accurate: prove \(W'(u_c-\delta)>0\) for a \emph{small}
\(\delta\propto\sigma\) (one Gaussian width, \(\delta\approx0.25\)), where
the moving saddle \(y_2^*\approx\delta/\sigma^2\approx3\) is still well
below the first pair singularity distance \(\gamma=14\); the Laplace
remainder there needs \(\Lambda'''\), \(\Lambda''''\) bounds on
\(|w|\le\sigma^{-1}\ln M\) (available from the convergent sums) and a
phase bound \(|\operatorname{Im}\Lambda(y_2^*+iw)-\text{cubic}|\) that
the cubic term already controls to \(0.2\) rad at \(w=4\).  If that
closes, \(U_{M,d}\ge u_c-\delta=u_1+\ln M-C\) for \(M\ge M_0\).  V85 is
the mandatory sweep.

\subsection*{Assumption ledger}
\label{ass:f2v83}
Exact rational arithmetic for the certificates; the Laplace and modulus
values are floating-point diagnostics; the devices use only the
intermediate value theorem and Rolle.
\endgroup

\section*{Version V83 append-only record}
\addcontentsline{toc}{section}{Version V83 append-only record}

V83 was appended byte-for-byte before the terminal marker of validated
V82, whose installed SHA--256 was
\[
 \texttt{03AEED3120AD42001861A554ADA2C6FE}
 \allowbreak\texttt{A87DCD0C58138F4AC19D3C2C8A9F8630}
\]
with \(21763\) lines and \(787925\) bytes.  It withdraws the
``routine'' characterization of the transfer bookkeeping with
quantitative evidence, states the two escape devices, and certifies the
pole-and-pairs filter's escape exactly up to \(M=96\).  After validation
V83 replaces V82 as the sole active numeric GRH note; the frozen
\texttt{V100\_f2} archive is unchanged.

% =============================================================================
% END APPENDED V83 MATERIAL
% =============================================================================




\clearpage
\section{V84: verdict on the equispaced filter family, and the escape law certified to \texorpdfstring{$M=192$}{M=192}}
\label{sec:f2v84-verdict}
\begingroup
\setcounter{equation}{0}
\renewcommand{\theequation}{N84.\arabic{equation}}
\renewcommand{\theHequation}{f2v84.\arabic{equation}}

V83 left the finite-\(M\) transfer as a two-point phase-controlled
asymptotic.  This iteration tested the two Fourier-side contours on
which such an asymptotic would have to live and found that neither is
absolutely dominated by its central region: on the line through the
moving saddle the multiplier's modulus is unbounded in the frequency,
and on the line \(\operatorname{Re}y=0\) the multiplier's growth nearly
cancels the Beta decay beyond one Gaussian width, leaving an oscillatory
plateau whose absolute mass exceeds the central part by a factor growing
with \(M\).  The escape is therefore a cancellation phenomenon in every
Fourier representation, and the only cancellation-free representation
is the one available in the limit (V80).  With that, the programme's
arithmetic side is given a verdict: the equispaced two-channel filter of
the archive is not Euler-positive beyond any fixed conductor at any
critical scale, proved for the Gaussian limit and certified exactly for
\(M\le192\) with the excess following \(\ln M-2.4\).  Checked by
\path{scripts/verify_grh_v84_verdict.py} (SHA--256
\texttt{36DD172CC3F28777DE930F62FD63F33B}\allowbreak\texttt{8796BC7D9A400A7E07C3F107C0E21420}).

\subsection{The two contours}

\begin{proposition}[No absolutely dominated Fourier contour]
\label{prop:f2v84-contours}
For \(\kappa=6\), \(\kappa_\chi=0\), \(u_2=u_c-\tfrac12\):
\begin{enumerate}[label=\textup{(\roman*)}]
\item on the vertical line through the real saddle \(y_2^*\) of
\(\Lambda(\cdot;u_2)\) (\(y_2^*=5.9,6.9,7.1\) for \(M=24,96,384\)), the
absolute integral of \(|(y_2^*+iw)e^{\Lambda}|\) over \(|w|>3\) exceeds the
central sine-positive part by factors of \(10^4\) to \(10^6\): the line
lies inside the ladder \(a_m=1,3,5,\dots\), within one height \(\gamma\)
of the pair zeros, and \(\prod_m|(y_2^*-a_m+iw)^2+\gamma^2|\) is unbounded
in \(|w|\) faster than the Beta factor decays;
\item on the line \(\operatorname{Re}y=0\), where the modulus is maximal
at \(w=0\) (Proposition~\ref{prop:f2v82-dominance}), the integrand of
\(W'(u_c-\delta)\), \(\delta=0.3\), has a central part
\(\int_0^{w_1}we^{R(w)}\sin(w\delta-\varphi_3(w))\dd w\approx9\)--\(12\)
(\(w_1\approx9\)) and an outer absolute part \(\int_{w_1}^{2d}we^{R(w)}\dd w\)
of \(26,233,339\) for \(M=24,96,384\): beyond \(|w|\approx10\) the
function \(R(w)=\operatorname{Re}\Lambda_0(iw)-\Lambda_0(0)\) flattens near
\(-3\) instead of decaying (the pair growth balances the Beta decay), so
the tail is an oscillatory plateau of large absolute mass.
\end{enumerate}
\end{proposition}

\begin{proof}
Numerical quadrature of the exact modulus and phase in the verifier;
(i) and (ii) are asserted there as inequalities between the computed
parts.
\end{proof}

So the finite-\(M\) last zero cannot be certified by absolute bounds on
either contour; the cancellation in the plateau must be proved, which is
exactly the kind of estimate the Gaussian-kernel representation of V80
avoids and that the finite kernel \(V\) (Proposition~\ref{prop:f2v81-mixture})
does not admit.  The transfer Problem~\ref{prob:f2v80-transfer} is
genuinely harder than the limit theorem, and it is left open.

\subsection{Certified escape to \(M=192\)}

\begin{proposition}[Escape law on certificates]
\label{prop:f2v84-law}
At \(\kappa=6\), \(\kappa_\chi=0\), the pole-and-pairs filter is negative
exactly at a rational point at least \(u_c-\tfrac12\), for
\(M=24,48,96,192\) (\(n=75,147,291,579\)); hence
\begin{equation}
\begin{array}{c|cccc}
 M & 24 & 48 & 96 & 192\\ \hline
 U_{M,d}-u_1\ \ge & 0.802 & 1.457 & 2.136 & 2.823\\
 \ln M-2.4 & 0.78 & 1.47 & 2.16 & 2.86
\end{array}
\label{eq:f2v84-law}
\end{equation}
and \(u_1=\varepsilon^{-1}\ln(1+n\varepsilon/d)\) grows like
\(0.405\sqrt M\).  At \(M=192\) the filter is positive again at
\(u_c-1\), so the certified point is within half a unit of the last zero.
\end{proposition}

\begin{proof}
Exact rational evaluation (V83 for \(M\le96\), the verifier for
\(M=192\)) and Lemma~\ref{lem:f2v83-devices}(i) with
Lemma~\ref{lem:f2v71-rolle}.
\end{proof}

\subsection{Verdict}

\begin{theorem}[Verdict on the equispaced family]
\label{thm:f2v84-verdict}
For the archive's equispaced two-channel filter \(K_{M,d}\) on a critical
line \(d=\kappa\sqrt M\) with \(\kappa\) below the barrier:
\begin{enumerate}[label=\textup{(\roman*)}]
\item its last zero satisfies \(U_{M,d}\ge u_1=\varepsilon^{-1}\ln(1+n\varepsilon/d)\)
whenever the real ladder alone is present (Theorem~\ref{thm:f2v71-real}),
and \(u_1\sim R\ln(1+\theta)\sqrt M/(\theta\kappa)\);
\item with the target pair ladder, the Gaussian-limit symbol's escape
point is \((\ln M)/s+O(1)\) for all \(M\ge M_0\) (Theorems~\ref{thm:f2v80-escape},
\ref{thm:f2v82-laplace}), and the finite filter's last zero exceeds
\(u_1+\ln M-2.4\) exactly for \(M=24,48,96,192\);
\item consequently the real-only filter (principal ladder with any
finite real padding) fails Euler positivity beyond every fixed conductor
scale \(L_f\) for all \(M>(\theta\kappa L_f/(R\ln(1+\theta)))^2\)
(unconditional), and the full two-channel filter fails it for
\(M\le192\) (exact) and for all \(M\ge M_0\) in the Gaussian limit
(Theorem~\ref{thm:f2v80-escape}); the finite-\(M\) transfer for
\(M>192\) is Problem~\ref{prob:f2v80-transfer}.
\end{enumerate}
In the cases covered, the family cannot host the contradiction with the
forced compensation (V80.13) at any critical scale, above, inside, or
below the band; the band theorems (V52--V61, (H61)) are not the
obstruction, the pole normalization is.
\end{theorem}

\begin{proof}
(i) is Theorem~\ref{thm:f2v71-real}: the pole factor applied first
gives the zero at \(u_1\), and the real factors push it out.  (ii) is
V80 and V82 for the limit and Proposition~\ref{prop:f2v84-law} for the
certificates.  (iii) follows from (i) for the real-only filter since
\(U_{M,d}\ge u_1\to\infty\); for the full filter it is (ii) restated.
Nothing is claimed for the full filter with \(M>192\) beyond the limit
theorem.
\end{proof}

The theorem's status is thus: exact for \(M\le192\), rigorous for the
real-only filter at every \(M\) and for the limit symbol at every large
\(M\), and numerically established to within a few percent at every
tested size for the full finite filter.
The practical consequence for the programme is unchanged from V70: the
arithmetic side must find a filter family whose alternating polynomial
keeps its real roots within \(O(1/\sqrt M)\) of \(q=1\), which no
pole-normalized rational filter of this type does, or replace Euler
positivity by a size bound on the prime side.

\subsection*{Firewall and V85 direction}
\label{fire:f2v84}
Propositions~\ref{prop:f2v84-contours} and \ref{prop:f2v84-law} are
established as stated (numerical inequalities and exact certificates);
Theorem~\ref{thm:f2v84-verdict}(i) is proved, (ii) is proved for the
limit and certified to \(M=192\), (iii) is proved from (i) for the
real-only filter and conditional on the transfer for the full filter.
GRH remains open.

V85 is the mandatory companion sweep.  After it the programme should
return upstream: with the equispaced family retired, examine whether the
archive's forced compensation (V80.13) can be tested by a
\emph{signed} weight with a prime-side bound in place of Euler
positivity, or whether the pole normalization can be replaced by a
smoothed target response that does not create a zero at \(u_1\).  The
finite-block band work is not to be resumed as a primary route.

\subsection*{Assumption ledger}
\label{ass:f2v84}
Exact rational arithmetic for the certificates; floating quadrature for
the contour diagnostics; the verdict's clause (iii) for the full filter
beyond \(M=192\) rests on the limit theorem and the numerical transfer.
\endgroup

\section*{Version V84 append-only record}
\addcontentsline{toc}{section}{Version V84 append-only record}

V84 was appended byte-for-byte before the terminal marker of validated
V83, whose installed SHA--256 was
\[
 \texttt{3124646C0F3F3B27DCDA870C29A26D14}
 \allowbreak\texttt{321BE81E20B630F2C1504ED1F439FB86}
\]
with \(21919\) lines and \(795190\) bytes.  It shows that no
Fourier-side contour is absolutely dominated at finite \(M\), certifies
the escape law to \(M=192\), and states the verdict on the equispaced
family.  After validation V84 replaces V83 as the sole active numeric
GRH note; the frozen \texttt{V100\_f2} archive is unchanged.

% =============================================================================
% END APPENDED V84 MATERIAL
% =============================================================================




\clearpage
\section{V85: companion sweep; the two upstream branches after the verdict, and the real-factor zero law}
\label{sec:f2v85-upstream}
\begingroup
\setcounter{equation}{0}
\renewcommand{\theequation}{N85.\arabic{equation}}
\renewcommand{\theHequation}{f2v85.\arabic{equation}}

V85 is the mandatory companion sweep.  Its mathematical part tests the
two alternatives that the verdict of V84 left standing: replacing Euler
positivity by a size bound on a signed weight, and replacing the pole
normalization by a smoothed target response.  Both are closed for the
archive's filter class, the first quantitatively (the weight is larger
than its target response by a factor \(e^{10M}\) on the prime
logarithms), the second exactly (any real factor at distance
\(O(\sqrt M)\) from the origin creates the \(\sqrt M\) zero, and the
principal ladder itself sits there).  The second closure isolates an
invariant of the whole class, the real-factor zero law, which says in
one line what a filter family must change to avoid the verdict.  Checked
by \path{scripts/verify_grh_v85_upstream_branches.py} (SHA--256
\texttt{2D1E2960CB5D69D8A898DFDCCDB4CBB8}\allowbreak\texttt{11FE7BF62F9AF4675F63566892C8223F}).

\subsection{Companion sweep}

The companion heads at this sweep are
\begin{equation}
\begin{array}{l|c|c}
 \text{programme} & \text{head (lines)} & \text{SHA--256}\\ \hline
 \text{SM--Einstein} & \text{V37 } (6840) & \texttt{D4E8D9124B3BBA09492C03E31AE8DE13}\allowbreak\texttt{8F04F4B80BD154A488CF067C1F6EA1D9}\\
 \text{Yang--Mills} & \text{V41 } (14452) & \texttt{2EEC02EB12520BAA0651AB07A1719708}\allowbreak\texttt{09451FB0DEB83B596B084312F9EE4961}\\
 \text{Navier--Stokes} & \text{V26 } (5319) & \texttt{82C887F8C2C69E4F28FAD7C3DC8DE5B9}\allowbreak\texttt{099CB5A4BF431EBA1FFF3E91935978AD}
\end{array}
\label{eq:f2v85-sources}
\end{equation}
Yang--Mills and Navier--Stokes are byte-identical to the V80 record.
SM--Einstein advanced from V33 to V37 (four sections: the wave-kinetic
equation of the decay channels, an audit, the gravitational reading of
the conversions, and modes in an expanding metric).

\begin{proposition}[Sweep decision]
\label{prop:f2v85-sweep}
No statement in SM--Einstein V34--V37 concerns real-rootedness, Hermite
or Appell families, Gamma ladders, backward heat flow, Laplace transforms
of rational filters, Euler-product positivity, signed prime weights,
Dirichlet \(L\)-functions, or exact rational positivity certificates;
the only occurrences of these words are in the SM audit record's
description of this note.  Nothing is imported.  (The SM audit's ledger
lists this note at V81; that is a bookkeeping lag on the SM side, to be
corrected at its next audit, V40.)
\end{proposition}

\subsection{The signed-weight branch}

The archive's Euler-admissible weight (V80.7) is normalized by a unit
target response, \(F(x_d)=1\), where \(F(y)=\sum_jb_j/(y+j\varepsilon)\)
is the Laplace transform of \(K_{M,d}(u)=\sum_jb_je^{-j\varepsilon u}\).
Writing \(D(y)=\prod_{j=0}^n(y+j\varepsilon)\), one has
\(F=C\,P/D\) with \(C=D(x_d)/P(x_d)\), so
\(b_j=C\,P(-j\varepsilon)/D'(-j\varepsilon)\) and in particular
\begin{equation}
 b_0=\frac{D(x_d)}{P(x_d)}\cdot\frac{P(0)}{D'(0)},\qquad
 K_{M,d}(u)=b_0\,H(e^{-\varepsilon u}).
\label{eq:f2v85-b0}
\end{equation}
If Euler positivity is dropped, the forced compensation
\(\Re\ge m_\rho L_x[W]>0\) (V80.12--13) can only be contradicted by an
upper bound on the prime side, and any such bound is in terms of
\(\sup_{u\ge\ln2}|K_{M,d}(u)|\) times a prime sum, against a target
response of size \(1\).

\begin{proposition}[Size of the weight on the prime logarithms]
\label{prop:f2v85-size}
At \(\kappa=6\), \(\kappa_\chi=0\), with \(b_0\) given by
\eqref{eq:f2v85-b0} (checked exactly against the partial-fraction
definition for \(M=6,12\), where \(\sum_jb_j/(x_d+j\varepsilon)=1\)
holds in rational arithmetic and \(b_0>0\)):
\begin{equation}
\begin{array}{c|cccc}
 M & 6 & 12 & 24 & 48\\ \hline
 \ln b_0 & 65.4 & 116.5 & 214.8 & 401.5\\
 \max_{u\ge\ln2}\ln|H(e^{-\varepsilon u})| & 0.0 & 0.0 & 24.6 & 79.8\\
 \ln\bigl(b_0\max|H|\bigr) & 65.4 & 116.5 & 239.4 & 481.2\\
 \text{per unit of }M & 10.9 & 9.7 & 10.0 & 10.0
\end{array}
\label{eq:f2v85-size}
\end{equation}
The weight exceeds its unit target response by a factor \(e^{10M(1+o(1))}\)
on \(u\ge\ln2\).
\end{proposition}

\begin{proof}
Direct evaluation in the verifier: \(b_0\) by the residue formula in
logarithms, \(H\) by the signed-logarithm sum of the coefficients
\(c_j=\binom njP(-j\varepsilon)/P(0)\) on a grid of \(3000\) points in
\(q\in(0,2^{-\varepsilon}]\).
\end{proof}

So a signed-weight test would need the prime sum
\(\sum_p\sum_k\chi(p^k)\ln p\,p^{-kx}K_{M,d}(k\ln p)\) to cancel to
relative precision \(e^{-10M}\), for every \(M\), with no positivity to
enforce the cancellation.  The branch is closed for this family:
the normalization that makes the target response \(1\) is what makes
the weight exponentially large in \(M\), because
\(P(0)/P(x_d)\) is a product of \(3M+2\) ratios each of order
\(M/\sqrt M\) or \(1\), and \(D(x_d)/D'(0)\) contributes
\(\prod_j(1+x_d/(j\varepsilon))\).

\subsection{The smoothed-pole branch and the real-factor zero law}

\begin{lemma}[Real-factor zero law]
\label{lem:f2v85-zero-law}
Let \(E(u)=(1-e^{-\varepsilon u})^n\) and let
\(\lambda_1\le\lambda_2\le\dots\le\lambda_r\) be positive reals.  The
function \(\prod_i(1-\partial_u/\lambda_i)E\) has its last zero at
\begin{equation}
 U\ \ge\ \max_i\ \varepsilon^{-1}\ln\Bigl(1+\frac{n\varepsilon}{\lambda_i}\Bigr)
 \;=\;\varepsilon^{-1}\ln\Bigl(1+\frac{n\varepsilon}{\lambda_1}\Bigr).
\label{eq:f2v85-law}
\end{equation}
\end{lemma}

\begin{proof}
The operators commute, so apply \((1-\partial_u/\lambda_i)\) first for
any \(i\): the result is \(E\,[1-n\varepsilon e^{-\varepsilon u}/(\lambda_i(1-e^{-\varepsilon u}))]\),
whose only zero on \(u>0\) is at \(e^{\varepsilon u}=1+n\varepsilon/\lambda_i\).
Each remaining real factor pushes the last zero outward
(Lemma~\ref{lem:f2v71-rolle}).  Taking the maximum over \(i\) gives
\eqref{eq:f2v85-law}; the maximum is attained at the smallest \(\lambda\).
\end{proof}

\begin{proposition}[No pole normalization avoids the \(\sqrt M\) zero]
\label{prop:f2v85-smoothed}
In the archive's filter the real factors are \(\lambda\in\{d,\xi_d\}\cup\{\sigma+2m\}\),
all of size \(\kappa\sqrt M+O(1)\) for the principal ladder's first
entries.  Removing or smoothing the pole factor \(d\) replaces the bound
\(u_1=\varepsilon^{-1}\ln(1+n\varepsilon/d)\) by the one from the next
factor \(\xi_d=d+1-\eta\):
\begin{equation}
\begin{array}{c|ccc}
 M & 6 & 24 & 96\\ \hline
 \varepsilon^{-1}\ln(1+n\varepsilon/\xi_d) & 1.161 & 2.068 & 4.011\\
 u_1 & 1.168 & 2.074 & 4.016\\
 \text{ratio} & 0.9944 & 0.9972 & 0.9986
\end{array}
\label{eq:f2v85-xi}
\end{equation}
and removing every factor below the principal ladder replaces it by
\(\varepsilon^{-1}\ln(1+n\varepsilon/\sigma)=u_1(1-O(1/d))\).  The
\(\sqrt M\) growth of the last zero is a property of the principal
zero ladder at distance \(\sigma=\kappa\sqrt M+1\) together with the
lattice spacing \(\varepsilon=\theta d/N\sim\theta\kappa/(3\sqrt M)\):
\begin{equation}
 \varepsilon^{-1}\ln\Bigl(1+\frac{n\varepsilon}{\lambda}\Bigr)
 \;\sim\;\frac{3\sqrt M}{\theta\kappa}\ln\Bigl(1+\frac{\theta\kappa\sqrt M}{\lambda}\Bigr),
\label{eq:f2v85-asym}
\end{equation}
which is \(\asymp\sqrt M\) for \(\lambda\asymp\sqrt M\) and \(O(1)\) only
for \(\lambda\gtrsim M\).
\end{proposition}

\begin{proof}
Lemma~\ref{lem:f2v85-zero-law} with \(\lambda_1=\xi_d\), or
\(\lambda_1=\sigma\); the table is the exact one-factor formula
evaluated in the verifier, and \eqref{eq:f2v85-asym} is
\(n\varepsilon=\theta d\,n/N\).
\end{proof}

The smoothed-pole branch is therefore closed as well: the pole is not
the first real factor by much, and the principal ladder (which cannot be
moved, since its zeros are what cancel the principal-channel poles at
the target) is at the same distance.  Theorem~\ref{thm:f2v84-verdict}'s
attribution ``the pole normalization is the obstruction'' is thereby
sharpened: the obstruction is the \emph{real-factor zero law} for the
pair (principal ladder at distance \(\sqrt M\), lattice spacing
\(1/\sqrt M\)).

\begin{proposition}[Sign at the first prime]
\label{prop:f2v85-first-prime}
For the \(24\) V70 instances, \(K_{M,d}(\ln2)<0\) exactly in \(11\) and
\(>0\) in \(13\); in all \(24\) the V70 certificate point
\(u_0=-\varepsilon^{-1}\ln q_0\) exceeds \(\ln2\), so the weight changes
sign inside the range of prime logarithms in every instance.
\end{proposition}

\begin{proof}
Exact rational evaluation at \(x=e^{-\varepsilon\ln2}\) rounded to a
rational with denominator \(\le10^9\) (the sign is the exact sign at
that rational point), and the V70 search re-run.
\end{proof}

\subsection{What the law says a family must change}

By \eqref{eq:f2v85-asym} the last zero is \(O(1)\) only if
\(n\varepsilon/\lambda_1=O(\varepsilon)\), i.e.\ if the exponent lattice
reaches at most a bounded multiple of \(\varepsilon\lambda_1\).  With the
archive's kernel this reads \(\theta d\lesssim\varepsilon\sigma\), false
by a factor \(\sqrt M\).  Two changes are conceivable and neither is
free:
\begin{enumerate}[label=\textup{(\alph*)}]
\item lattice spacing \(\varepsilon\asymp1\) with \(n\asymp M\): the
one-factor zero moves to \(\ln(1+n\varepsilon/\lambda)\asymp\tfrac12\ln M\),
the same order as the pair escape \((\ln M)/s\) of V80; the band
structure of the archive (\(\varepsilon=\theta d/N\) is what places the
kernel's exponents in \([0,\theta d]\)) would have to be rebuilt;
\item a kernel that is not a real-rooted alternating sum, so that
Lemma~\ref{lem:f2v71-rolle} does not apply; then the real-rootedness
inputs of V70--V72 are lost and Euler positivity has to be certified
differently.
\end{enumerate}

\begin{problem}[Filter class beyond the zero law]
\label{prob:f2v85-class}
Construct a family of Euler-admissible weights (V80.7: \(W\ge0\) on the
prime logarithms, \(L^{(k)}_d[W]=0\), \(L^{(k)}_x[W]>0\)) at critical
scales \(d=\kappa\sqrt M\) whose last sign change on \(u>0\) stays
bounded as \(M\to\infty\), or prove that the two conditions
\(L^{(k)}_d[W]=0\) (which forces real factors at distance \(d\) from the
origin) and \(W\ge0\) on \(u\ge\ln2\) are incompatible for every
finite exponential sum with exponents in an arithmetic progression of
spacing \(\varepsilon\lesssim1/\sqrt M\).  Lemma~\ref{lem:f2v85-zero-law}
proves the second alternative for real-rooted kernels
\(E=(1-e^{-\varepsilon u})^n\) with \(n\varepsilon\gg\varepsilon d\).
\end{problem}

\subsection*{Firewall and V86 direction}
\label{fire:f2v85}
Lemma~\ref{lem:f2v85-zero-law} is proved.
Propositions~\ref{prop:f2v85-size} and \ref{prop:f2v85-first-prime} are
established as stated (the first is a floating-point table with an exact
check of the residue formula at \(M=6,12\); the second is exact);
Proposition~\ref{prop:f2v85-smoothed} is the lemma applied.  No new
positivity is claimed; no group of (H61) is touched; GRH remains open.

V86 should take Problem~\ref{prob:f2v85-class} at alternative (a):
compute, for a Beta kernel with \(\varepsilon\asymp1\) and the same
principal and target ladders, whether the conditions
\(L^{(k)}_d[W]=0\) can still be met (the exponents then extend to
\(n\varepsilon\asymp M\gg d\)), and if so what the last zero is at
\(M=6,\dots,48\).  A negative answer there would strengthen the
incompatibility alternative; a positive one gives a candidate family
whose escape is \(O(\ln M)\) and for which the arithmetic forcing
(bridge items (A5)--(A6)) can be re-examined.  Per the programme rule,
this is a test of a proposed general theorem, not a certificate search.

\subsection*{Assumption ledger}
\label{ass:f2v85}
Exact rational arithmetic for the partial-fraction check and the
first-prime signs; floating logarithms for the size table; the zero
law uses only Rolle's theorem and the commutation of the ladder
operators.
\endgroup

\section*{Version V85 append-only record}
\addcontentsline{toc}{section}{Version V85 append-only record}

V85 was appended byte-for-byte before the terminal marker of validated
V84, whose installed SHA--256 was
\[
 \texttt{0A3B10506C96AF6E9FB17A7F0850AEE5}
 \allowbreak\texttt{330D483DAB2B7CC3DBF1639FB8C63693}
\]
with \(22103\) lines and \(804066\) bytes.  It records the companion
sweep, closes the signed-weight and smoothed-pole branches for the
archive's filter class, states the real-factor zero law, and poses the
filter-class problem.  After validation V85 replaces V84 as the sole
active numeric GRH note; the frozen \texttt{V100\_f2} archive is
unchanged.

% =============================================================================
% END APPENDED V85 MATERIAL
% =============================================================================


\clearpage
\section{V86: audit of V85 and an explicit three-term filter at critical scales}
\label{sec:f2v86-three-term}
\begingroup
\setcounter{equation}{0}
\renewcommand{\theequation}{N86.\arabic{equation}}
\renewcommand{\theHequation}{f2v86.\arabic{equation}}

This continuation starts from the installed V85, not from the superseded
height-twelve frontier.  V51 already completed the finite three-pair
problem, and V69--V85 returned to the arithmetic filter.  The proposed
universal incompatibility in Problem~\ref{prob:f2v85-class} is too broad:
the archive's V76 already proves qualitative feasibility, and its proof
works at any prescribed positive lattice spacing by rescaling the time
variable.  We do not repeat that existence theorem.  Instead we give a
closed three-term formula, a uniform bound on its last sign change, and
explicit upper and lower bounds on the normalized target cost.  These
quantitative statements distinguish the bare pole-cancellation conditions
from the additional growing-ladder interpolation conditions.

\subsection{Audit of inherited claims and companion changes}

The following corrections govern the interpretation of the retained V85.
They do not change its source text or the exact rational evaluations.
\begin{enumerate}
\item The grid and floating-point logarithms in
Proposition~\ref{prop:f2v85-size} are diagnostics.  They prove neither
the asserted supremum nor the asymptotic
\(e^{10M(1+o(1))}\).  A large supremum would itself not be a lower
bound for a signed prime sum.  The conclusion that the entire signed-weight
branch is closed therefore does not follow from that table.
\item The verifier of Proposition~\ref{prop:f2v85-first-prime} evaluates
a nearby rational \(q\), not \(2^{-\varepsilon}\).  Without an
enclosure of the latter and a sign bound on that enclosure, this does not
certify the sign at \(\log2\).  Negativity at some real \(u\ge\log2\)
also need not be negativity at a discrete prime-power logarithm.
\item The real-factor zero law is a valid statement about products of
real first-order factors.  It does not by itself propagate a lower bound
through a complex conjugate-pair factor.  Conclusions about the full
ladder require their own sign argument or exact certificates.
\item The numerical contour diagnostics of V83--V84 and the finite
lists of certificates give no all-size finite-filter theorem.  In
particular, two signs do not locate the last zero unless all later zeros
have been excluded.  The finite-to-limit transfer remains open.
\end{enumerate}

The user-requested restart sweep found SM--Einstein V41 in place of the
V37 recorded in V85.  Its V38--V41 concern the two-polarization metric,
matter stresses, a companion audit, and circle-length scaling of Yukawa
triads.  No identity about the present Laplace functionals is supplied.
The source inspected has SHA--256
\begin{center}
\texttt{8ACA6F9361691EF65AF81A82F4356BCC9}\allowbreak
\texttt{9C01E4F323060F65AB3D559545055B4}.
\end{center}
The Yang--Mills V41 and Navier--Stokes V26 heads still have the hashes
recorded in V85; no companion estimate is imported.  The next scheduled
ten-version sweep is V90.  The pre-existing exploratory file
\path{scripts/verify_grh_v86_unit_lattice_test.py} is not used as a proof:
its prime-point step has the same rounded-input issue just identified.

\subsection{Exact construction}

For an integer \(k\ge1\), \(d>0\), and \(\varepsilon>0\), put
\begin{equation}
 r=k+1,\qquad a=\frac{\varepsilon}{d},\qquad
 A=(1+a)^r,\qquad B=(1+2a)^r,
 \qquad W(u)=1-2A e^{-\varepsilon u}+B e^{-2\varepsilon u}.
 \label{eq:f2v86-weight}
\end{equation}
Retain the archive's convention
\begin{equation}
 F(y)=\mathcal L_y^{(k)}[W]
 :=\frac1{(k-1)!}\int_0^\infty u^{k-1}e^{-yu}W(u)\,du
 =\frac1{y^k}-\frac{2A}{(y+\varepsilon)^k}
       +\frac{B}{(y+2\varepsilon)^k}.
 \label{eq:f2v86-response}
\end{equation}

\begin{theorem}[Double pole cancellation and a positive real target]
\label{thm:f2v86-response}
The weight in \eqref{eq:f2v86-weight} satisfies
\begin{equation}
 F(d)=F'(d)=0,\qquad F(y)>0\quad(y>0,\ y\ne d).
 \label{eq:f2v86-cancellation}
\end{equation}
For \(x=d+\delta\), \(\delta>0\), its target response obeys
\begin{equation}
 \frac{k(k+1)\delta^2\varepsilon^2}
      {d^2(x+2\varepsilon)^{k+2}}
 \le F(x)\le
 \frac{k(k+1)\delta^2\varepsilon^2(d+2\varepsilon)^{k-1}}
      {d^{k+1}x^{k+2}}.
 \label{eq:f2v86-target-bounds}
\end{equation}
\end{theorem}

\begin{proof}
Substitution at \(d\) gives, respectively,
\(d^{-k}[1-2(1+a)+(1+2a)]=0\) and
\(-kd^{-k-1}(1-2+1)=0\).  For fixed real \(y>0\), define
\(f_y(t)=(d+t)^{k+1}(y+t)^{-k}\).  Direct differentiation gives
\begin{equation}
 f_y''(t)=k(k+1)(y-d)^2
               \frac{(d+t)^{k-1}}{(y+t)^{k+2}}.
 \label{eq:f2v86-derivative}
\end{equation}
Twice applying the fundamental theorem of calculus yields the exact
identity
\begin{equation}
 F(y)=\frac{k(k+1)(y-d)^2}{d^{k+1}}
       \int_0^{\varepsilon}\!\int_0^{\varepsilon}
       \frac{(d+s+t)^{k-1}}{(y+s+t)^{k+2}}\,ds\,dt.
 \label{eq:f2v86-integral}
\end{equation}
The integrand is positive.  Replacing its numerator and denominator by
their respective endpoint bounds at \(y=x\) proves
\eqref{eq:f2v86-target-bounds}.
\end{proof}

\begin{theorem}[Two hidden sign changes and a uniform last-zero bound]
\label{thm:f2v86-last-zero}
The weight \(W\) has exactly two simple zeros on \((0,\infty)\), is
positive before the first and after the second, and is negative between
them.  Its last zero is
\begin{equation}
 U=\varepsilon^{-1}\log\bigl(A+\sqrt{A^2-B}\bigr),\qquad
 0<U<\frac{k+1+\sqrt{k+1}}{d}.
 \label{eq:f2v86-last-zero}
\end{equation}
For fixed \(k\), the upper constant is sharp as
\(\varepsilon/d\to0\): \(dU\to k+1+\sqrt{k+1}\).
\end{theorem}

\begin{proof}
The polynomial \(1-2Aq+Bq^2\) has discriminant
\(4(A^2-B)>0\), since \((1+a)^2>1+2a\).  Strict convexity of
\(t\mapsto t^r\), with \(r\ge2\), gives \(B-2A+1>0\).
Consequently its two roots are
\(q_-=1/(A+\sqrt{A^2-B})\) and
\(q_+=1/(A-\sqrt{A^2-B})\), with
\(0<q_-<q_+<1\).  The substitution \(q=e^{-\varepsilon u}\)
proves the sign statement and exact last-zero formula.

Set \(z=a^2/(1+a)^2\).  Then
\[
 U=\varepsilon^{-1}\left\{
 r\log(1+a)+\log\bigl(1+\sqrt{1-(1-z)^r}\bigr)\right\}.
\]
For \(0<z<1\), Bernoulli's inequality gives
\(1-(1-z)^r\le rz\).  Using \(\log(1+t)<t\) for \(t>0\)
therefore bounds the braces strictly by
\(ra+\sqrt r\,a/(1+a)<(r+\sqrt r)a\).
This proves \eqref{eq:f2v86-last-zero}.  Dividing the braces by \(a\)
and using \(1-(1-z)^r=rz+O(z^2)\) for fixed \(r\) proves the limit.
\end{proof}

\begin{corollary}[Admissibility on an arithmetic progression at critical scales]
\label{cor:f2v86-critical}
Fix \(k\ge2\), \(\kappa,c,\delta,L>0\), and let
\(d=\kappa\sqrt M\), \(\varepsilon=c/\sqrt M\), \(x=d+\delta\).
For
\begin{equation}
 \kappa\sqrt M>\frac{k+1+\sqrt{k+1}}{L},
 \label{eq:f2v86-cutoff}
\end{equation}
the normalized weight \(\widehat W=W/F(x)\) has
\(\mathcal L_d^{(k)}[\widehat W]=0\),
\(\mathcal L_x^{(k)}[\widehat W]=1\), and
\(\widehat W(u)>0\) for all \(u\ge L\).  It uses just the
three exponents \(0,\varepsilon,2\varepsilon\), and its coefficient
variation has the explicit upper bound
\begin{equation}
 \sum_{j=0}^{2}|\widehat b_j|
 \le \frac{(1+2A+B)d^2(x+2\varepsilon)^{k+2}}
              {k(k+1)\delta^2\varepsilon^2}
 =O\bigl(M^{(k+6)/2}\bigr)
 \label{eq:f2v86-cost}
\end{equation}
with \(k,\kappa,c,\delta\) fixed.
\end{corollary}

\begin{proof}
Apply the two theorems and divide by the positive target response.
The norm is exactly \((1+2A+B)/F(x)\).  Equation
\eqref{eq:f2v86-target-bounds} bounds it, and substitution of the stated
scales proves the polynomial order.  Taking \(L=\log2\) implies
positivity at every prime-power logarithm, including ramified primes if
they were retained.
\end{proof}

For example, \(k=2\), \(d=6\sqrt M\), \(\varepsilon=1/\sqrt M\),
and \(\delta=1/4\) give \(U<5/d\).  For every real \(M\ge4\),
\(5/d\le5/12<1/2<\log2\); the last inequality follows by integrating
\(1/t>1/2\) on \(1<t<2\).  This is a parameter-uniform result, not
a sign test at a rounded prime coordinate.  It disproves the unrestricted
incompatibility alternative of Problem~\ref{prob:f2v85-class}.

\subsection{Scope and the next arithmetic step}

The archive already established that the three bare admissibility
constraints are compatible.  The advance here is their explicit
realization with three terms, a sharp last-zero scale, and a target-cost
bound.  No growing gamma ladder is annihilated: indeed
Theorem~\ref{thm:f2v86-response} gives \(F(y)>0\) at every other
positive real response node.  No signed remainder estimate follows from
admissibility alone.  Also, the stated critical-scale conclusion fixes
\(k\); it must not be used silently when \(k\) grows with \(M\).

V87 should use the exact integral \eqref{eq:f2v86-integral} at the
complex response nodes and determine the effect of increasing \(k\).
This tests the actual spectral response of the explicit weight, instead
of repeating either qualitative cone feasibility or unit-lattice root
searches.  GRH remains open.

\subsection*{Verification and append record}

The verifier \path{scripts/verify_grh_v86_three_term_filter.py} checks
\eqref{eq:f2v86-derivative} symbolically, verifies cancellation and
the response bounds in \(135\) exact rational cases, and checks six
normalized critical-scale instances.  It terminates with
\texttt{VERIFIED}; its SHA--256 is
\begin{center}
\texttt{F86F5B4E78A7EAB64466EF68C5D9E159E}\allowbreak
\texttt{776C72C4EBD9A83594BDB72BA3F346D}.
\end{center}
The proofs above supply the parameter-uniform assertions.  V86 appends
to V85, whose source has \(816931\) bytes and SHA--256
\begin{center}
\texttt{7FF60585A8BFD46822FF1015F0BFFE449}\allowbreak
\texttt{DC4123449F2E67A3AD3CEE3BE40ADEF}.
\end{center}
The predecessor prefix is retained byte-for-byte up to the unique final
document marker.  After validation V86 replaces V85; the frozen archive
and companion notes are unchanged.
\endgroup


\clearpage
\section{V87: the order budget and the complex response of the explicit filter}
\label{sec:f2v87-order-response}
\begingroup
\setcounter{equation}{0}
\renewcommand{\theequation}{N87.\arabic{equation}}
\renewcommand{\theHequation}{f2v87.\arabic{equation}}

V86 makes the bare admissibility constraints explicit at fixed derivative
order.  This iteration tests the two questions needed before applying
that construction to the arithmetic remainder: how far the order may
grow while preserving the positive prime tail, and what the same weight
does at the complex response nodes.  Both admit exact answers.  The
limiting response also exhibits a specific failure of absolute dominated
convergence on an infinite zero family.

\subsection{The derivative order is a time-location budget}

Retain \(r=k+1\), \(a=\varepsilon/d\), and the last zero \(U\) of
V86.  The exact formula and its proof give the two-sided estimate
\begin{equation}
 \frac{r}{d+\varepsilon}\le
 \frac{r}{\varepsilon}\log(1+\varepsilon/d)
 \le U<\frac{r+\sqrt r}{d}.
 \label{eq:f2v87-order-sandwich}
\end{equation}
The first inequality is \(\log(1+a)\ge a/(1+a)\).

\begin{theorem}[Sharp scale of the last zero at growing order]
\label{thm:f2v87-order}
Along any sequence with \(\varepsilon/d\to0\) and \(k\to\infty\),
\begin{equation}
 \frac{dU}{k+1}\longrightarrow1.
 \label{eq:f2v87-order-limit}
\end{equation}
Along any sequence with \(d\to\infty\), \(\varepsilon/d\to0\), and
\(k/d\to\tau\in[0,\infty)\), both positive zeros of \(W\)
converge to \(\tau\).  In particular, if \(\tau<L\), the weights
are positive on \([L,\infty)\) eventually.  If \(\tau>L\), their
negative interval lies inside \((L,\infty)\) eventually.
\end{theorem}

\begin{proof}
Divide \eqref{eq:f2v87-order-sandwich} by \(r/d\).  Its lower and
upper bounds tend to one when \(a\to0\) and \(r\to\infty\).
For the second assertion, the same sandwich gives \(U\to\tau\),
because \(\sqrt r/d\to0\).  If \(V<U\) is the first zero, the
quadratic root product from V86 gives
\[
 U+V=\frac{\log B}{\varepsilon}
      =\frac{r}{\varepsilon}\log(1+2\varepsilon/d)
      \longrightarrow2\tau.
\]
Thus \(V\to\tau\).  The sign pattern proved in V86 gives the last
two assertions.  They concern positivity on a continuous half-line;
failure of that stronger condition is not automatically failure on the
discrete prime-power set.
\end{proof}

For \(d=\kappa\sqrt M\), \(\varepsilon=c/\sqrt M\), and
\(k\sim\eta M\) with \(\eta>0\), the theorem yields
\begin{equation}
 U\sim\frac{\eta}{\kappa}\sqrt M.
 \label{eq:f2v87-high-order}
\end{equation}
Thus fixed order in V86 cannot be replaced by order comparable to the
growing block size.  The useful intermediate regime for this explicit
family is \(k/d\to\tau<L\).  This statement restricts the present
three-term family only, not every admissible exponential polynomial.

\subsection{An exact complex response and a uniform error bound}

Take \(x=d+\delta\), \(\delta>0\), and let \(F\) be the V86 response.
For \(\operatorname{Re}y\ge d\), define
\begin{equation}
 J_y=\frac1{\varepsilon^2}\int_0^\varepsilon\!\int_0^\varepsilon
 (1+(s+t)/d)^{k-1}(1+(s+t)/y)^{-k-2}\,ds\,dt.
 \label{eq:f2v87-J}
\end{equation}
The V86 differentiation identity, now with complex \(y\), gives
\begin{equation}
 \frac{F(y)}{F(x)}=
 \left(\frac{y-d}{\delta}\right)^2
 \left(\frac{x}{y}\right)^{k+2}\frac{J_y}{J_x}.
 \label{eq:f2v87-complex-response}
\end{equation}
At \(y=d\) this identity asserts zero on both sides; no division by
\(y-d\) is involved.

\begin{theorem}[Uniform relative control of the integral factor]
\label{thm:f2v87-uniform}
Put \(t_*=2(2k+1)\varepsilon/d\).  If \(0<t_*<1/2\), then,
uniformly on \(\operatorname{Re}y\ge d\),
\begin{equation}
 \left|\frac{J_y}{J_x}-1\right|
 \le\frac{2t_*}{1-2t_*}.
 \label{eq:f2v87-error}
\end{equation}
\end{theorem}

\begin{proof}
For \(v\in[0,2\varepsilon]\) let
\(g_y(v)=(1+v/d)^{k-1}(1+v/y)^{-k-2}\).  Its continuous logarithm,
chosen to vanish at zero, satisfies
\[
 (\log g_y)'(v)=\frac{k-1}{d+v}-\frac{k+2}{y+v},\qquad
 |(\log g_y)'(v)|\le\frac{2k+1}{d}.
\]
Here \(|y+v|\ge\operatorname{Re}(y+v)\ge d\).
Therefore \(|g_y(v)-1|\le e^{t_*}-1\).  For \(0<t_*<1\), the
power series gives \(e^{t_*}\le(1-t_*)^{-1}\), and hence
\( |J_y-1|\le t_* /(1-t_*)\).  The real number \(J_x\) is
positive and at least \((1-2t_*)/(1-t_*)\).  Applying the triangle
inequality to \((J_y-J_x)/J_x\) proves \eqref{eq:f2v87-error}.
\end{proof}

\begin{corollary}[Pointwise response at order proportional to distance]
\label{cor:f2v87-limit}
Suppose \(d\to\infty\), \(k/d\to\tau\in[0,\infty)\), and
\(k\varepsilon/d\to0\), with fixed \(\delta>0\).  For each fixed
\(w\) with \(\operatorname{Re}w\ge0\),
\begin{equation}
 \frac{F(d+w)}{F(d+\delta)}
 \longrightarrow
 G_\tau(w):=\left(\frac w\delta\right)^2
                 e^{-\tau(w-\delta)}.
 \label{eq:f2v87-limit}
\end{equation}
\end{corollary}

\begin{proof}
Since \(k\ge1\), the hypothesis implies \(t_*\to0\), so
\(J_{d+w}/J_x\to1\) by the theorem.  For fixed \(w\),
\[
 (k+2)\left[\log(1+\delta/d)-\log(1+w/d)\right]
 \longrightarrow\tau(\delta-w).
\]
Equation \eqref{eq:f2v87-complex-response} now gives the conclusion;
\(w=0\) follows directly from the exact double zero.
\end{proof}

At a zero \(\rho\) viewed from \(s=\sigma+i\gamma\), where
\(d=\sigma-1\), the corresponding \(w\) is
\(1+i\gamma-\rho\).  Thus \eqref{eq:f2v87-limit} explicitly
identifies the response that would have to be summed against the
remaining two-channel zeros.  The following obstruction prevents a
termwise absolute passage to that limit.

\begin{proposition}[No summable absolute majorant for the limiting zero family]
\label{prop:f2v87-no-dominated-limit}
Let \((w_j)\) be any infinite sequence with
\(0\le\operatorname{Re}w_j\le C\) and \(|\operatorname{Im}w_j|\to\infty\).
Then the terms \(G_\tau(w_j)\) do not tend to zero in modulus.
Consequently there is no summable sequence of nonnegative numbers
dominating all the normalized responses in
\eqref{eq:f2v87-limit}, for all sufficiently large parameters and all
these nodes.
\end{proposition}

\begin{proof}
The exact modulus obeys
\[
 |G_\tau(w_j)|
 =\frac{|w_j|^2}{\delta^2}e^{\tau\delta-\tau\operatorname{Re}w_j}
 \ge\frac{e^{\tau\delta-\tau C}}{\delta^2}|w_j|^2\longrightarrow\infty.
\]
If a summable majorant existed, taking the pointwise limit at each
fixed \(j\) would majorize these same limiting moduli, a contradiction.
\end{proof}

\subsection*{Status, verification, and next direction}

V86 supplies admissible weights without gamma-ladder cancellation.  V87
shows precisely how their derivative order consumes the allowed time
interval, and supplies the complex response with a uniform finite-parameter
error.  The resulting pointwise limit is not a convergent infinite zero
sum.  Any continuation must keep a finite-parameter tail estimate or use
a proved grouped identity before taking the limit.  The proposition does
not exclude such signed or grouped convergence.  It also supplies no upper
bound contradicting the archive's forced-compensation identity.

The next step is to bound the two-channel contribution with
\(|\operatorname{Im}w|\) comparable to the growing distance \(d\),
using \eqref{eq:f2v87-complex-response} before the limit, and determine
whether any signed tail cancellation follows from the actual arithmetic
data.  Another finite-block positivity calculation would not answer that
question.  GRH remains open.

The verifier \path{scripts/verify_grh_v87_order_and_complex_response.py}
checks \(36\) exact complex rational instances of
\eqref{eq:f2v87-error}, plus exact growing-order lower bounds.  Its
SHA--256 is
\begin{center}
\texttt{A45B20B3D68C292CDDD0E3A4EFB18CA0}\allowbreak
\texttt{856321592BBD735EF7B4E4DBE24E48B9}.
\end{center}
It terminates with \texttt{VERIFIED}; the parameter-uniform conclusions
are proved above and are not inferred from the finite checks.

V87 appends to the validated V86, whose \(826947\) bytes have SHA--256
\begin{center}
\texttt{861EDDE2BF49A177A45A4A7E3EE05580}\allowbreak
\texttt{5D438A5A1488EE555BA7A3DA31A2AF76}.
\end{center}
The complete predecessor prefix is preserved.  After validation V87 is
the sole active numeric GRH note.  The next scheduled companion sweep is
V90, and the frozen archive is unchanged.
\endgroup


\clearpage
\section{V88: the effective ordinate window is square-root sized}
\label{sec:f2v88-ordinate}
\begingroup
\setcounter{equation}{0}
\renewcommand{\theequation}{N88.\arabic{equation}}
\renewcommand{\theHequation}{f2v88.\arabic{equation}}

The V87 direction asked first about ordinates comparable to \(d\).
The finite-parameter response shows that this region is already
exponentially small when \(k\) is proportional to \(d\).  The growing
absolute mass is instead at ordinates of order \(\sqrt d\).  This
iteration proves both statements for the actual Dirichlet zero multiset,
without placing its zeros on the critical line.

Fix a primitive character \(\chi\) of conductor \(q>1\), a real
ordinate shift \(t_0\), and positive constants
\(\delta,c,\tau_0,\tau_1\).  Throughout this section,
\begin{equation}
 \varepsilon=c/d,\qquad \tau_0d\le k\le\tau_1d,
 \qquad x=d+\delta,\qquad
 R_d(w)=\frac{F(d+w)}{F(x)},
 \label{eq:f2v88-regime}
\end{equation}
where \(k\) is an integer and \(F\) is the V86 response.  Constants
below may depend on these fixed data and on \(\chi,t_0\).  No
uniformity in a growing conductor or target ordinate is asserted.

\subsection{An exact envelope before the limit}

Write \(w=b+iv\), \(0\le b\le1\).  For all sufficiently large
\(d\), the V87 bound gives \(1/2\le|J_{d+w}/J_x|\le2\).
Consequently
\begin{equation}
 \frac{v^2}{2\delta^2}
 \left(\frac{d^2}{(d+1)^2+v^2}\right)^{(k+2)/2}
 \le |R_d(b+iv)|\le
 \frac{2(1+v^2)}{\delta^2}
 \left(\frac{(d+\delta)^2}{d^2+v^2}\right)^{(k+2)/2}.
 \label{eq:f2v88-envelope}
\end{equation}
These are direct modulus inequalities in the exact V87 identity, not
approximations of the complex phase.

\begin{lemma}[Gaussian upper bound and a middle-window lower bound]
\label{lem:f2v88-window}
There are constants \(C_1,C_2>0\) such that, for large \(d\),
\begin{align}
 |R_d(b+iv)|&\le C_1(1+v^2)
       \exp\left(-\frac{\tau_0v^2}{4d}\right)
       &&(|v|\le d),\label{eq:f2v88-gaussian}\\
 |R_d(b+iv)|&\ge C_2d
       &&(\sqrt d\le |v|\le2\sqrt d).
       \label{eq:f2v88-annulus}
\end{align}
\end{lemma}

\begin{proof}
The factor \((1+\delta/d)^{k+2}\) is bounded by a fixed constant.
For \(0\le z\le1\), \(\log(1+z)\ge z/2\).  Applying this to
\(z=v^2/d^2\) in the upper envelope proves
\eqref{eq:f2v88-gaussian}.  In the specified annulus the denominator
ratio in the lower envelope is at least \((1+7/d)^{-1}\), for
\(d\ge1\).  Its power is bounded below by a positive constant since
\(k\le\tau_1d\), while \(v^2\ge d\).  This proves
\eqref{eq:f2v88-annulus}.
\end{proof}

\subsection{The zero-counting input}

Zeros are counted with multiplicity.  We use the unconditional estimate
\begin{equation}
 N_\chi(T)=\frac{T}{\pi}\log\frac{qT}{2\pi e}
                  +O_\chi(\log(2+T)),
 \qquad N_\chi(T)=\#\{\rho:0<\Re\rho<1,\ |\Im\rho|\le T\}.
 \label{eq:f2v88-count}
\end{equation}
For the stated \(q>1\) convention, a primary source with explicit
errors is Bennett, Martin, O'Bryant and Rechnitzer, \emph{Counting zeros
of Dirichlet \(L\)-functions}, Mathematics of Computation 90 (2021),
Theorem 1.1 and Corollary 1.2,
\url{https://doi.org/10.1090/mcom/3599}; author copy:
\url{https://personal.math.ubc.ca/~bennett/BeMaObRe-MathComp-2021.pdf}.
The count implies the unit-interval upper bound used in archived V82,
and also supplies the annular lower count needed here.

\begin{theorem}[Absolute mass and exponentially small far ordinates]
\label{thm:f2v88-absolute-mass}
For the family \eqref{eq:f2v88-regime},
\begin{equation}
 \sum_{\rho\in Z(\chi)}m_\rho
       |R_d(1+it_0-\rho)|\asymp d^{3/2}\log d.
 \label{eq:f2v88-total}
\end{equation}
Moreover,
\begin{equation}
 \sum_{|\Im\rho-t_0|\ge d}m_\rho
       |R_d(1+it_0-\rho)|
 \ll d^3\log d\;2^{-(k+2)/2}.
 \label{eq:f2v88-far}
\end{equation}
Each sum is finite for all sufficiently large \(d\).
\end{theorem}

\begin{proof}
Put \(v=t_0-\Im\rho\).  The zero count bounds the number of zeros in
each unit interval of \(v\) by \(O_{\chi,t_0}(\log(2+|v|))\).
Using \eqref{eq:f2v88-gaussian} and grouping these intervals gives
\[
 \sum_{|v|\le d}m_\rho|R_d(1-\Re\rho+iv)|
 \ll\sum_{j=0}^{\lceil d\rceil}
       (1+j)^2 e^{-c_0j^2/d}\log(2+j)
 \ll d^{3/2}\log d.
\]
Changing a fixed number of small intervals only changes the constant.
The final estimate follows by integral comparison after the substitution
\(j=\sqrt d\,z\).

For the far part put \(p=(k+2)/2\) and partition into
\(2^jd\le |v|<2^{j+1}d\), \(j\ge0\).  The upper envelope and the
zero count bound its \(j\)-th contribution by
\[
 C d^3\,2^{3j}(1+4^j)^{-p}\log(2+2^jd)
 \le C d^3\,2^{-p}2^{j(3-p)}(\log d+j+1),
\]
since \(1+4^j\ge2^{j+1}\).  For \(p\ge4\) this is summable with
a uniform geometric bound, proving \eqref{eq:f2v88-far} and the upper
bound in \eqref{eq:f2v88-total}.

Finally \eqref{eq:f2v88-count} shows that the annulus
\(5\sqrt d/4\le|\Im\rho|\le7\sqrt d/4\) contains
\(\gg_\chi\sqrt d\log d\) zeros for large \(d\).  Since \(t_0\)
is fixed, all of these have \(\sqrt d\le|v|\le2\sqrt d\)
eventually.  Equation \eqref{eq:f2v88-annulus} gives a contribution
\(\gg d\) from each zero, proving the lower bound.
\end{proof}

\subsection*{What the estimate resolves}

The remote part \(|\Im\rho-t_0|\ge d\) can be controlled in absolute
value without losing the target scale.  It is the window of width
\(\sqrt d\) that requires signed information: its absolute mass
diverges like \(d^{3/2}\log d\).  This is a property of the actual
zero count and the new explicit response, not a freely chosen masking
example.  The theorem supplies no sign for that mass.  The zeta channel
has the analogous conclusion from its usual zero-counting formula, but
the cited conductor-\(q>1\) theorem is not silently applied to it.

V89 should use the exact aggregate Euler identity to determine the signed
compensation size, including the archimedean terms, rather than take
absolute values of the middle window.  This does not authorize passing
the V87 pointwise limit through the zero sum.

\subsection*{Verification and append record}

\path{scripts/verify_grh_v88_ordinate_envelope.py} checks the squared
upper and lower envelopes at \(45\) exact complex rational points and
the dyadic inequality.  It terminates with \texttt{VERIFIED}.  The
infinite-sum assertions are proved above using the stated external
zero-counting theorem; finite checks are not evidence for those counts.

V88 appends to V87 with SHA--256
\begin{center}
\texttt{E9158DAD1AA5D1F2281A098939185A244}\allowbreak
\texttt{7D94F8E393CFAC0D0151B62B19D95E2}.
\end{center}
Its predecessor prefix is preserved before the final document marker.
The next scheduled companion sweep is V90.  GRH remains open.
\endgroup


\clearpage
\section{V89: exact signed compensation and its first-prime-power rate}
\label{sec:f2v89-compensation}
\begingroup
\setcounter{equation}{0}
\renewcommand{\theequation}{N89.\arabic{equation}}
\renewcommand{\theHequation}{f2v89.\arabic{equation}}

V88 located the large absolute zero mass but did not determine its sign.
Here the full Euler identity determines the signed remainder, with all
completed contributions retained.  For the explicit V86 weight at an
admissible growing order, the compensation exceeds the selected-zero
response by a strictly positive amount whose exact exponential rate is
set by the first nonzero prime-power coefficient.  This is more specific
than the archive's qualitative forced-compensation inequality; it does
not turn that necessary identity into a contradiction.

Fix a character \(\chi\) modulo \(f\), a real phase ordinate \(\gamma\),
and \(\delta,c>0\).  Let \(\ell\) be the least prime not dividing
\(f\), and put \(L_f=\log\ell\).  In this section
\begin{equation}
 d\longrightarrow\infty,\qquad \varepsilon=c/d,\qquad
 k/d\longrightarrow\tau\in(0,L_f),\qquad x=d+\delta.
 \label{eq:f2v89-regime}
\end{equation}
The integer \(k\) is at least two eventually.  Write
\(T_d=F(x)>0\) and \(\widehat W_d=W/T_d\).  V87 proves positivity
on the whole arithmetic half-line for all sufficiently large \(d\).

\subsection{The first nonzero arithmetic frequency}

For \(n=p^m\), \(p\nmid f\), define
\begin{equation}
 a(n)=1+\Re\bigl(\chi(n)e^{-i\gamma\log n}\bigr)\in[0,2],
 \qquad n_*:=\min\{p^m:p\nmid f,\ a(p^m)>0\},\quad L_*:=\log n_*.
 \label{eq:f2v89-support}
\end{equation}
This set is nonempty.  Indeed, for any unramified prime \(p\), if
\(a(p)=0\) then \(\chi(p)e^{-i\gamma\log p}=-1\), so
\(a(p^2)=2\).  In particular,
\begin{equation}
 \ell\le n_*\le\ell^2,\qquad L_f\le L_*\le2L_f.
 \label{eq:f2v89-first-range}
\end{equation}
This definition uses the exact phase; it requires no rounded test at a
prime logarithm.

Define the normalized full two-channel Euler sum by
\begin{equation}
 \mathfrak E_d=
 \frac1{\Gamma(k)}\sum_{\substack{n=p^m\\p\nmid f}}
 \Lambda(n)(\log n)^{k-1}n^{-d-1}a(n)\widehat W_d(\log n).
 \label{eq:f2v89-Euler}
\end{equation}
It converges absolutely at each finite parameter.  Every term is
nonnegative eventually, and the \(n_*\) term is positive, so
\(\mathfrak E_d>0\).

\subsection{The kernel estimates needed for the complete sum}

\begin{lemma}[Pointwise kernel limit and a global bound]
\label{lem:f2v89-kernel}
For every fixed real \(u\ge0\),
\begin{equation}
 \frac{\widehat W_d(u)}{d^{k+2}}
 \longrightarrow C_\tau(u):=
       \frac{e^{\tau\delta}(u-\tau)^2}{\tau^2\delta^2}.
 \label{eq:f2v89-kernel-limit}
\end{equation}
There is a constant \(C\), independent of \(u\ge0\) and large \(d\),
such that
\begin{equation}
 |\widehat W_d(u)|\le C d^{k+2}(1+u)^2.
 \label{eq:f2v89-kernel-bound}
\end{equation}
\end{lemma}

\begin{proof}
For fixed \(u\), put
\(h_u(v)=(1+v/d)^{k+1}e^{-vu}\).  The exact second-difference
identity is
\[
 W(u)=\int_0^\varepsilon\!\int_0^\varepsilon h_u''(s+t)\,ds\,dt,
 \qquad
 h_u''(v)=h_u(v)\left[
       \left(\frac{k+1}{d+v}-u\right)^2
       -\frac{k+1}{(d+v)^2}\right].
\]
Uniformly for \(0\le v\le2\varepsilon\), the bracket tends to
\((\tau-u)^2\) and \(h_u(v)\to1\).  Thus
\(W(u)/\varepsilon^2\to(u-\tau)^2\).  The V87 integral formula at
the real target gives
\[
 \frac{T_d}{\varepsilon^2d^{-k-2}}
 =\frac{k(k+1)}{d^2}\delta^2
        \left(\frac{d}{d+\delta}\right)^{k+2}J_x
 \longrightarrow\tau^2\delta^2e^{-\tau\delta}.
\]
Their quotient proves \eqref{eq:f2v89-kernel-limit}.

For all \(u\ge0\), \(h_u(v)\le(1+2\varepsilon/d)^{k+1}\)
and \((k+1)/d\) is bounded.  Hence
\(|h_u''(v)|\le C_1(1+u)^2\), uniformly on that interval.  The same
identity gives \(|W(u)|\le C_1\varepsilon^2(1+u)^2\).
The positive target limit just proved implies
\(T_d\ge C_2\varepsilon^2d^{-k-2}\) for large \(d\), establishing
\eqref{eq:f2v89-kernel-bound}.
\end{proof}

\begin{lemma}[Uniform arithmetic tail at a fixed logarithmic cutoff]
\label{lem:f2v89-arithmetic-tail}
For any fixed integer \(N\ge2\) with \(L=\log N>\tau\),
\begin{equation}
 \sum_{n\ge N}\frac{\Lambda(n)(\log n)^{k-1}}{\Gamma(k)}
       n^{-d-1}a(n)|\widehat W_d(\log n)|
 \ll \frac{d^{k+2}L^k}{\Gamma(k)}e^{-dL}.
 \label{eq:f2v89-tail}
\end{equation}
Here unsupported or ramified terms are interpreted as zero.
\end{lemma}

\begin{proof}
Use \(\Lambda(n)\le\log n\), \(a(n)\le2\), and
\eqref{eq:f2v89-kernel-bound}.  If \(u\ge L\), then
\[
 u^k(1+u)^2\le
 L^k(1+L)^2\exp\left[
 \left(\frac{k}{L}+\frac2{1+L}\right)(u-L)\right].
\]
Set \(p=d-k/L-2/(1+L)\), which is greater than one eventually.
Monotone integral comparison gives
\(\sum_{n\ge N}n^{-p-1}\le N^{-p-1}+N^{-p}/p\le2N^{-p}\).
Combining the exponents leaves exactly \(e^{-dL}\).  The remaining
factors give \eqref{eq:f2v89-tail}.
\end{proof}

\begin{theorem}[First-term asymptotic and exact exponential rate]
\label{thm:f2v89-first-term}
In the regime \eqref{eq:f2v89-regime},
\begin{equation}
 \mathfrak E_d\sim
 \frac{\Lambda(n_*)a(n_*)}{n_*\Gamma(k)}
 \frac{e^{\tau\delta}(L_*-\tau)^2}{\tau^2\delta^2}
 d^{k+2}L_*^{k-1}e^{-dL_*}.
 \label{eq:f2v89-asymptotic}
\end{equation}
In particular,
\begin{equation}
 \lim_{d\to\infty}\frac{\log\mathfrak E_d}{d}
 =-I_\tau(L_*),\qquad
 I_\tau(L)=L-\tau-\tau\log(L/\tau)>0\quad(L>\tau).
 \label{eq:f2v89-rate}
\end{equation}
\end{theorem}

\begin{proof}
The \(n_*\) term has the asymptotic displayed, by
\eqref{eq:f2v89-kernel-limit}; its constant is strictly positive since
\(L_*\ge L_f>\tau\).  Bound the rest by
Lemma~\ref{lem:f2v89-arithmetic-tail} with
\(L_+=\log(n_*+1)>L_*\).  Dividing that bound by the first term gives
a fixed constant times
\[
 \exp\left[-d\left\{L_+-L_*
                 -\frac{k}{d}\log(L_+/L_*)\right\}\right],
\]
which tends to zero exponentially.  Indeed
\(L_+-L_* > L_*\log(L_+/L_*)>\tau\log(L_+/L_*)\).
This proves \eqref{eq:f2v89-asymptotic}.

Integral comparison of \(\sum_{j<k}\log j\) with \(\int\log t\,dt\)
gives \(\log\Gamma(k)=k\log k-k+O(\log k)\).  Applying this to
\eqref{eq:f2v89-asymptotic} and using \(k/d\to\tau\) proves
\eqref{eq:f2v89-rate}.  Finally \(I_\tau(\tau)=0\) and
\(I_\tau'(L)=1-\tau/L>0\) for \(L>\tau\).
\end{proof}

\subsection{The completed signed remainder}

Suppose now that \(\rho=\beta+i\gamma\) is a zero of \(L(s,\chi)\)
with \(1/2<\beta<1\), of multiplicity \(m_\rho\), and take
\(\delta=1-\beta\).  Use the archive's exact completed remainder
\(\mathfrak R_{\star,k}^{\widehat W_d}(\rho)\), including all its
archimedean and endpoint terms.

\begin{corollary}[Strict saturation from above]
\label{cor:f2v89-saturation}
For all sufficiently large \(d\),
\begin{equation}
 \mathfrak R_{\star,k}^{\widehat W_d}(\rho)-m_\rho
 =\mathfrak E_d>0,
 \qquad
 \frac1d\log\left(
 \mathfrak R_{\star,k}^{\widehat W_d}(\rho)-m_\rho\right)
 \longrightarrow-I_\tau(L_*).
 \label{eq:f2v89-saturation}
\end{equation}
\end{corollary}

\begin{proof}
The archive's (V80.12) gives
\(\mathfrak S_\star^{(1)}=-m_\rho+\mathfrak R_\star+\mathfrak H_\star\),
since the normalized target is one.  Its (V79.7) identifies
\(-\mathfrak H_\star\) with the positive repeated-prime part.  Hence
\(\mathfrak S_\star^{(1)}-\mathfrak H_\star=\mathfrak E_d\), exactly
with the factor \(\Gamma(k)^{-1}\) in \eqref{eq:f2v89-Euler}.
Rearrange and apply the theorem.
\end{proof}

This determines the signed completed compensation for this family more
sharply than an absolute zero bound.  The excess is exponentially small
and has the noncontradictory sign.  Individual zero and gamma parts may
still be large; neither their separate limits nor a cancellation-free
estimate has been asserted.  The formula is conditional on the selected
zero when interpreted as \eqref{eq:f2v89-saturation}, whereas the
Euler-sum theorem itself needs no zero hypothesis.

\subsection*{Verification and next step}

\path{scripts/verify_grh_v89_compensation_kernel.py} verifies the exact
second-derivative identity, the quadratic coefficients of the time and
response kernels, and the rate-function derivative.  It also checks the
illustrative real mod-four phase, for which \(n_*=5\).  It terminates
with \texttt{VERIFIED}; no numerical prime sum or zero list is used to
infer the asymptotic theorem.

V89 appends to V88 with SHA--256
\begin{center}
\texttt{FBB07E957F8DBFCDCA238BA44DF40ED7}\allowbreak
\texttt{E89D90CCBF161C2579FB0EC76EE56F73}.
\end{center}
V90 must perform the scheduled companion sweep.  Its mathematical task
should examine the first-prime boundary \(k-dL_f=O(\sqrt d)\), where
the fixed positive rate \(I_\tau(L_f)\) vanishes and the V89
asymptotic is not uniform.  This tests an actual boundary of the result,
not another finite list of filter signs.  GRH remains open.
\endgroup


\clearpage
\section{V90: companion sweep and the first-prime boundary layer}
\label{sec:f2v90-boundary}
\begingroup
\setcounter{equation}{0}
\renewcommand{\theequation}{N90.\arabic{equation}}
\renewcommand{\theHequation}{f2v90.\arabic{equation}}

V89 gives a strictly positive exponential excess when the limiting
order ratio is strictly below the first arithmetic frequency.  Its
constant and exponential rate degenerate as that ratio approaches the
first prime logarithm.  This section analyzes the resulting
square-root boundary layer directly.  It does not substitute a moving
parameter into a fixed-parameter asymptotic.  The companion sweep is
also due at this version.

\subsection{Read-only companion sweep and provenance}

The following are the inspected snapshots, not a claim that parallel
writers have stopped.  Main SM advanced during the sweep from V49 to
V50, and main YM from V47 to V48.  The newly appended sections were
read before closing this checkpoint.
\begin{itemize}
\item SM V50, 401869 bytes, SHA--256
\begin{center}
\texttt{49919AFC25EBAA98A48C6BC6EF2A5F8B}\allowbreak
\texttt{E777DD8DCFD1E2F4467D54FF96771B5D}.
\end{center}
\item Main YM V48, 660156 bytes, SHA--256
\begin{center}
\texttt{FABA6EBF3DFF78225DB8723BB5682CAA}\allowbreak
\texttt{D94062DD7A540B1DD5C7E35CDFFFC555}.
\end{center}
\item NS V26, 278283 bytes, unchanged SHA--256
\begin{center}
\texttt{82C887F8C2C69E4F28FAD7C3DC8DE5B9}\allowbreak
\texttt{099CB5A4BF431EBA1FFF3E91935978AD}.
\end{center}
\item Parallel YM V5, 54174 bytes, unchanged SHA--256
\begin{center}
\texttt{306F1BB84CFA8A8D47EA569EC17E9545F}\allowbreak
\texttt{9867C8D32B548EC9D9C444B4F3C0512}.
\end{center}
\end{itemize}

SM V42--V50 distinguishes parent from daughter counting, repairs a
missing standing-wave channel, derives exact CAR occupation factors,
and constructs a subtracted one-parent resolvent and spectral measure.
Its threshold analysis distinguishes fixed-parameter algebraic tails
from a proposed weak-coupling limit.  Those are statements about a
projected model, not an arithmetic explicit formula.  One local display
requires care: at \texttt{thm:v20v49-tails} in the inspected V50, the
atom and continuous leading term in the noncritical amplitude have no
plus sign between them.  The stated spectral decomposition and its
proof require their \emph{sum}; reading the display as a product would
incorrectly kill the continuous leading term when there is no atom.
No such formula is imported here, and the companion file was not
modified.

Main YM V42--V48 repairs tree-gauge regularity, uses exact conditioned
covariances and graph precision, proves one-step and composed-lift
estimates, and exhibits a longitudinal component lost by a recursively
projected lift.  Its full-link repair changes the map being composed.
These results do not estimate a Dirichlet zero sum.  The transferable
methodological lessons are to preserve the exact combined object and
to prove a separate theorem at a singular boundary.  Neither a YM
constant nor an SM spectral assertion is assumed below.

The GRH ten-version block has accordingly moved from the finite-filter
escape and saddle diagnostics of V81--V85 to the exact three-term
kernel of V86, the nonuniform zero-window analysis of V87--V88, and
the signed arithmetic compensation in V89.  This does not re-prove
the archived finite three-pair separation or the previously established
bare cone feasibility.  The new issue is the moving first-prime
boundary for this specific normalized kernel.

\subsection{Scaling and the two time roots}

Keep the character, phase, and positive constants \(\delta,c\) fixed
as in V89.  Write \(\ell\) for the least unramified prime and
\(L=\log\ell>0\).  For a fixed real \(s\), let
\begin{equation}
 d\to\infty,\qquad \varepsilon=c/d,\qquad
 k=dL-s\sqrt d+O(1)\in\mathbb N,\qquad x=d+\delta.
 \label{eq:f2v90-scaling}
\end{equation}
The bounded error includes any integer rounding.  All error constants
below may depend on its bound and on the fixed parameters.  Set
\[
 r=k+1,\quad a=\varepsilon/d,\quad
 A=(1+a)^r,\quad B=(1+2a)^r,\qquad
 W(u)=1-2Ae^{-\varepsilon u}+Be^{-2\varepsilon u}.
\]
As in V86--V89, let \(T_d=F(d+\delta)>0\) and
\(\widehat W_d=W/T_d\).  The two positive simple roots of \(W\)
are denoted by \(V_d<U_d\); its negative interval is exactly
\((V_d,U_d)\).

\begin{theorem}[Boundary roots and discrete arithmetic positivity]
\label{thm:f2v90-roots}
In \eqref{eq:f2v90-scaling},
\begin{equation}
 V_d=L-\frac{s+\sqrt L}{\sqrt d}+O(d^{-1}),\qquad
 U_d=L+\frac{\sqrt L-s}{\sqrt d}+O(d^{-1}).
 \label{eq:f2v90-roots}
\end{equation}
Consequently, for sufficiently large \(d\):
\begin{enumerate}
\item If \(s>\sqrt L\), then \(U_d<L\) and \(W(u)>0\) for
every \(u\ge L\).
\item If \(-\sqrt L<s<\sqrt L\), then \(W(L)<0\).  Thus
the bare prime-power positivity condition fails at \(\ell\).
\item If \(s<-\sqrt L\), then
\(L<V_d<U_d<\log(\ell+1)\).  Positivity on the whole half-line
fails, but \(W(\log n)>0\) at every unramified prime power \(n\).
\end{enumerate}
No endpoint sign conclusion is asserted at \(s=\pm\sqrt L\).
\end{theorem}

\begin{proof}
The exact quadratic formula gives
\[
 U_d=\varepsilon^{-1}\log\bigl(A+\sqrt{A^2-B}\bigr),\qquad
 V_d+U_d=r\varepsilon^{-1}\log(1+2a).
\]
Put \(D=\log(A^2/B)\).  Since \(a=c/d^2\), \(r\asymp d\),
Taylor's theorem on a fixed small positive interval gives
\[
 D=r a^2(1+O(a)),\qquad
 \sqrt{1-e^{-D}}=\sqrt r\,a(1+O(a+D)).
\]
Use
\(\log(A+\sqrt{A^2-B})=r\log(1+a)
  +\log(1+\sqrt{1-e^{-D}})\).
The quadratic error in the last logarithm divided by
\(\varepsilon\) is \(O(r a^2/\varepsilon)=O(d^{-2})\).
The other Taylor errors are no larger.  Hence
\[
 U_d=(r+\sqrt r)/d+O(d^{-2}),\qquad
 V_d=(r-\sqrt r)/d+O(d^{-2}).
\]
Now \(r/d=L-s/\sqrt d+O(d^{-1})\) and
\(\sqrt r/d=\sqrt L/\sqrt d+O(d^{-1})\), proving the expansion.
Its strict signs prove the first two assertions.  In the third case
both roots approach \(L\) from above, so the fixed integer gap
\(\log(\ell+1)-L>0\) contains their entire negative interval.
Every unramified prime power is either \(\ell\) or an integer at
least \(\ell+1\); all therefore lie outside that interval.
\end{proof}

The third branch is important: the continuous-half-line condition used
as a convenient sufficient condition earlier is not necessary for
positivity on the actual discrete arithmetic support.  The factor
\(a(\ell)\) from V89 is not involved in this bare weight statement.
If it vanishes, even a negative value of \(W(L)\) makes no contribution
to the two-channel Euler sum.

\subsection{The normalized value at the first frequency}

\begin{lemma}[Boundary kernel coefficient]
\label{lem:f2v90-kernel}
For every fixed real \(s\) in \eqref{eq:f2v90-scaling},
\begin{equation}
 \frac{dW(L)}{\varepsilon^2}\longrightarrow s^2-L,
 \qquad
 \frac{\widehat W_d(L)}{d^{k+1}}
 \longrightarrow
 C_L(s):=\frac{e^{L\delta}(s^2-L)}{L^2\delta^2}.
 \label{eq:f2v90-kernel}
\end{equation}
Moreover \(|\widehat W_d(u)|\le C d^{k+2}(1+u)^2\) for all
\(u\ge0\), uniformly for large \(d\).
\end{lemma}

\begin{proof}
Apply the V89 exact second-difference identity to
\(h_L(v)=(1+v/d)^r e^{-vL}\).  Uniformly for
\(0\le v\le2\varepsilon\),
\[
 \frac r{d+v}-L=-s d^{-1/2}+O(d^{-1}),\qquad
 \frac r{(d+v)^2}=L/d+O(d^{-3/2}),\qquad h_L(v)=1+O(d^{-1}).
\]
Thus
\(h_L''(v)=(s^2-L)/d+O(d^{-3/2})\), and its double integral
gives the first assertion.  The V89 proof of the target asymptotic
requires only \(k/d\to L>0\), not \(L<L_f\), and gives directly
\[
 T_d/(\varepsilon^2d^{-k-2})\longrightarrow
 L^2\delta^2e^{-L\delta}>0.
\]
Division proves the normalized limit.  The proof of the global kernel
bound in V89 also uses only bounded positive \(k/d\) and this positive
target limit, so the same second-derivative bound applies here.
\end{proof}

\subsection{A full Euler-sum asymptotic in the boundary layer}

Let \(a_\ell=a(\ell)\in[0,2]\), with the exact phase definition
in \eqref{eq:f2v89-support}, and retain the full sum
\(\mathfrak E_d\) in \eqref{eq:f2v89-Euler}.  Its definition
makes sense also when the weight is not nonnegative.

\begin{theorem}[First-prime boundary law]
\label{thm:f2v90-Euler}
For every fixed real \(s\),
\begin{equation}
 \lim_{d\to\infty}d^{-3/2}\mathfrak E_d
 =\frac{a_\ell}{\ell}\,
 \frac{e^{L\delta}(s^2-L)e^{-s^2/(2L)}}
      {\sqrt{2\pi}\,L^{3/2}\delta^2}.
 \label{eq:f2v90-Euler-limit}
\end{equation}
In particular, if \(a_\ell>0\) and \(|s|>\sqrt L\), then
\(\mathfrak E_d\sim C d^{3/2}\) with \(C>0\), on both
discretely admissible branches.  If \(a_\ell>0\) and
\(|s|<\sqrt L\), the analogous constant is negative, consistently
with the failed first-prime sign condition.  If \(a_\ell=0\),
there is \(\eta>0\) such that
\(0<\mathfrak E_d=O(e^{-\eta d})\) for all sufficiently large
\(d\), for any fixed \(s\).  At \(s=\pm\sqrt L\) with
\(a_\ell>0\), the theorem gives a zero rescaled limit, not a sign
or an unscaled decay assertion.
\end{theorem}

\begin{proof}
First separate the term at \(\ell\).  Put
\(L_+=\log(\ell+1)>L\).  The proof of the V89 arithmetic tail
bound applies with limiting order ratio \(L<L_+\), and yields
\[
 \left|\mathfrak E_d-
   \frac{L a_\ell}{\ell\Gamma(k)}L^{k-1}e^{-dL}
       \widehat W_d(L)\right|
 \le C\frac{d^{k+2}L_+^k}{\Gamma(k)}e^{-dL_+}.
\]
Stirling's formula
\(\Gamma(k)=\sqrt{2\pi}\,k^{k-1/2}e^{-k}(1+O(k^{-1}))\)
is used in its positive-real regime; see
\href{https://dlmf.nist.gov/5.11.E3}{NIST DLMF, equation 5.11.3}.
For \(h=k-dL=-s\sqrt d+O(1)\), Taylor's theorem gives
\[
 -dL+k+k\log(dL/k)
 =-\frac{h^2}{2dL}+O(|h|^3/d^2)
 \longrightarrow-\frac{s^2}{2L}.
\]
Consequently
\begin{equation}
 \frac{d^kL^ke^{-dL}}{\Gamma(k)}
 \sim\sqrt{\frac{dL}{2\pi}}\,e^{-s^2/(2L)}.
 \label{eq:f2v90-Stirling}
\end{equation}
The tail bound divided by this last expression is bounded by
\[
 C d^2\exp\left[-d\left\{L_+-L
                 -\frac{k}{d}\log(L_+/L)\right\}\right].
\]
The quantity in braces tends to
\(L_+-L-L\log(L_+/L)>0\).  Absorbing all polynomial factors
shows that the absolute tail is \(O(e^{-\eta d})\), after decreasing
\(\eta>0\) if necessary.  This is a bound for the whole tail, not
an exchange of a pointwise limit with an infinite sum.

For the first term, substitute Lemma~\ref{lem:f2v90-kernel} and
\eqref{eq:f2v90-Stirling}.  Explicitly its ratio to \(d^{3/2}\) is
\[
 \frac{L a_\ell}{\ell}\,
 \frac{\widehat W_d(L)}{d^{k+1}}\,
 \frac1{L\sqrt d}\,
 \frac{d^kL^ke^{-dL}}{\Gamma(k)},
\]
which tends to the displayed constant, including when it is zero.
The exponentially bounded tail proves \eqref{eq:f2v90-Euler-limit}.
The nonzero-sign assertions follow immediately.  When \(a_\ell=0\),
the first term vanishes exactly and the tail bound applies to the whole
sum.  Since \(U_d\to L<L_+\), every remaining supported arithmetic
frequency has strictly positive weight.  There is a nonzero coefficient
at or before \(\ell^2\), by V89.  Hence the full sum is strictly
positive even in the middle branch in this exceptional phase case.
\end{proof}

For a selected zero as in Corollary~\ref{cor:f2v89-saturation}, the
exact identity
\(\mathfrak R_{\star,k}^{\widehat W_d}(\rho)-m_\rho
=\mathfrak E_d\) is algebraic and remains valid here.  On the two
admissible branches with \(a_\ell>0\), its excess grows like
\(d^{3/2}\), rather than becoming small.  The middle branch's negative
excess is not a contradiction because the premise of the positivity
argument fails at the first prime.  The exceptional phase case has a
small excess of the same noncontradictory positive sign.

\subsection*{Verification, limitations, and the next upstream step}

\path{scripts/verify_grh_v90_first_prime_boundary.py} checks the exact
formal coefficient \(c^2(s^2-L)\) after putting \(t=d^{-1/2}\),
the Stirling exponential coefficient, and the normalization of the
first-prime constant.  It ends with \texttt{VERIFIED}.  The uniform
Taylor and arithmetic-tail estimates are proved above, not inferred
from this symbolic check.

This closes the noncritical first-prime boundary for the explicit
three-term family.  The critical endpoints require a finer scaling.
The next substantive question is whether one can place the last time
root \emph{exactly} at an arithmetic frequency while retaining the
double pole cancellation and normalized target.  Such a root-pinning
theorem must be proved for a genuine unbounded integer-order sequence;
an approximate root or a rounded logarithm is insufficient.  Even an
exact first-prime annihilation would only shift the first surviving
Euler frequency, not by itself contradict the completed identity.
GRH remains unresolved, and no estimate below multiplicity for an
admissible completed remainder has been obtained.

V90 appends to V89 with SHA--256
\begin{center}
\texttt{E4A87232F9E6624C54DA6EFAADED5672}\allowbreak
\texttt{120C1C9444B7A6FCE59ADB95002EB3BF}.
\end{center}
The next scheduled ten-version companion sweep is V100.  On reaching
that checkpoint, retain the completed batch as the next frozen
\texttt{\_f3} archive and begin a new active V1, following the established
rolling convention.  No companion or frozen archive is replaced here.
\endgroup

\clearpage
\section{V91: exact first-prime pinning on both boundary branches}
\label{sec:f2v91-pinning}
\begingroup
\setcounter{equation}{0}
\renewcommand{\theequation}{N91.\arabic{equation}}
\renewcommand{\theHequation}{f2v91.\arabic{equation}}

V90 left the two critical boundary endpoints unresolved.  Here the
remaining continuous parameter is chosen so that the time root equals
the first prime logarithm exactly.  This produces unbounded sequences
with integer order, not approximate prime samples.  Both sequences
retain arithmetic positivity and double pole cancellation.  Their
completed excess remains positive, with its first surviving frequency
strictly beyond the pinned prime.

Fix the V89 character and phase, \(c,\delta>0\), and
\(L=\log\ell\).  For \(r=k+1\), \(\varepsilon=c/d\), use
exactly the three-term weight, response \(F\), and time roots
\(V_d<U_d\) of V86.  Dependence of these roots on the fixed order
\(k\) is suppressed in the notation.

\begin{theorem}[Two exact integer-order pinning sequences]
\label{thm:f2v91-existence}
There are constants \(C,k_0\) such that, for every integer
\(k\ge k_0\) and each \(\sigma\in\{+1,-1\}\), one can choose
\(d=d_{k,\sigma}>0\) satisfying
\begin{equation}
 \left|d_{k,\sigma}-\frac{r+\sigma\sqrt r}{L}\right|
 \le \frac C r,
 \qquad
 \begin{cases}
 U_{d_{k,+}}=L,&\sigma=+1,\\
 V_{d_{k,-}}=L,&\sigma=-1.
 \end{cases}
 \label{eq:f2v91-pinning}
\end{equation}
At either choice,
\begin{equation}
 F(d)=F'(d)=0,\quad F(d+\delta)>0,\quad W(L)=0,
 \quad W(\log n)>0\quad
 (n>\ell\text{ an unramified prime power}).
 \label{eq:f2v91-constraints}
\end{equation}
On the plus branch \(W(u)>0\) for all \(u>L\).  On the minus
branch the only negative interval above \(L\) lies strictly between
\(L\) and \(\log(\ell+1)\).  No uniqueness of the chosen parameter
is needed or asserted.
\end{theorem}

\begin{proof}
The uniform root calculation in the proof of
Theorem~\ref{thm:f2v90-roots}, before imposing its particular boundary
scaling, gives
\[
 U_d=(r+\sqrt r)/d+O(r^{-2}),\qquad
 V_d=(r-\sqrt r)/d+O(r^{-2})
\]
uniformly when \(d/r\) ranges in a fixed compact subset of
\((0,\infty)\).  Indeed \(a=c/d^2=O(r^{-2})\), and each
discarded logarithmic error divided by \(\varepsilon\) is
\(O(r a^2/\varepsilon)=O(r^{-2})\); the square-root errors are
smaller.  This uniformity is what permits varying \(d\).

Put \(d_0=(r+\sigma\sqrt r)/L\).  At
\(d=d_0\pm C/r\),
\[
 \frac{r+\sigma\sqrt r}{d}-L
 =\mp\frac{CL^2}{r(r+\sigma\sqrt r)}+O(r^{-4}).
\]
Choosing \(C\) larger than a fixed bound for the root-error constant
divided by \(L^2\) (and then increasing \(k_0\)) makes the signs
opposite.  The exact root functions are continuous for positive \(d\),
since their discriminant is strictly positive.  The intermediate value
theorem therefore supplies a parameter in the stated interval.  This
argument uses the real number \(L=\log\ell\) itself, not a decimal
surrogate.

The V86 identities give both pole cancellations and the positive
target for every such parameter.  On the plus branch the last root is
\(L\).  On the minus branch the first root is \(L\), and
\(U_d-V_d=2\sqrt r/d+O(r^{-2})\to0\).  Hence the last root is
less than \(\log(\ell+1)\) for large \(k\).  The V86 sign pattern
and the integer gap prove all the positivity statements.
\end{proof}

\begin{lemma}[The critical order shift]
\label{lem:f2v91-shift}
For any choices in \eqref{eq:f2v91-pinning},
\begin{equation}
 k-dL=-\sigma\sqrt{dL}-\tfrac12+o(1),\qquad k/d\to L.
 \label{eq:f2v91-shift}
\end{equation}
\end{lemma}
\begin{proof}
The pinning interval gives \(dL=r+\sigma\sqrt r+O(r^{-1})\).
Thus \(\sqrt{dL}=\sqrt r+\sigma/2+O(r^{-1/2})\), and
\(k-dL=-1-\sigma\sqrt r+O(r^{-1})\).  Combining the two
equalities gives the claim.  In particular the endpoints of V90 are
reached with a specific bounded order correction, not arbitrary
rounding at a fixed \(d\).
\end{proof}

\subsection{The first survivor after exact annihilation}

Define, using the exact V89 phase coefficients,
\begin{equation}
 n_\dagger=\min\{p^m>\ell:p\nmid f,\ a(p^m)>0\},
 \qquad u_\dagger=\log n_\dagger>L.
 \label{eq:f2v91-survivor}
\end{equation}
This is a fixed finite integer independent of \(k\) and the branch.
Indeed, if \(z=\chi(\ell)e^{-i\gamma L}\), \(|z|=1\), the two
coefficients \(1+\Re z^2\) and \(1+\Re z^3\) cannot both be
zero: that would imply \(z^2=z^3=-1\), which is impossible.
Consequently \(\ell<n_\dagger\le\ell^3\).

\begin{theorem}[Exact pinning shifts, but does not reverse, compensation]
\label{thm:f2v91-compensation}
Normalize by \(\widehat W_d=W/F(d+\delta)\), and let
\(\mathfrak E_d\) be the full V89 Euler sum.  On either pinning
sequence \(\mathfrak E_d>0\), and
\begin{equation}
 \frac1d\log\mathfrak E_d\longrightarrow
 -I_L(u_\dagger)<0,\qquad
 I_L(u)=u-L-L\log(u/L).
 \label{eq:f2v91-rate}
\end{equation}
More precisely, put
\begin{equation}
 K_\dagger=
 \frac{\Lambda(n_\dagger)a(n_\dagger)}{n_\dagger}
 \frac{e^{L\delta-1/2}(u_\dagger-L)^2}
      {\sqrt{2\pi}\,L\delta^2u_\dagger^{3/2}}>0.
 \label{eq:f2v91-constant}
\end{equation}
Then
\begin{equation}
 \mathfrak E_d\sim K_\dagger d^{5/2}
 \exp\left[-dI_L(u_\dagger)
       -\sigma\sqrt{Ld}\log(u_\dagger/L)\right].
 \label{eq:f2v91-sharp}
\end{equation}
\end{theorem}

\begin{proof}
The term at \(\ell\) is exactly zero.  All other weights are
positive by Theorem~\ref{thm:f2v91-existence}, and the survivor has
a positive phase coefficient, proving strict positivity.
For each fixed \(u>L\), the V89 kernel proof with \(k/d\to L\)
gives
\[
 \widehat W_d(u)/d^{k+2}\longrightarrow
 \frac{e^{L\delta}(u-L)^2}{L^2\delta^2}>0,
 \qquad |\widehat W_d(u)|\le C d^{k+2}(1+u)^2.
\]
The arithmetic-tail argument of V89, now with cutoff
\(\log(n_\dagger+1)>u_\dagger>L\), shows that the entire tail
after \(n_\dagger\) is exponentially smaller than that term.  Thus
\[
 \mathfrak E_d\sim
 \frac{\Lambda(n_\dagger)a(n_\dagger)}{n_\dagger}
 \frac{e^{L\delta}(u_\dagger-L)^2}{L^2\delta^2}
 \frac{d^{k+2}u_\dagger^{k-1}e^{-du_\dagger}}{\Gamma(k)}.
\]
This already proves \eqref{eq:f2v91-rate} by Stirling's formula.
For the sharper statement write \(h=k-dL\).  The same positive-real
Stirling expansion used in V90 yields, at a fixed \(u>L\),
\[
 \frac{d^ku^ke^{-du}}{\Gamma(k)}
 =\sqrt{\frac{dL}{2\pi}}(1+o(1))
 \exp\left[-dI_L(u)+h\log(u/L)-\frac{h^2}{2dL}+o(1)\right].
\]
By \eqref{eq:f2v91-shift},
\(h\log(u/L)=-\sigma\sqrt{dL}\log(u/L)-\frac12\log(u/L)+o(1)\)
and \(h^2/(2dL)\to1/2\).  Multiply this expression by the
remaining factor \(d^2/u\) in the preceding first-term asymptotic.
The resulting constant is exactly \eqref{eq:f2v91-constant}.
\end{proof}

Conditional on a selected off-line zero of multiplicity \(m_\rho\),
the archive's exact completed identity still reads
\(\mathfrak R_{\star,k}^{\widehat W_d}(\rho)-m_\rho
=\mathfrak E_d\).  It therefore saturates from above at the rate just
proved.  Exact pinning removes the polynomial boundary excess of V90,
but it does not supply an upper bound below multiplicity.  None of the
other zero or gamma contributions has been dropped.

\subsection*{Verification and upstream direction}

\path{scripts/verify_grh_v91_exact_prime_pinning.py} verifies the two
formal root expansions, the critical order correction, and the sharp
first-survivor normalization.  The existence theorem is the uniform
intermediate-value argument above; no root returned by a floating-point
solver is used as a certificate.  The script reports \texttt{VERIFIED}.

Repeating one-prime tuning cannot change the sign of the exact identity.
The next useful upstream question is quantitative: can a fixed-degree,
explicit finite exponential polynomial cancel an arbitrary prescribed
finite set of arithmetic samples while controlling its normalized
coefficient cost?  The archive already proves qualitative tail-positive
spectral interpolation.  A new construction must therefore supply its
degree, exact moment formulas, and parameter-dependent normalization,
not another approximation-based existence proof.

V91 appends to V90 with SHA--256
\begin{center}
\texttt{41A01F736CFBC65CE8D50FA6423BE88D}\allowbreak
\texttt{1099E1B8D3B895F81B91018F624F3F30}.
\end{center}
The next scheduled companion sweep remains V100.  GRH is not proved.
\endgroup

\clearpage
\section{V92: prescribed arithmetic zeros with an explicit quadratic conditioning cost}
\label{sec:f2v92-finite-pins}
\begingroup
\setcounter{equation}{0}
\renewcommand{\theequation}{N92.\arabic{equation}}
\renewcommand{\theHequation}{f2v92.\arabic{equation}}

V91 shows that exact cancellation of the first prime is possible but
only shifts the positive arithmetic excess.  The archive's V76 and V81
already establish qualitative cone feasibility and finite spectral
interpolation.  The additional content here is a closed finite-moment
construction with \emph{prescribed time-side zeros}, degree exactly
\(2m+2\) for \(m\) pins, an exact target-normalization asymptotic,
and an explicit coefficient-cost product.  The lattice spacing is now
fixed, unlike \(\varepsilon=c/d\) in V86--V91.  No gamma-ladder or
other-zero interpolation is claimed.

\subsection{Construction by two weighted moment equations}

Fix \(\varepsilon>0\), \(\delta>0\), and distinct finite real
numbers \(L_1,\ldots,L_m\ge L_f>0\).  The empty list is allowed.
Take \(d\to\infty\), integer \(k\ge1\), with
\begin{equation}
 k/d\to\tau\in(0,L_f),\qquad q(u)=e^{-\varepsilon u},\quad
 q_i=e^{-\varepsilon L_i},\quad
 P(q)=\prod_{i=1}^m(q-q_i),\quad f(u)=P(q(u))^2.
 \label{eq:f2v92-data}
\end{equation}
Thus \(f\ge0\) and \(f(\tau)>0\).  Define
\begin{equation}
 M_j=\frac1{\Gamma(k)}\int_0^\infty
       u^{k-1}e^{-du}f(u)q(u)^j\,du,\qquad
 N_j=\frac1{\Gamma(k)}\int_0^\infty
       u^ke^{-du}f(u)q(u)^j\,du\quad(0\le j\le2).
 \label{eq:f2v92-moments}
\end{equation}
These are explicit finite sums.  If
\(P(q)^2=\sum_{a=0}^{2m}c_aq^a\), then
\begin{equation}
 M_j=\sum_a\frac{c_a}{(d+(a+j)\varepsilon)^k},\qquad
 N_j=k\sum_a\frac{c_a}{(d+(a+j)\varepsilon)^{k+1}}.
 \label{eq:f2v92-finite-sums}
\end{equation}
Set
\begin{equation}
 \alpha_d=\frac{N_2M_0-N_0M_2}{N_1M_0-N_0M_1},\qquad
 \beta_d=\frac{M_2-\alpha_dM_1}{M_0},\quad
 Q_d(q)=q^2-\alpha_dq-\beta_d,
 \quad W_d(u)=f(u)Q_d(q(u)).
 \label{eq:f2v92-construction}
\end{equation}
The denominator is strictly negative.  To see this, let \(\mu_d\)
be the probability measure obtained by dividing the density defining
\(M_0\) by \(M_0\).  Then the denominator equals
\(M_0^2\operatorname{Cov}_{\mu_d}(u,q(u))\).  For two independent
copies \(U,V\),
\[
 2\operatorname{Cov}_{\mu_d}(u,q(u))
 =\mathbb E[(U-V)(q(U)-q(V))]<0,
\]
because \(q\) is strictly decreasing and the measure has a density
positive except at finitely many points.  Thus the construction is
well-defined at every finite \(d>0,k\ge1\).

Let
\begin{equation}
 F_d(y)=\frac1{\Gamma(k)}\int_0^\infty
             u^{k-1}e^{-yu}W_d(u)\,du\quad(\Re y>0).
 \label{eq:f2v92-response}
\end{equation}
Here \(d\) indexes the weight; it is held fixed when differentiating
with respect to \(y\).  If \(W_d(u)=\sum_{j=0}^{2m+2}b_{j,d}
e^{-j\varepsilon u}\), then
\(F_d(y)=\sum_jb_{j,d}(y+j\varepsilon)^{-k}\).

\begin{theorem}[Exact pins, degree, and real target sign]
\label{thm:f2v92-exact}
For every finite parameter choice above,
\begin{equation}
 F_d(d)=F_d'(d)=0,\qquad
 F_d(y)>0\quad(y>0,\ y\ne d),\qquad
 W_d(L_i)=W_d'(L_i)=0.
 \label{eq:f2v92-exact}
\end{equation}
The polynomial defining \(W_d\) has degree exactly \(2m+2\),
and hence at most \(2m+3\) radial terms.  The factor \(Q_d(q(u))\)
has exactly two simple time roots \(a_d<b_d\) in \((0,\infty)\).
It is positive outside \((a_d,b_d)\) and negative inside.  Multiplying
by \(f\) introduces no additional sign changes.
\end{theorem}

\begin{proof}
The two definitions in \eqref{eq:f2v92-construction} solve
\(M_2-\alpha_dM_1-\beta_dM_0=0\) and
\(N_2-\alpha_dN_1-\beta_dN_0=0\) exactly.  These are
\(\mathbb E_{\mu_d}Q_d=\mathbb E_{\mu_d}[uQ_d]=0\), and
give both response cancellations.  The squared factors give the
prescribed time zeros, including their first derivatives; the leading
polynomial coefficient is one.

There must be at least two sign changes of \(Q_d(q(u))\) in
\((0,\infty)\).  With none, its integral against \(\mu_d\)
could not vanish.  With exactly one, say at \(t\), the function
\((u-t)Q_d(q(u))\) would have one constant nonzero sign except at
isolated points, contradicting its zero integral.  A quadratic in the
strictly monotone coordinate \(q\in(0,1)\) has at most two such
zeros.  They are therefore simple and interior.  Since its leading
coefficient is positive, its sign is positive outside the root interval.

Fix real \(y>0\), \(y\ne d\), and put \(g(u)=e^{-(y-d)u}\).
Its second derivative is strictly positive.  Let \(h\) be the affine
function agreeing with \(g\) at \(a_d,b_d\).  Strict convexity,
or the second divided-difference formula, shows that \(g-h\) has
the sign of \((u-a_d)(u-b_d)\) for all other \(u\ge0\).
It therefore has the same sign as \(Q_d(q(u))\).  Orthogonality to
\(1,u\) gives
\[
 F_d(y)=M_0\mathbb E_{\mu_d}[Q_d(q(u))(g(u)-h(u))]>0.
\]
All integrals converge because \(y>0\), and the integrand is
strictly positive off finitely many points.  This proves the real
target sign, without a limit argument.
\end{proof}

\subsection{Concentration estimates and admissibility}

Put \(v=k/d\), and let \(\nu_d\) be the probability law with
density \(d^ku^{k-1}e^{-du}/\Gamma(k)\).  Then
\(\mu_d=f\nu_d/\mathbb E_{\nu_d}f\) and
\(M_0=d^{-k}\mathbb E_{\nu_d}f\).

\begin{lemma}[Elementary weighted gamma expansion]
\label{lem:f2v92-gamma}
If \(g\) and the products needed below have bounded derivatives
through order four on \([0,\infty)\), then, uniformly for \(v\)
in a compact interval on which \(f(v)>0\),
\begin{equation}
 \mathbb E_{\mu_d}g
 =g(v)+\frac{v}{2d}
     \left(g''(v)+2\frac{f'(v)}{f(v)}g'(v)\right)+O(d^{-2}).
 \label{eq:f2v92-gamma}
\end{equation}
The formula also applies to \(g(u)=u\) and \(u q(u)^j\).
Writing \(t=u-v\), the further estimates used here are
\begin{equation}
\begin{aligned}
 \mathbb E_{\mu_d}t&=O(d^{-1}),&
 \mathbb E_{\mu_d}t^2&=v/d+O(d^{-2}),\\
 \mathbb E_{\mu_d}t^3&=O(d^{-2}),&
 \mathbb E_{\mu_d}t^4&=3v^2/d^2+O(d^{-5/2}),\\
 \mathbb E_{\mu_d}|t|^j&=O(d^{-j/2})&& (1\le j\le6).
\end{aligned}
 \label{eq:f2v92-centered}
\end{equation}
\end{lemma}

\begin{proof}
The gamma integral gives
\(\mathbb E_{\nu_d}u^j=(k)_j/d^j\).  Expanding around
\(v=k/d\) yields exactly
\[
 \mathbb E_{\nu_d}t=0,\quad
 \mathbb E_{\nu_d}t^2=v/d,\quad
 \mathbb E_{\nu_d}t^3=2v/d^2,\quad
 \mathbb E_{\nu_d}t^4=3v^2/d^2+6v/d^3,
\]
and
\(\mathbb E_{\nu_d}t^6=15v^3/d^3+130v^2/d^4+120v/d^5\).
The sixth moment and H\"older's inequality give the absolute moment
bounds through order six.  Third-order Taylor expansion with bounded
fourth derivative gives
\(\mathbb E_{\nu_d}h=h(v)+vh''(v)/(2d)+O(d^{-2})\).
Apply this to \(h=fg\) and \(h=f\) and divide; the terms
involving \(f''g\) cancel.  The exceptional stated functions grow
at most linearly and have the required bounded higher derivatives, so
the same Taylor argument applies.

For the centered estimates, multiply powers of \(t\) by the Taylor
expansion of \(f(v+t)\), use the displayed exact moments, and divide
by \(\mathbb E_{\nu_d}f=f(v)+O(d^{-1})\).  In particular the
error for the fourth moment is bounded by a constant times
\(\mathbb E_{\nu_d}|t|^5=O(d^{-5/2})\).  Since \(f\) is
bounded and its normalizing expectation is bounded below, all the
absolute moment estimates also transfer to \(\mu_d\).
\end{proof}

\begin{theorem}[The two free roots lie below all arithmetic frequencies]
\label{thm:f2v92-tail}
In the fixed-list regime \eqref{eq:f2v92-data},
\begin{equation}
 \alpha_d=2q(v)+O(d^{-1}),\qquad
 \beta_d=-q(v)^2+O(d^{-1}),\qquad a_d,b_d\to\tau.
 \label{eq:f2v92-root-limit}
\end{equation}
Consequently, for sufficiently large \(d\),
\(W_d(u)\ge0\) for every \(u\ge L_f\), with equality there
exactly at the prescribed points \(L_i\).
\end{theorem}

\begin{proof}
Apply \eqref{eq:f2v92-gamma} to \(u,q,q^2,uq,uq^2\).
Subtraction gives
\[
 \operatorname{Cov}_{\mu_d}(u,q)
 =\frac vd q'(v)+O(d^{-2}),\qquad
 \operatorname{Cov}_{\mu_d}(u,q^2)
 =\frac vd (q^2)'(v)+O(d^{-2}).
\]
Since \(q'(v)=-\varepsilon q(v)\ne0\), their quotient is
\(2q(v)+O(d^{-1})\).  The formula for \(\beta_d\) then gives
its expansion.  The two real polynomial roots already proved to exist
therefore both converge to \(q(\tau)\); taking their logarithms
gives \(a_d,b_d\to\tau<L_f\).  The exact sign pattern proves
the conclusion, including the precise equality set above \(L_f\).
\end{proof}

\subsection{Normalization and the surviving pointwise zero response}

Set \(T_d=F_d(d+\delta)>0\), \(\widehat W_d=W_d/T_d\), and
\(q_\tau=e^{-\varepsilon\tau}\).  All limits in this subsection
hold for the fixed finite list in \eqref{eq:f2v92-data}.

\begin{theorem}[Target scale and fixed complex response]
\label{thm:f2v92-normalization}
One has
\begin{equation}
 d^{k+2}T_d\longrightarrow
 A_*:=P(q_\tau)^2\varepsilon^2q_\tau^2
          \tau^2\delta^2e^{-\tau\delta}>0.
 \label{eq:f2v92-target}
\end{equation}
For every fixed \(u\ge0\),
\begin{equation}
 \frac{\widehat W_d(u)}{d^{k+2}}
 \longrightarrow
 \frac{P(q(u))^2(q(u)-q_\tau)^2}{A_*},
 \qquad |\widehat W_d(u)|\le C d^{k+2}\quad(u\ge0).
 \label{eq:f2v92-time-limit}
\end{equation}
For every fixed complex \(w\) with \(\Re w\ge0\),
\begin{equation}
 \frac{F_d(d+w)}{F_d(d+\delta)}
 \longrightarrow (w/\delta)^2e^{-\tau(w-\delta)}.
 \label{eq:f2v92-complex-limit}
\end{equation}
The latter limit is independent of the finite pin list and the fixed
lattice spacing.  It is not asserted to be uniform in \(w\).
\end{theorem}

\begin{proof}
Use \(t=u-v\) and expand the bounded function \(Q_d(q(u))\)
around \(v\).  The root-coefficient estimates give
\[
 Q_d(q(u))=C_0+C_1t+C_2t^2+O(|t|^3),\quad
 C_1=O(d^{-1}),\quad C_2=q'(v)^2+O(d^{-1}),
\]
with a global remainder bound on \(u\ge0\).  Orthogonality to
one, together with Lemma~\ref{lem:f2v92-gamma}, gives
\(C_0=-q'(v)^2v/d+O(d^{-3/2})\).  Hence
\begin{equation}
 \mathbb E_{\mu_d}[Q_d(q(u))t^2]
 =2q'(v)^2v^2/d^2+O(d^{-5/2}),\qquad
 \mathbb E_{\mu_d}[|Q_d(q(u))||t|^3]=O(d^{-5/2}).
 \label{eq:f2v92-quadratic-response}
\end{equation}
For the first equality the leading terms are
\(-q'(v)^2v^2/d^2+3q'(v)^2v^2/d^2\); every other term is
bounded using the centered moments through order five.  The absolute
bound follows from the displayed expansion and moments through order six.

For fixed \(w\), \(\Re w\ge0\), Taylor expansion of
\(e^{-wu}\) around \(v\) has global remainder at most a fixed
constant times \(|t|^3\).  The constant and linear terms integrate
to zero by the exact moment equations.  Therefore
\[
 F_d(d+w)=M_0\left[
 \frac{w^2e^{-wv}}2\mathbb E_{\mu_d}[Q_d(q(u))t^2]
              +O_w(d^{-5/2})\right].
\]
Since \(M_0=d^{-k}(f(v)+O(d^{-1}))\), substituting
\eqref{eq:f2v92-quadratic-response} proves
\[
 d^{k+2}F_d(d+w)\longrightarrow
 f(\tau)q'(\tau)^2\tau^2 w^2 e^{-\tau w}.
\]
Taking \(w=\delta\), and then dividing, proves the target and
complex-response assertions, also when \(w=0\).  The time limit
follows from coefficient convergence of \(Q_d\).  Its coefficients
are bounded and \(P\) is fixed on \([0,1]\), so \(|W_d(u)|\le C\)
uniformly; the positive target asymptotic proves the global bound.
\end{proof}

The same pointwise zero response as in V87 survives every fixed finite
set of arithmetic cancellations.  In particular its quadratic growth
on a vertical sequence is still present.  This is not a license to sum
the limits over all zeros, and it supplies no improved zero-tail bound.

\subsection{Exact asymptotic coefficient cost}

Consider the normalized absolute radial cost at the pole abscissa,
\begin{equation}
 \mathcal C_d:=\frac1{T_d}
       \sum_{j=0}^{2m+2}\frac{|b_{j,d}|}{(d+j\varepsilon)^k}.
 \label{eq:f2v92-cost}
\end{equation}
This is a cost of the finite radial combination, not a bound for the
entire completed zero and archimedean remainder.

\begin{theorem}[Quadratic cost with an explicit pin product]
\label{thm:f2v92-cost}
In the fixed-list regime,
\begin{equation}
 \frac{\mathcal C_d}{d^2}\longrightarrow
 \frac{4e^{\tau\delta}}{\varepsilon^2\tau^2\delta^2}
 \prod_{i=1}^m
 \left(\frac{1+e^{-\varepsilon(L_i-\tau)}}
            {1-e^{-\varepsilon(L_i-\tau)}}\right)^2.
 \label{eq:f2v92-cost-product}
\end{equation}
\end{theorem}

\begin{proof}
The coefficient limit is the polynomial
\(B_\infty(q)=P(q)^2(q-q_\tau)^2\).  Every root is strictly
positive, so its nonzero coefficients alternate in sign.  Since
\(d^k(d+j\varepsilon)^{-k}\to q_\tau^j\), finite summation
and continuity of absolute value give
\[
 d^k\sum_j\frac{|b_{j,d}|}{(d+j\varepsilon)^k}
 \longrightarrow B_\infty(-q_\tau)
 =4q_\tau^2P(-q_\tau)^2.
\]
Divide by \eqref{eq:f2v92-target} and expand the quotient
\(P(-q_\tau)^2/P(q_\tau)^2\) to obtain the product.
\end{proof}

\begin{corollary}[Uniform bound for the iterated finite-pin cost constants]
\label{cor:f2v92-iterated}
Fix \(0<\tau<\log2\), \(\varepsilon>1\), and \(\delta>0\).
When each pin is \(L_i=\log n_i\) for a distinct integer \(n_i\ge2\),
the limit constants in \eqref{eq:f2v92-cost-product} have a common
finite upper bound independent of the finite pin set.  In particular,
for the sets \(\{2,\ldots,N\}\),
\begin{equation}
 \sup_{N\ge2}\ \lim_{\substack{d\to\infty\\ k/d\to\tau}}
                   \frac{\mathcal C_{d,N}}{d^2}<\infty.
 \label{eq:f2v92-iterated}
\end{equation}
This is an iterated-limit statement; no common finite-\(d\) threshold
for all \(N\) is included.
\end{corollary}

\begin{proof}
Put \(t_i=e^{\varepsilon\tau}n_i^{-\varepsilon}\) and
\(\theta=e^{\varepsilon\tau}2^{-\varepsilon}<1\).  For
\(0\le t\le\theta\), integrating \(2/(1-t^2)\) gives
\[
 2\log\frac{1+t}{1-t}\le\frac{4t}{1-\theta^2}.
\]
The logarithm of the product is therefore at most
\[
 \frac{4e^{\varepsilon\tau}}{1-\theta^2}
 \sum_{n=2}^{\infty}n^{-\varepsilon}
 \le\frac{4e^{\varepsilon\tau}}{1-\theta^2}
       \left(2^{-\varepsilon}
                  +\frac{2^{1-\varepsilon}}{\varepsilon-1}\right).
\]
The final inequality is elementary decreasing-function integral
comparison.  Exponentiate and include the fixed prefactor in
\eqref{eq:f2v92-cost-product}.
\end{proof}

\subsection{What the full arithmetic identity still says}

For a fixed finite pin list, define \(n_*>1\) to be the first
unramified prime power with \(a(n_*)>0\) that is not pinned, and put
\(L_*=\log n_*\).  It exists: along powers of any unramified prime,
infinitely many of the coefficients \(1+\Re z^j\) are positive,
whereas the pin list is finite.  For example, a zero coefficient at
exponent \(j\) forces a positive coefficient at exponent \(2j\).
By Theorem~\ref{thm:f2v92-tail}, the full Euler sum built from
\(\widehat W_d\) is strictly positive for large \(d\).  Its first
surviving term has the positive constant in
\eqref{eq:f2v92-time-limit}, since \(L_*>\tau\) and it is not a pin.
The global bound there and the same integer-tail comparison as in V89
give
\begin{equation}
 \frac1d\log\mathfrak E_d\longrightarrow
 -\bigl(L_*-\tau-\tau\log(L_*/\tau)\bigr)<0.
 \label{eq:f2v92-Euler-rate}
\end{equation}
For completeness, bounding \(\Lambda(n)\le\log n\) reduces a
tail beyond \(u_+=\log(n_*+1)\) to a constant times
\(d^{k+2}u_+^ke^{-du_+}/\Gamma(k)\); its ratio to the first term
decays exponentially because
\(u_+-L_*>\tau\log(u_+/L_*)\).  This proves the first-term
dominance used in \eqref{eq:f2v92-Euler-rate}, with no interchange
of an infinite sum and a pointwise limit.

The exact completed remainder at a hypothesized target zero again
satisfies \(\mathfrak R-m_\rho=\mathfrak E_d>0\).  Thus the
construction quantitatively rules out a supposed superpolynomial
radial-cost obstruction for every fixed finite pin list at this
abscissa.  It does not rule out the hypothesized zero.  In particular
the quadratic cost \(\mathcal C_d\) still diverges, the degree
grows with the pin count, and the common iterated-limit bound is not
a controlled infinite-pin analytic kernel.

\subsection*{Verification and next quantitative question}

\path{scripts/verify_grh_v92_finite_arithmetic_pins.py} verifies exact
rational moment cancellation, double pins, real target positivity,
and free-root placement for the pin sets
\(\varnothing,\{2\},\{2,3\},\{2,3,5\}\), with
\(d=64,k=8,\varepsilon=1,\delta=1/4\).  All sign tests are at
the exact rational coordinate \(q=1/2\), not a rounded logarithm.
It checks recovery of the V86 three-term coefficients for the empty
list and the coefficient-cost product identity.  It reports
\texttt{VERIFIED}; the general asymptotic and positivity proofs are
the analytic arguments above.

The next useful target is a \emph{simultaneous} pin-count/order estimate:
quantify the errors when the pin set grows with \(d\), rather than
silently replacing \eqref{eq:f2v92-iterated} by a joint limit.
The sub-prime factor \(P(e^{-\varepsilon u})^2\) itself changes
the concentration measure as the number of pins grows, so this is a
real moving-measure issue.  Any such theorem must also keep track of
the completed zero remainder, not only its arithmetic excess.  The
present results leave GRH unresolved.

V92 appends to V91 with SHA--256
\begin{center}
\texttt{813B3290C65C5127C6040DEBACC960732}\allowbreak
\texttt{4C1753B07362C9CAA34FFC67BA1BCE4}.
\end{center}
The next scheduled companion sweep is V100.
\endgroup

\clearpage
\section{V93: a joint growing-pin theorem and a sharp superexponential Euler rate}
\label{sec:f2v93-joint-pins}
\begingroup
\setcounter{equation}{0}
\renewcommand{\theequation}{N93.\arabic{equation}}
\renewcommand{\theHequation}{f2v93.\arabic{equation}}

V92 proves a fixed-list asymptotic and an iterated bound for its cost
constants.  That does not establish a joint limit.  This section closes
that specific gap under an explicit positive pin-load allowance.  The
essential adjustment is to keep the rate shift contributed by all
squared pin factors.  A uniform sub-prime product estimate and a global
tail bound then justify the moving-measure calculation.  We obtain
linearly growing pin lists, quadratic radial cost, and the sharp
logarithmic rate of the full surviving Euler sum.  Its sign remains
positive, so the completed remainder does not yield a GRH contradiction.

\subsection{Effective rate and a quantitative pin-load condition}

Fix
\begin{equation}
 \varepsilon>1,\quad \delta>0,\quad
 0<\tau_0<\tau_1<b<\log2,
 \qquad I_v(b)=b-v-v\log(b/v).
 \label{eq:f2v93-fixed}
\end{equation}
At each parameter choose a finite set \(S\) of distinct integers
at least two, let \(m=|S|\), and set
\begin{equation}
 P_S(q)=\prod_{n\in S}(q-n^{-\varepsilon}),\quad
 q(u)=e^{-\varepsilon u},\quad
 D=d+2m\varepsilon,\quad v=k/D\in[\tau_0,\tau_1],
 \quad k\in\mathbb N.
 \label{eq:f2v93-effective}
\end{equation}
The joint regime is \(D\to\infty\), subject to
\begin{equation}
 \lambda:=\frac{2m\varepsilon}{D},\qquad
 \lambda b\le\tfrac12 I_{\tau_1}(b).
 \label{eq:f2v93-load}
\end{equation}
In particular \(0\le\lambda<1/2\), \(d=(1-\lambda)D>0\),
and \(d\asymp D\), uniformly.  Define \(Q_{d,S}\), \(W_{d,S}\),
and \(F_{d,S}\) by the exact finite-moment formulas of V92, using
the original pole parameter \(d\), this \(k\), and \(P_S\).
All their finite-parameter identities continue to hold.  Pins at
ramified or non-prime integers simply add harmless time-side zeros;
we shall prove positivity above \(\log2\), which is no larger than
the first unramified prime logarithm.

Introduce the two finite products
\begin{equation}
 H_S^\pm(u)=\prod_{n\in S}(1\pm e^{\varepsilon u}n^{-\varepsilon})^2.
 \label{eq:f2v93-products}
\end{equation}
The exact, global identities are
\begin{equation}
 P_S(q(u))^2=e^{-2m\varepsilon u}H_S^-(u),\qquad
 P_S(-q(u))^2=e^{-2m\varepsilon u}H_S^+(u).
 \label{eq:f2v93-factorization}
\end{equation}
Thus the local gamma rate is \(D\), not \(d\).  Neither identity
asserts that \(H_S^\pm\) is bounded on the whole half-line.

\subsection{Uniform local products and global tails}

\begin{lemma}[Product bounds independent of the finite set]
\label{lem:f2v93-products}
For any fixed \(b'<\log2\), there are positive constants independent
of \(S\) such that on \([0,b']\)
\begin{equation}
 0<h_*\le H_S^-(u)\le1,\qquad
 1\le H_S^+(u)\le h^*,\qquad
 |\partial_u^jH_S^\pm(u)|\le C_j\quad(0\le j\le6).
 \label{eq:f2v93-product-bounds}
\end{equation}
On the full half-line one instead has the uniform bounds
\begin{equation}
 |P_S(q(u))|^2\le1,\qquad
 |P_S(-q(u))|^2\le
 C_*:=\prod_{n=2}^{\infty}(1+n^{-\varepsilon})^2<\infty.
 \label{eq:f2v93-global-P}
\end{equation}
\end{lemma}

\begin{proof}
For \(u\le b'\), put \(t_n=e^{\varepsilon u}n^{-\varepsilon}\).
Then \(0\le t_n\le\theta=e^{\varepsilon b'}2^{-\varepsilon}<1\),
and \(\sum_{n\ge2}t_n\) is uniformly bounded by decreasing-function
integral comparison.  The inequalities
\(-\log(1-t)\le t/(1-\theta)\) and \(\log(1+t)\le t\)
give the upper and lower product bounds.  Each derivative through
order six of \(\log(1\pm t_n(u))\) is bounded by a fixed constant
times \(t_n(u)\), since \(1\pm t_n\) stays away from zero.
Summing these derivatives and differentiating the exponential proves
the derivative estimates.  Finally \(q(u),n^{-\varepsilon}\in(0,1]\)
implies \(|q(u)-n^{-\varepsilon}|\le1\) and
\(q(u)+n^{-\varepsilon}\le1+n^{-\varepsilon}\).  The latter
infinite product converges by the same logarithmic bound.
\end{proof}

\begin{lemma}[Uniform remote-mass bound]
\label{lem:f2v93-tail}
There is \(\eta>0\), depending only on the fixed parameters, such
that for each fixed nonnegative integer \(j\), uniformly under
\eqref{eq:f2v93-load},
\begin{equation}
 \frac{D^k}{\Gamma(k)}\int_b^\infty
 u^{k+j-1}e^{-du}|P_S(\pm q(u))|^2\,du
 =O_j(e^{-\eta D}).
 \label{eq:f2v93-remote}
\end{equation}
Consequently the probability measure proportional to
\(u^{k-1}e^{-du}P_S(q(u))^2\,du\) has the weighted gamma
expansion and centered-moment estimates of V92 with \(D\) in place
of \(d\) and \(H_S^-\) in place of its fixed multiplier, with
constants uniform in \(S\).
\end{lemma}

\begin{proof}
Use \eqref{eq:f2v93-global-P}.  For \(j=0\), the gamma Chernoff
bound, obtained by multiplying the density by
\(e^{(d-k/b)(u-b)}\ge1\) on \(u\ge b\), gives
\[
 \frac{D^k}{\Gamma(k)}\int_b^\infty u^{k-1}e^{-du}\,du
 \le \exp[-db+k+k\log(Db/k)]
 =\exp[-D\{I_v(b)-\lambda b\}].
\]
Here \(d-k/b>0\): \eqref{eq:f2v93-load} implies
\((1-\lambda)b>\tau_1\ge v\).  Also
\(\partial_v I_v(b)=-\log(b/v)<0\) for \(v<b\), so
\(I_v(b)-\lambda b\ge I_{\tau_1}(b)/2>0\).
For fixed \(j\), apply the same calculation with shape \(k+j\),
including the factor \(\Gamma(k+j)/\Gamma(k)\).  Its logarithm
changes the displayed exponent by at most \(O_j(\log D)\),
uniformly for \(v\in[\tau_0,\tau_1]\).  For example any fixed
\(\eta<I_{\tau_1}(b)/2\) is valid after increasing the threshold.

For the final assertion choose \(b<b'<\log2\), and a fixed smooth
cutoff equal to one on \([0,b]\) and zero beyond \(b'\).
Multiply \(H_S^-\) by this cutoff.  Its derivatives through order
six are uniformly bounded and its value near \(v\) is bounded below.
By \eqref{eq:f2v93-factorization}, the original integrals and those
with the cut-off multiplier and gamma rate \(D\) differ only above
\(b\).  The former tails satisfy \eqref{eq:f2v93-remote}; the
latter satisfy the usual gamma Chernoff bound with rate \(D\).
The same holds with any fixed polynomial factor through the moment
orders used in V92.  Their common normalizer is bounded below after
multiplication by \(D^k\).  Taylor expansion with the exact gamma
moments, followed by division, therefore gives those expansions with
uniform errors.  This cutoff argument controls the actual remote
mass; local derivative bounds alone would not suffice.
\end{proof}

\subsection{Joint normalization, positivity, and absolute cost}

Write \(Q_{d,S}(q)=q^2-\alpha q-\beta\),
\(T_{d,S}=F_{d,S}(d+\delta)>0\), and
\(W_{d,S}(u)=\sum_{j=0}^{2m+2}b_{j,d,S}e^{-j\varepsilon u}\).
Define the same original-abscissa cost as before:
\begin{equation}
 \mathcal C_{d,S}=\frac1{T_{d,S}}
     \sum_{j=0}^{2m+2}\frac{|b_{j,d,S}|}{(d+j\varepsilon)^k}.
 \label{eq:f2v93-cost-definition}
\end{equation}

\begin{theorem}[Uniform joint estimates]
\label{thm:f2v93-joint}
In the joint regime \eqref{eq:f2v93-effective}--\eqref{eq:f2v93-load},
the following estimates have constants independent of \(S\):
\begin{align}
 \alpha&=2q(v)+O(D^{-1}),&
 \beta&=-q(v)^2+O(D^{-1}),
 \label{eq:f2v93-coefficients}\\
 T_{d,S}&=D^{-k-2}H_S^-(v)\varepsilon^2q(v)^2
       v^2\delta^2e^{-v\delta}\bigl(1+O(D^{-1/2})\bigr).
 \label{eq:f2v93-target}
\end{align}
The two free time roots are \(v+O(D^{-1/2})\).  Thus for all
sufficiently large \(D\), uniformly in the allowed pin sets,
\begin{equation}
 W_{d,S}(u)\ge0\quad(u\ge\log2),
 \qquad W_{d,S}(u)=0\text{ there exactly when }u=\log n,\ n\in S.
 \label{eq:f2v93-admissibility}
\end{equation}
For every fixed complex \(w\) with \(\Re w\ge0\),
\begin{equation}
 \frac{F_{d,S}(d+w)}{F_{d,S}(d+\delta)}
 =(w/\delta)^2e^{-v(w-\delta)}+O_w(D^{-1/2}).
 \label{eq:f2v93-response}
\end{equation}
Finally,
\begin{equation}
 \frac{\mathcal C_{d,S}}{D^2}
 =\frac{4e^{v\delta}}{\varepsilon^2v^2\delta^2}
       \frac{H_S^+(v)}{H_S^-(v)}
          \bigl(1+O(D^{-1/2})\bigr).
 \label{eq:f2v93-cost}
\end{equation}
In particular \(\mathcal C_{d,S}\asymp D^2\) uniformly.
The complex-response error is uniform in the pin set, not in unbounded
\(w\).
\end{theorem}

\begin{proof}
The moment expansion in Lemma~\ref{lem:f2v93-tail} gives
\[
 \operatorname{Cov}(u,q)=\frac vD q'(v)+O(D^{-2}),\qquad
 \operatorname{Cov}(u,q^2)=\frac vD (q^2)'(v)+O(D^{-2}).
\]
Their ratio proves the \(\alpha\) estimate; its defining mean
equation proves the \(\beta\) estimate.  All constants are uniform
because \(v\) stays in the specified positive compact interval.
The exact V92 theorem places the two roots in \(q\in(0,1)\).
Their discriminant is \(O(D^{-1})\); their center is
\(q(v)+O(D^{-1})\).  Taking logarithms proves the stated free-root
estimate and then \eqref{eq:f2v93-admissibility}.

Let \(\mu\) denote the exact normalized positive measure used in
those covariances, and put \(t=u-v\).  Its total unnormalized mass is
\(M_0=D^{-k}(H_S^-(v)+O(D^{-1}))\).  The V92 quadratic-response
argument now has uniform centered moments and uniform coefficients:
\[
 \mathbb E_\mu[Q_{d,S}(q(u))t^2]
     =2q'(v)^2v^2/D^2+O(D^{-5/2}),\quad
 \mathbb E_\mu[|Q_{d,S}(q(u))||t|^3]=O(D^{-5/2}).
\]
For clarity, its local expansion has linear coefficient \(O(D^{-1})\),
quadratic coefficient \(q'(v)^2+O(D^{-1})\), and constant coefficient
\(-q'(v)^2v/D+O(D^{-3/2})\), fixed by exact orthogonality to one.
The fourth centered moment is \(3v^2/D^2+O(D^{-5/2})\); subtracting
the constant contribution leaves the factor two above.  The global
Taylor remainder is bounded by moments through order six.  These are
precisely the estimates justified uniformly by the remote-mass lemma.

Expand \(e^{-wu}\) to second order at \(v\).  The constant and
linear terms cancel by exact orthogonality to \(1,u\); for fixed
\(w\) with \(\Re w\ge0\), the global cubic remainder is bounded
by \(C_w|t|^3\).  Hence
\[
 F_{d,S}(d+w)=D^{-k-2}H_S^-(v)q'(v)^2v^2w^2e^{-vw}
                  +O_w(D^{-k-5/2}).
\]
Set \(w=\delta\), use the uniform positive lower bound for
\(H_S^-(v)\), and divide to prove the target and response formulas.

For the cost, no sum of a growing number of coefficient limits is
taken.  The polynomial \(P_S(q)^2Q_{d,S}(q)\) is monic of even
degree, and all its roots are strictly positive.  Its coefficients
therefore alternate in sign exactly.  Consequently
\begin{equation}
 \sum_j\frac{|b_{j,d,S}|}{(d+j\varepsilon)^k}
 =\frac1{\Gamma(k)}\int_0^\infty
    u^{k-1}e^{-du}P_S(-q(u))^2Q_{d,S}(-q(u))\,du.
 \label{eq:f2v93-absolute-integral}
\end{equation}
Use the plus factorization in \eqref{eq:f2v93-factorization}.
The coefficients of \(Q\) are uniformly bounded, so the remote
integral is exponentially small by Lemma~\ref{lem:f2v93-tail}.
Locally the multiplier has uniformly bounded derivatives and
\(Q_{d,S}(-q(v))=4q(v)^2+O(D^{-1})\).  The gamma expansion thus
makes the right-hand side
\(D^{-k}(4q(v)^2H_S^+(v)+O(D^{-1}))\).  Dividing by
\eqref{eq:f2v93-target} proves the cost formula.  Both product
factors and all fixed prefactors are uniformly bounded above and away
from zero, proving the final two-sided estimate.
\end{proof}

\begin{corollary}[A true growing-integer-pin limit]
\label{cor:f2v93-infinite-product}
Take \(S_N=\{2,\ldots,N\}\), with \(N\to\infty\) and the
load condition maintained, and suppose \(v\to\tau\in(\tau_0,\tau_1)\).
Then
\begin{equation}
 \frac{\mathcal C_{d,S_N}}{D^2}\longrightarrow
 \frac{4e^{\tau\delta}}{\varepsilon^2\tau^2\delta^2}
 \prod_{n=2}^{\infty}
 \left(\frac{1+e^{\varepsilon\tau}n^{-\varepsilon}}
            {1-e^{\varepsilon\tau}n^{-\varepsilon}}\right)^2<\infty.
 \label{eq:f2v93-product-limit}
\end{equation}
This is a joint limit.  It allows, for example,
\(N=\lfloor aD\rfloor\), \(k=\lfloor\tau D\rfloor\),
\(d=D-2\varepsilon(N-1)\), for any fixed
\(0<a<I_{\tau_1}(b)/(4\varepsilon b)\).
\end{corollary}
\begin{proof}
The local infinite product converges uniformly on the compact interval
of possible \(v\), by the bounds in Lemma~\ref{lem:f2v93-products}.
Apply the uniform error estimate \eqref{eq:f2v93-cost}.  The example
satisfies \eqref{eq:f2v93-load} eventually and has positive \(d\).
Each member is still a finite exponential polynomial; this does not
construct an infinite-degree Laplace kernel at fixed \(d\).
\end{proof}

\subsection{Sharp rate for the remaining full Euler sum}

Fix a character \(\chi\) modulo \(f\) and a real phase ordinate
\(\gamma\).  Let \(a(n)=1+\Re(\chi(n)e^{-i\gamma\log n})\)
on unramified prime powers, and use \(S_N=\{2,\ldots,N\}\).
With \(\widehat W=W_{d,S_N}/T_{d,S_N}\), define
\begin{equation}
 \mathfrak E_{d,N}
 =\frac1{\Gamma(k)}\sum_{\substack{n=p^j>N\\p\nmid f}}
     \Lambda(n)(\log n)^{k-1}n^{-d-1}a(n)\widehat W(\log n).
 \label{eq:f2v93-Euler}
\end{equation}
This is the full Euler sum: every term at \(n\le N\) vanishes
exactly.  For large \(D\) every remaining term is nonnegative.

\begin{theorem}[Joint superexponential scale]
\label{thm:f2v93-Euler}
Under the joint load condition, with \(N\to\infty\),
\begin{equation}
 \mathfrak E_{d,N}>0,\qquad
 \lim\frac{\log\mathfrak E_{d,N}}{D\log N}=-1.
 \label{eq:f2v93-Euler-rate}
\end{equation}
The character and phase are fixed; constants in the lower bound may
depend on them.  In particular, for the linear pin regime of
Corollary~\ref{cor:f2v93-infinite-product},
\(\mathfrak E_{d,N}=\exp[-(1+o(1))D\log D]\).
\end{theorem}

\begin{proof}
For an integer \(n>N\), every factor has the exact form
\begin{equation}
 P_{S_N}(n^{-\varepsilon})^2
 =(N!)^{-2\varepsilon}
       \prod_{j=2}^N(1-(j/n)^{\varepsilon})^2
 \le(N!)^{-2\varepsilon}.
 \label{eq:f2v93-factorial-upper}
\end{equation}
The coefficients of \(Q\) are uniformly bounded and
\(T\asymp D^{-k-2}\), by Theorem~\ref{thm:f2v93-joint}.
Put \(L_N=\log(N+1)\).  Using \(a(n)\le2\) and
\(\Lambda(n)\le\log n\), the same monotone integer-tail
comparison as in V89 gives
\begin{equation}
 \mathfrak E_{d,N}\le
 C\frac{D^{k+2}}{\Gamma(k)}(N!)^{-2\varepsilon}
                      L_N^k e^{-dL_N}.
 \label{eq:f2v93-Euler-upper}
\end{equation}
Specifically, \(u^k\le L_N^ke^{(k/L_N)(u-L_N)}\) for
\(u\ge L_N\), and \(p=d-k/L_N>1\) eventually, since
\(d\asymp D\), \(k\asymp D\), and \(L_N\to\infty\).
The bound \(\sum_{n\ge N+1}n^{-p-1}\le2(N+1)^{-p}\)
then proves the displayed inequality with a uniform constant.

For a lower bound choose any one fixed unramified prime \(\ell\),
and put \(z=\chi(\ell)e^{-i\gamma\log\ell}\), \(|z|=1\).
If \(z=1\), all its power coefficients equal two.  Otherwise
\[
 a(\ell^j)+a(\ell^{j+1})
 =2+\Re(z^j(1+z))\ge2-|1+z|>0.
\]
Thus in each pair of consecutive exponents there is a coefficient
at least \(c_z=1-|1+z|/2>0\); take \(c_z=2\) when \(z=1\).
Choose the least exponent \(j\) for which \(\ell^j\ge2N\).
One of \(n=\ell^j,\ell^{j+1}\) then satisfies
\begin{equation}
 2N\le n<2\ell^2N,\qquad a(n)\ge c_z>0.
 \label{eq:f2v93-survivor}
\end{equation}
This uses no prime-distribution theorem or approximation to the phase.

For that \(n\), every \(j/n\) in
\eqref{eq:f2v93-factorial-upper} is at most \(1/2\), so
\[
 P_{S_N}(n^{-\varepsilon})^2
 \ge(N!)^{-2\varepsilon}(1-2^{-\varepsilon})^{2(N-1)}.
\]
Also \(Q(n^{-\varepsilon})\ge c>0\) uniformly for large \(N,D\):
its constant term \(-\beta=q(v)^2+O(D^{-1})\) is bounded below,
whereas \(n^{-\varepsilon}\to0\).  Keeping this single Euler term
therefore yields
\begin{equation}
 \mathfrak E_{d,N}\ge
 c\frac{D^{k+2}}{\Gamma(k)}(N!)^{-2\varepsilon}
 (1-2^{-\varepsilon})^{2(N-1)}
       (\log n)^{k-1}n^{-d-1}.
 \label{eq:f2v93-Euler-lower}
\end{equation}
Here \(\Lambda(n)=\log\ell\) and the fixed positive phase bound
are absorbed into \(c\).  This proves strict positivity as well.

Finally \(\log n=\log N+O(1)\),
\(\log(N!)=N\log N-N+O(\log N)\), and
\(\log\Gamma(k)=k\log k-k+O(\log k)\), the latter two
estimates also following by elementary integral comparison.
Apply them to the upper and lower bounds.  Their common leading
logarithm is
\[
 -d\log N-2\varepsilon N\log N
 =-D\log N+O(\log N).
\]
Every remaining term is
\(O(D\log\log N+D+N+\log D+\log N)\), together with a
possible \(O(d/N)\) from \(\log(N+1)\).  The load condition
gives \(N=O(D)\), and all these errors are \(o(D\log N)\)
as \(D,N\to\infty\).  The two bounds prove
\eqref{eq:f2v93-Euler-rate}.
\end{proof}

Conditional on a selected off-line zero \(\rho\) of multiplicity
\(m_\rho\), with \(\delta=1-\Re\rho\) and
\(\gamma=\Im\rho\), the same exact completed identity gives
\begin{equation}
 \mathfrak R_{\star,k}^{\widehat W}(\rho)-m_\rho
 =\mathfrak E_{d,N}>0,
 \qquad
 \frac{\log(\mathfrak R_{\star,k}^{\widehat W}(\rho)-m_\rho)}
      {D\log N}\longrightarrow-1.
 \label{eq:f2v93-completed}
\end{equation}
All zero, archimedean, and endpoint terms remain part of
\(\mathfrak R\).  Many exact arithmetic cancellations and a
superexponentially small Euler sum are thus compatible with saturation
of the completed remainder from above.  They do not make that remainder
smaller than the selected multiplicity.

\subsection*{Verification and next upstream question}

\path{scripts/verify_grh_v93_joint_pin_scaling.py} checks the exact
rate factorization, the absolute-coefficient integral identity for
finite rational pin examples, the Chernoff exponent algebra, and the
consecutive-power trigonometric identity.  It terminates with
\texttt{VERIFIED}.  The uniform local-product and remote-mass estimates,
the joint normalization, and the two-sided Euler asymptotic are proved
above; finite examples are not used to infer uniformity.

This resolves the joint-limit question posed in V92 under a stated
nonzero load allowance.  There is no reason to keep cancelling finite
arithmetic prefixes as if the exact positive compensation could change
sign.  The next upstream calculation should instead track the full
zero response on its moving ordinate scale.  Formula
\eqref{eq:f2v93-response} still has the same fixed-\(w\) limit as
V87, even while the number of pins grows.  A uniform central-limit
analysis for \(\Im w\asymp\sqrt D\) can decide whether these new
jointly controlled filters actually reduce the large absolute zero
mass or merely preserve its cancellation.  No such growing-ordinate
estimate is claimed here.  GRH remains unresolved.

V93 appends to V92 with SHA--256
\begin{center}
\texttt{56ABF47B43075F7A554A4FA4A3FB72DE}\allowbreak
\texttt{91E0377ADA1AAF1AD17CCA6C798B8EA4}.
\end{center}
The next scheduled companion sweep is V100.
\endgroup

\clearpage
\section{V94: the moving zero window survives jointly growing arithmetic pinning}
\label{sec:f2v94-zero-window}
\begingroup
\setcounter{equation}{0}
\renewcommand{\theequation}{N94.\arabic{equation}}
\renewcommand{\theHequation}{f2v94.\arabic{equation}}

V93 simultaneously controls the pin count, the normalized radial cost,
and the full Euler excess.  Its response theorem still fixes the
ordinate.  Here we prove a uniform response profile at ordinates of
order \(\sqrt D\), throughout the same joint pin-load regime.
This is not a substitution of a growing parameter into its fixed-point
limit.  An exact gamma characteristic function and the uniform
remote-mass bound supply the required new estimate.  The result shows
that even the jointly controlled filters retain an absolute zero mass
of order \(D^{3/2}\log D\) in a moving annulus.

Retain all assumptions and notation of V93, including
\(D=d+2m\varepsilon\), \(v=k/D\in[\tau_0,\tau_1]\), the
pin-load condition \eqref{eq:f2v93-load}, the fixed positive \(\delta\),
and \(T_{d,S}=F_{d,S}(d+\delta)\).  Put
\begin{equation}
 R_{d,S}(w)=\frac{F_{d,S}(d+w)}{T_{d,S}}.
 \label{eq:f2v94-response}
\end{equation}
The pin set may vary with \(D\); it need not be an initial integer
interval until the concluding comparison with the Euler rate.

\subsection{Exact centered gamma transform}

Let \(U\) have the gamma probability density
\(D^ku^{k-1}e^{-Du}/\Gamma(k)\), and write \(t=U-v\).
For a real scaled frequency \(\nu\), its exact centered transform is
\begin{equation}
 \phi_D(\nu)=\mathbb E e^{-i\nu\sqrt D(U-v)}
 =e^{i\nu v\sqrt D}(1+i\nu/\sqrt D)^{-k}.
 \label{eq:f2v94-gamma-transform}
\end{equation}
The logarithm uses the branch continuous from \(\nu=0\); its
argument always has positive real part.  If \(\ell_D=\log\phi_D\),
direct differentiation gives
\begin{equation}
 \ell_D'(\nu)=-\frac{v\nu}{1+i\nu/\sqrt D},\qquad
 \ell_D''(\nu)=-\frac{v}{(1+i\nu/\sqrt D)^2}.
 \label{eq:f2v94-log-derivatives}
\end{equation}
Taylor's theorem for the logarithm, or these identities together with
\(\phi_D(0)=1\), implies uniformly on each fixed compact frequency
interval and for \(v\in[\tau_0,\tau_1]\) that
\begin{equation}
 \phi_D(\nu)=e^{-v\nu^2/2}+O(D^{-1/2}),\qquad
 -\phi_D''(\nu)-v\phi_D(\nu)
 =-v^2\nu^2e^{-v\nu^2/2}+O(D^{-1/2}).
 \label{eq:f2v94-Hermite-transform}
\end{equation}
For the second identity use
\(\phi_D''=((\ell_D')^2+\ell_D'')\phi_D\); this proves
derivative control rather than assuming that differentiating a
pointwise limit is valid.  Differentiating under the gamma integral
is justified by its finite second moment, and gives
\begin{equation}
 D\mathbb E\left[(t^2-v/D)e^{-i\nu\sqrt D t}\right]
 =-\phi_D''(\nu)-v\phi_D(\nu).
 \label{eq:f2v94-centered-quadratic}
\end{equation}

\subsection{A uniform profile on the square-root ordinate scale}

\begin{theorem}[Jointly uniform Gaussian--quadratic response]
\label{thm:f2v94-profile}
Fix finite \(B,M>0\).  Uniformly over allowed pin sets, over
\(0\le\xi\le B\), and over \(|\nu|\le M\),
\begin{equation}
 \frac{e^{i\nu v\sqrt D}}{D}
 R_{d,S}(\xi+i\nu\sqrt D)
 =-\frac{\nu^2}{\delta^2}
      \exp\left(v(\delta-\xi)-\frac{v\nu^2}{2}\right)
       +O_{B,M}(D^{-1/2}).
 \label{eq:f2v94-profile}
\end{equation}
The phase uses the actual \(v=k/D\).  Replacing it by its limit
without a rate for \(v-\tau\) would generally not be justified.
\end{theorem}

\begin{proof}
Write \(H=H_S^-\), \(Q=Q_{d,S}\), and \(q'=q'(v)\).
The uniform coefficient and centered-moment estimates from V93 imply
the global expansion on \(u\ge0\), with \(t=u-v\),
\begin{equation}
 Q(q(u))=q'^2(t^2-v/D)+E_D(u),
 \qquad
 |E_D(u)|\le C\left(D^{-3/2}+D^{-1}|t|
                         +D^{-1}t^2+|t|^3\right).
 \label{eq:f2v94-Q-expansion}
\end{equation}
Indeed its constant, linear and quadratic coefficients were computed
in the proof of \eqref{eq:f2v93-target}, and the third derivative is
uniformly bounded because \(Q\)'s coefficients are bounded and
\(q(u)\) and all its derivatives are bounded.  In the unweighted
gamma law of \eqref{eq:f2v94-gamma-transform}, the expectation of
the right-hand error is \(O(D^{-3/2})\), using its absolute moments
through order three.

Use the fixed cutoff from Lemma~\ref{lem:f2v93-tail}.  Outside it,
the original integral with the factor
\(e^{-\xi u-i\nu\sqrt D u}Q(q(u))\) has exponentially small
mass after multiplying by \(D^k\), uniformly in \(\xi\ge0\)
and real \(\nu\), because the oscillatory factor has modulus one
and \(Q\) is uniformly bounded.  The corresponding cut-off gamma
polynomial moments have exponentially small tails as well.
Within the cutoff, \(H(u)e^{-\xi u}\) has uniformly bounded
derivative, for \(0\le\xi\le B\), and therefore
\[
 H(u)e^{-\xi u}=H(v)e^{-\xi v}+O_B(|t|).
\]
Multiplying this error by \(t^2-v/D\) contributes at most
\(C_B(\mathbb E|t|^3+(v/D)\mathbb E|t|)=O_B(D^{-3/2})\).
The error \eqref{eq:f2v94-Q-expansion} obeys the same bound after
multiplication by the bounded cut-off multiplier.  No oscillatory
cancellation is needed for these errors.

Using the exact rate factorization, and restoring the phase at \(v\),
we obtain uniformly for the stated parameter ranges
\[
 D^{k+1}e^{i\nu v\sqrt D}
 F_{d,S}(d+\xi+i\nu\sqrt D)
 =H(v)e^{-\xi v}q'^2
       D\mathbb E[(t^2-v/D)e^{-i\nu\sqrt D t}]
       +O_{B,M}(D^{-1/2}).
\]
Apply \eqref{eq:f2v94-Hermite-transform} and divide by the uniformly
positive target asymptotic
\[
 D^{k+2}T_{d,S}=H(v)q'^2v^2\delta^2e^{-v\delta}
                       (1+O(D^{-1/2})).
\]
All factors in the denominator are bounded away from zero under the
fixed parameter assumptions.  The resulting identity is
\eqref{eq:f2v94-profile}.
\end{proof}

\begin{corollary}[Uniform response size on any fixed scaled annulus]
\label{cor:f2v94-annulus}
For fixed \(0<a<M<\infty\) and \(B>0\), there are constants
\(c,C>0\), independent of the allowed pin set, such that for large
\(D\)
\begin{equation}
 cD\le |R_{d,S}(\xi+iy)|\le CD
 \qquad(0\le\xi\le B,\ a\sqrt D\le|y|\le M\sqrt D).
 \label{eq:f2v94-annulus}
\end{equation}
\end{corollary}
\begin{proof}
The modulus of the leading profile in
\eqref{eq:f2v94-profile} is bounded above and away from zero on the
specified compact ranges, since \(|\nu|\ge a>0\).  Absorb the
uniform error for large \(D\).
\end{proof}

\subsection{The actual Dirichlet zero count forces large absolute mass}

Fix a primitive Dirichlet character \(\chi\) of conductor \(q_\chi>1\),
and a fixed real center \(t_0\).  Let \(\mathcal Z_\chi\) denote
the set of distinct zeros in \(0<\Re\rho<1\), with multiplicities
\(m_\rho\).  As in V88, the unconditional input is
\begin{equation}
 N_\chi(T)=\sum_{\substack{\rho\in\mathcal Z_\chi\\|\Im\rho|\le T}}
                  m_\rho
 =\frac{T}{\pi}\log\frac{q_\chi T}{2\pi e}
                     +O_\chi(\log(2+T)).
 \label{eq:f2v94-zero-count}
\end{equation}
A primary source with explicit errors for this convention is Bennett,
Martin, O'Bryant and Rechnitzer, \emph{Counting zeros of Dirichlet
\(L\)-functions}, Mathematics of Computation 90 (2021), Corollary 1.2,
\url{https://doi.org/10.1090/mcom/3599}.  This conductor-\(q_\chi>1\)
result is not being applied to the zeta channel.

\begin{theorem}[Joint arithmetic pinning does not remove the absolute zero window]
\label{thm:f2v94-zero-mass}
For every fixed \(0<a<M<\infty\), uniformly over the allowed pin sets,
\begin{equation}
 \sum_{\substack{\rho\in\mathcal Z_\chi\\
       a\sqrt D\le|\Im\rho-t_0|\le M\sqrt D}}
       m_\rho|R_{d,S}(1+it_0-\rho)|
 \asymp_{\chi,t_0,a,M}D^{3/2}\log D.
 \label{eq:f2v94-zero-mass}
\end{equation}
All fixed filter parameters may enter the implied constants.  In
particular the full absolute zero sum is at least this large.  No
matching upper bound for the full sum, or sign for the annular sum,
is asserted here.
\end{theorem}

\begin{proof}
The number of zeros in the stated shifted annulus, counted with
multiplicity, is
\[
 \frac{M-a}{2\pi}\sqrt D\log D+O_{\chi,t_0,a,M}(\sqrt D).
\]
Indeed sandwich it between the differences of \(N_\chi\) at
\(M\sqrt D\mp|t_0|\) and \(a\sqrt D\pm|t_0|\), and apply
\eqref{eq:f2v94-zero-count}.  A fixed shift changes the main term by
only \(O(\log D)\); expanding the two unshifted main terms gives
the coefficient displayed.

For each such zero put \(\xi=1-\Re\rho\in(0,1)\) and
\(y=t_0-\Im\rho\).  Corollary~\ref{cor:f2v94-annulus}, with
\(B=1\), bounds its absolute response between \(cD\) and \(CD\).
Multiply these bounds by the annular count to prove the theorem.
At each finite parameter the full absolute series also converges:
\(F_{d,S}(d+w)\) is a finite sum of inverse \(k\)-th powers
with real shifts, hence is \(O_{d,S}(|\Im w|^{-k})\) at large
ordinate; the zero count makes this summable when \(k>1\), as holds
eventually.  This last argument has no uniform coefficient constant.
\end{proof}

\subsection{What is and is not closed by this calculation}

For the initial-integer pin sets of V93, the same family satisfies
simultaneously
\begin{equation}
 \mathcal C_{d,S_N}\asymp D^2,\qquad
 \log\mathfrak E_{d,N}\sim-D\log N,\qquad
 \sum_{\text{scaled zero annulus}}m_\rho|R_{d,S_N}(1+it_0-\rho)|
       \asymp D^{3/2}\log D.
 \label{eq:f2v94-comparison}
\end{equation}
The second quantity is the logarithm of a strictly positive Euler sum.
These three conclusions concern different objects.  They do not
conflict: the completed signed remainder retains large cancellations,
and, at a hypothesized selected zero,
\(\mathfrak R-m_\rho=\mathfrak E_{d,N}>0\).

In particular there is no single summable majorant on the actual zero
set that bounds all these normalized responses: such a majorant would
give a uniform bound for their full absolute sums, contradicting
\eqref{eq:f2v94-zero-mass}.  This obstruction now covers the jointly
growing arithmetic pin lists, not just V88's three-term small-spacing
family.  Passing the fixed-\(w\) limit through the zero sum is still
invalid without a separate signed convergence theorem.

This closes the proposed use of growing arithmetic-prefix annihilation
to make the absolute zero remainder small for this explicit construction.
It does not close all possible filters or solve GRH.  Further changes
to the pin count alone would revisit the same mechanism.  The next
upstream step should return to the programme's unresolved \emph{signed}
spectral or arithmetic realization condition and compare it with these
new quantitative bounds, after checking the existing signed-reader
notes for duplication.  A useful new estimate must constrain the
completed signed trace independently of its exact Euler identity;
neither a smaller positive excess nor another absolute zero count
supplies such information.

\subsection*{Verification and append record}

\path{scripts/verify_grh_v94_moving_zero_window.py} checks the exact
centered gamma logarithmic derivatives, their quadratic Fourier factor,
and the target normalization of the limiting profile.  It reports
\texttt{VERIFIED}.  Uniform error bounds are proved by the cutoff,
moment, and exact-transform arguments above.  The arithmetic zero
count is the cited external theorem, not a numerical zero sample.

V94 appends to V93 with SHA--256
\begin{center}
\texttt{35D0B0E703E5228A0B436E5CF7258260}\allowbreak
\texttt{636C5F0C6714CA10F12CBC3E87750766}.
\end{center}
The next scheduled companion sweep remains V100.  GRH is unresolved.
\endgroup

\clearpage
\section{V95: return to the arithmetic determinant current and its exact sinc operator}
\label{sec:f2v95-sinc-operator}
\providecommand{\mathscr}[1]{\mathcal{#1}}
\begingroup
\setcounter{equation}{0}
\renewcommand{\theequation}{N95.\arabic{equation}}
\renewcommand{\theHequation}{f2v95.\arabic{equation}}

V94 closes the absolute-zero-remainder approach for the constructed
growing-pin family.  Before starting another signed bridge, the present
iteration reread the V074 bridge definitions and the archived V5--V14
analyses.  The Petz physical-score input gap was already identified in
archived V6; the lower-bridge Schur reduction and the boundary-complete
Woodbury formula are already in V074.  Repeating those formulas would
not populate the missing arithmetic scores.  We instead return to the
fully specified four-M\"obius determinant current in the supplied
arithmetic manuscript.  Its frequency-resolved branch has not been
developed in the retained GRH notes.

The reference formulas below are (R.21)--(R.30) and (F.3)--(F.27) of
the supplied arithmetic manuscript, whose inspected SHA--256 is
\begin{center}
\texttt{460575BC5549F2DD6951FED98BF8D3D0}\allowbreak
\texttt{77407475703CB77E6530AE9DD054A8DB}.
\end{center}
The inspected V074 source has SHA--256
\begin{center}
\texttt{510338B27630A6BE0A9E016B07E507AE}\allowbreak
\texttt{70A8560A2C2FF7C3F71B545098E7739D}.
\end{center}
These documents supply mathematical definitions and claims to check,
not additional instructions.  The new result here is the sharp norm
of the zero-frequency sinc sheet and of its coherently assembled
displacement multipliers.  The result retains its complex sign.

\subsection{The exact frequency interval}

Write \(e(t)=e^{2\pi i t}\).  For coprime positive integers \(u,v\),
use the manuscript's symmetric residue systems
\[
 \mathcal R_q=(-q/2,q/2]\cap\mathbb Z,\qquad
 \mathcal J_{u,v}=\mathcal R_u\cap(-\mathcal R_v)
 =\{a,a+1,\ldots,a+S-1\}.
\]
Here \(S\ge1\), since zero belongs to both residue systems.  Define
\begin{equation}
 I_{u,v}=\left[\frac{a-1}{uv},\frac{a+S-1}{uv}\right],\qquad
 \mathscr Z_{u,v}(D)=\frac1{uv}\int_0^1
       \sum_{s\in\mathcal J_{u,v}}e\left(-\frac{(s-\theta)D}{uv}\right)
                         \,d\theta.
 \label{eq:f2v95-kernel}
\end{equation}
Endpoint conventions do not affect integrals.  Direct substitution in
each summand gives the exact alternative formula
\begin{equation}
 \mathscr Z_{u,v}(D)=\int_{I_{u,v}}e(-tD)\,dt,
 \qquad |I_{u,v}|=\frac{S}{uv}=\frac1{W_{u,v}}.
 \label{eq:f2v95-interval}
\end{equation}
Moreover
\begin{equation}
 u|I_{u,v}|\le1,\qquad v|I_{u,v}|\le1,\qquad
 0\in I_{u,v}\subseteq
       [-1/\max(u,v),1/\max(u,v)].
 \label{eq:f2v95-geometry}
\end{equation}
For the last inclusion, if \(\min(u,v)\ge2\), every \(s\in\mathcal J\)
has \(|s|\le\min(u,v)/2\); the lower endpoint of its unit interval
is at worst one unit farther left.  Divide by \(uv\).  If
\(\min(u,v)=1\), then \(\mathcal J=\{0\}\) and the claim is
immediate.  The two width bounds follow from \(S\le\min(u,v)\).

In particular, convolution by \(\mathscr Z_{u,v}\) on
\(\ell^2(\mathbb Z)\) is the Fourier projection to the arc
\(I_{u,v}\) modulo one.  This does not make every physical
two-index determinant operator positive: the two dilations and the
displacement phase must still be retained.

\subsection{The rectangular operator is a scaled partial isometry}

For \(\ell\in\mathbb Z\), initially on finitely supported sequences,
define
\begin{equation}
 (T_{u,v,\ell}d)_b
 =\sum_{n\in\mathbb Z}\mathscr Z_{u,v}(ub-vn-\ell)d_n.
 \label{eq:f2v95-T}
\end{equation}
For \(r=u\) or \(v\), let
\[
 (A_rc)(t)=\sum_{j\in\mathbb Z}c_j e(rjt),\quad t\in I_{u,v},
 \qquad M_\ell h(t)=e(\ell t)h(t).
\]

\begin{theorem}[Exact dilation norms and singular values]
\label{thm:f2v95-partial-isometry}
The operators extend boundedly, and
\begin{equation}
 A_rA_r^*=r^{-1}I_{L^2(I_{u,v})},\qquad
 P_r:=rA_r^*A_r\text{ is an orthogonal projection},\qquad
 T_{u,v,\ell}=A_u^*M_\ell A_v.
 \label{eq:f2v95-coisometries}
\end{equation}
Consequently
\begin{equation}
 T_{u,v,\ell}^*T_{u,v,\ell}=\frac1{uv}P_v,\qquad
 T_{u,v,\ell}T_{u,v,\ell}^*=\frac1{uv}P_u,\qquad
 \|T_{u,v,\ell}\|=\frac1{\sqrt{uv}}.
 \label{eq:f2v95-sharp-norm}
\end{equation}
All displacement dependence is in the unitary multiplier \(M_\ell\).
In particular the norm does not decrease when \(|\ell|\) increases.
\end{theorem}

\begin{proof}
The Fourier series map from \(\ell^2(\mathbb Z)\) to
\(L^2(\mathbb R/\mathbb Z)\) is unitary.  Change variables
\(x=rt\) and use \(r|I_{u,v}|\le1\):
\(\|A_rc\|_{L^2(I)}^2\le r^{-1}\|c\|_2^2\).
To compute the adjoint norm exactly, extend \(h(x/r)\), defined
on the arc \(rI\), by zero to the circle.  This extension has squared
norm \(r\|h\|_2^2\), and its \(j\)-th Fourier coefficient is
\(r\int_I h(t)e(-rjt)\,dt=r(A_r^*h)_j\).
Parseval gives \(\|A_r^*h\|_2^2=r^{-1}\|h\|_2^2\).
Polarization proves the first identity in
\eqref{eq:f2v95-coisometries}, and then \(P_r^2=P_r=P_r^*\).
The map \(t\mapsto rt\pmod1\) is injective almost everywhere
on the interval in question, including the case of width one, so the
extension argument loses no multiplicity.

Expanding the kernel of \(A_u^*M_\ell A_v\) gives precisely
\eqref{eq:f2v95-T}.  Since \(M_\ell\) is unitary, multiplying
the operators and using the two coisometry identities gives
\eqref{eq:f2v95-sharp-norm}.  The projections are nonzero because
\(|I|>0\), so the norm equality follows.  By density, the same norm
is the supremum over all finite input and output banks.
\end{proof}

For the bilinear rather than sesquilinear convention in the source
manuscript this yields
\begin{equation}
 \left|\sum_{b,n}c_bd_n\mathscr Z_{u,v}(ub-vn-\ell)\right|
 \le\frac{\|c\|_2\|d\|_2}{\sqrt{uv}},
 \label{eq:f2v95-bilinear}
\end{equation}
by applying the Hilbert-space bound to \(\langle\overline c,Td\rangle\).
It controls the signed sinc itself, without replacing it by an absolute
strip kernel or by its square.

\subsection{Coherent displacement sums and a conductor saving}

For a finitely supported scalar sequence \(\rho\), put
\(p_\rho(t)=\sum_\ell\rho_\ell e(\ell t)\).  The same factorization
gives an exact norm identity:
\begin{equation}
 \left\|\sum_\ell\rho_\ell T_{u,v,\ell}\right\|
 =\frac1{\sqrt{uv}}\,
        \|p_\rho\|_{L^\infty(I_{u,v})}.
 \label{eq:f2v95-multiplier-norm}
\end{equation}
Indeed \(\sqrt u A_u^*\) is an isometry and
\(\sqrt v A_v\) a coisometry onto the same \(L^2(I)\) space.
Their composition with a multiplier has exactly its multiplier norm.

For \(\rho\) equal to one on an interval of \(H\) integers,
\(\|p_\rho\|_{L^\infty(I)}=H\), since zero is in the closure
of \(I\).  Thus unweighted long-displacement averaging has no
universal saving at this level.  There is, however, an arithmetic
saving when a nonprincipal conductor phase genuinely occurs in that
displacement coefficient.

\begin{lemma}[Mean-zero periodic displacement multiplier]
\label{lem:f2v95-periodic}
Let \(a_j\) have period \(f\ge2\), \(\sum_{j=1}^f a_j=0\),
and \(|a_j|\le B_0\).  For any interval \(J\) of integers,
\begin{equation}
 \sup_{|t|\le1/(2f)}
     \left|\sum_{j\in J}a_j e(jt)\right|
 \le(1+\pi)B_0f.
 \label{eq:f2v95-periodic-bound}
\end{equation}
Hence if \(\max(u,v)\ge2f\),
\begin{equation}
 \left\|\sum_{\ell\in J}a_\ell T_{u,v,\ell}\right\|
 \le\frac{(1+\pi)B_0f}{\sqrt{uv}}.
 \label{eq:f2v95-periodic-operator}
\end{equation}
The bound is independent of the length and location of \(J\).
\end{lemma}

\begin{proof}
Split \(J\) into consecutive complete blocks of length \(f\)
and one remainder of fewer than \(f\) terms.  After removing an
irrelevant initial phase, a complete block polynomial is
\(A(t)=\sum_{j=1}^f a_{j+s}e(jt)\) for a fixed shift \(s\),
and \(A(0)=0\).  Therefore
\[
 |A(t)|\le2\pi B_0|t|\sum_{j=1}^f j
          =\pi B_0f(f+1)|t|.
\]
The sum of any number of complete blocks has modulus at most
\(2|A(t)|/|1-e(ft)|\).  For \(0<|ft|\le1/2\),
\(|1-e(ft)|=2|\sin(\pi ft)|\ge4f|t|\), giving the bound
\(\pi B_0(f+1)/2\le\pi B_0f\).  At zero the complete blocks
vanish exactly.  Add the remainder bound \(B_0f\) to prove
\eqref{eq:f2v95-periodic-bound}.  Now use the interval containment
\eqref{eq:f2v95-geometry} and the exact multiplier norm.
\end{proof}

This lemma cannot be used by inserting a new character weight into the
determinant current.  The next iteration must check whether its existing
phase can be transferred to the displacement by an exact arithmetic
identity while retaining every conductor record and every completion
frequency.  In particular, the reciprocal phase in the nonzero channel
must not be dropped.  GRH is not proved by the operator norm alone.

\subsection*{Verification and append record}

\path{scripts/verify_grh_v95_sinc_operator.py} verifies the residue
interval, its width and location, and the telescoping Fourier identity
for coprime \(u,v\le11\).  Exact small finite Fourier models check
the coisometry and partial-isometry normalizations.  The continuous
operator theorem and the periodic multiplier bound are proved above;
finite checks are not their substitute.  The script reports
\texttt{VERIFIED}.

V95 appends to V94 with SHA--256
\begin{center}
\texttt{522B0FEC333F1810CD2BE04C8BBE79C7}\allowbreak
\texttt{5A615C9D2E66D8BD6FB98F5D2649E206}.
\end{center}
This is an upstream branch audit, not the scheduled ten-version companion
sweep, which remains due at V100.
\endgroup

\clearpage
\section{V96: exact conductor routing closes the large-modulus zero sheet}
\label{sec:f2v96-conductor-routing}
\begingroup
\setcounter{equation}{0}
\renewcommand{\theequation}{N96.\arabic{equation}}
\renewcommand{\theHequation}{f2v96.\arabic{equation}}

The mean-zero displacement hypothesis in V95 can be obtained from the
\emph{existing} character phase, without inserting a new weight.
On the exact determinant equation, fix the conductor residue of one
product and transfer the other character value to the displacement.
Complete only after making this exact identity.  The resulting
zero-frequency contribution with large reduced moduli is bounded
unconditionally by \(O_\chi(K\log K)\).  This closes one part of
the source-defined rectangular current, not the whole GRH criterion.

An important distinction is recorded at the outset: relocating a phase
on the determinant equation changes its extension away from that
equation.  Individual completion-frequency pieces can therefore change.
We denote the new pieces with tildes, prove their exact sum, and do not
identify them with the unmodified (F.8) pieces in the supplied manuscript.
The rearrangement and its displacement envelope are fixed explicitly
below, independently of the value of the target sum.

\subsection{The protected connected square}

Fix a primitive nonprincipal Dirichlet character \(\chi\) modulo
\(f\ge3\), and positive integer cutoffs \(K,R,Y\), with \(R\ge2\).
For application to GRH take the corner-matched parameters of (R.6):
\(R=K^{1/2-o(1)}\), \(Y=K^{1/2+o(1)}\), and \(RY<K\).
Put
\[
 W_K(p)=S_\chi(\lfloor K/p\rfloor),\qquad
 C_\chi=\max_{0\le a<f}|S_\chi(a)|,
\]
so \(|W_K(p)|\le C_\chi\) by complete-period cancellation.
Define the low-reader scalar
\begin{equation}
 Q_{\chi,K}^{\rm low}
 =\sum_{\substack{m\le R,\ n\le Y\\ mn>K/R}}
       \mu(m)\mu(n)\chi(mn)W_K(mn).
 \label{eq:f2v96-low-reader}
\end{equation}
Its complete same-product diagonal \(\mathcal D_{\chi,K}\), not
merely its literal entry diagonal, satisfies the source bound
\(\mathcal D_{\chi,K}\ll_\chi K(1+\log K)^3\) at the
corner-matched cutoffs.  Let \(\mathcal O_{\chi,K}\) be the
connected remainder, so
\begin{equation}
 |Q_{\chi,K}^{\rm low}|^2
       =\mathcal D_{\chi,K}+\mathcal O_{\chi,K}.
 \label{eq:f2v96-square}
\end{equation}
Thus \(\mathcal O\) is the sum over ordered pairs
\((m,n),(a,b)\) with \(mn\ne ab\), both products above \(K/R\),
of
\[
 \mu(m)\mu(n)\mu(a)\mu(b)
 \chi(mn)\overline{\chi(ab)}W_K(mn)\overline{W_K(ab)}.
\]

Write \(m=gv\), \(a=gu\), with \((u,v)=1\).  Nonzero terms
require \((g,uv)=1\), \((guv,f)=1\), and contribute the prefactor
\(\mu(g)^2\mu(u)\mu(v)\).  Let \(\mathcal A_R\) denote all
triples satisfying these coprimality conditions and \(gu,gv\le R\);
zero M\"obius prefactors may be left in the sum.  For each such triple
define the fully specified sequences, supported on \(1\le b,n\le Y\),
\begin{align}
 c_b^{g,u}&=\mu(b)\mathbf1_{(b,f)=1}
       \mathbf1_{gub>K/R}\,\overline{W_K(gub)},
 \label{eq:f2v96-c}\\
 d_n^{g,v,\omega}&=\mu(n)\mathbf1_{(n,f)=1}
       \mathbf1_{gvn>K/R}\,W_K(gvn)
       \mathbf1_{vn\equiv\omega\ ({\rm mod}\ f)},
 \qquad \omega\in\mathbb Z/f\mathbb Z.
 \label{eq:f2v96-d}
\end{align}
The conductor residue is retained as an explicit source label.

\begin{theorem}[Equality-preserving character relocation]
\label{thm:f2v96-relocation}
For \(ub-vn=\ell\) and \(vn\equiv\omega\pmod f\),
\begin{equation}
 \chi(vn)\overline{\chi(ub)}
       =\chi(\omega)\overline{\chi(\omega+\ell)}
       =:\rho_\omega(\ell).
 \label{eq:f2v96-phase}
\end{equation}
For each \(\omega\), the sequence \(\rho_\omega\) has period
\(f\), modulus at most one, and mean zero over a period.  Set the
predeclared displacement envelope
\begin{equation}
 J_{u,v,Y}=\{-vY,\ldots,-1\}\ \cup\ \{1,\ldots,uY\}.
 \label{eq:f2v96-envelope}
\end{equation}
Then the connected current is exactly
\begin{equation}
 \mathcal O_{\chi,K}
 =\sum_{(g,u,v)\in\mathcal A_R}\mu(g)^2\mu(u)\mu(v)
   \sum_{\omega\ ({\rm mod}\ f)}\sum_{\ell\in J_{u,v,Y}}
       \rho_\omega(\ell)
       \sum_{ub-vn=\ell}c_b^{g,u}d_n^{g,v,\omega}.
 \label{eq:f2v96-exact-current}
\end{equation}
\end{theorem}

\begin{proof}
The phase identity is just substitution of the determinant equality
followed by reduction modulo \(f\).  Periodicity is immediate, and
\(\sum_{\ell\ ({\rm mod}\ f)}\rho_\omega(\ell)
=\chi(\omega)\sum_{a\ ({\rm mod}\ f)}\overline{\chi(a)}=0\)
because the character is nonprincipal.  Every original ordered pair
has a unique gcd triple \((g,u,v)\) and a unique retained residue
\(\omega\).  Its nonzero displacement lies in
\eqref{eq:f2v96-envelope}; the value zero is excluded exactly because
the whole same-product diagonal has been removed.  Conversely every
term satisfying the determinant equality in the displayed sum gives
such an ordered pair.  The squarefree and conductor coprimality
conditions give
\(\mu(gv)\mu(gu)=\mu(g)^2\mu(u)\mu(v)\) and cancel the common
character factor.  The definitions of \(c,d\) supply all remaining
reader, unit, and cutoff conditions.  This proves the exact identity.
\end{proof}

\subsection{The new exact frequency decomposition}

Apply the equality-preserving double completion (F.6) separately to
each inner determinant sum in \eqref{eq:f2v96-exact-current}, with
the shifted Fourier transforms of these \(c,d\).  Retain all
nonzero total completion frequencies \(\nu=r+s\) in one term
\(\widetilde{\mathcal O}_{\chi,K}^{(\ne0)}\).  The zero-frequency
term is explicitly
\begin{equation}
\begin{split}
 \widetilde{\mathcal O}_{\chi,K}^{(0)}
 ={}&\sum_{(g,u,v)\in\mathcal A_R}\mu(g)^2\mu(u)\mu(v)
    \sum_{\omega\ ({\rm mod}\ f)}\sum_{\ell\in J_{u,v,Y}}
       \rho_\omega(\ell)\\
 &\hspace{1em}\times\sum_{b,n}c_b^{g,u}d_n^{g,v,\omega}
               \mathscr Z_{u,v}(ub-vn-\ell).
\end{split}
 \label{eq:f2v96-zero-sheet}
\end{equation}
Since completion is exact,
\begin{equation}
 \mathcal O_{\chi,K}
 =\widetilde{\mathcal O}_{\chi,K}^{(0)}
  +\widetilde{\mathcal O}_{\chi,K}^{(\ne0)}.
 \label{eq:f2v96-frequency-sum}
\end{equation}
The envelope includes some displacements having no exact solutions;
their complete contributions are zero, but their separate frequency
pieces need not be.  They remain in both terms.  Likewise,
\eqref{eq:f2v96-phase} is not imposed away from the determinant
equation.  These observations are why neither new term is silently
equated with its old (F.8) counterpart.

\subsection{Unconditional closure for large reduced moduli}

Split \(\widetilde{\mathcal O}^{(0)}\) according to
\(\max(u,v)\ge2f\) or \(\max(u,v)<2f\).  These are reduced
gcd moduli; they are not the physical displacement or the original
left factor \(gu\).

\begin{theorem}[Large-modulus zero-sheet bound]
\label{thm:f2v96-large-zero}
Uniformly in \(K,R,Y\),
\begin{equation}
 \left|\widetilde{\mathcal O}_{\chi,K}^{(0,\mathrm{large})}\right|
 \le 8(1+\pi)f^{3/2}C_\chi^2\,RY
                  \sum_{g=1}^{R}\frac1g.
 \label{eq:f2v96-large-bound}
\end{equation}
Consequently, at the corner-matched cutoffs,
\begin{equation}
 \widetilde{\mathcal O}_{\chi,K}^{(0,\mathrm{large})}
 =O_\chi(K\log(2K))=K^{1+o(1)}.
 \label{eq:f2v96-large-scale}
\end{equation}
No Mertens estimate or M\"obius-correlation hypothesis is used.
\end{theorem}

\begin{proof}
Fix an admissible triple with \(\max(u,v)\ge2f\).
The displacement envelope is the union of two integer intervals.
By Lemma~\ref{lem:f2v95-periodic}, its coherently weighted operator
for each conductor residue has norm at most
\(2(1+\pi)f/\sqrt{uv}\).  Thus the magnitude of this triple's
zero-sheet contribution is at most
\[
 \frac{2(1+\pi)f}{\sqrt{uv}}
       \|c^{g,u}\|_2\sum_\omega\|d^{g,v,\omega}\|_2.
\]
The \(d\) vectors have disjoint supports as \(\omega\) varies.
Since their entries, and those of \(c\), have modulus at most
\(C_\chi\),
\[
 \|c^{g,u}\|_2\le C_\chi\sqrt Y,\qquad
 \sum_\omega\|d^{g,v,\omega}\|_2
 \le\sqrt f\left(\sum_\omega\|d^{g,v,\omega}\|_2^2\right)^{1/2}
 \le C_\chi\sqrt{fY}.
\]
After this arithmetic cancellation in the displacement variable, it
is safe to take absolute values in the remaining \(g,u,v\) sums.
Drop their coprimality restrictions and use
\(\sum_{j\le H}j^{-1/2}\le2\sqrt H\).  The result is bounded by
\[
 2(1+\pi)f^{3/2}C_\chi^2Y
       \sum_{g\le R}
          \left(\sum_{u\le R/g}u^{-1/2}\right)^2
 \le8(1+\pi)f^{3/2}C_\chi^2RY\sum_{g\le R}g^{-1},
\]
as stated.  The corner-matched scale follows from \(RY<K\) and
the elementary harmonic-sum bound.
\end{proof}

The gain is due to the actual nonprincipal character after exact
conductor-record conditioning.  It is not a generic cancellation of
the sinc kernel, and it does not use a norm bound for an arbitrary
ambient prime reader.  With unweighted displacement coefficients the
V95 norm would instead grow with the displacement interval length.

\subsection{The small-modulus bank is not deleted}

Only finitely many reduced pairs satisfy \(\max(u,v)<2f\), but
their common factor \(g\) and both right-hand indices still grow.
They cannot be called a bounded finite error.  In particular the pair
\(u=v=1\) is present.  Its sinc interval is the full circle, so
\(\mathscr Z_{1,1}(D)=\mathbf1_{D=0}\) for integer \(D\).
It has no nonzero completion frequency at all.  Its exact contribution
to the new zero sheet is
\begin{equation}
 \sum_{\substack{g\le R\\(g,f)=1}}\mu(g)^2
       \left(|B_g|^2-E_g\right),
 \label{eq:f2v96-unit-pair}
\end{equation}
where
\begin{align}
 B_g&=\sum_{\substack{n\le Y\\gn>K/R}}
                 \mu(n)\chi(n)W_K(gn),
 \label{eq:f2v96-Bg}\\
 E_g&=\sum_{\substack{n\le Y\\gn>K/R}}
                 \mu(n)^2|\chi(n)|^2|W_K(gn)|^2.
 \label{eq:f2v96-Eg}
\end{align}
This follows by expanding \(|B_g|^2\) and removing \(b=n\).
The summed diagonal \(\sum_g\mu(g)^2E_g\le C_\chi^2RY\) is
already harmless.  A separate absolute bound for this particular
subbank would therefore amount to a substantive mean-square estimate
for the actual rows \(B_g\).  It is not being assumed.  Other
small-modulus subbanks may cancel with it, so imposing such a separate
bound is stronger than the remaining coherent target.

\begin{corollary}[The reduced signed target]
\label{cor:f2v96-reduced-target}
At the corner-matched cutoffs, a sufficient remaining estimate is
\begin{equation}
 \left|
 \widetilde{\mathcal O}_{\chi,K}^{(0,\mathrm{small})}
 +\widetilde{\mathcal O}_{\chi,K}^{(\ne0)}
 \right|\le K^{1+o(1)}.
 \label{eq:f2v96-residual}
\end{equation}
If it holds for every fixed primitive nonprincipal \(\chi\), the
source rectangular descent gives GRH for those Dirichlet channels.
The stronger sum of the two absolute values is sufficient but is not
required by this reduction.
\end{corollary}

\begin{proof}
Combine \eqref{eq:f2v96-frequency-sum}, the proved large-modulus
bound, and \eqref{eq:f2v96-residual}.  Equation
\eqref{eq:f2v96-square} and the complete same-product diagonal bound
give \(|Q_{\chi,K}^{\rm low}|\le K^{1/2+o(1)}\).
The supplied manuscript's high-reader estimate and exact two-threshold
identity then give the square-root Mertens bound, with its stated
reciprocal-Dirichlet-series implication.  The hypothesis
\eqref{eq:f2v96-residual} remains open; this is a conditional closure,
not a proof of that hypothesis.
\end{proof}

\subsection*{Verification and next step}

\path{scripts/verify_grh_v96_conductor_routing.py} checks the phase
relocation term by term for the primitive real characters modulo three
and four at \((K,R,Y)=(91,8,10)\) and \((150,11,12)\).
It verifies the full connected-square identity with the entire
same-product diagonal removed, the displacement envelope, the gcd
conditions, and the zero-mean conductor periods.  Both examples contain
large-reduced-modulus terms.  The script reports \texttt{VERIFIED}.
It does not identify the old and new zero channels or use a numerical
frequency integral to infer the analytic norm bound.

The next target is the remaining coherent source in
\eqref{eq:f2v96-residual}.  Before using a Kloosterman estimate, retain
the new displacement character and the \(vn\pmod f\) record in the
nonzero-frequency completion; before treating the small-modulus bank,
retain the exact \(B_g\) reader and its cross-bank terms.  The present
result is an unconditional estimate for the large-modulus part of a
new, equality-preserving zero sheet.  It neither closes the nonzero
channel nor resolves GRH.

V96 appends to V95 with SHA--256
\begin{center}
\texttt{D5EA724CF801B085733FC90E91D108CF}\allowbreak
\texttt{66EE18A13F8244B63755352C5B8A1D0E}.
\end{center}
The next scheduled companion sweep is V100.
\endgroup

\clearpage
\section{V97: the complete frequency tiling and the conductor resonances}
\label{sec:f2v97-frequency-tiling}
\begingroup
\setcounter{equation}{0}
\renewcommand{\theequation}{N97.\arabic{equation}}
\renewcommand{\theHequation}{f2v97.\arabic{equation}}

V96 bounded the large-reduced-modulus part of a new zero sheet.  To
advance on its remaining channel, this iteration first retains and
unfolds every reciprocal-residue phase in the original exact completion.
The result is a partition of the Fourier circle, not a removal of the
reciprocal phase.  It locates the conductor resonances that the V96
bound avoids.  A sharp finite-bank rank calculation then shows why an
unrestricted operator-norm argument cannot close the remaining channel.
The calculation does not replace the actual M\"obius coefficients by
an extremizer in the GRH claim.

\subsection{All completion cells, with the reciprocal phase retained}

Use the residue systems and notation of V95.  Choose an integer
\(a\) with \(ua\equiv1\pmod v\); when \(v=1\), take \(a=0\).
For \(r\in\mathcal R_v\), \(s\in\mathcal R_u\), put
\begin{equation}
 \nu=r+s,\qquad k_{r,s}=s-\nu ua,\qquad
 I_{r,s}=\left[\frac{k_{r,s}-1}{uv},\frac{k_{r,s}}{uv}\right]
                         \pmod {\mathbb Z}.
 \label{eq:f2v97-cell}
\end{equation}
Intervals on the circle are counted up to null endpoints.

\begin{theorem}[Exact completion-cell tiling]
\label{thm:f2v97-tiling}
The \(uv\) cells \(I_{r,s}\) partition \(\mathbb R/\mathbb Z\).
For finitely supported \(c_b,d_n\), the individual \((r,s)\) term
in the manuscript's completion (F.6) is exactly
\begin{equation}
 \sum_{b,n}c_bd_n\int_{I_{r,s}}e\bigl(-t(ub-vn-\ell)\bigr)\,dt.
 \label{eq:f2v97-cell-term}
\end{equation}
The union of the \(\nu=0\) cells is the arc \(I_{u,v}\) of V95.
Consequently the aggregate \(\nu\ne0\) kernel is precisely
\begin{equation}
 \mathbf1_{D=0}-\mathscr Z_{u,v}(D)
 =\int_{(\mathbb R/\mathbb Z)\setminus I_{u,v}}e(-tD)\,dt,
 \qquad D\in\mathbb Z.
 \label{eq:f2v97-complement}
\end{equation}
\end{theorem}

\begin{proof}
The two congruences
\[
 k_{r,s}\equiv-r\pmod v,\qquad k_{r,s}\equiv s\pmod u
\]
show by the Chinese remainder theorem that \(k_{r,s}\pmod {uv}\)
runs through every residue once.  This proves the tiling, including
the cases where one modulus is one.

Expand the shifted transforms in (F.3).  For fixed \(b,n,\ell\),
the phase exponent in (F.6) is
\[
 (r+\theta)b/v+(s-\theta)n/u
       -\nu\ell a/v+(s-\theta)\ell/(uv).
\]
Set \(t=(s-\theta-\nu ua)/(uv)\).  The displayed exponent is
congruent modulo integers to \(-t(ub-vn-\ell)\): the differences
in the \(b,n\) coefficients are integral because \(ua\equiv1\pmod v\).
The factor \((uv)^{-1}d\theta\), with the reversed endpoints, becomes
\(dt\) on \(I_{r,s}\).  This proves \eqref{eq:f2v97-cell-term}.
When \(\nu=0\), \(k_{r,s}=s\) and the permitted \(s\)'s are the
consecutive set \(\mathcal J_{u,v}\); their cells give exactly
\(I_{u,v}\).  Finally the integral on the whole circle is
\(\mathbf1_{D=0}\), which proves \eqref{eq:f2v97-complement}.
\end{proof}

This derivation uses the reciprocal factor to determine the cells.
It does not justify dropping that factor in a sum over varying \(u,v\).
It also applies unchanged to the relocated V96 coefficients and its
entire displacement envelope, including displacements with no solutions.

\subsection{Where the relocated displacement multiplier is large}

Fix a unit residue \(\omega\pmod f\).  Recall
\(\rho_\omega(j)=\chi(\omega)\overline{\chi(\omega+j)}\).
Define its finite Fourier coefficients by
\begin{equation}
 \alpha_{\omega,h}=\frac1f\sum_{j=0}^{f-1}
           \rho_\omega(j)e(-hj/f),\qquad 0\le h<f.
 \label{eq:f2v97-alpha}
\end{equation}
Finite Fourier inversion and Parseval give
\begin{equation}
 \alpha_{\omega,0}=0,\quad
 \rho_\omega(j)=\sum_{h=1}^{f-1}\alpha_{\omega,h}e(hj/f),\quad
 \sum_{h=1}^{f-1}|\alpha_{\omega,h}|^2=\frac{\varphi(f)}f,
 \quad \sum_h|\alpha_{\omega,h}|\le\sqrt f.
 \label{eq:f2v97-dft}
\end{equation}
No Gauss-sum magnitude theorem is needed here: these identities follow
by summing the elementary roots-of-unity orthogonality relation.

For the two-interval envelope \(J=J_{u,v,Y}\), set
\(p(t)=\sum_{\ell\in J}\rho_\omega(\ell)e(\ell t)\), and write
\(D_{[A,B]}(x)=\sum_{j=A}^B e(jx)\).  Then
\begin{equation}
 p(t)=\sum_{h=1}^{f-1}\alpha_{\omega,h}
       \bigl(D_{[-vY,-1]}(t+h/f)+D_{[1,uY]}(t+h/f)\bigr).
 \label{eq:f2v97-resonance-form}
\end{equation}
In particular, if \(\|t+h/f\|_{\mathbb R/\mathbb Z}\ge\eta>0\)
for every nonzero coefficient, then
\begin{equation}
 |p(t)|\le\frac{\sqrt f}{\eta}.
 \label{eq:f2v97-off-resonance}
\end{equation}
For every \(1\le h<f\), at its exact resonance one has
\begin{equation}
 p(-h/f)=(u+v)Y\,\alpha_{\omega,h}+E_h,
                         \qquad |E_h|\le f^{3/2}.
 \label{eq:f2v97-resonance-value}
\end{equation}
Indeed an interval sum has modulus at most
\(1/|\sin(\pi x)|\le1/(2\|x\|)\) away from an integer.
Use this for each of the two intervals.  Distinct conductor frequencies
are separated by at least \(1/f\), so the sum of their contributions
at \(-h/f\) is bounded by \(f\sum_{h'\ne h}|\alpha_{\omega,h'}|\).
The matching frequency contributes the full length \((u+v)Y\).

At least one coefficient in \eqref{eq:f2v97-dft} is nonzero.  Thus these
are genuine growing resonances, not a bound caused solely by a triangle
inequality.  If \(\max(u,v)\ge2f\), the V95 zero arc lies within
distance \(1/(2f)\) of zero and contains none of the nonzero conductor
centres.  All those centres lie in the nonzero completion channel.
The label ``nonzero completion'' therefore does not, by itself, mean
that its coherently assembled arithmetic multiplier is small.

\subsection{An exact finite-bank rank calculation}

Let \(u,v\) also be units modulo \(f\).  Let \(B\subseteq\{1,\ldots,Y\}\)
contain only units modulo \(f\), and let
\(N\subseteq\{1,\ldots,Y\}\) satisfy \(vn\equiv\omega\pmod f\).
These sets may include the exact low-reader cutoffs of V96.  Define
operators from \(\ell^2(N)\) to \(\ell^2(B)\) by their matrices
\begin{align}
 K^{\rm all}_{bn}&=\sum_{\ell\in J}\rho_\omega(\ell)
                       \mathbf1_{ub-vn=\ell},
 \label{eq:f2v97-Kall}\\
 K^0_{bn}&=\sum_{\ell\in J}\rho_\omega(\ell)
                       \mathscr Z_{u,v}(ub-vn-\ell),\qquad
 K^{\ne0}=K^{\rm all}-K^0.
 \label{eq:f2v97-Kparts}
\end{align}
The last definition is exactly the aggregate nonzero completion, by
Theorem~\ref{thm:f2v97-tiling}.

\begin{theorem}[The remaining operator retains a large rank-one component]
\label{thm:f2v97-rank-one}
Put \(z_b=\chi(\omega)\overline{\chi(ub)}\), and
\(E_{bn}=\mathbf1_{ub=vn}\).  Then
\begin{equation}
 K^{\rm all}=z\mathbf1_N^*-E,\qquad \|E\|\le1,
 \qquad \|K^{\rm all}\|\ge\sqrt{|B||N|}-1.
 \label{eq:f2v97-rank-one}
\end{equation}
If \(\max(u,v)\ge2f\), then
\begin{equation}
 \|K^{\ne0}\|\ge\sqrt{|B||N|}-1
                    -\frac{2(1+\pi)f}{\sqrt{uv}}.
 \label{eq:f2v97-nonzero-lower}
\end{equation}
\end{theorem}

\begin{proof}
Every nonzero \(ub-vn\) for the admitted indices lies in \(J\).
On this exact displacement, the conductor identity gives
\(\rho_\omega(ub-vn)=z_b\), independently of \(n\in N\).
When \(ub=vn\), \(z_b=\rho_\omega(0)=1\), but the envelope
excludes zero.  This proves the matrix identity.  Since \((u,v)=1\),
the equality entries have the form \((b,n)=(vq,uq)\).  Each row and
column has at most one such entry; hence \(E\) is a restricted partial
permutation and has norm at most one.  All entries of \(z\) have
modulus one, giving rank-one norm \(\sqrt{|B||N|}\).
The reverse triangle inequality proves the first lower bound.
For the second, the two intervals of \(J\) and V95 give
\(\|K^0\|\le2(1+\pi)f/\sqrt{uv}\), also after compression to
\(B,N\).  Apply the reverse triangle inequality once more.
\end{proof}

For banks of positive conductor-dependent density, the lower bound is
of order \(Y\), rather than \((uv)^{-1/2}\).  This is an obstruction
to a uniform small operator norm, not a lower bound for the actual
M\"obius pairing.  The latter retains exactly the factors
\(\mu(b)\overline{W_K(gub)}\) and \(\mu(n)W_K(gvn)\), together
with their cutoffs and cross-triple assembly.  Their pairing with the
rank-one component is a product of their signed row sums; it is not
determined by their norms.  This is the arithmetic information the next
step must use.

\subsection*{Verification and direction}

\path{scripts/verify_grh_v97_frequency_tiling.py} checks the CRT tiling,
the phase congruence including the reciprocal factor, and the zero-cell
union using exact rational arithmetic.  It checks the finite-bank
rank identity with the low-reader cutoff for real and complex primitive
characters of moduli three, four, and five, and verifies finite Fourier
inversion in a cyclotomic polynomial quotient.  These finite checks
support the identities; the uniform statements are proved above.

This rules out applying the V96 norm saving unchanged to its complement.
The next move is an exact convolution analysis of the actual row sums,
starting with the \(u=v=1\) bank already isolated in V96.  A bound for
only one growing-frequency sample or only a surrogate extremizer would
not answer that source question.  GRH and the residual target
\eqref{eq:f2v96-residual} remain open.

V97 appends to V96 with SHA--256
\begin{center}
\texttt{69C71DDF642226D4258F34B2A280575F}\allowbreak
\texttt{B1626A88E1D9E8DABC23A2775FEAB5F5}.
\end{center}
The scheduled companion sweep remains due at V100.
\endgroup

\clearpage
\section{V98: actual row completion and a closed terminal band of the unit bank}
\label{sec:f2v98-terminal-unit-bank}
\begingroup
\setcounter{equation}{0}
\renewcommand{\theequation}{N98.\arabic{equation}}
\renewcommand{\theHequation}{f2v98.\arabic{equation}}

The preceding frequency audit requires arithmetic row information, not
another uniform operator estimate.  Here the complete convolution
\(\mu*\mathbf1=\delta_1\) is applied to the actual row \(B_g\) of
V96.  It proves an unconditional estimate for a growing terminal band
of the unit reduced-pair bank.  Its contribution is part of the original
relocated zero sheet, so it can be removed from the residual target
without changing the decomposition or dropping any cross-bank term.
The other rows are not asserted to satisfy the same estimate.

\subsection{An exact row identity, including all floors}

Fix \(K\in\mathbb N\), \(4\le R\le\sqrt K\), and
\(Y=\lfloor K/(R+1)\rfloor\).  Continue to use the fixed primitive
nonprincipal \(\chi\) and \(C_\chi\) from V96.  For \(2\le g\le R\),
put
\begin{equation}
 X_g=\lfloor K/g\rfloor,\qquad A_g=\lfloor K/(gR)\rfloor,
 \qquad
 L_g=\sum_{n\le A_g}\mu(n)\chi(n)S_\chi(\lfloor X_g/n\rfloor).
 \label{eq:f2v98-XAL}
\end{equation}
The row previously denoted by \(B_g\) is exactly
\begin{equation}
 B_g=\sum_{A_g<n\le Y}\mu(n)\chi(n)
                                  S_\chi(\lfloor X_g/n\rfloor).
 \label{eq:f2v98-row}
\end{equation}
Indeed \(gn>K/R\) is equivalent to \(n>A_g\), and
\(\lfloor K/(gn)\rfloor=\lfloor X_g/n\rfloor\).
One has \(A_g\le Y\le X_g\), since \(gR\ge R+1\) and \(g\le R\).

\begin{theorem}[Completion of the actual unit-bank row]
\label{thm:f2v98-row-completion}
For every \(2\le g\le R\),
\begin{equation}
 B_g=1-L_g-
       \sum_{Y<n\le X_g}\mu(n)\chi(n)S_\chi(\lfloor X_g/n\rfloor).
 \label{eq:f2v98-row-completion}
\end{equation}
If in addition \(g\ge(R+1)/2\), this simplifies to
\begin{equation}
 B_g=1-L_g-\bigl(M_\chi(X_g)-M_\chi(Y)\bigr).
 \label{eq:f2v98-short-tail}
\end{equation}
Neither identity assumes cancellation in a Mertens sum.
\end{theorem}

\begin{proof}
For every integer \(X\ge1\), finite divisor grouping gives
\[
 \sum_{n\le X}\mu(n)\chi(n)S_\chi(\lfloor X/n\rfloor)
 =\sum_{r\le X}\chi(r)\sum_{n\mid r}\mu(n)=1.
\]
Partition the range at \(A_g\) and \(Y\) to prove the first identity.
For the second, \(K<(R+1)(Y+1)\) and \(2g\ge R+1\) imply
\(X_g<2(Y+1)\).  Thus \(Y<n\le X_g\) forces
\(\lfloor X_g/n\rfloor=1\), and \(S_\chi(1)=1\).
The tail is precisely the stated Mertens difference.  No extension
of a finite sum past its actual endpoint has been made.
\end{proof}

The constant term in this identity is supplied by complete divisor
cancellation.  Bounding the original row term by term would give a
bound of order \(Y\) and miss the terminal gain below.  For rows
away from the endpoint, however, the Mertens difference in
\eqref{eq:f2v98-short-tail} is an unknown signed quantity, not a
closed estimate.

\subsection{A quantitative endpoint mean-square bound}

Write \(\Lambda=K/R^2\ge1\), and for an integer
\(1\le H\le\lfloor R/2\rfloor\) let
\[
 \mathcal G_H=\{R-H+1,\ldots,R\}.
\]
For \(g=R-j\in\mathcal G_H\), one has
\begin{equation}
 |L_g|\le2C_\chi\Lambda,\qquad
 0\le X_g-Y\le2\Lambda(j+1)+1,
 \qquad
 |B_g|\le2(C_\chi+2)\Lambda(j+1).
 \label{eq:f2v98-pointwise-edge}
\end{equation}
To verify the middle bound, use
\[
 X_g-Y\le\frac K g-\frac K{R+1}+1
 =\frac{K(R+1-g)}{g(R+1)}+1
 \le2\Lambda(j+1)+1.
\]
Also \(A_g\le K/(gR)\le2\Lambda\), proving the first bound.
Apply \eqref{eq:f2v98-short-tail} and the trivial inequality
\(|M_\chi(X_g)-M_\chi(Y)|\le X_g-Y\).  Since \(\Lambda\ge1\),
the resulting bound
\(2+2C_\chi\Lambda+2\Lambda(j+1)\) is at most the last expression
in \eqref{eq:f2v98-pointwise-edge}.

\begin{theorem}[An unconditional growing unit-bank sector]
\label{thm:f2v98-edge-sector}
Use the actual diagonal \(E_g\) from \eqref{eq:f2v96-Eg}, and set
\begin{equation}
 \mathcal U_H=\sum_{\substack{g\in\mathcal G_H\\(g,f)=1}}
                         \mu(g)^2\bigl(|B_g|^2-E_g\bigr).
 \label{eq:f2v98-unit-edge}
\end{equation}
Then
\begin{align}
 \sum_{g\in\mathcal G_H}|B_g|^2
     &\le4(C_\chi+2)^2\Lambda^2 H^3,
 \label{eq:f2v98-edge-ms}\\
 |\mathcal U_H|
     &\le4(C_\chi+2)^2\Lambda^2 H^3+C_\chi^2HY.
 \label{eq:f2v98-edge-current}
\end{align}
These are estimates for the prescribed M\"obius rows, not arbitrary
coefficient vectors.
\end{theorem}

\begin{proof}
Square \eqref{eq:f2v98-pointwise-edge}, sum over \(j=0,\ldots,H-1\),
and use \(\sum_{j=1}^Hj^2\le H^3\).
The diagonal definition gives \(0\le E_g\le C_\chi^2Y\).
Now take the absolute value of the displayed finite sector, use
\(\mu(g)^2\le1\), and drop only the nonnegative coprimality and
squarefree restrictions in this upper bound.  This proves the second
estimate.
\end{proof}

Choose
\begin{equation}
 H_* =\left\lfloor\min\left\{\frac R2,
                       \left(\frac{RY}{\Lambda^2}\right)^{1/3}
                          \right\}\right\rfloor.
 \label{eq:f2v98-Hstar}
\end{equation}
If this is zero, take \(\mathcal U_{H_*}=0\).  Otherwise the theorem
proves
\begin{equation}
 |\mathcal U_{H_*}|\le
          \bigl(4(C_\chi+2)^2+C_\chi^2\bigr)RY=O_\chi(K).
 \label{eq:f2v98-closed-sector}
\end{equation}
At the original corner-matched scale of (R.6),
\(\Lambda\sim E_K^2\) and
\begin{equation}
 H_*\sim K^{1/3}E_K^{-4/3}=K^{1/3-o(1)},\qquad
                  H_*/R\longrightarrow0.
 \label{eq:f2v98-edge-scale}
\end{equation}
In particular this is a growing set of common factors, but not the
whole small-modulus bank.  In the Mertens-difference coordinate its
tail length is at most \(O(K^{1/3}E_K^{2/3})\); the proof used
only the trivial bound on that tail.

\subsection{Update of the remaining signed target}

The unit-pair bank has \(u=v=1\), hence no nonzero completion channel.
Therefore \(\mathcal U_{H_*}\) is a literal summand of
\(\widetilde{\mathcal O}^{(0,\mathrm{small})}_{\chi,K}\) in V96,
with its same-product diagonal already removed.  Define the residual
\begin{equation}
 \mathcal R_{\chi,K}^{\rm bulk}
 =\widetilde{\mathcal O}^{(0,\mathrm{small})}_{\chi,K}
      -\mathcal U_{H_*}
      +\widetilde{\mathcal O}^{(\ne0)}_{\chi,K}.
 \label{eq:f2v98-bulk}
\end{equation}
The exact current now reads
\begin{equation}
 \mathcal O_{\chi,K}
 =\widetilde{\mathcal O}^{(0,\mathrm{large})}_{\chi,K}
        +\mathcal U_{H_*}+\mathcal R_{\chi,K}^{\rm bulk}.
 \label{eq:f2v98-current}
\end{equation}
The first two terms have been bounded unconditionally, respectively
by \(O_\chi(K\log K)\) and \(O_\chi(K)\).
At corner-matched cutoffs, the V96 residual estimate is equivalent to
\begin{equation}
 |\mathcal R_{\chi,K}^{\rm bulk}|\le K^{1+o(1)},
 \label{eq:f2v98-open-bulk}
\end{equation}
by the triangle inequality in both directions and
\eqref{eq:f2v98-closed-sector}.  This remains an unproved condition.
The unit bank outside \(\mathcal G_{H_*}\), all other small reduced
pairs, and the entire nonzero completion contribution are still present
coherently in \eqref{eq:f2v98-bulk}.  No separate bounds for them have
been inferred from the terminal estimate.  In particular, the statement
does not assume a square-root Mertens bound to prove one.

\subsection*{Verification and continuation record}

\path{scripts/verify_grh_v98_terminal_rows.py} verifies the completed
convolution, floor identities, short-tail simplification, and unit-bank
sector identity in exact arithmetic, including the complex character
modulo five.  It checks the explicit inequalities using the harmless
integer majorant \(C_\chi\le f\), and selects \(H_*\) by an integer
inequality rather than floating-point cube roots.  The analytic proof
above, not finite testing, supplies uniformity in \(K\).

The remaining task is to obtain signed information on the actual bulk
rows or their cross-bank assembly.  Extending the terminal estimate to
all \(g\) without a new argument would replace a short tail by a
long unknown Mertens difference and is not justified.  The current
iteration closes a specified source sector; GRH remains unresolved.

V98 appends to V97 with SHA--256
\begin{center}
\texttt{AC84C3CC09605D5F1F379225226E7AFC}\allowbreak
\texttt{29E6BCD8C8FFA298A5F2252F6D2B4AA3}.
\end{center}
The next companion-note sweep is V100, before the next volume rollover.
\endgroup

\clearpage
\section{V99: sparse squarefree row samples detect anchored Mertens energy}
\label{sec:f2v99-sparse-sampling}
\begingroup
\setcounter{equation}{0}
\renewcommand{\theequation}{N99.\arabic{equation}}
\renewcommand{\theHequation}{f2v99.\arabic{equation}}

V98 controls a terminal band by completing the actual rows.  Extending
that estimate to a positive proportion of the common factors is not
a routine enlargement of the same argument: it exposes long Mertens
increments.  The present iteration proves the sampling statement
needed to assess the strength of that proposed input.  In particular,
it does not silently remove the factor \(\mu(g)^2\mathbf1_{(g,f)=1}\)
or replace sparse reciprocal endpoints by a continuous average.

\subsection{Uniform local density of the admitted common factors}

For a fixed integer \(f\ge1\), put
\[
 \mathcal S_f=\{g\ge1:\mu(g)^2=1,\ (g,f)=1\},\qquad
 \delta_f=\frac{\varphi(f)}f\prod_{p\nmid f}(1-p^{-2})>0.
\]

\begin{lemma}[A local squarefree density estimate]
\label{lem:f2v99-density}
Uniformly for an interval \(J\subseteq[1,R]\),
\begin{equation}
 |J\cap\mathcal S_f|=\delta_f|J|+O_f(\sqrt R+1).
 \label{eq:f2v99-density}
\end{equation}
Here \(|J|\) denotes its real length; changing endpoint conventions
changes the count by at most two.  Consequently, for a sufficiently
large fixed \(c_f\), intervals of length at least
\(c_f\sqrt R\) contained in \([1,R]\) contain at least a fixed
positive \(f\)-dependent multiple of their length in \(\mathcal S_f\).
\end{lemma}

\begin{proof}
Use \(\mu(g)^2=\sum_{d^2\mid g}\mu(d)\).  If \((g,f)=1\),
only \((d,f)=1\) can contribute.  For \(d\le\sqrt R\),
inclusion--exclusion over the prime divisors of \(f\) gives
\[
 \#\{k:d^2k\in J,\ (k,f)=1\}
       =\frac{\varphi(f)}f\frac{|J|}{d^2}+O_f(1).
\]
Sum with coefficient \(\mu(d)\).  The total counting error is
\(O_f(\sqrt R+1)\); extending the main sum to all \(d\) costs
\(O(|J|/\sqrt R)=O(\sqrt R)\).  The absolutely convergent Euler
product is the displayed \(\delta_f\), and is positive because
\(\sum_p p^{-2}<\infty\).  Choosing the interval length large
enough relative to the error proves the last assertion.
\end{proof}

\subsection{A deterministic transfer from reciprocal samples}

Let \(a_n\) be any complex sequence with \(|a_n|\le1\), and put
\(M(x)=\sum_{n\le x}a_n\).  For integers \(Y\ge R\) and
\(R\) sufficiently large depending on \(f\), define
\begin{equation}
 K=(R+1)Y,\quad X_g=\lfloor K/g\rfloor,\quad
 \mathcal G=\mathcal S_f\cap
       \{\lceil(R+1)/2\rceil,\ldots,R\},
 \quad
 P_a(Y,R)=\sum_{g\in\mathcal G}|M(X_g)-M(Y)|^2.
 \label{eq:f2v99-samples}
\end{equation}
The anchoring at \(M(Y)\) is retained.

\begin{theorem}[Sparse squarefree sampling transfer]
\label{thm:f2v99-sampling}
For the above deterministic sequence and grid,
\begin{equation}
 \sum_{x=Y}^{2Y-1}|M(x)-M(Y)|^2
 \ll_f \frac YR P_a(Y,R)+\frac{Y^3}{R}.
 \label{eq:f2v99-transfer}
\end{equation}
The constant is independent of the sequence, \(Y\), and \(R\).
\end{theorem}

\begin{proof}
Choose an integer \(H\asymp_f\sqrt R\) large enough for
Lemma~\ref{lem:f2v99-density}; for large \(R\), also \(H\le R/8\).
For each integer \(Y\le x<2Y\), let
\[
 \mathcal G_x=\{g\in\mathcal G:|g-K/x|\le2H\}.
\]
The interval \([K/x-2H,K/x+2H]\), clipped to the common-factor
range, has length at least \(H\) for large \(R\).  To check the
endpoints, \(K/x\) lies between \((R+1)/2\) and \(R+1\);
at most one unit of the centre lies beyond the upper endpoint.
The density lemma therefore gives \(|\mathcal G_x|\gg_f H\).

For \(g\in\mathcal G_x\),
\[
 |X_g-x|\le\frac{x|K/x-g|}{g}+1
             \le\frac{8YH}{R}+1=:D.
\]
Both \(X_g\) and \(x\) are integers, so bounded increments give
\(|M(X_g)-M(x)|\le D\).  Average the squared triangle inequality:
\begin{equation}
 |M(x)-M(Y)|^2
 \le\frac2{|\mathcal G_x|}\sum_{g\in\mathcal G_x}
                |M(X_g)-M(Y)|^2+2D^2.
 \label{eq:f2v99-local-average}
\end{equation}
For fixed \(g\), the admitted \(x\)'s lie in
\([K/(g+2H),K/(g-2H)]\).  Its length is
\(4KH/(g^2-4H^2)\ll YH/R\), since \(g\ge R/2\) and
\(H\le R/8\).  The number of integer \(x\)'s is also
\(O(YH/R)\), as \(Y\ge R\).  Summing
\eqref{eq:f2v99-local-average} over \(x\) thus gives
\[
 \sum_{x=Y}^{2Y-1}|M(x)-M(Y)|^2
 \ll_f (Y/R)P_a(Y,R)+Y(1+YH/R)^2.
\]
Finally \(H^2\ll_f R\) and \(Y\ge R\), so the error is
\(O_f(Y^3/R)\), as asserted.
\end{proof}

The \(Y^3/R\) term records the cost of filling gaps between admitted
squarefree common factors.  It is not zero.  For
\(R\asymp Y/(\log Y)^2\), it is only \(Y^{2+o(1)}\), so at
that scale the sparse weighting cannot hide a larger power of anchored
energy.

\subsection{Anchored energy controls the anchor itself}

\begin{lemma}[Overlapping anchors]
\label{lem:f2v99-anchors}
For a sequence with \(|a_n|\le1\), suppose that, for every
\(\varepsilon>0\),
\begin{equation}
 I(Y):=\sum_{x=Y}^{2Y-1}|M(x)-M(Y)|^2
                  \ll_\varepsilon Y^{2+\varepsilon}
 \label{eq:f2v99-anchor-energy}
\end{equation}
for all sufficiently large integers \(Y\).  Then, for every
\(\varepsilon>0\),
\begin{equation}
 |M(X)|\ll_\varepsilon X^{1/2+\varepsilon}.
 \label{eq:f2v99-M-bound}
\end{equation}
Conversely, \eqref{eq:f2v99-M-bound} implies
\eqref{eq:f2v99-anchor-energy}, after relabelling \(\varepsilon\).
\end{lemma}

\begin{proof}
Set \(Z=\lfloor3Y/2\rfloor\).  The intervals of integers
\([Y,2Y)\) and \([Z,2Z)\) overlap in \([Z,2Y)\), which
contains at least \(Y/2\) integers.  For each such \(x\),
\[
 |M(Z)-M(Y)|^2
 \le2|M(x)-M(Y)|^2+2|M(x)-M(Z)|^2.
\]
Summing over the overlap and using the two assumed energy bounds
gives \(|M(Z)-M(Y)|\ll_\varepsilon Y^{1/2+\varepsilon}\).

To bound an arbitrary integer \(N\), descend through
\(N_{j+1}=\lfloor2N_j/3\rfloor\) until a bounded value is reached.
For \(Y=N_{j+1}\), the associated \(\lfloor3Y/2\rfloor\) differs
from \(N_j\) by at most two, costing at most two bounded increments.
The above estimate therefore bounds
\(|M(N_j)-M(N_{j+1})|\) by
\(O_\varepsilon(N_j^{1/2+\varepsilon}+1)\).
The geometric sum is \(O_\varepsilon(N^{1/2+\varepsilon})\).
Real upper limits follow by taking floors.  Conversely, bound each
anchored difference by the two pointwise bounds and sum \(Y\) terms.
\end{proof}

This overlap argument is essential: a single anchored average does
not determine \(M(Y)\).  Bounds at every large anchor, however, do
determine its growth.  No independence, random model, or sign condition
on \(a_n\) has been used.

\subsection*{Verification and use in the next iteration}

\path{scripts/verify_grh_v99_sparse_sampling.py} checks exact finite
neighbourhood multiplicities, the squared local averaging inequality,
and the overlapping-anchor identities on real and complex character
M\"obius sequences.  The asymptotic density and gap estimates are
proved above, not inferred from these finite examples.

The next iteration applies these lemmas to the completed rows of V98
at an explicit corner-matched family.  It must account for the low
prefix \(L_g\) and the entire unit-bank diagonal before identifying
the strength of a proposed positive-bank bound.  GRH remains open.

V99 appends to V98 with SHA--256
\begin{center}
\texttt{509EF5E3DA420068BEE54748A01CE10F}\allowbreak
\texttt{F0D982BB946490F743D6CB4C33151BB8}.
\end{center}
The next iteration includes the scheduled companion sweep.
\endgroup

\clearpage
\section{V100: the positive unit bank is already a Mertens criterion; companion sweep}
\label{sec:f2v100-unit-criterion}
\begingroup
\setcounter{equation}{0}
\renewcommand{\theequation}{N100.\arabic{equation}}
\renewcommand{\theHequation}{f2v100.\arabic{equation}}

The V99 sparse-sampling theorem now applies to the actual completed
rows of V98.  It proves that controlling even the upper-half unit
bank separately would already give the square-root Mertens estimate.
This does not rule out proving that estimate directly.  It does rule
out treating it as a previously established or substantially weaker
input to the remaining signed-current argument.  The distinction
changes the direction of the next volume: retain cross-bank compensation
instead of extending the terminal estimate by an unproved mean-square
assumption.

\subsection{An explicit corner-matched family}

For every sufficiently large integer \(Y\), set
\begin{equation}
 R=\left\lfloor\frac{Y}{(\log Y)^2}\right\rfloor,
 \qquad K=(R+1)Y,\qquad
 \mathcal G=\mathcal S_f\cap
                  \{\lceil(R+1)/2\rceil,\ldots,R\}.
 \label{eq:f2v100-family}
\end{equation}
Then \(Y=\lfloor K/(R+1)\rfloor\), \(RY<K<(R+1)(Y+1)\),
and \(R\le\sqrt K\).  It is a legitimate corner-matched family:
with \(E_Y=\sqrt K/(R+1/2)\), one has
\(\lfloor\sqrt K/E_Y\rfloor=R\), \(E_Y\sim\log Y\to\infty\),
and \(E_Y=K^{o(1)}\).  Every large integer anchor \(Y\) occurs;
no density-of-anchors assertion for an arbitrary cutoff rule is used.

For a fixed primitive nonprincipal \(\chi\) modulo \(f\), define
\begin{equation}
 \mathcal P_\chi(Y)=\sum_{g\in\mathcal G}|B_g|^2,\qquad
 \mathcal U_\chi(Y)=\sum_{g\in\mathcal G}(|B_g|^2-E_g),
 \label{eq:f2v100-PU}
\end{equation}
where \(B_g,E_g\) are the exact V96 rows and same-entry diagonal
at the cutoffs \eqref{eq:f2v100-family}.  Squarefreeness and conductor
coprimality are built into \(\mathcal G\), not removed.
The second quantity is a literal subbank of the relocated zero sheet.

\begin{theorem}[Strength of the isolated upper-half unit bank]
\label{thm:f2v100-unit-equivalence}
For this fixed character, the following three statements are equivalent:
\begin{enumerate}[label=\textup{(\roman*)}]
\item For every \(\varepsilon>0\),
\(\mathcal P_\chi(Y)\ll_{\chi,\varepsilon}RY^{1+\varepsilon}\)
for all sufficiently large integers \(Y\).
\item For every \(\varepsilon>0\),
\(|\mathcal U_\chi(Y)|\ll_{\chi,\varepsilon}RY^{1+\varepsilon}\)
for all sufficiently large integers \(Y\).
\item For every \(\varepsilon>0\),
\(|M_\chi(X)|\ll_{\chi,\varepsilon}X^{1/2+\varepsilon}\).
\end{enumerate}
Equivalently, either bank bound may be written as
\(K^{1+o(1)}\) on this family.  These statements are not proved
unconditionally here; their equivalence is the theorem.
\end{theorem}

\begin{proof}
Since \(0\le E_g\le C_\chi^2Y\),
\[
 0\le\mathcal P_\chi(Y)-\mathcal U_\chi(Y)
                         \le C_\chi^2RY.
\]
This proves the equivalence of (i) and (ii), including the absolute
value in (ii).

For every \(g\in\mathcal G\), V98 gives
\[
 M_\chi(X_g)-M_\chi(Y)=1-L_g-B_g,\qquad
 |L_g|\le C_\chi A_g\le 2C_\chi Y/R.
\]
The last bound uses \(K=(R+1)Y\) and \(2g\ge R+1\).
Thus the sampled energy of V99 satisfies
\begin{equation}
 \sum_{g\in\mathcal G}|M_\chi(X_g)-M_\chi(Y)|^2
             \le2\mathcal P_\chi(Y)+O_\chi(Y^2/R).
 \label{eq:f2v100-row-transfer}
\end{equation}
Apply Theorem~\ref{thm:f2v99-sampling} to
\(a_n=\mu(n)\chi(n)\).  Since \(R\asymp Y/(\log Y)^2\),
\begin{equation}
 \sum_{x=Y}^{2Y-1}|M_\chi(x)-M_\chi(Y)|^2
 \ll_\chi (Y/R)\mathcal P_\chi(Y)+Y^3/R.
 \label{eq:f2v100-anchored-transfer}
\end{equation}
Under (i), this is \(Y^{2+o(1)}\).  The overlapping-anchor
Lemma~\ref{lem:f2v99-anchors} proves (iii).  In particular no unknown
constant anchor value has been discarded.

Conversely suppose (iii).  For \(g\in\mathcal G\),
\(Y\le X_g\le2Y\).  The exact row identity gives, for any
fixed \(\eta>0\),
\[
 |B_g|\ll_{\chi,\eta}1+Y/R+Y^{1/2+\eta}.
\]
Square and sum at most \(R\) terms.  The low-prefix contribution
is \(O_\chi(Y^2/R)\), while the Mertens contribution is
\(O_{\chi,\eta}(RY^{1+2\eta})\).  Choosing \(\eta\) smaller
than half the desired exponent proves (i).  Finally
\(K\asymp RY=Y^{2+o(1)}\), so the two exponent conventions agree.
\end{proof}

\begin{corollary}[The terminal removal does not lower the bulk criterion]
\label{cor:f2v100-edge-removal}
Let \(\mathcal U_{H_*}\) be the already bounded terminal sector
from V98, at the family \eqref{eq:f2v100-family}.  Replacing
\(\mathcal U_\chi(Y)\) in assertion (ii) by
\(\mathcal U_\chi(Y)-\mathcal U_{H_*}\) gives the same equivalent
criterion.
\end{corollary}
\begin{proof}
The terminal common factors lie in the upper half, and
\(|\mathcal U_{H_*}|\ll_\chi RY\) is already proved.
Use the triangle inequality in both directions.  This does not
enlarge the terminal estimate to the other common factors.
\end{proof}

For completeness, the implication from (iii) to the Dirichlet GRH
channel uses the classical analytic continuation and functional equation
already part of the supplied arithmetic loading.  Partial summation
makes
\[
 F(s)=s\int_1^\infty M_\chi(x)x^{-s-1}\,dx
\]
holomorphic for \(\operatorname{Re}s>1/2\): on each compact set choose
the exponent in (iii) smaller than its distance from that boundary.
For \(\operatorname{Re}s>1\), the Euler product gives
\(F(s)=1/L(s,\chi)\).  The identity theorem then gives
\(L(s,\chi)F(s)=1\) throughout the right half-plane, so there are
no zeros there.  The functional equation together with conjugation
reflects a nontrivial zero \(\rho\) of this character to
\(1-\overline\rho\) of the same character, excluding zeros to the
left as well.  This explains the term ``GRH-strength'' here without
assuming the bank estimate itself or importing a converse analytic
criterion into the new proof.

The signed residual \eqref{eq:f2v98-open-bulk} does \emph{not} assert
a separate bound for \(\mathcal U_\chi\).  Its other small-pair and
nonzero-frequency terms can compensate that subbank.  The present
equivalence is therefore not a new obstruction to GRH and not a proof
that the full signed route fails.  It identifies an overly strong
intermediate isolation which would simply relocate the main problem.

\subsection{Scheduled companion sweep and the newer V075 manuscript}

The sweep compared the live lineages with the V90 checkpoint.  These
are inspected snapshots, not a claim that other writers have stopped.
Headings and scope statements were reviewed across the new sections;
the proofs most relevant to coherent correlations, conditional measures,
and moving-reference estimates were read in full.
\begin{itemize}
\item The workspace Grand Tensor is now V075, 3478023 bytes, SHA--256
\begin{center}
\texttt{9A374B408818ECB37E95798A5199A3CB}\allowbreak
\texttt{86CC2267F132DD740E4C67E5C5611888}.
\end{center}
The attached V074 remains unchanged.  A direct comparison locates
the V075 additions in the spacetime/reference-relative blocking branch,
not a new GRH arithmetic estimate.
\item SM advanced from V50 to V72, 607372 bytes, SHA--256
\begin{center}
\texttt{F44F9745D43379E1672C5550411B58E9}\allowbreak
\texttt{7195A4C9278592D221DF5D3F644DC1F5}.
\end{center}
\item Main YM advanced from V48 through V81 during the sweep.
The inspected V81 has 911339 bytes and SHA--256
\begin{center}
\texttt{3E06CC81F5E729029FD5C67AB1247F22}\allowbreak
\texttt{6D99B2F3264B976AEC197F574AA73BC8}.
\end{center}
\item NS V26, 278283 bytes, remains at SHA--256
\begin{center}
\texttt{82C887F8C2C69E4F28FAD7C3DC8DE5B9}\allowbreak
\texttt{099CB5A4BF431EBA1FFF3E91935978AD}.
\end{center}
Parallel YM V5, 54174 bytes, remains at SHA--256
\begin{center}
\texttt{306F1BB84CFA8A8D47EA569EC17E9545F}\allowbreak
\texttt{9867C8D32B548EC9D9C444B4F3C0512}.
\end{center}
\end{itemize}

V075's \texttt{sec:GT-blocking-matching} distinguishes conservative
elimination, normalized reference-relative integration, and a driven
law.  They do not preserve the same matching defect.  Its
\texttt{sec:GT-coupled-scalar-energy} is more directly relevant as
a research direction: for specified scalar--shear interactions it
retains reciprocal outputs and constructs a positive coupled energy
where diagonal positive weights cannot satisfy the required identity.
This is a result for its declared action and gauge, not a M\"obius
correlation theorem.  No such energy can be imported into GRH without
deriving the corresponding full arithmetic identity and its source
norms.  It nevertheless supports keeping cross-bank terms in the
next volume instead of demanding positivity estimates block by block.

SM V56 makes a coherent weighted cross-pair term explicit even when
individual covariance entries tend to zero.  V68 shows how compression
loses intermediate momentum transfers.  V71 identifies the mixed
correlation residual omitted by a consistent covariance approximation;
V72 actually derives its vanishing in a specified many-copy scaling,
for concentrated initial data and fixed physical mode sets and times.
No copy limit, concentration hypothesis, or independent arithmetic
state has been supplied for the present Dirichlet source.  Thus its
positive-state transport theorem is not an estimate for our residual.

YM V49--V68 repairs representative and nonlinear-fibre identities,
constructs the actual local density and compact extension, and separates
local curvature from control of the full measure.  V69 exhibits a
positive-measure obstruction to a proposed conditional repair; a
relative action gain on the specified fibre is required.  V78--V81
distinguish an exact section from a constrained minimum, construct a
volume-uniform local moving minimum, and prove spatial decay using a
gapped finite-range fibre inverse.  The retained Schur complement still
has gauge modes.  The GRH current has not been given either that
finite-range model or that uniform gap.  The lesson is to prove a
bound on the actual conditioned and retained object, not to transfer
an eliminated-coordinate estimate by analogy.  None of these companion
results closes a Dirichlet zero channel.

\subsection*{Verification, ten-version review, and rollover}

\path{scripts/verify_grh_v100_unit_criterion.py} checks the explicit
cutoff family, exact row-to-Mertens identities, signed/positive unit-bank
relation, low-prefix energy inequality, and subtraction of the V98
terminal sector for real and complex characters.  The asymptotic
implication uses the two uniform proofs in V99, not finite tests.

V91--V94 resolved the behavior of the growing pinned-filter family
and its absolute zero remainder.  V95--V97 moved upstream to the exact
determinant current, bounded a relocated large-modulus zero sheet, and
located its remaining conductor resonances.  V98 closed a genuine
terminal unit-bank sector.  V99--V100 now show why a separate
positive-bank extension is already Mertens-strength.  The next volume
must seek cancellation in the full prescribed signed current or a
new source-derived coupled identity.  Repeating the sampled mean-square
criterion as an assumed lemma would be circular.

This V100 is frozen as
\path{Renewal_GRH_Cofinal_Visibility_Attack_V100_f3.tex}; the previous
\texttt{f2} archive is unchanged.  The new active cycle restarts at
\texttt{V1}, with a short source and dependency handoff rather than
duplicating this volume.  The next scheduled companion sweep is the
new cycle's V10.  No companion document is edited or retired.

V100 appends to V99 with SHA--256
\begin{center}
\texttt{E1DC82685D7B0B60B68D51E202C1AE20}\allowbreak
\texttt{7A2F903358C63CE5EA15569207701CCC}.
\end{center}
The full GRH objective remains open.
\endgroup

\end{document}
